%=====================================================================
% ĐỀ THI THỬ TỐT NGHIỆP THPT 2026 -- MÔN VẬT LÍ
% Nguồn: THPT Nguyễn Thị Minh Khai, Hà Nội -- Lần 1 (PDF SieuGioi, không đáp án). Gõ lại theo template Eduzy.
%=====================================================================
\def\CheDo{solcolor}% loigiai | dethi
\def\tenchuong{THPT Nguyễn Thị Minh Khai, Hà Nội -- Lần 1}
\def\thuonghieu{VẬT LÍ 12 -- Eduzy.com.vn}
\def\TachCauDS{1}
\def\TachCauLG{1}
%============================================================
%  PREAMBLE EDUZY v2 - dùng chung cho đề thi, bài học, chuyên đề (gói ex_test)
%  Khác bản preamble-eduzy.tex ở đầu trang, chân trang: dải nền nâu đồng tràn ngang,
%  logo + Eduzy.com.vn, ô chữ vàng kem cạnh xiên, chân trang "Trang (số)";
%  logo Eduzy chìm giữa mọi trang (tắt: \logochimtat, bật lại: \logochimbat).
%  Cách dùng ở file chính (xem mau-chuan.tex, mau-chuyen-de.tex):
%      \def\CheDo{loigiai}          % loigiai | dethi | solcolor | color | book
%      \def\tenchuong{<dòng chữ trên đầu trang>}
%      \def\thuonghieu{<chữ ở chân trang lẻ>}   % không khai báo thì dùng mặc định
%      \def\tacgia{<chữ ở chân trang chẵn>}      % không khai báo thì dùng mặc định
%      \def\khauhieu{<chữ cạnh logo ở đầu trang>}  % tùy chọn, mặc định Eduzy.com.vn
%      %============================================================
%  PREAMBLE EDUZY v2 - dùng chung cho đề thi, bài học, chuyên đề (gói ex_test)
%  Khác bản preamble-eduzy.tex ở đầu trang, chân trang: dải nền nâu đồng tràn ngang,
%  logo + Eduzy.com.vn, ô chữ vàng kem cạnh xiên, chân trang "Trang (số)";
%  logo Eduzy chìm giữa mọi trang (tắt: \logochimtat, bật lại: \logochimbat).
%  Cách dùng ở file chính (xem mau-chuan.tex, mau-chuyen-de.tex):
%      \def\CheDo{loigiai}          % loigiai | dethi | solcolor | color | book
%      \def\tenchuong{<dòng chữ trên đầu trang>}
%      \def\thuonghieu{<chữ ở chân trang lẻ>}   % không khai báo thì dùng mặc định
%      \def\tacgia{<chữ ở chân trang chẵn>}      % không khai báo thì dùng mặc định
%      \def\khauhieu{<chữ cạnh logo ở đầu trang>}  % tùy chọn, mặc định Eduzy.com.vn
%      %============================================================
%  PREAMBLE EDUZY v2 - dùng chung cho đề thi, bài học, chuyên đề (gói ex_test)
%  Khác bản preamble-eduzy.tex ở đầu trang, chân trang: dải nền nâu đồng tràn ngang,
%  logo + Eduzy.com.vn, ô chữ vàng kem cạnh xiên, chân trang "Trang (số)";
%  logo Eduzy chìm giữa mọi trang (tắt: \logochimtat, bật lại: \logochimbat).
%  Cách dùng ở file chính (xem mau-chuan.tex, mau-chuyen-de.tex):
%      \def\CheDo{loigiai}          % loigiai | dethi | solcolor | color | book
%      \def\tenchuong{<dòng chữ trên đầu trang>}
%      \def\thuonghieu{<chữ ở chân trang lẻ>}   % không khai báo thì dùng mặc định
%      \def\tacgia{<chữ ở chân trang chẵn>}      % không khai báo thì dùng mặc định
%      \def\khauhieu{<chữ cạnh logo ở đầu trang>}  % tùy chọn, mặc định Eduzy.com.vn
%      \input{preamble-eduzy-v2}
%      \TLdapso                      % tài liệu: \TLfull | \TLdapso | \TLan
%      \begin{document} ... \end{document}
%  Logo Eduzy vẽ bằng TikZ ngay trong file (\eduzylogo), không cần file ảnh.
%  Biên dịch: pdfLaTeX (máy có gói vietnam) hoặc XeLaTeX.
%============================================================
\documentclass[12pt,a4paper,twoside]{article}%<-- twoside: bat dau trang/chan trang chan-le
\usepackage{iftex}
\usepackage{amsmath,amssymb,tikz,mathrsfs,tkz-tab}
\usepackage{graphicx}
\usepackage{pgfplots}% đồ thị hàm số
\pgfplotsset{compat=1.17}
%\usepackage{mathptmx}
\usepackage{mathpazo}
\ifXeTeX% XeLaTeX: font chữ có đủ dấu tiếng Việt, giữ nguyên font toán mathpazo
	\usepackage[no-math]{fontspec}
	\setmainfont{TeX Gyre Pagella}
	\newfontfamily\EZqag{TeX Gyre Adventor}[NFSSFamily=qag,Ligatures=TeX]
\fi
\usepackage[many]{tcolorbox}
\usetikzlibrary{math,through,calc,intersections,angles,quotes,shapes,shapes.geometric,arrows,patterns,matrix,chains,arrows.meta,decorations.shapes,decorations.fractals,decorations.markings,shadows,decorations.pathreplacing}
\usepackage[top=1.6cm, bottom=1.8cm, left=1.8cm, right=1.8cm]{geometry}% v2: top 1.6cm (bản cũ 1.5cm) để chữ không sát dải đầu trang
\providecommand{\CheDo}{loigiai}% loigiai | dethi | solcolor | color | book
\expandafter\usepackage\expandafter[\CheDo]{ex_test}
% \tenchuong: khai báo ở file chính (trước \input) hoặc \renewcommand sau \input
%============================================================
%  [GỘP TỪ eduzy-de.tex] EDUZY - MÀU NÂU ĐỒNG / VÀNG KIM
%  + ĐẦU TRANG, CHÂN TRANG KIỂU DẢI NGANG (v2), ĐỐI XỨNG TRANG CHẴN / TRANG LẺ.
%  (không cần file eduzy-de.tex riêng nữa)
%  Khối này phải đặt SAU \usepackage[...]{ex_test} và SAU \def\tenchuong{...}
%============================================================
\usepackage{lastpage}
\usepackage{needspace}
%--- DÀN TRANG ĐỀ THI: tách câu (THÊM 05/11/2026; chỉ thêm lệnh mới, KHÔNG đổi hành vi mặc định) ---
%  Hai công tắc, mặc định tắt (= preamble chạy y như cũ). Bật bằng cách đặt TRƯỚC \input{preamble-eduzy-v2}
%  trong file đề (tools/dan_trang/dan_trang.py tự ghi khi đặt \NeedspaceCau):
%    \def\TachCauDS{1}   câu Đúng/Sai được tách: câu dẫn (kèm hình) ở cuối trang, bảng Phát biểu sang trang sau
%                        (câu dẫn và hình không tách nhau, bảng không tách giữa hai trang)
%    \def\TachCauLG{1}   bản lời giải (\CheDo khác dethi): phần đề của câu cao hơn 1/3 trang được phép tách
%  Lệnh dùng trong file đề (giá trị do tools/dan_trang/dan_trang.py đo ở bản dethi):
%    \NeedspaceCau[<cao câu dẫn + hình>]{<cao cả khối câu>}   ngay trên \begin{ex} (câu đầu phần: trên tiêu đề phần)
%        công tắc tắt: đúng bằng \Needspace{<cao cả khối câu>} như cũ;
%        \TachCauDS=1 và có [..] (câu Đúng/Sai được chọn tách): chỉ giữ câu dẫn + hình;
%        \TachCauLG=1 ở bản lời giải: đòi tối đa 1/3 \textheight.
%    \NeedspaceBang{<cao bảng Phát biểu>}%   ngay trước \choiceTFt: \TachCauDS=1 thì giữ trọn bảng trên một trang;
%        công tắc tắt: không làm gì.
\providecommand{\TachCauDS}{0}
\providecommand{\TachCauLG}{0}
\newlength{\EZnsCau}
\newcommand{\NeedspaceCau}[2][]{%
	\setlength{\EZnsCau}{#2}%
	\ifnum\TachCauDS=1 \if\relax\detokenize{#1}\relax\else\setlength{\EZnsCau}{#1}\fi\fi
	\ifnum\TachCauLG=1 \ifdefstring{\kindOPTN}{dethi}{}{\ifdim\EZnsCau>0.3333\textheight\setlength{\EZnsCau}{0.3333\textheight}\fi}\fi
	\Needspace{\EZnsCau}}
\newcommand{\NeedspaceBang}[1]{\ifnum\TachCauDS=1 \Needspace{#1}\fi}

%--- Thang màu: đậm -> nhạt
\definecolor{gold1}{RGB}{126,96,0}    % nâu đồng đậm  #7E6000
\definecolor{gold2}{RGB}{192,144,3}   % vàng kim      #C09003
\definecolor{gold3}{RGB}{255,215,99}  % vàng nhạt     #FFD763
\definecolor{gold4}{RGB}{255,231,153} % vàng rất nhạt #FFE799
\definecolor{gold5}{RGB}{255,243,205} % vàng kem      #FFF3CD
%--- Ghi đè tông xanh/đỏ cũ
\definecolor{myblue}{RGB}{126,96,0}
\colorlet{blue}{gold1}
\colorlet{cyan}{gold2}
\colorlet{teal}{gold2}

%--- Tên chương mặc định (file đề nên tự khai báo lại)
\providecommand{\tenchuong}{}
%--- Chân trang
\providecommand{\thuonghieu}{TOÁN 12 -- Eduzy.com.vn}
\providecommand{\tacgia}{Tô Gia Bảo -- 0353703783}

%--- Chữ cạnh logo ở đầu trang
\providecommand{\khauhieu}{Eduzy.com.vn}

%====================== THÔNG SỐ ĐẦU TRANG, CHÂN TRANG (v2) ======================
%  Dải nền đậm tràn ngang trang; ô chữ vàng kem có cạnh xiên, kèm 2 vạch xiên chuyển màu.
%  Đầu trang lẻ : logo + \khauhieu (trái)   | ô chữ \tenchuong (phải)
%  Chân trang lẻ: ô chữ \thuonghieu (trái)  | Trang (số) (phải)
%  Trang chẵn soi gương trang lẻ; chân trang chẵn ghi \tacgia.
%--- Màu: chỉ dùng thang vàng kim Eduzy
\colorlet{ezNen}{gold1}% dải nền đậm
\colorlet{ezVach1}{gold3}% vạch xiên sát dải nền (đậm hơn)
\colorlet{ezVach2}{gold4}% vạch xiên sát ô chữ (nhạt hơn)
\colorlet{ezO}{gold5}% ô chữ vàng kem
%--- Kích thước
\newlength{\EZdai}\setlength{\EZdai}{1.1cm}% chiều cao dải nền
\newlength{\EZo}\setlength{\EZo}{0.68cm}% chiều cao ô chữ
\newlength{\EZxien}\setlength{\EZxien}{0.55cm}% độ lệch ngang của cạnh xiên
\newlength{\EZvach}\setlength{\EZvach}{0.75cm}% bề rộng mỗi vạch xiên
\newlength{\EZkhe}\setlength{\EZkhe}{0.12cm}% khe trắng giữa vạch và ô chữ
\newlength{\EZdem}\setlength{\EZdem}{0.4cm}% khoảng đệm từ chữ tới cạnh xiên
\newlength{\EZle}\setlength{\EZle}{1.8cm}% lề trái/phải của chữ (bằng lề thân bài)
\newlength{\EZrongmin}\setlength{\EZrongmin}{4.5cm}% ô chữ rộng tối thiểu
\newlength{\EZrong}% bề rộng chữ trong ô (tính tự động)
%--- Font
\newcommand{\EZfontO}{\fontsize{10.5pt}{12pt}\selectfont\bfseries}% chữ trong ô
\newcommand{\EZfontKH}{\fontfamily{qag}\fontsize{11pt}{12pt}\selectfont\bfseries}% \khauhieu cạnh logo
\newcommand{\EZfontTrang}{\fontsize{11.5pt}{12pt}\selectfont\bfseries}% chữ "Trang" và số trang
%--- Đo bề rộng chữ trong ô
\newcommand{\EZdorong}[1]{%
	\settowidth{\EZrong}{\EZfontO #1}%
	\ifdim\EZrong<\EZrongmin\setlength{\EZrong}{\EZrongmin}\fi}
%--- Logo + chữ trắng \khauhieu (dựng sẵn một lần khi vào \begin{document})
\newsavebox{\EZlogobox}
\AtBeginDocument{\sbox{\EZlogobox}{%
		\raisebox{-0.5\height}{\eduzylogo[1cm]}\hspace{0.15cm}%
		\raisebox{-0.5\height}{\color{white}\EZfontKH\khauhieu}}}
%--- Hai vạch xiên + khe trắng. #1: điểm gốc ở mép dưới, #2: hướng (1 sang phải, -1 sang trái),
%    #3: hướng nghiêng (1 kiểu "/", -1 kiểu "\")
\newcommand{\EZvachxien}[3]{%
	\fill[white] (#1) --++({(#3)*\EZxien},\EZo) --++({(#2)*\EZkhe},0) --++({-(#3)*\EZxien},-\EZo) --cycle;
	\fill[ezVach2] ([xshift=(#2)*\EZkhe]#1) --++({(#3)*\EZxien},\EZo) --++({(#2)*\EZvach},0) --++({-(#3)*\EZxien},-\EZo) --cycle;
	\fill[ezVach1] ([xshift=(#2)*(\EZkhe+\EZvach)]#1) --++({(#3)*\EZxien},\EZo) --++({(#2)*\EZvach},0) --++({-(#3)*\EZxien},-\EZo) --cycle;
	\draw[white,line width=0.6pt] ([xshift=(#2)*(\EZkhe+\EZvach)]#1) --++({(#3)*\EZxien},\EZo);
}
%--- Chữ trong ô (căn giữa theo chiều cao ô, không lệch vì dấu tiếng Việt)
\tikzset{ezchu/.style={inner sep=0pt,text=ezNen,text height=0.72em,text depth=0.22em}}

%====================== LOGO CHÌM GIỮA TRANG ======================
%  Vẽ cùng đầu trang nên nằm dưới chữ của thân bài, có trên mọi trang.
%  Tắt: \logochimtat (từ trang đó trở đi), bật lại: \logochimbat.
%  Đổi cỡ, độ đậm: khai báo trước \input, ví dụ \def\EZchimcao{9cm} \def\EZchimdo{0.1}
\providecommand{\EZchimcao}{10cm}% chiều cao logo chìm
\providecommand{\EZchimdo}{0.12}% độ đậm (0 = trong suốt, 1 = rõ hẳn)
\newif\ifEZchim \EZchimtrue
\newcommand{\logochimtat}{\EZchimfalse}
\newcommand{\logochimbat}{\EZchimtrue}
%  Dựng sẵn logo mờ một lần khi vào \begin{document}; transparency group làm mờ cả logo
%  một lượt nên không lộ vệt ở chỗ các dải màu chồng nhau.
\newsavebox{\EZchimtho}
\newsavebox{\EZchimbox}
\AtBeginDocument{%
	\sbox{\EZchimtho}{\eduzylogo[\EZchimcao]}%
	\sbox{\EZchimbox}{\begin{tikzpicture}
			\begin{scope}[transparency group,opacity=\EZchimdo]
				\node[inner sep=0pt] {\usebox{\EZchimtho}};
			\end{scope}
	\end{tikzpicture}}}
\newcommand{\EZvechim}{%
	\ifEZchim
		\node[inner sep=0pt] at (current page.center) {\usebox{\EZchimbox}};
	\fi}

%====================== ĐẦU TRANG - TRANG LẺ ======================
%  (logo bên trái, ô chữ \tenchuong bên phải)
\newcommand{\myfancyheadOdd}{%
	\EZdorong{\tenchuong}%
	\begin{tikzpicture}[remember picture,overlay]
		\coordinate (T) at (current page.north west);
		\coordinate (P) at (current page.north east);
		%--- logo chìm giữa trang
		\EZvechim
		%--- dải nền
		\fill[ezNen] (T) rectangle ([yshift=-\EZdai]P);
		%--- ô chữ: X là góc dưới bên trái, cạnh trái xiên kiểu "/"
		\coordinate (X) at ([xshift=-\EZle-\EZrong-\EZdem-\EZxien,yshift=-\EZdai]P);
		\fill[ezO] (X) --++(\EZxien,\EZo) -- ([yshift=-\EZdai+\EZo]P) -- ([yshift=-\EZdai]P) --cycle;
		\EZvachxien{X}{-1}{1}
		\node[ezchu,anchor=east] at ([xshift=-\EZle,yshift=-\EZdai+0.5*\EZo]P) {\EZfontO\tenchuong};
		%--- logo + tên
		\node[inner sep=0pt,anchor=west] at ([xshift=\EZle,yshift=-0.5*\EZdai]T) {\usebox{\EZlogobox}};
	\end{tikzpicture}%
}

%====================== ĐẦU TRANG - TRANG CHẴN =====================
%  (soi gương: ô chữ \tenchuong bên trái, logo bên phải)
\newcommand{\myfancyheadEven}{%
	\EZdorong{\tenchuong}%
	\begin{tikzpicture}[remember picture,overlay]
		\coordinate (T) at (current page.north west);
		\coordinate (P) at (current page.north east);
		%--- logo chìm giữa trang
		\EZvechim
		%--- dải nền
		\fill[ezNen] (T) rectangle ([yshift=-\EZdai]P);
		%--- ô chữ: X là góc dưới bên phải, cạnh phải xiên kiểu "\"
		\coordinate (X) at ([xshift=\EZle+\EZrong+\EZdem+\EZxien,yshift=-\EZdai]T);
		\fill[ezO] (X) --++(-\EZxien,\EZo) -- ([yshift=-\EZdai+\EZo]T) -- ([yshift=-\EZdai]T) --cycle;
		\EZvachxien{X}{1}{-1}
		\node[ezchu,anchor=west] at ([xshift=\EZle,yshift=-\EZdai+0.5*\EZo]T) {\EZfontO\tenchuong};
		%--- logo + tên
		\node[inner sep=0pt,anchor=east] at ([xshift=-\EZle,yshift=-0.5*\EZdai]P) {\usebox{\EZlogobox}};
	\end{tikzpicture}%
}

%====================== CHÂN TRANG - TRANG LẺ ======================
%  (ô chữ \thuonghieu bên trái, "Trang (số)" bên phải)
\newcommand{\myfancyfootOdd}{%
	\EZdorong{\thuonghieu}%
	\begin{tikzpicture}[remember picture,overlay]
		\coordinate (T) at (current page.south west);
		\coordinate (P) at (current page.south east);
		%--- dải nền
		\fill[ezNen] (T) rectangle ([yshift=\EZdai]P);
		%--- ô chữ sát mép trên dải: X là góc dưới bên phải, cạnh phải xiên kiểu "/"
		\coordinate (X) at ([xshift=\EZle+\EZrong+\EZdem,yshift=\EZdai-\EZo]T);
		\fill[ezO] ([yshift=\EZdai-\EZo]T) -- (X) --++(\EZxien,\EZo) -- ([yshift=\EZdai]T) --cycle;
		\EZvachxien{X}{1}{1}
		\node[ezchu,anchor=west] at ([xshift=\EZle,yshift=\EZdai-0.5*\EZo]T) {\EZfontO\thuonghieu};
		%--- số trang
		\node[circle,draw=white,line width=0.9pt,text=white,inner sep=0pt,minimum size=0.62cm,anchor=east]
		(so) at ([xshift=-\EZle,yshift=0.5*\EZdai]P) {\EZfontTrang\thepage};
		\node[inner sep=0pt,text=white,anchor=base east] (tr) at ([xshift=-0.15cm]so.base west) {\EZfontTrang Trang};
	\end{tikzpicture}%
}

%====================== CHÂN TRANG - TRANG CHẴN ====================
%  (soi gương: "Trang (số)" bên trái, ô chữ \tacgia bên phải)
\newcommand{\myfancyfootEven}{%
	\EZdorong{\tacgia}%
	\begin{tikzpicture}[remember picture,overlay]
		\coordinate (T) at (current page.south west);
		\coordinate (P) at (current page.south east);
		%--- dải nền
		\fill[ezNen] (T) rectangle ([yshift=\EZdai]P);
		%--- ô chữ sát mép trên dải: X là góc dưới bên trái, cạnh trái xiên kiểu "\"
		\coordinate (X) at ([xshift=-\EZle-\EZrong-\EZdem,yshift=\EZdai-\EZo]P);
		\fill[ezO] ([yshift=\EZdai-\EZo]P) -- (X) --++(-\EZxien,\EZo) -- ([yshift=\EZdai]P) --cycle;
		\EZvachxien{X}{-1}{-1}
		\node[ezchu,anchor=east] at ([xshift=-\EZle,yshift=\EZdai-0.5*\EZo]P) {\EZfontO\tacgia};
		%--- số trang
		%    (vòng tròn ẩn so0 nằm giữa dải để lấy đường chân chữ, giống trang lẻ)
		\node[circle,inner sep=0pt,minimum size=0.62cm] (so0) at ([yshift=0.5*\EZdai]T) {\EZfontTrang\phantom{\thepage}};
		\node[inner sep=0pt,text=white,anchor=base west] (tr) at ($(so0.base -| T)+(\EZle,0)$) {\EZfontTrang Trang};
		\node[circle,draw=white,line width=0.9pt,text=white,inner sep=0pt,minimum size=0.62cm,anchor=base west]
		(so) at ([xshift=0.15cm]tr.base east) {\EZfontTrang\thepage};
	\end{tikzpicture}%
}

%--- Giữ tương thích với các file cũ đang gọi \myfancyhead, \myfancyfoot
\newcommand{\myfancyhead}{\myfancyheadOdd}
\newcommand{\myfancyfoot}{\myfancyfootOdd}

%======================= GẮN VÀO TRANG =============================
\pagestyle{fancy}
\fancyhf{}
\fancyhead[LO]{\myfancyheadOdd}
\fancyhead[RE]{\myfancyheadEven}
\fancyfoot[LO]{\myfancyfootOdd}
\fancyfoot[RE]{\myfancyfootEven}
\renewcommand{\headrulewidth}{0pt}
\renewcommand{\footrulewidth}{0pt}
\fancypagestyle{plain}{%
	\fancyhf{}
	\fancyhead[LO]{\myfancyheadOdd}
	\fancyhead[RE]{\myfancyheadEven}
	\fancyfoot[LO]{\myfancyfootOdd}
	\fancyfoot[RE]{\myfancyfootEven}
	\renewcommand{\headrulewidth}{0pt}
	\renewcommand{\footrulewidth}{0pt}
}
%=================== HẾT PHẦN GỘP TỪ eduzy-de.tex ===================
\newcommand{\indam}[1]{
	{\fontsize{12pt}{9pt}\sffamily \selectfont\bfseries\color{gold1} #1}
}
\newcommand{\ind}[1]{
	{\fontsize{14pt}{9pt}\fontfamily{qag} \selectfont\bfseries\color{gold1} #1}
}
\renewcommand{\FalseEX}{\stepcounter{dapan}\textbf{\selectfont\color{gold1}\Alph{dapan}.}\ignorespaces}
\renewtheorem{ex}{\color{gold1}Câu}
\renewcommand{\baselinestretch}{1.1}
\newcommand{\hoac}[1]{
	\left[\begin{aligned}#1\end{aligned}\right.}
\newcommand{\heva}[1]{
	\left\{\begin{aligned}#1\end{aligned}\right.}
\usepackage{esvect}
\def\vec{\vv}
\def\vt{\vv}
\def\overrightarrow{\vv}
%====================== LOGO EDUZY VẼ BẰNG TIKZ ======================
%  Đường viền dò từ logo gốc; dải màu vàng kim đo trực tiếp trên ảnh gốc
%  (nghiêng 21 độ, tô bằng 80 dải hẹp nên hiển thị giống nhau trên mọi trình xem PDF).
%  Dùng: \eduzylogo          -> cao 2.4cm (mặc định)
%        \eduzylogo[1.5cm]   -> cao 1.5cm
\definecolor{ezText}{RGB}{245,200,46}
\newcommand{\eduzylogo}[1][2.4cm]{%
	\resizebox{!}{#1}{%
		\begin{tikzpicture}[x=0.01cm,y=0.01cm]
			\path[use as bounding box] (0,0) rectangle (224,-224);
			% mũ tốt nghiệp + vương miện
			\begin{scope}
				\clip
				(119.4,-176.1)..controls(118.9,-176.1)and(116.8,-176)..(114.6,-176)
				..controls(101.4,-176)and(85.8,-175.9)..(83.7,-175.7)
				..controls(81.8,-175.6)and(81,-175.6)..(80.2,-175.7)
				..controls(78.8,-175.9)and(77.7,-175.9)..(76.2,-175.7)
				..controls(74.9,-175.5)and(70.8,-175.5)..(64.3,-175.7)
				..controls(59.1,-175.9)and(58.3,-175.7)..(58.6,-173.7)
				..controls(58.6,-172.8)and(59.6,-169.6)..(60.3,-168)
				..controls(60.8,-166.9)and(61.4,-165.8)..(62.3,-165)
				..controls(63.1,-164.2)and(64.3,-163.4)..(65,-163.2)
				..controls(65.5,-163.1)and(72.3,-163)..(76.6,-163.1)
				..controls(77.9,-163.1)and(80.1,-163.1)..(81.4,-163)
				..controls(82.8,-162.9)and(86.3,-162.9)..(89.2,-162.9)
				..controls(93,-162.9)and(94.9,-162.8)..(96,-162.7)
				..controls(97.2,-162.5)and(101,-162.5)..(114,-162.4)
				..controls(123,-162.3)and(134.9,-162.2)..(140.2,-162.1)
				..controls(150.1,-162)and(152.2,-162.1)..(153.6,-162.4)
				..controls(154.4,-162.6)and(157.1,-163.8)..(158.1,-164.4)
				..controls(158.8,-164.8)and(159.7,-165.7)..(159.9,-166.3)
				..controls(160,-166.5)and(160.2,-167.1)..(160.2,-167.6)
				..controls(160.3,-168)and(160.6,-169.4)..(160.9,-170.6)
				..controls(161.2,-171.7)and(161.5,-173.1)..(161.6,-173.6)
				..controls(161.8,-174.4)and(161.8,-174.5)..(161.5,-174.9)
				..controls(161.2,-175.6)and(160.8,-175.9)..(159.9,-176)
				..controls(159.2,-176.1)and(154.1,-176.2)..(151.5,-176.2)
				..controls(148.1,-176.1)and(122.6,-176.1)..(121.6,-176.2)
				..controls(120.8,-176.2)and(119.8,-176.2)..(119.4,-176.1)--cycle
				(130.4,-160)..controls(129,-160)and(119.8,-160)..(110.1,-159.9)
				..controls(100.3,-159.8)and(88.5,-159.7)..(83.9,-159.6)
				..controls(77.8,-159.5)and(75.3,-159.5)..(74.5,-159.6)
				..controls(73.9,-159.7)and(72.5,-159.8)..(71.3,-159.8)
				..controls(69.6,-159.8)and(69,-159.8)..(68.5,-159.6)
				..controls(67.6,-159.3)and(66.8,-158.6)..(66.2,-157.6)
				..controls(66,-157.3)and(65.6,-156.5)..(65.3,-156.1)
				..controls(64.6,-155.1)and(64.5,-154.9)..(63.7,-152.9)
				..controls(63.2,-151.6)and(62.9,-150.9)..(61.8,-149.6)
				..controls(61.5,-149.2)and(61.1,-148.5)..(60.9,-147.9)
				..controls(60.7,-147.4)and(60.3,-146.6)..(60.1,-146.1)
				..controls(59.9,-145.5)and(59.5,-144.9)..(59,-144.3)
				..controls(58.5,-143.7)and(58.2,-143.2)..(57.9,-142.5)
				..controls(57.7,-141.9)and(57.3,-141.1)..(57,-140.7)
				..controls(56.6,-140.3)and(56.1,-139.3)..(55.8,-138.5)
				..controls(55.5,-137.8)and(55.1,-136.8)..(54.9,-136.5)
				..controls(54.7,-136.1)and(54.3,-135.2)..(53.9,-134.5)
				..controls(53.6,-133.8)and(53,-132.7)..(52.6,-132)
				..controls(51.6,-130.5)and(51.4,-130.2)..(51.1,-129)
				..controls(50.6,-127.5)and(50.6,-127.4)..(49.7,-126.3)
				..controls(49.2,-125.7)and(48.6,-124.8)..(48.3,-124.2)
				..controls(48,-123.6)and(47.5,-122.7)..(47.3,-122.2)
				..controls(47,-121.8)and(46.5,-120.7)..(46.1,-119.8)
				..controls(45.8,-118.9)and(45.2,-117.8)..(44.9,-117.2)
				..controls(44.2,-116.1)and(43.3,-114.4)..(42.3,-112.2)
				..controls(41.8,-111.3)and(41.3,-110.4)..(41.1,-110.1)
				..controls(40.5,-109.4)and(40.3,-109)..(40,-107.8)
				..controls(39.8,-107)and(39.6,-106.5)..(39.2,-106)
				..controls(38.9,-105.6)and(38.4,-104.7)..(38,-104)
				..controls(37.7,-103.4)and(37.2,-102.5)..(37,-102.1)
				..controls(36.7,-101.7)and(36.3,-100.8)..(36,-100.1)
				..controls(35.6,-98.9)and(35.3,-98.4)..(34.6,-97.6)
				..controls(34.2,-97.1)and(34.2,-96.3)..(34.6,-96)
				..controls(35,-95.7)and(35.5,-95.9)..(36.2,-96.7)
				..controls(36.5,-97)and(37.2,-97.5)..(37.6,-97.7)..controls(38,-98)and(38.6,-98.4)..(39.1,-98.9)
				..controls(39.6,-99.4)and(40.3,-100)..(40.6,-100.3)
				..controls(41.3,-100.8)and(42.4,-101.8)..(43.7,-102.9)
				..controls(44,-103.2)and(44.6,-103.7)..(45,-103.9)
				..controls(45.4,-104.2)and(45.9,-104.6)..(46.1,-104.8)
				..controls(47.1,-105.9)and(48.8,-107.1)..(50.5,-107.9)
				..controls(51.8,-108.6)and(52.3,-109)..(53.4,-110)
				..controls(54.8,-111.2)and(55.6,-111.6)..(57,-112.1)
				..controls(57.5,-112.3)and(58.5,-112.8)..(59.4,-113.3)
				..controls(60.2,-113.7)and(61.3,-114.3)..(61.9,-114.5)
				..controls(62.4,-114.7)and(63.3,-115.1)..(63.8,-115.3)
				..controls(64.4,-115.6)and(65.4,-115.9)..(66.2,-116.1)
				..controls(67,-116.3)and(68.4,-116.6)..(69.2,-116.9)
				..controls(72.3,-117.9)and(72.8,-118)..(78,-118.1)
				..controls(84,-118.3)and(86.1,-118.1)..(88.1,-117.1)
				..controls(89.4,-116.5)and(93,-115)..(93.7,-114.9)
				..controls(94.8,-114.6)and(95.4,-114.2)..(95.9,-113.4)
				..controls(96.2,-113)and(96.7,-112.5)..(97.1,-112.2)
				..controls(98,-111.6)and(101.3,-108.4)..(102.1,-107.2)
				..controls(102.4,-106.8)and(102.9,-106.2)..(103.2,-105.8)
				..controls(103.5,-105.5)and(104,-104.8)..(104.2,-104.2)
				..controls(104.5,-103.6)and(105,-102.8)..(105.3,-102.4)
				..controls(105.6,-102)and(106.1,-101.2)..(106.5,-100.6)
				..controls(107.1,-99.5)and(108.6,-97.5)..(109.5,-96.5)
				..controls(109.9,-96.1)and(110.1,-96)..(110.4,-96)
				..controls(111.1,-96)and(111.6,-96.5)..(111.8,-97.6)
				..controls(112,-98.6)and(112.3,-99.3)..(112.8,-99.9)
				..controls(113.7,-100.9)and(115.2,-103.1)..(115.7,-104.2)
				..controls(116.1,-105)and(116.5,-105.5)..(116.9,-106)
				..controls(117.3,-106.4)and(117.8,-106.9)..(118,-107.3)
				..controls(118.2,-107.6)and(118.6,-108)..(119,-108.3)
				..controls(119.6,-108.8)and(121.1,-110.3)..(121.8,-111.2)
				..controls(122.1,-111.7)and(122.5,-112)..(123.4,-112.4)
				..controls(124.8,-113.2)and(125.8,-113.8)..(127,-114.6)
				..controls(127.5,-114.9)and(128.1,-115.3)..(128.2,-115.4)
				..controls(128.4,-115.5)and(129.3,-115.8)..(130.2,-116)
				..controls(131.1,-116.2)and(132.5,-116.6)..(133.3,-116.9)
				..controls(135.7,-117.9)and(136.1,-117.9)..(142.5,-117.8)
				..controls(145.8,-117.8)and(148.4,-117.7)..(148.8,-117.6)
				..controls(149.2,-117.6)and(150.2,-117.2)..(151,-116.8)
				..controls(152,-116.4)and(152.8,-116)..(153.6,-115.9)
				..controls(154.5,-115.7)and(155.2,-115.4)..(156.3,-114.9)
				..controls(157.1,-114.5)and(158.2,-114)..(158.7,-113.8)
				..controls(159.1,-113.7)and(159.8,-113.3)..(160.2,-113)
				..controls(160.6,-112.7)and(161.3,-112.4)..(161.7,-112.2)
				..controls(162.7,-111.9)and(163.2,-111.7)..(164.1,-111.1)
				..controls(164.5,-110.8)and(165.3,-110.4)..(165.9,-110)
				..controls(166.5,-109.7)and(167.3,-109.2)..(167.6,-108.9)
				..controls(167.9,-108.6)and(168.6,-108.2)..(169,-108)
				..controls(169.4,-107.8)and(170.2,-107.3)..(170.7,-106.9)
				..controls(171.2,-106.5)and(172.1,-105.8)..(172.7,-105.4)
				..controls(173.2,-105)and(174.1,-104.4)..(174.5,-104)
				..controls(174.9,-103.7)and(175.7,-103.1)..(176.2,-102.6)
				..controls(176.7,-102.2)and(177.4,-101.5)..(177.9,-101.2)
				..controls(178.3,-100.8)and(178.8,-100.4)..(179,-100.2)
				..controls(179.2,-100)and(179.6,-99.6)..(180,-99.3)
				..controls(180.4,-99.1)and(180.9,-98.6)..(181.2,-98.4)
				..controls(181.5,-98.1)and(182.1,-97.7)..(182.5,-97.5)
				..controls(183.2,-97.2)and(183.5,-96.9)..(184.3,-95.9)
				..controls(184.6,-95.6)and(184.6,-95.6)..(184.8,-95.7)
				..controls(185,-95.8)and(185.1,-96)..(185.1,-96.2)
				..controls(185.1,-96.5)and(185,-96.7)..(184.6,-97.2)
				..controls(184.1,-97.8)and(184,-97.9)..(183.9,-98.8)
				..controls(183.7,-100.3)and(183.3,-101.3)..(182.9,-101.8)
				..controls(182.6,-102.1)and(182.3,-102.6)..(182.2,-102.9)
				..controls(182,-103.3)and(181.6,-104)..(181.2,-104.6)
				..controls(180.9,-105.2)and(180.5,-106.1)..(180.4,-106.6)
				..controls(179.9,-108)and(179.7,-108.5)..(179,-109.4)
				..controls(178.6,-109.8)and(178.2,-110.5)..(177.9,-111)
				..controls(177.7,-111.5)and(177.3,-112.3)..(177,-112.7)
				..controls(176.6,-113.2)and(176.4,-113.7)..(176.3,-114.2)
				..controls(176.2,-114.7)and(175.9,-115.2)..(175.4,-116)
				..controls(175,-116.6)and(174.4,-117.7)..(174,-118.5)
				..controls(173.7,-119.2)and(173.3,-120.1)..(173,-120.5)
				..controls(172.8,-120.9)and(172.4,-121.8)..(172.1,-122.5)
				..controls(171.8,-123.2)and(171.3,-124.1)..(171,-124.6)
				..controls(170.8,-125.1)and(170.3,-126)..(169.9,-126.6)
				..controls(169.6,-127.3)and(169.2,-128.1)..(169,-128.3)
				..controls(168.8,-128.6)and(168.5,-129.3)..(168.3,-129.8)
				..controls(168.1,-130.3)and(167.8,-131.1)..(167.5,-131.6)
				..controls(166.9,-132.5)and(164.2,-137.9)..(164.2,-138.1)
				..controls(164.2,-138.1)and(163.9,-138.8)..(163.4,-139.7)
				..controls(162.9,-140.5)and(162.4,-141.5)..(162.3,-141.9)
				..controls(161.9,-143)and(161.6,-143.4)..(160.9,-144.1)
				..controls(160.4,-144.6)and(160.3,-144.8)..(160.1,-145.4)
				..controls(160,-145.8)and(159.6,-146.8)..(159.2,-147.8)
				..controls(158.7,-148.8)and(158.1,-150)..(157.9,-150.6)
				..controls(157.6,-151.1)and(157.2,-151.9)..(156.9,-152.4)
				..controls(156.5,-152.8)and(156.1,-153.7)..(155.8,-154.4)
				..controls(155.6,-155)and(155.2,-155.9)..(154.9,-156.3)
				..controls(154.7,-156.7)and(154.4,-157.4)..(154.2,-157.9)
				..controls(154.1,-158.4)and(153.8,-159)..(153.5,-159.3)
				..controls(152.8,-160.1)and(153,-160.1)..(142.4,-160.1)
				..controls(137.2,-160.1)and(131.8,-160.1)..(130.4,-160)--cycle
				(141.5,-103.7)..controls(140.8,-103)and(139.8,-102.1)..(139.2,-101.6)
				..controls(138.5,-101.2)and(137.9,-100.6)..(137.8,-100.4)
				..controls(137,-99.5)and(134,-96.5)..(133.2,-96)
				..controls(132.7,-95.7)and(131.8,-94.9)..(131.1,-94.3)
				..controls(129.4,-92.7)and(127.9,-91.6)..(126.6,-91.1)
				..controls(125.2,-90.5)and(124.1,-90)..(122.9,-89.1)
				..controls(122.1,-88.6)and(121.3,-88.2)..(119.9,-87.7)
				..controls(118.8,-87.3)and(117.5,-86.8)..(116.9,-86.5)
				..controls(114.3,-85.5)and(113.7,-85.4)..(109.7,-85.5)
				..controls(105.3,-85.7)and(103.9,-85.9)..(101.6,-87)
				..controls(100.7,-87.4)and(99.4,-87.8)..(97.9,-88)
				..controls(97,-88.2)and(96.6,-88.3)..(95.8,-88.9)
				..controls(95.4,-89.2)and(94.7,-89.6)..(94.2,-89.8)
				..controls(92.9,-90.3)and(92.3,-90.7)..(90.9,-92.1)
				..controls(90.2,-92.7)and(89.4,-93.3)..(89,-93.5)..controls(87,-94.5)and(85.2,-96)..(83.6,-97.9)
				..controls(83.3,-98.2)and(82.5,-98.9)..(81.8,-99.4)
				..controls(80.7,-100.1)and(80.5,-100.3)..(80,-101.1)
				..controls(78.7,-103.1)and(77.2,-103.9)..(76.4,-103)
				..controls(75.7,-102.2)and(75.7,-103.2)..(75.7,-83.6)
				..controls(75.6,-73.6)and(75.6,-64.8)..(75.5,-64)
				..controls(75.5,-63.2)and(75.5,-61.4)..(75.6,-59.8)--(75.7,-57.1)--(76,-56.8)
				..controls(76.8,-56.2)and(78,-56.3)..(79.8,-57.2)
				..controls(80.4,-57.5)and(81.3,-57.8)..(81.8,-57.9)
				..controls(82.3,-58.1)and(83,-58.4)..(83.6,-58.8)..controls(84.3,-59.2)and(85,-59.5)..(86,-59.8)
				..controls(86.8,-60.1)and(88,-60.6)..(88.7,-60.9)
				..controls(89.3,-61.2)and(90.6,-61.7)..(91.4,-62)
				..controls(92.2,-62.3)and(93.7,-62.9)..(94.6,-63.3)
				..controls(95.6,-63.7)and(96.8,-64.1)..(97.4,-64.3)
				..controls(98,-64.5)and(99,-64.9)..(99.7,-65.2)
				..controls(100.4,-65.6)and(101.5,-66)..(102.3,-66.2)
				..controls(104.2,-66.8)and(105,-67.1)..(106.2,-67.8)
				..controls(107.9,-68.7)and(110.2,-69.4)..(111.3,-69.2)
				..controls(112.5,-69)and(115,-68.2)..(115.6,-67.8)
				..controls(117.3,-66.7)and(117.6,-66.6)..(119.5,-66.2)
				..controls(121.9,-65.8)and(122.1,-65.7)..(123,-65.1)
				..controls(124.1,-64.4)and(124.4,-64.3)..(126,-63.9)
				..controls(127,-63.7)and(127.7,-63.4)..(128.4,-63.1)
				..controls(129,-62.8)and(129.7,-62.5)..(130.6,-62.4)
				..controls(132.1,-62.1)and(134,-61.5)..(135.2,-61)
				..controls(135.7,-60.8)and(136.7,-60.5)..(137.3,-60.3)
				..controls(138.4,-60)and(139.3,-59.7)..(141.1,-59)
				..controls(141.6,-58.7)and(142,-58.6)..(142.5,-58.6)
				..controls(143.2,-58.6)and(143.3,-58.7)..(143.6,-59)
				..controls(144,-59.5)and(144.1,-59.9)..(144,-60.8)
				..controls(144,-61.2)and(144,-71.1)..(144,-82.8)--(143.9,-104.1)--(143.7,-104.5)
				..controls(143.5,-104.7)and(143.3,-104.8)..(143,-104.9)
				..controls(142.6,-104.9)and(142.5,-104.8)..(141.5,-103.7)--cycle
				(27.8,-91.4)..controls(27.1,-90.9)and(26.3,-90)..(26.1,-89.5)
				..controls(26,-89.2)and(25.8,-88.6)..(25.7,-87.9)
				..controls(25.4,-86.2)and(25.6,-84.8)..(26.5,-83.5)
				..controls(27.4,-82.3)and(28.3,-82)..(29.4,-82.3)
				..controls(31.1,-82.8)and(32.1,-84.9)..(32,-87.4)
				..controls(31.9,-89.3)and(31.2,-90.6)..(29.9,-91.4)
				..controls(29.1,-91.9)and(28.6,-91.9)..(27.8,-91.4)--cycle
				(190.1,-91.3)..controls(189.7,-91.2)and(189.3,-90.9)..(189.1,-90.7)
				..controls(188.3,-89.8)and(187.6,-87.7)..(187.6,-86.2)
				..controls(187.6,-85.5)and(187.7,-85.2)..(188.2,-84.4)
				..controls(189.1,-82.6)and(190.1,-81.8)..(191.3,-82.1)
				..controls(192,-82.3)and(193.5,-83.7)..(193.8,-84.4)
				..controls(194.1,-85.1)and(194.2,-87.1)..(194,-88.3)
				..controls(193.8,-89.8)and(193.3,-90.5)..(192.3,-91.3)
				..controls(191.7,-91.7)and(190.9,-91.7)..(190.1,-91.3)--cycle
				(173.7,-83.4)..controls(172.2,-82.5)and(172,-82.5)..(170.6,-82.4)
				..controls(169.7,-82.4)and(168.8,-82.4)..(168.3,-82.5)
				..controls(167.2,-82.7)and(166.9,-82.7)..(166.6,-82.2)
				..controls(166.3,-81.8)and(166.3,-81)..(166.8,-80)
				..controls(167.2,-78.9)and(167.4,-77.4)..(167.7,-74.1)
				..controls(167.8,-72.8)and(167.9,-71.8)..(167.9,-71.7)
				..controls(168,-71.7)and(168,-70.2)..(168,-68.4)..controls(168,-65.4)and(168,-64.9)..(167.7,-63)
				..controls(167.3,-60.5)and(167.1,-59.9)..(166.4,-58.4)
				..controls(166.1,-57.7)and(165.8,-56.9)..(165.7,-56.6)
				..controls(165.7,-56.2)and(165.4,-55.6)..(165.2,-55.2)
				..controls(164.6,-53.8)and(161.9,-50.7)..(161,-50.3)
				..controls(159.8,-49.7)and(157.9,-50)..(155.2,-51.2)
				..controls(154.7,-51.4)and(153.4,-51.8)..(152.2,-52.1)
				..controls(150.8,-52.4)and(149.4,-52.9)..(148.6,-53.2)
				..controls(147.8,-53.6)and(146.7,-54)..(146,-54.2)
				..controls(145.4,-54.4)and(144.5,-54.7)..(144,-54.9)
				..controls(143.6,-55.2)and(142.7,-55.5)..(141.8,-55.7)
				..controls(141,-55.9)and(139.6,-56.3)..(138.7,-56.6)
				..controls(137.8,-57)and(136.3,-57.4)..(135.3,-57.6)
				..controls(133.8,-57.9)and(133.2,-58.1)..(132.1,-58.7)
				..controls(131.1,-59.2)and(130.3,-59.4)..(129.4,-59.6)
				..controls(128.7,-59.8)and(127.8,-60)..(127.4,-60.1)
				..controls(127,-60.3)and(126,-60.5)..(125.1,-60.7)
				..controls(123.2,-61.1)and(121.7,-61.5)..(120.2,-62.1)
				..controls(119.7,-62.3)and(118.1,-62.8)..(116.7,-63.2)
				..controls(114.6,-63.9)and(114,-64.1)..(113,-64.7)
				..controls(111.9,-65.4)and(111.8,-65.4)..(111,-65.5)
				..controls(109.9,-65.5)and(109.3,-65.4)..(107.7,-64.7)
				..controls(107,-64.4)and(106,-64)..(105.3,-63.9)
				..controls(103.5,-63.5)and(102.2,-63)..(101.2,-62.5)
				..controls(100.1,-61.8)and(99.3,-61.5)..(97.8,-61.1)
				..controls(97.2,-61)and(96.2,-60.7)..(95.7,-60.4)
				..controls(95.1,-60.1)and(94.2,-59.8)..(93.5,-59.6)
				..controls(92.9,-59.4)and(91.7,-59)..(91,-58.7)..controls(90.2,-58.4)and(88.9,-58)..(88,-57.7)
				..controls(87.2,-57.5)and(86.1,-57.1)..(85.5,-56.8)
				..controls(85,-56.6)and(84.1,-56.2)..(83.5,-56)..controls(82.9,-55.7)and(81.8,-55.2)..(81,-54.9)
				..controls(80.1,-54.5)and(79.1,-54.2)..(78.7,-54.1)
				..controls(77.9,-53.8)and(77.3,-53.6)..(76.2,-52.9)
				..controls(75.7,-52.6)and(75.2,-52.4)..(74.1,-52.2)
				..controls(72.5,-51.9)and(72.4,-51.8)..(71.3,-51.1)
				..controls(69.9,-50.1)and(69.4,-49.9)..(67.9,-49.8)
				..controls(67,-49.7)and(66.6,-49.6)..(66.3,-49.4)
				..controls(66.1,-49.3)and(65.7,-49)..(65.4,-48.8)
				..controls(65.1,-48.7)and(64.2,-48.4)..(63.6,-48.2)
				..controls(62.9,-48)and(61.7,-47.6)..(60.9,-47.2)
				..controls(60.1,-46.8)and(58.9,-46.4)..(58.2,-46.2)
				..controls(57.6,-46)and(56.8,-45.7)..(56.5,-45.5)
				..controls(54.9,-44.6)and(54.6,-44.5)..(53.2,-44.2)
				..controls(51.7,-43.8)and(51,-43.6)..(49.9,-42.9)
				..controls(49.5,-42.6)and(48.9,-42.4)..(47.5,-42)
				..controls(45.7,-41.6)and(45.6,-41.5)..(44.3,-40.6)
				..controls(43.5,-40.1)and(42.8,-39.6)..(42.7,-39.4)
				..controls(42.4,-38.7)and(43,-38.1)..(44.5,-37.6)..controls(44.9,-37.5)and(46,-37)..(46.9,-36.6)
				..controls(47.9,-36.1)and(48.8,-35.8)..(49.6,-35.6)
				..controls(50.5,-35.4)and(51.3,-35.2)..(52.4,-34.6)
				..controls(53.3,-34.1)and(54.5,-33.7)..(54.9,-33.5)
				..controls(55.4,-33.4)and(56.4,-33)..(57.2,-32.6)
				..controls(58.6,-32)and(60.1,-31.5)..(61.6,-31.2)
				..controls(62.3,-31.1)and(62.7,-31)..(63.2,-30.7)
				..controls(64.1,-30)and(64.3,-29.9)..(65.3,-29.7)
				..controls(65.7,-29.6)and(66.5,-29.2)..(67.1,-28.9)
				..controls(67.6,-28.6)and(68.8,-28.2)..(69.6,-27.9)
				..controls(70.5,-27.7)and(71.6,-27.3)..(72.1,-27)
				..controls(72.6,-26.7)and(73.8,-26.2)..(74.8,-25.9)
				..controls(75.7,-25.5)and(76.7,-25.1)..(77.1,-25)
				..controls(77.4,-24.8)and(78.5,-24.4)..(79.6,-24.1)
				..controls(80.6,-23.8)and(82,-23.2)..(82.8,-22.9)
				..controls(83.5,-22.5)and(84.8,-22)..(85.8,-21.8)
				..controls(86.7,-21.5)and(87.8,-21.1)..(88.1,-20.9)
				..controls(88.4,-20.7)and(89.3,-20.3)..(90.1,-20.1)
				..controls(90.8,-19.8)and(91.9,-19.4)..(92.4,-19.1)
				..controls(93.5,-18.5)and(94.4,-18.2)..(95.4,-18.1)
				..controls(95.7,-18)and(96.4,-17.8)..(96.9,-17.5)
				..controls(97.5,-17.3)and(98.4,-16.9)..(99,-16.6)
				..controls(99.7,-16.3)and(100.9,-15.9)..(101.9,-15.6)
				..controls(102.9,-15.3)and(104.4,-14.7)..(105.3,-14.3)
				..controls(108.1,-13.2)and(108.5,-13.1)..(110.1,-13.1)
				..controls(111.4,-13.1)and(111.5,-13.1)..(112.4,-13.5)
				..controls(113.2,-14)and(115,-14.8)..(116.6,-15.4)
				..controls(116.9,-15.5)and(117.9,-15.8)..(118.8,-16.1)
				..controls(120.9,-16.6)and(121.5,-16.9)..(122.6,-17.9)
				..controls(123,-18.3)and(123.7,-18.8)..(123.9,-18.9)
				..controls(124.6,-19.4)and(126,-19.8)..(127.9,-20.2)
				..controls(129.6,-20.5)and(129.8,-20.6)..(131.1,-21.4)
				..controls(131.5,-21.7)and(132.4,-22.1)..(133.1,-22.3)
				..controls(136.1,-23.2)and(138.3,-24.1)..(140.2,-25)
				..controls(141,-25.5)and(141.5,-25.6)..(142.5,-25.8)
				..controls(145.1,-26.3)and(146,-26.6)..(149.4,-28.1)
				..controls(151.3,-28.9)and(153.5,-29.8)..(154.3,-30.1)
				..controls(155.1,-30.4)and(156.2,-30.8)..(156.6,-31)
				..controls(157.1,-31.2)and(158.2,-31.6)..(159.1,-31.8)
				..controls(160.6,-32.2)and(160.9,-32.3)..(161.7,-32.9)
				..controls(162.4,-33.3)and(162.9,-33.6)..(163.7,-33.8)
				..controls(164.3,-34)and(164.8,-34.1)..(164.9,-34.1)
				..controls(165,-34.1)and(165.6,-34.4)..(166.4,-34.8)
				..controls(167.3,-35.2)and(168.1,-35.4)..(169,-35.6)
				..controls(170,-35.8)and(170.5,-36)..(171.7,-36.6)
				..controls(172.4,-37)and(173.3,-37.4)..(173.6,-37.4)
				..controls(175.6,-37.7)and(176.2,-38)..(176.7,-39.1)
				..controls(176.9,-39.6)and(176.9,-39.7)..(176.7,-40.4)
				..controls(176.4,-41.6)and(175.8,-42.7)..(175.2,-43.4)
				..controls(174.5,-44.1)and(172.6,-45.5)..(171.7,-45.8)
				..controls(170.6,-46.2)and(169.1,-46)..(168.2,-45.4)
				..controls(167.9,-45.2)and(167.3,-44.8)..(166.9,-44.4)
				..controls(166.4,-44.1)and(165.8,-43.6)..(165.4,-43.2)
				..controls(165.1,-42.9)and(164.4,-42.3)..(163.9,-41.9)
				..controls(163.3,-41.5)and(162.7,-40.9)..(162.4,-40.5)
				..controls(161.4,-39.2)and(159.9,-38)..(158.8,-37.8)
				..controls(158.4,-37.8)and(158.3,-37.7)..(158.2,-37.2)
				..controls(158,-36.8)and(157.4,-36.1)..(157.1,-36.1)
				..controls(157,-36.1)and(156.7,-36.1)..(156.5,-36)
				..controls(156.1,-35.9)and(155.9,-35.7)..(155.7,-35.2)
				..controls(155.3,-34.6)and(155.1,-34.4)..(154,-34.1)
				..controls(153.6,-34)and(153.2,-33.7)..(152.9,-33.5)
				..controls(152.4,-32.9)and(151.6,-32.4)..(151.3,-32.4)
				..controls(150.9,-32.4)and(150.4,-32.6)..(150.4,-32.9)
				..controls(150.4,-33.2)and(150.8,-33.4)..(151.4,-33.6)
				..controls(152.1,-33.7)and(152.8,-34.1)..(153.5,-34.9)
				..controls(154.1,-35.5)and(154.3,-35.6)..(155,-35.8)
				..controls(155.7,-36)and(155.7,-36)..(156,-36.6)
				..controls(156.3,-37.3)and(156.7,-37.7)..(157.5,-37.8)
				..controls(158.1,-37.9)and(158.2,-38)..(158.4,-38.8)
				..controls(158.5,-39.4)and(158.6,-39.6)..(159.6,-40.4)
				..controls(160.2,-41)and(160.9,-41.6)..(161.3,-41.8)
				..controls(161.7,-42)and(161.9,-42.3)..(162.2,-42.8)
				..controls(162.4,-43.2)and(162.8,-43.8)..(163.1,-44.1)
				..controls(163.9,-44.8)and(165.1,-46.7)..(165.7,-48)
				..controls(166.2,-49.3)and(166.4,-49.7)..(167.2,-50.5)
				..controls(167.5,-50.8)and(167.9,-51.4)..(168.1,-51.8)
				..controls(168.3,-52.2)and(168.8,-53)..(169.1,-53.6)
				..controls(169.6,-54.3)and(169.8,-54.9)..(169.9,-55.6)
				..controls(170.4,-57.8)and(170.5,-58.1)..(171.2,-59.5)
				..controls(171.6,-60.3)and(172,-61.5)..(172.2,-62.1)
				..controls(172.3,-62.8)and(172.7,-63.7)..(172.9,-64.1)
				..controls(173.2,-64.6)and(173.4,-65.3)..(173.4,-65.6)
				..controls(173.8,-68.8)and(174.2,-70.6)..(174.9,-72.6)
				..controls(175.5,-74.5)and(175.9,-77.1)..(176.1,-80.2)
				..controls(176.3,-83.4)and(175.5,-84.4)..(173.7,-83.4)--cycle;
				\definecolor{ezs}{RGB}{247,179,26}\fill[ezs](-5.7,-156.6)--(-3.4,-157.5)--(61.9,12.7)--(59.7,13.6)--cycle;
				\definecolor{ezs}{RGB}{247,179,22}\fill[ezs](-3.8,-157.3)--(-1.5,-158.2)--(63.9,12)--(61.6,12.9)--cycle;
				\definecolor{ezs}{RGB}{248,179,19}\fill[ezs](-1.8,-158.1)--(0.4,-158.9)--(65.8,11.2)--(63.5,12.1)--cycle;
				\definecolor{ezs}{RGB}{248,178,16}\fill[ezs](0.1,-158.8)--(2.4,-159.7)--(67.7,10.5)--(65.4,11.4)--cycle;
				\definecolor{ezs}{RGB}{249,178,15}\fill[ezs](2,-159.5)--(4.3,-160.4)--(69.6,9.8)--(67.3,10.7)--cycle;
				\definecolor{ezs}{RGB}{251,177,17}\fill[ezs](3.9,-160.3)--(6.2,-161.1)--(71.5,9)--(69.2,9.9)--cycle;
				\definecolor{ezs}{RGB}{252,179,18}\fill[ezs](5.8,-161)--(8.1,-161.9)--(73.4,8.3)--(71.1,9.2)--cycle;
				\definecolor{ezs}{RGB}{252,182,17}\fill[ezs](7.7,-161.7)--(10,-162.6)--(75.3,7.6)--(73,8.5)--cycle;
				\definecolor{ezs}{RGB}{252,185,17}\fill[ezs](9.6,-162.5)--(11.9,-163.4)--(77.2,6.8)--(75,7.7)--cycle;
				\definecolor{ezs}{RGB}{251,189,18}\fill[ezs](11.5,-163.2)--(13.8,-164.1)--(79.2,6.1)--(76.9,7)--cycle;
				\definecolor{ezs}{RGB}{251,191,22}\fill[ezs](13.4,-163.9)--(15.7,-164.8)--(81.1,5.4)--(78.8,6.2)--cycle;
				\definecolor{ezs}{RGB}{252,193,25}\fill[ezs](15.4,-164.7)--(17.6,-165.6)--(83,4.6)--(80.7,5.5)--cycle;
				\definecolor{ezs}{RGB}{253,195,27}\fill[ezs](17.3,-165.4)--(19.6,-166.3)--(84.9,3.9)--(82.6,4.8)--cycle;
				\definecolor{ezs}{RGB}{253,196,31}\fill[ezs](19.2,-166.1)--(21.5,-167)--(86.8,3.2)--(84.5,4)--cycle;
				\definecolor{ezs}{RGB}{254,199,35}\fill[ezs](21.1,-166.9)--(23.4,-167.8)--(88.7,2.4)--(86.4,3.3)--cycle;
				\definecolor{ezs}{RGB}{253,201,38}\fill[ezs](23,-167.6)--(25.3,-168.5)--(90.6,1.7)--(88.3,2.6)--cycle;
				\definecolor{ezs}{RGB}{253,203,40}\fill[ezs](24.9,-168.3)--(27.2,-169.2)--(92.5,1)--(90.3,1.8)--cycle;
				\definecolor{ezs}{RGB}{253,205,42}\fill[ezs](26.8,-169.1)--(29.1,-170)--(94.4,0.2)--(92.2,1.1)--cycle;
				\definecolor{ezs}{RGB}{253,207,45}\fill[ezs](28.7,-169.8)--(31,-170.7)--(96.4,-0.5)--(94.1,0.4)--cycle;
				\definecolor{ezs}{RGB}{253,210,49}\fill[ezs](30.7,-170.5)--(32.9,-171.4)--(98.3,-1.2)--(96,-0.4)--cycle;
				\definecolor{ezs}{RGB}{253,212,54}\fill[ezs](32.6,-171.3)--(34.9,-172.2)--(100.2,-2)--(97.9,-1.1)--cycle;
				\definecolor{ezs}{RGB}{253,214,59}\fill[ezs](34.5,-172)--(36.8,-172.9)--(102.1,-2.7)--(99.8,-1.8)--cycle;
				\definecolor{ezs}{RGB}{253,216,63}\fill[ezs](36.4,-172.7)--(38.7,-173.6)--(104,-3.4)--(101.7,-2.6)--cycle;
				\definecolor{ezs}{RGB}{253,218,65}\fill[ezs](38.3,-173.5)--(40.6,-174.4)--(105.9,-4.2)--(103.6,-3.3)--cycle;
				\definecolor{ezs}{RGB}{254,220,66}\fill[ezs](40.2,-174.2)--(42.5,-175.1)--(107.8,-4.9)--(105.5,-4)--cycle;
				\definecolor{ezs}{RGB}{254,222,68}\fill[ezs](42.1,-175)--(44.4,-175.8)--(109.7,-5.6)--(107.5,-4.8)--cycle;
				\definecolor{ezs}{RGB}{254,224,71}\fill[ezs](44,-175.7)--(46.3,-176.6)--(111.7,-6.4)--(109.4,-5.5)--cycle;
				\definecolor{ezs}{RGB}{254,226,75}\fill[ezs](46,-176.4)--(48.2,-177.3)--(113.6,-7.1)--(111.3,-6.2)--cycle;
				\definecolor{ezs}{RGB}{254,228,79}\fill[ezs](47.9,-177.2)--(50.1,-178)--(115.5,-7.8)--(113.2,-7)--cycle;
				\definecolor{ezs}{RGB}{254,230,82}\fill[ezs](49.8,-177.9)--(52.1,-178.8)--(117.4,-8.6)--(115.1,-7.7)--cycle;
				\definecolor{ezs}{RGB}{254,232,85}\fill[ezs](51.7,-178.6)--(54,-179.5)--(119.3,-9.3)--(117,-8.4)--cycle;
				\definecolor{ezs}{RGB}{254,234,88}\fill[ezs](53.6,-179.4)--(55.9,-180.2)--(121.2,-10)--(118.9,-9.2)--cycle;
				\definecolor{ezs}{RGB}{254,237,91}\fill[ezs](55.5,-180.1)--(57.8,-181)--(123.1,-10.8)--(120.8,-9.9)--cycle;
				\definecolor{ezs}{RGB}{254,239,94}\fill[ezs](57.4,-180.8)--(59.7,-181.7)--(125,-11.5)--(122.8,-10.6)--cycle;
				\definecolor{ezs}{RGB}{254,241,97}\fill[ezs](59.3,-181.6)--(61.6,-182.4)--(127,-12.2)--(124.7,-11.4)--cycle;
				\definecolor{ezs}{RGB}{254,243,101}\fill[ezs](61.2,-182.3)--(63.5,-183.2)--(128.9,-13)--(126.6,-12.1)--cycle;
				\definecolor{ezs}{RGB}{254,245,103}\fill[ezs](63.2,-183)--(65.4,-183.9)--(130.8,-13.7)--(128.5,-12.8)--cycle;
				\definecolor{ezs}{RGB}{254,245,105}\fill[ezs](65.1,-183.8)--(67.4,-184.6)--(132.7,-14.4)--(130.4,-13.6)--cycle;
				\definecolor{ezs}{RGB}{254,245,105}\fill[ezs](67,-184.5)--(69.3,-185.4)--(134.6,-15.2)--(132.3,-14.3)--cycle;
				\definecolor{ezs}{RGB}{254,244,103}\fill[ezs](68.9,-185.2)--(71.2,-186.1)--(136.5,-15.9)--(134.2,-15)--cycle;
				\definecolor{ezs}{RGB}{253,243,101}\fill[ezs](70.8,-186)--(73.1,-186.8)--(138.4,-16.6)--(136.1,-15.8)--cycle;
				\definecolor{ezs}{RGB}{253,241,98}\fill[ezs](72.7,-186.7)--(75,-187.6)--(140.3,-17.4)--(138,-16.5)--cycle;
				\definecolor{ezs}{RGB}{253,239,95}\fill[ezs](74.6,-187.4)--(76.9,-188.3)--(142.2,-18.1)--(140,-17.2)--cycle;
				\definecolor{ezs}{RGB}{253,236,92}\fill[ezs](76.5,-188.2)--(78.8,-189)--(144.2,-18.8)--(141.9,-18)--cycle;
				\definecolor{ezs}{RGB}{253,234,88}\fill[ezs](78.5,-188.9)--(80.7,-189.8)--(146.1,-19.6)--(143.8,-18.7)--cycle;
				\definecolor{ezs}{RGB}{253,232,85}\fill[ezs](80.4,-189.6)--(82.7,-190.5)--(148,-20.3)--(145.7,-19.4)--cycle;
				\definecolor{ezs}{RGB}{253,230,82}\fill[ezs](82.3,-190.4)--(84.6,-191.2)--(149.9,-21)--(147.6,-20.2)--cycle;
				\definecolor{ezs}{RGB}{253,228,78}\fill[ezs](84.2,-191.1)--(86.5,-192)--(151.8,-21.8)--(149.5,-20.9)--cycle;
				\definecolor{ezs}{RGB}{253,226,74}\fill[ezs](86.1,-191.8)--(88.4,-192.7)--(153.7,-22.5)--(151.4,-21.6)--cycle;
				\definecolor{ezs}{RGB}{253,224,70}\fill[ezs](88,-192.6)--(90.3,-193.4)--(155.6,-23.3)--(153.3,-22.4)--cycle;
				\definecolor{ezs}{RGB}{253,222,67}\fill[ezs](89.9,-193.3)--(92.2,-194.2)--(157.5,-24)--(155.3,-23.1)--cycle;
				\definecolor{ezs}{RGB}{253,220,65}\fill[ezs](91.8,-194)--(94.1,-194.9)--(159.5,-24.7)--(157.2,-23.8)--cycle;
				\definecolor{ezs}{RGB}{253,217,63}\fill[ezs](93.7,-194.8)--(96,-195.6)--(161.4,-25.5)--(159.1,-24.6)--cycle;
				\definecolor{ezs}{RGB}{253,215,60}\fill[ezs](95.7,-195.5)--(97.9,-196.4)--(163.3,-26.2)--(161,-25.3)--cycle;
				\definecolor{ezs}{RGB}{253,213,56}\fill[ezs](97.6,-196.2)--(99.9,-197.1)--(165.2,-26.9)--(162.9,-26)--cycle;
				\definecolor{ezs}{RGB}{252,211,53}\fill[ezs](99.5,-197)--(101.8,-197.8)--(167.1,-27.7)--(164.8,-26.8)--cycle;
				\definecolor{ezs}{RGB}{252,209,49}\fill[ezs](101.4,-197.7)--(103.7,-198.6)--(169,-28.4)--(166.7,-27.5)--cycle;
				\definecolor{ezs}{RGB}{252,207,45}\fill[ezs](103.3,-198.4)--(105.6,-199.3)--(170.9,-29.1)--(168.6,-28.2)--cycle;
				\definecolor{ezs}{RGB}{253,205,42}\fill[ezs](105.2,-199.2)--(107.5,-200)--(172.8,-29.9)--(170.6,-29)--cycle;
				\definecolor{ezs}{RGB}{253,202,40}\fill[ezs](107.1,-199.9)--(109.4,-200.8)--(174.7,-30.6)--(172.5,-29.7)--cycle;
				\definecolor{ezs}{RGB}{253,200,38}\fill[ezs](109,-200.6)--(111.3,-201.5)--(176.7,-31.3)--(174.4,-30.4)--cycle;
				\definecolor{ezs}{RGB}{253,198,35}\fill[ezs](111,-201.4)--(113.2,-202.2)--(178.6,-32.1)--(176.3,-31.2)--cycle;
				\definecolor{ezs}{RGB}{253,196,31}\fill[ezs](112.9,-202.1)--(115.2,-203)--(180.5,-32.8)--(178.2,-31.9)--cycle;
				\definecolor{ezs}{RGB}{253,194,28}\fill[ezs](114.8,-202.8)--(117.1,-203.7)--(182.4,-33.5)--(180.1,-32.6)--cycle;
				\definecolor{ezs}{RGB}{253,192,26}\fill[ezs](116.7,-203.6)--(119,-204.5)--(184.3,-34.3)--(182,-33.4)--cycle;
				\definecolor{ezs}{RGB}{253,189,24}\fill[ezs](118.6,-204.3)--(120.9,-205.2)--(186.2,-35)--(183.9,-34.1)--cycle;
				\definecolor{ezs}{RGB}{253,186,20}\fill[ezs](120.5,-205)--(122.8,-205.9)--(188.1,-35.7)--(185.8,-34.8)--cycle;
				\definecolor{ezs}{RGB}{253,184,17}\fill[ezs](122.4,-205.8)--(124.7,-206.7)--(190,-36.5)--(187.8,-35.6)--cycle;
				\definecolor{ezs}{RGB}{252,183,15}\fill[ezs](124.3,-206.5)--(126.6,-207.4)--(192,-37.2)--(189.7,-36.3)--cycle;
				\definecolor{ezs}{RGB}{252,181,14}\fill[ezs](126.3,-207.2)--(128.5,-208.1)--(193.9,-37.9)--(191.6,-37.1)--cycle;
				\definecolor{ezs}{RGB}{252,179,13}\fill[ezs](128.2,-208)--(130.4,-208.9)--(195.8,-38.7)--(193.5,-37.8)--cycle;
				\definecolor{ezs}{RGB}{252,178,12}\fill[ezs](130.1,-208.7)--(132.4,-209.6)--(197.7,-39.4)--(195.4,-38.5)--cycle;
				\definecolor{ezs}{RGB}{251,178,10}\fill[ezs](132,-209.4)--(134.3,-210.3)--(199.6,-40.1)--(197.3,-39.3)--cycle;
				\definecolor{ezs}{RGB}{251,178,10}\fill[ezs](133.9,-210.2)--(136.2,-211.1)--(201.5,-40.9)--(199.2,-40)--cycle;
				\definecolor{ezs}{RGB}{250,177,11}\fill[ezs](135.8,-210.9)--(138.1,-211.8)--(203.4,-41.6)--(201.1,-40.7)--cycle;
				\definecolor{ezs}{RGB}{250,178,11}\fill[ezs](137.7,-211.6)--(140,-212.5)--(205.3,-42.3)--(203.1,-41.5)--cycle;
				\definecolor{ezs}{RGB}{248,180,10}\fill[ezs](139.6,-212.4)--(141.9,-213.3)--(207.3,-43.1)--(205,-42.2)--cycle;
				\definecolor{ezs}{RGB}{247,181,9}\fill[ezs](141.5,-213.1)--(143.8,-214)--(209.2,-43.8)--(206.9,-42.9)--cycle;
				\definecolor{ezs}{RGB}{246,182,9}\fill[ezs](143.5,-213.8)--(145.7,-214.7)--(211.1,-44.5)--(208.8,-43.7)--cycle;
				\definecolor{ezs}{RGB}{246,182,10}\fill[ezs](145.4,-214.6)--(147.7,-215.5)--(213,-45.3)--(210.7,-44.4)--cycle;
			\end{scope}
			% chữ EDUZY (Times đậm, làm dày nét bằng 5 lớp lệch nhẹ; căn đúng khung chữ logo gốc)
			\foreach \dx/\dy in {-0.7/0,0.7/0,0/-0.3,0/0.3,0/0}
			\node[anchor=center,inner sep=0pt,outer sep=0pt,text=ezText,align=center] at (109.6+\dx,-193.6+\dy)
			{\resizebox{1.01cm}{0.22cm}{\fontencoding{T1}\fontfamily{ptm}\fontseries{b}\selectfont EDUZY}};
		\end{tikzpicture}%
	}%
}
%====================== THIẾT LẬP THÂN ĐỀ ======================
%--- Môi trường tự luận {bt}: ex_test 3.3.4 chưa khai báo sẵn nên tạo ở đây
%    (ex_test đã tự móc \fixshowans và \listenumerate cho bt)
\makeatletter
\@ifundefined{bt}{\newtheorem{bt}{\color{gold1}\namebt}}{}
\makeatother
\renewcommand{\sectionmark}[1]{\markright{\thesection\ #1}}
%--- Khung tiêu đề (logo + tối đa 4 dòng; dòng 3, 4 để trống {} thì bỏ qua)
%  Đề thi:    \khungtieude{THI THỬ THPT QUỐC GIA NĂM 2026}{Môn Toán -- Đề số 21}
%                         {(Đề thi gồm: ...)}{Thời gian làm bài: 90 phút (không kể thời gian phát đề).}
%  Chuyên đề: \khungtieude{CHUYÊN ĐỀ 3 -- DÃY SỐ VÀ CÔNG THỨC TRUY HỒI}
%                         {Bồi dưỡng học sinh giỏi Toán THPT}{...}{}
\newcommand{\khungtieude}[4]{%
	\begin{tcolorbox}[colframe=gold1,colback=gold5!45,boxrule=0.5mm]
		\begin{minipage}[c]{3cm}
			\centerline{\eduzylogo[2.4cm]}
		\end{minipage}\hspace{0.5cm}
		\begin{minipage}[c]{12.5cm}
			\vskip 0.15 cm
			\centerline{\ind{#1}}
			\vskip 0.15 cm
			\centerline{\textbf{#2}}
			\if\relax\detokenize{#3}\relax\else
				\vskip 0.15 cm
				\centerline{\textit{#3}}
			\fi
			\if\relax\detokenize{#4}\relax\else
				\vskip 0.15 cm
				\centerline{\textit{#4}}
			\fi
			\vskip 0.15 cm
		\end{minipage}
	\end{tcolorbox}
}
%====================== BỔ SUNG CHO BÀI HỌC, CHUYÊN ĐỀ ======================
\usepackage[shortlabels]{enumitem}
\usepackage{tabularx}
%--- Danh sách gọn: \begin{enumerate}[gon] ... \end{enumerate}
\SetEnumitemKey{gon}{leftmargin=*,itemsep=1pt,parsep=0pt,topsep=3pt}
%--- Khung lý thuyết: \begin{khung}{} ... \end{khung}; tiêu đề chỉ dùng khi thật ngắn ("Mẹo", "Lỗi thường gặp")
\newtcolorbox{khung}[1]{breakable,enhanced,colframe=gold1,colback=gold5!45,
	boxrule=0.4mm,fonttitle=\bfseries\sffamily,coltitle=white,colbacktitle=gold1,
	attach boxed title to top left={xshift=6mm,yshift=-3mm},
	boxed title style={arc=1mm,boxrule=0.3mm},title={#1}}
%--- Tiêu đề mục nhỏ trong một dạng: \muc{Dấu hiệu}, \muc{Phương pháp}, \muc{Bài tự luyện}
\makeatletter
\newcommand{\muc}[1]{\par\medskip\noindent\indam{#1}\par\nobreak\smallskip\@afterindentfalse\@afterheading}
\makeatother
%--- Ô công thức cần nhớ (căn giữa, viền vàng, vừa khít nội dung)
%     \begin{center}\congthuc{$Q(n)=(n-1)!$}\end{center}
\newtcbox{\congthuc}{on line,colframe=gold1,colback=gold5!45,boxrule=0.5pt,arc=2pt,
	left=8pt,right=8pt,top=4pt,bottom=4pt}
%--- Công tắc lời giải bài tự luyện (ví dụ mẫu luôn dùng \loigiai)
%     \loigiaiTL[{đáp số}]{lời giải đầy đủ}
%     \TLfull -> in lời giải đầy đủ;  \TLdapso -> chỉ in một dòng "Đáp số";  \TLan -> ẩn hoàn toàn
%     (\CheDo là dethi thì mọi lời giải và đáp số đều ẩn)
\makeatletter
\newif\ifTLhien   \TLhientrue
\newif\ifTLdapso  \TLdapsofalse
\newcommand{\TLfull}{\TLhientrue\TLdapsofalse}
\newcommand{\TLdapso}{\TLhientrue\TLdapsotrue}
\newcommand{\TLan}{\TLhienfalse}
\newcommand{\loigiaiTL}[2][]{%
	\ifTLhien
		\ifTLdapso
			\def\@tlarg{#1}%
			\ifx\@tlarg\@empty\else
				\ifdefstring{\kindOPTN}{dethi}{}{\par\noindent{\bfseries\color{gold1}Đáp số:} #1\par}%
			\fi
		\else
			\loigiai{#2}%
		\fi
	\fi}
\makeatother
\TLdapso% mặc định: chỉ in đáp số bài tự luyện
%====================== BẢNG ĐÚNG/SAI (\choiceTFt) ======================
%  Hàng tiêu đề nền nâu đồng chữ trắng (Phát biểu | Đ | S), các hàng tô xen kẽ kem / trắng,
%  đường kẻ vàng kim, nhãn a) b) c) d) màu nâu đồng. Bản có đáp án: cột được chọn có vòng tròn
%  nâu đồng với dấu tích trắng (câu đúng tích ở cột Đúng, câu sai tích ở cột Sai).
%  Tiêu đề ghi đủ "Đúng | Sai": \OPTN{viettat=0} sau \input (cột Đúng/Sai rộng ra, phát biểu hẹp lại).
%  Dùng bảng gốc của ex_test: \def\BangDSgoc{1} trước \input{preamble-eduzy-v2}.
\providecommand{\BangDSgoc}{0}
\ifnum\BangDSgoc=0
\usepackage{colortbl,pifont}
\colorlet{ezDSke}{gold2!60!gold1}% đường kẻ bảng (nét 0,4pt như ex_test để không đổi chiều cao bảng)
\colorlet{ezDSle}{gold5!60}% hàng a), c)
\colorlet{ezDSchan}{white}% hàng b), d)
%--- dấu chọn ở cột Đúng/Sai
%    (bề rộng 0, căn giữa ô: không làm đổi bề rộng cột của ex_test, không báo tràn lề)
\newcommand{\EZdauchon}{\makebox[0pt]{\tikz[baseline=-1ex]\node[circle,fill=gold1,inner sep=0pt,
	minimum size=1.2em,text=white,font=\small]{\ding{51}};}}
\def\tickT{\EZdauchon}\def\tickF{\EZdauchon}
\AtBeginDocument{\sbox{\ccTFTab}{\EZdauchon}}
%--- nhãn a) b) c) d) màu nâu đồng (ở bản lời giải, nhãn câu đúng vẫn đóng khung như ex_test)
\renewcommand{\FalseTF}{\stepcounter{dapan}\textbf{\color{gold1}\DapAnTF\sepTF}\ignorespaces}
%--- trước khi dựng bảng: kẻ màu vàng kim (bề rộng cột giữ nguyên như ex_test); sau bảng trả lại màu đen
\newcount\EZhangTF
\newcommand{\EZbatdauTF}{%
	\global\EZhangTF=0\relax
	\arrayrulecolor{ezDSke}}
\newcommand{\EZketthucTF}{\arrayrulecolor{black}}
%--- hàng tiêu đề nền nâu đồng, chữ trắng
\makeatletter
\newcommand{\EZo@TF}[1]{\centering\bfseries\color{white}#1}
\newcommand{\EZtieudeTF}{%
	\ifcase\cotDS
		\ifnum\cotSTT=1\relax
			\def\rowHeader{\hline\rowcolor{gold1}\EZo@TF{TT}&\EZo@TF{\phatbieu}\tabularnewline}%
		\else
			\def\rowHeader{\hline\rowcolor{gold1}\EZo@TF{\phatbieu}\tabularnewline}%
		\fi
	\or
		\ifnum\cotSTT=1\relax
			\def\rowHeader{\hline\rowcolor{gold1}\EZo@TF{TT}&\EZo@TF{\phatbieu}&\EZo@TF{\Dung}\tabularnewline}%
		\else
			\def\rowHeader{\hline\rowcolor{gold1}\EZo@TF{\phatbieu}&\EZo@TF{\Dung}\tabularnewline}%
		\fi
	\or
		\ifnum\cotSTT=1\relax
			\def\rowHeader{\hline\rowcolor{gold1}\EZo@TF{TT}&\EZo@TF{\phatbieu}&\EZo@TF{\Dung}&\EZo@TF{\Sai}\tabularnewline}%
		\else
			\def\rowHeader{\hline\rowcolor{gold1}\EZo@TF{\phatbieu}&\EZo@TF{\Dung}&\EZo@TF{\Sai}\tabularnewline}%
		\fi
	\fi}
%--- tô xen kẽ từng hàng (gọi ngay sau \hline đầu mỗi hàng)
\newcommand{\EZmauhangTF}{%
	\noalign{\global\advance\EZhangTF\@ne}%
	\ifodd\EZhangTF\rowcolor{ezDSle}\else\rowcolor{ezDSchan}\fi}
%--- gắn vào \choiceTFt / \writeChoiceTFt của ex_test
\xpretocmd{\choiceTFt}{\EZbatdauTF}{}{\typeout{[preamble-eduzy-v2] Không gắn được bảng Đúng/Sai (đầu bảng)}}
\xpatchcmd{\choiceTFt}{\ifnum\dongPB=0\def\rowHeader{}\fi}{\EZtieudeTF\ifnum\dongPB=0\def\rowHeader{}\fi}
	{}{\typeout{[preamble-eduzy-v2] Không gắn được bảng Đúng/Sai (hàng tiêu đề)}}
\xpatchcmd{\writeChoiceTFt}{\hline\xdef\dotEX{\sdotEX}}{\hline\EZmauhangTF\xdef\dotEX{\sdotEX}}
	{}{\typeout{[preamble-eduzy-v2] Không gắn được bảng Đúng/Sai (tô hàng)}}
\xpatchcmd{\writeChoiceTFt}{{\end{longtable}}}{{\end{longtable}\EZketthucTF}}
	{}{\typeout{[preamble-eduzy-v2] Không gắn được bảng Đúng/Sai (cuối longtable)}}
\xpatchcmd{\writeChoiceTFt}{{\end{xtabular}\end{center}\medskip\noindent}}{{\end{xtabular}\end{center}\EZketthucTF\medskip\noindent}}
	{}{\typeout{[preamble-eduzy-v2] Không gắn được bảng Đúng/Sai (cuối xtabular)}}
\fi
%=====================================================================
\renewcommand{\parallel}{\mathrel{/\!\!/}}

%      \TLdapso                      % tài liệu: \TLfull | \TLdapso | \TLan
%      \begin{document} ... \end{document}
%  Logo Eduzy vẽ bằng TikZ ngay trong file (\eduzylogo), không cần file ảnh.
%  Biên dịch: pdfLaTeX (máy có gói vietnam) hoặc XeLaTeX.
%============================================================
\documentclass[12pt,a4paper,twoside]{article}%<-- twoside: bat dau trang/chan trang chan-le
\usepackage{iftex}
\usepackage{amsmath,amssymb,tikz,mathrsfs,tkz-tab}
\usepackage{graphicx}
\usepackage{pgfplots}% đồ thị hàm số
\pgfplotsset{compat=1.17}
%\usepackage{mathptmx}
\usepackage{mathpazo}
\ifXeTeX% XeLaTeX: font chữ có đủ dấu tiếng Việt, giữ nguyên font toán mathpazo
	\usepackage[no-math]{fontspec}
	\setmainfont{TeX Gyre Pagella}
	\newfontfamily\EZqag{TeX Gyre Adventor}[NFSSFamily=qag,Ligatures=TeX]
\fi
\usepackage[many]{tcolorbox}
\usetikzlibrary{math,through,calc,intersections,angles,quotes,shapes,shapes.geometric,arrows,patterns,matrix,chains,arrows.meta,decorations.shapes,decorations.fractals,decorations.markings,shadows,decorations.pathreplacing}
\usepackage[top=1.6cm, bottom=1.8cm, left=1.8cm, right=1.8cm]{geometry}% v2: top 1.6cm (bản cũ 1.5cm) để chữ không sát dải đầu trang
\providecommand{\CheDo}{loigiai}% loigiai | dethi | solcolor | color | book
\expandafter\usepackage\expandafter[\CheDo]{ex_test}
% \tenchuong: khai báo ở file chính (trước \input) hoặc \renewcommand sau \input
%============================================================
%  [GỘP TỪ eduzy-de.tex] EDUZY - MÀU NÂU ĐỒNG / VÀNG KIM
%  + ĐẦU TRANG, CHÂN TRANG KIỂU DẢI NGANG (v2), ĐỐI XỨNG TRANG CHẴN / TRANG LẺ.
%  (không cần file eduzy-de.tex riêng nữa)
%  Khối này phải đặt SAU \usepackage[...]{ex_test} và SAU \def\tenchuong{...}
%============================================================
\usepackage{lastpage}
\usepackage{needspace}
%--- DÀN TRANG ĐỀ THI: tách câu (THÊM 05/11/2026; chỉ thêm lệnh mới, KHÔNG đổi hành vi mặc định) ---
%  Hai công tắc, mặc định tắt (= preamble chạy y như cũ). Bật bằng cách đặt TRƯỚC %============================================================
%  PREAMBLE EDUZY v2 - dùng chung cho đề thi, bài học, chuyên đề (gói ex_test)
%  Khác bản preamble-eduzy.tex ở đầu trang, chân trang: dải nền nâu đồng tràn ngang,
%  logo + Eduzy.com.vn, ô chữ vàng kem cạnh xiên, chân trang "Trang (số)";
%  logo Eduzy chìm giữa mọi trang (tắt: \logochimtat, bật lại: \logochimbat).
%  Cách dùng ở file chính (xem mau-chuan.tex, mau-chuyen-de.tex):
%      \def\CheDo{loigiai}          % loigiai | dethi | solcolor | color | book
%      \def\tenchuong{<dòng chữ trên đầu trang>}
%      \def\thuonghieu{<chữ ở chân trang lẻ>}   % không khai báo thì dùng mặc định
%      \def\tacgia{<chữ ở chân trang chẵn>}      % không khai báo thì dùng mặc định
%      \def\khauhieu{<chữ cạnh logo ở đầu trang>}  % tùy chọn, mặc định Eduzy.com.vn
%      \input{preamble-eduzy-v2}
%      \TLdapso                      % tài liệu: \TLfull | \TLdapso | \TLan
%      \begin{document} ... \end{document}
%  Logo Eduzy vẽ bằng TikZ ngay trong file (\eduzylogo), không cần file ảnh.
%  Biên dịch: pdfLaTeX (máy có gói vietnam) hoặc XeLaTeX.
%============================================================
\documentclass[12pt,a4paper,twoside]{article}%<-- twoside: bat dau trang/chan trang chan-le
\usepackage{iftex}
\usepackage{amsmath,amssymb,tikz,mathrsfs,tkz-tab}
\usepackage{graphicx}
\usepackage{pgfplots}% đồ thị hàm số
\pgfplotsset{compat=1.17}
%\usepackage{mathptmx}
\usepackage{mathpazo}
\ifXeTeX% XeLaTeX: font chữ có đủ dấu tiếng Việt, giữ nguyên font toán mathpazo
	\usepackage[no-math]{fontspec}
	\setmainfont{TeX Gyre Pagella}
	\newfontfamily\EZqag{TeX Gyre Adventor}[NFSSFamily=qag,Ligatures=TeX]
\fi
\usepackage[many]{tcolorbox}
\usetikzlibrary{math,through,calc,intersections,angles,quotes,shapes,shapes.geometric,arrows,patterns,matrix,chains,arrows.meta,decorations.shapes,decorations.fractals,decorations.markings,shadows,decorations.pathreplacing}
\usepackage[top=1.6cm, bottom=1.8cm, left=1.8cm, right=1.8cm]{geometry}% v2: top 1.6cm (bản cũ 1.5cm) để chữ không sát dải đầu trang
\providecommand{\CheDo}{loigiai}% loigiai | dethi | solcolor | color | book
\expandafter\usepackage\expandafter[\CheDo]{ex_test}
% \tenchuong: khai báo ở file chính (trước \input) hoặc \renewcommand sau \input
%============================================================
%  [GỘP TỪ eduzy-de.tex] EDUZY - MÀU NÂU ĐỒNG / VÀNG KIM
%  + ĐẦU TRANG, CHÂN TRANG KIỂU DẢI NGANG (v2), ĐỐI XỨNG TRANG CHẴN / TRANG LẺ.
%  (không cần file eduzy-de.tex riêng nữa)
%  Khối này phải đặt SAU \usepackage[...]{ex_test} và SAU \def\tenchuong{...}
%============================================================
\usepackage{lastpage}
\usepackage{needspace}
%--- DÀN TRANG ĐỀ THI: tách câu (THÊM 05/11/2026; chỉ thêm lệnh mới, KHÔNG đổi hành vi mặc định) ---
%  Hai công tắc, mặc định tắt (= preamble chạy y như cũ). Bật bằng cách đặt TRƯỚC \input{preamble-eduzy-v2}
%  trong file đề (tools/dan_trang/dan_trang.py tự ghi khi đặt \NeedspaceCau):
%    \def\TachCauDS{1}   câu Đúng/Sai được tách: câu dẫn (kèm hình) ở cuối trang, bảng Phát biểu sang trang sau
%                        (câu dẫn và hình không tách nhau, bảng không tách giữa hai trang)
%    \def\TachCauLG{1}   bản lời giải (\CheDo khác dethi): phần đề của câu cao hơn 1/3 trang được phép tách
%  Lệnh dùng trong file đề (giá trị do tools/dan_trang/dan_trang.py đo ở bản dethi):
%    \NeedspaceCau[<cao câu dẫn + hình>]{<cao cả khối câu>}   ngay trên \begin{ex} (câu đầu phần: trên tiêu đề phần)
%        công tắc tắt: đúng bằng \Needspace{<cao cả khối câu>} như cũ;
%        \TachCauDS=1 và có [..] (câu Đúng/Sai được chọn tách): chỉ giữ câu dẫn + hình;
%        \TachCauLG=1 ở bản lời giải: đòi tối đa 1/3 \textheight.
%    \NeedspaceBang{<cao bảng Phát biểu>}%   ngay trước \choiceTFt: \TachCauDS=1 thì giữ trọn bảng trên một trang;
%        công tắc tắt: không làm gì.
\providecommand{\TachCauDS}{0}
\providecommand{\TachCauLG}{0}
\newlength{\EZnsCau}
\newcommand{\NeedspaceCau}[2][]{%
	\setlength{\EZnsCau}{#2}%
	\ifnum\TachCauDS=1 \if\relax\detokenize{#1}\relax\else\setlength{\EZnsCau}{#1}\fi\fi
	\ifnum\TachCauLG=1 \ifdefstring{\kindOPTN}{dethi}{}{\ifdim\EZnsCau>0.3333\textheight\setlength{\EZnsCau}{0.3333\textheight}\fi}\fi
	\Needspace{\EZnsCau}}
\newcommand{\NeedspaceBang}[1]{\ifnum\TachCauDS=1 \Needspace{#1}\fi}

%--- Thang màu: đậm -> nhạt
\definecolor{gold1}{RGB}{126,96,0}    % nâu đồng đậm  #7E6000
\definecolor{gold2}{RGB}{192,144,3}   % vàng kim      #C09003
\definecolor{gold3}{RGB}{255,215,99}  % vàng nhạt     #FFD763
\definecolor{gold4}{RGB}{255,231,153} % vàng rất nhạt #FFE799
\definecolor{gold5}{RGB}{255,243,205} % vàng kem      #FFF3CD
%--- Ghi đè tông xanh/đỏ cũ
\definecolor{myblue}{RGB}{126,96,0}
\colorlet{blue}{gold1}
\colorlet{cyan}{gold2}
\colorlet{teal}{gold2}

%--- Tên chương mặc định (file đề nên tự khai báo lại)
\providecommand{\tenchuong}{}
%--- Chân trang
\providecommand{\thuonghieu}{TOÁN 12 -- Eduzy.com.vn}
\providecommand{\tacgia}{Tô Gia Bảo -- 0353703783}

%--- Chữ cạnh logo ở đầu trang
\providecommand{\khauhieu}{Eduzy.com.vn}

%====================== THÔNG SỐ ĐẦU TRANG, CHÂN TRANG (v2) ======================
%  Dải nền đậm tràn ngang trang; ô chữ vàng kem có cạnh xiên, kèm 2 vạch xiên chuyển màu.
%  Đầu trang lẻ : logo + \khauhieu (trái)   | ô chữ \tenchuong (phải)
%  Chân trang lẻ: ô chữ \thuonghieu (trái)  | Trang (số) (phải)
%  Trang chẵn soi gương trang lẻ; chân trang chẵn ghi \tacgia.
%--- Màu: chỉ dùng thang vàng kim Eduzy
\colorlet{ezNen}{gold1}% dải nền đậm
\colorlet{ezVach1}{gold3}% vạch xiên sát dải nền (đậm hơn)
\colorlet{ezVach2}{gold4}% vạch xiên sát ô chữ (nhạt hơn)
\colorlet{ezO}{gold5}% ô chữ vàng kem
%--- Kích thước
\newlength{\EZdai}\setlength{\EZdai}{1.1cm}% chiều cao dải nền
\newlength{\EZo}\setlength{\EZo}{0.68cm}% chiều cao ô chữ
\newlength{\EZxien}\setlength{\EZxien}{0.55cm}% độ lệch ngang của cạnh xiên
\newlength{\EZvach}\setlength{\EZvach}{0.75cm}% bề rộng mỗi vạch xiên
\newlength{\EZkhe}\setlength{\EZkhe}{0.12cm}% khe trắng giữa vạch và ô chữ
\newlength{\EZdem}\setlength{\EZdem}{0.4cm}% khoảng đệm từ chữ tới cạnh xiên
\newlength{\EZle}\setlength{\EZle}{1.8cm}% lề trái/phải của chữ (bằng lề thân bài)
\newlength{\EZrongmin}\setlength{\EZrongmin}{4.5cm}% ô chữ rộng tối thiểu
\newlength{\EZrong}% bề rộng chữ trong ô (tính tự động)
%--- Font
\newcommand{\EZfontO}{\fontsize{10.5pt}{12pt}\selectfont\bfseries}% chữ trong ô
\newcommand{\EZfontKH}{\fontfamily{qag}\fontsize{11pt}{12pt}\selectfont\bfseries}% \khauhieu cạnh logo
\newcommand{\EZfontTrang}{\fontsize{11.5pt}{12pt}\selectfont\bfseries}% chữ "Trang" và số trang
%--- Đo bề rộng chữ trong ô
\newcommand{\EZdorong}[1]{%
	\settowidth{\EZrong}{\EZfontO #1}%
	\ifdim\EZrong<\EZrongmin\setlength{\EZrong}{\EZrongmin}\fi}
%--- Logo + chữ trắng \khauhieu (dựng sẵn một lần khi vào \begin{document})
\newsavebox{\EZlogobox}
\AtBeginDocument{\sbox{\EZlogobox}{%
		\raisebox{-0.5\height}{\eduzylogo[1cm]}\hspace{0.15cm}%
		\raisebox{-0.5\height}{\color{white}\EZfontKH\khauhieu}}}
%--- Hai vạch xiên + khe trắng. #1: điểm gốc ở mép dưới, #2: hướng (1 sang phải, -1 sang trái),
%    #3: hướng nghiêng (1 kiểu "/", -1 kiểu "\")
\newcommand{\EZvachxien}[3]{%
	\fill[white] (#1) --++({(#3)*\EZxien},\EZo) --++({(#2)*\EZkhe},0) --++({-(#3)*\EZxien},-\EZo) --cycle;
	\fill[ezVach2] ([xshift=(#2)*\EZkhe]#1) --++({(#3)*\EZxien},\EZo) --++({(#2)*\EZvach},0) --++({-(#3)*\EZxien},-\EZo) --cycle;
	\fill[ezVach1] ([xshift=(#2)*(\EZkhe+\EZvach)]#1) --++({(#3)*\EZxien},\EZo) --++({(#2)*\EZvach},0) --++({-(#3)*\EZxien},-\EZo) --cycle;
	\draw[white,line width=0.6pt] ([xshift=(#2)*(\EZkhe+\EZvach)]#1) --++({(#3)*\EZxien},\EZo);
}
%--- Chữ trong ô (căn giữa theo chiều cao ô, không lệch vì dấu tiếng Việt)
\tikzset{ezchu/.style={inner sep=0pt,text=ezNen,text height=0.72em,text depth=0.22em}}

%====================== LOGO CHÌM GIỮA TRANG ======================
%  Vẽ cùng đầu trang nên nằm dưới chữ của thân bài, có trên mọi trang.
%  Tắt: \logochimtat (từ trang đó trở đi), bật lại: \logochimbat.
%  Đổi cỡ, độ đậm: khai báo trước \input, ví dụ \def\EZchimcao{9cm} \def\EZchimdo{0.1}
\providecommand{\EZchimcao}{10cm}% chiều cao logo chìm
\providecommand{\EZchimdo}{0.12}% độ đậm (0 = trong suốt, 1 = rõ hẳn)
\newif\ifEZchim \EZchimtrue
\newcommand{\logochimtat}{\EZchimfalse}
\newcommand{\logochimbat}{\EZchimtrue}
%  Dựng sẵn logo mờ một lần khi vào \begin{document}; transparency group làm mờ cả logo
%  một lượt nên không lộ vệt ở chỗ các dải màu chồng nhau.
\newsavebox{\EZchimtho}
\newsavebox{\EZchimbox}
\AtBeginDocument{%
	\sbox{\EZchimtho}{\eduzylogo[\EZchimcao]}%
	\sbox{\EZchimbox}{\begin{tikzpicture}
			\begin{scope}[transparency group,opacity=\EZchimdo]
				\node[inner sep=0pt] {\usebox{\EZchimtho}};
			\end{scope}
	\end{tikzpicture}}}
\newcommand{\EZvechim}{%
	\ifEZchim
		\node[inner sep=0pt] at (current page.center) {\usebox{\EZchimbox}};
	\fi}

%====================== ĐẦU TRANG - TRANG LẺ ======================
%  (logo bên trái, ô chữ \tenchuong bên phải)
\newcommand{\myfancyheadOdd}{%
	\EZdorong{\tenchuong}%
	\begin{tikzpicture}[remember picture,overlay]
		\coordinate (T) at (current page.north west);
		\coordinate (P) at (current page.north east);
		%--- logo chìm giữa trang
		\EZvechim
		%--- dải nền
		\fill[ezNen] (T) rectangle ([yshift=-\EZdai]P);
		%--- ô chữ: X là góc dưới bên trái, cạnh trái xiên kiểu "/"
		\coordinate (X) at ([xshift=-\EZle-\EZrong-\EZdem-\EZxien,yshift=-\EZdai]P);
		\fill[ezO] (X) --++(\EZxien,\EZo) -- ([yshift=-\EZdai+\EZo]P) -- ([yshift=-\EZdai]P) --cycle;
		\EZvachxien{X}{-1}{1}
		\node[ezchu,anchor=east] at ([xshift=-\EZle,yshift=-\EZdai+0.5*\EZo]P) {\EZfontO\tenchuong};
		%--- logo + tên
		\node[inner sep=0pt,anchor=west] at ([xshift=\EZle,yshift=-0.5*\EZdai]T) {\usebox{\EZlogobox}};
	\end{tikzpicture}%
}

%====================== ĐẦU TRANG - TRANG CHẴN =====================
%  (soi gương: ô chữ \tenchuong bên trái, logo bên phải)
\newcommand{\myfancyheadEven}{%
	\EZdorong{\tenchuong}%
	\begin{tikzpicture}[remember picture,overlay]
		\coordinate (T) at (current page.north west);
		\coordinate (P) at (current page.north east);
		%--- logo chìm giữa trang
		\EZvechim
		%--- dải nền
		\fill[ezNen] (T) rectangle ([yshift=-\EZdai]P);
		%--- ô chữ: X là góc dưới bên phải, cạnh phải xiên kiểu "\"
		\coordinate (X) at ([xshift=\EZle+\EZrong+\EZdem+\EZxien,yshift=-\EZdai]T);
		\fill[ezO] (X) --++(-\EZxien,\EZo) -- ([yshift=-\EZdai+\EZo]T) -- ([yshift=-\EZdai]T) --cycle;
		\EZvachxien{X}{1}{-1}
		\node[ezchu,anchor=west] at ([xshift=\EZle,yshift=-\EZdai+0.5*\EZo]T) {\EZfontO\tenchuong};
		%--- logo + tên
		\node[inner sep=0pt,anchor=east] at ([xshift=-\EZle,yshift=-0.5*\EZdai]P) {\usebox{\EZlogobox}};
	\end{tikzpicture}%
}

%====================== CHÂN TRANG - TRANG LẺ ======================
%  (ô chữ \thuonghieu bên trái, "Trang (số)" bên phải)
\newcommand{\myfancyfootOdd}{%
	\EZdorong{\thuonghieu}%
	\begin{tikzpicture}[remember picture,overlay]
		\coordinate (T) at (current page.south west);
		\coordinate (P) at (current page.south east);
		%--- dải nền
		\fill[ezNen] (T) rectangle ([yshift=\EZdai]P);
		%--- ô chữ sát mép trên dải: X là góc dưới bên phải, cạnh phải xiên kiểu "/"
		\coordinate (X) at ([xshift=\EZle+\EZrong+\EZdem,yshift=\EZdai-\EZo]T);
		\fill[ezO] ([yshift=\EZdai-\EZo]T) -- (X) --++(\EZxien,\EZo) -- ([yshift=\EZdai]T) --cycle;
		\EZvachxien{X}{1}{1}
		\node[ezchu,anchor=west] at ([xshift=\EZle,yshift=\EZdai-0.5*\EZo]T) {\EZfontO\thuonghieu};
		%--- số trang
		\node[circle,draw=white,line width=0.9pt,text=white,inner sep=0pt,minimum size=0.62cm,anchor=east]
		(so) at ([xshift=-\EZle,yshift=0.5*\EZdai]P) {\EZfontTrang\thepage};
		\node[inner sep=0pt,text=white,anchor=base east] (tr) at ([xshift=-0.15cm]so.base west) {\EZfontTrang Trang};
	\end{tikzpicture}%
}

%====================== CHÂN TRANG - TRANG CHẴN ====================
%  (soi gương: "Trang (số)" bên trái, ô chữ \tacgia bên phải)
\newcommand{\myfancyfootEven}{%
	\EZdorong{\tacgia}%
	\begin{tikzpicture}[remember picture,overlay]
		\coordinate (T) at (current page.south west);
		\coordinate (P) at (current page.south east);
		%--- dải nền
		\fill[ezNen] (T) rectangle ([yshift=\EZdai]P);
		%--- ô chữ sát mép trên dải: X là góc dưới bên trái, cạnh trái xiên kiểu "\"
		\coordinate (X) at ([xshift=-\EZle-\EZrong-\EZdem,yshift=\EZdai-\EZo]P);
		\fill[ezO] ([yshift=\EZdai-\EZo]P) -- (X) --++(-\EZxien,\EZo) -- ([yshift=\EZdai]P) --cycle;
		\EZvachxien{X}{-1}{-1}
		\node[ezchu,anchor=east] at ([xshift=-\EZle,yshift=\EZdai-0.5*\EZo]P) {\EZfontO\tacgia};
		%--- số trang
		%    (vòng tròn ẩn so0 nằm giữa dải để lấy đường chân chữ, giống trang lẻ)
		\node[circle,inner sep=0pt,minimum size=0.62cm] (so0) at ([yshift=0.5*\EZdai]T) {\EZfontTrang\phantom{\thepage}};
		\node[inner sep=0pt,text=white,anchor=base west] (tr) at ($(so0.base -| T)+(\EZle,0)$) {\EZfontTrang Trang};
		\node[circle,draw=white,line width=0.9pt,text=white,inner sep=0pt,minimum size=0.62cm,anchor=base west]
		(so) at ([xshift=0.15cm]tr.base east) {\EZfontTrang\thepage};
	\end{tikzpicture}%
}

%--- Giữ tương thích với các file cũ đang gọi \myfancyhead, \myfancyfoot
\newcommand{\myfancyhead}{\myfancyheadOdd}
\newcommand{\myfancyfoot}{\myfancyfootOdd}

%======================= GẮN VÀO TRANG =============================
\pagestyle{fancy}
\fancyhf{}
\fancyhead[LO]{\myfancyheadOdd}
\fancyhead[RE]{\myfancyheadEven}
\fancyfoot[LO]{\myfancyfootOdd}
\fancyfoot[RE]{\myfancyfootEven}
\renewcommand{\headrulewidth}{0pt}
\renewcommand{\footrulewidth}{0pt}
\fancypagestyle{plain}{%
	\fancyhf{}
	\fancyhead[LO]{\myfancyheadOdd}
	\fancyhead[RE]{\myfancyheadEven}
	\fancyfoot[LO]{\myfancyfootOdd}
	\fancyfoot[RE]{\myfancyfootEven}
	\renewcommand{\headrulewidth}{0pt}
	\renewcommand{\footrulewidth}{0pt}
}
%=================== HẾT PHẦN GỘP TỪ eduzy-de.tex ===================
\newcommand{\indam}[1]{
	{\fontsize{12pt}{9pt}\sffamily \selectfont\bfseries\color{gold1} #1}
}
\newcommand{\ind}[1]{
	{\fontsize{14pt}{9pt}\fontfamily{qag} \selectfont\bfseries\color{gold1} #1}
}
\renewcommand{\FalseEX}{\stepcounter{dapan}\textbf{\selectfont\color{gold1}\Alph{dapan}.}\ignorespaces}
\renewtheorem{ex}{\color{gold1}Câu}
\renewcommand{\baselinestretch}{1.1}
\newcommand{\hoac}[1]{
	\left[\begin{aligned}#1\end{aligned}\right.}
\newcommand{\heva}[1]{
	\left\{\begin{aligned}#1\end{aligned}\right.}
\usepackage{esvect}
\def\vec{\vv}
\def\vt{\vv}
\def\overrightarrow{\vv}
%====================== LOGO EDUZY VẼ BẰNG TIKZ ======================
%  Đường viền dò từ logo gốc; dải màu vàng kim đo trực tiếp trên ảnh gốc
%  (nghiêng 21 độ, tô bằng 80 dải hẹp nên hiển thị giống nhau trên mọi trình xem PDF).
%  Dùng: \eduzylogo          -> cao 2.4cm (mặc định)
%        \eduzylogo[1.5cm]   -> cao 1.5cm
\definecolor{ezText}{RGB}{245,200,46}
\newcommand{\eduzylogo}[1][2.4cm]{%
	\resizebox{!}{#1}{%
		\begin{tikzpicture}[x=0.01cm,y=0.01cm]
			\path[use as bounding box] (0,0) rectangle (224,-224);
			% mũ tốt nghiệp + vương miện
			\begin{scope}
				\clip
				(119.4,-176.1)..controls(118.9,-176.1)and(116.8,-176)..(114.6,-176)
				..controls(101.4,-176)and(85.8,-175.9)..(83.7,-175.7)
				..controls(81.8,-175.6)and(81,-175.6)..(80.2,-175.7)
				..controls(78.8,-175.9)and(77.7,-175.9)..(76.2,-175.7)
				..controls(74.9,-175.5)and(70.8,-175.5)..(64.3,-175.7)
				..controls(59.1,-175.9)and(58.3,-175.7)..(58.6,-173.7)
				..controls(58.6,-172.8)and(59.6,-169.6)..(60.3,-168)
				..controls(60.8,-166.9)and(61.4,-165.8)..(62.3,-165)
				..controls(63.1,-164.2)and(64.3,-163.4)..(65,-163.2)
				..controls(65.5,-163.1)and(72.3,-163)..(76.6,-163.1)
				..controls(77.9,-163.1)and(80.1,-163.1)..(81.4,-163)
				..controls(82.8,-162.9)and(86.3,-162.9)..(89.2,-162.9)
				..controls(93,-162.9)and(94.9,-162.8)..(96,-162.7)
				..controls(97.2,-162.5)and(101,-162.5)..(114,-162.4)
				..controls(123,-162.3)and(134.9,-162.2)..(140.2,-162.1)
				..controls(150.1,-162)and(152.2,-162.1)..(153.6,-162.4)
				..controls(154.4,-162.6)and(157.1,-163.8)..(158.1,-164.4)
				..controls(158.8,-164.8)and(159.7,-165.7)..(159.9,-166.3)
				..controls(160,-166.5)and(160.2,-167.1)..(160.2,-167.6)
				..controls(160.3,-168)and(160.6,-169.4)..(160.9,-170.6)
				..controls(161.2,-171.7)and(161.5,-173.1)..(161.6,-173.6)
				..controls(161.8,-174.4)and(161.8,-174.5)..(161.5,-174.9)
				..controls(161.2,-175.6)and(160.8,-175.9)..(159.9,-176)
				..controls(159.2,-176.1)and(154.1,-176.2)..(151.5,-176.2)
				..controls(148.1,-176.1)and(122.6,-176.1)..(121.6,-176.2)
				..controls(120.8,-176.2)and(119.8,-176.2)..(119.4,-176.1)--cycle
				(130.4,-160)..controls(129,-160)and(119.8,-160)..(110.1,-159.9)
				..controls(100.3,-159.8)and(88.5,-159.7)..(83.9,-159.6)
				..controls(77.8,-159.5)and(75.3,-159.5)..(74.5,-159.6)
				..controls(73.9,-159.7)and(72.5,-159.8)..(71.3,-159.8)
				..controls(69.6,-159.8)and(69,-159.8)..(68.5,-159.6)
				..controls(67.6,-159.3)and(66.8,-158.6)..(66.2,-157.6)
				..controls(66,-157.3)and(65.6,-156.5)..(65.3,-156.1)
				..controls(64.6,-155.1)and(64.5,-154.9)..(63.7,-152.9)
				..controls(63.2,-151.6)and(62.9,-150.9)..(61.8,-149.6)
				..controls(61.5,-149.2)and(61.1,-148.5)..(60.9,-147.9)
				..controls(60.7,-147.4)and(60.3,-146.6)..(60.1,-146.1)
				..controls(59.9,-145.5)and(59.5,-144.9)..(59,-144.3)
				..controls(58.5,-143.7)and(58.2,-143.2)..(57.9,-142.5)
				..controls(57.7,-141.9)and(57.3,-141.1)..(57,-140.7)
				..controls(56.6,-140.3)and(56.1,-139.3)..(55.8,-138.5)
				..controls(55.5,-137.8)and(55.1,-136.8)..(54.9,-136.5)
				..controls(54.7,-136.1)and(54.3,-135.2)..(53.9,-134.5)
				..controls(53.6,-133.8)and(53,-132.7)..(52.6,-132)
				..controls(51.6,-130.5)and(51.4,-130.2)..(51.1,-129)
				..controls(50.6,-127.5)and(50.6,-127.4)..(49.7,-126.3)
				..controls(49.2,-125.7)and(48.6,-124.8)..(48.3,-124.2)
				..controls(48,-123.6)and(47.5,-122.7)..(47.3,-122.2)
				..controls(47,-121.8)and(46.5,-120.7)..(46.1,-119.8)
				..controls(45.8,-118.9)and(45.2,-117.8)..(44.9,-117.2)
				..controls(44.2,-116.1)and(43.3,-114.4)..(42.3,-112.2)
				..controls(41.8,-111.3)and(41.3,-110.4)..(41.1,-110.1)
				..controls(40.5,-109.4)and(40.3,-109)..(40,-107.8)
				..controls(39.8,-107)and(39.6,-106.5)..(39.2,-106)
				..controls(38.9,-105.6)and(38.4,-104.7)..(38,-104)
				..controls(37.7,-103.4)and(37.2,-102.5)..(37,-102.1)
				..controls(36.7,-101.7)and(36.3,-100.8)..(36,-100.1)
				..controls(35.6,-98.9)and(35.3,-98.4)..(34.6,-97.6)
				..controls(34.2,-97.1)and(34.2,-96.3)..(34.6,-96)
				..controls(35,-95.7)and(35.5,-95.9)..(36.2,-96.7)
				..controls(36.5,-97)and(37.2,-97.5)..(37.6,-97.7)..controls(38,-98)and(38.6,-98.4)..(39.1,-98.9)
				..controls(39.6,-99.4)and(40.3,-100)..(40.6,-100.3)
				..controls(41.3,-100.8)and(42.4,-101.8)..(43.7,-102.9)
				..controls(44,-103.2)and(44.6,-103.7)..(45,-103.9)
				..controls(45.4,-104.2)and(45.9,-104.6)..(46.1,-104.8)
				..controls(47.1,-105.9)and(48.8,-107.1)..(50.5,-107.9)
				..controls(51.8,-108.6)and(52.3,-109)..(53.4,-110)
				..controls(54.8,-111.2)and(55.6,-111.6)..(57,-112.1)
				..controls(57.5,-112.3)and(58.5,-112.8)..(59.4,-113.3)
				..controls(60.2,-113.7)and(61.3,-114.3)..(61.9,-114.5)
				..controls(62.4,-114.7)and(63.3,-115.1)..(63.8,-115.3)
				..controls(64.4,-115.6)and(65.4,-115.9)..(66.2,-116.1)
				..controls(67,-116.3)and(68.4,-116.6)..(69.2,-116.9)
				..controls(72.3,-117.9)and(72.8,-118)..(78,-118.1)
				..controls(84,-118.3)and(86.1,-118.1)..(88.1,-117.1)
				..controls(89.4,-116.5)and(93,-115)..(93.7,-114.9)
				..controls(94.8,-114.6)and(95.4,-114.2)..(95.9,-113.4)
				..controls(96.2,-113)and(96.7,-112.5)..(97.1,-112.2)
				..controls(98,-111.6)and(101.3,-108.4)..(102.1,-107.2)
				..controls(102.4,-106.8)and(102.9,-106.2)..(103.2,-105.8)
				..controls(103.5,-105.5)and(104,-104.8)..(104.2,-104.2)
				..controls(104.5,-103.6)and(105,-102.8)..(105.3,-102.4)
				..controls(105.6,-102)and(106.1,-101.2)..(106.5,-100.6)
				..controls(107.1,-99.5)and(108.6,-97.5)..(109.5,-96.5)
				..controls(109.9,-96.1)and(110.1,-96)..(110.4,-96)
				..controls(111.1,-96)and(111.6,-96.5)..(111.8,-97.6)
				..controls(112,-98.6)and(112.3,-99.3)..(112.8,-99.9)
				..controls(113.7,-100.9)and(115.2,-103.1)..(115.7,-104.2)
				..controls(116.1,-105)and(116.5,-105.5)..(116.9,-106)
				..controls(117.3,-106.4)and(117.8,-106.9)..(118,-107.3)
				..controls(118.2,-107.6)and(118.6,-108)..(119,-108.3)
				..controls(119.6,-108.8)and(121.1,-110.3)..(121.8,-111.2)
				..controls(122.1,-111.7)and(122.5,-112)..(123.4,-112.4)
				..controls(124.8,-113.2)and(125.8,-113.8)..(127,-114.6)
				..controls(127.5,-114.9)and(128.1,-115.3)..(128.2,-115.4)
				..controls(128.4,-115.5)and(129.3,-115.8)..(130.2,-116)
				..controls(131.1,-116.2)and(132.5,-116.6)..(133.3,-116.9)
				..controls(135.7,-117.9)and(136.1,-117.9)..(142.5,-117.8)
				..controls(145.8,-117.8)and(148.4,-117.7)..(148.8,-117.6)
				..controls(149.2,-117.6)and(150.2,-117.2)..(151,-116.8)
				..controls(152,-116.4)and(152.8,-116)..(153.6,-115.9)
				..controls(154.5,-115.7)and(155.2,-115.4)..(156.3,-114.9)
				..controls(157.1,-114.5)and(158.2,-114)..(158.7,-113.8)
				..controls(159.1,-113.7)and(159.8,-113.3)..(160.2,-113)
				..controls(160.6,-112.7)and(161.3,-112.4)..(161.7,-112.2)
				..controls(162.7,-111.9)and(163.2,-111.7)..(164.1,-111.1)
				..controls(164.5,-110.8)and(165.3,-110.4)..(165.9,-110)
				..controls(166.5,-109.7)and(167.3,-109.2)..(167.6,-108.9)
				..controls(167.9,-108.6)and(168.6,-108.2)..(169,-108)
				..controls(169.4,-107.8)and(170.2,-107.3)..(170.7,-106.9)
				..controls(171.2,-106.5)and(172.1,-105.8)..(172.7,-105.4)
				..controls(173.2,-105)and(174.1,-104.4)..(174.5,-104)
				..controls(174.9,-103.7)and(175.7,-103.1)..(176.2,-102.6)
				..controls(176.7,-102.2)and(177.4,-101.5)..(177.9,-101.2)
				..controls(178.3,-100.8)and(178.8,-100.4)..(179,-100.2)
				..controls(179.2,-100)and(179.6,-99.6)..(180,-99.3)
				..controls(180.4,-99.1)and(180.9,-98.6)..(181.2,-98.4)
				..controls(181.5,-98.1)and(182.1,-97.7)..(182.5,-97.5)
				..controls(183.2,-97.2)and(183.5,-96.9)..(184.3,-95.9)
				..controls(184.6,-95.6)and(184.6,-95.6)..(184.8,-95.7)
				..controls(185,-95.8)and(185.1,-96)..(185.1,-96.2)
				..controls(185.1,-96.5)and(185,-96.7)..(184.6,-97.2)
				..controls(184.1,-97.8)and(184,-97.9)..(183.9,-98.8)
				..controls(183.7,-100.3)and(183.3,-101.3)..(182.9,-101.8)
				..controls(182.6,-102.1)and(182.3,-102.6)..(182.2,-102.9)
				..controls(182,-103.3)and(181.6,-104)..(181.2,-104.6)
				..controls(180.9,-105.2)and(180.5,-106.1)..(180.4,-106.6)
				..controls(179.9,-108)and(179.7,-108.5)..(179,-109.4)
				..controls(178.6,-109.8)and(178.2,-110.5)..(177.9,-111)
				..controls(177.7,-111.5)and(177.3,-112.3)..(177,-112.7)
				..controls(176.6,-113.2)and(176.4,-113.7)..(176.3,-114.2)
				..controls(176.2,-114.7)and(175.9,-115.2)..(175.4,-116)
				..controls(175,-116.6)and(174.4,-117.7)..(174,-118.5)
				..controls(173.7,-119.2)and(173.3,-120.1)..(173,-120.5)
				..controls(172.8,-120.9)and(172.4,-121.8)..(172.1,-122.5)
				..controls(171.8,-123.2)and(171.3,-124.1)..(171,-124.6)
				..controls(170.8,-125.1)and(170.3,-126)..(169.9,-126.6)
				..controls(169.6,-127.3)and(169.2,-128.1)..(169,-128.3)
				..controls(168.8,-128.6)and(168.5,-129.3)..(168.3,-129.8)
				..controls(168.1,-130.3)and(167.8,-131.1)..(167.5,-131.6)
				..controls(166.9,-132.5)and(164.2,-137.9)..(164.2,-138.1)
				..controls(164.2,-138.1)and(163.9,-138.8)..(163.4,-139.7)
				..controls(162.9,-140.5)and(162.4,-141.5)..(162.3,-141.9)
				..controls(161.9,-143)and(161.6,-143.4)..(160.9,-144.1)
				..controls(160.4,-144.6)and(160.3,-144.8)..(160.1,-145.4)
				..controls(160,-145.8)and(159.6,-146.8)..(159.2,-147.8)
				..controls(158.7,-148.8)and(158.1,-150)..(157.9,-150.6)
				..controls(157.6,-151.1)and(157.2,-151.9)..(156.9,-152.4)
				..controls(156.5,-152.8)and(156.1,-153.7)..(155.8,-154.4)
				..controls(155.6,-155)and(155.2,-155.9)..(154.9,-156.3)
				..controls(154.7,-156.7)and(154.4,-157.4)..(154.2,-157.9)
				..controls(154.1,-158.4)and(153.8,-159)..(153.5,-159.3)
				..controls(152.8,-160.1)and(153,-160.1)..(142.4,-160.1)
				..controls(137.2,-160.1)and(131.8,-160.1)..(130.4,-160)--cycle
				(141.5,-103.7)..controls(140.8,-103)and(139.8,-102.1)..(139.2,-101.6)
				..controls(138.5,-101.2)and(137.9,-100.6)..(137.8,-100.4)
				..controls(137,-99.5)and(134,-96.5)..(133.2,-96)
				..controls(132.7,-95.7)and(131.8,-94.9)..(131.1,-94.3)
				..controls(129.4,-92.7)and(127.9,-91.6)..(126.6,-91.1)
				..controls(125.2,-90.5)and(124.1,-90)..(122.9,-89.1)
				..controls(122.1,-88.6)and(121.3,-88.2)..(119.9,-87.7)
				..controls(118.8,-87.3)and(117.5,-86.8)..(116.9,-86.5)
				..controls(114.3,-85.5)and(113.7,-85.4)..(109.7,-85.5)
				..controls(105.3,-85.7)and(103.9,-85.9)..(101.6,-87)
				..controls(100.7,-87.4)and(99.4,-87.8)..(97.9,-88)
				..controls(97,-88.2)and(96.6,-88.3)..(95.8,-88.9)
				..controls(95.4,-89.2)and(94.7,-89.6)..(94.2,-89.8)
				..controls(92.9,-90.3)and(92.3,-90.7)..(90.9,-92.1)
				..controls(90.2,-92.7)and(89.4,-93.3)..(89,-93.5)..controls(87,-94.5)and(85.2,-96)..(83.6,-97.9)
				..controls(83.3,-98.2)and(82.5,-98.9)..(81.8,-99.4)
				..controls(80.7,-100.1)and(80.5,-100.3)..(80,-101.1)
				..controls(78.7,-103.1)and(77.2,-103.9)..(76.4,-103)
				..controls(75.7,-102.2)and(75.7,-103.2)..(75.7,-83.6)
				..controls(75.6,-73.6)and(75.6,-64.8)..(75.5,-64)
				..controls(75.5,-63.2)and(75.5,-61.4)..(75.6,-59.8)--(75.7,-57.1)--(76,-56.8)
				..controls(76.8,-56.2)and(78,-56.3)..(79.8,-57.2)
				..controls(80.4,-57.5)and(81.3,-57.8)..(81.8,-57.9)
				..controls(82.3,-58.1)and(83,-58.4)..(83.6,-58.8)..controls(84.3,-59.2)and(85,-59.5)..(86,-59.8)
				..controls(86.8,-60.1)and(88,-60.6)..(88.7,-60.9)
				..controls(89.3,-61.2)and(90.6,-61.7)..(91.4,-62)
				..controls(92.2,-62.3)and(93.7,-62.9)..(94.6,-63.3)
				..controls(95.6,-63.7)and(96.8,-64.1)..(97.4,-64.3)
				..controls(98,-64.5)and(99,-64.9)..(99.7,-65.2)
				..controls(100.4,-65.6)and(101.5,-66)..(102.3,-66.2)
				..controls(104.2,-66.8)and(105,-67.1)..(106.2,-67.8)
				..controls(107.9,-68.7)and(110.2,-69.4)..(111.3,-69.2)
				..controls(112.5,-69)and(115,-68.2)..(115.6,-67.8)
				..controls(117.3,-66.7)and(117.6,-66.6)..(119.5,-66.2)
				..controls(121.9,-65.8)and(122.1,-65.7)..(123,-65.1)
				..controls(124.1,-64.4)and(124.4,-64.3)..(126,-63.9)
				..controls(127,-63.7)and(127.7,-63.4)..(128.4,-63.1)
				..controls(129,-62.8)and(129.7,-62.5)..(130.6,-62.4)
				..controls(132.1,-62.1)and(134,-61.5)..(135.2,-61)
				..controls(135.7,-60.8)and(136.7,-60.5)..(137.3,-60.3)
				..controls(138.4,-60)and(139.3,-59.7)..(141.1,-59)
				..controls(141.6,-58.7)and(142,-58.6)..(142.5,-58.6)
				..controls(143.2,-58.6)and(143.3,-58.7)..(143.6,-59)
				..controls(144,-59.5)and(144.1,-59.9)..(144,-60.8)
				..controls(144,-61.2)and(144,-71.1)..(144,-82.8)--(143.9,-104.1)--(143.7,-104.5)
				..controls(143.5,-104.7)and(143.3,-104.8)..(143,-104.9)
				..controls(142.6,-104.9)and(142.5,-104.8)..(141.5,-103.7)--cycle
				(27.8,-91.4)..controls(27.1,-90.9)and(26.3,-90)..(26.1,-89.5)
				..controls(26,-89.2)and(25.8,-88.6)..(25.7,-87.9)
				..controls(25.4,-86.2)and(25.6,-84.8)..(26.5,-83.5)
				..controls(27.4,-82.3)and(28.3,-82)..(29.4,-82.3)
				..controls(31.1,-82.8)and(32.1,-84.9)..(32,-87.4)
				..controls(31.9,-89.3)and(31.2,-90.6)..(29.9,-91.4)
				..controls(29.1,-91.9)and(28.6,-91.9)..(27.8,-91.4)--cycle
				(190.1,-91.3)..controls(189.7,-91.2)and(189.3,-90.9)..(189.1,-90.7)
				..controls(188.3,-89.8)and(187.6,-87.7)..(187.6,-86.2)
				..controls(187.6,-85.5)and(187.7,-85.2)..(188.2,-84.4)
				..controls(189.1,-82.6)and(190.1,-81.8)..(191.3,-82.1)
				..controls(192,-82.3)and(193.5,-83.7)..(193.8,-84.4)
				..controls(194.1,-85.1)and(194.2,-87.1)..(194,-88.3)
				..controls(193.8,-89.8)and(193.3,-90.5)..(192.3,-91.3)
				..controls(191.7,-91.7)and(190.9,-91.7)..(190.1,-91.3)--cycle
				(173.7,-83.4)..controls(172.2,-82.5)and(172,-82.5)..(170.6,-82.4)
				..controls(169.7,-82.4)and(168.8,-82.4)..(168.3,-82.5)
				..controls(167.2,-82.7)and(166.9,-82.7)..(166.6,-82.2)
				..controls(166.3,-81.8)and(166.3,-81)..(166.8,-80)
				..controls(167.2,-78.9)and(167.4,-77.4)..(167.7,-74.1)
				..controls(167.8,-72.8)and(167.9,-71.8)..(167.9,-71.7)
				..controls(168,-71.7)and(168,-70.2)..(168,-68.4)..controls(168,-65.4)and(168,-64.9)..(167.7,-63)
				..controls(167.3,-60.5)and(167.1,-59.9)..(166.4,-58.4)
				..controls(166.1,-57.7)and(165.8,-56.9)..(165.7,-56.6)
				..controls(165.7,-56.2)and(165.4,-55.6)..(165.2,-55.2)
				..controls(164.6,-53.8)and(161.9,-50.7)..(161,-50.3)
				..controls(159.8,-49.7)and(157.9,-50)..(155.2,-51.2)
				..controls(154.7,-51.4)and(153.4,-51.8)..(152.2,-52.1)
				..controls(150.8,-52.4)and(149.4,-52.9)..(148.6,-53.2)
				..controls(147.8,-53.6)and(146.7,-54)..(146,-54.2)
				..controls(145.4,-54.4)and(144.5,-54.7)..(144,-54.9)
				..controls(143.6,-55.2)and(142.7,-55.5)..(141.8,-55.7)
				..controls(141,-55.9)and(139.6,-56.3)..(138.7,-56.6)
				..controls(137.8,-57)and(136.3,-57.4)..(135.3,-57.6)
				..controls(133.8,-57.9)and(133.2,-58.1)..(132.1,-58.7)
				..controls(131.1,-59.2)and(130.3,-59.4)..(129.4,-59.6)
				..controls(128.7,-59.8)and(127.8,-60)..(127.4,-60.1)
				..controls(127,-60.3)and(126,-60.5)..(125.1,-60.7)
				..controls(123.2,-61.1)and(121.7,-61.5)..(120.2,-62.1)
				..controls(119.7,-62.3)and(118.1,-62.8)..(116.7,-63.2)
				..controls(114.6,-63.9)and(114,-64.1)..(113,-64.7)
				..controls(111.9,-65.4)and(111.8,-65.4)..(111,-65.5)
				..controls(109.9,-65.5)and(109.3,-65.4)..(107.7,-64.7)
				..controls(107,-64.4)and(106,-64)..(105.3,-63.9)
				..controls(103.5,-63.5)and(102.2,-63)..(101.2,-62.5)
				..controls(100.1,-61.8)and(99.3,-61.5)..(97.8,-61.1)
				..controls(97.2,-61)and(96.2,-60.7)..(95.7,-60.4)
				..controls(95.1,-60.1)and(94.2,-59.8)..(93.5,-59.6)
				..controls(92.9,-59.4)and(91.7,-59)..(91,-58.7)..controls(90.2,-58.4)and(88.9,-58)..(88,-57.7)
				..controls(87.2,-57.5)and(86.1,-57.1)..(85.5,-56.8)
				..controls(85,-56.6)and(84.1,-56.2)..(83.5,-56)..controls(82.9,-55.7)and(81.8,-55.2)..(81,-54.9)
				..controls(80.1,-54.5)and(79.1,-54.2)..(78.7,-54.1)
				..controls(77.9,-53.8)and(77.3,-53.6)..(76.2,-52.9)
				..controls(75.7,-52.6)and(75.2,-52.4)..(74.1,-52.2)
				..controls(72.5,-51.9)and(72.4,-51.8)..(71.3,-51.1)
				..controls(69.9,-50.1)and(69.4,-49.9)..(67.9,-49.8)
				..controls(67,-49.7)and(66.6,-49.6)..(66.3,-49.4)
				..controls(66.1,-49.3)and(65.7,-49)..(65.4,-48.8)
				..controls(65.1,-48.7)and(64.2,-48.4)..(63.6,-48.2)
				..controls(62.9,-48)and(61.7,-47.6)..(60.9,-47.2)
				..controls(60.1,-46.8)and(58.9,-46.4)..(58.2,-46.2)
				..controls(57.6,-46)and(56.8,-45.7)..(56.5,-45.5)
				..controls(54.9,-44.6)and(54.6,-44.5)..(53.2,-44.2)
				..controls(51.7,-43.8)and(51,-43.6)..(49.9,-42.9)
				..controls(49.5,-42.6)and(48.9,-42.4)..(47.5,-42)
				..controls(45.7,-41.6)and(45.6,-41.5)..(44.3,-40.6)
				..controls(43.5,-40.1)and(42.8,-39.6)..(42.7,-39.4)
				..controls(42.4,-38.7)and(43,-38.1)..(44.5,-37.6)..controls(44.9,-37.5)and(46,-37)..(46.9,-36.6)
				..controls(47.9,-36.1)and(48.8,-35.8)..(49.6,-35.6)
				..controls(50.5,-35.4)and(51.3,-35.2)..(52.4,-34.6)
				..controls(53.3,-34.1)and(54.5,-33.7)..(54.9,-33.5)
				..controls(55.4,-33.4)and(56.4,-33)..(57.2,-32.6)
				..controls(58.6,-32)and(60.1,-31.5)..(61.6,-31.2)
				..controls(62.3,-31.1)and(62.7,-31)..(63.2,-30.7)
				..controls(64.1,-30)and(64.3,-29.9)..(65.3,-29.7)
				..controls(65.7,-29.6)and(66.5,-29.2)..(67.1,-28.9)
				..controls(67.6,-28.6)and(68.8,-28.2)..(69.6,-27.9)
				..controls(70.5,-27.7)and(71.6,-27.3)..(72.1,-27)
				..controls(72.6,-26.7)and(73.8,-26.2)..(74.8,-25.9)
				..controls(75.7,-25.5)and(76.7,-25.1)..(77.1,-25)
				..controls(77.4,-24.8)and(78.5,-24.4)..(79.6,-24.1)
				..controls(80.6,-23.8)and(82,-23.2)..(82.8,-22.9)
				..controls(83.5,-22.5)and(84.8,-22)..(85.8,-21.8)
				..controls(86.7,-21.5)and(87.8,-21.1)..(88.1,-20.9)
				..controls(88.4,-20.7)and(89.3,-20.3)..(90.1,-20.1)
				..controls(90.8,-19.8)and(91.9,-19.4)..(92.4,-19.1)
				..controls(93.5,-18.5)and(94.4,-18.2)..(95.4,-18.1)
				..controls(95.7,-18)and(96.4,-17.8)..(96.9,-17.5)
				..controls(97.5,-17.3)and(98.4,-16.9)..(99,-16.6)
				..controls(99.7,-16.3)and(100.9,-15.9)..(101.9,-15.6)
				..controls(102.9,-15.3)and(104.4,-14.7)..(105.3,-14.3)
				..controls(108.1,-13.2)and(108.5,-13.1)..(110.1,-13.1)
				..controls(111.4,-13.1)and(111.5,-13.1)..(112.4,-13.5)
				..controls(113.2,-14)and(115,-14.8)..(116.6,-15.4)
				..controls(116.9,-15.5)and(117.9,-15.8)..(118.8,-16.1)
				..controls(120.9,-16.6)and(121.5,-16.9)..(122.6,-17.9)
				..controls(123,-18.3)and(123.7,-18.8)..(123.9,-18.9)
				..controls(124.6,-19.4)and(126,-19.8)..(127.9,-20.2)
				..controls(129.6,-20.5)and(129.8,-20.6)..(131.1,-21.4)
				..controls(131.5,-21.7)and(132.4,-22.1)..(133.1,-22.3)
				..controls(136.1,-23.2)and(138.3,-24.1)..(140.2,-25)
				..controls(141,-25.5)and(141.5,-25.6)..(142.5,-25.8)
				..controls(145.1,-26.3)and(146,-26.6)..(149.4,-28.1)
				..controls(151.3,-28.9)and(153.5,-29.8)..(154.3,-30.1)
				..controls(155.1,-30.4)and(156.2,-30.8)..(156.6,-31)
				..controls(157.1,-31.2)and(158.2,-31.6)..(159.1,-31.8)
				..controls(160.6,-32.2)and(160.9,-32.3)..(161.7,-32.9)
				..controls(162.4,-33.3)and(162.9,-33.6)..(163.7,-33.8)
				..controls(164.3,-34)and(164.8,-34.1)..(164.9,-34.1)
				..controls(165,-34.1)and(165.6,-34.4)..(166.4,-34.8)
				..controls(167.3,-35.2)and(168.1,-35.4)..(169,-35.6)
				..controls(170,-35.8)and(170.5,-36)..(171.7,-36.6)
				..controls(172.4,-37)and(173.3,-37.4)..(173.6,-37.4)
				..controls(175.6,-37.7)and(176.2,-38)..(176.7,-39.1)
				..controls(176.9,-39.6)and(176.9,-39.7)..(176.7,-40.4)
				..controls(176.4,-41.6)and(175.8,-42.7)..(175.2,-43.4)
				..controls(174.5,-44.1)and(172.6,-45.5)..(171.7,-45.8)
				..controls(170.6,-46.2)and(169.1,-46)..(168.2,-45.4)
				..controls(167.9,-45.2)and(167.3,-44.8)..(166.9,-44.4)
				..controls(166.4,-44.1)and(165.8,-43.6)..(165.4,-43.2)
				..controls(165.1,-42.9)and(164.4,-42.3)..(163.9,-41.9)
				..controls(163.3,-41.5)and(162.7,-40.9)..(162.4,-40.5)
				..controls(161.4,-39.2)and(159.9,-38)..(158.8,-37.8)
				..controls(158.4,-37.8)and(158.3,-37.7)..(158.2,-37.2)
				..controls(158,-36.8)and(157.4,-36.1)..(157.1,-36.1)
				..controls(157,-36.1)and(156.7,-36.1)..(156.5,-36)
				..controls(156.1,-35.9)and(155.9,-35.7)..(155.7,-35.2)
				..controls(155.3,-34.6)and(155.1,-34.4)..(154,-34.1)
				..controls(153.6,-34)and(153.2,-33.7)..(152.9,-33.5)
				..controls(152.4,-32.9)and(151.6,-32.4)..(151.3,-32.4)
				..controls(150.9,-32.4)and(150.4,-32.6)..(150.4,-32.9)
				..controls(150.4,-33.2)and(150.8,-33.4)..(151.4,-33.6)
				..controls(152.1,-33.7)and(152.8,-34.1)..(153.5,-34.9)
				..controls(154.1,-35.5)and(154.3,-35.6)..(155,-35.8)
				..controls(155.7,-36)and(155.7,-36)..(156,-36.6)
				..controls(156.3,-37.3)and(156.7,-37.7)..(157.5,-37.8)
				..controls(158.1,-37.9)and(158.2,-38)..(158.4,-38.8)
				..controls(158.5,-39.4)and(158.6,-39.6)..(159.6,-40.4)
				..controls(160.2,-41)and(160.9,-41.6)..(161.3,-41.8)
				..controls(161.7,-42)and(161.9,-42.3)..(162.2,-42.8)
				..controls(162.4,-43.2)and(162.8,-43.8)..(163.1,-44.1)
				..controls(163.9,-44.8)and(165.1,-46.7)..(165.7,-48)
				..controls(166.2,-49.3)and(166.4,-49.7)..(167.2,-50.5)
				..controls(167.5,-50.8)and(167.9,-51.4)..(168.1,-51.8)
				..controls(168.3,-52.2)and(168.8,-53)..(169.1,-53.6)
				..controls(169.6,-54.3)and(169.8,-54.9)..(169.9,-55.6)
				..controls(170.4,-57.8)and(170.5,-58.1)..(171.2,-59.5)
				..controls(171.6,-60.3)and(172,-61.5)..(172.2,-62.1)
				..controls(172.3,-62.8)and(172.7,-63.7)..(172.9,-64.1)
				..controls(173.2,-64.6)and(173.4,-65.3)..(173.4,-65.6)
				..controls(173.8,-68.8)and(174.2,-70.6)..(174.9,-72.6)
				..controls(175.5,-74.5)and(175.9,-77.1)..(176.1,-80.2)
				..controls(176.3,-83.4)and(175.5,-84.4)..(173.7,-83.4)--cycle;
				\definecolor{ezs}{RGB}{247,179,26}\fill[ezs](-5.7,-156.6)--(-3.4,-157.5)--(61.9,12.7)--(59.7,13.6)--cycle;
				\definecolor{ezs}{RGB}{247,179,22}\fill[ezs](-3.8,-157.3)--(-1.5,-158.2)--(63.9,12)--(61.6,12.9)--cycle;
				\definecolor{ezs}{RGB}{248,179,19}\fill[ezs](-1.8,-158.1)--(0.4,-158.9)--(65.8,11.2)--(63.5,12.1)--cycle;
				\definecolor{ezs}{RGB}{248,178,16}\fill[ezs](0.1,-158.8)--(2.4,-159.7)--(67.7,10.5)--(65.4,11.4)--cycle;
				\definecolor{ezs}{RGB}{249,178,15}\fill[ezs](2,-159.5)--(4.3,-160.4)--(69.6,9.8)--(67.3,10.7)--cycle;
				\definecolor{ezs}{RGB}{251,177,17}\fill[ezs](3.9,-160.3)--(6.2,-161.1)--(71.5,9)--(69.2,9.9)--cycle;
				\definecolor{ezs}{RGB}{252,179,18}\fill[ezs](5.8,-161)--(8.1,-161.9)--(73.4,8.3)--(71.1,9.2)--cycle;
				\definecolor{ezs}{RGB}{252,182,17}\fill[ezs](7.7,-161.7)--(10,-162.6)--(75.3,7.6)--(73,8.5)--cycle;
				\definecolor{ezs}{RGB}{252,185,17}\fill[ezs](9.6,-162.5)--(11.9,-163.4)--(77.2,6.8)--(75,7.7)--cycle;
				\definecolor{ezs}{RGB}{251,189,18}\fill[ezs](11.5,-163.2)--(13.8,-164.1)--(79.2,6.1)--(76.9,7)--cycle;
				\definecolor{ezs}{RGB}{251,191,22}\fill[ezs](13.4,-163.9)--(15.7,-164.8)--(81.1,5.4)--(78.8,6.2)--cycle;
				\definecolor{ezs}{RGB}{252,193,25}\fill[ezs](15.4,-164.7)--(17.6,-165.6)--(83,4.6)--(80.7,5.5)--cycle;
				\definecolor{ezs}{RGB}{253,195,27}\fill[ezs](17.3,-165.4)--(19.6,-166.3)--(84.9,3.9)--(82.6,4.8)--cycle;
				\definecolor{ezs}{RGB}{253,196,31}\fill[ezs](19.2,-166.1)--(21.5,-167)--(86.8,3.2)--(84.5,4)--cycle;
				\definecolor{ezs}{RGB}{254,199,35}\fill[ezs](21.1,-166.9)--(23.4,-167.8)--(88.7,2.4)--(86.4,3.3)--cycle;
				\definecolor{ezs}{RGB}{253,201,38}\fill[ezs](23,-167.6)--(25.3,-168.5)--(90.6,1.7)--(88.3,2.6)--cycle;
				\definecolor{ezs}{RGB}{253,203,40}\fill[ezs](24.9,-168.3)--(27.2,-169.2)--(92.5,1)--(90.3,1.8)--cycle;
				\definecolor{ezs}{RGB}{253,205,42}\fill[ezs](26.8,-169.1)--(29.1,-170)--(94.4,0.2)--(92.2,1.1)--cycle;
				\definecolor{ezs}{RGB}{253,207,45}\fill[ezs](28.7,-169.8)--(31,-170.7)--(96.4,-0.5)--(94.1,0.4)--cycle;
				\definecolor{ezs}{RGB}{253,210,49}\fill[ezs](30.7,-170.5)--(32.9,-171.4)--(98.3,-1.2)--(96,-0.4)--cycle;
				\definecolor{ezs}{RGB}{253,212,54}\fill[ezs](32.6,-171.3)--(34.9,-172.2)--(100.2,-2)--(97.9,-1.1)--cycle;
				\definecolor{ezs}{RGB}{253,214,59}\fill[ezs](34.5,-172)--(36.8,-172.9)--(102.1,-2.7)--(99.8,-1.8)--cycle;
				\definecolor{ezs}{RGB}{253,216,63}\fill[ezs](36.4,-172.7)--(38.7,-173.6)--(104,-3.4)--(101.7,-2.6)--cycle;
				\definecolor{ezs}{RGB}{253,218,65}\fill[ezs](38.3,-173.5)--(40.6,-174.4)--(105.9,-4.2)--(103.6,-3.3)--cycle;
				\definecolor{ezs}{RGB}{254,220,66}\fill[ezs](40.2,-174.2)--(42.5,-175.1)--(107.8,-4.9)--(105.5,-4)--cycle;
				\definecolor{ezs}{RGB}{254,222,68}\fill[ezs](42.1,-175)--(44.4,-175.8)--(109.7,-5.6)--(107.5,-4.8)--cycle;
				\definecolor{ezs}{RGB}{254,224,71}\fill[ezs](44,-175.7)--(46.3,-176.6)--(111.7,-6.4)--(109.4,-5.5)--cycle;
				\definecolor{ezs}{RGB}{254,226,75}\fill[ezs](46,-176.4)--(48.2,-177.3)--(113.6,-7.1)--(111.3,-6.2)--cycle;
				\definecolor{ezs}{RGB}{254,228,79}\fill[ezs](47.9,-177.2)--(50.1,-178)--(115.5,-7.8)--(113.2,-7)--cycle;
				\definecolor{ezs}{RGB}{254,230,82}\fill[ezs](49.8,-177.9)--(52.1,-178.8)--(117.4,-8.6)--(115.1,-7.7)--cycle;
				\definecolor{ezs}{RGB}{254,232,85}\fill[ezs](51.7,-178.6)--(54,-179.5)--(119.3,-9.3)--(117,-8.4)--cycle;
				\definecolor{ezs}{RGB}{254,234,88}\fill[ezs](53.6,-179.4)--(55.9,-180.2)--(121.2,-10)--(118.9,-9.2)--cycle;
				\definecolor{ezs}{RGB}{254,237,91}\fill[ezs](55.5,-180.1)--(57.8,-181)--(123.1,-10.8)--(120.8,-9.9)--cycle;
				\definecolor{ezs}{RGB}{254,239,94}\fill[ezs](57.4,-180.8)--(59.7,-181.7)--(125,-11.5)--(122.8,-10.6)--cycle;
				\definecolor{ezs}{RGB}{254,241,97}\fill[ezs](59.3,-181.6)--(61.6,-182.4)--(127,-12.2)--(124.7,-11.4)--cycle;
				\definecolor{ezs}{RGB}{254,243,101}\fill[ezs](61.2,-182.3)--(63.5,-183.2)--(128.9,-13)--(126.6,-12.1)--cycle;
				\definecolor{ezs}{RGB}{254,245,103}\fill[ezs](63.2,-183)--(65.4,-183.9)--(130.8,-13.7)--(128.5,-12.8)--cycle;
				\definecolor{ezs}{RGB}{254,245,105}\fill[ezs](65.1,-183.8)--(67.4,-184.6)--(132.7,-14.4)--(130.4,-13.6)--cycle;
				\definecolor{ezs}{RGB}{254,245,105}\fill[ezs](67,-184.5)--(69.3,-185.4)--(134.6,-15.2)--(132.3,-14.3)--cycle;
				\definecolor{ezs}{RGB}{254,244,103}\fill[ezs](68.9,-185.2)--(71.2,-186.1)--(136.5,-15.9)--(134.2,-15)--cycle;
				\definecolor{ezs}{RGB}{253,243,101}\fill[ezs](70.8,-186)--(73.1,-186.8)--(138.4,-16.6)--(136.1,-15.8)--cycle;
				\definecolor{ezs}{RGB}{253,241,98}\fill[ezs](72.7,-186.7)--(75,-187.6)--(140.3,-17.4)--(138,-16.5)--cycle;
				\definecolor{ezs}{RGB}{253,239,95}\fill[ezs](74.6,-187.4)--(76.9,-188.3)--(142.2,-18.1)--(140,-17.2)--cycle;
				\definecolor{ezs}{RGB}{253,236,92}\fill[ezs](76.5,-188.2)--(78.8,-189)--(144.2,-18.8)--(141.9,-18)--cycle;
				\definecolor{ezs}{RGB}{253,234,88}\fill[ezs](78.5,-188.9)--(80.7,-189.8)--(146.1,-19.6)--(143.8,-18.7)--cycle;
				\definecolor{ezs}{RGB}{253,232,85}\fill[ezs](80.4,-189.6)--(82.7,-190.5)--(148,-20.3)--(145.7,-19.4)--cycle;
				\definecolor{ezs}{RGB}{253,230,82}\fill[ezs](82.3,-190.4)--(84.6,-191.2)--(149.9,-21)--(147.6,-20.2)--cycle;
				\definecolor{ezs}{RGB}{253,228,78}\fill[ezs](84.2,-191.1)--(86.5,-192)--(151.8,-21.8)--(149.5,-20.9)--cycle;
				\definecolor{ezs}{RGB}{253,226,74}\fill[ezs](86.1,-191.8)--(88.4,-192.7)--(153.7,-22.5)--(151.4,-21.6)--cycle;
				\definecolor{ezs}{RGB}{253,224,70}\fill[ezs](88,-192.6)--(90.3,-193.4)--(155.6,-23.3)--(153.3,-22.4)--cycle;
				\definecolor{ezs}{RGB}{253,222,67}\fill[ezs](89.9,-193.3)--(92.2,-194.2)--(157.5,-24)--(155.3,-23.1)--cycle;
				\definecolor{ezs}{RGB}{253,220,65}\fill[ezs](91.8,-194)--(94.1,-194.9)--(159.5,-24.7)--(157.2,-23.8)--cycle;
				\definecolor{ezs}{RGB}{253,217,63}\fill[ezs](93.7,-194.8)--(96,-195.6)--(161.4,-25.5)--(159.1,-24.6)--cycle;
				\definecolor{ezs}{RGB}{253,215,60}\fill[ezs](95.7,-195.5)--(97.9,-196.4)--(163.3,-26.2)--(161,-25.3)--cycle;
				\definecolor{ezs}{RGB}{253,213,56}\fill[ezs](97.6,-196.2)--(99.9,-197.1)--(165.2,-26.9)--(162.9,-26)--cycle;
				\definecolor{ezs}{RGB}{252,211,53}\fill[ezs](99.5,-197)--(101.8,-197.8)--(167.1,-27.7)--(164.8,-26.8)--cycle;
				\definecolor{ezs}{RGB}{252,209,49}\fill[ezs](101.4,-197.7)--(103.7,-198.6)--(169,-28.4)--(166.7,-27.5)--cycle;
				\definecolor{ezs}{RGB}{252,207,45}\fill[ezs](103.3,-198.4)--(105.6,-199.3)--(170.9,-29.1)--(168.6,-28.2)--cycle;
				\definecolor{ezs}{RGB}{253,205,42}\fill[ezs](105.2,-199.2)--(107.5,-200)--(172.8,-29.9)--(170.6,-29)--cycle;
				\definecolor{ezs}{RGB}{253,202,40}\fill[ezs](107.1,-199.9)--(109.4,-200.8)--(174.7,-30.6)--(172.5,-29.7)--cycle;
				\definecolor{ezs}{RGB}{253,200,38}\fill[ezs](109,-200.6)--(111.3,-201.5)--(176.7,-31.3)--(174.4,-30.4)--cycle;
				\definecolor{ezs}{RGB}{253,198,35}\fill[ezs](111,-201.4)--(113.2,-202.2)--(178.6,-32.1)--(176.3,-31.2)--cycle;
				\definecolor{ezs}{RGB}{253,196,31}\fill[ezs](112.9,-202.1)--(115.2,-203)--(180.5,-32.8)--(178.2,-31.9)--cycle;
				\definecolor{ezs}{RGB}{253,194,28}\fill[ezs](114.8,-202.8)--(117.1,-203.7)--(182.4,-33.5)--(180.1,-32.6)--cycle;
				\definecolor{ezs}{RGB}{253,192,26}\fill[ezs](116.7,-203.6)--(119,-204.5)--(184.3,-34.3)--(182,-33.4)--cycle;
				\definecolor{ezs}{RGB}{253,189,24}\fill[ezs](118.6,-204.3)--(120.9,-205.2)--(186.2,-35)--(183.9,-34.1)--cycle;
				\definecolor{ezs}{RGB}{253,186,20}\fill[ezs](120.5,-205)--(122.8,-205.9)--(188.1,-35.7)--(185.8,-34.8)--cycle;
				\definecolor{ezs}{RGB}{253,184,17}\fill[ezs](122.4,-205.8)--(124.7,-206.7)--(190,-36.5)--(187.8,-35.6)--cycle;
				\definecolor{ezs}{RGB}{252,183,15}\fill[ezs](124.3,-206.5)--(126.6,-207.4)--(192,-37.2)--(189.7,-36.3)--cycle;
				\definecolor{ezs}{RGB}{252,181,14}\fill[ezs](126.3,-207.2)--(128.5,-208.1)--(193.9,-37.9)--(191.6,-37.1)--cycle;
				\definecolor{ezs}{RGB}{252,179,13}\fill[ezs](128.2,-208)--(130.4,-208.9)--(195.8,-38.7)--(193.5,-37.8)--cycle;
				\definecolor{ezs}{RGB}{252,178,12}\fill[ezs](130.1,-208.7)--(132.4,-209.6)--(197.7,-39.4)--(195.4,-38.5)--cycle;
				\definecolor{ezs}{RGB}{251,178,10}\fill[ezs](132,-209.4)--(134.3,-210.3)--(199.6,-40.1)--(197.3,-39.3)--cycle;
				\definecolor{ezs}{RGB}{251,178,10}\fill[ezs](133.9,-210.2)--(136.2,-211.1)--(201.5,-40.9)--(199.2,-40)--cycle;
				\definecolor{ezs}{RGB}{250,177,11}\fill[ezs](135.8,-210.9)--(138.1,-211.8)--(203.4,-41.6)--(201.1,-40.7)--cycle;
				\definecolor{ezs}{RGB}{250,178,11}\fill[ezs](137.7,-211.6)--(140,-212.5)--(205.3,-42.3)--(203.1,-41.5)--cycle;
				\definecolor{ezs}{RGB}{248,180,10}\fill[ezs](139.6,-212.4)--(141.9,-213.3)--(207.3,-43.1)--(205,-42.2)--cycle;
				\definecolor{ezs}{RGB}{247,181,9}\fill[ezs](141.5,-213.1)--(143.8,-214)--(209.2,-43.8)--(206.9,-42.9)--cycle;
				\definecolor{ezs}{RGB}{246,182,9}\fill[ezs](143.5,-213.8)--(145.7,-214.7)--(211.1,-44.5)--(208.8,-43.7)--cycle;
				\definecolor{ezs}{RGB}{246,182,10}\fill[ezs](145.4,-214.6)--(147.7,-215.5)--(213,-45.3)--(210.7,-44.4)--cycle;
			\end{scope}
			% chữ EDUZY (Times đậm, làm dày nét bằng 5 lớp lệch nhẹ; căn đúng khung chữ logo gốc)
			\foreach \dx/\dy in {-0.7/0,0.7/0,0/-0.3,0/0.3,0/0}
			\node[anchor=center,inner sep=0pt,outer sep=0pt,text=ezText,align=center] at (109.6+\dx,-193.6+\dy)
			{\resizebox{1.01cm}{0.22cm}{\fontencoding{T1}\fontfamily{ptm}\fontseries{b}\selectfont EDUZY}};
		\end{tikzpicture}%
	}%
}
%====================== THIẾT LẬP THÂN ĐỀ ======================
%--- Môi trường tự luận {bt}: ex_test 3.3.4 chưa khai báo sẵn nên tạo ở đây
%    (ex_test đã tự móc \fixshowans và \listenumerate cho bt)
\makeatletter
\@ifundefined{bt}{\newtheorem{bt}{\color{gold1}\namebt}}{}
\makeatother
\renewcommand{\sectionmark}[1]{\markright{\thesection\ #1}}
%--- Khung tiêu đề (logo + tối đa 4 dòng; dòng 3, 4 để trống {} thì bỏ qua)
%  Đề thi:    \khungtieude{THI THỬ THPT QUỐC GIA NĂM 2026}{Môn Toán -- Đề số 21}
%                         {(Đề thi gồm: ...)}{Thời gian làm bài: 90 phút (không kể thời gian phát đề).}
%  Chuyên đề: \khungtieude{CHUYÊN ĐỀ 3 -- DÃY SỐ VÀ CÔNG THỨC TRUY HỒI}
%                         {Bồi dưỡng học sinh giỏi Toán THPT}{...}{}
\newcommand{\khungtieude}[4]{%
	\begin{tcolorbox}[colframe=gold1,colback=gold5!45,boxrule=0.5mm]
		\begin{minipage}[c]{3cm}
			\centerline{\eduzylogo[2.4cm]}
		\end{minipage}\hspace{0.5cm}
		\begin{minipage}[c]{12.5cm}
			\vskip 0.15 cm
			\centerline{\ind{#1}}
			\vskip 0.15 cm
			\centerline{\textbf{#2}}
			\if\relax\detokenize{#3}\relax\else
				\vskip 0.15 cm
				\centerline{\textit{#3}}
			\fi
			\if\relax\detokenize{#4}\relax\else
				\vskip 0.15 cm
				\centerline{\textit{#4}}
			\fi
			\vskip 0.15 cm
		\end{minipage}
	\end{tcolorbox}
}
%====================== BỔ SUNG CHO BÀI HỌC, CHUYÊN ĐỀ ======================
\usepackage[shortlabels]{enumitem}
\usepackage{tabularx}
%--- Danh sách gọn: \begin{enumerate}[gon] ... \end{enumerate}
\SetEnumitemKey{gon}{leftmargin=*,itemsep=1pt,parsep=0pt,topsep=3pt}
%--- Khung lý thuyết: \begin{khung}{} ... \end{khung}; tiêu đề chỉ dùng khi thật ngắn ("Mẹo", "Lỗi thường gặp")
\newtcolorbox{khung}[1]{breakable,enhanced,colframe=gold1,colback=gold5!45,
	boxrule=0.4mm,fonttitle=\bfseries\sffamily,coltitle=white,colbacktitle=gold1,
	attach boxed title to top left={xshift=6mm,yshift=-3mm},
	boxed title style={arc=1mm,boxrule=0.3mm},title={#1}}
%--- Tiêu đề mục nhỏ trong một dạng: \muc{Dấu hiệu}, \muc{Phương pháp}, \muc{Bài tự luyện}
\makeatletter
\newcommand{\muc}[1]{\par\medskip\noindent\indam{#1}\par\nobreak\smallskip\@afterindentfalse\@afterheading}
\makeatother
%--- Ô công thức cần nhớ (căn giữa, viền vàng, vừa khít nội dung)
%     \begin{center}\congthuc{$Q(n)=(n-1)!$}\end{center}
\newtcbox{\congthuc}{on line,colframe=gold1,colback=gold5!45,boxrule=0.5pt,arc=2pt,
	left=8pt,right=8pt,top=4pt,bottom=4pt}
%--- Công tắc lời giải bài tự luyện (ví dụ mẫu luôn dùng \loigiai)
%     \loigiaiTL[{đáp số}]{lời giải đầy đủ}
%     \TLfull -> in lời giải đầy đủ;  \TLdapso -> chỉ in một dòng "Đáp số";  \TLan -> ẩn hoàn toàn
%     (\CheDo là dethi thì mọi lời giải và đáp số đều ẩn)
\makeatletter
\newif\ifTLhien   \TLhientrue
\newif\ifTLdapso  \TLdapsofalse
\newcommand{\TLfull}{\TLhientrue\TLdapsofalse}
\newcommand{\TLdapso}{\TLhientrue\TLdapsotrue}
\newcommand{\TLan}{\TLhienfalse}
\newcommand{\loigiaiTL}[2][]{%
	\ifTLhien
		\ifTLdapso
			\def\@tlarg{#1}%
			\ifx\@tlarg\@empty\else
				\ifdefstring{\kindOPTN}{dethi}{}{\par\noindent{\bfseries\color{gold1}Đáp số:} #1\par}%
			\fi
		\else
			\loigiai{#2}%
		\fi
	\fi}
\makeatother
\TLdapso% mặc định: chỉ in đáp số bài tự luyện
%====================== BẢNG ĐÚNG/SAI (\choiceTFt) ======================
%  Hàng tiêu đề nền nâu đồng chữ trắng (Phát biểu | Đ | S), các hàng tô xen kẽ kem / trắng,
%  đường kẻ vàng kim, nhãn a) b) c) d) màu nâu đồng. Bản có đáp án: cột được chọn có vòng tròn
%  nâu đồng với dấu tích trắng (câu đúng tích ở cột Đúng, câu sai tích ở cột Sai).
%  Tiêu đề ghi đủ "Đúng | Sai": \OPTN{viettat=0} sau \input (cột Đúng/Sai rộng ra, phát biểu hẹp lại).
%  Dùng bảng gốc của ex_test: \def\BangDSgoc{1} trước \input{preamble-eduzy-v2}.
\providecommand{\BangDSgoc}{0}
\ifnum\BangDSgoc=0
\usepackage{colortbl,pifont}
\colorlet{ezDSke}{gold2!60!gold1}% đường kẻ bảng (nét 0,4pt như ex_test để không đổi chiều cao bảng)
\colorlet{ezDSle}{gold5!60}% hàng a), c)
\colorlet{ezDSchan}{white}% hàng b), d)
%--- dấu chọn ở cột Đúng/Sai
%    (bề rộng 0, căn giữa ô: không làm đổi bề rộng cột của ex_test, không báo tràn lề)
\newcommand{\EZdauchon}{\makebox[0pt]{\tikz[baseline=-1ex]\node[circle,fill=gold1,inner sep=0pt,
	minimum size=1.2em,text=white,font=\small]{\ding{51}};}}
\def\tickT{\EZdauchon}\def\tickF{\EZdauchon}
\AtBeginDocument{\sbox{\ccTFTab}{\EZdauchon}}
%--- nhãn a) b) c) d) màu nâu đồng (ở bản lời giải, nhãn câu đúng vẫn đóng khung như ex_test)
\renewcommand{\FalseTF}{\stepcounter{dapan}\textbf{\color{gold1}\DapAnTF\sepTF}\ignorespaces}
%--- trước khi dựng bảng: kẻ màu vàng kim (bề rộng cột giữ nguyên như ex_test); sau bảng trả lại màu đen
\newcount\EZhangTF
\newcommand{\EZbatdauTF}{%
	\global\EZhangTF=0\relax
	\arrayrulecolor{ezDSke}}
\newcommand{\EZketthucTF}{\arrayrulecolor{black}}
%--- hàng tiêu đề nền nâu đồng, chữ trắng
\makeatletter
\newcommand{\EZo@TF}[1]{\centering\bfseries\color{white}#1}
\newcommand{\EZtieudeTF}{%
	\ifcase\cotDS
		\ifnum\cotSTT=1\relax
			\def\rowHeader{\hline\rowcolor{gold1}\EZo@TF{TT}&\EZo@TF{\phatbieu}\tabularnewline}%
		\else
			\def\rowHeader{\hline\rowcolor{gold1}\EZo@TF{\phatbieu}\tabularnewline}%
		\fi
	\or
		\ifnum\cotSTT=1\relax
			\def\rowHeader{\hline\rowcolor{gold1}\EZo@TF{TT}&\EZo@TF{\phatbieu}&\EZo@TF{\Dung}\tabularnewline}%
		\else
			\def\rowHeader{\hline\rowcolor{gold1}\EZo@TF{\phatbieu}&\EZo@TF{\Dung}\tabularnewline}%
		\fi
	\or
		\ifnum\cotSTT=1\relax
			\def\rowHeader{\hline\rowcolor{gold1}\EZo@TF{TT}&\EZo@TF{\phatbieu}&\EZo@TF{\Dung}&\EZo@TF{\Sai}\tabularnewline}%
		\else
			\def\rowHeader{\hline\rowcolor{gold1}\EZo@TF{\phatbieu}&\EZo@TF{\Dung}&\EZo@TF{\Sai}\tabularnewline}%
		\fi
	\fi}
%--- tô xen kẽ từng hàng (gọi ngay sau \hline đầu mỗi hàng)
\newcommand{\EZmauhangTF}{%
	\noalign{\global\advance\EZhangTF\@ne}%
	\ifodd\EZhangTF\rowcolor{ezDSle}\else\rowcolor{ezDSchan}\fi}
%--- gắn vào \choiceTFt / \writeChoiceTFt của ex_test
\xpretocmd{\choiceTFt}{\EZbatdauTF}{}{\typeout{[preamble-eduzy-v2] Không gắn được bảng Đúng/Sai (đầu bảng)}}
\xpatchcmd{\choiceTFt}{\ifnum\dongPB=0\def\rowHeader{}\fi}{\EZtieudeTF\ifnum\dongPB=0\def\rowHeader{}\fi}
	{}{\typeout{[preamble-eduzy-v2] Không gắn được bảng Đúng/Sai (hàng tiêu đề)}}
\xpatchcmd{\writeChoiceTFt}{\hline\xdef\dotEX{\sdotEX}}{\hline\EZmauhangTF\xdef\dotEX{\sdotEX}}
	{}{\typeout{[preamble-eduzy-v2] Không gắn được bảng Đúng/Sai (tô hàng)}}
\xpatchcmd{\writeChoiceTFt}{{\end{longtable}}}{{\end{longtable}\EZketthucTF}}
	{}{\typeout{[preamble-eduzy-v2] Không gắn được bảng Đúng/Sai (cuối longtable)}}
\xpatchcmd{\writeChoiceTFt}{{\end{xtabular}\end{center}\medskip\noindent}}{{\end{xtabular}\end{center}\EZketthucTF\medskip\noindent}}
	{}{\typeout{[preamble-eduzy-v2] Không gắn được bảng Đúng/Sai (cuối xtabular)}}
\fi
%=====================================================================
\renewcommand{\parallel}{\mathrel{/\!\!/}}

%  trong file đề (tools/dan_trang/dan_trang.py tự ghi khi đặt \NeedspaceCau):
%    \def\TachCauDS{1}   câu Đúng/Sai được tách: câu dẫn (kèm hình) ở cuối trang, bảng Phát biểu sang trang sau
%                        (câu dẫn và hình không tách nhau, bảng không tách giữa hai trang)
%    \def\TachCauLG{1}   bản lời giải (\CheDo khác dethi): phần đề của câu cao hơn 1/3 trang được phép tách
%  Lệnh dùng trong file đề (giá trị do tools/dan_trang/dan_trang.py đo ở bản dethi):
%    \NeedspaceCau[<cao câu dẫn + hình>]{<cao cả khối câu>}   ngay trên \begin{ex} (câu đầu phần: trên tiêu đề phần)
%        công tắc tắt: đúng bằng \Needspace{<cao cả khối câu>} như cũ;
%        \TachCauDS=1 và có [..] (câu Đúng/Sai được chọn tách): chỉ giữ câu dẫn + hình;
%        \TachCauLG=1 ở bản lời giải: đòi tối đa 1/3 \textheight.
%    \NeedspaceBang{<cao bảng Phát biểu>}%   ngay trước \choiceTFt: \TachCauDS=1 thì giữ trọn bảng trên một trang;
%        công tắc tắt: không làm gì.
\providecommand{\TachCauDS}{0}
\providecommand{\TachCauLG}{0}
\newlength{\EZnsCau}
\newcommand{\NeedspaceCau}[2][]{%
	\setlength{\EZnsCau}{#2}%
	\ifnum\TachCauDS=1 \if\relax\detokenize{#1}\relax\else\setlength{\EZnsCau}{#1}\fi\fi
	\ifnum\TachCauLG=1 \ifdefstring{\kindOPTN}{dethi}{}{\ifdim\EZnsCau>0.3333\textheight\setlength{\EZnsCau}{0.3333\textheight}\fi}\fi
	\Needspace{\EZnsCau}}
\newcommand{\NeedspaceBang}[1]{\ifnum\TachCauDS=1 \Needspace{#1}\fi}

%--- Thang màu: đậm -> nhạt
\definecolor{gold1}{RGB}{126,96,0}    % nâu đồng đậm  #7E6000
\definecolor{gold2}{RGB}{192,144,3}   % vàng kim      #C09003
\definecolor{gold3}{RGB}{255,215,99}  % vàng nhạt     #FFD763
\definecolor{gold4}{RGB}{255,231,153} % vàng rất nhạt #FFE799
\definecolor{gold5}{RGB}{255,243,205} % vàng kem      #FFF3CD
%--- Ghi đè tông xanh/đỏ cũ
\definecolor{myblue}{RGB}{126,96,0}
\colorlet{blue}{gold1}
\colorlet{cyan}{gold2}
\colorlet{teal}{gold2}

%--- Tên chương mặc định (file đề nên tự khai báo lại)
\providecommand{\tenchuong}{}
%--- Chân trang
\providecommand{\thuonghieu}{TOÁN 12 -- Eduzy.com.vn}
\providecommand{\tacgia}{Tô Gia Bảo -- 0353703783}

%--- Chữ cạnh logo ở đầu trang
\providecommand{\khauhieu}{Eduzy.com.vn}

%====================== THÔNG SỐ ĐẦU TRANG, CHÂN TRANG (v2) ======================
%  Dải nền đậm tràn ngang trang; ô chữ vàng kem có cạnh xiên, kèm 2 vạch xiên chuyển màu.
%  Đầu trang lẻ : logo + \khauhieu (trái)   | ô chữ \tenchuong (phải)
%  Chân trang lẻ: ô chữ \thuonghieu (trái)  | Trang (số) (phải)
%  Trang chẵn soi gương trang lẻ; chân trang chẵn ghi \tacgia.
%--- Màu: chỉ dùng thang vàng kim Eduzy
\colorlet{ezNen}{gold1}% dải nền đậm
\colorlet{ezVach1}{gold3}% vạch xiên sát dải nền (đậm hơn)
\colorlet{ezVach2}{gold4}% vạch xiên sát ô chữ (nhạt hơn)
\colorlet{ezO}{gold5}% ô chữ vàng kem
%--- Kích thước
\newlength{\EZdai}\setlength{\EZdai}{1.1cm}% chiều cao dải nền
\newlength{\EZo}\setlength{\EZo}{0.68cm}% chiều cao ô chữ
\newlength{\EZxien}\setlength{\EZxien}{0.55cm}% độ lệch ngang của cạnh xiên
\newlength{\EZvach}\setlength{\EZvach}{0.75cm}% bề rộng mỗi vạch xiên
\newlength{\EZkhe}\setlength{\EZkhe}{0.12cm}% khe trắng giữa vạch và ô chữ
\newlength{\EZdem}\setlength{\EZdem}{0.4cm}% khoảng đệm từ chữ tới cạnh xiên
\newlength{\EZle}\setlength{\EZle}{1.8cm}% lề trái/phải của chữ (bằng lề thân bài)
\newlength{\EZrongmin}\setlength{\EZrongmin}{4.5cm}% ô chữ rộng tối thiểu
\newlength{\EZrong}% bề rộng chữ trong ô (tính tự động)
%--- Font
\newcommand{\EZfontO}{\fontsize{10.5pt}{12pt}\selectfont\bfseries}% chữ trong ô
\newcommand{\EZfontKH}{\fontfamily{qag}\fontsize{11pt}{12pt}\selectfont\bfseries}% \khauhieu cạnh logo
\newcommand{\EZfontTrang}{\fontsize{11.5pt}{12pt}\selectfont\bfseries}% chữ "Trang" và số trang
%--- Đo bề rộng chữ trong ô
\newcommand{\EZdorong}[1]{%
	\settowidth{\EZrong}{\EZfontO #1}%
	\ifdim\EZrong<\EZrongmin\setlength{\EZrong}{\EZrongmin}\fi}
%--- Logo + chữ trắng \khauhieu (dựng sẵn một lần khi vào \begin{document})
\newsavebox{\EZlogobox}
\AtBeginDocument{\sbox{\EZlogobox}{%
		\raisebox{-0.5\height}{\eduzylogo[1cm]}\hspace{0.15cm}%
		\raisebox{-0.5\height}{\color{white}\EZfontKH\khauhieu}}}
%--- Hai vạch xiên + khe trắng. #1: điểm gốc ở mép dưới, #2: hướng (1 sang phải, -1 sang trái),
%    #3: hướng nghiêng (1 kiểu "/", -1 kiểu "\")
\newcommand{\EZvachxien}[3]{%
	\fill[white] (#1) --++({(#3)*\EZxien},\EZo) --++({(#2)*\EZkhe},0) --++({-(#3)*\EZxien},-\EZo) --cycle;
	\fill[ezVach2] ([xshift=(#2)*\EZkhe]#1) --++({(#3)*\EZxien},\EZo) --++({(#2)*\EZvach},0) --++({-(#3)*\EZxien},-\EZo) --cycle;
	\fill[ezVach1] ([xshift=(#2)*(\EZkhe+\EZvach)]#1) --++({(#3)*\EZxien},\EZo) --++({(#2)*\EZvach},0) --++({-(#3)*\EZxien},-\EZo) --cycle;
	\draw[white,line width=0.6pt] ([xshift=(#2)*(\EZkhe+\EZvach)]#1) --++({(#3)*\EZxien},\EZo);
}
%--- Chữ trong ô (căn giữa theo chiều cao ô, không lệch vì dấu tiếng Việt)
\tikzset{ezchu/.style={inner sep=0pt,text=ezNen,text height=0.72em,text depth=0.22em}}

%====================== LOGO CHÌM GIỮA TRANG ======================
%  Vẽ cùng đầu trang nên nằm dưới chữ của thân bài, có trên mọi trang.
%  Tắt: \logochimtat (từ trang đó trở đi), bật lại: \logochimbat.
%  Đổi cỡ, độ đậm: khai báo trước \input, ví dụ \def\EZchimcao{9cm} \def\EZchimdo{0.1}
\providecommand{\EZchimcao}{10cm}% chiều cao logo chìm
\providecommand{\EZchimdo}{0.12}% độ đậm (0 = trong suốt, 1 = rõ hẳn)
\newif\ifEZchim \EZchimtrue
\newcommand{\logochimtat}{\EZchimfalse}
\newcommand{\logochimbat}{\EZchimtrue}
%  Dựng sẵn logo mờ một lần khi vào \begin{document}; transparency group làm mờ cả logo
%  một lượt nên không lộ vệt ở chỗ các dải màu chồng nhau.
\newsavebox{\EZchimtho}
\newsavebox{\EZchimbox}
\AtBeginDocument{%
	\sbox{\EZchimtho}{\eduzylogo[\EZchimcao]}%
	\sbox{\EZchimbox}{\begin{tikzpicture}
			\begin{scope}[transparency group,opacity=\EZchimdo]
				\node[inner sep=0pt] {\usebox{\EZchimtho}};
			\end{scope}
	\end{tikzpicture}}}
\newcommand{\EZvechim}{%
	\ifEZchim
		\node[inner sep=0pt] at (current page.center) {\usebox{\EZchimbox}};
	\fi}

%====================== ĐẦU TRANG - TRANG LẺ ======================
%  (logo bên trái, ô chữ \tenchuong bên phải)
\newcommand{\myfancyheadOdd}{%
	\EZdorong{\tenchuong}%
	\begin{tikzpicture}[remember picture,overlay]
		\coordinate (T) at (current page.north west);
		\coordinate (P) at (current page.north east);
		%--- logo chìm giữa trang
		\EZvechim
		%--- dải nền
		\fill[ezNen] (T) rectangle ([yshift=-\EZdai]P);
		%--- ô chữ: X là góc dưới bên trái, cạnh trái xiên kiểu "/"
		\coordinate (X) at ([xshift=-\EZle-\EZrong-\EZdem-\EZxien,yshift=-\EZdai]P);
		\fill[ezO] (X) --++(\EZxien,\EZo) -- ([yshift=-\EZdai+\EZo]P) -- ([yshift=-\EZdai]P) --cycle;
		\EZvachxien{X}{-1}{1}
		\node[ezchu,anchor=east] at ([xshift=-\EZle,yshift=-\EZdai+0.5*\EZo]P) {\EZfontO\tenchuong};
		%--- logo + tên
		\node[inner sep=0pt,anchor=west] at ([xshift=\EZle,yshift=-0.5*\EZdai]T) {\usebox{\EZlogobox}};
	\end{tikzpicture}%
}

%====================== ĐẦU TRANG - TRANG CHẴN =====================
%  (soi gương: ô chữ \tenchuong bên trái, logo bên phải)
\newcommand{\myfancyheadEven}{%
	\EZdorong{\tenchuong}%
	\begin{tikzpicture}[remember picture,overlay]
		\coordinate (T) at (current page.north west);
		\coordinate (P) at (current page.north east);
		%--- logo chìm giữa trang
		\EZvechim
		%--- dải nền
		\fill[ezNen] (T) rectangle ([yshift=-\EZdai]P);
		%--- ô chữ: X là góc dưới bên phải, cạnh phải xiên kiểu "\"
		\coordinate (X) at ([xshift=\EZle+\EZrong+\EZdem+\EZxien,yshift=-\EZdai]T);
		\fill[ezO] (X) --++(-\EZxien,\EZo) -- ([yshift=-\EZdai+\EZo]T) -- ([yshift=-\EZdai]T) --cycle;
		\EZvachxien{X}{1}{-1}
		\node[ezchu,anchor=west] at ([xshift=\EZle,yshift=-\EZdai+0.5*\EZo]T) {\EZfontO\tenchuong};
		%--- logo + tên
		\node[inner sep=0pt,anchor=east] at ([xshift=-\EZle,yshift=-0.5*\EZdai]P) {\usebox{\EZlogobox}};
	\end{tikzpicture}%
}

%====================== CHÂN TRANG - TRANG LẺ ======================
%  (ô chữ \thuonghieu bên trái, "Trang (số)" bên phải)
\newcommand{\myfancyfootOdd}{%
	\EZdorong{\thuonghieu}%
	\begin{tikzpicture}[remember picture,overlay]
		\coordinate (T) at (current page.south west);
		\coordinate (P) at (current page.south east);
		%--- dải nền
		\fill[ezNen] (T) rectangle ([yshift=\EZdai]P);
		%--- ô chữ sát mép trên dải: X là góc dưới bên phải, cạnh phải xiên kiểu "/"
		\coordinate (X) at ([xshift=\EZle+\EZrong+\EZdem,yshift=\EZdai-\EZo]T);
		\fill[ezO] ([yshift=\EZdai-\EZo]T) -- (X) --++(\EZxien,\EZo) -- ([yshift=\EZdai]T) --cycle;
		\EZvachxien{X}{1}{1}
		\node[ezchu,anchor=west] at ([xshift=\EZle,yshift=\EZdai-0.5*\EZo]T) {\EZfontO\thuonghieu};
		%--- số trang
		\node[circle,draw=white,line width=0.9pt,text=white,inner sep=0pt,minimum size=0.62cm,anchor=east]
		(so) at ([xshift=-\EZle,yshift=0.5*\EZdai]P) {\EZfontTrang\thepage};
		\node[inner sep=0pt,text=white,anchor=base east] (tr) at ([xshift=-0.15cm]so.base west) {\EZfontTrang Trang};
	\end{tikzpicture}%
}

%====================== CHÂN TRANG - TRANG CHẴN ====================
%  (soi gương: "Trang (số)" bên trái, ô chữ \tacgia bên phải)
\newcommand{\myfancyfootEven}{%
	\EZdorong{\tacgia}%
	\begin{tikzpicture}[remember picture,overlay]
		\coordinate (T) at (current page.south west);
		\coordinate (P) at (current page.south east);
		%--- dải nền
		\fill[ezNen] (T) rectangle ([yshift=\EZdai]P);
		%--- ô chữ sát mép trên dải: X là góc dưới bên trái, cạnh trái xiên kiểu "\"
		\coordinate (X) at ([xshift=-\EZle-\EZrong-\EZdem,yshift=\EZdai-\EZo]P);
		\fill[ezO] ([yshift=\EZdai-\EZo]P) -- (X) --++(-\EZxien,\EZo) -- ([yshift=\EZdai]P) --cycle;
		\EZvachxien{X}{-1}{-1}
		\node[ezchu,anchor=east] at ([xshift=-\EZle,yshift=\EZdai-0.5*\EZo]P) {\EZfontO\tacgia};
		%--- số trang
		%    (vòng tròn ẩn so0 nằm giữa dải để lấy đường chân chữ, giống trang lẻ)
		\node[circle,inner sep=0pt,minimum size=0.62cm] (so0) at ([yshift=0.5*\EZdai]T) {\EZfontTrang\phantom{\thepage}};
		\node[inner sep=0pt,text=white,anchor=base west] (tr) at ($(so0.base -| T)+(\EZle,0)$) {\EZfontTrang Trang};
		\node[circle,draw=white,line width=0.9pt,text=white,inner sep=0pt,minimum size=0.62cm,anchor=base west]
		(so) at ([xshift=0.15cm]tr.base east) {\EZfontTrang\thepage};
	\end{tikzpicture}%
}

%--- Giữ tương thích với các file cũ đang gọi \myfancyhead, \myfancyfoot
\newcommand{\myfancyhead}{\myfancyheadOdd}
\newcommand{\myfancyfoot}{\myfancyfootOdd}

%======================= GẮN VÀO TRANG =============================
\pagestyle{fancy}
\fancyhf{}
\fancyhead[LO]{\myfancyheadOdd}
\fancyhead[RE]{\myfancyheadEven}
\fancyfoot[LO]{\myfancyfootOdd}
\fancyfoot[RE]{\myfancyfootEven}
\renewcommand{\headrulewidth}{0pt}
\renewcommand{\footrulewidth}{0pt}
\fancypagestyle{plain}{%
	\fancyhf{}
	\fancyhead[LO]{\myfancyheadOdd}
	\fancyhead[RE]{\myfancyheadEven}
	\fancyfoot[LO]{\myfancyfootOdd}
	\fancyfoot[RE]{\myfancyfootEven}
	\renewcommand{\headrulewidth}{0pt}
	\renewcommand{\footrulewidth}{0pt}
}
%=================== HẾT PHẦN GỘP TỪ eduzy-de.tex ===================
\newcommand{\indam}[1]{
	{\fontsize{12pt}{9pt}\sffamily \selectfont\bfseries\color{gold1} #1}
}
\newcommand{\ind}[1]{
	{\fontsize{14pt}{9pt}\fontfamily{qag} \selectfont\bfseries\color{gold1} #1}
}
\renewcommand{\FalseEX}{\stepcounter{dapan}\textbf{\selectfont\color{gold1}\Alph{dapan}.}\ignorespaces}
\renewtheorem{ex}{\color{gold1}Câu}
\renewcommand{\baselinestretch}{1.1}
\newcommand{\hoac}[1]{
	\left[\begin{aligned}#1\end{aligned}\right.}
\newcommand{\heva}[1]{
	\left\{\begin{aligned}#1\end{aligned}\right.}
\usepackage{esvect}
\def\vec{\vv}
\def\vt{\vv}
\def\overrightarrow{\vv}
%====================== LOGO EDUZY VẼ BẰNG TIKZ ======================
%  Đường viền dò từ logo gốc; dải màu vàng kim đo trực tiếp trên ảnh gốc
%  (nghiêng 21 độ, tô bằng 80 dải hẹp nên hiển thị giống nhau trên mọi trình xem PDF).
%  Dùng: \eduzylogo          -> cao 2.4cm (mặc định)
%        \eduzylogo[1.5cm]   -> cao 1.5cm
\definecolor{ezText}{RGB}{245,200,46}
\newcommand{\eduzylogo}[1][2.4cm]{%
	\resizebox{!}{#1}{%
		\begin{tikzpicture}[x=0.01cm,y=0.01cm]
			\path[use as bounding box] (0,0) rectangle (224,-224);
			% mũ tốt nghiệp + vương miện
			\begin{scope}
				\clip
				(119.4,-176.1)..controls(118.9,-176.1)and(116.8,-176)..(114.6,-176)
				..controls(101.4,-176)and(85.8,-175.9)..(83.7,-175.7)
				..controls(81.8,-175.6)and(81,-175.6)..(80.2,-175.7)
				..controls(78.8,-175.9)and(77.7,-175.9)..(76.2,-175.7)
				..controls(74.9,-175.5)and(70.8,-175.5)..(64.3,-175.7)
				..controls(59.1,-175.9)and(58.3,-175.7)..(58.6,-173.7)
				..controls(58.6,-172.8)and(59.6,-169.6)..(60.3,-168)
				..controls(60.8,-166.9)and(61.4,-165.8)..(62.3,-165)
				..controls(63.1,-164.2)and(64.3,-163.4)..(65,-163.2)
				..controls(65.5,-163.1)and(72.3,-163)..(76.6,-163.1)
				..controls(77.9,-163.1)and(80.1,-163.1)..(81.4,-163)
				..controls(82.8,-162.9)and(86.3,-162.9)..(89.2,-162.9)
				..controls(93,-162.9)and(94.9,-162.8)..(96,-162.7)
				..controls(97.2,-162.5)and(101,-162.5)..(114,-162.4)
				..controls(123,-162.3)and(134.9,-162.2)..(140.2,-162.1)
				..controls(150.1,-162)and(152.2,-162.1)..(153.6,-162.4)
				..controls(154.4,-162.6)and(157.1,-163.8)..(158.1,-164.4)
				..controls(158.8,-164.8)and(159.7,-165.7)..(159.9,-166.3)
				..controls(160,-166.5)and(160.2,-167.1)..(160.2,-167.6)
				..controls(160.3,-168)and(160.6,-169.4)..(160.9,-170.6)
				..controls(161.2,-171.7)and(161.5,-173.1)..(161.6,-173.6)
				..controls(161.8,-174.4)and(161.8,-174.5)..(161.5,-174.9)
				..controls(161.2,-175.6)and(160.8,-175.9)..(159.9,-176)
				..controls(159.2,-176.1)and(154.1,-176.2)..(151.5,-176.2)
				..controls(148.1,-176.1)and(122.6,-176.1)..(121.6,-176.2)
				..controls(120.8,-176.2)and(119.8,-176.2)..(119.4,-176.1)--cycle
				(130.4,-160)..controls(129,-160)and(119.8,-160)..(110.1,-159.9)
				..controls(100.3,-159.8)and(88.5,-159.7)..(83.9,-159.6)
				..controls(77.8,-159.5)and(75.3,-159.5)..(74.5,-159.6)
				..controls(73.9,-159.7)and(72.5,-159.8)..(71.3,-159.8)
				..controls(69.6,-159.8)and(69,-159.8)..(68.5,-159.6)
				..controls(67.6,-159.3)and(66.8,-158.6)..(66.2,-157.6)
				..controls(66,-157.3)and(65.6,-156.5)..(65.3,-156.1)
				..controls(64.6,-155.1)and(64.5,-154.9)..(63.7,-152.9)
				..controls(63.2,-151.6)and(62.9,-150.9)..(61.8,-149.6)
				..controls(61.5,-149.2)and(61.1,-148.5)..(60.9,-147.9)
				..controls(60.7,-147.4)and(60.3,-146.6)..(60.1,-146.1)
				..controls(59.9,-145.5)and(59.5,-144.9)..(59,-144.3)
				..controls(58.5,-143.7)and(58.2,-143.2)..(57.9,-142.5)
				..controls(57.7,-141.9)and(57.3,-141.1)..(57,-140.7)
				..controls(56.6,-140.3)and(56.1,-139.3)..(55.8,-138.5)
				..controls(55.5,-137.8)and(55.1,-136.8)..(54.9,-136.5)
				..controls(54.7,-136.1)and(54.3,-135.2)..(53.9,-134.5)
				..controls(53.6,-133.8)and(53,-132.7)..(52.6,-132)
				..controls(51.6,-130.5)and(51.4,-130.2)..(51.1,-129)
				..controls(50.6,-127.5)and(50.6,-127.4)..(49.7,-126.3)
				..controls(49.2,-125.7)and(48.6,-124.8)..(48.3,-124.2)
				..controls(48,-123.6)and(47.5,-122.7)..(47.3,-122.2)
				..controls(47,-121.8)and(46.5,-120.7)..(46.1,-119.8)
				..controls(45.8,-118.9)and(45.2,-117.8)..(44.9,-117.2)
				..controls(44.2,-116.1)and(43.3,-114.4)..(42.3,-112.2)
				..controls(41.8,-111.3)and(41.3,-110.4)..(41.1,-110.1)
				..controls(40.5,-109.4)and(40.3,-109)..(40,-107.8)
				..controls(39.8,-107)and(39.6,-106.5)..(39.2,-106)
				..controls(38.9,-105.6)and(38.4,-104.7)..(38,-104)
				..controls(37.7,-103.4)and(37.2,-102.5)..(37,-102.1)
				..controls(36.7,-101.7)and(36.3,-100.8)..(36,-100.1)
				..controls(35.6,-98.9)and(35.3,-98.4)..(34.6,-97.6)
				..controls(34.2,-97.1)and(34.2,-96.3)..(34.6,-96)
				..controls(35,-95.7)and(35.5,-95.9)..(36.2,-96.7)
				..controls(36.5,-97)and(37.2,-97.5)..(37.6,-97.7)..controls(38,-98)and(38.6,-98.4)..(39.1,-98.9)
				..controls(39.6,-99.4)and(40.3,-100)..(40.6,-100.3)
				..controls(41.3,-100.8)and(42.4,-101.8)..(43.7,-102.9)
				..controls(44,-103.2)and(44.6,-103.7)..(45,-103.9)
				..controls(45.4,-104.2)and(45.9,-104.6)..(46.1,-104.8)
				..controls(47.1,-105.9)and(48.8,-107.1)..(50.5,-107.9)
				..controls(51.8,-108.6)and(52.3,-109)..(53.4,-110)
				..controls(54.8,-111.2)and(55.6,-111.6)..(57,-112.1)
				..controls(57.5,-112.3)and(58.5,-112.8)..(59.4,-113.3)
				..controls(60.2,-113.7)and(61.3,-114.3)..(61.9,-114.5)
				..controls(62.4,-114.7)and(63.3,-115.1)..(63.8,-115.3)
				..controls(64.4,-115.6)and(65.4,-115.9)..(66.2,-116.1)
				..controls(67,-116.3)and(68.4,-116.6)..(69.2,-116.9)
				..controls(72.3,-117.9)and(72.8,-118)..(78,-118.1)
				..controls(84,-118.3)and(86.1,-118.1)..(88.1,-117.1)
				..controls(89.4,-116.5)and(93,-115)..(93.7,-114.9)
				..controls(94.8,-114.6)and(95.4,-114.2)..(95.9,-113.4)
				..controls(96.2,-113)and(96.7,-112.5)..(97.1,-112.2)
				..controls(98,-111.6)and(101.3,-108.4)..(102.1,-107.2)
				..controls(102.4,-106.8)and(102.9,-106.2)..(103.2,-105.8)
				..controls(103.5,-105.5)and(104,-104.8)..(104.2,-104.2)
				..controls(104.5,-103.6)and(105,-102.8)..(105.3,-102.4)
				..controls(105.6,-102)and(106.1,-101.2)..(106.5,-100.6)
				..controls(107.1,-99.5)and(108.6,-97.5)..(109.5,-96.5)
				..controls(109.9,-96.1)and(110.1,-96)..(110.4,-96)
				..controls(111.1,-96)and(111.6,-96.5)..(111.8,-97.6)
				..controls(112,-98.6)and(112.3,-99.3)..(112.8,-99.9)
				..controls(113.7,-100.9)and(115.2,-103.1)..(115.7,-104.2)
				..controls(116.1,-105)and(116.5,-105.5)..(116.9,-106)
				..controls(117.3,-106.4)and(117.8,-106.9)..(118,-107.3)
				..controls(118.2,-107.6)and(118.6,-108)..(119,-108.3)
				..controls(119.6,-108.8)and(121.1,-110.3)..(121.8,-111.2)
				..controls(122.1,-111.7)and(122.5,-112)..(123.4,-112.4)
				..controls(124.8,-113.2)and(125.8,-113.8)..(127,-114.6)
				..controls(127.5,-114.9)and(128.1,-115.3)..(128.2,-115.4)
				..controls(128.4,-115.5)and(129.3,-115.8)..(130.2,-116)
				..controls(131.1,-116.2)and(132.5,-116.6)..(133.3,-116.9)
				..controls(135.7,-117.9)and(136.1,-117.9)..(142.5,-117.8)
				..controls(145.8,-117.8)and(148.4,-117.7)..(148.8,-117.6)
				..controls(149.2,-117.6)and(150.2,-117.2)..(151,-116.8)
				..controls(152,-116.4)and(152.8,-116)..(153.6,-115.9)
				..controls(154.5,-115.7)and(155.2,-115.4)..(156.3,-114.9)
				..controls(157.1,-114.5)and(158.2,-114)..(158.7,-113.8)
				..controls(159.1,-113.7)and(159.8,-113.3)..(160.2,-113)
				..controls(160.6,-112.7)and(161.3,-112.4)..(161.7,-112.2)
				..controls(162.7,-111.9)and(163.2,-111.7)..(164.1,-111.1)
				..controls(164.5,-110.8)and(165.3,-110.4)..(165.9,-110)
				..controls(166.5,-109.7)and(167.3,-109.2)..(167.6,-108.9)
				..controls(167.9,-108.6)and(168.6,-108.2)..(169,-108)
				..controls(169.4,-107.8)and(170.2,-107.3)..(170.7,-106.9)
				..controls(171.2,-106.5)and(172.1,-105.8)..(172.7,-105.4)
				..controls(173.2,-105)and(174.1,-104.4)..(174.5,-104)
				..controls(174.9,-103.7)and(175.7,-103.1)..(176.2,-102.6)
				..controls(176.7,-102.2)and(177.4,-101.5)..(177.9,-101.2)
				..controls(178.3,-100.8)and(178.8,-100.4)..(179,-100.2)
				..controls(179.2,-100)and(179.6,-99.6)..(180,-99.3)
				..controls(180.4,-99.1)and(180.9,-98.6)..(181.2,-98.4)
				..controls(181.5,-98.1)and(182.1,-97.7)..(182.5,-97.5)
				..controls(183.2,-97.2)and(183.5,-96.9)..(184.3,-95.9)
				..controls(184.6,-95.6)and(184.6,-95.6)..(184.8,-95.7)
				..controls(185,-95.8)and(185.1,-96)..(185.1,-96.2)
				..controls(185.1,-96.5)and(185,-96.7)..(184.6,-97.2)
				..controls(184.1,-97.8)and(184,-97.9)..(183.9,-98.8)
				..controls(183.7,-100.3)and(183.3,-101.3)..(182.9,-101.8)
				..controls(182.6,-102.1)and(182.3,-102.6)..(182.2,-102.9)
				..controls(182,-103.3)and(181.6,-104)..(181.2,-104.6)
				..controls(180.9,-105.2)and(180.5,-106.1)..(180.4,-106.6)
				..controls(179.9,-108)and(179.7,-108.5)..(179,-109.4)
				..controls(178.6,-109.8)and(178.2,-110.5)..(177.9,-111)
				..controls(177.7,-111.5)and(177.3,-112.3)..(177,-112.7)
				..controls(176.6,-113.2)and(176.4,-113.7)..(176.3,-114.2)
				..controls(176.2,-114.7)and(175.9,-115.2)..(175.4,-116)
				..controls(175,-116.6)and(174.4,-117.7)..(174,-118.5)
				..controls(173.7,-119.2)and(173.3,-120.1)..(173,-120.5)
				..controls(172.8,-120.9)and(172.4,-121.8)..(172.1,-122.5)
				..controls(171.8,-123.2)and(171.3,-124.1)..(171,-124.6)
				..controls(170.8,-125.1)and(170.3,-126)..(169.9,-126.6)
				..controls(169.6,-127.3)and(169.2,-128.1)..(169,-128.3)
				..controls(168.8,-128.6)and(168.5,-129.3)..(168.3,-129.8)
				..controls(168.1,-130.3)and(167.8,-131.1)..(167.5,-131.6)
				..controls(166.9,-132.5)and(164.2,-137.9)..(164.2,-138.1)
				..controls(164.2,-138.1)and(163.9,-138.8)..(163.4,-139.7)
				..controls(162.9,-140.5)and(162.4,-141.5)..(162.3,-141.9)
				..controls(161.9,-143)and(161.6,-143.4)..(160.9,-144.1)
				..controls(160.4,-144.6)and(160.3,-144.8)..(160.1,-145.4)
				..controls(160,-145.8)and(159.6,-146.8)..(159.2,-147.8)
				..controls(158.7,-148.8)and(158.1,-150)..(157.9,-150.6)
				..controls(157.6,-151.1)and(157.2,-151.9)..(156.9,-152.4)
				..controls(156.5,-152.8)and(156.1,-153.7)..(155.8,-154.4)
				..controls(155.6,-155)and(155.2,-155.9)..(154.9,-156.3)
				..controls(154.7,-156.7)and(154.4,-157.4)..(154.2,-157.9)
				..controls(154.1,-158.4)and(153.8,-159)..(153.5,-159.3)
				..controls(152.8,-160.1)and(153,-160.1)..(142.4,-160.1)
				..controls(137.2,-160.1)and(131.8,-160.1)..(130.4,-160)--cycle
				(141.5,-103.7)..controls(140.8,-103)and(139.8,-102.1)..(139.2,-101.6)
				..controls(138.5,-101.2)and(137.9,-100.6)..(137.8,-100.4)
				..controls(137,-99.5)and(134,-96.5)..(133.2,-96)
				..controls(132.7,-95.7)and(131.8,-94.9)..(131.1,-94.3)
				..controls(129.4,-92.7)and(127.9,-91.6)..(126.6,-91.1)
				..controls(125.2,-90.5)and(124.1,-90)..(122.9,-89.1)
				..controls(122.1,-88.6)and(121.3,-88.2)..(119.9,-87.7)
				..controls(118.8,-87.3)and(117.5,-86.8)..(116.9,-86.5)
				..controls(114.3,-85.5)and(113.7,-85.4)..(109.7,-85.5)
				..controls(105.3,-85.7)and(103.9,-85.9)..(101.6,-87)
				..controls(100.7,-87.4)and(99.4,-87.8)..(97.9,-88)
				..controls(97,-88.2)and(96.6,-88.3)..(95.8,-88.9)
				..controls(95.4,-89.2)and(94.7,-89.6)..(94.2,-89.8)
				..controls(92.9,-90.3)and(92.3,-90.7)..(90.9,-92.1)
				..controls(90.2,-92.7)and(89.4,-93.3)..(89,-93.5)..controls(87,-94.5)and(85.2,-96)..(83.6,-97.9)
				..controls(83.3,-98.2)and(82.5,-98.9)..(81.8,-99.4)
				..controls(80.7,-100.1)and(80.5,-100.3)..(80,-101.1)
				..controls(78.7,-103.1)and(77.2,-103.9)..(76.4,-103)
				..controls(75.7,-102.2)and(75.7,-103.2)..(75.7,-83.6)
				..controls(75.6,-73.6)and(75.6,-64.8)..(75.5,-64)
				..controls(75.5,-63.2)and(75.5,-61.4)..(75.6,-59.8)--(75.7,-57.1)--(76,-56.8)
				..controls(76.8,-56.2)and(78,-56.3)..(79.8,-57.2)
				..controls(80.4,-57.5)and(81.3,-57.8)..(81.8,-57.9)
				..controls(82.3,-58.1)and(83,-58.4)..(83.6,-58.8)..controls(84.3,-59.2)and(85,-59.5)..(86,-59.8)
				..controls(86.8,-60.1)and(88,-60.6)..(88.7,-60.9)
				..controls(89.3,-61.2)and(90.6,-61.7)..(91.4,-62)
				..controls(92.2,-62.3)and(93.7,-62.9)..(94.6,-63.3)
				..controls(95.6,-63.7)and(96.8,-64.1)..(97.4,-64.3)
				..controls(98,-64.5)and(99,-64.9)..(99.7,-65.2)
				..controls(100.4,-65.6)and(101.5,-66)..(102.3,-66.2)
				..controls(104.2,-66.8)and(105,-67.1)..(106.2,-67.8)
				..controls(107.9,-68.7)and(110.2,-69.4)..(111.3,-69.2)
				..controls(112.5,-69)and(115,-68.2)..(115.6,-67.8)
				..controls(117.3,-66.7)and(117.6,-66.6)..(119.5,-66.2)
				..controls(121.9,-65.8)and(122.1,-65.7)..(123,-65.1)
				..controls(124.1,-64.4)and(124.4,-64.3)..(126,-63.9)
				..controls(127,-63.7)and(127.7,-63.4)..(128.4,-63.1)
				..controls(129,-62.8)and(129.7,-62.5)..(130.6,-62.4)
				..controls(132.1,-62.1)and(134,-61.5)..(135.2,-61)
				..controls(135.7,-60.8)and(136.7,-60.5)..(137.3,-60.3)
				..controls(138.4,-60)and(139.3,-59.7)..(141.1,-59)
				..controls(141.6,-58.7)and(142,-58.6)..(142.5,-58.6)
				..controls(143.2,-58.6)and(143.3,-58.7)..(143.6,-59)
				..controls(144,-59.5)and(144.1,-59.9)..(144,-60.8)
				..controls(144,-61.2)and(144,-71.1)..(144,-82.8)--(143.9,-104.1)--(143.7,-104.5)
				..controls(143.5,-104.7)and(143.3,-104.8)..(143,-104.9)
				..controls(142.6,-104.9)and(142.5,-104.8)..(141.5,-103.7)--cycle
				(27.8,-91.4)..controls(27.1,-90.9)and(26.3,-90)..(26.1,-89.5)
				..controls(26,-89.2)and(25.8,-88.6)..(25.7,-87.9)
				..controls(25.4,-86.2)and(25.6,-84.8)..(26.5,-83.5)
				..controls(27.4,-82.3)and(28.3,-82)..(29.4,-82.3)
				..controls(31.1,-82.8)and(32.1,-84.9)..(32,-87.4)
				..controls(31.9,-89.3)and(31.2,-90.6)..(29.9,-91.4)
				..controls(29.1,-91.9)and(28.6,-91.9)..(27.8,-91.4)--cycle
				(190.1,-91.3)..controls(189.7,-91.2)and(189.3,-90.9)..(189.1,-90.7)
				..controls(188.3,-89.8)and(187.6,-87.7)..(187.6,-86.2)
				..controls(187.6,-85.5)and(187.7,-85.2)..(188.2,-84.4)
				..controls(189.1,-82.6)and(190.1,-81.8)..(191.3,-82.1)
				..controls(192,-82.3)and(193.5,-83.7)..(193.8,-84.4)
				..controls(194.1,-85.1)and(194.2,-87.1)..(194,-88.3)
				..controls(193.8,-89.8)and(193.3,-90.5)..(192.3,-91.3)
				..controls(191.7,-91.7)and(190.9,-91.7)..(190.1,-91.3)--cycle
				(173.7,-83.4)..controls(172.2,-82.5)and(172,-82.5)..(170.6,-82.4)
				..controls(169.7,-82.4)and(168.8,-82.4)..(168.3,-82.5)
				..controls(167.2,-82.7)and(166.9,-82.7)..(166.6,-82.2)
				..controls(166.3,-81.8)and(166.3,-81)..(166.8,-80)
				..controls(167.2,-78.9)and(167.4,-77.4)..(167.7,-74.1)
				..controls(167.8,-72.8)and(167.9,-71.8)..(167.9,-71.7)
				..controls(168,-71.7)and(168,-70.2)..(168,-68.4)..controls(168,-65.4)and(168,-64.9)..(167.7,-63)
				..controls(167.3,-60.5)and(167.1,-59.9)..(166.4,-58.4)
				..controls(166.1,-57.7)and(165.8,-56.9)..(165.7,-56.6)
				..controls(165.7,-56.2)and(165.4,-55.6)..(165.2,-55.2)
				..controls(164.6,-53.8)and(161.9,-50.7)..(161,-50.3)
				..controls(159.8,-49.7)and(157.9,-50)..(155.2,-51.2)
				..controls(154.7,-51.4)and(153.4,-51.8)..(152.2,-52.1)
				..controls(150.8,-52.4)and(149.4,-52.9)..(148.6,-53.2)
				..controls(147.8,-53.6)and(146.7,-54)..(146,-54.2)
				..controls(145.4,-54.4)and(144.5,-54.7)..(144,-54.9)
				..controls(143.6,-55.2)and(142.7,-55.5)..(141.8,-55.7)
				..controls(141,-55.9)and(139.6,-56.3)..(138.7,-56.6)
				..controls(137.8,-57)and(136.3,-57.4)..(135.3,-57.6)
				..controls(133.8,-57.9)and(133.2,-58.1)..(132.1,-58.7)
				..controls(131.1,-59.2)and(130.3,-59.4)..(129.4,-59.6)
				..controls(128.7,-59.8)and(127.8,-60)..(127.4,-60.1)
				..controls(127,-60.3)and(126,-60.5)..(125.1,-60.7)
				..controls(123.2,-61.1)and(121.7,-61.5)..(120.2,-62.1)
				..controls(119.7,-62.3)and(118.1,-62.8)..(116.7,-63.2)
				..controls(114.6,-63.9)and(114,-64.1)..(113,-64.7)
				..controls(111.9,-65.4)and(111.8,-65.4)..(111,-65.5)
				..controls(109.9,-65.5)and(109.3,-65.4)..(107.7,-64.7)
				..controls(107,-64.4)and(106,-64)..(105.3,-63.9)
				..controls(103.5,-63.5)and(102.2,-63)..(101.2,-62.5)
				..controls(100.1,-61.8)and(99.3,-61.5)..(97.8,-61.1)
				..controls(97.2,-61)and(96.2,-60.7)..(95.7,-60.4)
				..controls(95.1,-60.1)and(94.2,-59.8)..(93.5,-59.6)
				..controls(92.9,-59.4)and(91.7,-59)..(91,-58.7)..controls(90.2,-58.4)and(88.9,-58)..(88,-57.7)
				..controls(87.2,-57.5)and(86.1,-57.1)..(85.5,-56.8)
				..controls(85,-56.6)and(84.1,-56.2)..(83.5,-56)..controls(82.9,-55.7)and(81.8,-55.2)..(81,-54.9)
				..controls(80.1,-54.5)and(79.1,-54.2)..(78.7,-54.1)
				..controls(77.9,-53.8)and(77.3,-53.6)..(76.2,-52.9)
				..controls(75.7,-52.6)and(75.2,-52.4)..(74.1,-52.2)
				..controls(72.5,-51.9)and(72.4,-51.8)..(71.3,-51.1)
				..controls(69.9,-50.1)and(69.4,-49.9)..(67.9,-49.8)
				..controls(67,-49.7)and(66.6,-49.6)..(66.3,-49.4)
				..controls(66.1,-49.3)and(65.7,-49)..(65.4,-48.8)
				..controls(65.1,-48.7)and(64.2,-48.4)..(63.6,-48.2)
				..controls(62.9,-48)and(61.7,-47.6)..(60.9,-47.2)
				..controls(60.1,-46.8)and(58.9,-46.4)..(58.2,-46.2)
				..controls(57.6,-46)and(56.8,-45.7)..(56.5,-45.5)
				..controls(54.9,-44.6)and(54.6,-44.5)..(53.2,-44.2)
				..controls(51.7,-43.8)and(51,-43.6)..(49.9,-42.9)
				..controls(49.5,-42.6)and(48.9,-42.4)..(47.5,-42)
				..controls(45.7,-41.6)and(45.6,-41.5)..(44.3,-40.6)
				..controls(43.5,-40.1)and(42.8,-39.6)..(42.7,-39.4)
				..controls(42.4,-38.7)and(43,-38.1)..(44.5,-37.6)..controls(44.9,-37.5)and(46,-37)..(46.9,-36.6)
				..controls(47.9,-36.1)and(48.8,-35.8)..(49.6,-35.6)
				..controls(50.5,-35.4)and(51.3,-35.2)..(52.4,-34.6)
				..controls(53.3,-34.1)and(54.5,-33.7)..(54.9,-33.5)
				..controls(55.4,-33.4)and(56.4,-33)..(57.2,-32.6)
				..controls(58.6,-32)and(60.1,-31.5)..(61.6,-31.2)
				..controls(62.3,-31.1)and(62.7,-31)..(63.2,-30.7)
				..controls(64.1,-30)and(64.3,-29.9)..(65.3,-29.7)
				..controls(65.7,-29.6)and(66.5,-29.2)..(67.1,-28.9)
				..controls(67.6,-28.6)and(68.8,-28.2)..(69.6,-27.9)
				..controls(70.5,-27.7)and(71.6,-27.3)..(72.1,-27)
				..controls(72.6,-26.7)and(73.8,-26.2)..(74.8,-25.9)
				..controls(75.7,-25.5)and(76.7,-25.1)..(77.1,-25)
				..controls(77.4,-24.8)and(78.5,-24.4)..(79.6,-24.1)
				..controls(80.6,-23.8)and(82,-23.2)..(82.8,-22.9)
				..controls(83.5,-22.5)and(84.8,-22)..(85.8,-21.8)
				..controls(86.7,-21.5)and(87.8,-21.1)..(88.1,-20.9)
				..controls(88.4,-20.7)and(89.3,-20.3)..(90.1,-20.1)
				..controls(90.8,-19.8)and(91.9,-19.4)..(92.4,-19.1)
				..controls(93.5,-18.5)and(94.4,-18.2)..(95.4,-18.1)
				..controls(95.7,-18)and(96.4,-17.8)..(96.9,-17.5)
				..controls(97.5,-17.3)and(98.4,-16.9)..(99,-16.6)
				..controls(99.7,-16.3)and(100.9,-15.9)..(101.9,-15.6)
				..controls(102.9,-15.3)and(104.4,-14.7)..(105.3,-14.3)
				..controls(108.1,-13.2)and(108.5,-13.1)..(110.1,-13.1)
				..controls(111.4,-13.1)and(111.5,-13.1)..(112.4,-13.5)
				..controls(113.2,-14)and(115,-14.8)..(116.6,-15.4)
				..controls(116.9,-15.5)and(117.9,-15.8)..(118.8,-16.1)
				..controls(120.9,-16.6)and(121.5,-16.9)..(122.6,-17.9)
				..controls(123,-18.3)and(123.7,-18.8)..(123.9,-18.9)
				..controls(124.6,-19.4)and(126,-19.8)..(127.9,-20.2)
				..controls(129.6,-20.5)and(129.8,-20.6)..(131.1,-21.4)
				..controls(131.5,-21.7)and(132.4,-22.1)..(133.1,-22.3)
				..controls(136.1,-23.2)and(138.3,-24.1)..(140.2,-25)
				..controls(141,-25.5)and(141.5,-25.6)..(142.5,-25.8)
				..controls(145.1,-26.3)and(146,-26.6)..(149.4,-28.1)
				..controls(151.3,-28.9)and(153.5,-29.8)..(154.3,-30.1)
				..controls(155.1,-30.4)and(156.2,-30.8)..(156.6,-31)
				..controls(157.1,-31.2)and(158.2,-31.6)..(159.1,-31.8)
				..controls(160.6,-32.2)and(160.9,-32.3)..(161.7,-32.9)
				..controls(162.4,-33.3)and(162.9,-33.6)..(163.7,-33.8)
				..controls(164.3,-34)and(164.8,-34.1)..(164.9,-34.1)
				..controls(165,-34.1)and(165.6,-34.4)..(166.4,-34.8)
				..controls(167.3,-35.2)and(168.1,-35.4)..(169,-35.6)
				..controls(170,-35.8)and(170.5,-36)..(171.7,-36.6)
				..controls(172.4,-37)and(173.3,-37.4)..(173.6,-37.4)
				..controls(175.6,-37.7)and(176.2,-38)..(176.7,-39.1)
				..controls(176.9,-39.6)and(176.9,-39.7)..(176.7,-40.4)
				..controls(176.4,-41.6)and(175.8,-42.7)..(175.2,-43.4)
				..controls(174.5,-44.1)and(172.6,-45.5)..(171.7,-45.8)
				..controls(170.6,-46.2)and(169.1,-46)..(168.2,-45.4)
				..controls(167.9,-45.2)and(167.3,-44.8)..(166.9,-44.4)
				..controls(166.4,-44.1)and(165.8,-43.6)..(165.4,-43.2)
				..controls(165.1,-42.9)and(164.4,-42.3)..(163.9,-41.9)
				..controls(163.3,-41.5)and(162.7,-40.9)..(162.4,-40.5)
				..controls(161.4,-39.2)and(159.9,-38)..(158.8,-37.8)
				..controls(158.4,-37.8)and(158.3,-37.7)..(158.2,-37.2)
				..controls(158,-36.8)and(157.4,-36.1)..(157.1,-36.1)
				..controls(157,-36.1)and(156.7,-36.1)..(156.5,-36)
				..controls(156.1,-35.9)and(155.9,-35.7)..(155.7,-35.2)
				..controls(155.3,-34.6)and(155.1,-34.4)..(154,-34.1)
				..controls(153.6,-34)and(153.2,-33.7)..(152.9,-33.5)
				..controls(152.4,-32.9)and(151.6,-32.4)..(151.3,-32.4)
				..controls(150.9,-32.4)and(150.4,-32.6)..(150.4,-32.9)
				..controls(150.4,-33.2)and(150.8,-33.4)..(151.4,-33.6)
				..controls(152.1,-33.7)and(152.8,-34.1)..(153.5,-34.9)
				..controls(154.1,-35.5)and(154.3,-35.6)..(155,-35.8)
				..controls(155.7,-36)and(155.7,-36)..(156,-36.6)
				..controls(156.3,-37.3)and(156.7,-37.7)..(157.5,-37.8)
				..controls(158.1,-37.9)and(158.2,-38)..(158.4,-38.8)
				..controls(158.5,-39.4)and(158.6,-39.6)..(159.6,-40.4)
				..controls(160.2,-41)and(160.9,-41.6)..(161.3,-41.8)
				..controls(161.7,-42)and(161.9,-42.3)..(162.2,-42.8)
				..controls(162.4,-43.2)and(162.8,-43.8)..(163.1,-44.1)
				..controls(163.9,-44.8)and(165.1,-46.7)..(165.7,-48)
				..controls(166.2,-49.3)and(166.4,-49.7)..(167.2,-50.5)
				..controls(167.5,-50.8)and(167.9,-51.4)..(168.1,-51.8)
				..controls(168.3,-52.2)and(168.8,-53)..(169.1,-53.6)
				..controls(169.6,-54.3)and(169.8,-54.9)..(169.9,-55.6)
				..controls(170.4,-57.8)and(170.5,-58.1)..(171.2,-59.5)
				..controls(171.6,-60.3)and(172,-61.5)..(172.2,-62.1)
				..controls(172.3,-62.8)and(172.7,-63.7)..(172.9,-64.1)
				..controls(173.2,-64.6)and(173.4,-65.3)..(173.4,-65.6)
				..controls(173.8,-68.8)and(174.2,-70.6)..(174.9,-72.6)
				..controls(175.5,-74.5)and(175.9,-77.1)..(176.1,-80.2)
				..controls(176.3,-83.4)and(175.5,-84.4)..(173.7,-83.4)--cycle;
				\definecolor{ezs}{RGB}{247,179,26}\fill[ezs](-5.7,-156.6)--(-3.4,-157.5)--(61.9,12.7)--(59.7,13.6)--cycle;
				\definecolor{ezs}{RGB}{247,179,22}\fill[ezs](-3.8,-157.3)--(-1.5,-158.2)--(63.9,12)--(61.6,12.9)--cycle;
				\definecolor{ezs}{RGB}{248,179,19}\fill[ezs](-1.8,-158.1)--(0.4,-158.9)--(65.8,11.2)--(63.5,12.1)--cycle;
				\definecolor{ezs}{RGB}{248,178,16}\fill[ezs](0.1,-158.8)--(2.4,-159.7)--(67.7,10.5)--(65.4,11.4)--cycle;
				\definecolor{ezs}{RGB}{249,178,15}\fill[ezs](2,-159.5)--(4.3,-160.4)--(69.6,9.8)--(67.3,10.7)--cycle;
				\definecolor{ezs}{RGB}{251,177,17}\fill[ezs](3.9,-160.3)--(6.2,-161.1)--(71.5,9)--(69.2,9.9)--cycle;
				\definecolor{ezs}{RGB}{252,179,18}\fill[ezs](5.8,-161)--(8.1,-161.9)--(73.4,8.3)--(71.1,9.2)--cycle;
				\definecolor{ezs}{RGB}{252,182,17}\fill[ezs](7.7,-161.7)--(10,-162.6)--(75.3,7.6)--(73,8.5)--cycle;
				\definecolor{ezs}{RGB}{252,185,17}\fill[ezs](9.6,-162.5)--(11.9,-163.4)--(77.2,6.8)--(75,7.7)--cycle;
				\definecolor{ezs}{RGB}{251,189,18}\fill[ezs](11.5,-163.2)--(13.8,-164.1)--(79.2,6.1)--(76.9,7)--cycle;
				\definecolor{ezs}{RGB}{251,191,22}\fill[ezs](13.4,-163.9)--(15.7,-164.8)--(81.1,5.4)--(78.8,6.2)--cycle;
				\definecolor{ezs}{RGB}{252,193,25}\fill[ezs](15.4,-164.7)--(17.6,-165.6)--(83,4.6)--(80.7,5.5)--cycle;
				\definecolor{ezs}{RGB}{253,195,27}\fill[ezs](17.3,-165.4)--(19.6,-166.3)--(84.9,3.9)--(82.6,4.8)--cycle;
				\definecolor{ezs}{RGB}{253,196,31}\fill[ezs](19.2,-166.1)--(21.5,-167)--(86.8,3.2)--(84.5,4)--cycle;
				\definecolor{ezs}{RGB}{254,199,35}\fill[ezs](21.1,-166.9)--(23.4,-167.8)--(88.7,2.4)--(86.4,3.3)--cycle;
				\definecolor{ezs}{RGB}{253,201,38}\fill[ezs](23,-167.6)--(25.3,-168.5)--(90.6,1.7)--(88.3,2.6)--cycle;
				\definecolor{ezs}{RGB}{253,203,40}\fill[ezs](24.9,-168.3)--(27.2,-169.2)--(92.5,1)--(90.3,1.8)--cycle;
				\definecolor{ezs}{RGB}{253,205,42}\fill[ezs](26.8,-169.1)--(29.1,-170)--(94.4,0.2)--(92.2,1.1)--cycle;
				\definecolor{ezs}{RGB}{253,207,45}\fill[ezs](28.7,-169.8)--(31,-170.7)--(96.4,-0.5)--(94.1,0.4)--cycle;
				\definecolor{ezs}{RGB}{253,210,49}\fill[ezs](30.7,-170.5)--(32.9,-171.4)--(98.3,-1.2)--(96,-0.4)--cycle;
				\definecolor{ezs}{RGB}{253,212,54}\fill[ezs](32.6,-171.3)--(34.9,-172.2)--(100.2,-2)--(97.9,-1.1)--cycle;
				\definecolor{ezs}{RGB}{253,214,59}\fill[ezs](34.5,-172)--(36.8,-172.9)--(102.1,-2.7)--(99.8,-1.8)--cycle;
				\definecolor{ezs}{RGB}{253,216,63}\fill[ezs](36.4,-172.7)--(38.7,-173.6)--(104,-3.4)--(101.7,-2.6)--cycle;
				\definecolor{ezs}{RGB}{253,218,65}\fill[ezs](38.3,-173.5)--(40.6,-174.4)--(105.9,-4.2)--(103.6,-3.3)--cycle;
				\definecolor{ezs}{RGB}{254,220,66}\fill[ezs](40.2,-174.2)--(42.5,-175.1)--(107.8,-4.9)--(105.5,-4)--cycle;
				\definecolor{ezs}{RGB}{254,222,68}\fill[ezs](42.1,-175)--(44.4,-175.8)--(109.7,-5.6)--(107.5,-4.8)--cycle;
				\definecolor{ezs}{RGB}{254,224,71}\fill[ezs](44,-175.7)--(46.3,-176.6)--(111.7,-6.4)--(109.4,-5.5)--cycle;
				\definecolor{ezs}{RGB}{254,226,75}\fill[ezs](46,-176.4)--(48.2,-177.3)--(113.6,-7.1)--(111.3,-6.2)--cycle;
				\definecolor{ezs}{RGB}{254,228,79}\fill[ezs](47.9,-177.2)--(50.1,-178)--(115.5,-7.8)--(113.2,-7)--cycle;
				\definecolor{ezs}{RGB}{254,230,82}\fill[ezs](49.8,-177.9)--(52.1,-178.8)--(117.4,-8.6)--(115.1,-7.7)--cycle;
				\definecolor{ezs}{RGB}{254,232,85}\fill[ezs](51.7,-178.6)--(54,-179.5)--(119.3,-9.3)--(117,-8.4)--cycle;
				\definecolor{ezs}{RGB}{254,234,88}\fill[ezs](53.6,-179.4)--(55.9,-180.2)--(121.2,-10)--(118.9,-9.2)--cycle;
				\definecolor{ezs}{RGB}{254,237,91}\fill[ezs](55.5,-180.1)--(57.8,-181)--(123.1,-10.8)--(120.8,-9.9)--cycle;
				\definecolor{ezs}{RGB}{254,239,94}\fill[ezs](57.4,-180.8)--(59.7,-181.7)--(125,-11.5)--(122.8,-10.6)--cycle;
				\definecolor{ezs}{RGB}{254,241,97}\fill[ezs](59.3,-181.6)--(61.6,-182.4)--(127,-12.2)--(124.7,-11.4)--cycle;
				\definecolor{ezs}{RGB}{254,243,101}\fill[ezs](61.2,-182.3)--(63.5,-183.2)--(128.9,-13)--(126.6,-12.1)--cycle;
				\definecolor{ezs}{RGB}{254,245,103}\fill[ezs](63.2,-183)--(65.4,-183.9)--(130.8,-13.7)--(128.5,-12.8)--cycle;
				\definecolor{ezs}{RGB}{254,245,105}\fill[ezs](65.1,-183.8)--(67.4,-184.6)--(132.7,-14.4)--(130.4,-13.6)--cycle;
				\definecolor{ezs}{RGB}{254,245,105}\fill[ezs](67,-184.5)--(69.3,-185.4)--(134.6,-15.2)--(132.3,-14.3)--cycle;
				\definecolor{ezs}{RGB}{254,244,103}\fill[ezs](68.9,-185.2)--(71.2,-186.1)--(136.5,-15.9)--(134.2,-15)--cycle;
				\definecolor{ezs}{RGB}{253,243,101}\fill[ezs](70.8,-186)--(73.1,-186.8)--(138.4,-16.6)--(136.1,-15.8)--cycle;
				\definecolor{ezs}{RGB}{253,241,98}\fill[ezs](72.7,-186.7)--(75,-187.6)--(140.3,-17.4)--(138,-16.5)--cycle;
				\definecolor{ezs}{RGB}{253,239,95}\fill[ezs](74.6,-187.4)--(76.9,-188.3)--(142.2,-18.1)--(140,-17.2)--cycle;
				\definecolor{ezs}{RGB}{253,236,92}\fill[ezs](76.5,-188.2)--(78.8,-189)--(144.2,-18.8)--(141.9,-18)--cycle;
				\definecolor{ezs}{RGB}{253,234,88}\fill[ezs](78.5,-188.9)--(80.7,-189.8)--(146.1,-19.6)--(143.8,-18.7)--cycle;
				\definecolor{ezs}{RGB}{253,232,85}\fill[ezs](80.4,-189.6)--(82.7,-190.5)--(148,-20.3)--(145.7,-19.4)--cycle;
				\definecolor{ezs}{RGB}{253,230,82}\fill[ezs](82.3,-190.4)--(84.6,-191.2)--(149.9,-21)--(147.6,-20.2)--cycle;
				\definecolor{ezs}{RGB}{253,228,78}\fill[ezs](84.2,-191.1)--(86.5,-192)--(151.8,-21.8)--(149.5,-20.9)--cycle;
				\definecolor{ezs}{RGB}{253,226,74}\fill[ezs](86.1,-191.8)--(88.4,-192.7)--(153.7,-22.5)--(151.4,-21.6)--cycle;
				\definecolor{ezs}{RGB}{253,224,70}\fill[ezs](88,-192.6)--(90.3,-193.4)--(155.6,-23.3)--(153.3,-22.4)--cycle;
				\definecolor{ezs}{RGB}{253,222,67}\fill[ezs](89.9,-193.3)--(92.2,-194.2)--(157.5,-24)--(155.3,-23.1)--cycle;
				\definecolor{ezs}{RGB}{253,220,65}\fill[ezs](91.8,-194)--(94.1,-194.9)--(159.5,-24.7)--(157.2,-23.8)--cycle;
				\definecolor{ezs}{RGB}{253,217,63}\fill[ezs](93.7,-194.8)--(96,-195.6)--(161.4,-25.5)--(159.1,-24.6)--cycle;
				\definecolor{ezs}{RGB}{253,215,60}\fill[ezs](95.7,-195.5)--(97.9,-196.4)--(163.3,-26.2)--(161,-25.3)--cycle;
				\definecolor{ezs}{RGB}{253,213,56}\fill[ezs](97.6,-196.2)--(99.9,-197.1)--(165.2,-26.9)--(162.9,-26)--cycle;
				\definecolor{ezs}{RGB}{252,211,53}\fill[ezs](99.5,-197)--(101.8,-197.8)--(167.1,-27.7)--(164.8,-26.8)--cycle;
				\definecolor{ezs}{RGB}{252,209,49}\fill[ezs](101.4,-197.7)--(103.7,-198.6)--(169,-28.4)--(166.7,-27.5)--cycle;
				\definecolor{ezs}{RGB}{252,207,45}\fill[ezs](103.3,-198.4)--(105.6,-199.3)--(170.9,-29.1)--(168.6,-28.2)--cycle;
				\definecolor{ezs}{RGB}{253,205,42}\fill[ezs](105.2,-199.2)--(107.5,-200)--(172.8,-29.9)--(170.6,-29)--cycle;
				\definecolor{ezs}{RGB}{253,202,40}\fill[ezs](107.1,-199.9)--(109.4,-200.8)--(174.7,-30.6)--(172.5,-29.7)--cycle;
				\definecolor{ezs}{RGB}{253,200,38}\fill[ezs](109,-200.6)--(111.3,-201.5)--(176.7,-31.3)--(174.4,-30.4)--cycle;
				\definecolor{ezs}{RGB}{253,198,35}\fill[ezs](111,-201.4)--(113.2,-202.2)--(178.6,-32.1)--(176.3,-31.2)--cycle;
				\definecolor{ezs}{RGB}{253,196,31}\fill[ezs](112.9,-202.1)--(115.2,-203)--(180.5,-32.8)--(178.2,-31.9)--cycle;
				\definecolor{ezs}{RGB}{253,194,28}\fill[ezs](114.8,-202.8)--(117.1,-203.7)--(182.4,-33.5)--(180.1,-32.6)--cycle;
				\definecolor{ezs}{RGB}{253,192,26}\fill[ezs](116.7,-203.6)--(119,-204.5)--(184.3,-34.3)--(182,-33.4)--cycle;
				\definecolor{ezs}{RGB}{253,189,24}\fill[ezs](118.6,-204.3)--(120.9,-205.2)--(186.2,-35)--(183.9,-34.1)--cycle;
				\definecolor{ezs}{RGB}{253,186,20}\fill[ezs](120.5,-205)--(122.8,-205.9)--(188.1,-35.7)--(185.8,-34.8)--cycle;
				\definecolor{ezs}{RGB}{253,184,17}\fill[ezs](122.4,-205.8)--(124.7,-206.7)--(190,-36.5)--(187.8,-35.6)--cycle;
				\definecolor{ezs}{RGB}{252,183,15}\fill[ezs](124.3,-206.5)--(126.6,-207.4)--(192,-37.2)--(189.7,-36.3)--cycle;
				\definecolor{ezs}{RGB}{252,181,14}\fill[ezs](126.3,-207.2)--(128.5,-208.1)--(193.9,-37.9)--(191.6,-37.1)--cycle;
				\definecolor{ezs}{RGB}{252,179,13}\fill[ezs](128.2,-208)--(130.4,-208.9)--(195.8,-38.7)--(193.5,-37.8)--cycle;
				\definecolor{ezs}{RGB}{252,178,12}\fill[ezs](130.1,-208.7)--(132.4,-209.6)--(197.7,-39.4)--(195.4,-38.5)--cycle;
				\definecolor{ezs}{RGB}{251,178,10}\fill[ezs](132,-209.4)--(134.3,-210.3)--(199.6,-40.1)--(197.3,-39.3)--cycle;
				\definecolor{ezs}{RGB}{251,178,10}\fill[ezs](133.9,-210.2)--(136.2,-211.1)--(201.5,-40.9)--(199.2,-40)--cycle;
				\definecolor{ezs}{RGB}{250,177,11}\fill[ezs](135.8,-210.9)--(138.1,-211.8)--(203.4,-41.6)--(201.1,-40.7)--cycle;
				\definecolor{ezs}{RGB}{250,178,11}\fill[ezs](137.7,-211.6)--(140,-212.5)--(205.3,-42.3)--(203.1,-41.5)--cycle;
				\definecolor{ezs}{RGB}{248,180,10}\fill[ezs](139.6,-212.4)--(141.9,-213.3)--(207.3,-43.1)--(205,-42.2)--cycle;
				\definecolor{ezs}{RGB}{247,181,9}\fill[ezs](141.5,-213.1)--(143.8,-214)--(209.2,-43.8)--(206.9,-42.9)--cycle;
				\definecolor{ezs}{RGB}{246,182,9}\fill[ezs](143.5,-213.8)--(145.7,-214.7)--(211.1,-44.5)--(208.8,-43.7)--cycle;
				\definecolor{ezs}{RGB}{246,182,10}\fill[ezs](145.4,-214.6)--(147.7,-215.5)--(213,-45.3)--(210.7,-44.4)--cycle;
			\end{scope}
			% chữ EDUZY (Times đậm, làm dày nét bằng 5 lớp lệch nhẹ; căn đúng khung chữ logo gốc)
			\foreach \dx/\dy in {-0.7/0,0.7/0,0/-0.3,0/0.3,0/0}
			\node[anchor=center,inner sep=0pt,outer sep=0pt,text=ezText,align=center] at (109.6+\dx,-193.6+\dy)
			{\resizebox{1.01cm}{0.22cm}{\fontencoding{T1}\fontfamily{ptm}\fontseries{b}\selectfont EDUZY}};
		\end{tikzpicture}%
	}%
}
%====================== THIẾT LẬP THÂN ĐỀ ======================
%--- Môi trường tự luận {bt}: ex_test 3.3.4 chưa khai báo sẵn nên tạo ở đây
%    (ex_test đã tự móc \fixshowans và \listenumerate cho bt)
\makeatletter
\@ifundefined{bt}{\newtheorem{bt}{\color{gold1}\namebt}}{}
\makeatother
\renewcommand{\sectionmark}[1]{\markright{\thesection\ #1}}
%--- Khung tiêu đề (logo + tối đa 4 dòng; dòng 3, 4 để trống {} thì bỏ qua)
%  Đề thi:    \khungtieude{THI THỬ THPT QUỐC GIA NĂM 2026}{Môn Toán -- Đề số 21}
%                         {(Đề thi gồm: ...)}{Thời gian làm bài: 90 phút (không kể thời gian phát đề).}
%  Chuyên đề: \khungtieude{CHUYÊN ĐỀ 3 -- DÃY SỐ VÀ CÔNG THỨC TRUY HỒI}
%                         {Bồi dưỡng học sinh giỏi Toán THPT}{...}{}
\newcommand{\khungtieude}[4]{%
	\begin{tcolorbox}[colframe=gold1,colback=gold5!45,boxrule=0.5mm]
		\begin{minipage}[c]{3cm}
			\centerline{\eduzylogo[2.4cm]}
		\end{minipage}\hspace{0.5cm}
		\begin{minipage}[c]{12.5cm}
			\vskip 0.15 cm
			\centerline{\ind{#1}}
			\vskip 0.15 cm
			\centerline{\textbf{#2}}
			\if\relax\detokenize{#3}\relax\else
				\vskip 0.15 cm
				\centerline{\textit{#3}}
			\fi
			\if\relax\detokenize{#4}\relax\else
				\vskip 0.15 cm
				\centerline{\textit{#4}}
			\fi
			\vskip 0.15 cm
		\end{minipage}
	\end{tcolorbox}
}
%====================== BỔ SUNG CHO BÀI HỌC, CHUYÊN ĐỀ ======================
\usepackage[shortlabels]{enumitem}
\usepackage{tabularx}
%--- Danh sách gọn: \begin{enumerate}[gon] ... \end{enumerate}
\SetEnumitemKey{gon}{leftmargin=*,itemsep=1pt,parsep=0pt,topsep=3pt}
%--- Khung lý thuyết: \begin{khung}{} ... \end{khung}; tiêu đề chỉ dùng khi thật ngắn ("Mẹo", "Lỗi thường gặp")
\newtcolorbox{khung}[1]{breakable,enhanced,colframe=gold1,colback=gold5!45,
	boxrule=0.4mm,fonttitle=\bfseries\sffamily,coltitle=white,colbacktitle=gold1,
	attach boxed title to top left={xshift=6mm,yshift=-3mm},
	boxed title style={arc=1mm,boxrule=0.3mm},title={#1}}
%--- Tiêu đề mục nhỏ trong một dạng: \muc{Dấu hiệu}, \muc{Phương pháp}, \muc{Bài tự luyện}
\makeatletter
\newcommand{\muc}[1]{\par\medskip\noindent\indam{#1}\par\nobreak\smallskip\@afterindentfalse\@afterheading}
\makeatother
%--- Ô công thức cần nhớ (căn giữa, viền vàng, vừa khít nội dung)
%     \begin{center}\congthuc{$Q(n)=(n-1)!$}\end{center}
\newtcbox{\congthuc}{on line,colframe=gold1,colback=gold5!45,boxrule=0.5pt,arc=2pt,
	left=8pt,right=8pt,top=4pt,bottom=4pt}
%--- Công tắc lời giải bài tự luyện (ví dụ mẫu luôn dùng \loigiai)
%     \loigiaiTL[{đáp số}]{lời giải đầy đủ}
%     \TLfull -> in lời giải đầy đủ;  \TLdapso -> chỉ in một dòng "Đáp số";  \TLan -> ẩn hoàn toàn
%     (\CheDo là dethi thì mọi lời giải và đáp số đều ẩn)
\makeatletter
\newif\ifTLhien   \TLhientrue
\newif\ifTLdapso  \TLdapsofalse
\newcommand{\TLfull}{\TLhientrue\TLdapsofalse}
\newcommand{\TLdapso}{\TLhientrue\TLdapsotrue}
\newcommand{\TLan}{\TLhienfalse}
\newcommand{\loigiaiTL}[2][]{%
	\ifTLhien
		\ifTLdapso
			\def\@tlarg{#1}%
			\ifx\@tlarg\@empty\else
				\ifdefstring{\kindOPTN}{dethi}{}{\par\noindent{\bfseries\color{gold1}Đáp số:} #1\par}%
			\fi
		\else
			\loigiai{#2}%
		\fi
	\fi}
\makeatother
\TLdapso% mặc định: chỉ in đáp số bài tự luyện
%====================== BẢNG ĐÚNG/SAI (\choiceTFt) ======================
%  Hàng tiêu đề nền nâu đồng chữ trắng (Phát biểu | Đ | S), các hàng tô xen kẽ kem / trắng,
%  đường kẻ vàng kim, nhãn a) b) c) d) màu nâu đồng. Bản có đáp án: cột được chọn có vòng tròn
%  nâu đồng với dấu tích trắng (câu đúng tích ở cột Đúng, câu sai tích ở cột Sai).
%  Tiêu đề ghi đủ "Đúng | Sai": \OPTN{viettat=0} sau \input (cột Đúng/Sai rộng ra, phát biểu hẹp lại).
%  Dùng bảng gốc của ex_test: \def\BangDSgoc{1} trước %============================================================
%  PREAMBLE EDUZY v2 - dùng chung cho đề thi, bài học, chuyên đề (gói ex_test)
%  Khác bản preamble-eduzy.tex ở đầu trang, chân trang: dải nền nâu đồng tràn ngang,
%  logo + Eduzy.com.vn, ô chữ vàng kem cạnh xiên, chân trang "Trang (số)";
%  logo Eduzy chìm giữa mọi trang (tắt: \logochimtat, bật lại: \logochimbat).
%  Cách dùng ở file chính (xem mau-chuan.tex, mau-chuyen-de.tex):
%      \def\CheDo{loigiai}          % loigiai | dethi | solcolor | color | book
%      \def\tenchuong{<dòng chữ trên đầu trang>}
%      \def\thuonghieu{<chữ ở chân trang lẻ>}   % không khai báo thì dùng mặc định
%      \def\tacgia{<chữ ở chân trang chẵn>}      % không khai báo thì dùng mặc định
%      \def\khauhieu{<chữ cạnh logo ở đầu trang>}  % tùy chọn, mặc định Eduzy.com.vn
%      \input{preamble-eduzy-v2}
%      \TLdapso                      % tài liệu: \TLfull | \TLdapso | \TLan
%      \begin{document} ... \end{document}
%  Logo Eduzy vẽ bằng TikZ ngay trong file (\eduzylogo), không cần file ảnh.
%  Biên dịch: pdfLaTeX (máy có gói vietnam) hoặc XeLaTeX.
%============================================================
\documentclass[12pt,a4paper,twoside]{article}%<-- twoside: bat dau trang/chan trang chan-le
\usepackage{iftex}
\usepackage{amsmath,amssymb,tikz,mathrsfs,tkz-tab}
\usepackage{graphicx}
\usepackage{pgfplots}% đồ thị hàm số
\pgfplotsset{compat=1.17}
%\usepackage{mathptmx}
\usepackage{mathpazo}
\ifXeTeX% XeLaTeX: font chữ có đủ dấu tiếng Việt, giữ nguyên font toán mathpazo
	\usepackage[no-math]{fontspec}
	\setmainfont{TeX Gyre Pagella}
	\newfontfamily\EZqag{TeX Gyre Adventor}[NFSSFamily=qag,Ligatures=TeX]
\fi
\usepackage[many]{tcolorbox}
\usetikzlibrary{math,through,calc,intersections,angles,quotes,shapes,shapes.geometric,arrows,patterns,matrix,chains,arrows.meta,decorations.shapes,decorations.fractals,decorations.markings,shadows,decorations.pathreplacing}
\usepackage[top=1.6cm, bottom=1.8cm, left=1.8cm, right=1.8cm]{geometry}% v2: top 1.6cm (bản cũ 1.5cm) để chữ không sát dải đầu trang
\providecommand{\CheDo}{loigiai}% loigiai | dethi | solcolor | color | book
\expandafter\usepackage\expandafter[\CheDo]{ex_test}
% \tenchuong: khai báo ở file chính (trước \input) hoặc \renewcommand sau \input
%============================================================
%  [GỘP TỪ eduzy-de.tex] EDUZY - MÀU NÂU ĐỒNG / VÀNG KIM
%  + ĐẦU TRANG, CHÂN TRANG KIỂU DẢI NGANG (v2), ĐỐI XỨNG TRANG CHẴN / TRANG LẺ.
%  (không cần file eduzy-de.tex riêng nữa)
%  Khối này phải đặt SAU \usepackage[...]{ex_test} và SAU \def\tenchuong{...}
%============================================================
\usepackage{lastpage}
\usepackage{needspace}
%--- DÀN TRANG ĐỀ THI: tách câu (THÊM 05/11/2026; chỉ thêm lệnh mới, KHÔNG đổi hành vi mặc định) ---
%  Hai công tắc, mặc định tắt (= preamble chạy y như cũ). Bật bằng cách đặt TRƯỚC \input{preamble-eduzy-v2}
%  trong file đề (tools/dan_trang/dan_trang.py tự ghi khi đặt \NeedspaceCau):
%    \def\TachCauDS{1}   câu Đúng/Sai được tách: câu dẫn (kèm hình) ở cuối trang, bảng Phát biểu sang trang sau
%                        (câu dẫn và hình không tách nhau, bảng không tách giữa hai trang)
%    \def\TachCauLG{1}   bản lời giải (\CheDo khác dethi): phần đề của câu cao hơn 1/3 trang được phép tách
%  Lệnh dùng trong file đề (giá trị do tools/dan_trang/dan_trang.py đo ở bản dethi):
%    \NeedspaceCau[<cao câu dẫn + hình>]{<cao cả khối câu>}   ngay trên \begin{ex} (câu đầu phần: trên tiêu đề phần)
%        công tắc tắt: đúng bằng \Needspace{<cao cả khối câu>} như cũ;
%        \TachCauDS=1 và có [..] (câu Đúng/Sai được chọn tách): chỉ giữ câu dẫn + hình;
%        \TachCauLG=1 ở bản lời giải: đòi tối đa 1/3 \textheight.
%    \NeedspaceBang{<cao bảng Phát biểu>}%   ngay trước \choiceTFt: \TachCauDS=1 thì giữ trọn bảng trên một trang;
%        công tắc tắt: không làm gì.
\providecommand{\TachCauDS}{0}
\providecommand{\TachCauLG}{0}
\newlength{\EZnsCau}
\newcommand{\NeedspaceCau}[2][]{%
	\setlength{\EZnsCau}{#2}%
	\ifnum\TachCauDS=1 \if\relax\detokenize{#1}\relax\else\setlength{\EZnsCau}{#1}\fi\fi
	\ifnum\TachCauLG=1 \ifdefstring{\kindOPTN}{dethi}{}{\ifdim\EZnsCau>0.3333\textheight\setlength{\EZnsCau}{0.3333\textheight}\fi}\fi
	\Needspace{\EZnsCau}}
\newcommand{\NeedspaceBang}[1]{\ifnum\TachCauDS=1 \Needspace{#1}\fi}

%--- Thang màu: đậm -> nhạt
\definecolor{gold1}{RGB}{126,96,0}    % nâu đồng đậm  #7E6000
\definecolor{gold2}{RGB}{192,144,3}   % vàng kim      #C09003
\definecolor{gold3}{RGB}{255,215,99}  % vàng nhạt     #FFD763
\definecolor{gold4}{RGB}{255,231,153} % vàng rất nhạt #FFE799
\definecolor{gold5}{RGB}{255,243,205} % vàng kem      #FFF3CD
%--- Ghi đè tông xanh/đỏ cũ
\definecolor{myblue}{RGB}{126,96,0}
\colorlet{blue}{gold1}
\colorlet{cyan}{gold2}
\colorlet{teal}{gold2}

%--- Tên chương mặc định (file đề nên tự khai báo lại)
\providecommand{\tenchuong}{}
%--- Chân trang
\providecommand{\thuonghieu}{TOÁN 12 -- Eduzy.com.vn}
\providecommand{\tacgia}{Tô Gia Bảo -- 0353703783}

%--- Chữ cạnh logo ở đầu trang
\providecommand{\khauhieu}{Eduzy.com.vn}

%====================== THÔNG SỐ ĐẦU TRANG, CHÂN TRANG (v2) ======================
%  Dải nền đậm tràn ngang trang; ô chữ vàng kem có cạnh xiên, kèm 2 vạch xiên chuyển màu.
%  Đầu trang lẻ : logo + \khauhieu (trái)   | ô chữ \tenchuong (phải)
%  Chân trang lẻ: ô chữ \thuonghieu (trái)  | Trang (số) (phải)
%  Trang chẵn soi gương trang lẻ; chân trang chẵn ghi \tacgia.
%--- Màu: chỉ dùng thang vàng kim Eduzy
\colorlet{ezNen}{gold1}% dải nền đậm
\colorlet{ezVach1}{gold3}% vạch xiên sát dải nền (đậm hơn)
\colorlet{ezVach2}{gold4}% vạch xiên sát ô chữ (nhạt hơn)
\colorlet{ezO}{gold5}% ô chữ vàng kem
%--- Kích thước
\newlength{\EZdai}\setlength{\EZdai}{1.1cm}% chiều cao dải nền
\newlength{\EZo}\setlength{\EZo}{0.68cm}% chiều cao ô chữ
\newlength{\EZxien}\setlength{\EZxien}{0.55cm}% độ lệch ngang của cạnh xiên
\newlength{\EZvach}\setlength{\EZvach}{0.75cm}% bề rộng mỗi vạch xiên
\newlength{\EZkhe}\setlength{\EZkhe}{0.12cm}% khe trắng giữa vạch và ô chữ
\newlength{\EZdem}\setlength{\EZdem}{0.4cm}% khoảng đệm từ chữ tới cạnh xiên
\newlength{\EZle}\setlength{\EZle}{1.8cm}% lề trái/phải của chữ (bằng lề thân bài)
\newlength{\EZrongmin}\setlength{\EZrongmin}{4.5cm}% ô chữ rộng tối thiểu
\newlength{\EZrong}% bề rộng chữ trong ô (tính tự động)
%--- Font
\newcommand{\EZfontO}{\fontsize{10.5pt}{12pt}\selectfont\bfseries}% chữ trong ô
\newcommand{\EZfontKH}{\fontfamily{qag}\fontsize{11pt}{12pt}\selectfont\bfseries}% \khauhieu cạnh logo
\newcommand{\EZfontTrang}{\fontsize{11.5pt}{12pt}\selectfont\bfseries}% chữ "Trang" và số trang
%--- Đo bề rộng chữ trong ô
\newcommand{\EZdorong}[1]{%
	\settowidth{\EZrong}{\EZfontO #1}%
	\ifdim\EZrong<\EZrongmin\setlength{\EZrong}{\EZrongmin}\fi}
%--- Logo + chữ trắng \khauhieu (dựng sẵn một lần khi vào \begin{document})
\newsavebox{\EZlogobox}
\AtBeginDocument{\sbox{\EZlogobox}{%
		\raisebox{-0.5\height}{\eduzylogo[1cm]}\hspace{0.15cm}%
		\raisebox{-0.5\height}{\color{white}\EZfontKH\khauhieu}}}
%--- Hai vạch xiên + khe trắng. #1: điểm gốc ở mép dưới, #2: hướng (1 sang phải, -1 sang trái),
%    #3: hướng nghiêng (1 kiểu "/", -1 kiểu "\")
\newcommand{\EZvachxien}[3]{%
	\fill[white] (#1) --++({(#3)*\EZxien},\EZo) --++({(#2)*\EZkhe},0) --++({-(#3)*\EZxien},-\EZo) --cycle;
	\fill[ezVach2] ([xshift=(#2)*\EZkhe]#1) --++({(#3)*\EZxien},\EZo) --++({(#2)*\EZvach},0) --++({-(#3)*\EZxien},-\EZo) --cycle;
	\fill[ezVach1] ([xshift=(#2)*(\EZkhe+\EZvach)]#1) --++({(#3)*\EZxien},\EZo) --++({(#2)*\EZvach},0) --++({-(#3)*\EZxien},-\EZo) --cycle;
	\draw[white,line width=0.6pt] ([xshift=(#2)*(\EZkhe+\EZvach)]#1) --++({(#3)*\EZxien},\EZo);
}
%--- Chữ trong ô (căn giữa theo chiều cao ô, không lệch vì dấu tiếng Việt)
\tikzset{ezchu/.style={inner sep=0pt,text=ezNen,text height=0.72em,text depth=0.22em}}

%====================== LOGO CHÌM GIỮA TRANG ======================
%  Vẽ cùng đầu trang nên nằm dưới chữ của thân bài, có trên mọi trang.
%  Tắt: \logochimtat (từ trang đó trở đi), bật lại: \logochimbat.
%  Đổi cỡ, độ đậm: khai báo trước \input, ví dụ \def\EZchimcao{9cm} \def\EZchimdo{0.1}
\providecommand{\EZchimcao}{10cm}% chiều cao logo chìm
\providecommand{\EZchimdo}{0.12}% độ đậm (0 = trong suốt, 1 = rõ hẳn)
\newif\ifEZchim \EZchimtrue
\newcommand{\logochimtat}{\EZchimfalse}
\newcommand{\logochimbat}{\EZchimtrue}
%  Dựng sẵn logo mờ một lần khi vào \begin{document}; transparency group làm mờ cả logo
%  một lượt nên không lộ vệt ở chỗ các dải màu chồng nhau.
\newsavebox{\EZchimtho}
\newsavebox{\EZchimbox}
\AtBeginDocument{%
	\sbox{\EZchimtho}{\eduzylogo[\EZchimcao]}%
	\sbox{\EZchimbox}{\begin{tikzpicture}
			\begin{scope}[transparency group,opacity=\EZchimdo]
				\node[inner sep=0pt] {\usebox{\EZchimtho}};
			\end{scope}
	\end{tikzpicture}}}
\newcommand{\EZvechim}{%
	\ifEZchim
		\node[inner sep=0pt] at (current page.center) {\usebox{\EZchimbox}};
	\fi}

%====================== ĐẦU TRANG - TRANG LẺ ======================
%  (logo bên trái, ô chữ \tenchuong bên phải)
\newcommand{\myfancyheadOdd}{%
	\EZdorong{\tenchuong}%
	\begin{tikzpicture}[remember picture,overlay]
		\coordinate (T) at (current page.north west);
		\coordinate (P) at (current page.north east);
		%--- logo chìm giữa trang
		\EZvechim
		%--- dải nền
		\fill[ezNen] (T) rectangle ([yshift=-\EZdai]P);
		%--- ô chữ: X là góc dưới bên trái, cạnh trái xiên kiểu "/"
		\coordinate (X) at ([xshift=-\EZle-\EZrong-\EZdem-\EZxien,yshift=-\EZdai]P);
		\fill[ezO] (X) --++(\EZxien,\EZo) -- ([yshift=-\EZdai+\EZo]P) -- ([yshift=-\EZdai]P) --cycle;
		\EZvachxien{X}{-1}{1}
		\node[ezchu,anchor=east] at ([xshift=-\EZle,yshift=-\EZdai+0.5*\EZo]P) {\EZfontO\tenchuong};
		%--- logo + tên
		\node[inner sep=0pt,anchor=west] at ([xshift=\EZle,yshift=-0.5*\EZdai]T) {\usebox{\EZlogobox}};
	\end{tikzpicture}%
}

%====================== ĐẦU TRANG - TRANG CHẴN =====================
%  (soi gương: ô chữ \tenchuong bên trái, logo bên phải)
\newcommand{\myfancyheadEven}{%
	\EZdorong{\tenchuong}%
	\begin{tikzpicture}[remember picture,overlay]
		\coordinate (T) at (current page.north west);
		\coordinate (P) at (current page.north east);
		%--- logo chìm giữa trang
		\EZvechim
		%--- dải nền
		\fill[ezNen] (T) rectangle ([yshift=-\EZdai]P);
		%--- ô chữ: X là góc dưới bên phải, cạnh phải xiên kiểu "\"
		\coordinate (X) at ([xshift=\EZle+\EZrong+\EZdem+\EZxien,yshift=-\EZdai]T);
		\fill[ezO] (X) --++(-\EZxien,\EZo) -- ([yshift=-\EZdai+\EZo]T) -- ([yshift=-\EZdai]T) --cycle;
		\EZvachxien{X}{1}{-1}
		\node[ezchu,anchor=west] at ([xshift=\EZle,yshift=-\EZdai+0.5*\EZo]T) {\EZfontO\tenchuong};
		%--- logo + tên
		\node[inner sep=0pt,anchor=east] at ([xshift=-\EZle,yshift=-0.5*\EZdai]P) {\usebox{\EZlogobox}};
	\end{tikzpicture}%
}

%====================== CHÂN TRANG - TRANG LẺ ======================
%  (ô chữ \thuonghieu bên trái, "Trang (số)" bên phải)
\newcommand{\myfancyfootOdd}{%
	\EZdorong{\thuonghieu}%
	\begin{tikzpicture}[remember picture,overlay]
		\coordinate (T) at (current page.south west);
		\coordinate (P) at (current page.south east);
		%--- dải nền
		\fill[ezNen] (T) rectangle ([yshift=\EZdai]P);
		%--- ô chữ sát mép trên dải: X là góc dưới bên phải, cạnh phải xiên kiểu "/"
		\coordinate (X) at ([xshift=\EZle+\EZrong+\EZdem,yshift=\EZdai-\EZo]T);
		\fill[ezO] ([yshift=\EZdai-\EZo]T) -- (X) --++(\EZxien,\EZo) -- ([yshift=\EZdai]T) --cycle;
		\EZvachxien{X}{1}{1}
		\node[ezchu,anchor=west] at ([xshift=\EZle,yshift=\EZdai-0.5*\EZo]T) {\EZfontO\thuonghieu};
		%--- số trang
		\node[circle,draw=white,line width=0.9pt,text=white,inner sep=0pt,minimum size=0.62cm,anchor=east]
		(so) at ([xshift=-\EZle,yshift=0.5*\EZdai]P) {\EZfontTrang\thepage};
		\node[inner sep=0pt,text=white,anchor=base east] (tr) at ([xshift=-0.15cm]so.base west) {\EZfontTrang Trang};
	\end{tikzpicture}%
}

%====================== CHÂN TRANG - TRANG CHẴN ====================
%  (soi gương: "Trang (số)" bên trái, ô chữ \tacgia bên phải)
\newcommand{\myfancyfootEven}{%
	\EZdorong{\tacgia}%
	\begin{tikzpicture}[remember picture,overlay]
		\coordinate (T) at (current page.south west);
		\coordinate (P) at (current page.south east);
		%--- dải nền
		\fill[ezNen] (T) rectangle ([yshift=\EZdai]P);
		%--- ô chữ sát mép trên dải: X là góc dưới bên trái, cạnh trái xiên kiểu "\"
		\coordinate (X) at ([xshift=-\EZle-\EZrong-\EZdem,yshift=\EZdai-\EZo]P);
		\fill[ezO] ([yshift=\EZdai-\EZo]P) -- (X) --++(-\EZxien,\EZo) -- ([yshift=\EZdai]P) --cycle;
		\EZvachxien{X}{-1}{-1}
		\node[ezchu,anchor=east] at ([xshift=-\EZle,yshift=\EZdai-0.5*\EZo]P) {\EZfontO\tacgia};
		%--- số trang
		%    (vòng tròn ẩn so0 nằm giữa dải để lấy đường chân chữ, giống trang lẻ)
		\node[circle,inner sep=0pt,minimum size=0.62cm] (so0) at ([yshift=0.5*\EZdai]T) {\EZfontTrang\phantom{\thepage}};
		\node[inner sep=0pt,text=white,anchor=base west] (tr) at ($(so0.base -| T)+(\EZle,0)$) {\EZfontTrang Trang};
		\node[circle,draw=white,line width=0.9pt,text=white,inner sep=0pt,minimum size=0.62cm,anchor=base west]
		(so) at ([xshift=0.15cm]tr.base east) {\EZfontTrang\thepage};
	\end{tikzpicture}%
}

%--- Giữ tương thích với các file cũ đang gọi \myfancyhead, \myfancyfoot
\newcommand{\myfancyhead}{\myfancyheadOdd}
\newcommand{\myfancyfoot}{\myfancyfootOdd}

%======================= GẮN VÀO TRANG =============================
\pagestyle{fancy}
\fancyhf{}
\fancyhead[LO]{\myfancyheadOdd}
\fancyhead[RE]{\myfancyheadEven}
\fancyfoot[LO]{\myfancyfootOdd}
\fancyfoot[RE]{\myfancyfootEven}
\renewcommand{\headrulewidth}{0pt}
\renewcommand{\footrulewidth}{0pt}
\fancypagestyle{plain}{%
	\fancyhf{}
	\fancyhead[LO]{\myfancyheadOdd}
	\fancyhead[RE]{\myfancyheadEven}
	\fancyfoot[LO]{\myfancyfootOdd}
	\fancyfoot[RE]{\myfancyfootEven}
	\renewcommand{\headrulewidth}{0pt}
	\renewcommand{\footrulewidth}{0pt}
}
%=================== HẾT PHẦN GỘP TỪ eduzy-de.tex ===================
\newcommand{\indam}[1]{
	{\fontsize{12pt}{9pt}\sffamily \selectfont\bfseries\color{gold1} #1}
}
\newcommand{\ind}[1]{
	{\fontsize{14pt}{9pt}\fontfamily{qag} \selectfont\bfseries\color{gold1} #1}
}
\renewcommand{\FalseEX}{\stepcounter{dapan}\textbf{\selectfont\color{gold1}\Alph{dapan}.}\ignorespaces}
\renewtheorem{ex}{\color{gold1}Câu}
\renewcommand{\baselinestretch}{1.1}
\newcommand{\hoac}[1]{
	\left[\begin{aligned}#1\end{aligned}\right.}
\newcommand{\heva}[1]{
	\left\{\begin{aligned}#1\end{aligned}\right.}
\usepackage{esvect}
\def\vec{\vv}
\def\vt{\vv}
\def\overrightarrow{\vv}
%====================== LOGO EDUZY VẼ BẰNG TIKZ ======================
%  Đường viền dò từ logo gốc; dải màu vàng kim đo trực tiếp trên ảnh gốc
%  (nghiêng 21 độ, tô bằng 80 dải hẹp nên hiển thị giống nhau trên mọi trình xem PDF).
%  Dùng: \eduzylogo          -> cao 2.4cm (mặc định)
%        \eduzylogo[1.5cm]   -> cao 1.5cm
\definecolor{ezText}{RGB}{245,200,46}
\newcommand{\eduzylogo}[1][2.4cm]{%
	\resizebox{!}{#1}{%
		\begin{tikzpicture}[x=0.01cm,y=0.01cm]
			\path[use as bounding box] (0,0) rectangle (224,-224);
			% mũ tốt nghiệp + vương miện
			\begin{scope}
				\clip
				(119.4,-176.1)..controls(118.9,-176.1)and(116.8,-176)..(114.6,-176)
				..controls(101.4,-176)and(85.8,-175.9)..(83.7,-175.7)
				..controls(81.8,-175.6)and(81,-175.6)..(80.2,-175.7)
				..controls(78.8,-175.9)and(77.7,-175.9)..(76.2,-175.7)
				..controls(74.9,-175.5)and(70.8,-175.5)..(64.3,-175.7)
				..controls(59.1,-175.9)and(58.3,-175.7)..(58.6,-173.7)
				..controls(58.6,-172.8)and(59.6,-169.6)..(60.3,-168)
				..controls(60.8,-166.9)and(61.4,-165.8)..(62.3,-165)
				..controls(63.1,-164.2)and(64.3,-163.4)..(65,-163.2)
				..controls(65.5,-163.1)and(72.3,-163)..(76.6,-163.1)
				..controls(77.9,-163.1)and(80.1,-163.1)..(81.4,-163)
				..controls(82.8,-162.9)and(86.3,-162.9)..(89.2,-162.9)
				..controls(93,-162.9)and(94.9,-162.8)..(96,-162.7)
				..controls(97.2,-162.5)and(101,-162.5)..(114,-162.4)
				..controls(123,-162.3)and(134.9,-162.2)..(140.2,-162.1)
				..controls(150.1,-162)and(152.2,-162.1)..(153.6,-162.4)
				..controls(154.4,-162.6)and(157.1,-163.8)..(158.1,-164.4)
				..controls(158.8,-164.8)and(159.7,-165.7)..(159.9,-166.3)
				..controls(160,-166.5)and(160.2,-167.1)..(160.2,-167.6)
				..controls(160.3,-168)and(160.6,-169.4)..(160.9,-170.6)
				..controls(161.2,-171.7)and(161.5,-173.1)..(161.6,-173.6)
				..controls(161.8,-174.4)and(161.8,-174.5)..(161.5,-174.9)
				..controls(161.2,-175.6)and(160.8,-175.9)..(159.9,-176)
				..controls(159.2,-176.1)and(154.1,-176.2)..(151.5,-176.2)
				..controls(148.1,-176.1)and(122.6,-176.1)..(121.6,-176.2)
				..controls(120.8,-176.2)and(119.8,-176.2)..(119.4,-176.1)--cycle
				(130.4,-160)..controls(129,-160)and(119.8,-160)..(110.1,-159.9)
				..controls(100.3,-159.8)and(88.5,-159.7)..(83.9,-159.6)
				..controls(77.8,-159.5)and(75.3,-159.5)..(74.5,-159.6)
				..controls(73.9,-159.7)and(72.5,-159.8)..(71.3,-159.8)
				..controls(69.6,-159.8)and(69,-159.8)..(68.5,-159.6)
				..controls(67.6,-159.3)and(66.8,-158.6)..(66.2,-157.6)
				..controls(66,-157.3)and(65.6,-156.5)..(65.3,-156.1)
				..controls(64.6,-155.1)and(64.5,-154.9)..(63.7,-152.9)
				..controls(63.2,-151.6)and(62.9,-150.9)..(61.8,-149.6)
				..controls(61.5,-149.2)and(61.1,-148.5)..(60.9,-147.9)
				..controls(60.7,-147.4)and(60.3,-146.6)..(60.1,-146.1)
				..controls(59.9,-145.5)and(59.5,-144.9)..(59,-144.3)
				..controls(58.5,-143.7)and(58.2,-143.2)..(57.9,-142.5)
				..controls(57.7,-141.9)and(57.3,-141.1)..(57,-140.7)
				..controls(56.6,-140.3)and(56.1,-139.3)..(55.8,-138.5)
				..controls(55.5,-137.8)and(55.1,-136.8)..(54.9,-136.5)
				..controls(54.7,-136.1)and(54.3,-135.2)..(53.9,-134.5)
				..controls(53.6,-133.8)and(53,-132.7)..(52.6,-132)
				..controls(51.6,-130.5)and(51.4,-130.2)..(51.1,-129)
				..controls(50.6,-127.5)and(50.6,-127.4)..(49.7,-126.3)
				..controls(49.2,-125.7)and(48.6,-124.8)..(48.3,-124.2)
				..controls(48,-123.6)and(47.5,-122.7)..(47.3,-122.2)
				..controls(47,-121.8)and(46.5,-120.7)..(46.1,-119.8)
				..controls(45.8,-118.9)and(45.2,-117.8)..(44.9,-117.2)
				..controls(44.2,-116.1)and(43.3,-114.4)..(42.3,-112.2)
				..controls(41.8,-111.3)and(41.3,-110.4)..(41.1,-110.1)
				..controls(40.5,-109.4)and(40.3,-109)..(40,-107.8)
				..controls(39.8,-107)and(39.6,-106.5)..(39.2,-106)
				..controls(38.9,-105.6)and(38.4,-104.7)..(38,-104)
				..controls(37.7,-103.4)and(37.2,-102.5)..(37,-102.1)
				..controls(36.7,-101.7)and(36.3,-100.8)..(36,-100.1)
				..controls(35.6,-98.9)and(35.3,-98.4)..(34.6,-97.6)
				..controls(34.2,-97.1)and(34.2,-96.3)..(34.6,-96)
				..controls(35,-95.7)and(35.5,-95.9)..(36.2,-96.7)
				..controls(36.5,-97)and(37.2,-97.5)..(37.6,-97.7)..controls(38,-98)and(38.6,-98.4)..(39.1,-98.9)
				..controls(39.6,-99.4)and(40.3,-100)..(40.6,-100.3)
				..controls(41.3,-100.8)and(42.4,-101.8)..(43.7,-102.9)
				..controls(44,-103.2)and(44.6,-103.7)..(45,-103.9)
				..controls(45.4,-104.2)and(45.9,-104.6)..(46.1,-104.8)
				..controls(47.1,-105.9)and(48.8,-107.1)..(50.5,-107.9)
				..controls(51.8,-108.6)and(52.3,-109)..(53.4,-110)
				..controls(54.8,-111.2)and(55.6,-111.6)..(57,-112.1)
				..controls(57.5,-112.3)and(58.5,-112.8)..(59.4,-113.3)
				..controls(60.2,-113.7)and(61.3,-114.3)..(61.9,-114.5)
				..controls(62.4,-114.7)and(63.3,-115.1)..(63.8,-115.3)
				..controls(64.4,-115.6)and(65.4,-115.9)..(66.2,-116.1)
				..controls(67,-116.3)and(68.4,-116.6)..(69.2,-116.9)
				..controls(72.3,-117.9)and(72.8,-118)..(78,-118.1)
				..controls(84,-118.3)and(86.1,-118.1)..(88.1,-117.1)
				..controls(89.4,-116.5)and(93,-115)..(93.7,-114.9)
				..controls(94.8,-114.6)and(95.4,-114.2)..(95.9,-113.4)
				..controls(96.2,-113)and(96.7,-112.5)..(97.1,-112.2)
				..controls(98,-111.6)and(101.3,-108.4)..(102.1,-107.2)
				..controls(102.4,-106.8)and(102.9,-106.2)..(103.2,-105.8)
				..controls(103.5,-105.5)and(104,-104.8)..(104.2,-104.2)
				..controls(104.5,-103.6)and(105,-102.8)..(105.3,-102.4)
				..controls(105.6,-102)and(106.1,-101.2)..(106.5,-100.6)
				..controls(107.1,-99.5)and(108.6,-97.5)..(109.5,-96.5)
				..controls(109.9,-96.1)and(110.1,-96)..(110.4,-96)
				..controls(111.1,-96)and(111.6,-96.5)..(111.8,-97.6)
				..controls(112,-98.6)and(112.3,-99.3)..(112.8,-99.9)
				..controls(113.7,-100.9)and(115.2,-103.1)..(115.7,-104.2)
				..controls(116.1,-105)and(116.5,-105.5)..(116.9,-106)
				..controls(117.3,-106.4)and(117.8,-106.9)..(118,-107.3)
				..controls(118.2,-107.6)and(118.6,-108)..(119,-108.3)
				..controls(119.6,-108.8)and(121.1,-110.3)..(121.8,-111.2)
				..controls(122.1,-111.7)and(122.5,-112)..(123.4,-112.4)
				..controls(124.8,-113.2)and(125.8,-113.8)..(127,-114.6)
				..controls(127.5,-114.9)and(128.1,-115.3)..(128.2,-115.4)
				..controls(128.4,-115.5)and(129.3,-115.8)..(130.2,-116)
				..controls(131.1,-116.2)and(132.5,-116.6)..(133.3,-116.9)
				..controls(135.7,-117.9)and(136.1,-117.9)..(142.5,-117.8)
				..controls(145.8,-117.8)and(148.4,-117.7)..(148.8,-117.6)
				..controls(149.2,-117.6)and(150.2,-117.2)..(151,-116.8)
				..controls(152,-116.4)and(152.8,-116)..(153.6,-115.9)
				..controls(154.5,-115.7)and(155.2,-115.4)..(156.3,-114.9)
				..controls(157.1,-114.5)and(158.2,-114)..(158.7,-113.8)
				..controls(159.1,-113.7)and(159.8,-113.3)..(160.2,-113)
				..controls(160.6,-112.7)and(161.3,-112.4)..(161.7,-112.2)
				..controls(162.7,-111.9)and(163.2,-111.7)..(164.1,-111.1)
				..controls(164.5,-110.8)and(165.3,-110.4)..(165.9,-110)
				..controls(166.5,-109.7)and(167.3,-109.2)..(167.6,-108.9)
				..controls(167.9,-108.6)and(168.6,-108.2)..(169,-108)
				..controls(169.4,-107.8)and(170.2,-107.3)..(170.7,-106.9)
				..controls(171.2,-106.5)and(172.1,-105.8)..(172.7,-105.4)
				..controls(173.2,-105)and(174.1,-104.4)..(174.5,-104)
				..controls(174.9,-103.7)and(175.7,-103.1)..(176.2,-102.6)
				..controls(176.7,-102.2)and(177.4,-101.5)..(177.9,-101.2)
				..controls(178.3,-100.8)and(178.8,-100.4)..(179,-100.2)
				..controls(179.2,-100)and(179.6,-99.6)..(180,-99.3)
				..controls(180.4,-99.1)and(180.9,-98.6)..(181.2,-98.4)
				..controls(181.5,-98.1)and(182.1,-97.7)..(182.5,-97.5)
				..controls(183.2,-97.2)and(183.5,-96.9)..(184.3,-95.9)
				..controls(184.6,-95.6)and(184.6,-95.6)..(184.8,-95.7)
				..controls(185,-95.8)and(185.1,-96)..(185.1,-96.2)
				..controls(185.1,-96.5)and(185,-96.7)..(184.6,-97.2)
				..controls(184.1,-97.8)and(184,-97.9)..(183.9,-98.8)
				..controls(183.7,-100.3)and(183.3,-101.3)..(182.9,-101.8)
				..controls(182.6,-102.1)and(182.3,-102.6)..(182.2,-102.9)
				..controls(182,-103.3)and(181.6,-104)..(181.2,-104.6)
				..controls(180.9,-105.2)and(180.5,-106.1)..(180.4,-106.6)
				..controls(179.9,-108)and(179.7,-108.5)..(179,-109.4)
				..controls(178.6,-109.8)and(178.2,-110.5)..(177.9,-111)
				..controls(177.7,-111.5)and(177.3,-112.3)..(177,-112.7)
				..controls(176.6,-113.2)and(176.4,-113.7)..(176.3,-114.2)
				..controls(176.2,-114.7)and(175.9,-115.2)..(175.4,-116)
				..controls(175,-116.6)and(174.4,-117.7)..(174,-118.5)
				..controls(173.7,-119.2)and(173.3,-120.1)..(173,-120.5)
				..controls(172.8,-120.9)and(172.4,-121.8)..(172.1,-122.5)
				..controls(171.8,-123.2)and(171.3,-124.1)..(171,-124.6)
				..controls(170.8,-125.1)and(170.3,-126)..(169.9,-126.6)
				..controls(169.6,-127.3)and(169.2,-128.1)..(169,-128.3)
				..controls(168.8,-128.6)and(168.5,-129.3)..(168.3,-129.8)
				..controls(168.1,-130.3)and(167.8,-131.1)..(167.5,-131.6)
				..controls(166.9,-132.5)and(164.2,-137.9)..(164.2,-138.1)
				..controls(164.2,-138.1)and(163.9,-138.8)..(163.4,-139.7)
				..controls(162.9,-140.5)and(162.4,-141.5)..(162.3,-141.9)
				..controls(161.9,-143)and(161.6,-143.4)..(160.9,-144.1)
				..controls(160.4,-144.6)and(160.3,-144.8)..(160.1,-145.4)
				..controls(160,-145.8)and(159.6,-146.8)..(159.2,-147.8)
				..controls(158.7,-148.8)and(158.1,-150)..(157.9,-150.6)
				..controls(157.6,-151.1)and(157.2,-151.9)..(156.9,-152.4)
				..controls(156.5,-152.8)and(156.1,-153.7)..(155.8,-154.4)
				..controls(155.6,-155)and(155.2,-155.9)..(154.9,-156.3)
				..controls(154.7,-156.7)and(154.4,-157.4)..(154.2,-157.9)
				..controls(154.1,-158.4)and(153.8,-159)..(153.5,-159.3)
				..controls(152.8,-160.1)and(153,-160.1)..(142.4,-160.1)
				..controls(137.2,-160.1)and(131.8,-160.1)..(130.4,-160)--cycle
				(141.5,-103.7)..controls(140.8,-103)and(139.8,-102.1)..(139.2,-101.6)
				..controls(138.5,-101.2)and(137.9,-100.6)..(137.8,-100.4)
				..controls(137,-99.5)and(134,-96.5)..(133.2,-96)
				..controls(132.7,-95.7)and(131.8,-94.9)..(131.1,-94.3)
				..controls(129.4,-92.7)and(127.9,-91.6)..(126.6,-91.1)
				..controls(125.2,-90.5)and(124.1,-90)..(122.9,-89.1)
				..controls(122.1,-88.6)and(121.3,-88.2)..(119.9,-87.7)
				..controls(118.8,-87.3)and(117.5,-86.8)..(116.9,-86.5)
				..controls(114.3,-85.5)and(113.7,-85.4)..(109.7,-85.5)
				..controls(105.3,-85.7)and(103.9,-85.9)..(101.6,-87)
				..controls(100.7,-87.4)and(99.4,-87.8)..(97.9,-88)
				..controls(97,-88.2)and(96.6,-88.3)..(95.8,-88.9)
				..controls(95.4,-89.2)and(94.7,-89.6)..(94.2,-89.8)
				..controls(92.9,-90.3)and(92.3,-90.7)..(90.9,-92.1)
				..controls(90.2,-92.7)and(89.4,-93.3)..(89,-93.5)..controls(87,-94.5)and(85.2,-96)..(83.6,-97.9)
				..controls(83.3,-98.2)and(82.5,-98.9)..(81.8,-99.4)
				..controls(80.7,-100.1)and(80.5,-100.3)..(80,-101.1)
				..controls(78.7,-103.1)and(77.2,-103.9)..(76.4,-103)
				..controls(75.7,-102.2)and(75.7,-103.2)..(75.7,-83.6)
				..controls(75.6,-73.6)and(75.6,-64.8)..(75.5,-64)
				..controls(75.5,-63.2)and(75.5,-61.4)..(75.6,-59.8)--(75.7,-57.1)--(76,-56.8)
				..controls(76.8,-56.2)and(78,-56.3)..(79.8,-57.2)
				..controls(80.4,-57.5)and(81.3,-57.8)..(81.8,-57.9)
				..controls(82.3,-58.1)and(83,-58.4)..(83.6,-58.8)..controls(84.3,-59.2)and(85,-59.5)..(86,-59.8)
				..controls(86.8,-60.1)and(88,-60.6)..(88.7,-60.9)
				..controls(89.3,-61.2)and(90.6,-61.7)..(91.4,-62)
				..controls(92.2,-62.3)and(93.7,-62.9)..(94.6,-63.3)
				..controls(95.6,-63.7)and(96.8,-64.1)..(97.4,-64.3)
				..controls(98,-64.5)and(99,-64.9)..(99.7,-65.2)
				..controls(100.4,-65.6)and(101.5,-66)..(102.3,-66.2)
				..controls(104.2,-66.8)and(105,-67.1)..(106.2,-67.8)
				..controls(107.9,-68.7)and(110.2,-69.4)..(111.3,-69.2)
				..controls(112.5,-69)and(115,-68.2)..(115.6,-67.8)
				..controls(117.3,-66.7)and(117.6,-66.6)..(119.5,-66.2)
				..controls(121.9,-65.8)and(122.1,-65.7)..(123,-65.1)
				..controls(124.1,-64.4)and(124.4,-64.3)..(126,-63.9)
				..controls(127,-63.7)and(127.7,-63.4)..(128.4,-63.1)
				..controls(129,-62.8)and(129.7,-62.5)..(130.6,-62.4)
				..controls(132.1,-62.1)and(134,-61.5)..(135.2,-61)
				..controls(135.7,-60.8)and(136.7,-60.5)..(137.3,-60.3)
				..controls(138.4,-60)and(139.3,-59.7)..(141.1,-59)
				..controls(141.6,-58.7)and(142,-58.6)..(142.5,-58.6)
				..controls(143.2,-58.6)and(143.3,-58.7)..(143.6,-59)
				..controls(144,-59.5)and(144.1,-59.9)..(144,-60.8)
				..controls(144,-61.2)and(144,-71.1)..(144,-82.8)--(143.9,-104.1)--(143.7,-104.5)
				..controls(143.5,-104.7)and(143.3,-104.8)..(143,-104.9)
				..controls(142.6,-104.9)and(142.5,-104.8)..(141.5,-103.7)--cycle
				(27.8,-91.4)..controls(27.1,-90.9)and(26.3,-90)..(26.1,-89.5)
				..controls(26,-89.2)and(25.8,-88.6)..(25.7,-87.9)
				..controls(25.4,-86.2)and(25.6,-84.8)..(26.5,-83.5)
				..controls(27.4,-82.3)and(28.3,-82)..(29.4,-82.3)
				..controls(31.1,-82.8)and(32.1,-84.9)..(32,-87.4)
				..controls(31.9,-89.3)and(31.2,-90.6)..(29.9,-91.4)
				..controls(29.1,-91.9)and(28.6,-91.9)..(27.8,-91.4)--cycle
				(190.1,-91.3)..controls(189.7,-91.2)and(189.3,-90.9)..(189.1,-90.7)
				..controls(188.3,-89.8)and(187.6,-87.7)..(187.6,-86.2)
				..controls(187.6,-85.5)and(187.7,-85.2)..(188.2,-84.4)
				..controls(189.1,-82.6)and(190.1,-81.8)..(191.3,-82.1)
				..controls(192,-82.3)and(193.5,-83.7)..(193.8,-84.4)
				..controls(194.1,-85.1)and(194.2,-87.1)..(194,-88.3)
				..controls(193.8,-89.8)and(193.3,-90.5)..(192.3,-91.3)
				..controls(191.7,-91.7)and(190.9,-91.7)..(190.1,-91.3)--cycle
				(173.7,-83.4)..controls(172.2,-82.5)and(172,-82.5)..(170.6,-82.4)
				..controls(169.7,-82.4)and(168.8,-82.4)..(168.3,-82.5)
				..controls(167.2,-82.7)and(166.9,-82.7)..(166.6,-82.2)
				..controls(166.3,-81.8)and(166.3,-81)..(166.8,-80)
				..controls(167.2,-78.9)and(167.4,-77.4)..(167.7,-74.1)
				..controls(167.8,-72.8)and(167.9,-71.8)..(167.9,-71.7)
				..controls(168,-71.7)and(168,-70.2)..(168,-68.4)..controls(168,-65.4)and(168,-64.9)..(167.7,-63)
				..controls(167.3,-60.5)and(167.1,-59.9)..(166.4,-58.4)
				..controls(166.1,-57.7)and(165.8,-56.9)..(165.7,-56.6)
				..controls(165.7,-56.2)and(165.4,-55.6)..(165.2,-55.2)
				..controls(164.6,-53.8)and(161.9,-50.7)..(161,-50.3)
				..controls(159.8,-49.7)and(157.9,-50)..(155.2,-51.2)
				..controls(154.7,-51.4)and(153.4,-51.8)..(152.2,-52.1)
				..controls(150.8,-52.4)and(149.4,-52.9)..(148.6,-53.2)
				..controls(147.8,-53.6)and(146.7,-54)..(146,-54.2)
				..controls(145.4,-54.4)and(144.5,-54.7)..(144,-54.9)
				..controls(143.6,-55.2)and(142.7,-55.5)..(141.8,-55.7)
				..controls(141,-55.9)and(139.6,-56.3)..(138.7,-56.6)
				..controls(137.8,-57)and(136.3,-57.4)..(135.3,-57.6)
				..controls(133.8,-57.9)and(133.2,-58.1)..(132.1,-58.7)
				..controls(131.1,-59.2)and(130.3,-59.4)..(129.4,-59.6)
				..controls(128.7,-59.8)and(127.8,-60)..(127.4,-60.1)
				..controls(127,-60.3)and(126,-60.5)..(125.1,-60.7)
				..controls(123.2,-61.1)and(121.7,-61.5)..(120.2,-62.1)
				..controls(119.7,-62.3)and(118.1,-62.8)..(116.7,-63.2)
				..controls(114.6,-63.9)and(114,-64.1)..(113,-64.7)
				..controls(111.9,-65.4)and(111.8,-65.4)..(111,-65.5)
				..controls(109.9,-65.5)and(109.3,-65.4)..(107.7,-64.7)
				..controls(107,-64.4)and(106,-64)..(105.3,-63.9)
				..controls(103.5,-63.5)and(102.2,-63)..(101.2,-62.5)
				..controls(100.1,-61.8)and(99.3,-61.5)..(97.8,-61.1)
				..controls(97.2,-61)and(96.2,-60.7)..(95.7,-60.4)
				..controls(95.1,-60.1)and(94.2,-59.8)..(93.5,-59.6)
				..controls(92.9,-59.4)and(91.7,-59)..(91,-58.7)..controls(90.2,-58.4)and(88.9,-58)..(88,-57.7)
				..controls(87.2,-57.5)and(86.1,-57.1)..(85.5,-56.8)
				..controls(85,-56.6)and(84.1,-56.2)..(83.5,-56)..controls(82.9,-55.7)and(81.8,-55.2)..(81,-54.9)
				..controls(80.1,-54.5)and(79.1,-54.2)..(78.7,-54.1)
				..controls(77.9,-53.8)and(77.3,-53.6)..(76.2,-52.9)
				..controls(75.7,-52.6)and(75.2,-52.4)..(74.1,-52.2)
				..controls(72.5,-51.9)and(72.4,-51.8)..(71.3,-51.1)
				..controls(69.9,-50.1)and(69.4,-49.9)..(67.9,-49.8)
				..controls(67,-49.7)and(66.6,-49.6)..(66.3,-49.4)
				..controls(66.1,-49.3)and(65.7,-49)..(65.4,-48.8)
				..controls(65.1,-48.7)and(64.2,-48.4)..(63.6,-48.2)
				..controls(62.9,-48)and(61.7,-47.6)..(60.9,-47.2)
				..controls(60.1,-46.8)and(58.9,-46.4)..(58.2,-46.2)
				..controls(57.6,-46)and(56.8,-45.7)..(56.5,-45.5)
				..controls(54.9,-44.6)and(54.6,-44.5)..(53.2,-44.2)
				..controls(51.7,-43.8)and(51,-43.6)..(49.9,-42.9)
				..controls(49.5,-42.6)and(48.9,-42.4)..(47.5,-42)
				..controls(45.7,-41.6)and(45.6,-41.5)..(44.3,-40.6)
				..controls(43.5,-40.1)and(42.8,-39.6)..(42.7,-39.4)
				..controls(42.4,-38.7)and(43,-38.1)..(44.5,-37.6)..controls(44.9,-37.5)and(46,-37)..(46.9,-36.6)
				..controls(47.9,-36.1)and(48.8,-35.8)..(49.6,-35.6)
				..controls(50.5,-35.4)and(51.3,-35.2)..(52.4,-34.6)
				..controls(53.3,-34.1)and(54.5,-33.7)..(54.9,-33.5)
				..controls(55.4,-33.4)and(56.4,-33)..(57.2,-32.6)
				..controls(58.6,-32)and(60.1,-31.5)..(61.6,-31.2)
				..controls(62.3,-31.1)and(62.7,-31)..(63.2,-30.7)
				..controls(64.1,-30)and(64.3,-29.9)..(65.3,-29.7)
				..controls(65.7,-29.6)and(66.5,-29.2)..(67.1,-28.9)
				..controls(67.6,-28.6)and(68.8,-28.2)..(69.6,-27.9)
				..controls(70.5,-27.7)and(71.6,-27.3)..(72.1,-27)
				..controls(72.6,-26.7)and(73.8,-26.2)..(74.8,-25.9)
				..controls(75.7,-25.5)and(76.7,-25.1)..(77.1,-25)
				..controls(77.4,-24.8)and(78.5,-24.4)..(79.6,-24.1)
				..controls(80.6,-23.8)and(82,-23.2)..(82.8,-22.9)
				..controls(83.5,-22.5)and(84.8,-22)..(85.8,-21.8)
				..controls(86.7,-21.5)and(87.8,-21.1)..(88.1,-20.9)
				..controls(88.4,-20.7)and(89.3,-20.3)..(90.1,-20.1)
				..controls(90.8,-19.8)and(91.9,-19.4)..(92.4,-19.1)
				..controls(93.5,-18.5)and(94.4,-18.2)..(95.4,-18.1)
				..controls(95.7,-18)and(96.4,-17.8)..(96.9,-17.5)
				..controls(97.5,-17.3)and(98.4,-16.9)..(99,-16.6)
				..controls(99.7,-16.3)and(100.9,-15.9)..(101.9,-15.6)
				..controls(102.9,-15.3)and(104.4,-14.7)..(105.3,-14.3)
				..controls(108.1,-13.2)and(108.5,-13.1)..(110.1,-13.1)
				..controls(111.4,-13.1)and(111.5,-13.1)..(112.4,-13.5)
				..controls(113.2,-14)and(115,-14.8)..(116.6,-15.4)
				..controls(116.9,-15.5)and(117.9,-15.8)..(118.8,-16.1)
				..controls(120.9,-16.6)and(121.5,-16.9)..(122.6,-17.9)
				..controls(123,-18.3)and(123.7,-18.8)..(123.9,-18.9)
				..controls(124.6,-19.4)and(126,-19.8)..(127.9,-20.2)
				..controls(129.6,-20.5)and(129.8,-20.6)..(131.1,-21.4)
				..controls(131.5,-21.7)and(132.4,-22.1)..(133.1,-22.3)
				..controls(136.1,-23.2)and(138.3,-24.1)..(140.2,-25)
				..controls(141,-25.5)and(141.5,-25.6)..(142.5,-25.8)
				..controls(145.1,-26.3)and(146,-26.6)..(149.4,-28.1)
				..controls(151.3,-28.9)and(153.5,-29.8)..(154.3,-30.1)
				..controls(155.1,-30.4)and(156.2,-30.8)..(156.6,-31)
				..controls(157.1,-31.2)and(158.2,-31.6)..(159.1,-31.8)
				..controls(160.6,-32.2)and(160.9,-32.3)..(161.7,-32.9)
				..controls(162.4,-33.3)and(162.9,-33.6)..(163.7,-33.8)
				..controls(164.3,-34)and(164.8,-34.1)..(164.9,-34.1)
				..controls(165,-34.1)and(165.6,-34.4)..(166.4,-34.8)
				..controls(167.3,-35.2)and(168.1,-35.4)..(169,-35.6)
				..controls(170,-35.8)and(170.5,-36)..(171.7,-36.6)
				..controls(172.4,-37)and(173.3,-37.4)..(173.6,-37.4)
				..controls(175.6,-37.7)and(176.2,-38)..(176.7,-39.1)
				..controls(176.9,-39.6)and(176.9,-39.7)..(176.7,-40.4)
				..controls(176.4,-41.6)and(175.8,-42.7)..(175.2,-43.4)
				..controls(174.5,-44.1)and(172.6,-45.5)..(171.7,-45.8)
				..controls(170.6,-46.2)and(169.1,-46)..(168.2,-45.4)
				..controls(167.9,-45.2)and(167.3,-44.8)..(166.9,-44.4)
				..controls(166.4,-44.1)and(165.8,-43.6)..(165.4,-43.2)
				..controls(165.1,-42.9)and(164.4,-42.3)..(163.9,-41.9)
				..controls(163.3,-41.5)and(162.7,-40.9)..(162.4,-40.5)
				..controls(161.4,-39.2)and(159.9,-38)..(158.8,-37.8)
				..controls(158.4,-37.8)and(158.3,-37.7)..(158.2,-37.2)
				..controls(158,-36.8)and(157.4,-36.1)..(157.1,-36.1)
				..controls(157,-36.1)and(156.7,-36.1)..(156.5,-36)
				..controls(156.1,-35.9)and(155.9,-35.7)..(155.7,-35.2)
				..controls(155.3,-34.6)and(155.1,-34.4)..(154,-34.1)
				..controls(153.6,-34)and(153.2,-33.7)..(152.9,-33.5)
				..controls(152.4,-32.9)and(151.6,-32.4)..(151.3,-32.4)
				..controls(150.9,-32.4)and(150.4,-32.6)..(150.4,-32.9)
				..controls(150.4,-33.2)and(150.8,-33.4)..(151.4,-33.6)
				..controls(152.1,-33.7)and(152.8,-34.1)..(153.5,-34.9)
				..controls(154.1,-35.5)and(154.3,-35.6)..(155,-35.8)
				..controls(155.7,-36)and(155.7,-36)..(156,-36.6)
				..controls(156.3,-37.3)and(156.7,-37.7)..(157.5,-37.8)
				..controls(158.1,-37.9)and(158.2,-38)..(158.4,-38.8)
				..controls(158.5,-39.4)and(158.6,-39.6)..(159.6,-40.4)
				..controls(160.2,-41)and(160.9,-41.6)..(161.3,-41.8)
				..controls(161.7,-42)and(161.9,-42.3)..(162.2,-42.8)
				..controls(162.4,-43.2)and(162.8,-43.8)..(163.1,-44.1)
				..controls(163.9,-44.8)and(165.1,-46.7)..(165.7,-48)
				..controls(166.2,-49.3)and(166.4,-49.7)..(167.2,-50.5)
				..controls(167.5,-50.8)and(167.9,-51.4)..(168.1,-51.8)
				..controls(168.3,-52.2)and(168.8,-53)..(169.1,-53.6)
				..controls(169.6,-54.3)and(169.8,-54.9)..(169.9,-55.6)
				..controls(170.4,-57.8)and(170.5,-58.1)..(171.2,-59.5)
				..controls(171.6,-60.3)and(172,-61.5)..(172.2,-62.1)
				..controls(172.3,-62.8)and(172.7,-63.7)..(172.9,-64.1)
				..controls(173.2,-64.6)and(173.4,-65.3)..(173.4,-65.6)
				..controls(173.8,-68.8)and(174.2,-70.6)..(174.9,-72.6)
				..controls(175.5,-74.5)and(175.9,-77.1)..(176.1,-80.2)
				..controls(176.3,-83.4)and(175.5,-84.4)..(173.7,-83.4)--cycle;
				\definecolor{ezs}{RGB}{247,179,26}\fill[ezs](-5.7,-156.6)--(-3.4,-157.5)--(61.9,12.7)--(59.7,13.6)--cycle;
				\definecolor{ezs}{RGB}{247,179,22}\fill[ezs](-3.8,-157.3)--(-1.5,-158.2)--(63.9,12)--(61.6,12.9)--cycle;
				\definecolor{ezs}{RGB}{248,179,19}\fill[ezs](-1.8,-158.1)--(0.4,-158.9)--(65.8,11.2)--(63.5,12.1)--cycle;
				\definecolor{ezs}{RGB}{248,178,16}\fill[ezs](0.1,-158.8)--(2.4,-159.7)--(67.7,10.5)--(65.4,11.4)--cycle;
				\definecolor{ezs}{RGB}{249,178,15}\fill[ezs](2,-159.5)--(4.3,-160.4)--(69.6,9.8)--(67.3,10.7)--cycle;
				\definecolor{ezs}{RGB}{251,177,17}\fill[ezs](3.9,-160.3)--(6.2,-161.1)--(71.5,9)--(69.2,9.9)--cycle;
				\definecolor{ezs}{RGB}{252,179,18}\fill[ezs](5.8,-161)--(8.1,-161.9)--(73.4,8.3)--(71.1,9.2)--cycle;
				\definecolor{ezs}{RGB}{252,182,17}\fill[ezs](7.7,-161.7)--(10,-162.6)--(75.3,7.6)--(73,8.5)--cycle;
				\definecolor{ezs}{RGB}{252,185,17}\fill[ezs](9.6,-162.5)--(11.9,-163.4)--(77.2,6.8)--(75,7.7)--cycle;
				\definecolor{ezs}{RGB}{251,189,18}\fill[ezs](11.5,-163.2)--(13.8,-164.1)--(79.2,6.1)--(76.9,7)--cycle;
				\definecolor{ezs}{RGB}{251,191,22}\fill[ezs](13.4,-163.9)--(15.7,-164.8)--(81.1,5.4)--(78.8,6.2)--cycle;
				\definecolor{ezs}{RGB}{252,193,25}\fill[ezs](15.4,-164.7)--(17.6,-165.6)--(83,4.6)--(80.7,5.5)--cycle;
				\definecolor{ezs}{RGB}{253,195,27}\fill[ezs](17.3,-165.4)--(19.6,-166.3)--(84.9,3.9)--(82.6,4.8)--cycle;
				\definecolor{ezs}{RGB}{253,196,31}\fill[ezs](19.2,-166.1)--(21.5,-167)--(86.8,3.2)--(84.5,4)--cycle;
				\definecolor{ezs}{RGB}{254,199,35}\fill[ezs](21.1,-166.9)--(23.4,-167.8)--(88.7,2.4)--(86.4,3.3)--cycle;
				\definecolor{ezs}{RGB}{253,201,38}\fill[ezs](23,-167.6)--(25.3,-168.5)--(90.6,1.7)--(88.3,2.6)--cycle;
				\definecolor{ezs}{RGB}{253,203,40}\fill[ezs](24.9,-168.3)--(27.2,-169.2)--(92.5,1)--(90.3,1.8)--cycle;
				\definecolor{ezs}{RGB}{253,205,42}\fill[ezs](26.8,-169.1)--(29.1,-170)--(94.4,0.2)--(92.2,1.1)--cycle;
				\definecolor{ezs}{RGB}{253,207,45}\fill[ezs](28.7,-169.8)--(31,-170.7)--(96.4,-0.5)--(94.1,0.4)--cycle;
				\definecolor{ezs}{RGB}{253,210,49}\fill[ezs](30.7,-170.5)--(32.9,-171.4)--(98.3,-1.2)--(96,-0.4)--cycle;
				\definecolor{ezs}{RGB}{253,212,54}\fill[ezs](32.6,-171.3)--(34.9,-172.2)--(100.2,-2)--(97.9,-1.1)--cycle;
				\definecolor{ezs}{RGB}{253,214,59}\fill[ezs](34.5,-172)--(36.8,-172.9)--(102.1,-2.7)--(99.8,-1.8)--cycle;
				\definecolor{ezs}{RGB}{253,216,63}\fill[ezs](36.4,-172.7)--(38.7,-173.6)--(104,-3.4)--(101.7,-2.6)--cycle;
				\definecolor{ezs}{RGB}{253,218,65}\fill[ezs](38.3,-173.5)--(40.6,-174.4)--(105.9,-4.2)--(103.6,-3.3)--cycle;
				\definecolor{ezs}{RGB}{254,220,66}\fill[ezs](40.2,-174.2)--(42.5,-175.1)--(107.8,-4.9)--(105.5,-4)--cycle;
				\definecolor{ezs}{RGB}{254,222,68}\fill[ezs](42.1,-175)--(44.4,-175.8)--(109.7,-5.6)--(107.5,-4.8)--cycle;
				\definecolor{ezs}{RGB}{254,224,71}\fill[ezs](44,-175.7)--(46.3,-176.6)--(111.7,-6.4)--(109.4,-5.5)--cycle;
				\definecolor{ezs}{RGB}{254,226,75}\fill[ezs](46,-176.4)--(48.2,-177.3)--(113.6,-7.1)--(111.3,-6.2)--cycle;
				\definecolor{ezs}{RGB}{254,228,79}\fill[ezs](47.9,-177.2)--(50.1,-178)--(115.5,-7.8)--(113.2,-7)--cycle;
				\definecolor{ezs}{RGB}{254,230,82}\fill[ezs](49.8,-177.9)--(52.1,-178.8)--(117.4,-8.6)--(115.1,-7.7)--cycle;
				\definecolor{ezs}{RGB}{254,232,85}\fill[ezs](51.7,-178.6)--(54,-179.5)--(119.3,-9.3)--(117,-8.4)--cycle;
				\definecolor{ezs}{RGB}{254,234,88}\fill[ezs](53.6,-179.4)--(55.9,-180.2)--(121.2,-10)--(118.9,-9.2)--cycle;
				\definecolor{ezs}{RGB}{254,237,91}\fill[ezs](55.5,-180.1)--(57.8,-181)--(123.1,-10.8)--(120.8,-9.9)--cycle;
				\definecolor{ezs}{RGB}{254,239,94}\fill[ezs](57.4,-180.8)--(59.7,-181.7)--(125,-11.5)--(122.8,-10.6)--cycle;
				\definecolor{ezs}{RGB}{254,241,97}\fill[ezs](59.3,-181.6)--(61.6,-182.4)--(127,-12.2)--(124.7,-11.4)--cycle;
				\definecolor{ezs}{RGB}{254,243,101}\fill[ezs](61.2,-182.3)--(63.5,-183.2)--(128.9,-13)--(126.6,-12.1)--cycle;
				\definecolor{ezs}{RGB}{254,245,103}\fill[ezs](63.2,-183)--(65.4,-183.9)--(130.8,-13.7)--(128.5,-12.8)--cycle;
				\definecolor{ezs}{RGB}{254,245,105}\fill[ezs](65.1,-183.8)--(67.4,-184.6)--(132.7,-14.4)--(130.4,-13.6)--cycle;
				\definecolor{ezs}{RGB}{254,245,105}\fill[ezs](67,-184.5)--(69.3,-185.4)--(134.6,-15.2)--(132.3,-14.3)--cycle;
				\definecolor{ezs}{RGB}{254,244,103}\fill[ezs](68.9,-185.2)--(71.2,-186.1)--(136.5,-15.9)--(134.2,-15)--cycle;
				\definecolor{ezs}{RGB}{253,243,101}\fill[ezs](70.8,-186)--(73.1,-186.8)--(138.4,-16.6)--(136.1,-15.8)--cycle;
				\definecolor{ezs}{RGB}{253,241,98}\fill[ezs](72.7,-186.7)--(75,-187.6)--(140.3,-17.4)--(138,-16.5)--cycle;
				\definecolor{ezs}{RGB}{253,239,95}\fill[ezs](74.6,-187.4)--(76.9,-188.3)--(142.2,-18.1)--(140,-17.2)--cycle;
				\definecolor{ezs}{RGB}{253,236,92}\fill[ezs](76.5,-188.2)--(78.8,-189)--(144.2,-18.8)--(141.9,-18)--cycle;
				\definecolor{ezs}{RGB}{253,234,88}\fill[ezs](78.5,-188.9)--(80.7,-189.8)--(146.1,-19.6)--(143.8,-18.7)--cycle;
				\definecolor{ezs}{RGB}{253,232,85}\fill[ezs](80.4,-189.6)--(82.7,-190.5)--(148,-20.3)--(145.7,-19.4)--cycle;
				\definecolor{ezs}{RGB}{253,230,82}\fill[ezs](82.3,-190.4)--(84.6,-191.2)--(149.9,-21)--(147.6,-20.2)--cycle;
				\definecolor{ezs}{RGB}{253,228,78}\fill[ezs](84.2,-191.1)--(86.5,-192)--(151.8,-21.8)--(149.5,-20.9)--cycle;
				\definecolor{ezs}{RGB}{253,226,74}\fill[ezs](86.1,-191.8)--(88.4,-192.7)--(153.7,-22.5)--(151.4,-21.6)--cycle;
				\definecolor{ezs}{RGB}{253,224,70}\fill[ezs](88,-192.6)--(90.3,-193.4)--(155.6,-23.3)--(153.3,-22.4)--cycle;
				\definecolor{ezs}{RGB}{253,222,67}\fill[ezs](89.9,-193.3)--(92.2,-194.2)--(157.5,-24)--(155.3,-23.1)--cycle;
				\definecolor{ezs}{RGB}{253,220,65}\fill[ezs](91.8,-194)--(94.1,-194.9)--(159.5,-24.7)--(157.2,-23.8)--cycle;
				\definecolor{ezs}{RGB}{253,217,63}\fill[ezs](93.7,-194.8)--(96,-195.6)--(161.4,-25.5)--(159.1,-24.6)--cycle;
				\definecolor{ezs}{RGB}{253,215,60}\fill[ezs](95.7,-195.5)--(97.9,-196.4)--(163.3,-26.2)--(161,-25.3)--cycle;
				\definecolor{ezs}{RGB}{253,213,56}\fill[ezs](97.6,-196.2)--(99.9,-197.1)--(165.2,-26.9)--(162.9,-26)--cycle;
				\definecolor{ezs}{RGB}{252,211,53}\fill[ezs](99.5,-197)--(101.8,-197.8)--(167.1,-27.7)--(164.8,-26.8)--cycle;
				\definecolor{ezs}{RGB}{252,209,49}\fill[ezs](101.4,-197.7)--(103.7,-198.6)--(169,-28.4)--(166.7,-27.5)--cycle;
				\definecolor{ezs}{RGB}{252,207,45}\fill[ezs](103.3,-198.4)--(105.6,-199.3)--(170.9,-29.1)--(168.6,-28.2)--cycle;
				\definecolor{ezs}{RGB}{253,205,42}\fill[ezs](105.2,-199.2)--(107.5,-200)--(172.8,-29.9)--(170.6,-29)--cycle;
				\definecolor{ezs}{RGB}{253,202,40}\fill[ezs](107.1,-199.9)--(109.4,-200.8)--(174.7,-30.6)--(172.5,-29.7)--cycle;
				\definecolor{ezs}{RGB}{253,200,38}\fill[ezs](109,-200.6)--(111.3,-201.5)--(176.7,-31.3)--(174.4,-30.4)--cycle;
				\definecolor{ezs}{RGB}{253,198,35}\fill[ezs](111,-201.4)--(113.2,-202.2)--(178.6,-32.1)--(176.3,-31.2)--cycle;
				\definecolor{ezs}{RGB}{253,196,31}\fill[ezs](112.9,-202.1)--(115.2,-203)--(180.5,-32.8)--(178.2,-31.9)--cycle;
				\definecolor{ezs}{RGB}{253,194,28}\fill[ezs](114.8,-202.8)--(117.1,-203.7)--(182.4,-33.5)--(180.1,-32.6)--cycle;
				\definecolor{ezs}{RGB}{253,192,26}\fill[ezs](116.7,-203.6)--(119,-204.5)--(184.3,-34.3)--(182,-33.4)--cycle;
				\definecolor{ezs}{RGB}{253,189,24}\fill[ezs](118.6,-204.3)--(120.9,-205.2)--(186.2,-35)--(183.9,-34.1)--cycle;
				\definecolor{ezs}{RGB}{253,186,20}\fill[ezs](120.5,-205)--(122.8,-205.9)--(188.1,-35.7)--(185.8,-34.8)--cycle;
				\definecolor{ezs}{RGB}{253,184,17}\fill[ezs](122.4,-205.8)--(124.7,-206.7)--(190,-36.5)--(187.8,-35.6)--cycle;
				\definecolor{ezs}{RGB}{252,183,15}\fill[ezs](124.3,-206.5)--(126.6,-207.4)--(192,-37.2)--(189.7,-36.3)--cycle;
				\definecolor{ezs}{RGB}{252,181,14}\fill[ezs](126.3,-207.2)--(128.5,-208.1)--(193.9,-37.9)--(191.6,-37.1)--cycle;
				\definecolor{ezs}{RGB}{252,179,13}\fill[ezs](128.2,-208)--(130.4,-208.9)--(195.8,-38.7)--(193.5,-37.8)--cycle;
				\definecolor{ezs}{RGB}{252,178,12}\fill[ezs](130.1,-208.7)--(132.4,-209.6)--(197.7,-39.4)--(195.4,-38.5)--cycle;
				\definecolor{ezs}{RGB}{251,178,10}\fill[ezs](132,-209.4)--(134.3,-210.3)--(199.6,-40.1)--(197.3,-39.3)--cycle;
				\definecolor{ezs}{RGB}{251,178,10}\fill[ezs](133.9,-210.2)--(136.2,-211.1)--(201.5,-40.9)--(199.2,-40)--cycle;
				\definecolor{ezs}{RGB}{250,177,11}\fill[ezs](135.8,-210.9)--(138.1,-211.8)--(203.4,-41.6)--(201.1,-40.7)--cycle;
				\definecolor{ezs}{RGB}{250,178,11}\fill[ezs](137.7,-211.6)--(140,-212.5)--(205.3,-42.3)--(203.1,-41.5)--cycle;
				\definecolor{ezs}{RGB}{248,180,10}\fill[ezs](139.6,-212.4)--(141.9,-213.3)--(207.3,-43.1)--(205,-42.2)--cycle;
				\definecolor{ezs}{RGB}{247,181,9}\fill[ezs](141.5,-213.1)--(143.8,-214)--(209.2,-43.8)--(206.9,-42.9)--cycle;
				\definecolor{ezs}{RGB}{246,182,9}\fill[ezs](143.5,-213.8)--(145.7,-214.7)--(211.1,-44.5)--(208.8,-43.7)--cycle;
				\definecolor{ezs}{RGB}{246,182,10}\fill[ezs](145.4,-214.6)--(147.7,-215.5)--(213,-45.3)--(210.7,-44.4)--cycle;
			\end{scope}
			% chữ EDUZY (Times đậm, làm dày nét bằng 5 lớp lệch nhẹ; căn đúng khung chữ logo gốc)
			\foreach \dx/\dy in {-0.7/0,0.7/0,0/-0.3,0/0.3,0/0}
			\node[anchor=center,inner sep=0pt,outer sep=0pt,text=ezText,align=center] at (109.6+\dx,-193.6+\dy)
			{\resizebox{1.01cm}{0.22cm}{\fontencoding{T1}\fontfamily{ptm}\fontseries{b}\selectfont EDUZY}};
		\end{tikzpicture}%
	}%
}
%====================== THIẾT LẬP THÂN ĐỀ ======================
%--- Môi trường tự luận {bt}: ex_test 3.3.4 chưa khai báo sẵn nên tạo ở đây
%    (ex_test đã tự móc \fixshowans và \listenumerate cho bt)
\makeatletter
\@ifundefined{bt}{\newtheorem{bt}{\color{gold1}\namebt}}{}
\makeatother
\renewcommand{\sectionmark}[1]{\markright{\thesection\ #1}}
%--- Khung tiêu đề (logo + tối đa 4 dòng; dòng 3, 4 để trống {} thì bỏ qua)
%  Đề thi:    \khungtieude{THI THỬ THPT QUỐC GIA NĂM 2026}{Môn Toán -- Đề số 21}
%                         {(Đề thi gồm: ...)}{Thời gian làm bài: 90 phút (không kể thời gian phát đề).}
%  Chuyên đề: \khungtieude{CHUYÊN ĐỀ 3 -- DÃY SỐ VÀ CÔNG THỨC TRUY HỒI}
%                         {Bồi dưỡng học sinh giỏi Toán THPT}{...}{}
\newcommand{\khungtieude}[4]{%
	\begin{tcolorbox}[colframe=gold1,colback=gold5!45,boxrule=0.5mm]
		\begin{minipage}[c]{3cm}
			\centerline{\eduzylogo[2.4cm]}
		\end{minipage}\hspace{0.5cm}
		\begin{minipage}[c]{12.5cm}
			\vskip 0.15 cm
			\centerline{\ind{#1}}
			\vskip 0.15 cm
			\centerline{\textbf{#2}}
			\if\relax\detokenize{#3}\relax\else
				\vskip 0.15 cm
				\centerline{\textit{#3}}
			\fi
			\if\relax\detokenize{#4}\relax\else
				\vskip 0.15 cm
				\centerline{\textit{#4}}
			\fi
			\vskip 0.15 cm
		\end{minipage}
	\end{tcolorbox}
}
%====================== BỔ SUNG CHO BÀI HỌC, CHUYÊN ĐỀ ======================
\usepackage[shortlabels]{enumitem}
\usepackage{tabularx}
%--- Danh sách gọn: \begin{enumerate}[gon] ... \end{enumerate}
\SetEnumitemKey{gon}{leftmargin=*,itemsep=1pt,parsep=0pt,topsep=3pt}
%--- Khung lý thuyết: \begin{khung}{} ... \end{khung}; tiêu đề chỉ dùng khi thật ngắn ("Mẹo", "Lỗi thường gặp")
\newtcolorbox{khung}[1]{breakable,enhanced,colframe=gold1,colback=gold5!45,
	boxrule=0.4mm,fonttitle=\bfseries\sffamily,coltitle=white,colbacktitle=gold1,
	attach boxed title to top left={xshift=6mm,yshift=-3mm},
	boxed title style={arc=1mm,boxrule=0.3mm},title={#1}}
%--- Tiêu đề mục nhỏ trong một dạng: \muc{Dấu hiệu}, \muc{Phương pháp}, \muc{Bài tự luyện}
\makeatletter
\newcommand{\muc}[1]{\par\medskip\noindent\indam{#1}\par\nobreak\smallskip\@afterindentfalse\@afterheading}
\makeatother
%--- Ô công thức cần nhớ (căn giữa, viền vàng, vừa khít nội dung)
%     \begin{center}\congthuc{$Q(n)=(n-1)!$}\end{center}
\newtcbox{\congthuc}{on line,colframe=gold1,colback=gold5!45,boxrule=0.5pt,arc=2pt,
	left=8pt,right=8pt,top=4pt,bottom=4pt}
%--- Công tắc lời giải bài tự luyện (ví dụ mẫu luôn dùng \loigiai)
%     \loigiaiTL[{đáp số}]{lời giải đầy đủ}
%     \TLfull -> in lời giải đầy đủ;  \TLdapso -> chỉ in một dòng "Đáp số";  \TLan -> ẩn hoàn toàn
%     (\CheDo là dethi thì mọi lời giải và đáp số đều ẩn)
\makeatletter
\newif\ifTLhien   \TLhientrue
\newif\ifTLdapso  \TLdapsofalse
\newcommand{\TLfull}{\TLhientrue\TLdapsofalse}
\newcommand{\TLdapso}{\TLhientrue\TLdapsotrue}
\newcommand{\TLan}{\TLhienfalse}
\newcommand{\loigiaiTL}[2][]{%
	\ifTLhien
		\ifTLdapso
			\def\@tlarg{#1}%
			\ifx\@tlarg\@empty\else
				\ifdefstring{\kindOPTN}{dethi}{}{\par\noindent{\bfseries\color{gold1}Đáp số:} #1\par}%
			\fi
		\else
			\loigiai{#2}%
		\fi
	\fi}
\makeatother
\TLdapso% mặc định: chỉ in đáp số bài tự luyện
%====================== BẢNG ĐÚNG/SAI (\choiceTFt) ======================
%  Hàng tiêu đề nền nâu đồng chữ trắng (Phát biểu | Đ | S), các hàng tô xen kẽ kem / trắng,
%  đường kẻ vàng kim, nhãn a) b) c) d) màu nâu đồng. Bản có đáp án: cột được chọn có vòng tròn
%  nâu đồng với dấu tích trắng (câu đúng tích ở cột Đúng, câu sai tích ở cột Sai).
%  Tiêu đề ghi đủ "Đúng | Sai": \OPTN{viettat=0} sau \input (cột Đúng/Sai rộng ra, phát biểu hẹp lại).
%  Dùng bảng gốc của ex_test: \def\BangDSgoc{1} trước \input{preamble-eduzy-v2}.
\providecommand{\BangDSgoc}{0}
\ifnum\BangDSgoc=0
\usepackage{colortbl,pifont}
\colorlet{ezDSke}{gold2!60!gold1}% đường kẻ bảng (nét 0,4pt như ex_test để không đổi chiều cao bảng)
\colorlet{ezDSle}{gold5!60}% hàng a), c)
\colorlet{ezDSchan}{white}% hàng b), d)
%--- dấu chọn ở cột Đúng/Sai
%    (bề rộng 0, căn giữa ô: không làm đổi bề rộng cột của ex_test, không báo tràn lề)
\newcommand{\EZdauchon}{\makebox[0pt]{\tikz[baseline=-1ex]\node[circle,fill=gold1,inner sep=0pt,
	minimum size=1.2em,text=white,font=\small]{\ding{51}};}}
\def\tickT{\EZdauchon}\def\tickF{\EZdauchon}
\AtBeginDocument{\sbox{\ccTFTab}{\EZdauchon}}
%--- nhãn a) b) c) d) màu nâu đồng (ở bản lời giải, nhãn câu đúng vẫn đóng khung như ex_test)
\renewcommand{\FalseTF}{\stepcounter{dapan}\textbf{\color{gold1}\DapAnTF\sepTF}\ignorespaces}
%--- trước khi dựng bảng: kẻ màu vàng kim (bề rộng cột giữ nguyên như ex_test); sau bảng trả lại màu đen
\newcount\EZhangTF
\newcommand{\EZbatdauTF}{%
	\global\EZhangTF=0\relax
	\arrayrulecolor{ezDSke}}
\newcommand{\EZketthucTF}{\arrayrulecolor{black}}
%--- hàng tiêu đề nền nâu đồng, chữ trắng
\makeatletter
\newcommand{\EZo@TF}[1]{\centering\bfseries\color{white}#1}
\newcommand{\EZtieudeTF}{%
	\ifcase\cotDS
		\ifnum\cotSTT=1\relax
			\def\rowHeader{\hline\rowcolor{gold1}\EZo@TF{TT}&\EZo@TF{\phatbieu}\tabularnewline}%
		\else
			\def\rowHeader{\hline\rowcolor{gold1}\EZo@TF{\phatbieu}\tabularnewline}%
		\fi
	\or
		\ifnum\cotSTT=1\relax
			\def\rowHeader{\hline\rowcolor{gold1}\EZo@TF{TT}&\EZo@TF{\phatbieu}&\EZo@TF{\Dung}\tabularnewline}%
		\else
			\def\rowHeader{\hline\rowcolor{gold1}\EZo@TF{\phatbieu}&\EZo@TF{\Dung}\tabularnewline}%
		\fi
	\or
		\ifnum\cotSTT=1\relax
			\def\rowHeader{\hline\rowcolor{gold1}\EZo@TF{TT}&\EZo@TF{\phatbieu}&\EZo@TF{\Dung}&\EZo@TF{\Sai}\tabularnewline}%
		\else
			\def\rowHeader{\hline\rowcolor{gold1}\EZo@TF{\phatbieu}&\EZo@TF{\Dung}&\EZo@TF{\Sai}\tabularnewline}%
		\fi
	\fi}
%--- tô xen kẽ từng hàng (gọi ngay sau \hline đầu mỗi hàng)
\newcommand{\EZmauhangTF}{%
	\noalign{\global\advance\EZhangTF\@ne}%
	\ifodd\EZhangTF\rowcolor{ezDSle}\else\rowcolor{ezDSchan}\fi}
%--- gắn vào \choiceTFt / \writeChoiceTFt của ex_test
\xpretocmd{\choiceTFt}{\EZbatdauTF}{}{\typeout{[preamble-eduzy-v2] Không gắn được bảng Đúng/Sai (đầu bảng)}}
\xpatchcmd{\choiceTFt}{\ifnum\dongPB=0\def\rowHeader{}\fi}{\EZtieudeTF\ifnum\dongPB=0\def\rowHeader{}\fi}
	{}{\typeout{[preamble-eduzy-v2] Không gắn được bảng Đúng/Sai (hàng tiêu đề)}}
\xpatchcmd{\writeChoiceTFt}{\hline\xdef\dotEX{\sdotEX}}{\hline\EZmauhangTF\xdef\dotEX{\sdotEX}}
	{}{\typeout{[preamble-eduzy-v2] Không gắn được bảng Đúng/Sai (tô hàng)}}
\xpatchcmd{\writeChoiceTFt}{{\end{longtable}}}{{\end{longtable}\EZketthucTF}}
	{}{\typeout{[preamble-eduzy-v2] Không gắn được bảng Đúng/Sai (cuối longtable)}}
\xpatchcmd{\writeChoiceTFt}{{\end{xtabular}\end{center}\medskip\noindent}}{{\end{xtabular}\end{center}\EZketthucTF\medskip\noindent}}
	{}{\typeout{[preamble-eduzy-v2] Không gắn được bảng Đúng/Sai (cuối xtabular)}}
\fi
%=====================================================================
\renewcommand{\parallel}{\mathrel{/\!\!/}}
.
\providecommand{\BangDSgoc}{0}
\ifnum\BangDSgoc=0
\usepackage{colortbl,pifont}
\colorlet{ezDSke}{gold2!60!gold1}% đường kẻ bảng (nét 0,4pt như ex_test để không đổi chiều cao bảng)
\colorlet{ezDSle}{gold5!60}% hàng a), c)
\colorlet{ezDSchan}{white}% hàng b), d)
%--- dấu chọn ở cột Đúng/Sai
%    (bề rộng 0, căn giữa ô: không làm đổi bề rộng cột của ex_test, không báo tràn lề)
\newcommand{\EZdauchon}{\makebox[0pt]{\tikz[baseline=-1ex]\node[circle,fill=gold1,inner sep=0pt,
	minimum size=1.2em,text=white,font=\small]{\ding{51}};}}
\def\tickT{\EZdauchon}\def\tickF{\EZdauchon}
\AtBeginDocument{\sbox{\ccTFTab}{\EZdauchon}}
%--- nhãn a) b) c) d) màu nâu đồng (ở bản lời giải, nhãn câu đúng vẫn đóng khung như ex_test)
\renewcommand{\FalseTF}{\stepcounter{dapan}\textbf{\color{gold1}\DapAnTF\sepTF}\ignorespaces}
%--- trước khi dựng bảng: kẻ màu vàng kim (bề rộng cột giữ nguyên như ex_test); sau bảng trả lại màu đen
\newcount\EZhangTF
\newcommand{\EZbatdauTF}{%
	\global\EZhangTF=0\relax
	\arrayrulecolor{ezDSke}}
\newcommand{\EZketthucTF}{\arrayrulecolor{black}}
%--- hàng tiêu đề nền nâu đồng, chữ trắng
\makeatletter
\newcommand{\EZo@TF}[1]{\centering\bfseries\color{white}#1}
\newcommand{\EZtieudeTF}{%
	\ifcase\cotDS
		\ifnum\cotSTT=1\relax
			\def\rowHeader{\hline\rowcolor{gold1}\EZo@TF{TT}&\EZo@TF{\phatbieu}\tabularnewline}%
		\else
			\def\rowHeader{\hline\rowcolor{gold1}\EZo@TF{\phatbieu}\tabularnewline}%
		\fi
	\or
		\ifnum\cotSTT=1\relax
			\def\rowHeader{\hline\rowcolor{gold1}\EZo@TF{TT}&\EZo@TF{\phatbieu}&\EZo@TF{\Dung}\tabularnewline}%
		\else
			\def\rowHeader{\hline\rowcolor{gold1}\EZo@TF{\phatbieu}&\EZo@TF{\Dung}\tabularnewline}%
		\fi
	\or
		\ifnum\cotSTT=1\relax
			\def\rowHeader{\hline\rowcolor{gold1}\EZo@TF{TT}&\EZo@TF{\phatbieu}&\EZo@TF{\Dung}&\EZo@TF{\Sai}\tabularnewline}%
		\else
			\def\rowHeader{\hline\rowcolor{gold1}\EZo@TF{\phatbieu}&\EZo@TF{\Dung}&\EZo@TF{\Sai}\tabularnewline}%
		\fi
	\fi}
%--- tô xen kẽ từng hàng (gọi ngay sau \hline đầu mỗi hàng)
\newcommand{\EZmauhangTF}{%
	\noalign{\global\advance\EZhangTF\@ne}%
	\ifodd\EZhangTF\rowcolor{ezDSle}\else\rowcolor{ezDSchan}\fi}
%--- gắn vào \choiceTFt / \writeChoiceTFt của ex_test
\xpretocmd{\choiceTFt}{\EZbatdauTF}{}{\typeout{[preamble-eduzy-v2] Không gắn được bảng Đúng/Sai (đầu bảng)}}
\xpatchcmd{\choiceTFt}{\ifnum\dongPB=0\def\rowHeader{}\fi}{\EZtieudeTF\ifnum\dongPB=0\def\rowHeader{}\fi}
	{}{\typeout{[preamble-eduzy-v2] Không gắn được bảng Đúng/Sai (hàng tiêu đề)}}
\xpatchcmd{\writeChoiceTFt}{\hline\xdef\dotEX{\sdotEX}}{\hline\EZmauhangTF\xdef\dotEX{\sdotEX}}
	{}{\typeout{[preamble-eduzy-v2] Không gắn được bảng Đúng/Sai (tô hàng)}}
\xpatchcmd{\writeChoiceTFt}{{\end{longtable}}}{{\end{longtable}\EZketthucTF}}
	{}{\typeout{[preamble-eduzy-v2] Không gắn được bảng Đúng/Sai (cuối longtable)}}
\xpatchcmd{\writeChoiceTFt}{{\end{xtabular}\end{center}\medskip\noindent}}{{\end{xtabular}\end{center}\EZketthucTF\medskip\noindent}}
	{}{\typeout{[preamble-eduzy-v2] Không gắn được bảng Đúng/Sai (cuối xtabular)}}
\fi
%=====================================================================
\renewcommand{\parallel}{\mathrel{/\!\!/}}

%      \TLdapso                      % tài liệu: \TLfull | \TLdapso | \TLan
%      \begin{document} ... \end{document}
%  Logo Eduzy vẽ bằng TikZ ngay trong file (\eduzylogo), không cần file ảnh.
%  Biên dịch: pdfLaTeX (máy có gói vietnam) hoặc XeLaTeX.
%============================================================
\documentclass[12pt,a4paper,twoside]{article}%<-- twoside: bat dau trang/chan trang chan-le
\usepackage{iftex}
\usepackage{amsmath,amssymb,tikz,mathrsfs,tkz-tab}
\usepackage{graphicx}
\usepackage{pgfplots}% đồ thị hàm số
\pgfplotsset{compat=1.17}
%\usepackage{mathptmx}
\usepackage{mathpazo}
\ifXeTeX% XeLaTeX: font chữ có đủ dấu tiếng Việt, giữ nguyên font toán mathpazo
	\usepackage[no-math]{fontspec}
	\setmainfont{TeX Gyre Pagella}
	\newfontfamily\EZqag{TeX Gyre Adventor}[NFSSFamily=qag,Ligatures=TeX]
\fi
\usepackage[many]{tcolorbox}
\usetikzlibrary{math,through,calc,intersections,angles,quotes,shapes,shapes.geometric,arrows,patterns,matrix,chains,arrows.meta,decorations.shapes,decorations.fractals,decorations.markings,shadows,decorations.pathreplacing}
\usepackage[top=1.6cm, bottom=1.8cm, left=1.8cm, right=1.8cm]{geometry}% v2: top 1.6cm (bản cũ 1.5cm) để chữ không sát dải đầu trang
\providecommand{\CheDo}{loigiai}% loigiai | dethi | solcolor | color | book
\expandafter\usepackage\expandafter[\CheDo]{ex_test}
% \tenchuong: khai báo ở file chính (trước \input) hoặc \renewcommand sau \input
%============================================================
%  [GỘP TỪ eduzy-de.tex] EDUZY - MÀU NÂU ĐỒNG / VÀNG KIM
%  + ĐẦU TRANG, CHÂN TRANG KIỂU DẢI NGANG (v2), ĐỐI XỨNG TRANG CHẴN / TRANG LẺ.
%  (không cần file eduzy-de.tex riêng nữa)
%  Khối này phải đặt SAU \usepackage[...]{ex_test} và SAU \def\tenchuong{...}
%============================================================
\usepackage{lastpage}
\usepackage{needspace}
%--- DÀN TRANG ĐỀ THI: tách câu (THÊM 05/11/2026; chỉ thêm lệnh mới, KHÔNG đổi hành vi mặc định) ---
%  Hai công tắc, mặc định tắt (= preamble chạy y như cũ). Bật bằng cách đặt TRƯỚC %============================================================
%  PREAMBLE EDUZY v2 - dùng chung cho đề thi, bài học, chuyên đề (gói ex_test)
%  Khác bản preamble-eduzy.tex ở đầu trang, chân trang: dải nền nâu đồng tràn ngang,
%  logo + Eduzy.com.vn, ô chữ vàng kem cạnh xiên, chân trang "Trang (số)";
%  logo Eduzy chìm giữa mọi trang (tắt: \logochimtat, bật lại: \logochimbat).
%  Cách dùng ở file chính (xem mau-chuan.tex, mau-chuyen-de.tex):
%      \def\CheDo{loigiai}          % loigiai | dethi | solcolor | color | book
%      \def\tenchuong{<dòng chữ trên đầu trang>}
%      \def\thuonghieu{<chữ ở chân trang lẻ>}   % không khai báo thì dùng mặc định
%      \def\tacgia{<chữ ở chân trang chẵn>}      % không khai báo thì dùng mặc định
%      \def\khauhieu{<chữ cạnh logo ở đầu trang>}  % tùy chọn, mặc định Eduzy.com.vn
%      %============================================================
%  PREAMBLE EDUZY v2 - dùng chung cho đề thi, bài học, chuyên đề (gói ex_test)
%  Khác bản preamble-eduzy.tex ở đầu trang, chân trang: dải nền nâu đồng tràn ngang,
%  logo + Eduzy.com.vn, ô chữ vàng kem cạnh xiên, chân trang "Trang (số)";
%  logo Eduzy chìm giữa mọi trang (tắt: \logochimtat, bật lại: \logochimbat).
%  Cách dùng ở file chính (xem mau-chuan.tex, mau-chuyen-de.tex):
%      \def\CheDo{loigiai}          % loigiai | dethi | solcolor | color | book
%      \def\tenchuong{<dòng chữ trên đầu trang>}
%      \def\thuonghieu{<chữ ở chân trang lẻ>}   % không khai báo thì dùng mặc định
%      \def\tacgia{<chữ ở chân trang chẵn>}      % không khai báo thì dùng mặc định
%      \def\khauhieu{<chữ cạnh logo ở đầu trang>}  % tùy chọn, mặc định Eduzy.com.vn
%      \input{preamble-eduzy-v2}
%      \TLdapso                      % tài liệu: \TLfull | \TLdapso | \TLan
%      \begin{document} ... \end{document}
%  Logo Eduzy vẽ bằng TikZ ngay trong file (\eduzylogo), không cần file ảnh.
%  Biên dịch: pdfLaTeX (máy có gói vietnam) hoặc XeLaTeX.
%============================================================
\documentclass[12pt,a4paper,twoside]{article}%<-- twoside: bat dau trang/chan trang chan-le
\usepackage{iftex}
\usepackage{amsmath,amssymb,tikz,mathrsfs,tkz-tab}
\usepackage{graphicx}
\usepackage{pgfplots}% đồ thị hàm số
\pgfplotsset{compat=1.17}
%\usepackage{mathptmx}
\usepackage{mathpazo}
\ifXeTeX% XeLaTeX: font chữ có đủ dấu tiếng Việt, giữ nguyên font toán mathpazo
	\usepackage[no-math]{fontspec}
	\setmainfont{TeX Gyre Pagella}
	\newfontfamily\EZqag{TeX Gyre Adventor}[NFSSFamily=qag,Ligatures=TeX]
\fi
\usepackage[many]{tcolorbox}
\usetikzlibrary{math,through,calc,intersections,angles,quotes,shapes,shapes.geometric,arrows,patterns,matrix,chains,arrows.meta,decorations.shapes,decorations.fractals,decorations.markings,shadows,decorations.pathreplacing}
\usepackage[top=1.6cm, bottom=1.8cm, left=1.8cm, right=1.8cm]{geometry}% v2: top 1.6cm (bản cũ 1.5cm) để chữ không sát dải đầu trang
\providecommand{\CheDo}{loigiai}% loigiai | dethi | solcolor | color | book
\expandafter\usepackage\expandafter[\CheDo]{ex_test}
% \tenchuong: khai báo ở file chính (trước \input) hoặc \renewcommand sau \input
%============================================================
%  [GỘP TỪ eduzy-de.tex] EDUZY - MÀU NÂU ĐỒNG / VÀNG KIM
%  + ĐẦU TRANG, CHÂN TRANG KIỂU DẢI NGANG (v2), ĐỐI XỨNG TRANG CHẴN / TRANG LẺ.
%  (không cần file eduzy-de.tex riêng nữa)
%  Khối này phải đặt SAU \usepackage[...]{ex_test} và SAU \def\tenchuong{...}
%============================================================
\usepackage{lastpage}
\usepackage{needspace}
%--- DÀN TRANG ĐỀ THI: tách câu (THÊM 05/11/2026; chỉ thêm lệnh mới, KHÔNG đổi hành vi mặc định) ---
%  Hai công tắc, mặc định tắt (= preamble chạy y như cũ). Bật bằng cách đặt TRƯỚC \input{preamble-eduzy-v2}
%  trong file đề (tools/dan_trang/dan_trang.py tự ghi khi đặt \NeedspaceCau):
%    \def\TachCauDS{1}   câu Đúng/Sai được tách: câu dẫn (kèm hình) ở cuối trang, bảng Phát biểu sang trang sau
%                        (câu dẫn và hình không tách nhau, bảng không tách giữa hai trang)
%    \def\TachCauLG{1}   bản lời giải (\CheDo khác dethi): phần đề của câu cao hơn 1/3 trang được phép tách
%  Lệnh dùng trong file đề (giá trị do tools/dan_trang/dan_trang.py đo ở bản dethi):
%    \NeedspaceCau[<cao câu dẫn + hình>]{<cao cả khối câu>}   ngay trên \begin{ex} (câu đầu phần: trên tiêu đề phần)
%        công tắc tắt: đúng bằng \Needspace{<cao cả khối câu>} như cũ;
%        \TachCauDS=1 và có [..] (câu Đúng/Sai được chọn tách): chỉ giữ câu dẫn + hình;
%        \TachCauLG=1 ở bản lời giải: đòi tối đa 1/3 \textheight.
%    \NeedspaceBang{<cao bảng Phát biểu>}%   ngay trước \choiceTFt: \TachCauDS=1 thì giữ trọn bảng trên một trang;
%        công tắc tắt: không làm gì.
\providecommand{\TachCauDS}{0}
\providecommand{\TachCauLG}{0}
\newlength{\EZnsCau}
\newcommand{\NeedspaceCau}[2][]{%
	\setlength{\EZnsCau}{#2}%
	\ifnum\TachCauDS=1 \if\relax\detokenize{#1}\relax\else\setlength{\EZnsCau}{#1}\fi\fi
	\ifnum\TachCauLG=1 \ifdefstring{\kindOPTN}{dethi}{}{\ifdim\EZnsCau>0.3333\textheight\setlength{\EZnsCau}{0.3333\textheight}\fi}\fi
	\Needspace{\EZnsCau}}
\newcommand{\NeedspaceBang}[1]{\ifnum\TachCauDS=1 \Needspace{#1}\fi}

%--- Thang màu: đậm -> nhạt
\definecolor{gold1}{RGB}{126,96,0}    % nâu đồng đậm  #7E6000
\definecolor{gold2}{RGB}{192,144,3}   % vàng kim      #C09003
\definecolor{gold3}{RGB}{255,215,99}  % vàng nhạt     #FFD763
\definecolor{gold4}{RGB}{255,231,153} % vàng rất nhạt #FFE799
\definecolor{gold5}{RGB}{255,243,205} % vàng kem      #FFF3CD
%--- Ghi đè tông xanh/đỏ cũ
\definecolor{myblue}{RGB}{126,96,0}
\colorlet{blue}{gold1}
\colorlet{cyan}{gold2}
\colorlet{teal}{gold2}

%--- Tên chương mặc định (file đề nên tự khai báo lại)
\providecommand{\tenchuong}{}
%--- Chân trang
\providecommand{\thuonghieu}{TOÁN 12 -- Eduzy.com.vn}
\providecommand{\tacgia}{Tô Gia Bảo -- 0353703783}

%--- Chữ cạnh logo ở đầu trang
\providecommand{\khauhieu}{Eduzy.com.vn}

%====================== THÔNG SỐ ĐẦU TRANG, CHÂN TRANG (v2) ======================
%  Dải nền đậm tràn ngang trang; ô chữ vàng kem có cạnh xiên, kèm 2 vạch xiên chuyển màu.
%  Đầu trang lẻ : logo + \khauhieu (trái)   | ô chữ \tenchuong (phải)
%  Chân trang lẻ: ô chữ \thuonghieu (trái)  | Trang (số) (phải)
%  Trang chẵn soi gương trang lẻ; chân trang chẵn ghi \tacgia.
%--- Màu: chỉ dùng thang vàng kim Eduzy
\colorlet{ezNen}{gold1}% dải nền đậm
\colorlet{ezVach1}{gold3}% vạch xiên sát dải nền (đậm hơn)
\colorlet{ezVach2}{gold4}% vạch xiên sát ô chữ (nhạt hơn)
\colorlet{ezO}{gold5}% ô chữ vàng kem
%--- Kích thước
\newlength{\EZdai}\setlength{\EZdai}{1.1cm}% chiều cao dải nền
\newlength{\EZo}\setlength{\EZo}{0.68cm}% chiều cao ô chữ
\newlength{\EZxien}\setlength{\EZxien}{0.55cm}% độ lệch ngang của cạnh xiên
\newlength{\EZvach}\setlength{\EZvach}{0.75cm}% bề rộng mỗi vạch xiên
\newlength{\EZkhe}\setlength{\EZkhe}{0.12cm}% khe trắng giữa vạch và ô chữ
\newlength{\EZdem}\setlength{\EZdem}{0.4cm}% khoảng đệm từ chữ tới cạnh xiên
\newlength{\EZle}\setlength{\EZle}{1.8cm}% lề trái/phải của chữ (bằng lề thân bài)
\newlength{\EZrongmin}\setlength{\EZrongmin}{4.5cm}% ô chữ rộng tối thiểu
\newlength{\EZrong}% bề rộng chữ trong ô (tính tự động)
%--- Font
\newcommand{\EZfontO}{\fontsize{10.5pt}{12pt}\selectfont\bfseries}% chữ trong ô
\newcommand{\EZfontKH}{\fontfamily{qag}\fontsize{11pt}{12pt}\selectfont\bfseries}% \khauhieu cạnh logo
\newcommand{\EZfontTrang}{\fontsize{11.5pt}{12pt}\selectfont\bfseries}% chữ "Trang" và số trang
%--- Đo bề rộng chữ trong ô
\newcommand{\EZdorong}[1]{%
	\settowidth{\EZrong}{\EZfontO #1}%
	\ifdim\EZrong<\EZrongmin\setlength{\EZrong}{\EZrongmin}\fi}
%--- Logo + chữ trắng \khauhieu (dựng sẵn một lần khi vào \begin{document})
\newsavebox{\EZlogobox}
\AtBeginDocument{\sbox{\EZlogobox}{%
		\raisebox{-0.5\height}{\eduzylogo[1cm]}\hspace{0.15cm}%
		\raisebox{-0.5\height}{\color{white}\EZfontKH\khauhieu}}}
%--- Hai vạch xiên + khe trắng. #1: điểm gốc ở mép dưới, #2: hướng (1 sang phải, -1 sang trái),
%    #3: hướng nghiêng (1 kiểu "/", -1 kiểu "\")
\newcommand{\EZvachxien}[3]{%
	\fill[white] (#1) --++({(#3)*\EZxien},\EZo) --++({(#2)*\EZkhe},0) --++({-(#3)*\EZxien},-\EZo) --cycle;
	\fill[ezVach2] ([xshift=(#2)*\EZkhe]#1) --++({(#3)*\EZxien},\EZo) --++({(#2)*\EZvach},0) --++({-(#3)*\EZxien},-\EZo) --cycle;
	\fill[ezVach1] ([xshift=(#2)*(\EZkhe+\EZvach)]#1) --++({(#3)*\EZxien},\EZo) --++({(#2)*\EZvach},0) --++({-(#3)*\EZxien},-\EZo) --cycle;
	\draw[white,line width=0.6pt] ([xshift=(#2)*(\EZkhe+\EZvach)]#1) --++({(#3)*\EZxien},\EZo);
}
%--- Chữ trong ô (căn giữa theo chiều cao ô, không lệch vì dấu tiếng Việt)
\tikzset{ezchu/.style={inner sep=0pt,text=ezNen,text height=0.72em,text depth=0.22em}}

%====================== LOGO CHÌM GIỮA TRANG ======================
%  Vẽ cùng đầu trang nên nằm dưới chữ của thân bài, có trên mọi trang.
%  Tắt: \logochimtat (từ trang đó trở đi), bật lại: \logochimbat.
%  Đổi cỡ, độ đậm: khai báo trước \input, ví dụ \def\EZchimcao{9cm} \def\EZchimdo{0.1}
\providecommand{\EZchimcao}{10cm}% chiều cao logo chìm
\providecommand{\EZchimdo}{0.12}% độ đậm (0 = trong suốt, 1 = rõ hẳn)
\newif\ifEZchim \EZchimtrue
\newcommand{\logochimtat}{\EZchimfalse}
\newcommand{\logochimbat}{\EZchimtrue}
%  Dựng sẵn logo mờ một lần khi vào \begin{document}; transparency group làm mờ cả logo
%  một lượt nên không lộ vệt ở chỗ các dải màu chồng nhau.
\newsavebox{\EZchimtho}
\newsavebox{\EZchimbox}
\AtBeginDocument{%
	\sbox{\EZchimtho}{\eduzylogo[\EZchimcao]}%
	\sbox{\EZchimbox}{\begin{tikzpicture}
			\begin{scope}[transparency group,opacity=\EZchimdo]
				\node[inner sep=0pt] {\usebox{\EZchimtho}};
			\end{scope}
	\end{tikzpicture}}}
\newcommand{\EZvechim}{%
	\ifEZchim
		\node[inner sep=0pt] at (current page.center) {\usebox{\EZchimbox}};
	\fi}

%====================== ĐẦU TRANG - TRANG LẺ ======================
%  (logo bên trái, ô chữ \tenchuong bên phải)
\newcommand{\myfancyheadOdd}{%
	\EZdorong{\tenchuong}%
	\begin{tikzpicture}[remember picture,overlay]
		\coordinate (T) at (current page.north west);
		\coordinate (P) at (current page.north east);
		%--- logo chìm giữa trang
		\EZvechim
		%--- dải nền
		\fill[ezNen] (T) rectangle ([yshift=-\EZdai]P);
		%--- ô chữ: X là góc dưới bên trái, cạnh trái xiên kiểu "/"
		\coordinate (X) at ([xshift=-\EZle-\EZrong-\EZdem-\EZxien,yshift=-\EZdai]P);
		\fill[ezO] (X) --++(\EZxien,\EZo) -- ([yshift=-\EZdai+\EZo]P) -- ([yshift=-\EZdai]P) --cycle;
		\EZvachxien{X}{-1}{1}
		\node[ezchu,anchor=east] at ([xshift=-\EZle,yshift=-\EZdai+0.5*\EZo]P) {\EZfontO\tenchuong};
		%--- logo + tên
		\node[inner sep=0pt,anchor=west] at ([xshift=\EZle,yshift=-0.5*\EZdai]T) {\usebox{\EZlogobox}};
	\end{tikzpicture}%
}

%====================== ĐẦU TRANG - TRANG CHẴN =====================
%  (soi gương: ô chữ \tenchuong bên trái, logo bên phải)
\newcommand{\myfancyheadEven}{%
	\EZdorong{\tenchuong}%
	\begin{tikzpicture}[remember picture,overlay]
		\coordinate (T) at (current page.north west);
		\coordinate (P) at (current page.north east);
		%--- logo chìm giữa trang
		\EZvechim
		%--- dải nền
		\fill[ezNen] (T) rectangle ([yshift=-\EZdai]P);
		%--- ô chữ: X là góc dưới bên phải, cạnh phải xiên kiểu "\"
		\coordinate (X) at ([xshift=\EZle+\EZrong+\EZdem+\EZxien,yshift=-\EZdai]T);
		\fill[ezO] (X) --++(-\EZxien,\EZo) -- ([yshift=-\EZdai+\EZo]T) -- ([yshift=-\EZdai]T) --cycle;
		\EZvachxien{X}{1}{-1}
		\node[ezchu,anchor=west] at ([xshift=\EZle,yshift=-\EZdai+0.5*\EZo]T) {\EZfontO\tenchuong};
		%--- logo + tên
		\node[inner sep=0pt,anchor=east] at ([xshift=-\EZle,yshift=-0.5*\EZdai]P) {\usebox{\EZlogobox}};
	\end{tikzpicture}%
}

%====================== CHÂN TRANG - TRANG LẺ ======================
%  (ô chữ \thuonghieu bên trái, "Trang (số)" bên phải)
\newcommand{\myfancyfootOdd}{%
	\EZdorong{\thuonghieu}%
	\begin{tikzpicture}[remember picture,overlay]
		\coordinate (T) at (current page.south west);
		\coordinate (P) at (current page.south east);
		%--- dải nền
		\fill[ezNen] (T) rectangle ([yshift=\EZdai]P);
		%--- ô chữ sát mép trên dải: X là góc dưới bên phải, cạnh phải xiên kiểu "/"
		\coordinate (X) at ([xshift=\EZle+\EZrong+\EZdem,yshift=\EZdai-\EZo]T);
		\fill[ezO] ([yshift=\EZdai-\EZo]T) -- (X) --++(\EZxien,\EZo) -- ([yshift=\EZdai]T) --cycle;
		\EZvachxien{X}{1}{1}
		\node[ezchu,anchor=west] at ([xshift=\EZle,yshift=\EZdai-0.5*\EZo]T) {\EZfontO\thuonghieu};
		%--- số trang
		\node[circle,draw=white,line width=0.9pt,text=white,inner sep=0pt,minimum size=0.62cm,anchor=east]
		(so) at ([xshift=-\EZle,yshift=0.5*\EZdai]P) {\EZfontTrang\thepage};
		\node[inner sep=0pt,text=white,anchor=base east] (tr) at ([xshift=-0.15cm]so.base west) {\EZfontTrang Trang};
	\end{tikzpicture}%
}

%====================== CHÂN TRANG - TRANG CHẴN ====================
%  (soi gương: "Trang (số)" bên trái, ô chữ \tacgia bên phải)
\newcommand{\myfancyfootEven}{%
	\EZdorong{\tacgia}%
	\begin{tikzpicture}[remember picture,overlay]
		\coordinate (T) at (current page.south west);
		\coordinate (P) at (current page.south east);
		%--- dải nền
		\fill[ezNen] (T) rectangle ([yshift=\EZdai]P);
		%--- ô chữ sát mép trên dải: X là góc dưới bên trái, cạnh trái xiên kiểu "\"
		\coordinate (X) at ([xshift=-\EZle-\EZrong-\EZdem,yshift=\EZdai-\EZo]P);
		\fill[ezO] ([yshift=\EZdai-\EZo]P) -- (X) --++(-\EZxien,\EZo) -- ([yshift=\EZdai]P) --cycle;
		\EZvachxien{X}{-1}{-1}
		\node[ezchu,anchor=east] at ([xshift=-\EZle,yshift=\EZdai-0.5*\EZo]P) {\EZfontO\tacgia};
		%--- số trang
		%    (vòng tròn ẩn so0 nằm giữa dải để lấy đường chân chữ, giống trang lẻ)
		\node[circle,inner sep=0pt,minimum size=0.62cm] (so0) at ([yshift=0.5*\EZdai]T) {\EZfontTrang\phantom{\thepage}};
		\node[inner sep=0pt,text=white,anchor=base west] (tr) at ($(so0.base -| T)+(\EZle,0)$) {\EZfontTrang Trang};
		\node[circle,draw=white,line width=0.9pt,text=white,inner sep=0pt,minimum size=0.62cm,anchor=base west]
		(so) at ([xshift=0.15cm]tr.base east) {\EZfontTrang\thepage};
	\end{tikzpicture}%
}

%--- Giữ tương thích với các file cũ đang gọi \myfancyhead, \myfancyfoot
\newcommand{\myfancyhead}{\myfancyheadOdd}
\newcommand{\myfancyfoot}{\myfancyfootOdd}

%======================= GẮN VÀO TRANG =============================
\pagestyle{fancy}
\fancyhf{}
\fancyhead[LO]{\myfancyheadOdd}
\fancyhead[RE]{\myfancyheadEven}
\fancyfoot[LO]{\myfancyfootOdd}
\fancyfoot[RE]{\myfancyfootEven}
\renewcommand{\headrulewidth}{0pt}
\renewcommand{\footrulewidth}{0pt}
\fancypagestyle{plain}{%
	\fancyhf{}
	\fancyhead[LO]{\myfancyheadOdd}
	\fancyhead[RE]{\myfancyheadEven}
	\fancyfoot[LO]{\myfancyfootOdd}
	\fancyfoot[RE]{\myfancyfootEven}
	\renewcommand{\headrulewidth}{0pt}
	\renewcommand{\footrulewidth}{0pt}
}
%=================== HẾT PHẦN GỘP TỪ eduzy-de.tex ===================
\newcommand{\indam}[1]{
	{\fontsize{12pt}{9pt}\sffamily \selectfont\bfseries\color{gold1} #1}
}
\newcommand{\ind}[1]{
	{\fontsize{14pt}{9pt}\fontfamily{qag} \selectfont\bfseries\color{gold1} #1}
}
\renewcommand{\FalseEX}{\stepcounter{dapan}\textbf{\selectfont\color{gold1}\Alph{dapan}.}\ignorespaces}
\renewtheorem{ex}{\color{gold1}Câu}
\renewcommand{\baselinestretch}{1.1}
\newcommand{\hoac}[1]{
	\left[\begin{aligned}#1\end{aligned}\right.}
\newcommand{\heva}[1]{
	\left\{\begin{aligned}#1\end{aligned}\right.}
\usepackage{esvect}
\def\vec{\vv}
\def\vt{\vv}
\def\overrightarrow{\vv}
%====================== LOGO EDUZY VẼ BẰNG TIKZ ======================
%  Đường viền dò từ logo gốc; dải màu vàng kim đo trực tiếp trên ảnh gốc
%  (nghiêng 21 độ, tô bằng 80 dải hẹp nên hiển thị giống nhau trên mọi trình xem PDF).
%  Dùng: \eduzylogo          -> cao 2.4cm (mặc định)
%        \eduzylogo[1.5cm]   -> cao 1.5cm
\definecolor{ezText}{RGB}{245,200,46}
\newcommand{\eduzylogo}[1][2.4cm]{%
	\resizebox{!}{#1}{%
		\begin{tikzpicture}[x=0.01cm,y=0.01cm]
			\path[use as bounding box] (0,0) rectangle (224,-224);
			% mũ tốt nghiệp + vương miện
			\begin{scope}
				\clip
				(119.4,-176.1)..controls(118.9,-176.1)and(116.8,-176)..(114.6,-176)
				..controls(101.4,-176)and(85.8,-175.9)..(83.7,-175.7)
				..controls(81.8,-175.6)and(81,-175.6)..(80.2,-175.7)
				..controls(78.8,-175.9)and(77.7,-175.9)..(76.2,-175.7)
				..controls(74.9,-175.5)and(70.8,-175.5)..(64.3,-175.7)
				..controls(59.1,-175.9)and(58.3,-175.7)..(58.6,-173.7)
				..controls(58.6,-172.8)and(59.6,-169.6)..(60.3,-168)
				..controls(60.8,-166.9)and(61.4,-165.8)..(62.3,-165)
				..controls(63.1,-164.2)and(64.3,-163.4)..(65,-163.2)
				..controls(65.5,-163.1)and(72.3,-163)..(76.6,-163.1)
				..controls(77.9,-163.1)and(80.1,-163.1)..(81.4,-163)
				..controls(82.8,-162.9)and(86.3,-162.9)..(89.2,-162.9)
				..controls(93,-162.9)and(94.9,-162.8)..(96,-162.7)
				..controls(97.2,-162.5)and(101,-162.5)..(114,-162.4)
				..controls(123,-162.3)and(134.9,-162.2)..(140.2,-162.1)
				..controls(150.1,-162)and(152.2,-162.1)..(153.6,-162.4)
				..controls(154.4,-162.6)and(157.1,-163.8)..(158.1,-164.4)
				..controls(158.8,-164.8)and(159.7,-165.7)..(159.9,-166.3)
				..controls(160,-166.5)and(160.2,-167.1)..(160.2,-167.6)
				..controls(160.3,-168)and(160.6,-169.4)..(160.9,-170.6)
				..controls(161.2,-171.7)and(161.5,-173.1)..(161.6,-173.6)
				..controls(161.8,-174.4)and(161.8,-174.5)..(161.5,-174.9)
				..controls(161.2,-175.6)and(160.8,-175.9)..(159.9,-176)
				..controls(159.2,-176.1)and(154.1,-176.2)..(151.5,-176.2)
				..controls(148.1,-176.1)and(122.6,-176.1)..(121.6,-176.2)
				..controls(120.8,-176.2)and(119.8,-176.2)..(119.4,-176.1)--cycle
				(130.4,-160)..controls(129,-160)and(119.8,-160)..(110.1,-159.9)
				..controls(100.3,-159.8)and(88.5,-159.7)..(83.9,-159.6)
				..controls(77.8,-159.5)and(75.3,-159.5)..(74.5,-159.6)
				..controls(73.9,-159.7)and(72.5,-159.8)..(71.3,-159.8)
				..controls(69.6,-159.8)and(69,-159.8)..(68.5,-159.6)
				..controls(67.6,-159.3)and(66.8,-158.6)..(66.2,-157.6)
				..controls(66,-157.3)and(65.6,-156.5)..(65.3,-156.1)
				..controls(64.6,-155.1)and(64.5,-154.9)..(63.7,-152.9)
				..controls(63.2,-151.6)and(62.9,-150.9)..(61.8,-149.6)
				..controls(61.5,-149.2)and(61.1,-148.5)..(60.9,-147.9)
				..controls(60.7,-147.4)and(60.3,-146.6)..(60.1,-146.1)
				..controls(59.9,-145.5)and(59.5,-144.9)..(59,-144.3)
				..controls(58.5,-143.7)and(58.2,-143.2)..(57.9,-142.5)
				..controls(57.7,-141.9)and(57.3,-141.1)..(57,-140.7)
				..controls(56.6,-140.3)and(56.1,-139.3)..(55.8,-138.5)
				..controls(55.5,-137.8)and(55.1,-136.8)..(54.9,-136.5)
				..controls(54.7,-136.1)and(54.3,-135.2)..(53.9,-134.5)
				..controls(53.6,-133.8)and(53,-132.7)..(52.6,-132)
				..controls(51.6,-130.5)and(51.4,-130.2)..(51.1,-129)
				..controls(50.6,-127.5)and(50.6,-127.4)..(49.7,-126.3)
				..controls(49.2,-125.7)and(48.6,-124.8)..(48.3,-124.2)
				..controls(48,-123.6)and(47.5,-122.7)..(47.3,-122.2)
				..controls(47,-121.8)and(46.5,-120.7)..(46.1,-119.8)
				..controls(45.8,-118.9)and(45.2,-117.8)..(44.9,-117.2)
				..controls(44.2,-116.1)and(43.3,-114.4)..(42.3,-112.2)
				..controls(41.8,-111.3)and(41.3,-110.4)..(41.1,-110.1)
				..controls(40.5,-109.4)and(40.3,-109)..(40,-107.8)
				..controls(39.8,-107)and(39.6,-106.5)..(39.2,-106)
				..controls(38.9,-105.6)and(38.4,-104.7)..(38,-104)
				..controls(37.7,-103.4)and(37.2,-102.5)..(37,-102.1)
				..controls(36.7,-101.7)and(36.3,-100.8)..(36,-100.1)
				..controls(35.6,-98.9)and(35.3,-98.4)..(34.6,-97.6)
				..controls(34.2,-97.1)and(34.2,-96.3)..(34.6,-96)
				..controls(35,-95.7)and(35.5,-95.9)..(36.2,-96.7)
				..controls(36.5,-97)and(37.2,-97.5)..(37.6,-97.7)..controls(38,-98)and(38.6,-98.4)..(39.1,-98.9)
				..controls(39.6,-99.4)and(40.3,-100)..(40.6,-100.3)
				..controls(41.3,-100.8)and(42.4,-101.8)..(43.7,-102.9)
				..controls(44,-103.2)and(44.6,-103.7)..(45,-103.9)
				..controls(45.4,-104.2)and(45.9,-104.6)..(46.1,-104.8)
				..controls(47.1,-105.9)and(48.8,-107.1)..(50.5,-107.9)
				..controls(51.8,-108.6)and(52.3,-109)..(53.4,-110)
				..controls(54.8,-111.2)and(55.6,-111.6)..(57,-112.1)
				..controls(57.5,-112.3)and(58.5,-112.8)..(59.4,-113.3)
				..controls(60.2,-113.7)and(61.3,-114.3)..(61.9,-114.5)
				..controls(62.4,-114.7)and(63.3,-115.1)..(63.8,-115.3)
				..controls(64.4,-115.6)and(65.4,-115.9)..(66.2,-116.1)
				..controls(67,-116.3)and(68.4,-116.6)..(69.2,-116.9)
				..controls(72.3,-117.9)and(72.8,-118)..(78,-118.1)
				..controls(84,-118.3)and(86.1,-118.1)..(88.1,-117.1)
				..controls(89.4,-116.5)and(93,-115)..(93.7,-114.9)
				..controls(94.8,-114.6)and(95.4,-114.2)..(95.9,-113.4)
				..controls(96.2,-113)and(96.7,-112.5)..(97.1,-112.2)
				..controls(98,-111.6)and(101.3,-108.4)..(102.1,-107.2)
				..controls(102.4,-106.8)and(102.9,-106.2)..(103.2,-105.8)
				..controls(103.5,-105.5)and(104,-104.8)..(104.2,-104.2)
				..controls(104.5,-103.6)and(105,-102.8)..(105.3,-102.4)
				..controls(105.6,-102)and(106.1,-101.2)..(106.5,-100.6)
				..controls(107.1,-99.5)and(108.6,-97.5)..(109.5,-96.5)
				..controls(109.9,-96.1)and(110.1,-96)..(110.4,-96)
				..controls(111.1,-96)and(111.6,-96.5)..(111.8,-97.6)
				..controls(112,-98.6)and(112.3,-99.3)..(112.8,-99.9)
				..controls(113.7,-100.9)and(115.2,-103.1)..(115.7,-104.2)
				..controls(116.1,-105)and(116.5,-105.5)..(116.9,-106)
				..controls(117.3,-106.4)and(117.8,-106.9)..(118,-107.3)
				..controls(118.2,-107.6)and(118.6,-108)..(119,-108.3)
				..controls(119.6,-108.8)and(121.1,-110.3)..(121.8,-111.2)
				..controls(122.1,-111.7)and(122.5,-112)..(123.4,-112.4)
				..controls(124.8,-113.2)and(125.8,-113.8)..(127,-114.6)
				..controls(127.5,-114.9)and(128.1,-115.3)..(128.2,-115.4)
				..controls(128.4,-115.5)and(129.3,-115.8)..(130.2,-116)
				..controls(131.1,-116.2)and(132.5,-116.6)..(133.3,-116.9)
				..controls(135.7,-117.9)and(136.1,-117.9)..(142.5,-117.8)
				..controls(145.8,-117.8)and(148.4,-117.7)..(148.8,-117.6)
				..controls(149.2,-117.6)and(150.2,-117.2)..(151,-116.8)
				..controls(152,-116.4)and(152.8,-116)..(153.6,-115.9)
				..controls(154.5,-115.7)and(155.2,-115.4)..(156.3,-114.9)
				..controls(157.1,-114.5)and(158.2,-114)..(158.7,-113.8)
				..controls(159.1,-113.7)and(159.8,-113.3)..(160.2,-113)
				..controls(160.6,-112.7)and(161.3,-112.4)..(161.7,-112.2)
				..controls(162.7,-111.9)and(163.2,-111.7)..(164.1,-111.1)
				..controls(164.5,-110.8)and(165.3,-110.4)..(165.9,-110)
				..controls(166.5,-109.7)and(167.3,-109.2)..(167.6,-108.9)
				..controls(167.9,-108.6)and(168.6,-108.2)..(169,-108)
				..controls(169.4,-107.8)and(170.2,-107.3)..(170.7,-106.9)
				..controls(171.2,-106.5)and(172.1,-105.8)..(172.7,-105.4)
				..controls(173.2,-105)and(174.1,-104.4)..(174.5,-104)
				..controls(174.9,-103.7)and(175.7,-103.1)..(176.2,-102.6)
				..controls(176.7,-102.2)and(177.4,-101.5)..(177.9,-101.2)
				..controls(178.3,-100.8)and(178.8,-100.4)..(179,-100.2)
				..controls(179.2,-100)and(179.6,-99.6)..(180,-99.3)
				..controls(180.4,-99.1)and(180.9,-98.6)..(181.2,-98.4)
				..controls(181.5,-98.1)and(182.1,-97.7)..(182.5,-97.5)
				..controls(183.2,-97.2)and(183.5,-96.9)..(184.3,-95.9)
				..controls(184.6,-95.6)and(184.6,-95.6)..(184.8,-95.7)
				..controls(185,-95.8)and(185.1,-96)..(185.1,-96.2)
				..controls(185.1,-96.5)and(185,-96.7)..(184.6,-97.2)
				..controls(184.1,-97.8)and(184,-97.9)..(183.9,-98.8)
				..controls(183.7,-100.3)and(183.3,-101.3)..(182.9,-101.8)
				..controls(182.6,-102.1)and(182.3,-102.6)..(182.2,-102.9)
				..controls(182,-103.3)and(181.6,-104)..(181.2,-104.6)
				..controls(180.9,-105.2)and(180.5,-106.1)..(180.4,-106.6)
				..controls(179.9,-108)and(179.7,-108.5)..(179,-109.4)
				..controls(178.6,-109.8)and(178.2,-110.5)..(177.9,-111)
				..controls(177.7,-111.5)and(177.3,-112.3)..(177,-112.7)
				..controls(176.6,-113.2)and(176.4,-113.7)..(176.3,-114.2)
				..controls(176.2,-114.7)and(175.9,-115.2)..(175.4,-116)
				..controls(175,-116.6)and(174.4,-117.7)..(174,-118.5)
				..controls(173.7,-119.2)and(173.3,-120.1)..(173,-120.5)
				..controls(172.8,-120.9)and(172.4,-121.8)..(172.1,-122.5)
				..controls(171.8,-123.2)and(171.3,-124.1)..(171,-124.6)
				..controls(170.8,-125.1)and(170.3,-126)..(169.9,-126.6)
				..controls(169.6,-127.3)and(169.2,-128.1)..(169,-128.3)
				..controls(168.8,-128.6)and(168.5,-129.3)..(168.3,-129.8)
				..controls(168.1,-130.3)and(167.8,-131.1)..(167.5,-131.6)
				..controls(166.9,-132.5)and(164.2,-137.9)..(164.2,-138.1)
				..controls(164.2,-138.1)and(163.9,-138.8)..(163.4,-139.7)
				..controls(162.9,-140.5)and(162.4,-141.5)..(162.3,-141.9)
				..controls(161.9,-143)and(161.6,-143.4)..(160.9,-144.1)
				..controls(160.4,-144.6)and(160.3,-144.8)..(160.1,-145.4)
				..controls(160,-145.8)and(159.6,-146.8)..(159.2,-147.8)
				..controls(158.7,-148.8)and(158.1,-150)..(157.9,-150.6)
				..controls(157.6,-151.1)and(157.2,-151.9)..(156.9,-152.4)
				..controls(156.5,-152.8)and(156.1,-153.7)..(155.8,-154.4)
				..controls(155.6,-155)and(155.2,-155.9)..(154.9,-156.3)
				..controls(154.7,-156.7)and(154.4,-157.4)..(154.2,-157.9)
				..controls(154.1,-158.4)and(153.8,-159)..(153.5,-159.3)
				..controls(152.8,-160.1)and(153,-160.1)..(142.4,-160.1)
				..controls(137.2,-160.1)and(131.8,-160.1)..(130.4,-160)--cycle
				(141.5,-103.7)..controls(140.8,-103)and(139.8,-102.1)..(139.2,-101.6)
				..controls(138.5,-101.2)and(137.9,-100.6)..(137.8,-100.4)
				..controls(137,-99.5)and(134,-96.5)..(133.2,-96)
				..controls(132.7,-95.7)and(131.8,-94.9)..(131.1,-94.3)
				..controls(129.4,-92.7)and(127.9,-91.6)..(126.6,-91.1)
				..controls(125.2,-90.5)and(124.1,-90)..(122.9,-89.1)
				..controls(122.1,-88.6)and(121.3,-88.2)..(119.9,-87.7)
				..controls(118.8,-87.3)and(117.5,-86.8)..(116.9,-86.5)
				..controls(114.3,-85.5)and(113.7,-85.4)..(109.7,-85.5)
				..controls(105.3,-85.7)and(103.9,-85.9)..(101.6,-87)
				..controls(100.7,-87.4)and(99.4,-87.8)..(97.9,-88)
				..controls(97,-88.2)and(96.6,-88.3)..(95.8,-88.9)
				..controls(95.4,-89.2)and(94.7,-89.6)..(94.2,-89.8)
				..controls(92.9,-90.3)and(92.3,-90.7)..(90.9,-92.1)
				..controls(90.2,-92.7)and(89.4,-93.3)..(89,-93.5)..controls(87,-94.5)and(85.2,-96)..(83.6,-97.9)
				..controls(83.3,-98.2)and(82.5,-98.9)..(81.8,-99.4)
				..controls(80.7,-100.1)and(80.5,-100.3)..(80,-101.1)
				..controls(78.7,-103.1)and(77.2,-103.9)..(76.4,-103)
				..controls(75.7,-102.2)and(75.7,-103.2)..(75.7,-83.6)
				..controls(75.6,-73.6)and(75.6,-64.8)..(75.5,-64)
				..controls(75.5,-63.2)and(75.5,-61.4)..(75.6,-59.8)--(75.7,-57.1)--(76,-56.8)
				..controls(76.8,-56.2)and(78,-56.3)..(79.8,-57.2)
				..controls(80.4,-57.5)and(81.3,-57.8)..(81.8,-57.9)
				..controls(82.3,-58.1)and(83,-58.4)..(83.6,-58.8)..controls(84.3,-59.2)and(85,-59.5)..(86,-59.8)
				..controls(86.8,-60.1)and(88,-60.6)..(88.7,-60.9)
				..controls(89.3,-61.2)and(90.6,-61.7)..(91.4,-62)
				..controls(92.2,-62.3)and(93.7,-62.9)..(94.6,-63.3)
				..controls(95.6,-63.7)and(96.8,-64.1)..(97.4,-64.3)
				..controls(98,-64.5)and(99,-64.9)..(99.7,-65.2)
				..controls(100.4,-65.6)and(101.5,-66)..(102.3,-66.2)
				..controls(104.2,-66.8)and(105,-67.1)..(106.2,-67.8)
				..controls(107.9,-68.7)and(110.2,-69.4)..(111.3,-69.2)
				..controls(112.5,-69)and(115,-68.2)..(115.6,-67.8)
				..controls(117.3,-66.7)and(117.6,-66.6)..(119.5,-66.2)
				..controls(121.9,-65.8)and(122.1,-65.7)..(123,-65.1)
				..controls(124.1,-64.4)and(124.4,-64.3)..(126,-63.9)
				..controls(127,-63.7)and(127.7,-63.4)..(128.4,-63.1)
				..controls(129,-62.8)and(129.7,-62.5)..(130.6,-62.4)
				..controls(132.1,-62.1)and(134,-61.5)..(135.2,-61)
				..controls(135.7,-60.8)and(136.7,-60.5)..(137.3,-60.3)
				..controls(138.4,-60)and(139.3,-59.7)..(141.1,-59)
				..controls(141.6,-58.7)and(142,-58.6)..(142.5,-58.6)
				..controls(143.2,-58.6)and(143.3,-58.7)..(143.6,-59)
				..controls(144,-59.5)and(144.1,-59.9)..(144,-60.8)
				..controls(144,-61.2)and(144,-71.1)..(144,-82.8)--(143.9,-104.1)--(143.7,-104.5)
				..controls(143.5,-104.7)and(143.3,-104.8)..(143,-104.9)
				..controls(142.6,-104.9)and(142.5,-104.8)..(141.5,-103.7)--cycle
				(27.8,-91.4)..controls(27.1,-90.9)and(26.3,-90)..(26.1,-89.5)
				..controls(26,-89.2)and(25.8,-88.6)..(25.7,-87.9)
				..controls(25.4,-86.2)and(25.6,-84.8)..(26.5,-83.5)
				..controls(27.4,-82.3)and(28.3,-82)..(29.4,-82.3)
				..controls(31.1,-82.8)and(32.1,-84.9)..(32,-87.4)
				..controls(31.9,-89.3)and(31.2,-90.6)..(29.9,-91.4)
				..controls(29.1,-91.9)and(28.6,-91.9)..(27.8,-91.4)--cycle
				(190.1,-91.3)..controls(189.7,-91.2)and(189.3,-90.9)..(189.1,-90.7)
				..controls(188.3,-89.8)and(187.6,-87.7)..(187.6,-86.2)
				..controls(187.6,-85.5)and(187.7,-85.2)..(188.2,-84.4)
				..controls(189.1,-82.6)and(190.1,-81.8)..(191.3,-82.1)
				..controls(192,-82.3)and(193.5,-83.7)..(193.8,-84.4)
				..controls(194.1,-85.1)and(194.2,-87.1)..(194,-88.3)
				..controls(193.8,-89.8)and(193.3,-90.5)..(192.3,-91.3)
				..controls(191.7,-91.7)and(190.9,-91.7)..(190.1,-91.3)--cycle
				(173.7,-83.4)..controls(172.2,-82.5)and(172,-82.5)..(170.6,-82.4)
				..controls(169.7,-82.4)and(168.8,-82.4)..(168.3,-82.5)
				..controls(167.2,-82.7)and(166.9,-82.7)..(166.6,-82.2)
				..controls(166.3,-81.8)and(166.3,-81)..(166.8,-80)
				..controls(167.2,-78.9)and(167.4,-77.4)..(167.7,-74.1)
				..controls(167.8,-72.8)and(167.9,-71.8)..(167.9,-71.7)
				..controls(168,-71.7)and(168,-70.2)..(168,-68.4)..controls(168,-65.4)and(168,-64.9)..(167.7,-63)
				..controls(167.3,-60.5)and(167.1,-59.9)..(166.4,-58.4)
				..controls(166.1,-57.7)and(165.8,-56.9)..(165.7,-56.6)
				..controls(165.7,-56.2)and(165.4,-55.6)..(165.2,-55.2)
				..controls(164.6,-53.8)and(161.9,-50.7)..(161,-50.3)
				..controls(159.8,-49.7)and(157.9,-50)..(155.2,-51.2)
				..controls(154.7,-51.4)and(153.4,-51.8)..(152.2,-52.1)
				..controls(150.8,-52.4)and(149.4,-52.9)..(148.6,-53.2)
				..controls(147.8,-53.6)and(146.7,-54)..(146,-54.2)
				..controls(145.4,-54.4)and(144.5,-54.7)..(144,-54.9)
				..controls(143.6,-55.2)and(142.7,-55.5)..(141.8,-55.7)
				..controls(141,-55.9)and(139.6,-56.3)..(138.7,-56.6)
				..controls(137.8,-57)and(136.3,-57.4)..(135.3,-57.6)
				..controls(133.8,-57.9)and(133.2,-58.1)..(132.1,-58.7)
				..controls(131.1,-59.2)and(130.3,-59.4)..(129.4,-59.6)
				..controls(128.7,-59.8)and(127.8,-60)..(127.4,-60.1)
				..controls(127,-60.3)and(126,-60.5)..(125.1,-60.7)
				..controls(123.2,-61.1)and(121.7,-61.5)..(120.2,-62.1)
				..controls(119.7,-62.3)and(118.1,-62.8)..(116.7,-63.2)
				..controls(114.6,-63.9)and(114,-64.1)..(113,-64.7)
				..controls(111.9,-65.4)and(111.8,-65.4)..(111,-65.5)
				..controls(109.9,-65.5)and(109.3,-65.4)..(107.7,-64.7)
				..controls(107,-64.4)and(106,-64)..(105.3,-63.9)
				..controls(103.5,-63.5)and(102.2,-63)..(101.2,-62.5)
				..controls(100.1,-61.8)and(99.3,-61.5)..(97.8,-61.1)
				..controls(97.2,-61)and(96.2,-60.7)..(95.7,-60.4)
				..controls(95.1,-60.1)and(94.2,-59.8)..(93.5,-59.6)
				..controls(92.9,-59.4)and(91.7,-59)..(91,-58.7)..controls(90.2,-58.4)and(88.9,-58)..(88,-57.7)
				..controls(87.2,-57.5)and(86.1,-57.1)..(85.5,-56.8)
				..controls(85,-56.6)and(84.1,-56.2)..(83.5,-56)..controls(82.9,-55.7)and(81.8,-55.2)..(81,-54.9)
				..controls(80.1,-54.5)and(79.1,-54.2)..(78.7,-54.1)
				..controls(77.9,-53.8)and(77.3,-53.6)..(76.2,-52.9)
				..controls(75.7,-52.6)and(75.2,-52.4)..(74.1,-52.2)
				..controls(72.5,-51.9)and(72.4,-51.8)..(71.3,-51.1)
				..controls(69.9,-50.1)and(69.4,-49.9)..(67.9,-49.8)
				..controls(67,-49.7)and(66.6,-49.6)..(66.3,-49.4)
				..controls(66.1,-49.3)and(65.7,-49)..(65.4,-48.8)
				..controls(65.1,-48.7)and(64.2,-48.4)..(63.6,-48.2)
				..controls(62.9,-48)and(61.7,-47.6)..(60.9,-47.2)
				..controls(60.1,-46.8)and(58.9,-46.4)..(58.2,-46.2)
				..controls(57.6,-46)and(56.8,-45.7)..(56.5,-45.5)
				..controls(54.9,-44.6)and(54.6,-44.5)..(53.2,-44.2)
				..controls(51.7,-43.8)and(51,-43.6)..(49.9,-42.9)
				..controls(49.5,-42.6)and(48.9,-42.4)..(47.5,-42)
				..controls(45.7,-41.6)and(45.6,-41.5)..(44.3,-40.6)
				..controls(43.5,-40.1)and(42.8,-39.6)..(42.7,-39.4)
				..controls(42.4,-38.7)and(43,-38.1)..(44.5,-37.6)..controls(44.9,-37.5)and(46,-37)..(46.9,-36.6)
				..controls(47.9,-36.1)and(48.8,-35.8)..(49.6,-35.6)
				..controls(50.5,-35.4)and(51.3,-35.2)..(52.4,-34.6)
				..controls(53.3,-34.1)and(54.5,-33.7)..(54.9,-33.5)
				..controls(55.4,-33.4)and(56.4,-33)..(57.2,-32.6)
				..controls(58.6,-32)and(60.1,-31.5)..(61.6,-31.2)
				..controls(62.3,-31.1)and(62.7,-31)..(63.2,-30.7)
				..controls(64.1,-30)and(64.3,-29.9)..(65.3,-29.7)
				..controls(65.7,-29.6)and(66.5,-29.2)..(67.1,-28.9)
				..controls(67.6,-28.6)and(68.8,-28.2)..(69.6,-27.9)
				..controls(70.5,-27.7)and(71.6,-27.3)..(72.1,-27)
				..controls(72.6,-26.7)and(73.8,-26.2)..(74.8,-25.9)
				..controls(75.7,-25.5)and(76.7,-25.1)..(77.1,-25)
				..controls(77.4,-24.8)and(78.5,-24.4)..(79.6,-24.1)
				..controls(80.6,-23.8)and(82,-23.2)..(82.8,-22.9)
				..controls(83.5,-22.5)and(84.8,-22)..(85.8,-21.8)
				..controls(86.7,-21.5)and(87.8,-21.1)..(88.1,-20.9)
				..controls(88.4,-20.7)and(89.3,-20.3)..(90.1,-20.1)
				..controls(90.8,-19.8)and(91.9,-19.4)..(92.4,-19.1)
				..controls(93.5,-18.5)and(94.4,-18.2)..(95.4,-18.1)
				..controls(95.7,-18)and(96.4,-17.8)..(96.9,-17.5)
				..controls(97.5,-17.3)and(98.4,-16.9)..(99,-16.6)
				..controls(99.7,-16.3)and(100.9,-15.9)..(101.9,-15.6)
				..controls(102.9,-15.3)and(104.4,-14.7)..(105.3,-14.3)
				..controls(108.1,-13.2)and(108.5,-13.1)..(110.1,-13.1)
				..controls(111.4,-13.1)and(111.5,-13.1)..(112.4,-13.5)
				..controls(113.2,-14)and(115,-14.8)..(116.6,-15.4)
				..controls(116.9,-15.5)and(117.9,-15.8)..(118.8,-16.1)
				..controls(120.9,-16.6)and(121.5,-16.9)..(122.6,-17.9)
				..controls(123,-18.3)and(123.7,-18.8)..(123.9,-18.9)
				..controls(124.6,-19.4)and(126,-19.8)..(127.9,-20.2)
				..controls(129.6,-20.5)and(129.8,-20.6)..(131.1,-21.4)
				..controls(131.5,-21.7)and(132.4,-22.1)..(133.1,-22.3)
				..controls(136.1,-23.2)and(138.3,-24.1)..(140.2,-25)
				..controls(141,-25.5)and(141.5,-25.6)..(142.5,-25.8)
				..controls(145.1,-26.3)and(146,-26.6)..(149.4,-28.1)
				..controls(151.3,-28.9)and(153.5,-29.8)..(154.3,-30.1)
				..controls(155.1,-30.4)and(156.2,-30.8)..(156.6,-31)
				..controls(157.1,-31.2)and(158.2,-31.6)..(159.1,-31.8)
				..controls(160.6,-32.2)and(160.9,-32.3)..(161.7,-32.9)
				..controls(162.4,-33.3)and(162.9,-33.6)..(163.7,-33.8)
				..controls(164.3,-34)and(164.8,-34.1)..(164.9,-34.1)
				..controls(165,-34.1)and(165.6,-34.4)..(166.4,-34.8)
				..controls(167.3,-35.2)and(168.1,-35.4)..(169,-35.6)
				..controls(170,-35.8)and(170.5,-36)..(171.7,-36.6)
				..controls(172.4,-37)and(173.3,-37.4)..(173.6,-37.4)
				..controls(175.6,-37.7)and(176.2,-38)..(176.7,-39.1)
				..controls(176.9,-39.6)and(176.9,-39.7)..(176.7,-40.4)
				..controls(176.4,-41.6)and(175.8,-42.7)..(175.2,-43.4)
				..controls(174.5,-44.1)and(172.6,-45.5)..(171.7,-45.8)
				..controls(170.6,-46.2)and(169.1,-46)..(168.2,-45.4)
				..controls(167.9,-45.2)and(167.3,-44.8)..(166.9,-44.4)
				..controls(166.4,-44.1)and(165.8,-43.6)..(165.4,-43.2)
				..controls(165.1,-42.9)and(164.4,-42.3)..(163.9,-41.9)
				..controls(163.3,-41.5)and(162.7,-40.9)..(162.4,-40.5)
				..controls(161.4,-39.2)and(159.9,-38)..(158.8,-37.8)
				..controls(158.4,-37.8)and(158.3,-37.7)..(158.2,-37.2)
				..controls(158,-36.8)and(157.4,-36.1)..(157.1,-36.1)
				..controls(157,-36.1)and(156.7,-36.1)..(156.5,-36)
				..controls(156.1,-35.9)and(155.9,-35.7)..(155.7,-35.2)
				..controls(155.3,-34.6)and(155.1,-34.4)..(154,-34.1)
				..controls(153.6,-34)and(153.2,-33.7)..(152.9,-33.5)
				..controls(152.4,-32.9)and(151.6,-32.4)..(151.3,-32.4)
				..controls(150.9,-32.4)and(150.4,-32.6)..(150.4,-32.9)
				..controls(150.4,-33.2)and(150.8,-33.4)..(151.4,-33.6)
				..controls(152.1,-33.7)and(152.8,-34.1)..(153.5,-34.9)
				..controls(154.1,-35.5)and(154.3,-35.6)..(155,-35.8)
				..controls(155.7,-36)and(155.7,-36)..(156,-36.6)
				..controls(156.3,-37.3)and(156.7,-37.7)..(157.5,-37.8)
				..controls(158.1,-37.9)and(158.2,-38)..(158.4,-38.8)
				..controls(158.5,-39.4)and(158.6,-39.6)..(159.6,-40.4)
				..controls(160.2,-41)and(160.9,-41.6)..(161.3,-41.8)
				..controls(161.7,-42)and(161.9,-42.3)..(162.2,-42.8)
				..controls(162.4,-43.2)and(162.8,-43.8)..(163.1,-44.1)
				..controls(163.9,-44.8)and(165.1,-46.7)..(165.7,-48)
				..controls(166.2,-49.3)and(166.4,-49.7)..(167.2,-50.5)
				..controls(167.5,-50.8)and(167.9,-51.4)..(168.1,-51.8)
				..controls(168.3,-52.2)and(168.8,-53)..(169.1,-53.6)
				..controls(169.6,-54.3)and(169.8,-54.9)..(169.9,-55.6)
				..controls(170.4,-57.8)and(170.5,-58.1)..(171.2,-59.5)
				..controls(171.6,-60.3)and(172,-61.5)..(172.2,-62.1)
				..controls(172.3,-62.8)and(172.7,-63.7)..(172.9,-64.1)
				..controls(173.2,-64.6)and(173.4,-65.3)..(173.4,-65.6)
				..controls(173.8,-68.8)and(174.2,-70.6)..(174.9,-72.6)
				..controls(175.5,-74.5)and(175.9,-77.1)..(176.1,-80.2)
				..controls(176.3,-83.4)and(175.5,-84.4)..(173.7,-83.4)--cycle;
				\definecolor{ezs}{RGB}{247,179,26}\fill[ezs](-5.7,-156.6)--(-3.4,-157.5)--(61.9,12.7)--(59.7,13.6)--cycle;
				\definecolor{ezs}{RGB}{247,179,22}\fill[ezs](-3.8,-157.3)--(-1.5,-158.2)--(63.9,12)--(61.6,12.9)--cycle;
				\definecolor{ezs}{RGB}{248,179,19}\fill[ezs](-1.8,-158.1)--(0.4,-158.9)--(65.8,11.2)--(63.5,12.1)--cycle;
				\definecolor{ezs}{RGB}{248,178,16}\fill[ezs](0.1,-158.8)--(2.4,-159.7)--(67.7,10.5)--(65.4,11.4)--cycle;
				\definecolor{ezs}{RGB}{249,178,15}\fill[ezs](2,-159.5)--(4.3,-160.4)--(69.6,9.8)--(67.3,10.7)--cycle;
				\definecolor{ezs}{RGB}{251,177,17}\fill[ezs](3.9,-160.3)--(6.2,-161.1)--(71.5,9)--(69.2,9.9)--cycle;
				\definecolor{ezs}{RGB}{252,179,18}\fill[ezs](5.8,-161)--(8.1,-161.9)--(73.4,8.3)--(71.1,9.2)--cycle;
				\definecolor{ezs}{RGB}{252,182,17}\fill[ezs](7.7,-161.7)--(10,-162.6)--(75.3,7.6)--(73,8.5)--cycle;
				\definecolor{ezs}{RGB}{252,185,17}\fill[ezs](9.6,-162.5)--(11.9,-163.4)--(77.2,6.8)--(75,7.7)--cycle;
				\definecolor{ezs}{RGB}{251,189,18}\fill[ezs](11.5,-163.2)--(13.8,-164.1)--(79.2,6.1)--(76.9,7)--cycle;
				\definecolor{ezs}{RGB}{251,191,22}\fill[ezs](13.4,-163.9)--(15.7,-164.8)--(81.1,5.4)--(78.8,6.2)--cycle;
				\definecolor{ezs}{RGB}{252,193,25}\fill[ezs](15.4,-164.7)--(17.6,-165.6)--(83,4.6)--(80.7,5.5)--cycle;
				\definecolor{ezs}{RGB}{253,195,27}\fill[ezs](17.3,-165.4)--(19.6,-166.3)--(84.9,3.9)--(82.6,4.8)--cycle;
				\definecolor{ezs}{RGB}{253,196,31}\fill[ezs](19.2,-166.1)--(21.5,-167)--(86.8,3.2)--(84.5,4)--cycle;
				\definecolor{ezs}{RGB}{254,199,35}\fill[ezs](21.1,-166.9)--(23.4,-167.8)--(88.7,2.4)--(86.4,3.3)--cycle;
				\definecolor{ezs}{RGB}{253,201,38}\fill[ezs](23,-167.6)--(25.3,-168.5)--(90.6,1.7)--(88.3,2.6)--cycle;
				\definecolor{ezs}{RGB}{253,203,40}\fill[ezs](24.9,-168.3)--(27.2,-169.2)--(92.5,1)--(90.3,1.8)--cycle;
				\definecolor{ezs}{RGB}{253,205,42}\fill[ezs](26.8,-169.1)--(29.1,-170)--(94.4,0.2)--(92.2,1.1)--cycle;
				\definecolor{ezs}{RGB}{253,207,45}\fill[ezs](28.7,-169.8)--(31,-170.7)--(96.4,-0.5)--(94.1,0.4)--cycle;
				\definecolor{ezs}{RGB}{253,210,49}\fill[ezs](30.7,-170.5)--(32.9,-171.4)--(98.3,-1.2)--(96,-0.4)--cycle;
				\definecolor{ezs}{RGB}{253,212,54}\fill[ezs](32.6,-171.3)--(34.9,-172.2)--(100.2,-2)--(97.9,-1.1)--cycle;
				\definecolor{ezs}{RGB}{253,214,59}\fill[ezs](34.5,-172)--(36.8,-172.9)--(102.1,-2.7)--(99.8,-1.8)--cycle;
				\definecolor{ezs}{RGB}{253,216,63}\fill[ezs](36.4,-172.7)--(38.7,-173.6)--(104,-3.4)--(101.7,-2.6)--cycle;
				\definecolor{ezs}{RGB}{253,218,65}\fill[ezs](38.3,-173.5)--(40.6,-174.4)--(105.9,-4.2)--(103.6,-3.3)--cycle;
				\definecolor{ezs}{RGB}{254,220,66}\fill[ezs](40.2,-174.2)--(42.5,-175.1)--(107.8,-4.9)--(105.5,-4)--cycle;
				\definecolor{ezs}{RGB}{254,222,68}\fill[ezs](42.1,-175)--(44.4,-175.8)--(109.7,-5.6)--(107.5,-4.8)--cycle;
				\definecolor{ezs}{RGB}{254,224,71}\fill[ezs](44,-175.7)--(46.3,-176.6)--(111.7,-6.4)--(109.4,-5.5)--cycle;
				\definecolor{ezs}{RGB}{254,226,75}\fill[ezs](46,-176.4)--(48.2,-177.3)--(113.6,-7.1)--(111.3,-6.2)--cycle;
				\definecolor{ezs}{RGB}{254,228,79}\fill[ezs](47.9,-177.2)--(50.1,-178)--(115.5,-7.8)--(113.2,-7)--cycle;
				\definecolor{ezs}{RGB}{254,230,82}\fill[ezs](49.8,-177.9)--(52.1,-178.8)--(117.4,-8.6)--(115.1,-7.7)--cycle;
				\definecolor{ezs}{RGB}{254,232,85}\fill[ezs](51.7,-178.6)--(54,-179.5)--(119.3,-9.3)--(117,-8.4)--cycle;
				\definecolor{ezs}{RGB}{254,234,88}\fill[ezs](53.6,-179.4)--(55.9,-180.2)--(121.2,-10)--(118.9,-9.2)--cycle;
				\definecolor{ezs}{RGB}{254,237,91}\fill[ezs](55.5,-180.1)--(57.8,-181)--(123.1,-10.8)--(120.8,-9.9)--cycle;
				\definecolor{ezs}{RGB}{254,239,94}\fill[ezs](57.4,-180.8)--(59.7,-181.7)--(125,-11.5)--(122.8,-10.6)--cycle;
				\definecolor{ezs}{RGB}{254,241,97}\fill[ezs](59.3,-181.6)--(61.6,-182.4)--(127,-12.2)--(124.7,-11.4)--cycle;
				\definecolor{ezs}{RGB}{254,243,101}\fill[ezs](61.2,-182.3)--(63.5,-183.2)--(128.9,-13)--(126.6,-12.1)--cycle;
				\definecolor{ezs}{RGB}{254,245,103}\fill[ezs](63.2,-183)--(65.4,-183.9)--(130.8,-13.7)--(128.5,-12.8)--cycle;
				\definecolor{ezs}{RGB}{254,245,105}\fill[ezs](65.1,-183.8)--(67.4,-184.6)--(132.7,-14.4)--(130.4,-13.6)--cycle;
				\definecolor{ezs}{RGB}{254,245,105}\fill[ezs](67,-184.5)--(69.3,-185.4)--(134.6,-15.2)--(132.3,-14.3)--cycle;
				\definecolor{ezs}{RGB}{254,244,103}\fill[ezs](68.9,-185.2)--(71.2,-186.1)--(136.5,-15.9)--(134.2,-15)--cycle;
				\definecolor{ezs}{RGB}{253,243,101}\fill[ezs](70.8,-186)--(73.1,-186.8)--(138.4,-16.6)--(136.1,-15.8)--cycle;
				\definecolor{ezs}{RGB}{253,241,98}\fill[ezs](72.7,-186.7)--(75,-187.6)--(140.3,-17.4)--(138,-16.5)--cycle;
				\definecolor{ezs}{RGB}{253,239,95}\fill[ezs](74.6,-187.4)--(76.9,-188.3)--(142.2,-18.1)--(140,-17.2)--cycle;
				\definecolor{ezs}{RGB}{253,236,92}\fill[ezs](76.5,-188.2)--(78.8,-189)--(144.2,-18.8)--(141.9,-18)--cycle;
				\definecolor{ezs}{RGB}{253,234,88}\fill[ezs](78.5,-188.9)--(80.7,-189.8)--(146.1,-19.6)--(143.8,-18.7)--cycle;
				\definecolor{ezs}{RGB}{253,232,85}\fill[ezs](80.4,-189.6)--(82.7,-190.5)--(148,-20.3)--(145.7,-19.4)--cycle;
				\definecolor{ezs}{RGB}{253,230,82}\fill[ezs](82.3,-190.4)--(84.6,-191.2)--(149.9,-21)--(147.6,-20.2)--cycle;
				\definecolor{ezs}{RGB}{253,228,78}\fill[ezs](84.2,-191.1)--(86.5,-192)--(151.8,-21.8)--(149.5,-20.9)--cycle;
				\definecolor{ezs}{RGB}{253,226,74}\fill[ezs](86.1,-191.8)--(88.4,-192.7)--(153.7,-22.5)--(151.4,-21.6)--cycle;
				\definecolor{ezs}{RGB}{253,224,70}\fill[ezs](88,-192.6)--(90.3,-193.4)--(155.6,-23.3)--(153.3,-22.4)--cycle;
				\definecolor{ezs}{RGB}{253,222,67}\fill[ezs](89.9,-193.3)--(92.2,-194.2)--(157.5,-24)--(155.3,-23.1)--cycle;
				\definecolor{ezs}{RGB}{253,220,65}\fill[ezs](91.8,-194)--(94.1,-194.9)--(159.5,-24.7)--(157.2,-23.8)--cycle;
				\definecolor{ezs}{RGB}{253,217,63}\fill[ezs](93.7,-194.8)--(96,-195.6)--(161.4,-25.5)--(159.1,-24.6)--cycle;
				\definecolor{ezs}{RGB}{253,215,60}\fill[ezs](95.7,-195.5)--(97.9,-196.4)--(163.3,-26.2)--(161,-25.3)--cycle;
				\definecolor{ezs}{RGB}{253,213,56}\fill[ezs](97.6,-196.2)--(99.9,-197.1)--(165.2,-26.9)--(162.9,-26)--cycle;
				\definecolor{ezs}{RGB}{252,211,53}\fill[ezs](99.5,-197)--(101.8,-197.8)--(167.1,-27.7)--(164.8,-26.8)--cycle;
				\definecolor{ezs}{RGB}{252,209,49}\fill[ezs](101.4,-197.7)--(103.7,-198.6)--(169,-28.4)--(166.7,-27.5)--cycle;
				\definecolor{ezs}{RGB}{252,207,45}\fill[ezs](103.3,-198.4)--(105.6,-199.3)--(170.9,-29.1)--(168.6,-28.2)--cycle;
				\definecolor{ezs}{RGB}{253,205,42}\fill[ezs](105.2,-199.2)--(107.5,-200)--(172.8,-29.9)--(170.6,-29)--cycle;
				\definecolor{ezs}{RGB}{253,202,40}\fill[ezs](107.1,-199.9)--(109.4,-200.8)--(174.7,-30.6)--(172.5,-29.7)--cycle;
				\definecolor{ezs}{RGB}{253,200,38}\fill[ezs](109,-200.6)--(111.3,-201.5)--(176.7,-31.3)--(174.4,-30.4)--cycle;
				\definecolor{ezs}{RGB}{253,198,35}\fill[ezs](111,-201.4)--(113.2,-202.2)--(178.6,-32.1)--(176.3,-31.2)--cycle;
				\definecolor{ezs}{RGB}{253,196,31}\fill[ezs](112.9,-202.1)--(115.2,-203)--(180.5,-32.8)--(178.2,-31.9)--cycle;
				\definecolor{ezs}{RGB}{253,194,28}\fill[ezs](114.8,-202.8)--(117.1,-203.7)--(182.4,-33.5)--(180.1,-32.6)--cycle;
				\definecolor{ezs}{RGB}{253,192,26}\fill[ezs](116.7,-203.6)--(119,-204.5)--(184.3,-34.3)--(182,-33.4)--cycle;
				\definecolor{ezs}{RGB}{253,189,24}\fill[ezs](118.6,-204.3)--(120.9,-205.2)--(186.2,-35)--(183.9,-34.1)--cycle;
				\definecolor{ezs}{RGB}{253,186,20}\fill[ezs](120.5,-205)--(122.8,-205.9)--(188.1,-35.7)--(185.8,-34.8)--cycle;
				\definecolor{ezs}{RGB}{253,184,17}\fill[ezs](122.4,-205.8)--(124.7,-206.7)--(190,-36.5)--(187.8,-35.6)--cycle;
				\definecolor{ezs}{RGB}{252,183,15}\fill[ezs](124.3,-206.5)--(126.6,-207.4)--(192,-37.2)--(189.7,-36.3)--cycle;
				\definecolor{ezs}{RGB}{252,181,14}\fill[ezs](126.3,-207.2)--(128.5,-208.1)--(193.9,-37.9)--(191.6,-37.1)--cycle;
				\definecolor{ezs}{RGB}{252,179,13}\fill[ezs](128.2,-208)--(130.4,-208.9)--(195.8,-38.7)--(193.5,-37.8)--cycle;
				\definecolor{ezs}{RGB}{252,178,12}\fill[ezs](130.1,-208.7)--(132.4,-209.6)--(197.7,-39.4)--(195.4,-38.5)--cycle;
				\definecolor{ezs}{RGB}{251,178,10}\fill[ezs](132,-209.4)--(134.3,-210.3)--(199.6,-40.1)--(197.3,-39.3)--cycle;
				\definecolor{ezs}{RGB}{251,178,10}\fill[ezs](133.9,-210.2)--(136.2,-211.1)--(201.5,-40.9)--(199.2,-40)--cycle;
				\definecolor{ezs}{RGB}{250,177,11}\fill[ezs](135.8,-210.9)--(138.1,-211.8)--(203.4,-41.6)--(201.1,-40.7)--cycle;
				\definecolor{ezs}{RGB}{250,178,11}\fill[ezs](137.7,-211.6)--(140,-212.5)--(205.3,-42.3)--(203.1,-41.5)--cycle;
				\definecolor{ezs}{RGB}{248,180,10}\fill[ezs](139.6,-212.4)--(141.9,-213.3)--(207.3,-43.1)--(205,-42.2)--cycle;
				\definecolor{ezs}{RGB}{247,181,9}\fill[ezs](141.5,-213.1)--(143.8,-214)--(209.2,-43.8)--(206.9,-42.9)--cycle;
				\definecolor{ezs}{RGB}{246,182,9}\fill[ezs](143.5,-213.8)--(145.7,-214.7)--(211.1,-44.5)--(208.8,-43.7)--cycle;
				\definecolor{ezs}{RGB}{246,182,10}\fill[ezs](145.4,-214.6)--(147.7,-215.5)--(213,-45.3)--(210.7,-44.4)--cycle;
			\end{scope}
			% chữ EDUZY (Times đậm, làm dày nét bằng 5 lớp lệch nhẹ; căn đúng khung chữ logo gốc)
			\foreach \dx/\dy in {-0.7/0,0.7/0,0/-0.3,0/0.3,0/0}
			\node[anchor=center,inner sep=0pt,outer sep=0pt,text=ezText,align=center] at (109.6+\dx,-193.6+\dy)
			{\resizebox{1.01cm}{0.22cm}{\fontencoding{T1}\fontfamily{ptm}\fontseries{b}\selectfont EDUZY}};
		\end{tikzpicture}%
	}%
}
%====================== THIẾT LẬP THÂN ĐỀ ======================
%--- Môi trường tự luận {bt}: ex_test 3.3.4 chưa khai báo sẵn nên tạo ở đây
%    (ex_test đã tự móc \fixshowans và \listenumerate cho bt)
\makeatletter
\@ifundefined{bt}{\newtheorem{bt}{\color{gold1}\namebt}}{}
\makeatother
\renewcommand{\sectionmark}[1]{\markright{\thesection\ #1}}
%--- Khung tiêu đề (logo + tối đa 4 dòng; dòng 3, 4 để trống {} thì bỏ qua)
%  Đề thi:    \khungtieude{THI THỬ THPT QUỐC GIA NĂM 2026}{Môn Toán -- Đề số 21}
%                         {(Đề thi gồm: ...)}{Thời gian làm bài: 90 phút (không kể thời gian phát đề).}
%  Chuyên đề: \khungtieude{CHUYÊN ĐỀ 3 -- DÃY SỐ VÀ CÔNG THỨC TRUY HỒI}
%                         {Bồi dưỡng học sinh giỏi Toán THPT}{...}{}
\newcommand{\khungtieude}[4]{%
	\begin{tcolorbox}[colframe=gold1,colback=gold5!45,boxrule=0.5mm]
		\begin{minipage}[c]{3cm}
			\centerline{\eduzylogo[2.4cm]}
		\end{minipage}\hspace{0.5cm}
		\begin{minipage}[c]{12.5cm}
			\vskip 0.15 cm
			\centerline{\ind{#1}}
			\vskip 0.15 cm
			\centerline{\textbf{#2}}
			\if\relax\detokenize{#3}\relax\else
				\vskip 0.15 cm
				\centerline{\textit{#3}}
			\fi
			\if\relax\detokenize{#4}\relax\else
				\vskip 0.15 cm
				\centerline{\textit{#4}}
			\fi
			\vskip 0.15 cm
		\end{minipage}
	\end{tcolorbox}
}
%====================== BỔ SUNG CHO BÀI HỌC, CHUYÊN ĐỀ ======================
\usepackage[shortlabels]{enumitem}
\usepackage{tabularx}
%--- Danh sách gọn: \begin{enumerate}[gon] ... \end{enumerate}
\SetEnumitemKey{gon}{leftmargin=*,itemsep=1pt,parsep=0pt,topsep=3pt}
%--- Khung lý thuyết: \begin{khung}{} ... \end{khung}; tiêu đề chỉ dùng khi thật ngắn ("Mẹo", "Lỗi thường gặp")
\newtcolorbox{khung}[1]{breakable,enhanced,colframe=gold1,colback=gold5!45,
	boxrule=0.4mm,fonttitle=\bfseries\sffamily,coltitle=white,colbacktitle=gold1,
	attach boxed title to top left={xshift=6mm,yshift=-3mm},
	boxed title style={arc=1mm,boxrule=0.3mm},title={#1}}
%--- Tiêu đề mục nhỏ trong một dạng: \muc{Dấu hiệu}, \muc{Phương pháp}, \muc{Bài tự luyện}
\makeatletter
\newcommand{\muc}[1]{\par\medskip\noindent\indam{#1}\par\nobreak\smallskip\@afterindentfalse\@afterheading}
\makeatother
%--- Ô công thức cần nhớ (căn giữa, viền vàng, vừa khít nội dung)
%     \begin{center}\congthuc{$Q(n)=(n-1)!$}\end{center}
\newtcbox{\congthuc}{on line,colframe=gold1,colback=gold5!45,boxrule=0.5pt,arc=2pt,
	left=8pt,right=8pt,top=4pt,bottom=4pt}
%--- Công tắc lời giải bài tự luyện (ví dụ mẫu luôn dùng \loigiai)
%     \loigiaiTL[{đáp số}]{lời giải đầy đủ}
%     \TLfull -> in lời giải đầy đủ;  \TLdapso -> chỉ in một dòng "Đáp số";  \TLan -> ẩn hoàn toàn
%     (\CheDo là dethi thì mọi lời giải và đáp số đều ẩn)
\makeatletter
\newif\ifTLhien   \TLhientrue
\newif\ifTLdapso  \TLdapsofalse
\newcommand{\TLfull}{\TLhientrue\TLdapsofalse}
\newcommand{\TLdapso}{\TLhientrue\TLdapsotrue}
\newcommand{\TLan}{\TLhienfalse}
\newcommand{\loigiaiTL}[2][]{%
	\ifTLhien
		\ifTLdapso
			\def\@tlarg{#1}%
			\ifx\@tlarg\@empty\else
				\ifdefstring{\kindOPTN}{dethi}{}{\par\noindent{\bfseries\color{gold1}Đáp số:} #1\par}%
			\fi
		\else
			\loigiai{#2}%
		\fi
	\fi}
\makeatother
\TLdapso% mặc định: chỉ in đáp số bài tự luyện
%====================== BẢNG ĐÚNG/SAI (\choiceTFt) ======================
%  Hàng tiêu đề nền nâu đồng chữ trắng (Phát biểu | Đ | S), các hàng tô xen kẽ kem / trắng,
%  đường kẻ vàng kim, nhãn a) b) c) d) màu nâu đồng. Bản có đáp án: cột được chọn có vòng tròn
%  nâu đồng với dấu tích trắng (câu đúng tích ở cột Đúng, câu sai tích ở cột Sai).
%  Tiêu đề ghi đủ "Đúng | Sai": \OPTN{viettat=0} sau \input (cột Đúng/Sai rộng ra, phát biểu hẹp lại).
%  Dùng bảng gốc của ex_test: \def\BangDSgoc{1} trước \input{preamble-eduzy-v2}.
\providecommand{\BangDSgoc}{0}
\ifnum\BangDSgoc=0
\usepackage{colortbl,pifont}
\colorlet{ezDSke}{gold2!60!gold1}% đường kẻ bảng (nét 0,4pt như ex_test để không đổi chiều cao bảng)
\colorlet{ezDSle}{gold5!60}% hàng a), c)
\colorlet{ezDSchan}{white}% hàng b), d)
%--- dấu chọn ở cột Đúng/Sai
%    (bề rộng 0, căn giữa ô: không làm đổi bề rộng cột của ex_test, không báo tràn lề)
\newcommand{\EZdauchon}{\makebox[0pt]{\tikz[baseline=-1ex]\node[circle,fill=gold1,inner sep=0pt,
	minimum size=1.2em,text=white,font=\small]{\ding{51}};}}
\def\tickT{\EZdauchon}\def\tickF{\EZdauchon}
\AtBeginDocument{\sbox{\ccTFTab}{\EZdauchon}}
%--- nhãn a) b) c) d) màu nâu đồng (ở bản lời giải, nhãn câu đúng vẫn đóng khung như ex_test)
\renewcommand{\FalseTF}{\stepcounter{dapan}\textbf{\color{gold1}\DapAnTF\sepTF}\ignorespaces}
%--- trước khi dựng bảng: kẻ màu vàng kim (bề rộng cột giữ nguyên như ex_test); sau bảng trả lại màu đen
\newcount\EZhangTF
\newcommand{\EZbatdauTF}{%
	\global\EZhangTF=0\relax
	\arrayrulecolor{ezDSke}}
\newcommand{\EZketthucTF}{\arrayrulecolor{black}}
%--- hàng tiêu đề nền nâu đồng, chữ trắng
\makeatletter
\newcommand{\EZo@TF}[1]{\centering\bfseries\color{white}#1}
\newcommand{\EZtieudeTF}{%
	\ifcase\cotDS
		\ifnum\cotSTT=1\relax
			\def\rowHeader{\hline\rowcolor{gold1}\EZo@TF{TT}&\EZo@TF{\phatbieu}\tabularnewline}%
		\else
			\def\rowHeader{\hline\rowcolor{gold1}\EZo@TF{\phatbieu}\tabularnewline}%
		\fi
	\or
		\ifnum\cotSTT=1\relax
			\def\rowHeader{\hline\rowcolor{gold1}\EZo@TF{TT}&\EZo@TF{\phatbieu}&\EZo@TF{\Dung}\tabularnewline}%
		\else
			\def\rowHeader{\hline\rowcolor{gold1}\EZo@TF{\phatbieu}&\EZo@TF{\Dung}\tabularnewline}%
		\fi
	\or
		\ifnum\cotSTT=1\relax
			\def\rowHeader{\hline\rowcolor{gold1}\EZo@TF{TT}&\EZo@TF{\phatbieu}&\EZo@TF{\Dung}&\EZo@TF{\Sai}\tabularnewline}%
		\else
			\def\rowHeader{\hline\rowcolor{gold1}\EZo@TF{\phatbieu}&\EZo@TF{\Dung}&\EZo@TF{\Sai}\tabularnewline}%
		\fi
	\fi}
%--- tô xen kẽ từng hàng (gọi ngay sau \hline đầu mỗi hàng)
\newcommand{\EZmauhangTF}{%
	\noalign{\global\advance\EZhangTF\@ne}%
	\ifodd\EZhangTF\rowcolor{ezDSle}\else\rowcolor{ezDSchan}\fi}
%--- gắn vào \choiceTFt / \writeChoiceTFt của ex_test
\xpretocmd{\choiceTFt}{\EZbatdauTF}{}{\typeout{[preamble-eduzy-v2] Không gắn được bảng Đúng/Sai (đầu bảng)}}
\xpatchcmd{\choiceTFt}{\ifnum\dongPB=0\def\rowHeader{}\fi}{\EZtieudeTF\ifnum\dongPB=0\def\rowHeader{}\fi}
	{}{\typeout{[preamble-eduzy-v2] Không gắn được bảng Đúng/Sai (hàng tiêu đề)}}
\xpatchcmd{\writeChoiceTFt}{\hline\xdef\dotEX{\sdotEX}}{\hline\EZmauhangTF\xdef\dotEX{\sdotEX}}
	{}{\typeout{[preamble-eduzy-v2] Không gắn được bảng Đúng/Sai (tô hàng)}}
\xpatchcmd{\writeChoiceTFt}{{\end{longtable}}}{{\end{longtable}\EZketthucTF}}
	{}{\typeout{[preamble-eduzy-v2] Không gắn được bảng Đúng/Sai (cuối longtable)}}
\xpatchcmd{\writeChoiceTFt}{{\end{xtabular}\end{center}\medskip\noindent}}{{\end{xtabular}\end{center}\EZketthucTF\medskip\noindent}}
	{}{\typeout{[preamble-eduzy-v2] Không gắn được bảng Đúng/Sai (cuối xtabular)}}
\fi
%=====================================================================
\renewcommand{\parallel}{\mathrel{/\!\!/}}

%      \TLdapso                      % tài liệu: \TLfull | \TLdapso | \TLan
%      \begin{document} ... \end{document}
%  Logo Eduzy vẽ bằng TikZ ngay trong file (\eduzylogo), không cần file ảnh.
%  Biên dịch: pdfLaTeX (máy có gói vietnam) hoặc XeLaTeX.
%============================================================
\documentclass[12pt,a4paper,twoside]{article}%<-- twoside: bat dau trang/chan trang chan-le
\usepackage{iftex}
\usepackage{amsmath,amssymb,tikz,mathrsfs,tkz-tab}
\usepackage{graphicx}
\usepackage{pgfplots}% đồ thị hàm số
\pgfplotsset{compat=1.17}
%\usepackage{mathptmx}
\usepackage{mathpazo}
\ifXeTeX% XeLaTeX: font chữ có đủ dấu tiếng Việt, giữ nguyên font toán mathpazo
	\usepackage[no-math]{fontspec}
	\setmainfont{TeX Gyre Pagella}
	\newfontfamily\EZqag{TeX Gyre Adventor}[NFSSFamily=qag,Ligatures=TeX]
\fi
\usepackage[many]{tcolorbox}
\usetikzlibrary{math,through,calc,intersections,angles,quotes,shapes,shapes.geometric,arrows,patterns,matrix,chains,arrows.meta,decorations.shapes,decorations.fractals,decorations.markings,shadows,decorations.pathreplacing}
\usepackage[top=1.6cm, bottom=1.8cm, left=1.8cm, right=1.8cm]{geometry}% v2: top 1.6cm (bản cũ 1.5cm) để chữ không sát dải đầu trang
\providecommand{\CheDo}{loigiai}% loigiai | dethi | solcolor | color | book
\expandafter\usepackage\expandafter[\CheDo]{ex_test}
% \tenchuong: khai báo ở file chính (trước \input) hoặc \renewcommand sau \input
%============================================================
%  [GỘP TỪ eduzy-de.tex] EDUZY - MÀU NÂU ĐỒNG / VÀNG KIM
%  + ĐẦU TRANG, CHÂN TRANG KIỂU DẢI NGANG (v2), ĐỐI XỨNG TRANG CHẴN / TRANG LẺ.
%  (không cần file eduzy-de.tex riêng nữa)
%  Khối này phải đặt SAU \usepackage[...]{ex_test} và SAU \def\tenchuong{...}
%============================================================
\usepackage{lastpage}
\usepackage{needspace}
%--- DÀN TRANG ĐỀ THI: tách câu (THÊM 05/11/2026; chỉ thêm lệnh mới, KHÔNG đổi hành vi mặc định) ---
%  Hai công tắc, mặc định tắt (= preamble chạy y như cũ). Bật bằng cách đặt TRƯỚC %============================================================
%  PREAMBLE EDUZY v2 - dùng chung cho đề thi, bài học, chuyên đề (gói ex_test)
%  Khác bản preamble-eduzy.tex ở đầu trang, chân trang: dải nền nâu đồng tràn ngang,
%  logo + Eduzy.com.vn, ô chữ vàng kem cạnh xiên, chân trang "Trang (số)";
%  logo Eduzy chìm giữa mọi trang (tắt: \logochimtat, bật lại: \logochimbat).
%  Cách dùng ở file chính (xem mau-chuan.tex, mau-chuyen-de.tex):
%      \def\CheDo{loigiai}          % loigiai | dethi | solcolor | color | book
%      \def\tenchuong{<dòng chữ trên đầu trang>}
%      \def\thuonghieu{<chữ ở chân trang lẻ>}   % không khai báo thì dùng mặc định
%      \def\tacgia{<chữ ở chân trang chẵn>}      % không khai báo thì dùng mặc định
%      \def\khauhieu{<chữ cạnh logo ở đầu trang>}  % tùy chọn, mặc định Eduzy.com.vn
%      \input{preamble-eduzy-v2}
%      \TLdapso                      % tài liệu: \TLfull | \TLdapso | \TLan
%      \begin{document} ... \end{document}
%  Logo Eduzy vẽ bằng TikZ ngay trong file (\eduzylogo), không cần file ảnh.
%  Biên dịch: pdfLaTeX (máy có gói vietnam) hoặc XeLaTeX.
%============================================================
\documentclass[12pt,a4paper,twoside]{article}%<-- twoside: bat dau trang/chan trang chan-le
\usepackage{iftex}
\usepackage{amsmath,amssymb,tikz,mathrsfs,tkz-tab}
\usepackage{graphicx}
\usepackage{pgfplots}% đồ thị hàm số
\pgfplotsset{compat=1.17}
%\usepackage{mathptmx}
\usepackage{mathpazo}
\ifXeTeX% XeLaTeX: font chữ có đủ dấu tiếng Việt, giữ nguyên font toán mathpazo
	\usepackage[no-math]{fontspec}
	\setmainfont{TeX Gyre Pagella}
	\newfontfamily\EZqag{TeX Gyre Adventor}[NFSSFamily=qag,Ligatures=TeX]
\fi
\usepackage[many]{tcolorbox}
\usetikzlibrary{math,through,calc,intersections,angles,quotes,shapes,shapes.geometric,arrows,patterns,matrix,chains,arrows.meta,decorations.shapes,decorations.fractals,decorations.markings,shadows,decorations.pathreplacing}
\usepackage[top=1.6cm, bottom=1.8cm, left=1.8cm, right=1.8cm]{geometry}% v2: top 1.6cm (bản cũ 1.5cm) để chữ không sát dải đầu trang
\providecommand{\CheDo}{loigiai}% loigiai | dethi | solcolor | color | book
\expandafter\usepackage\expandafter[\CheDo]{ex_test}
% \tenchuong: khai báo ở file chính (trước \input) hoặc \renewcommand sau \input
%============================================================
%  [GỘP TỪ eduzy-de.tex] EDUZY - MÀU NÂU ĐỒNG / VÀNG KIM
%  + ĐẦU TRANG, CHÂN TRANG KIỂU DẢI NGANG (v2), ĐỐI XỨNG TRANG CHẴN / TRANG LẺ.
%  (không cần file eduzy-de.tex riêng nữa)
%  Khối này phải đặt SAU \usepackage[...]{ex_test} và SAU \def\tenchuong{...}
%============================================================
\usepackage{lastpage}
\usepackage{needspace}
%--- DÀN TRANG ĐỀ THI: tách câu (THÊM 05/11/2026; chỉ thêm lệnh mới, KHÔNG đổi hành vi mặc định) ---
%  Hai công tắc, mặc định tắt (= preamble chạy y như cũ). Bật bằng cách đặt TRƯỚC \input{preamble-eduzy-v2}
%  trong file đề (tools/dan_trang/dan_trang.py tự ghi khi đặt \NeedspaceCau):
%    \def\TachCauDS{1}   câu Đúng/Sai được tách: câu dẫn (kèm hình) ở cuối trang, bảng Phát biểu sang trang sau
%                        (câu dẫn và hình không tách nhau, bảng không tách giữa hai trang)
%    \def\TachCauLG{1}   bản lời giải (\CheDo khác dethi): phần đề của câu cao hơn 1/3 trang được phép tách
%  Lệnh dùng trong file đề (giá trị do tools/dan_trang/dan_trang.py đo ở bản dethi):
%    \NeedspaceCau[<cao câu dẫn + hình>]{<cao cả khối câu>}   ngay trên \begin{ex} (câu đầu phần: trên tiêu đề phần)
%        công tắc tắt: đúng bằng \Needspace{<cao cả khối câu>} như cũ;
%        \TachCauDS=1 và có [..] (câu Đúng/Sai được chọn tách): chỉ giữ câu dẫn + hình;
%        \TachCauLG=1 ở bản lời giải: đòi tối đa 1/3 \textheight.
%    \NeedspaceBang{<cao bảng Phát biểu>}%   ngay trước \choiceTFt: \TachCauDS=1 thì giữ trọn bảng trên một trang;
%        công tắc tắt: không làm gì.
\providecommand{\TachCauDS}{0}
\providecommand{\TachCauLG}{0}
\newlength{\EZnsCau}
\newcommand{\NeedspaceCau}[2][]{%
	\setlength{\EZnsCau}{#2}%
	\ifnum\TachCauDS=1 \if\relax\detokenize{#1}\relax\else\setlength{\EZnsCau}{#1}\fi\fi
	\ifnum\TachCauLG=1 \ifdefstring{\kindOPTN}{dethi}{}{\ifdim\EZnsCau>0.3333\textheight\setlength{\EZnsCau}{0.3333\textheight}\fi}\fi
	\Needspace{\EZnsCau}}
\newcommand{\NeedspaceBang}[1]{\ifnum\TachCauDS=1 \Needspace{#1}\fi}

%--- Thang màu: đậm -> nhạt
\definecolor{gold1}{RGB}{126,96,0}    % nâu đồng đậm  #7E6000
\definecolor{gold2}{RGB}{192,144,3}   % vàng kim      #C09003
\definecolor{gold3}{RGB}{255,215,99}  % vàng nhạt     #FFD763
\definecolor{gold4}{RGB}{255,231,153} % vàng rất nhạt #FFE799
\definecolor{gold5}{RGB}{255,243,205} % vàng kem      #FFF3CD
%--- Ghi đè tông xanh/đỏ cũ
\definecolor{myblue}{RGB}{126,96,0}
\colorlet{blue}{gold1}
\colorlet{cyan}{gold2}
\colorlet{teal}{gold2}

%--- Tên chương mặc định (file đề nên tự khai báo lại)
\providecommand{\tenchuong}{}
%--- Chân trang
\providecommand{\thuonghieu}{TOÁN 12 -- Eduzy.com.vn}
\providecommand{\tacgia}{Tô Gia Bảo -- 0353703783}

%--- Chữ cạnh logo ở đầu trang
\providecommand{\khauhieu}{Eduzy.com.vn}

%====================== THÔNG SỐ ĐẦU TRANG, CHÂN TRANG (v2) ======================
%  Dải nền đậm tràn ngang trang; ô chữ vàng kem có cạnh xiên, kèm 2 vạch xiên chuyển màu.
%  Đầu trang lẻ : logo + \khauhieu (trái)   | ô chữ \tenchuong (phải)
%  Chân trang lẻ: ô chữ \thuonghieu (trái)  | Trang (số) (phải)
%  Trang chẵn soi gương trang lẻ; chân trang chẵn ghi \tacgia.
%--- Màu: chỉ dùng thang vàng kim Eduzy
\colorlet{ezNen}{gold1}% dải nền đậm
\colorlet{ezVach1}{gold3}% vạch xiên sát dải nền (đậm hơn)
\colorlet{ezVach2}{gold4}% vạch xiên sát ô chữ (nhạt hơn)
\colorlet{ezO}{gold5}% ô chữ vàng kem
%--- Kích thước
\newlength{\EZdai}\setlength{\EZdai}{1.1cm}% chiều cao dải nền
\newlength{\EZo}\setlength{\EZo}{0.68cm}% chiều cao ô chữ
\newlength{\EZxien}\setlength{\EZxien}{0.55cm}% độ lệch ngang của cạnh xiên
\newlength{\EZvach}\setlength{\EZvach}{0.75cm}% bề rộng mỗi vạch xiên
\newlength{\EZkhe}\setlength{\EZkhe}{0.12cm}% khe trắng giữa vạch và ô chữ
\newlength{\EZdem}\setlength{\EZdem}{0.4cm}% khoảng đệm từ chữ tới cạnh xiên
\newlength{\EZle}\setlength{\EZle}{1.8cm}% lề trái/phải của chữ (bằng lề thân bài)
\newlength{\EZrongmin}\setlength{\EZrongmin}{4.5cm}% ô chữ rộng tối thiểu
\newlength{\EZrong}% bề rộng chữ trong ô (tính tự động)
%--- Font
\newcommand{\EZfontO}{\fontsize{10.5pt}{12pt}\selectfont\bfseries}% chữ trong ô
\newcommand{\EZfontKH}{\fontfamily{qag}\fontsize{11pt}{12pt}\selectfont\bfseries}% \khauhieu cạnh logo
\newcommand{\EZfontTrang}{\fontsize{11.5pt}{12pt}\selectfont\bfseries}% chữ "Trang" và số trang
%--- Đo bề rộng chữ trong ô
\newcommand{\EZdorong}[1]{%
	\settowidth{\EZrong}{\EZfontO #1}%
	\ifdim\EZrong<\EZrongmin\setlength{\EZrong}{\EZrongmin}\fi}
%--- Logo + chữ trắng \khauhieu (dựng sẵn một lần khi vào \begin{document})
\newsavebox{\EZlogobox}
\AtBeginDocument{\sbox{\EZlogobox}{%
		\raisebox{-0.5\height}{\eduzylogo[1cm]}\hspace{0.15cm}%
		\raisebox{-0.5\height}{\color{white}\EZfontKH\khauhieu}}}
%--- Hai vạch xiên + khe trắng. #1: điểm gốc ở mép dưới, #2: hướng (1 sang phải, -1 sang trái),
%    #3: hướng nghiêng (1 kiểu "/", -1 kiểu "\")
\newcommand{\EZvachxien}[3]{%
	\fill[white] (#1) --++({(#3)*\EZxien},\EZo) --++({(#2)*\EZkhe},0) --++({-(#3)*\EZxien},-\EZo) --cycle;
	\fill[ezVach2] ([xshift=(#2)*\EZkhe]#1) --++({(#3)*\EZxien},\EZo) --++({(#2)*\EZvach},0) --++({-(#3)*\EZxien},-\EZo) --cycle;
	\fill[ezVach1] ([xshift=(#2)*(\EZkhe+\EZvach)]#1) --++({(#3)*\EZxien},\EZo) --++({(#2)*\EZvach},0) --++({-(#3)*\EZxien},-\EZo) --cycle;
	\draw[white,line width=0.6pt] ([xshift=(#2)*(\EZkhe+\EZvach)]#1) --++({(#3)*\EZxien},\EZo);
}
%--- Chữ trong ô (căn giữa theo chiều cao ô, không lệch vì dấu tiếng Việt)
\tikzset{ezchu/.style={inner sep=0pt,text=ezNen,text height=0.72em,text depth=0.22em}}

%====================== LOGO CHÌM GIỮA TRANG ======================
%  Vẽ cùng đầu trang nên nằm dưới chữ của thân bài, có trên mọi trang.
%  Tắt: \logochimtat (từ trang đó trở đi), bật lại: \logochimbat.
%  Đổi cỡ, độ đậm: khai báo trước \input, ví dụ \def\EZchimcao{9cm} \def\EZchimdo{0.1}
\providecommand{\EZchimcao}{10cm}% chiều cao logo chìm
\providecommand{\EZchimdo}{0.12}% độ đậm (0 = trong suốt, 1 = rõ hẳn)
\newif\ifEZchim \EZchimtrue
\newcommand{\logochimtat}{\EZchimfalse}
\newcommand{\logochimbat}{\EZchimtrue}
%  Dựng sẵn logo mờ một lần khi vào \begin{document}; transparency group làm mờ cả logo
%  một lượt nên không lộ vệt ở chỗ các dải màu chồng nhau.
\newsavebox{\EZchimtho}
\newsavebox{\EZchimbox}
\AtBeginDocument{%
	\sbox{\EZchimtho}{\eduzylogo[\EZchimcao]}%
	\sbox{\EZchimbox}{\begin{tikzpicture}
			\begin{scope}[transparency group,opacity=\EZchimdo]
				\node[inner sep=0pt] {\usebox{\EZchimtho}};
			\end{scope}
	\end{tikzpicture}}}
\newcommand{\EZvechim}{%
	\ifEZchim
		\node[inner sep=0pt] at (current page.center) {\usebox{\EZchimbox}};
	\fi}

%====================== ĐẦU TRANG - TRANG LẺ ======================
%  (logo bên trái, ô chữ \tenchuong bên phải)
\newcommand{\myfancyheadOdd}{%
	\EZdorong{\tenchuong}%
	\begin{tikzpicture}[remember picture,overlay]
		\coordinate (T) at (current page.north west);
		\coordinate (P) at (current page.north east);
		%--- logo chìm giữa trang
		\EZvechim
		%--- dải nền
		\fill[ezNen] (T) rectangle ([yshift=-\EZdai]P);
		%--- ô chữ: X là góc dưới bên trái, cạnh trái xiên kiểu "/"
		\coordinate (X) at ([xshift=-\EZle-\EZrong-\EZdem-\EZxien,yshift=-\EZdai]P);
		\fill[ezO] (X) --++(\EZxien,\EZo) -- ([yshift=-\EZdai+\EZo]P) -- ([yshift=-\EZdai]P) --cycle;
		\EZvachxien{X}{-1}{1}
		\node[ezchu,anchor=east] at ([xshift=-\EZle,yshift=-\EZdai+0.5*\EZo]P) {\EZfontO\tenchuong};
		%--- logo + tên
		\node[inner sep=0pt,anchor=west] at ([xshift=\EZle,yshift=-0.5*\EZdai]T) {\usebox{\EZlogobox}};
	\end{tikzpicture}%
}

%====================== ĐẦU TRANG - TRANG CHẴN =====================
%  (soi gương: ô chữ \tenchuong bên trái, logo bên phải)
\newcommand{\myfancyheadEven}{%
	\EZdorong{\tenchuong}%
	\begin{tikzpicture}[remember picture,overlay]
		\coordinate (T) at (current page.north west);
		\coordinate (P) at (current page.north east);
		%--- logo chìm giữa trang
		\EZvechim
		%--- dải nền
		\fill[ezNen] (T) rectangle ([yshift=-\EZdai]P);
		%--- ô chữ: X là góc dưới bên phải, cạnh phải xiên kiểu "\"
		\coordinate (X) at ([xshift=\EZle+\EZrong+\EZdem+\EZxien,yshift=-\EZdai]T);
		\fill[ezO] (X) --++(-\EZxien,\EZo) -- ([yshift=-\EZdai+\EZo]T) -- ([yshift=-\EZdai]T) --cycle;
		\EZvachxien{X}{1}{-1}
		\node[ezchu,anchor=west] at ([xshift=\EZle,yshift=-\EZdai+0.5*\EZo]T) {\EZfontO\tenchuong};
		%--- logo + tên
		\node[inner sep=0pt,anchor=east] at ([xshift=-\EZle,yshift=-0.5*\EZdai]P) {\usebox{\EZlogobox}};
	\end{tikzpicture}%
}

%====================== CHÂN TRANG - TRANG LẺ ======================
%  (ô chữ \thuonghieu bên trái, "Trang (số)" bên phải)
\newcommand{\myfancyfootOdd}{%
	\EZdorong{\thuonghieu}%
	\begin{tikzpicture}[remember picture,overlay]
		\coordinate (T) at (current page.south west);
		\coordinate (P) at (current page.south east);
		%--- dải nền
		\fill[ezNen] (T) rectangle ([yshift=\EZdai]P);
		%--- ô chữ sát mép trên dải: X là góc dưới bên phải, cạnh phải xiên kiểu "/"
		\coordinate (X) at ([xshift=\EZle+\EZrong+\EZdem,yshift=\EZdai-\EZo]T);
		\fill[ezO] ([yshift=\EZdai-\EZo]T) -- (X) --++(\EZxien,\EZo) -- ([yshift=\EZdai]T) --cycle;
		\EZvachxien{X}{1}{1}
		\node[ezchu,anchor=west] at ([xshift=\EZle,yshift=\EZdai-0.5*\EZo]T) {\EZfontO\thuonghieu};
		%--- số trang
		\node[circle,draw=white,line width=0.9pt,text=white,inner sep=0pt,minimum size=0.62cm,anchor=east]
		(so) at ([xshift=-\EZle,yshift=0.5*\EZdai]P) {\EZfontTrang\thepage};
		\node[inner sep=0pt,text=white,anchor=base east] (tr) at ([xshift=-0.15cm]so.base west) {\EZfontTrang Trang};
	\end{tikzpicture}%
}

%====================== CHÂN TRANG - TRANG CHẴN ====================
%  (soi gương: "Trang (số)" bên trái, ô chữ \tacgia bên phải)
\newcommand{\myfancyfootEven}{%
	\EZdorong{\tacgia}%
	\begin{tikzpicture}[remember picture,overlay]
		\coordinate (T) at (current page.south west);
		\coordinate (P) at (current page.south east);
		%--- dải nền
		\fill[ezNen] (T) rectangle ([yshift=\EZdai]P);
		%--- ô chữ sát mép trên dải: X là góc dưới bên trái, cạnh trái xiên kiểu "\"
		\coordinate (X) at ([xshift=-\EZle-\EZrong-\EZdem,yshift=\EZdai-\EZo]P);
		\fill[ezO] ([yshift=\EZdai-\EZo]P) -- (X) --++(-\EZxien,\EZo) -- ([yshift=\EZdai]P) --cycle;
		\EZvachxien{X}{-1}{-1}
		\node[ezchu,anchor=east] at ([xshift=-\EZle,yshift=\EZdai-0.5*\EZo]P) {\EZfontO\tacgia};
		%--- số trang
		%    (vòng tròn ẩn so0 nằm giữa dải để lấy đường chân chữ, giống trang lẻ)
		\node[circle,inner sep=0pt,minimum size=0.62cm] (so0) at ([yshift=0.5*\EZdai]T) {\EZfontTrang\phantom{\thepage}};
		\node[inner sep=0pt,text=white,anchor=base west] (tr) at ($(so0.base -| T)+(\EZle,0)$) {\EZfontTrang Trang};
		\node[circle,draw=white,line width=0.9pt,text=white,inner sep=0pt,minimum size=0.62cm,anchor=base west]
		(so) at ([xshift=0.15cm]tr.base east) {\EZfontTrang\thepage};
	\end{tikzpicture}%
}

%--- Giữ tương thích với các file cũ đang gọi \myfancyhead, \myfancyfoot
\newcommand{\myfancyhead}{\myfancyheadOdd}
\newcommand{\myfancyfoot}{\myfancyfootOdd}

%======================= GẮN VÀO TRANG =============================
\pagestyle{fancy}
\fancyhf{}
\fancyhead[LO]{\myfancyheadOdd}
\fancyhead[RE]{\myfancyheadEven}
\fancyfoot[LO]{\myfancyfootOdd}
\fancyfoot[RE]{\myfancyfootEven}
\renewcommand{\headrulewidth}{0pt}
\renewcommand{\footrulewidth}{0pt}
\fancypagestyle{plain}{%
	\fancyhf{}
	\fancyhead[LO]{\myfancyheadOdd}
	\fancyhead[RE]{\myfancyheadEven}
	\fancyfoot[LO]{\myfancyfootOdd}
	\fancyfoot[RE]{\myfancyfootEven}
	\renewcommand{\headrulewidth}{0pt}
	\renewcommand{\footrulewidth}{0pt}
}
%=================== HẾT PHẦN GỘP TỪ eduzy-de.tex ===================
\newcommand{\indam}[1]{
	{\fontsize{12pt}{9pt}\sffamily \selectfont\bfseries\color{gold1} #1}
}
\newcommand{\ind}[1]{
	{\fontsize{14pt}{9pt}\fontfamily{qag} \selectfont\bfseries\color{gold1} #1}
}
\renewcommand{\FalseEX}{\stepcounter{dapan}\textbf{\selectfont\color{gold1}\Alph{dapan}.}\ignorespaces}
\renewtheorem{ex}{\color{gold1}Câu}
\renewcommand{\baselinestretch}{1.1}
\newcommand{\hoac}[1]{
	\left[\begin{aligned}#1\end{aligned}\right.}
\newcommand{\heva}[1]{
	\left\{\begin{aligned}#1\end{aligned}\right.}
\usepackage{esvect}
\def\vec{\vv}
\def\vt{\vv}
\def\overrightarrow{\vv}
%====================== LOGO EDUZY VẼ BẰNG TIKZ ======================
%  Đường viền dò từ logo gốc; dải màu vàng kim đo trực tiếp trên ảnh gốc
%  (nghiêng 21 độ, tô bằng 80 dải hẹp nên hiển thị giống nhau trên mọi trình xem PDF).
%  Dùng: \eduzylogo          -> cao 2.4cm (mặc định)
%        \eduzylogo[1.5cm]   -> cao 1.5cm
\definecolor{ezText}{RGB}{245,200,46}
\newcommand{\eduzylogo}[1][2.4cm]{%
	\resizebox{!}{#1}{%
		\begin{tikzpicture}[x=0.01cm,y=0.01cm]
			\path[use as bounding box] (0,0) rectangle (224,-224);
			% mũ tốt nghiệp + vương miện
			\begin{scope}
				\clip
				(119.4,-176.1)..controls(118.9,-176.1)and(116.8,-176)..(114.6,-176)
				..controls(101.4,-176)and(85.8,-175.9)..(83.7,-175.7)
				..controls(81.8,-175.6)and(81,-175.6)..(80.2,-175.7)
				..controls(78.8,-175.9)and(77.7,-175.9)..(76.2,-175.7)
				..controls(74.9,-175.5)and(70.8,-175.5)..(64.3,-175.7)
				..controls(59.1,-175.9)and(58.3,-175.7)..(58.6,-173.7)
				..controls(58.6,-172.8)and(59.6,-169.6)..(60.3,-168)
				..controls(60.8,-166.9)and(61.4,-165.8)..(62.3,-165)
				..controls(63.1,-164.2)and(64.3,-163.4)..(65,-163.2)
				..controls(65.5,-163.1)and(72.3,-163)..(76.6,-163.1)
				..controls(77.9,-163.1)and(80.1,-163.1)..(81.4,-163)
				..controls(82.8,-162.9)and(86.3,-162.9)..(89.2,-162.9)
				..controls(93,-162.9)and(94.9,-162.8)..(96,-162.7)
				..controls(97.2,-162.5)and(101,-162.5)..(114,-162.4)
				..controls(123,-162.3)and(134.9,-162.2)..(140.2,-162.1)
				..controls(150.1,-162)and(152.2,-162.1)..(153.6,-162.4)
				..controls(154.4,-162.6)and(157.1,-163.8)..(158.1,-164.4)
				..controls(158.8,-164.8)and(159.7,-165.7)..(159.9,-166.3)
				..controls(160,-166.5)and(160.2,-167.1)..(160.2,-167.6)
				..controls(160.3,-168)and(160.6,-169.4)..(160.9,-170.6)
				..controls(161.2,-171.7)and(161.5,-173.1)..(161.6,-173.6)
				..controls(161.8,-174.4)and(161.8,-174.5)..(161.5,-174.9)
				..controls(161.2,-175.6)and(160.8,-175.9)..(159.9,-176)
				..controls(159.2,-176.1)and(154.1,-176.2)..(151.5,-176.2)
				..controls(148.1,-176.1)and(122.6,-176.1)..(121.6,-176.2)
				..controls(120.8,-176.2)and(119.8,-176.2)..(119.4,-176.1)--cycle
				(130.4,-160)..controls(129,-160)and(119.8,-160)..(110.1,-159.9)
				..controls(100.3,-159.8)and(88.5,-159.7)..(83.9,-159.6)
				..controls(77.8,-159.5)and(75.3,-159.5)..(74.5,-159.6)
				..controls(73.9,-159.7)and(72.5,-159.8)..(71.3,-159.8)
				..controls(69.6,-159.8)and(69,-159.8)..(68.5,-159.6)
				..controls(67.6,-159.3)and(66.8,-158.6)..(66.2,-157.6)
				..controls(66,-157.3)and(65.6,-156.5)..(65.3,-156.1)
				..controls(64.6,-155.1)and(64.5,-154.9)..(63.7,-152.9)
				..controls(63.2,-151.6)and(62.9,-150.9)..(61.8,-149.6)
				..controls(61.5,-149.2)and(61.1,-148.5)..(60.9,-147.9)
				..controls(60.7,-147.4)and(60.3,-146.6)..(60.1,-146.1)
				..controls(59.9,-145.5)and(59.5,-144.9)..(59,-144.3)
				..controls(58.5,-143.7)and(58.2,-143.2)..(57.9,-142.5)
				..controls(57.7,-141.9)and(57.3,-141.1)..(57,-140.7)
				..controls(56.6,-140.3)and(56.1,-139.3)..(55.8,-138.5)
				..controls(55.5,-137.8)and(55.1,-136.8)..(54.9,-136.5)
				..controls(54.7,-136.1)and(54.3,-135.2)..(53.9,-134.5)
				..controls(53.6,-133.8)and(53,-132.7)..(52.6,-132)
				..controls(51.6,-130.5)and(51.4,-130.2)..(51.1,-129)
				..controls(50.6,-127.5)and(50.6,-127.4)..(49.7,-126.3)
				..controls(49.2,-125.7)and(48.6,-124.8)..(48.3,-124.2)
				..controls(48,-123.6)and(47.5,-122.7)..(47.3,-122.2)
				..controls(47,-121.8)and(46.5,-120.7)..(46.1,-119.8)
				..controls(45.8,-118.9)and(45.2,-117.8)..(44.9,-117.2)
				..controls(44.2,-116.1)and(43.3,-114.4)..(42.3,-112.2)
				..controls(41.8,-111.3)and(41.3,-110.4)..(41.1,-110.1)
				..controls(40.5,-109.4)and(40.3,-109)..(40,-107.8)
				..controls(39.8,-107)and(39.6,-106.5)..(39.2,-106)
				..controls(38.9,-105.6)and(38.4,-104.7)..(38,-104)
				..controls(37.7,-103.4)and(37.2,-102.5)..(37,-102.1)
				..controls(36.7,-101.7)and(36.3,-100.8)..(36,-100.1)
				..controls(35.6,-98.9)and(35.3,-98.4)..(34.6,-97.6)
				..controls(34.2,-97.1)and(34.2,-96.3)..(34.6,-96)
				..controls(35,-95.7)and(35.5,-95.9)..(36.2,-96.7)
				..controls(36.5,-97)and(37.2,-97.5)..(37.6,-97.7)..controls(38,-98)and(38.6,-98.4)..(39.1,-98.9)
				..controls(39.6,-99.4)and(40.3,-100)..(40.6,-100.3)
				..controls(41.3,-100.8)and(42.4,-101.8)..(43.7,-102.9)
				..controls(44,-103.2)and(44.6,-103.7)..(45,-103.9)
				..controls(45.4,-104.2)and(45.9,-104.6)..(46.1,-104.8)
				..controls(47.1,-105.9)and(48.8,-107.1)..(50.5,-107.9)
				..controls(51.8,-108.6)and(52.3,-109)..(53.4,-110)
				..controls(54.8,-111.2)and(55.6,-111.6)..(57,-112.1)
				..controls(57.5,-112.3)and(58.5,-112.8)..(59.4,-113.3)
				..controls(60.2,-113.7)and(61.3,-114.3)..(61.9,-114.5)
				..controls(62.4,-114.7)and(63.3,-115.1)..(63.8,-115.3)
				..controls(64.4,-115.6)and(65.4,-115.9)..(66.2,-116.1)
				..controls(67,-116.3)and(68.4,-116.6)..(69.2,-116.9)
				..controls(72.3,-117.9)and(72.8,-118)..(78,-118.1)
				..controls(84,-118.3)and(86.1,-118.1)..(88.1,-117.1)
				..controls(89.4,-116.5)and(93,-115)..(93.7,-114.9)
				..controls(94.8,-114.6)and(95.4,-114.2)..(95.9,-113.4)
				..controls(96.2,-113)and(96.7,-112.5)..(97.1,-112.2)
				..controls(98,-111.6)and(101.3,-108.4)..(102.1,-107.2)
				..controls(102.4,-106.8)and(102.9,-106.2)..(103.2,-105.8)
				..controls(103.5,-105.5)and(104,-104.8)..(104.2,-104.2)
				..controls(104.5,-103.6)and(105,-102.8)..(105.3,-102.4)
				..controls(105.6,-102)and(106.1,-101.2)..(106.5,-100.6)
				..controls(107.1,-99.5)and(108.6,-97.5)..(109.5,-96.5)
				..controls(109.9,-96.1)and(110.1,-96)..(110.4,-96)
				..controls(111.1,-96)and(111.6,-96.5)..(111.8,-97.6)
				..controls(112,-98.6)and(112.3,-99.3)..(112.8,-99.9)
				..controls(113.7,-100.9)and(115.2,-103.1)..(115.7,-104.2)
				..controls(116.1,-105)and(116.5,-105.5)..(116.9,-106)
				..controls(117.3,-106.4)and(117.8,-106.9)..(118,-107.3)
				..controls(118.2,-107.6)and(118.6,-108)..(119,-108.3)
				..controls(119.6,-108.8)and(121.1,-110.3)..(121.8,-111.2)
				..controls(122.1,-111.7)and(122.5,-112)..(123.4,-112.4)
				..controls(124.8,-113.2)and(125.8,-113.8)..(127,-114.6)
				..controls(127.5,-114.9)and(128.1,-115.3)..(128.2,-115.4)
				..controls(128.4,-115.5)and(129.3,-115.8)..(130.2,-116)
				..controls(131.1,-116.2)and(132.5,-116.6)..(133.3,-116.9)
				..controls(135.7,-117.9)and(136.1,-117.9)..(142.5,-117.8)
				..controls(145.8,-117.8)and(148.4,-117.7)..(148.8,-117.6)
				..controls(149.2,-117.6)and(150.2,-117.2)..(151,-116.8)
				..controls(152,-116.4)and(152.8,-116)..(153.6,-115.9)
				..controls(154.5,-115.7)and(155.2,-115.4)..(156.3,-114.9)
				..controls(157.1,-114.5)and(158.2,-114)..(158.7,-113.8)
				..controls(159.1,-113.7)and(159.8,-113.3)..(160.2,-113)
				..controls(160.6,-112.7)and(161.3,-112.4)..(161.7,-112.2)
				..controls(162.7,-111.9)and(163.2,-111.7)..(164.1,-111.1)
				..controls(164.5,-110.8)and(165.3,-110.4)..(165.9,-110)
				..controls(166.5,-109.7)and(167.3,-109.2)..(167.6,-108.9)
				..controls(167.9,-108.6)and(168.6,-108.2)..(169,-108)
				..controls(169.4,-107.8)and(170.2,-107.3)..(170.7,-106.9)
				..controls(171.2,-106.5)and(172.1,-105.8)..(172.7,-105.4)
				..controls(173.2,-105)and(174.1,-104.4)..(174.5,-104)
				..controls(174.9,-103.7)and(175.7,-103.1)..(176.2,-102.6)
				..controls(176.7,-102.2)and(177.4,-101.5)..(177.9,-101.2)
				..controls(178.3,-100.8)and(178.8,-100.4)..(179,-100.2)
				..controls(179.2,-100)and(179.6,-99.6)..(180,-99.3)
				..controls(180.4,-99.1)and(180.9,-98.6)..(181.2,-98.4)
				..controls(181.5,-98.1)and(182.1,-97.7)..(182.5,-97.5)
				..controls(183.2,-97.2)and(183.5,-96.9)..(184.3,-95.9)
				..controls(184.6,-95.6)and(184.6,-95.6)..(184.8,-95.7)
				..controls(185,-95.8)and(185.1,-96)..(185.1,-96.2)
				..controls(185.1,-96.5)and(185,-96.7)..(184.6,-97.2)
				..controls(184.1,-97.8)and(184,-97.9)..(183.9,-98.8)
				..controls(183.7,-100.3)and(183.3,-101.3)..(182.9,-101.8)
				..controls(182.6,-102.1)and(182.3,-102.6)..(182.2,-102.9)
				..controls(182,-103.3)and(181.6,-104)..(181.2,-104.6)
				..controls(180.9,-105.2)and(180.5,-106.1)..(180.4,-106.6)
				..controls(179.9,-108)and(179.7,-108.5)..(179,-109.4)
				..controls(178.6,-109.8)and(178.2,-110.5)..(177.9,-111)
				..controls(177.7,-111.5)and(177.3,-112.3)..(177,-112.7)
				..controls(176.6,-113.2)and(176.4,-113.7)..(176.3,-114.2)
				..controls(176.2,-114.7)and(175.9,-115.2)..(175.4,-116)
				..controls(175,-116.6)and(174.4,-117.7)..(174,-118.5)
				..controls(173.7,-119.2)and(173.3,-120.1)..(173,-120.5)
				..controls(172.8,-120.9)and(172.4,-121.8)..(172.1,-122.5)
				..controls(171.8,-123.2)and(171.3,-124.1)..(171,-124.6)
				..controls(170.8,-125.1)and(170.3,-126)..(169.9,-126.6)
				..controls(169.6,-127.3)and(169.2,-128.1)..(169,-128.3)
				..controls(168.8,-128.6)and(168.5,-129.3)..(168.3,-129.8)
				..controls(168.1,-130.3)and(167.8,-131.1)..(167.5,-131.6)
				..controls(166.9,-132.5)and(164.2,-137.9)..(164.2,-138.1)
				..controls(164.2,-138.1)and(163.9,-138.8)..(163.4,-139.7)
				..controls(162.9,-140.5)and(162.4,-141.5)..(162.3,-141.9)
				..controls(161.9,-143)and(161.6,-143.4)..(160.9,-144.1)
				..controls(160.4,-144.6)and(160.3,-144.8)..(160.1,-145.4)
				..controls(160,-145.8)and(159.6,-146.8)..(159.2,-147.8)
				..controls(158.7,-148.8)and(158.1,-150)..(157.9,-150.6)
				..controls(157.6,-151.1)and(157.2,-151.9)..(156.9,-152.4)
				..controls(156.5,-152.8)and(156.1,-153.7)..(155.8,-154.4)
				..controls(155.6,-155)and(155.2,-155.9)..(154.9,-156.3)
				..controls(154.7,-156.7)and(154.4,-157.4)..(154.2,-157.9)
				..controls(154.1,-158.4)and(153.8,-159)..(153.5,-159.3)
				..controls(152.8,-160.1)and(153,-160.1)..(142.4,-160.1)
				..controls(137.2,-160.1)and(131.8,-160.1)..(130.4,-160)--cycle
				(141.5,-103.7)..controls(140.8,-103)and(139.8,-102.1)..(139.2,-101.6)
				..controls(138.5,-101.2)and(137.9,-100.6)..(137.8,-100.4)
				..controls(137,-99.5)and(134,-96.5)..(133.2,-96)
				..controls(132.7,-95.7)and(131.8,-94.9)..(131.1,-94.3)
				..controls(129.4,-92.7)and(127.9,-91.6)..(126.6,-91.1)
				..controls(125.2,-90.5)and(124.1,-90)..(122.9,-89.1)
				..controls(122.1,-88.6)and(121.3,-88.2)..(119.9,-87.7)
				..controls(118.8,-87.3)and(117.5,-86.8)..(116.9,-86.5)
				..controls(114.3,-85.5)and(113.7,-85.4)..(109.7,-85.5)
				..controls(105.3,-85.7)and(103.9,-85.9)..(101.6,-87)
				..controls(100.7,-87.4)and(99.4,-87.8)..(97.9,-88)
				..controls(97,-88.2)and(96.6,-88.3)..(95.8,-88.9)
				..controls(95.4,-89.2)and(94.7,-89.6)..(94.2,-89.8)
				..controls(92.9,-90.3)and(92.3,-90.7)..(90.9,-92.1)
				..controls(90.2,-92.7)and(89.4,-93.3)..(89,-93.5)..controls(87,-94.5)and(85.2,-96)..(83.6,-97.9)
				..controls(83.3,-98.2)and(82.5,-98.9)..(81.8,-99.4)
				..controls(80.7,-100.1)and(80.5,-100.3)..(80,-101.1)
				..controls(78.7,-103.1)and(77.2,-103.9)..(76.4,-103)
				..controls(75.7,-102.2)and(75.7,-103.2)..(75.7,-83.6)
				..controls(75.6,-73.6)and(75.6,-64.8)..(75.5,-64)
				..controls(75.5,-63.2)and(75.5,-61.4)..(75.6,-59.8)--(75.7,-57.1)--(76,-56.8)
				..controls(76.8,-56.2)and(78,-56.3)..(79.8,-57.2)
				..controls(80.4,-57.5)and(81.3,-57.8)..(81.8,-57.9)
				..controls(82.3,-58.1)and(83,-58.4)..(83.6,-58.8)..controls(84.3,-59.2)and(85,-59.5)..(86,-59.8)
				..controls(86.8,-60.1)and(88,-60.6)..(88.7,-60.9)
				..controls(89.3,-61.2)and(90.6,-61.7)..(91.4,-62)
				..controls(92.2,-62.3)and(93.7,-62.9)..(94.6,-63.3)
				..controls(95.6,-63.7)and(96.8,-64.1)..(97.4,-64.3)
				..controls(98,-64.5)and(99,-64.9)..(99.7,-65.2)
				..controls(100.4,-65.6)and(101.5,-66)..(102.3,-66.2)
				..controls(104.2,-66.8)and(105,-67.1)..(106.2,-67.8)
				..controls(107.9,-68.7)and(110.2,-69.4)..(111.3,-69.2)
				..controls(112.5,-69)and(115,-68.2)..(115.6,-67.8)
				..controls(117.3,-66.7)and(117.6,-66.6)..(119.5,-66.2)
				..controls(121.9,-65.8)and(122.1,-65.7)..(123,-65.1)
				..controls(124.1,-64.4)and(124.4,-64.3)..(126,-63.9)
				..controls(127,-63.7)and(127.7,-63.4)..(128.4,-63.1)
				..controls(129,-62.8)and(129.7,-62.5)..(130.6,-62.4)
				..controls(132.1,-62.1)and(134,-61.5)..(135.2,-61)
				..controls(135.7,-60.8)and(136.7,-60.5)..(137.3,-60.3)
				..controls(138.4,-60)and(139.3,-59.7)..(141.1,-59)
				..controls(141.6,-58.7)and(142,-58.6)..(142.5,-58.6)
				..controls(143.2,-58.6)and(143.3,-58.7)..(143.6,-59)
				..controls(144,-59.5)and(144.1,-59.9)..(144,-60.8)
				..controls(144,-61.2)and(144,-71.1)..(144,-82.8)--(143.9,-104.1)--(143.7,-104.5)
				..controls(143.5,-104.7)and(143.3,-104.8)..(143,-104.9)
				..controls(142.6,-104.9)and(142.5,-104.8)..(141.5,-103.7)--cycle
				(27.8,-91.4)..controls(27.1,-90.9)and(26.3,-90)..(26.1,-89.5)
				..controls(26,-89.2)and(25.8,-88.6)..(25.7,-87.9)
				..controls(25.4,-86.2)and(25.6,-84.8)..(26.5,-83.5)
				..controls(27.4,-82.3)and(28.3,-82)..(29.4,-82.3)
				..controls(31.1,-82.8)and(32.1,-84.9)..(32,-87.4)
				..controls(31.9,-89.3)and(31.2,-90.6)..(29.9,-91.4)
				..controls(29.1,-91.9)and(28.6,-91.9)..(27.8,-91.4)--cycle
				(190.1,-91.3)..controls(189.7,-91.2)and(189.3,-90.9)..(189.1,-90.7)
				..controls(188.3,-89.8)and(187.6,-87.7)..(187.6,-86.2)
				..controls(187.6,-85.5)and(187.7,-85.2)..(188.2,-84.4)
				..controls(189.1,-82.6)and(190.1,-81.8)..(191.3,-82.1)
				..controls(192,-82.3)and(193.5,-83.7)..(193.8,-84.4)
				..controls(194.1,-85.1)and(194.2,-87.1)..(194,-88.3)
				..controls(193.8,-89.8)and(193.3,-90.5)..(192.3,-91.3)
				..controls(191.7,-91.7)and(190.9,-91.7)..(190.1,-91.3)--cycle
				(173.7,-83.4)..controls(172.2,-82.5)and(172,-82.5)..(170.6,-82.4)
				..controls(169.7,-82.4)and(168.8,-82.4)..(168.3,-82.5)
				..controls(167.2,-82.7)and(166.9,-82.7)..(166.6,-82.2)
				..controls(166.3,-81.8)and(166.3,-81)..(166.8,-80)
				..controls(167.2,-78.9)and(167.4,-77.4)..(167.7,-74.1)
				..controls(167.8,-72.8)and(167.9,-71.8)..(167.9,-71.7)
				..controls(168,-71.7)and(168,-70.2)..(168,-68.4)..controls(168,-65.4)and(168,-64.9)..(167.7,-63)
				..controls(167.3,-60.5)and(167.1,-59.9)..(166.4,-58.4)
				..controls(166.1,-57.7)and(165.8,-56.9)..(165.7,-56.6)
				..controls(165.7,-56.2)and(165.4,-55.6)..(165.2,-55.2)
				..controls(164.6,-53.8)and(161.9,-50.7)..(161,-50.3)
				..controls(159.8,-49.7)and(157.9,-50)..(155.2,-51.2)
				..controls(154.7,-51.4)and(153.4,-51.8)..(152.2,-52.1)
				..controls(150.8,-52.4)and(149.4,-52.9)..(148.6,-53.2)
				..controls(147.8,-53.6)and(146.7,-54)..(146,-54.2)
				..controls(145.4,-54.4)and(144.5,-54.7)..(144,-54.9)
				..controls(143.6,-55.2)and(142.7,-55.5)..(141.8,-55.7)
				..controls(141,-55.9)and(139.6,-56.3)..(138.7,-56.6)
				..controls(137.8,-57)and(136.3,-57.4)..(135.3,-57.6)
				..controls(133.8,-57.9)and(133.2,-58.1)..(132.1,-58.7)
				..controls(131.1,-59.2)and(130.3,-59.4)..(129.4,-59.6)
				..controls(128.7,-59.8)and(127.8,-60)..(127.4,-60.1)
				..controls(127,-60.3)and(126,-60.5)..(125.1,-60.7)
				..controls(123.2,-61.1)and(121.7,-61.5)..(120.2,-62.1)
				..controls(119.7,-62.3)and(118.1,-62.8)..(116.7,-63.2)
				..controls(114.6,-63.9)and(114,-64.1)..(113,-64.7)
				..controls(111.9,-65.4)and(111.8,-65.4)..(111,-65.5)
				..controls(109.9,-65.5)and(109.3,-65.4)..(107.7,-64.7)
				..controls(107,-64.4)and(106,-64)..(105.3,-63.9)
				..controls(103.5,-63.5)and(102.2,-63)..(101.2,-62.5)
				..controls(100.1,-61.8)and(99.3,-61.5)..(97.8,-61.1)
				..controls(97.2,-61)and(96.2,-60.7)..(95.7,-60.4)
				..controls(95.1,-60.1)and(94.2,-59.8)..(93.5,-59.6)
				..controls(92.9,-59.4)and(91.7,-59)..(91,-58.7)..controls(90.2,-58.4)and(88.9,-58)..(88,-57.7)
				..controls(87.2,-57.5)and(86.1,-57.1)..(85.5,-56.8)
				..controls(85,-56.6)and(84.1,-56.2)..(83.5,-56)..controls(82.9,-55.7)and(81.8,-55.2)..(81,-54.9)
				..controls(80.1,-54.5)and(79.1,-54.2)..(78.7,-54.1)
				..controls(77.9,-53.8)and(77.3,-53.6)..(76.2,-52.9)
				..controls(75.7,-52.6)and(75.2,-52.4)..(74.1,-52.2)
				..controls(72.5,-51.9)and(72.4,-51.8)..(71.3,-51.1)
				..controls(69.9,-50.1)and(69.4,-49.9)..(67.9,-49.8)
				..controls(67,-49.7)and(66.6,-49.6)..(66.3,-49.4)
				..controls(66.1,-49.3)and(65.7,-49)..(65.4,-48.8)
				..controls(65.1,-48.7)and(64.2,-48.4)..(63.6,-48.2)
				..controls(62.9,-48)and(61.7,-47.6)..(60.9,-47.2)
				..controls(60.1,-46.8)and(58.9,-46.4)..(58.2,-46.2)
				..controls(57.6,-46)and(56.8,-45.7)..(56.5,-45.5)
				..controls(54.9,-44.6)and(54.6,-44.5)..(53.2,-44.2)
				..controls(51.7,-43.8)and(51,-43.6)..(49.9,-42.9)
				..controls(49.5,-42.6)and(48.9,-42.4)..(47.5,-42)
				..controls(45.7,-41.6)and(45.6,-41.5)..(44.3,-40.6)
				..controls(43.5,-40.1)and(42.8,-39.6)..(42.7,-39.4)
				..controls(42.4,-38.7)and(43,-38.1)..(44.5,-37.6)..controls(44.9,-37.5)and(46,-37)..(46.9,-36.6)
				..controls(47.9,-36.1)and(48.8,-35.8)..(49.6,-35.6)
				..controls(50.5,-35.4)and(51.3,-35.2)..(52.4,-34.6)
				..controls(53.3,-34.1)and(54.5,-33.7)..(54.9,-33.5)
				..controls(55.4,-33.4)and(56.4,-33)..(57.2,-32.6)
				..controls(58.6,-32)and(60.1,-31.5)..(61.6,-31.2)
				..controls(62.3,-31.1)and(62.7,-31)..(63.2,-30.7)
				..controls(64.1,-30)and(64.3,-29.9)..(65.3,-29.7)
				..controls(65.7,-29.6)and(66.5,-29.2)..(67.1,-28.9)
				..controls(67.6,-28.6)and(68.8,-28.2)..(69.6,-27.9)
				..controls(70.5,-27.7)and(71.6,-27.3)..(72.1,-27)
				..controls(72.6,-26.7)and(73.8,-26.2)..(74.8,-25.9)
				..controls(75.7,-25.5)and(76.7,-25.1)..(77.1,-25)
				..controls(77.4,-24.8)and(78.5,-24.4)..(79.6,-24.1)
				..controls(80.6,-23.8)and(82,-23.2)..(82.8,-22.9)
				..controls(83.5,-22.5)and(84.8,-22)..(85.8,-21.8)
				..controls(86.7,-21.5)and(87.8,-21.1)..(88.1,-20.9)
				..controls(88.4,-20.7)and(89.3,-20.3)..(90.1,-20.1)
				..controls(90.8,-19.8)and(91.9,-19.4)..(92.4,-19.1)
				..controls(93.5,-18.5)and(94.4,-18.2)..(95.4,-18.1)
				..controls(95.7,-18)and(96.4,-17.8)..(96.9,-17.5)
				..controls(97.5,-17.3)and(98.4,-16.9)..(99,-16.6)
				..controls(99.7,-16.3)and(100.9,-15.9)..(101.9,-15.6)
				..controls(102.9,-15.3)and(104.4,-14.7)..(105.3,-14.3)
				..controls(108.1,-13.2)and(108.5,-13.1)..(110.1,-13.1)
				..controls(111.4,-13.1)and(111.5,-13.1)..(112.4,-13.5)
				..controls(113.2,-14)and(115,-14.8)..(116.6,-15.4)
				..controls(116.9,-15.5)and(117.9,-15.8)..(118.8,-16.1)
				..controls(120.9,-16.6)and(121.5,-16.9)..(122.6,-17.9)
				..controls(123,-18.3)and(123.7,-18.8)..(123.9,-18.9)
				..controls(124.6,-19.4)and(126,-19.8)..(127.9,-20.2)
				..controls(129.6,-20.5)and(129.8,-20.6)..(131.1,-21.4)
				..controls(131.5,-21.7)and(132.4,-22.1)..(133.1,-22.3)
				..controls(136.1,-23.2)and(138.3,-24.1)..(140.2,-25)
				..controls(141,-25.5)and(141.5,-25.6)..(142.5,-25.8)
				..controls(145.1,-26.3)and(146,-26.6)..(149.4,-28.1)
				..controls(151.3,-28.9)and(153.5,-29.8)..(154.3,-30.1)
				..controls(155.1,-30.4)and(156.2,-30.8)..(156.6,-31)
				..controls(157.1,-31.2)and(158.2,-31.6)..(159.1,-31.8)
				..controls(160.6,-32.2)and(160.9,-32.3)..(161.7,-32.9)
				..controls(162.4,-33.3)and(162.9,-33.6)..(163.7,-33.8)
				..controls(164.3,-34)and(164.8,-34.1)..(164.9,-34.1)
				..controls(165,-34.1)and(165.6,-34.4)..(166.4,-34.8)
				..controls(167.3,-35.2)and(168.1,-35.4)..(169,-35.6)
				..controls(170,-35.8)and(170.5,-36)..(171.7,-36.6)
				..controls(172.4,-37)and(173.3,-37.4)..(173.6,-37.4)
				..controls(175.6,-37.7)and(176.2,-38)..(176.7,-39.1)
				..controls(176.9,-39.6)and(176.9,-39.7)..(176.7,-40.4)
				..controls(176.4,-41.6)and(175.8,-42.7)..(175.2,-43.4)
				..controls(174.5,-44.1)and(172.6,-45.5)..(171.7,-45.8)
				..controls(170.6,-46.2)and(169.1,-46)..(168.2,-45.4)
				..controls(167.9,-45.2)and(167.3,-44.8)..(166.9,-44.4)
				..controls(166.4,-44.1)and(165.8,-43.6)..(165.4,-43.2)
				..controls(165.1,-42.9)and(164.4,-42.3)..(163.9,-41.9)
				..controls(163.3,-41.5)and(162.7,-40.9)..(162.4,-40.5)
				..controls(161.4,-39.2)and(159.9,-38)..(158.8,-37.8)
				..controls(158.4,-37.8)and(158.3,-37.7)..(158.2,-37.2)
				..controls(158,-36.8)and(157.4,-36.1)..(157.1,-36.1)
				..controls(157,-36.1)and(156.7,-36.1)..(156.5,-36)
				..controls(156.1,-35.9)and(155.9,-35.7)..(155.7,-35.2)
				..controls(155.3,-34.6)and(155.1,-34.4)..(154,-34.1)
				..controls(153.6,-34)and(153.2,-33.7)..(152.9,-33.5)
				..controls(152.4,-32.9)and(151.6,-32.4)..(151.3,-32.4)
				..controls(150.9,-32.4)and(150.4,-32.6)..(150.4,-32.9)
				..controls(150.4,-33.2)and(150.8,-33.4)..(151.4,-33.6)
				..controls(152.1,-33.7)and(152.8,-34.1)..(153.5,-34.9)
				..controls(154.1,-35.5)and(154.3,-35.6)..(155,-35.8)
				..controls(155.7,-36)and(155.7,-36)..(156,-36.6)
				..controls(156.3,-37.3)and(156.7,-37.7)..(157.5,-37.8)
				..controls(158.1,-37.9)and(158.2,-38)..(158.4,-38.8)
				..controls(158.5,-39.4)and(158.6,-39.6)..(159.6,-40.4)
				..controls(160.2,-41)and(160.9,-41.6)..(161.3,-41.8)
				..controls(161.7,-42)and(161.9,-42.3)..(162.2,-42.8)
				..controls(162.4,-43.2)and(162.8,-43.8)..(163.1,-44.1)
				..controls(163.9,-44.8)and(165.1,-46.7)..(165.7,-48)
				..controls(166.2,-49.3)and(166.4,-49.7)..(167.2,-50.5)
				..controls(167.5,-50.8)and(167.9,-51.4)..(168.1,-51.8)
				..controls(168.3,-52.2)and(168.8,-53)..(169.1,-53.6)
				..controls(169.6,-54.3)and(169.8,-54.9)..(169.9,-55.6)
				..controls(170.4,-57.8)and(170.5,-58.1)..(171.2,-59.5)
				..controls(171.6,-60.3)and(172,-61.5)..(172.2,-62.1)
				..controls(172.3,-62.8)and(172.7,-63.7)..(172.9,-64.1)
				..controls(173.2,-64.6)and(173.4,-65.3)..(173.4,-65.6)
				..controls(173.8,-68.8)and(174.2,-70.6)..(174.9,-72.6)
				..controls(175.5,-74.5)and(175.9,-77.1)..(176.1,-80.2)
				..controls(176.3,-83.4)and(175.5,-84.4)..(173.7,-83.4)--cycle;
				\definecolor{ezs}{RGB}{247,179,26}\fill[ezs](-5.7,-156.6)--(-3.4,-157.5)--(61.9,12.7)--(59.7,13.6)--cycle;
				\definecolor{ezs}{RGB}{247,179,22}\fill[ezs](-3.8,-157.3)--(-1.5,-158.2)--(63.9,12)--(61.6,12.9)--cycle;
				\definecolor{ezs}{RGB}{248,179,19}\fill[ezs](-1.8,-158.1)--(0.4,-158.9)--(65.8,11.2)--(63.5,12.1)--cycle;
				\definecolor{ezs}{RGB}{248,178,16}\fill[ezs](0.1,-158.8)--(2.4,-159.7)--(67.7,10.5)--(65.4,11.4)--cycle;
				\definecolor{ezs}{RGB}{249,178,15}\fill[ezs](2,-159.5)--(4.3,-160.4)--(69.6,9.8)--(67.3,10.7)--cycle;
				\definecolor{ezs}{RGB}{251,177,17}\fill[ezs](3.9,-160.3)--(6.2,-161.1)--(71.5,9)--(69.2,9.9)--cycle;
				\definecolor{ezs}{RGB}{252,179,18}\fill[ezs](5.8,-161)--(8.1,-161.9)--(73.4,8.3)--(71.1,9.2)--cycle;
				\definecolor{ezs}{RGB}{252,182,17}\fill[ezs](7.7,-161.7)--(10,-162.6)--(75.3,7.6)--(73,8.5)--cycle;
				\definecolor{ezs}{RGB}{252,185,17}\fill[ezs](9.6,-162.5)--(11.9,-163.4)--(77.2,6.8)--(75,7.7)--cycle;
				\definecolor{ezs}{RGB}{251,189,18}\fill[ezs](11.5,-163.2)--(13.8,-164.1)--(79.2,6.1)--(76.9,7)--cycle;
				\definecolor{ezs}{RGB}{251,191,22}\fill[ezs](13.4,-163.9)--(15.7,-164.8)--(81.1,5.4)--(78.8,6.2)--cycle;
				\definecolor{ezs}{RGB}{252,193,25}\fill[ezs](15.4,-164.7)--(17.6,-165.6)--(83,4.6)--(80.7,5.5)--cycle;
				\definecolor{ezs}{RGB}{253,195,27}\fill[ezs](17.3,-165.4)--(19.6,-166.3)--(84.9,3.9)--(82.6,4.8)--cycle;
				\definecolor{ezs}{RGB}{253,196,31}\fill[ezs](19.2,-166.1)--(21.5,-167)--(86.8,3.2)--(84.5,4)--cycle;
				\definecolor{ezs}{RGB}{254,199,35}\fill[ezs](21.1,-166.9)--(23.4,-167.8)--(88.7,2.4)--(86.4,3.3)--cycle;
				\definecolor{ezs}{RGB}{253,201,38}\fill[ezs](23,-167.6)--(25.3,-168.5)--(90.6,1.7)--(88.3,2.6)--cycle;
				\definecolor{ezs}{RGB}{253,203,40}\fill[ezs](24.9,-168.3)--(27.2,-169.2)--(92.5,1)--(90.3,1.8)--cycle;
				\definecolor{ezs}{RGB}{253,205,42}\fill[ezs](26.8,-169.1)--(29.1,-170)--(94.4,0.2)--(92.2,1.1)--cycle;
				\definecolor{ezs}{RGB}{253,207,45}\fill[ezs](28.7,-169.8)--(31,-170.7)--(96.4,-0.5)--(94.1,0.4)--cycle;
				\definecolor{ezs}{RGB}{253,210,49}\fill[ezs](30.7,-170.5)--(32.9,-171.4)--(98.3,-1.2)--(96,-0.4)--cycle;
				\definecolor{ezs}{RGB}{253,212,54}\fill[ezs](32.6,-171.3)--(34.9,-172.2)--(100.2,-2)--(97.9,-1.1)--cycle;
				\definecolor{ezs}{RGB}{253,214,59}\fill[ezs](34.5,-172)--(36.8,-172.9)--(102.1,-2.7)--(99.8,-1.8)--cycle;
				\definecolor{ezs}{RGB}{253,216,63}\fill[ezs](36.4,-172.7)--(38.7,-173.6)--(104,-3.4)--(101.7,-2.6)--cycle;
				\definecolor{ezs}{RGB}{253,218,65}\fill[ezs](38.3,-173.5)--(40.6,-174.4)--(105.9,-4.2)--(103.6,-3.3)--cycle;
				\definecolor{ezs}{RGB}{254,220,66}\fill[ezs](40.2,-174.2)--(42.5,-175.1)--(107.8,-4.9)--(105.5,-4)--cycle;
				\definecolor{ezs}{RGB}{254,222,68}\fill[ezs](42.1,-175)--(44.4,-175.8)--(109.7,-5.6)--(107.5,-4.8)--cycle;
				\definecolor{ezs}{RGB}{254,224,71}\fill[ezs](44,-175.7)--(46.3,-176.6)--(111.7,-6.4)--(109.4,-5.5)--cycle;
				\definecolor{ezs}{RGB}{254,226,75}\fill[ezs](46,-176.4)--(48.2,-177.3)--(113.6,-7.1)--(111.3,-6.2)--cycle;
				\definecolor{ezs}{RGB}{254,228,79}\fill[ezs](47.9,-177.2)--(50.1,-178)--(115.5,-7.8)--(113.2,-7)--cycle;
				\definecolor{ezs}{RGB}{254,230,82}\fill[ezs](49.8,-177.9)--(52.1,-178.8)--(117.4,-8.6)--(115.1,-7.7)--cycle;
				\definecolor{ezs}{RGB}{254,232,85}\fill[ezs](51.7,-178.6)--(54,-179.5)--(119.3,-9.3)--(117,-8.4)--cycle;
				\definecolor{ezs}{RGB}{254,234,88}\fill[ezs](53.6,-179.4)--(55.9,-180.2)--(121.2,-10)--(118.9,-9.2)--cycle;
				\definecolor{ezs}{RGB}{254,237,91}\fill[ezs](55.5,-180.1)--(57.8,-181)--(123.1,-10.8)--(120.8,-9.9)--cycle;
				\definecolor{ezs}{RGB}{254,239,94}\fill[ezs](57.4,-180.8)--(59.7,-181.7)--(125,-11.5)--(122.8,-10.6)--cycle;
				\definecolor{ezs}{RGB}{254,241,97}\fill[ezs](59.3,-181.6)--(61.6,-182.4)--(127,-12.2)--(124.7,-11.4)--cycle;
				\definecolor{ezs}{RGB}{254,243,101}\fill[ezs](61.2,-182.3)--(63.5,-183.2)--(128.9,-13)--(126.6,-12.1)--cycle;
				\definecolor{ezs}{RGB}{254,245,103}\fill[ezs](63.2,-183)--(65.4,-183.9)--(130.8,-13.7)--(128.5,-12.8)--cycle;
				\definecolor{ezs}{RGB}{254,245,105}\fill[ezs](65.1,-183.8)--(67.4,-184.6)--(132.7,-14.4)--(130.4,-13.6)--cycle;
				\definecolor{ezs}{RGB}{254,245,105}\fill[ezs](67,-184.5)--(69.3,-185.4)--(134.6,-15.2)--(132.3,-14.3)--cycle;
				\definecolor{ezs}{RGB}{254,244,103}\fill[ezs](68.9,-185.2)--(71.2,-186.1)--(136.5,-15.9)--(134.2,-15)--cycle;
				\definecolor{ezs}{RGB}{253,243,101}\fill[ezs](70.8,-186)--(73.1,-186.8)--(138.4,-16.6)--(136.1,-15.8)--cycle;
				\definecolor{ezs}{RGB}{253,241,98}\fill[ezs](72.7,-186.7)--(75,-187.6)--(140.3,-17.4)--(138,-16.5)--cycle;
				\definecolor{ezs}{RGB}{253,239,95}\fill[ezs](74.6,-187.4)--(76.9,-188.3)--(142.2,-18.1)--(140,-17.2)--cycle;
				\definecolor{ezs}{RGB}{253,236,92}\fill[ezs](76.5,-188.2)--(78.8,-189)--(144.2,-18.8)--(141.9,-18)--cycle;
				\definecolor{ezs}{RGB}{253,234,88}\fill[ezs](78.5,-188.9)--(80.7,-189.8)--(146.1,-19.6)--(143.8,-18.7)--cycle;
				\definecolor{ezs}{RGB}{253,232,85}\fill[ezs](80.4,-189.6)--(82.7,-190.5)--(148,-20.3)--(145.7,-19.4)--cycle;
				\definecolor{ezs}{RGB}{253,230,82}\fill[ezs](82.3,-190.4)--(84.6,-191.2)--(149.9,-21)--(147.6,-20.2)--cycle;
				\definecolor{ezs}{RGB}{253,228,78}\fill[ezs](84.2,-191.1)--(86.5,-192)--(151.8,-21.8)--(149.5,-20.9)--cycle;
				\definecolor{ezs}{RGB}{253,226,74}\fill[ezs](86.1,-191.8)--(88.4,-192.7)--(153.7,-22.5)--(151.4,-21.6)--cycle;
				\definecolor{ezs}{RGB}{253,224,70}\fill[ezs](88,-192.6)--(90.3,-193.4)--(155.6,-23.3)--(153.3,-22.4)--cycle;
				\definecolor{ezs}{RGB}{253,222,67}\fill[ezs](89.9,-193.3)--(92.2,-194.2)--(157.5,-24)--(155.3,-23.1)--cycle;
				\definecolor{ezs}{RGB}{253,220,65}\fill[ezs](91.8,-194)--(94.1,-194.9)--(159.5,-24.7)--(157.2,-23.8)--cycle;
				\definecolor{ezs}{RGB}{253,217,63}\fill[ezs](93.7,-194.8)--(96,-195.6)--(161.4,-25.5)--(159.1,-24.6)--cycle;
				\definecolor{ezs}{RGB}{253,215,60}\fill[ezs](95.7,-195.5)--(97.9,-196.4)--(163.3,-26.2)--(161,-25.3)--cycle;
				\definecolor{ezs}{RGB}{253,213,56}\fill[ezs](97.6,-196.2)--(99.9,-197.1)--(165.2,-26.9)--(162.9,-26)--cycle;
				\definecolor{ezs}{RGB}{252,211,53}\fill[ezs](99.5,-197)--(101.8,-197.8)--(167.1,-27.7)--(164.8,-26.8)--cycle;
				\definecolor{ezs}{RGB}{252,209,49}\fill[ezs](101.4,-197.7)--(103.7,-198.6)--(169,-28.4)--(166.7,-27.5)--cycle;
				\definecolor{ezs}{RGB}{252,207,45}\fill[ezs](103.3,-198.4)--(105.6,-199.3)--(170.9,-29.1)--(168.6,-28.2)--cycle;
				\definecolor{ezs}{RGB}{253,205,42}\fill[ezs](105.2,-199.2)--(107.5,-200)--(172.8,-29.9)--(170.6,-29)--cycle;
				\definecolor{ezs}{RGB}{253,202,40}\fill[ezs](107.1,-199.9)--(109.4,-200.8)--(174.7,-30.6)--(172.5,-29.7)--cycle;
				\definecolor{ezs}{RGB}{253,200,38}\fill[ezs](109,-200.6)--(111.3,-201.5)--(176.7,-31.3)--(174.4,-30.4)--cycle;
				\definecolor{ezs}{RGB}{253,198,35}\fill[ezs](111,-201.4)--(113.2,-202.2)--(178.6,-32.1)--(176.3,-31.2)--cycle;
				\definecolor{ezs}{RGB}{253,196,31}\fill[ezs](112.9,-202.1)--(115.2,-203)--(180.5,-32.8)--(178.2,-31.9)--cycle;
				\definecolor{ezs}{RGB}{253,194,28}\fill[ezs](114.8,-202.8)--(117.1,-203.7)--(182.4,-33.5)--(180.1,-32.6)--cycle;
				\definecolor{ezs}{RGB}{253,192,26}\fill[ezs](116.7,-203.6)--(119,-204.5)--(184.3,-34.3)--(182,-33.4)--cycle;
				\definecolor{ezs}{RGB}{253,189,24}\fill[ezs](118.6,-204.3)--(120.9,-205.2)--(186.2,-35)--(183.9,-34.1)--cycle;
				\definecolor{ezs}{RGB}{253,186,20}\fill[ezs](120.5,-205)--(122.8,-205.9)--(188.1,-35.7)--(185.8,-34.8)--cycle;
				\definecolor{ezs}{RGB}{253,184,17}\fill[ezs](122.4,-205.8)--(124.7,-206.7)--(190,-36.5)--(187.8,-35.6)--cycle;
				\definecolor{ezs}{RGB}{252,183,15}\fill[ezs](124.3,-206.5)--(126.6,-207.4)--(192,-37.2)--(189.7,-36.3)--cycle;
				\definecolor{ezs}{RGB}{252,181,14}\fill[ezs](126.3,-207.2)--(128.5,-208.1)--(193.9,-37.9)--(191.6,-37.1)--cycle;
				\definecolor{ezs}{RGB}{252,179,13}\fill[ezs](128.2,-208)--(130.4,-208.9)--(195.8,-38.7)--(193.5,-37.8)--cycle;
				\definecolor{ezs}{RGB}{252,178,12}\fill[ezs](130.1,-208.7)--(132.4,-209.6)--(197.7,-39.4)--(195.4,-38.5)--cycle;
				\definecolor{ezs}{RGB}{251,178,10}\fill[ezs](132,-209.4)--(134.3,-210.3)--(199.6,-40.1)--(197.3,-39.3)--cycle;
				\definecolor{ezs}{RGB}{251,178,10}\fill[ezs](133.9,-210.2)--(136.2,-211.1)--(201.5,-40.9)--(199.2,-40)--cycle;
				\definecolor{ezs}{RGB}{250,177,11}\fill[ezs](135.8,-210.9)--(138.1,-211.8)--(203.4,-41.6)--(201.1,-40.7)--cycle;
				\definecolor{ezs}{RGB}{250,178,11}\fill[ezs](137.7,-211.6)--(140,-212.5)--(205.3,-42.3)--(203.1,-41.5)--cycle;
				\definecolor{ezs}{RGB}{248,180,10}\fill[ezs](139.6,-212.4)--(141.9,-213.3)--(207.3,-43.1)--(205,-42.2)--cycle;
				\definecolor{ezs}{RGB}{247,181,9}\fill[ezs](141.5,-213.1)--(143.8,-214)--(209.2,-43.8)--(206.9,-42.9)--cycle;
				\definecolor{ezs}{RGB}{246,182,9}\fill[ezs](143.5,-213.8)--(145.7,-214.7)--(211.1,-44.5)--(208.8,-43.7)--cycle;
				\definecolor{ezs}{RGB}{246,182,10}\fill[ezs](145.4,-214.6)--(147.7,-215.5)--(213,-45.3)--(210.7,-44.4)--cycle;
			\end{scope}
			% chữ EDUZY (Times đậm, làm dày nét bằng 5 lớp lệch nhẹ; căn đúng khung chữ logo gốc)
			\foreach \dx/\dy in {-0.7/0,0.7/0,0/-0.3,0/0.3,0/0}
			\node[anchor=center,inner sep=0pt,outer sep=0pt,text=ezText,align=center] at (109.6+\dx,-193.6+\dy)
			{\resizebox{1.01cm}{0.22cm}{\fontencoding{T1}\fontfamily{ptm}\fontseries{b}\selectfont EDUZY}};
		\end{tikzpicture}%
	}%
}
%====================== THIẾT LẬP THÂN ĐỀ ======================
%--- Môi trường tự luận {bt}: ex_test 3.3.4 chưa khai báo sẵn nên tạo ở đây
%    (ex_test đã tự móc \fixshowans và \listenumerate cho bt)
\makeatletter
\@ifundefined{bt}{\newtheorem{bt}{\color{gold1}\namebt}}{}
\makeatother
\renewcommand{\sectionmark}[1]{\markright{\thesection\ #1}}
%--- Khung tiêu đề (logo + tối đa 4 dòng; dòng 3, 4 để trống {} thì bỏ qua)
%  Đề thi:    \khungtieude{THI THỬ THPT QUỐC GIA NĂM 2026}{Môn Toán -- Đề số 21}
%                         {(Đề thi gồm: ...)}{Thời gian làm bài: 90 phút (không kể thời gian phát đề).}
%  Chuyên đề: \khungtieude{CHUYÊN ĐỀ 3 -- DÃY SỐ VÀ CÔNG THỨC TRUY HỒI}
%                         {Bồi dưỡng học sinh giỏi Toán THPT}{...}{}
\newcommand{\khungtieude}[4]{%
	\begin{tcolorbox}[colframe=gold1,colback=gold5!45,boxrule=0.5mm]
		\begin{minipage}[c]{3cm}
			\centerline{\eduzylogo[2.4cm]}
		\end{minipage}\hspace{0.5cm}
		\begin{minipage}[c]{12.5cm}
			\vskip 0.15 cm
			\centerline{\ind{#1}}
			\vskip 0.15 cm
			\centerline{\textbf{#2}}
			\if\relax\detokenize{#3}\relax\else
				\vskip 0.15 cm
				\centerline{\textit{#3}}
			\fi
			\if\relax\detokenize{#4}\relax\else
				\vskip 0.15 cm
				\centerline{\textit{#4}}
			\fi
			\vskip 0.15 cm
		\end{minipage}
	\end{tcolorbox}
}
%====================== BỔ SUNG CHO BÀI HỌC, CHUYÊN ĐỀ ======================
\usepackage[shortlabels]{enumitem}
\usepackage{tabularx}
%--- Danh sách gọn: \begin{enumerate}[gon] ... \end{enumerate}
\SetEnumitemKey{gon}{leftmargin=*,itemsep=1pt,parsep=0pt,topsep=3pt}
%--- Khung lý thuyết: \begin{khung}{} ... \end{khung}; tiêu đề chỉ dùng khi thật ngắn ("Mẹo", "Lỗi thường gặp")
\newtcolorbox{khung}[1]{breakable,enhanced,colframe=gold1,colback=gold5!45,
	boxrule=0.4mm,fonttitle=\bfseries\sffamily,coltitle=white,colbacktitle=gold1,
	attach boxed title to top left={xshift=6mm,yshift=-3mm},
	boxed title style={arc=1mm,boxrule=0.3mm},title={#1}}
%--- Tiêu đề mục nhỏ trong một dạng: \muc{Dấu hiệu}, \muc{Phương pháp}, \muc{Bài tự luyện}
\makeatletter
\newcommand{\muc}[1]{\par\medskip\noindent\indam{#1}\par\nobreak\smallskip\@afterindentfalse\@afterheading}
\makeatother
%--- Ô công thức cần nhớ (căn giữa, viền vàng, vừa khít nội dung)
%     \begin{center}\congthuc{$Q(n)=(n-1)!$}\end{center}
\newtcbox{\congthuc}{on line,colframe=gold1,colback=gold5!45,boxrule=0.5pt,arc=2pt,
	left=8pt,right=8pt,top=4pt,bottom=4pt}
%--- Công tắc lời giải bài tự luyện (ví dụ mẫu luôn dùng \loigiai)
%     \loigiaiTL[{đáp số}]{lời giải đầy đủ}
%     \TLfull -> in lời giải đầy đủ;  \TLdapso -> chỉ in một dòng "Đáp số";  \TLan -> ẩn hoàn toàn
%     (\CheDo là dethi thì mọi lời giải và đáp số đều ẩn)
\makeatletter
\newif\ifTLhien   \TLhientrue
\newif\ifTLdapso  \TLdapsofalse
\newcommand{\TLfull}{\TLhientrue\TLdapsofalse}
\newcommand{\TLdapso}{\TLhientrue\TLdapsotrue}
\newcommand{\TLan}{\TLhienfalse}
\newcommand{\loigiaiTL}[2][]{%
	\ifTLhien
		\ifTLdapso
			\def\@tlarg{#1}%
			\ifx\@tlarg\@empty\else
				\ifdefstring{\kindOPTN}{dethi}{}{\par\noindent{\bfseries\color{gold1}Đáp số:} #1\par}%
			\fi
		\else
			\loigiai{#2}%
		\fi
	\fi}
\makeatother
\TLdapso% mặc định: chỉ in đáp số bài tự luyện
%====================== BẢNG ĐÚNG/SAI (\choiceTFt) ======================
%  Hàng tiêu đề nền nâu đồng chữ trắng (Phát biểu | Đ | S), các hàng tô xen kẽ kem / trắng,
%  đường kẻ vàng kim, nhãn a) b) c) d) màu nâu đồng. Bản có đáp án: cột được chọn có vòng tròn
%  nâu đồng với dấu tích trắng (câu đúng tích ở cột Đúng, câu sai tích ở cột Sai).
%  Tiêu đề ghi đủ "Đúng | Sai": \OPTN{viettat=0} sau \input (cột Đúng/Sai rộng ra, phát biểu hẹp lại).
%  Dùng bảng gốc của ex_test: \def\BangDSgoc{1} trước \input{preamble-eduzy-v2}.
\providecommand{\BangDSgoc}{0}
\ifnum\BangDSgoc=0
\usepackage{colortbl,pifont}
\colorlet{ezDSke}{gold2!60!gold1}% đường kẻ bảng (nét 0,4pt như ex_test để không đổi chiều cao bảng)
\colorlet{ezDSle}{gold5!60}% hàng a), c)
\colorlet{ezDSchan}{white}% hàng b), d)
%--- dấu chọn ở cột Đúng/Sai
%    (bề rộng 0, căn giữa ô: không làm đổi bề rộng cột của ex_test, không báo tràn lề)
\newcommand{\EZdauchon}{\makebox[0pt]{\tikz[baseline=-1ex]\node[circle,fill=gold1,inner sep=0pt,
	minimum size=1.2em,text=white,font=\small]{\ding{51}};}}
\def\tickT{\EZdauchon}\def\tickF{\EZdauchon}
\AtBeginDocument{\sbox{\ccTFTab}{\EZdauchon}}
%--- nhãn a) b) c) d) màu nâu đồng (ở bản lời giải, nhãn câu đúng vẫn đóng khung như ex_test)
\renewcommand{\FalseTF}{\stepcounter{dapan}\textbf{\color{gold1}\DapAnTF\sepTF}\ignorespaces}
%--- trước khi dựng bảng: kẻ màu vàng kim (bề rộng cột giữ nguyên như ex_test); sau bảng trả lại màu đen
\newcount\EZhangTF
\newcommand{\EZbatdauTF}{%
	\global\EZhangTF=0\relax
	\arrayrulecolor{ezDSke}}
\newcommand{\EZketthucTF}{\arrayrulecolor{black}}
%--- hàng tiêu đề nền nâu đồng, chữ trắng
\makeatletter
\newcommand{\EZo@TF}[1]{\centering\bfseries\color{white}#1}
\newcommand{\EZtieudeTF}{%
	\ifcase\cotDS
		\ifnum\cotSTT=1\relax
			\def\rowHeader{\hline\rowcolor{gold1}\EZo@TF{TT}&\EZo@TF{\phatbieu}\tabularnewline}%
		\else
			\def\rowHeader{\hline\rowcolor{gold1}\EZo@TF{\phatbieu}\tabularnewline}%
		\fi
	\or
		\ifnum\cotSTT=1\relax
			\def\rowHeader{\hline\rowcolor{gold1}\EZo@TF{TT}&\EZo@TF{\phatbieu}&\EZo@TF{\Dung}\tabularnewline}%
		\else
			\def\rowHeader{\hline\rowcolor{gold1}\EZo@TF{\phatbieu}&\EZo@TF{\Dung}\tabularnewline}%
		\fi
	\or
		\ifnum\cotSTT=1\relax
			\def\rowHeader{\hline\rowcolor{gold1}\EZo@TF{TT}&\EZo@TF{\phatbieu}&\EZo@TF{\Dung}&\EZo@TF{\Sai}\tabularnewline}%
		\else
			\def\rowHeader{\hline\rowcolor{gold1}\EZo@TF{\phatbieu}&\EZo@TF{\Dung}&\EZo@TF{\Sai}\tabularnewline}%
		\fi
	\fi}
%--- tô xen kẽ từng hàng (gọi ngay sau \hline đầu mỗi hàng)
\newcommand{\EZmauhangTF}{%
	\noalign{\global\advance\EZhangTF\@ne}%
	\ifodd\EZhangTF\rowcolor{ezDSle}\else\rowcolor{ezDSchan}\fi}
%--- gắn vào \choiceTFt / \writeChoiceTFt của ex_test
\xpretocmd{\choiceTFt}{\EZbatdauTF}{}{\typeout{[preamble-eduzy-v2] Không gắn được bảng Đúng/Sai (đầu bảng)}}
\xpatchcmd{\choiceTFt}{\ifnum\dongPB=0\def\rowHeader{}\fi}{\EZtieudeTF\ifnum\dongPB=0\def\rowHeader{}\fi}
	{}{\typeout{[preamble-eduzy-v2] Không gắn được bảng Đúng/Sai (hàng tiêu đề)}}
\xpatchcmd{\writeChoiceTFt}{\hline\xdef\dotEX{\sdotEX}}{\hline\EZmauhangTF\xdef\dotEX{\sdotEX}}
	{}{\typeout{[preamble-eduzy-v2] Không gắn được bảng Đúng/Sai (tô hàng)}}
\xpatchcmd{\writeChoiceTFt}{{\end{longtable}}}{{\end{longtable}\EZketthucTF}}
	{}{\typeout{[preamble-eduzy-v2] Không gắn được bảng Đúng/Sai (cuối longtable)}}
\xpatchcmd{\writeChoiceTFt}{{\end{xtabular}\end{center}\medskip\noindent}}{{\end{xtabular}\end{center}\EZketthucTF\medskip\noindent}}
	{}{\typeout{[preamble-eduzy-v2] Không gắn được bảng Đúng/Sai (cuối xtabular)}}
\fi
%=====================================================================
\renewcommand{\parallel}{\mathrel{/\!\!/}}

%  trong file đề (tools/dan_trang/dan_trang.py tự ghi khi đặt \NeedspaceCau):
%    \def\TachCauDS{1}   câu Đúng/Sai được tách: câu dẫn (kèm hình) ở cuối trang, bảng Phát biểu sang trang sau
%                        (câu dẫn và hình không tách nhau, bảng không tách giữa hai trang)
%    \def\TachCauLG{1}   bản lời giải (\CheDo khác dethi): phần đề của câu cao hơn 1/3 trang được phép tách
%  Lệnh dùng trong file đề (giá trị do tools/dan_trang/dan_trang.py đo ở bản dethi):
%    \NeedspaceCau[<cao câu dẫn + hình>]{<cao cả khối câu>}   ngay trên \begin{ex} (câu đầu phần: trên tiêu đề phần)
%        công tắc tắt: đúng bằng \Needspace{<cao cả khối câu>} như cũ;
%        \TachCauDS=1 và có [..] (câu Đúng/Sai được chọn tách): chỉ giữ câu dẫn + hình;
%        \TachCauLG=1 ở bản lời giải: đòi tối đa 1/3 \textheight.
%    \NeedspaceBang{<cao bảng Phát biểu>}%   ngay trước \choiceTFt: \TachCauDS=1 thì giữ trọn bảng trên một trang;
%        công tắc tắt: không làm gì.
\providecommand{\TachCauDS}{0}
\providecommand{\TachCauLG}{0}
\newlength{\EZnsCau}
\newcommand{\NeedspaceCau}[2][]{%
	\setlength{\EZnsCau}{#2}%
	\ifnum\TachCauDS=1 \if\relax\detokenize{#1}\relax\else\setlength{\EZnsCau}{#1}\fi\fi
	\ifnum\TachCauLG=1 \ifdefstring{\kindOPTN}{dethi}{}{\ifdim\EZnsCau>0.3333\textheight\setlength{\EZnsCau}{0.3333\textheight}\fi}\fi
	\Needspace{\EZnsCau}}
\newcommand{\NeedspaceBang}[1]{\ifnum\TachCauDS=1 \Needspace{#1}\fi}

%--- Thang màu: đậm -> nhạt
\definecolor{gold1}{RGB}{126,96,0}    % nâu đồng đậm  #7E6000
\definecolor{gold2}{RGB}{192,144,3}   % vàng kim      #C09003
\definecolor{gold3}{RGB}{255,215,99}  % vàng nhạt     #FFD763
\definecolor{gold4}{RGB}{255,231,153} % vàng rất nhạt #FFE799
\definecolor{gold5}{RGB}{255,243,205} % vàng kem      #FFF3CD
%--- Ghi đè tông xanh/đỏ cũ
\definecolor{myblue}{RGB}{126,96,0}
\colorlet{blue}{gold1}
\colorlet{cyan}{gold2}
\colorlet{teal}{gold2}

%--- Tên chương mặc định (file đề nên tự khai báo lại)
\providecommand{\tenchuong}{}
%--- Chân trang
\providecommand{\thuonghieu}{TOÁN 12 -- Eduzy.com.vn}
\providecommand{\tacgia}{Tô Gia Bảo -- 0353703783}

%--- Chữ cạnh logo ở đầu trang
\providecommand{\khauhieu}{Eduzy.com.vn}

%====================== THÔNG SỐ ĐẦU TRANG, CHÂN TRANG (v2) ======================
%  Dải nền đậm tràn ngang trang; ô chữ vàng kem có cạnh xiên, kèm 2 vạch xiên chuyển màu.
%  Đầu trang lẻ : logo + \khauhieu (trái)   | ô chữ \tenchuong (phải)
%  Chân trang lẻ: ô chữ \thuonghieu (trái)  | Trang (số) (phải)
%  Trang chẵn soi gương trang lẻ; chân trang chẵn ghi \tacgia.
%--- Màu: chỉ dùng thang vàng kim Eduzy
\colorlet{ezNen}{gold1}% dải nền đậm
\colorlet{ezVach1}{gold3}% vạch xiên sát dải nền (đậm hơn)
\colorlet{ezVach2}{gold4}% vạch xiên sát ô chữ (nhạt hơn)
\colorlet{ezO}{gold5}% ô chữ vàng kem
%--- Kích thước
\newlength{\EZdai}\setlength{\EZdai}{1.1cm}% chiều cao dải nền
\newlength{\EZo}\setlength{\EZo}{0.68cm}% chiều cao ô chữ
\newlength{\EZxien}\setlength{\EZxien}{0.55cm}% độ lệch ngang của cạnh xiên
\newlength{\EZvach}\setlength{\EZvach}{0.75cm}% bề rộng mỗi vạch xiên
\newlength{\EZkhe}\setlength{\EZkhe}{0.12cm}% khe trắng giữa vạch và ô chữ
\newlength{\EZdem}\setlength{\EZdem}{0.4cm}% khoảng đệm từ chữ tới cạnh xiên
\newlength{\EZle}\setlength{\EZle}{1.8cm}% lề trái/phải của chữ (bằng lề thân bài)
\newlength{\EZrongmin}\setlength{\EZrongmin}{4.5cm}% ô chữ rộng tối thiểu
\newlength{\EZrong}% bề rộng chữ trong ô (tính tự động)
%--- Font
\newcommand{\EZfontO}{\fontsize{10.5pt}{12pt}\selectfont\bfseries}% chữ trong ô
\newcommand{\EZfontKH}{\fontfamily{qag}\fontsize{11pt}{12pt}\selectfont\bfseries}% \khauhieu cạnh logo
\newcommand{\EZfontTrang}{\fontsize{11.5pt}{12pt}\selectfont\bfseries}% chữ "Trang" và số trang
%--- Đo bề rộng chữ trong ô
\newcommand{\EZdorong}[1]{%
	\settowidth{\EZrong}{\EZfontO #1}%
	\ifdim\EZrong<\EZrongmin\setlength{\EZrong}{\EZrongmin}\fi}
%--- Logo + chữ trắng \khauhieu (dựng sẵn một lần khi vào \begin{document})
\newsavebox{\EZlogobox}
\AtBeginDocument{\sbox{\EZlogobox}{%
		\raisebox{-0.5\height}{\eduzylogo[1cm]}\hspace{0.15cm}%
		\raisebox{-0.5\height}{\color{white}\EZfontKH\khauhieu}}}
%--- Hai vạch xiên + khe trắng. #1: điểm gốc ở mép dưới, #2: hướng (1 sang phải, -1 sang trái),
%    #3: hướng nghiêng (1 kiểu "/", -1 kiểu "\")
\newcommand{\EZvachxien}[3]{%
	\fill[white] (#1) --++({(#3)*\EZxien},\EZo) --++({(#2)*\EZkhe},0) --++({-(#3)*\EZxien},-\EZo) --cycle;
	\fill[ezVach2] ([xshift=(#2)*\EZkhe]#1) --++({(#3)*\EZxien},\EZo) --++({(#2)*\EZvach},0) --++({-(#3)*\EZxien},-\EZo) --cycle;
	\fill[ezVach1] ([xshift=(#2)*(\EZkhe+\EZvach)]#1) --++({(#3)*\EZxien},\EZo) --++({(#2)*\EZvach},0) --++({-(#3)*\EZxien},-\EZo) --cycle;
	\draw[white,line width=0.6pt] ([xshift=(#2)*(\EZkhe+\EZvach)]#1) --++({(#3)*\EZxien},\EZo);
}
%--- Chữ trong ô (căn giữa theo chiều cao ô, không lệch vì dấu tiếng Việt)
\tikzset{ezchu/.style={inner sep=0pt,text=ezNen,text height=0.72em,text depth=0.22em}}

%====================== LOGO CHÌM GIỮA TRANG ======================
%  Vẽ cùng đầu trang nên nằm dưới chữ của thân bài, có trên mọi trang.
%  Tắt: \logochimtat (từ trang đó trở đi), bật lại: \logochimbat.
%  Đổi cỡ, độ đậm: khai báo trước \input, ví dụ \def\EZchimcao{9cm} \def\EZchimdo{0.1}
\providecommand{\EZchimcao}{10cm}% chiều cao logo chìm
\providecommand{\EZchimdo}{0.12}% độ đậm (0 = trong suốt, 1 = rõ hẳn)
\newif\ifEZchim \EZchimtrue
\newcommand{\logochimtat}{\EZchimfalse}
\newcommand{\logochimbat}{\EZchimtrue}
%  Dựng sẵn logo mờ một lần khi vào \begin{document}; transparency group làm mờ cả logo
%  một lượt nên không lộ vệt ở chỗ các dải màu chồng nhau.
\newsavebox{\EZchimtho}
\newsavebox{\EZchimbox}
\AtBeginDocument{%
	\sbox{\EZchimtho}{\eduzylogo[\EZchimcao]}%
	\sbox{\EZchimbox}{\begin{tikzpicture}
			\begin{scope}[transparency group,opacity=\EZchimdo]
				\node[inner sep=0pt] {\usebox{\EZchimtho}};
			\end{scope}
	\end{tikzpicture}}}
\newcommand{\EZvechim}{%
	\ifEZchim
		\node[inner sep=0pt] at (current page.center) {\usebox{\EZchimbox}};
	\fi}

%====================== ĐẦU TRANG - TRANG LẺ ======================
%  (logo bên trái, ô chữ \tenchuong bên phải)
\newcommand{\myfancyheadOdd}{%
	\EZdorong{\tenchuong}%
	\begin{tikzpicture}[remember picture,overlay]
		\coordinate (T) at (current page.north west);
		\coordinate (P) at (current page.north east);
		%--- logo chìm giữa trang
		\EZvechim
		%--- dải nền
		\fill[ezNen] (T) rectangle ([yshift=-\EZdai]P);
		%--- ô chữ: X là góc dưới bên trái, cạnh trái xiên kiểu "/"
		\coordinate (X) at ([xshift=-\EZle-\EZrong-\EZdem-\EZxien,yshift=-\EZdai]P);
		\fill[ezO] (X) --++(\EZxien,\EZo) -- ([yshift=-\EZdai+\EZo]P) -- ([yshift=-\EZdai]P) --cycle;
		\EZvachxien{X}{-1}{1}
		\node[ezchu,anchor=east] at ([xshift=-\EZle,yshift=-\EZdai+0.5*\EZo]P) {\EZfontO\tenchuong};
		%--- logo + tên
		\node[inner sep=0pt,anchor=west] at ([xshift=\EZle,yshift=-0.5*\EZdai]T) {\usebox{\EZlogobox}};
	\end{tikzpicture}%
}

%====================== ĐẦU TRANG - TRANG CHẴN =====================
%  (soi gương: ô chữ \tenchuong bên trái, logo bên phải)
\newcommand{\myfancyheadEven}{%
	\EZdorong{\tenchuong}%
	\begin{tikzpicture}[remember picture,overlay]
		\coordinate (T) at (current page.north west);
		\coordinate (P) at (current page.north east);
		%--- logo chìm giữa trang
		\EZvechim
		%--- dải nền
		\fill[ezNen] (T) rectangle ([yshift=-\EZdai]P);
		%--- ô chữ: X là góc dưới bên phải, cạnh phải xiên kiểu "\"
		\coordinate (X) at ([xshift=\EZle+\EZrong+\EZdem+\EZxien,yshift=-\EZdai]T);
		\fill[ezO] (X) --++(-\EZxien,\EZo) -- ([yshift=-\EZdai+\EZo]T) -- ([yshift=-\EZdai]T) --cycle;
		\EZvachxien{X}{1}{-1}
		\node[ezchu,anchor=west] at ([xshift=\EZle,yshift=-\EZdai+0.5*\EZo]T) {\EZfontO\tenchuong};
		%--- logo + tên
		\node[inner sep=0pt,anchor=east] at ([xshift=-\EZle,yshift=-0.5*\EZdai]P) {\usebox{\EZlogobox}};
	\end{tikzpicture}%
}

%====================== CHÂN TRANG - TRANG LẺ ======================
%  (ô chữ \thuonghieu bên trái, "Trang (số)" bên phải)
\newcommand{\myfancyfootOdd}{%
	\EZdorong{\thuonghieu}%
	\begin{tikzpicture}[remember picture,overlay]
		\coordinate (T) at (current page.south west);
		\coordinate (P) at (current page.south east);
		%--- dải nền
		\fill[ezNen] (T) rectangle ([yshift=\EZdai]P);
		%--- ô chữ sát mép trên dải: X là góc dưới bên phải, cạnh phải xiên kiểu "/"
		\coordinate (X) at ([xshift=\EZle+\EZrong+\EZdem,yshift=\EZdai-\EZo]T);
		\fill[ezO] ([yshift=\EZdai-\EZo]T) -- (X) --++(\EZxien,\EZo) -- ([yshift=\EZdai]T) --cycle;
		\EZvachxien{X}{1}{1}
		\node[ezchu,anchor=west] at ([xshift=\EZle,yshift=\EZdai-0.5*\EZo]T) {\EZfontO\thuonghieu};
		%--- số trang
		\node[circle,draw=white,line width=0.9pt,text=white,inner sep=0pt,minimum size=0.62cm,anchor=east]
		(so) at ([xshift=-\EZle,yshift=0.5*\EZdai]P) {\EZfontTrang\thepage};
		\node[inner sep=0pt,text=white,anchor=base east] (tr) at ([xshift=-0.15cm]so.base west) {\EZfontTrang Trang};
	\end{tikzpicture}%
}

%====================== CHÂN TRANG - TRANG CHẴN ====================
%  (soi gương: "Trang (số)" bên trái, ô chữ \tacgia bên phải)
\newcommand{\myfancyfootEven}{%
	\EZdorong{\tacgia}%
	\begin{tikzpicture}[remember picture,overlay]
		\coordinate (T) at (current page.south west);
		\coordinate (P) at (current page.south east);
		%--- dải nền
		\fill[ezNen] (T) rectangle ([yshift=\EZdai]P);
		%--- ô chữ sát mép trên dải: X là góc dưới bên trái, cạnh trái xiên kiểu "\"
		\coordinate (X) at ([xshift=-\EZle-\EZrong-\EZdem,yshift=\EZdai-\EZo]P);
		\fill[ezO] ([yshift=\EZdai-\EZo]P) -- (X) --++(-\EZxien,\EZo) -- ([yshift=\EZdai]P) --cycle;
		\EZvachxien{X}{-1}{-1}
		\node[ezchu,anchor=east] at ([xshift=-\EZle,yshift=\EZdai-0.5*\EZo]P) {\EZfontO\tacgia};
		%--- số trang
		%    (vòng tròn ẩn so0 nằm giữa dải để lấy đường chân chữ, giống trang lẻ)
		\node[circle,inner sep=0pt,minimum size=0.62cm] (so0) at ([yshift=0.5*\EZdai]T) {\EZfontTrang\phantom{\thepage}};
		\node[inner sep=0pt,text=white,anchor=base west] (tr) at ($(so0.base -| T)+(\EZle,0)$) {\EZfontTrang Trang};
		\node[circle,draw=white,line width=0.9pt,text=white,inner sep=0pt,minimum size=0.62cm,anchor=base west]
		(so) at ([xshift=0.15cm]tr.base east) {\EZfontTrang\thepage};
	\end{tikzpicture}%
}

%--- Giữ tương thích với các file cũ đang gọi \myfancyhead, \myfancyfoot
\newcommand{\myfancyhead}{\myfancyheadOdd}
\newcommand{\myfancyfoot}{\myfancyfootOdd}

%======================= GẮN VÀO TRANG =============================
\pagestyle{fancy}
\fancyhf{}
\fancyhead[LO]{\myfancyheadOdd}
\fancyhead[RE]{\myfancyheadEven}
\fancyfoot[LO]{\myfancyfootOdd}
\fancyfoot[RE]{\myfancyfootEven}
\renewcommand{\headrulewidth}{0pt}
\renewcommand{\footrulewidth}{0pt}
\fancypagestyle{plain}{%
	\fancyhf{}
	\fancyhead[LO]{\myfancyheadOdd}
	\fancyhead[RE]{\myfancyheadEven}
	\fancyfoot[LO]{\myfancyfootOdd}
	\fancyfoot[RE]{\myfancyfootEven}
	\renewcommand{\headrulewidth}{0pt}
	\renewcommand{\footrulewidth}{0pt}
}
%=================== HẾT PHẦN GỘP TỪ eduzy-de.tex ===================
\newcommand{\indam}[1]{
	{\fontsize{12pt}{9pt}\sffamily \selectfont\bfseries\color{gold1} #1}
}
\newcommand{\ind}[1]{
	{\fontsize{14pt}{9pt}\fontfamily{qag} \selectfont\bfseries\color{gold1} #1}
}
\renewcommand{\FalseEX}{\stepcounter{dapan}\textbf{\selectfont\color{gold1}\Alph{dapan}.}\ignorespaces}
\renewtheorem{ex}{\color{gold1}Câu}
\renewcommand{\baselinestretch}{1.1}
\newcommand{\hoac}[1]{
	\left[\begin{aligned}#1\end{aligned}\right.}
\newcommand{\heva}[1]{
	\left\{\begin{aligned}#1\end{aligned}\right.}
\usepackage{esvect}
\def\vec{\vv}
\def\vt{\vv}
\def\overrightarrow{\vv}
%====================== LOGO EDUZY VẼ BẰNG TIKZ ======================
%  Đường viền dò từ logo gốc; dải màu vàng kim đo trực tiếp trên ảnh gốc
%  (nghiêng 21 độ, tô bằng 80 dải hẹp nên hiển thị giống nhau trên mọi trình xem PDF).
%  Dùng: \eduzylogo          -> cao 2.4cm (mặc định)
%        \eduzylogo[1.5cm]   -> cao 1.5cm
\definecolor{ezText}{RGB}{245,200,46}
\newcommand{\eduzylogo}[1][2.4cm]{%
	\resizebox{!}{#1}{%
		\begin{tikzpicture}[x=0.01cm,y=0.01cm]
			\path[use as bounding box] (0,0) rectangle (224,-224);
			% mũ tốt nghiệp + vương miện
			\begin{scope}
				\clip
				(119.4,-176.1)..controls(118.9,-176.1)and(116.8,-176)..(114.6,-176)
				..controls(101.4,-176)and(85.8,-175.9)..(83.7,-175.7)
				..controls(81.8,-175.6)and(81,-175.6)..(80.2,-175.7)
				..controls(78.8,-175.9)and(77.7,-175.9)..(76.2,-175.7)
				..controls(74.9,-175.5)and(70.8,-175.5)..(64.3,-175.7)
				..controls(59.1,-175.9)and(58.3,-175.7)..(58.6,-173.7)
				..controls(58.6,-172.8)and(59.6,-169.6)..(60.3,-168)
				..controls(60.8,-166.9)and(61.4,-165.8)..(62.3,-165)
				..controls(63.1,-164.2)and(64.3,-163.4)..(65,-163.2)
				..controls(65.5,-163.1)and(72.3,-163)..(76.6,-163.1)
				..controls(77.9,-163.1)and(80.1,-163.1)..(81.4,-163)
				..controls(82.8,-162.9)and(86.3,-162.9)..(89.2,-162.9)
				..controls(93,-162.9)and(94.9,-162.8)..(96,-162.7)
				..controls(97.2,-162.5)and(101,-162.5)..(114,-162.4)
				..controls(123,-162.3)and(134.9,-162.2)..(140.2,-162.1)
				..controls(150.1,-162)and(152.2,-162.1)..(153.6,-162.4)
				..controls(154.4,-162.6)and(157.1,-163.8)..(158.1,-164.4)
				..controls(158.8,-164.8)and(159.7,-165.7)..(159.9,-166.3)
				..controls(160,-166.5)and(160.2,-167.1)..(160.2,-167.6)
				..controls(160.3,-168)and(160.6,-169.4)..(160.9,-170.6)
				..controls(161.2,-171.7)and(161.5,-173.1)..(161.6,-173.6)
				..controls(161.8,-174.4)and(161.8,-174.5)..(161.5,-174.9)
				..controls(161.2,-175.6)and(160.8,-175.9)..(159.9,-176)
				..controls(159.2,-176.1)and(154.1,-176.2)..(151.5,-176.2)
				..controls(148.1,-176.1)and(122.6,-176.1)..(121.6,-176.2)
				..controls(120.8,-176.2)and(119.8,-176.2)..(119.4,-176.1)--cycle
				(130.4,-160)..controls(129,-160)and(119.8,-160)..(110.1,-159.9)
				..controls(100.3,-159.8)and(88.5,-159.7)..(83.9,-159.6)
				..controls(77.8,-159.5)and(75.3,-159.5)..(74.5,-159.6)
				..controls(73.9,-159.7)and(72.5,-159.8)..(71.3,-159.8)
				..controls(69.6,-159.8)and(69,-159.8)..(68.5,-159.6)
				..controls(67.6,-159.3)and(66.8,-158.6)..(66.2,-157.6)
				..controls(66,-157.3)and(65.6,-156.5)..(65.3,-156.1)
				..controls(64.6,-155.1)and(64.5,-154.9)..(63.7,-152.9)
				..controls(63.2,-151.6)and(62.9,-150.9)..(61.8,-149.6)
				..controls(61.5,-149.2)and(61.1,-148.5)..(60.9,-147.9)
				..controls(60.7,-147.4)and(60.3,-146.6)..(60.1,-146.1)
				..controls(59.9,-145.5)and(59.5,-144.9)..(59,-144.3)
				..controls(58.5,-143.7)and(58.2,-143.2)..(57.9,-142.5)
				..controls(57.7,-141.9)and(57.3,-141.1)..(57,-140.7)
				..controls(56.6,-140.3)and(56.1,-139.3)..(55.8,-138.5)
				..controls(55.5,-137.8)and(55.1,-136.8)..(54.9,-136.5)
				..controls(54.7,-136.1)and(54.3,-135.2)..(53.9,-134.5)
				..controls(53.6,-133.8)and(53,-132.7)..(52.6,-132)
				..controls(51.6,-130.5)and(51.4,-130.2)..(51.1,-129)
				..controls(50.6,-127.5)and(50.6,-127.4)..(49.7,-126.3)
				..controls(49.2,-125.7)and(48.6,-124.8)..(48.3,-124.2)
				..controls(48,-123.6)and(47.5,-122.7)..(47.3,-122.2)
				..controls(47,-121.8)and(46.5,-120.7)..(46.1,-119.8)
				..controls(45.8,-118.9)and(45.2,-117.8)..(44.9,-117.2)
				..controls(44.2,-116.1)and(43.3,-114.4)..(42.3,-112.2)
				..controls(41.8,-111.3)and(41.3,-110.4)..(41.1,-110.1)
				..controls(40.5,-109.4)and(40.3,-109)..(40,-107.8)
				..controls(39.8,-107)and(39.6,-106.5)..(39.2,-106)
				..controls(38.9,-105.6)and(38.4,-104.7)..(38,-104)
				..controls(37.7,-103.4)and(37.2,-102.5)..(37,-102.1)
				..controls(36.7,-101.7)and(36.3,-100.8)..(36,-100.1)
				..controls(35.6,-98.9)and(35.3,-98.4)..(34.6,-97.6)
				..controls(34.2,-97.1)and(34.2,-96.3)..(34.6,-96)
				..controls(35,-95.7)and(35.5,-95.9)..(36.2,-96.7)
				..controls(36.5,-97)and(37.2,-97.5)..(37.6,-97.7)..controls(38,-98)and(38.6,-98.4)..(39.1,-98.9)
				..controls(39.6,-99.4)and(40.3,-100)..(40.6,-100.3)
				..controls(41.3,-100.8)and(42.4,-101.8)..(43.7,-102.9)
				..controls(44,-103.2)and(44.6,-103.7)..(45,-103.9)
				..controls(45.4,-104.2)and(45.9,-104.6)..(46.1,-104.8)
				..controls(47.1,-105.9)and(48.8,-107.1)..(50.5,-107.9)
				..controls(51.8,-108.6)and(52.3,-109)..(53.4,-110)
				..controls(54.8,-111.2)and(55.6,-111.6)..(57,-112.1)
				..controls(57.5,-112.3)and(58.5,-112.8)..(59.4,-113.3)
				..controls(60.2,-113.7)and(61.3,-114.3)..(61.9,-114.5)
				..controls(62.4,-114.7)and(63.3,-115.1)..(63.8,-115.3)
				..controls(64.4,-115.6)and(65.4,-115.9)..(66.2,-116.1)
				..controls(67,-116.3)and(68.4,-116.6)..(69.2,-116.9)
				..controls(72.3,-117.9)and(72.8,-118)..(78,-118.1)
				..controls(84,-118.3)and(86.1,-118.1)..(88.1,-117.1)
				..controls(89.4,-116.5)and(93,-115)..(93.7,-114.9)
				..controls(94.8,-114.6)and(95.4,-114.2)..(95.9,-113.4)
				..controls(96.2,-113)and(96.7,-112.5)..(97.1,-112.2)
				..controls(98,-111.6)and(101.3,-108.4)..(102.1,-107.2)
				..controls(102.4,-106.8)and(102.9,-106.2)..(103.2,-105.8)
				..controls(103.5,-105.5)and(104,-104.8)..(104.2,-104.2)
				..controls(104.5,-103.6)and(105,-102.8)..(105.3,-102.4)
				..controls(105.6,-102)and(106.1,-101.2)..(106.5,-100.6)
				..controls(107.1,-99.5)and(108.6,-97.5)..(109.5,-96.5)
				..controls(109.9,-96.1)and(110.1,-96)..(110.4,-96)
				..controls(111.1,-96)and(111.6,-96.5)..(111.8,-97.6)
				..controls(112,-98.6)and(112.3,-99.3)..(112.8,-99.9)
				..controls(113.7,-100.9)and(115.2,-103.1)..(115.7,-104.2)
				..controls(116.1,-105)and(116.5,-105.5)..(116.9,-106)
				..controls(117.3,-106.4)and(117.8,-106.9)..(118,-107.3)
				..controls(118.2,-107.6)and(118.6,-108)..(119,-108.3)
				..controls(119.6,-108.8)and(121.1,-110.3)..(121.8,-111.2)
				..controls(122.1,-111.7)and(122.5,-112)..(123.4,-112.4)
				..controls(124.8,-113.2)and(125.8,-113.8)..(127,-114.6)
				..controls(127.5,-114.9)and(128.1,-115.3)..(128.2,-115.4)
				..controls(128.4,-115.5)and(129.3,-115.8)..(130.2,-116)
				..controls(131.1,-116.2)and(132.5,-116.6)..(133.3,-116.9)
				..controls(135.7,-117.9)and(136.1,-117.9)..(142.5,-117.8)
				..controls(145.8,-117.8)and(148.4,-117.7)..(148.8,-117.6)
				..controls(149.2,-117.6)and(150.2,-117.2)..(151,-116.8)
				..controls(152,-116.4)and(152.8,-116)..(153.6,-115.9)
				..controls(154.5,-115.7)and(155.2,-115.4)..(156.3,-114.9)
				..controls(157.1,-114.5)and(158.2,-114)..(158.7,-113.8)
				..controls(159.1,-113.7)and(159.8,-113.3)..(160.2,-113)
				..controls(160.6,-112.7)and(161.3,-112.4)..(161.7,-112.2)
				..controls(162.7,-111.9)and(163.2,-111.7)..(164.1,-111.1)
				..controls(164.5,-110.8)and(165.3,-110.4)..(165.9,-110)
				..controls(166.5,-109.7)and(167.3,-109.2)..(167.6,-108.9)
				..controls(167.9,-108.6)and(168.6,-108.2)..(169,-108)
				..controls(169.4,-107.8)and(170.2,-107.3)..(170.7,-106.9)
				..controls(171.2,-106.5)and(172.1,-105.8)..(172.7,-105.4)
				..controls(173.2,-105)and(174.1,-104.4)..(174.5,-104)
				..controls(174.9,-103.7)and(175.7,-103.1)..(176.2,-102.6)
				..controls(176.7,-102.2)and(177.4,-101.5)..(177.9,-101.2)
				..controls(178.3,-100.8)and(178.8,-100.4)..(179,-100.2)
				..controls(179.2,-100)and(179.6,-99.6)..(180,-99.3)
				..controls(180.4,-99.1)and(180.9,-98.6)..(181.2,-98.4)
				..controls(181.5,-98.1)and(182.1,-97.7)..(182.5,-97.5)
				..controls(183.2,-97.2)and(183.5,-96.9)..(184.3,-95.9)
				..controls(184.6,-95.6)and(184.6,-95.6)..(184.8,-95.7)
				..controls(185,-95.8)and(185.1,-96)..(185.1,-96.2)
				..controls(185.1,-96.5)and(185,-96.7)..(184.6,-97.2)
				..controls(184.1,-97.8)and(184,-97.9)..(183.9,-98.8)
				..controls(183.7,-100.3)and(183.3,-101.3)..(182.9,-101.8)
				..controls(182.6,-102.1)and(182.3,-102.6)..(182.2,-102.9)
				..controls(182,-103.3)and(181.6,-104)..(181.2,-104.6)
				..controls(180.9,-105.2)and(180.5,-106.1)..(180.4,-106.6)
				..controls(179.9,-108)and(179.7,-108.5)..(179,-109.4)
				..controls(178.6,-109.8)and(178.2,-110.5)..(177.9,-111)
				..controls(177.7,-111.5)and(177.3,-112.3)..(177,-112.7)
				..controls(176.6,-113.2)and(176.4,-113.7)..(176.3,-114.2)
				..controls(176.2,-114.7)and(175.9,-115.2)..(175.4,-116)
				..controls(175,-116.6)and(174.4,-117.7)..(174,-118.5)
				..controls(173.7,-119.2)and(173.3,-120.1)..(173,-120.5)
				..controls(172.8,-120.9)and(172.4,-121.8)..(172.1,-122.5)
				..controls(171.8,-123.2)and(171.3,-124.1)..(171,-124.6)
				..controls(170.8,-125.1)and(170.3,-126)..(169.9,-126.6)
				..controls(169.6,-127.3)and(169.2,-128.1)..(169,-128.3)
				..controls(168.8,-128.6)and(168.5,-129.3)..(168.3,-129.8)
				..controls(168.1,-130.3)and(167.8,-131.1)..(167.5,-131.6)
				..controls(166.9,-132.5)and(164.2,-137.9)..(164.2,-138.1)
				..controls(164.2,-138.1)and(163.9,-138.8)..(163.4,-139.7)
				..controls(162.9,-140.5)and(162.4,-141.5)..(162.3,-141.9)
				..controls(161.9,-143)and(161.6,-143.4)..(160.9,-144.1)
				..controls(160.4,-144.6)and(160.3,-144.8)..(160.1,-145.4)
				..controls(160,-145.8)and(159.6,-146.8)..(159.2,-147.8)
				..controls(158.7,-148.8)and(158.1,-150)..(157.9,-150.6)
				..controls(157.6,-151.1)and(157.2,-151.9)..(156.9,-152.4)
				..controls(156.5,-152.8)and(156.1,-153.7)..(155.8,-154.4)
				..controls(155.6,-155)and(155.2,-155.9)..(154.9,-156.3)
				..controls(154.7,-156.7)and(154.4,-157.4)..(154.2,-157.9)
				..controls(154.1,-158.4)and(153.8,-159)..(153.5,-159.3)
				..controls(152.8,-160.1)and(153,-160.1)..(142.4,-160.1)
				..controls(137.2,-160.1)and(131.8,-160.1)..(130.4,-160)--cycle
				(141.5,-103.7)..controls(140.8,-103)and(139.8,-102.1)..(139.2,-101.6)
				..controls(138.5,-101.2)and(137.9,-100.6)..(137.8,-100.4)
				..controls(137,-99.5)and(134,-96.5)..(133.2,-96)
				..controls(132.7,-95.7)and(131.8,-94.9)..(131.1,-94.3)
				..controls(129.4,-92.7)and(127.9,-91.6)..(126.6,-91.1)
				..controls(125.2,-90.5)and(124.1,-90)..(122.9,-89.1)
				..controls(122.1,-88.6)and(121.3,-88.2)..(119.9,-87.7)
				..controls(118.8,-87.3)and(117.5,-86.8)..(116.9,-86.5)
				..controls(114.3,-85.5)and(113.7,-85.4)..(109.7,-85.5)
				..controls(105.3,-85.7)and(103.9,-85.9)..(101.6,-87)
				..controls(100.7,-87.4)and(99.4,-87.8)..(97.9,-88)
				..controls(97,-88.2)and(96.6,-88.3)..(95.8,-88.9)
				..controls(95.4,-89.2)and(94.7,-89.6)..(94.2,-89.8)
				..controls(92.9,-90.3)and(92.3,-90.7)..(90.9,-92.1)
				..controls(90.2,-92.7)and(89.4,-93.3)..(89,-93.5)..controls(87,-94.5)and(85.2,-96)..(83.6,-97.9)
				..controls(83.3,-98.2)and(82.5,-98.9)..(81.8,-99.4)
				..controls(80.7,-100.1)and(80.5,-100.3)..(80,-101.1)
				..controls(78.7,-103.1)and(77.2,-103.9)..(76.4,-103)
				..controls(75.7,-102.2)and(75.7,-103.2)..(75.7,-83.6)
				..controls(75.6,-73.6)and(75.6,-64.8)..(75.5,-64)
				..controls(75.5,-63.2)and(75.5,-61.4)..(75.6,-59.8)--(75.7,-57.1)--(76,-56.8)
				..controls(76.8,-56.2)and(78,-56.3)..(79.8,-57.2)
				..controls(80.4,-57.5)and(81.3,-57.8)..(81.8,-57.9)
				..controls(82.3,-58.1)and(83,-58.4)..(83.6,-58.8)..controls(84.3,-59.2)and(85,-59.5)..(86,-59.8)
				..controls(86.8,-60.1)and(88,-60.6)..(88.7,-60.9)
				..controls(89.3,-61.2)and(90.6,-61.7)..(91.4,-62)
				..controls(92.2,-62.3)and(93.7,-62.9)..(94.6,-63.3)
				..controls(95.6,-63.7)and(96.8,-64.1)..(97.4,-64.3)
				..controls(98,-64.5)and(99,-64.9)..(99.7,-65.2)
				..controls(100.4,-65.6)and(101.5,-66)..(102.3,-66.2)
				..controls(104.2,-66.8)and(105,-67.1)..(106.2,-67.8)
				..controls(107.9,-68.7)and(110.2,-69.4)..(111.3,-69.2)
				..controls(112.5,-69)and(115,-68.2)..(115.6,-67.8)
				..controls(117.3,-66.7)and(117.6,-66.6)..(119.5,-66.2)
				..controls(121.9,-65.8)and(122.1,-65.7)..(123,-65.1)
				..controls(124.1,-64.4)and(124.4,-64.3)..(126,-63.9)
				..controls(127,-63.7)and(127.7,-63.4)..(128.4,-63.1)
				..controls(129,-62.8)and(129.7,-62.5)..(130.6,-62.4)
				..controls(132.1,-62.1)and(134,-61.5)..(135.2,-61)
				..controls(135.7,-60.8)and(136.7,-60.5)..(137.3,-60.3)
				..controls(138.4,-60)and(139.3,-59.7)..(141.1,-59)
				..controls(141.6,-58.7)and(142,-58.6)..(142.5,-58.6)
				..controls(143.2,-58.6)and(143.3,-58.7)..(143.6,-59)
				..controls(144,-59.5)and(144.1,-59.9)..(144,-60.8)
				..controls(144,-61.2)and(144,-71.1)..(144,-82.8)--(143.9,-104.1)--(143.7,-104.5)
				..controls(143.5,-104.7)and(143.3,-104.8)..(143,-104.9)
				..controls(142.6,-104.9)and(142.5,-104.8)..(141.5,-103.7)--cycle
				(27.8,-91.4)..controls(27.1,-90.9)and(26.3,-90)..(26.1,-89.5)
				..controls(26,-89.2)and(25.8,-88.6)..(25.7,-87.9)
				..controls(25.4,-86.2)and(25.6,-84.8)..(26.5,-83.5)
				..controls(27.4,-82.3)and(28.3,-82)..(29.4,-82.3)
				..controls(31.1,-82.8)and(32.1,-84.9)..(32,-87.4)
				..controls(31.9,-89.3)and(31.2,-90.6)..(29.9,-91.4)
				..controls(29.1,-91.9)and(28.6,-91.9)..(27.8,-91.4)--cycle
				(190.1,-91.3)..controls(189.7,-91.2)and(189.3,-90.9)..(189.1,-90.7)
				..controls(188.3,-89.8)and(187.6,-87.7)..(187.6,-86.2)
				..controls(187.6,-85.5)and(187.7,-85.2)..(188.2,-84.4)
				..controls(189.1,-82.6)and(190.1,-81.8)..(191.3,-82.1)
				..controls(192,-82.3)and(193.5,-83.7)..(193.8,-84.4)
				..controls(194.1,-85.1)and(194.2,-87.1)..(194,-88.3)
				..controls(193.8,-89.8)and(193.3,-90.5)..(192.3,-91.3)
				..controls(191.7,-91.7)and(190.9,-91.7)..(190.1,-91.3)--cycle
				(173.7,-83.4)..controls(172.2,-82.5)and(172,-82.5)..(170.6,-82.4)
				..controls(169.7,-82.4)and(168.8,-82.4)..(168.3,-82.5)
				..controls(167.2,-82.7)and(166.9,-82.7)..(166.6,-82.2)
				..controls(166.3,-81.8)and(166.3,-81)..(166.8,-80)
				..controls(167.2,-78.9)and(167.4,-77.4)..(167.7,-74.1)
				..controls(167.8,-72.8)and(167.9,-71.8)..(167.9,-71.7)
				..controls(168,-71.7)and(168,-70.2)..(168,-68.4)..controls(168,-65.4)and(168,-64.9)..(167.7,-63)
				..controls(167.3,-60.5)and(167.1,-59.9)..(166.4,-58.4)
				..controls(166.1,-57.7)and(165.8,-56.9)..(165.7,-56.6)
				..controls(165.7,-56.2)and(165.4,-55.6)..(165.2,-55.2)
				..controls(164.6,-53.8)and(161.9,-50.7)..(161,-50.3)
				..controls(159.8,-49.7)and(157.9,-50)..(155.2,-51.2)
				..controls(154.7,-51.4)and(153.4,-51.8)..(152.2,-52.1)
				..controls(150.8,-52.4)and(149.4,-52.9)..(148.6,-53.2)
				..controls(147.8,-53.6)and(146.7,-54)..(146,-54.2)
				..controls(145.4,-54.4)and(144.5,-54.7)..(144,-54.9)
				..controls(143.6,-55.2)and(142.7,-55.5)..(141.8,-55.7)
				..controls(141,-55.9)and(139.6,-56.3)..(138.7,-56.6)
				..controls(137.8,-57)and(136.3,-57.4)..(135.3,-57.6)
				..controls(133.8,-57.9)and(133.2,-58.1)..(132.1,-58.7)
				..controls(131.1,-59.2)and(130.3,-59.4)..(129.4,-59.6)
				..controls(128.7,-59.8)and(127.8,-60)..(127.4,-60.1)
				..controls(127,-60.3)and(126,-60.5)..(125.1,-60.7)
				..controls(123.2,-61.1)and(121.7,-61.5)..(120.2,-62.1)
				..controls(119.7,-62.3)and(118.1,-62.8)..(116.7,-63.2)
				..controls(114.6,-63.9)and(114,-64.1)..(113,-64.7)
				..controls(111.9,-65.4)and(111.8,-65.4)..(111,-65.5)
				..controls(109.9,-65.5)and(109.3,-65.4)..(107.7,-64.7)
				..controls(107,-64.4)and(106,-64)..(105.3,-63.9)
				..controls(103.5,-63.5)and(102.2,-63)..(101.2,-62.5)
				..controls(100.1,-61.8)and(99.3,-61.5)..(97.8,-61.1)
				..controls(97.2,-61)and(96.2,-60.7)..(95.7,-60.4)
				..controls(95.1,-60.1)and(94.2,-59.8)..(93.5,-59.6)
				..controls(92.9,-59.4)and(91.7,-59)..(91,-58.7)..controls(90.2,-58.4)and(88.9,-58)..(88,-57.7)
				..controls(87.2,-57.5)and(86.1,-57.1)..(85.5,-56.8)
				..controls(85,-56.6)and(84.1,-56.2)..(83.5,-56)..controls(82.9,-55.7)and(81.8,-55.2)..(81,-54.9)
				..controls(80.1,-54.5)and(79.1,-54.2)..(78.7,-54.1)
				..controls(77.9,-53.8)and(77.3,-53.6)..(76.2,-52.9)
				..controls(75.7,-52.6)and(75.2,-52.4)..(74.1,-52.2)
				..controls(72.5,-51.9)and(72.4,-51.8)..(71.3,-51.1)
				..controls(69.9,-50.1)and(69.4,-49.9)..(67.9,-49.8)
				..controls(67,-49.7)and(66.6,-49.6)..(66.3,-49.4)
				..controls(66.1,-49.3)and(65.7,-49)..(65.4,-48.8)
				..controls(65.1,-48.7)and(64.2,-48.4)..(63.6,-48.2)
				..controls(62.9,-48)and(61.7,-47.6)..(60.9,-47.2)
				..controls(60.1,-46.8)and(58.9,-46.4)..(58.2,-46.2)
				..controls(57.6,-46)and(56.8,-45.7)..(56.5,-45.5)
				..controls(54.9,-44.6)and(54.6,-44.5)..(53.2,-44.2)
				..controls(51.7,-43.8)and(51,-43.6)..(49.9,-42.9)
				..controls(49.5,-42.6)and(48.9,-42.4)..(47.5,-42)
				..controls(45.7,-41.6)and(45.6,-41.5)..(44.3,-40.6)
				..controls(43.5,-40.1)and(42.8,-39.6)..(42.7,-39.4)
				..controls(42.4,-38.7)and(43,-38.1)..(44.5,-37.6)..controls(44.9,-37.5)and(46,-37)..(46.9,-36.6)
				..controls(47.9,-36.1)and(48.8,-35.8)..(49.6,-35.6)
				..controls(50.5,-35.4)and(51.3,-35.2)..(52.4,-34.6)
				..controls(53.3,-34.1)and(54.5,-33.7)..(54.9,-33.5)
				..controls(55.4,-33.4)and(56.4,-33)..(57.2,-32.6)
				..controls(58.6,-32)and(60.1,-31.5)..(61.6,-31.2)
				..controls(62.3,-31.1)and(62.7,-31)..(63.2,-30.7)
				..controls(64.1,-30)and(64.3,-29.9)..(65.3,-29.7)
				..controls(65.7,-29.6)and(66.5,-29.2)..(67.1,-28.9)
				..controls(67.6,-28.6)and(68.8,-28.2)..(69.6,-27.9)
				..controls(70.5,-27.7)and(71.6,-27.3)..(72.1,-27)
				..controls(72.6,-26.7)and(73.8,-26.2)..(74.8,-25.9)
				..controls(75.7,-25.5)and(76.7,-25.1)..(77.1,-25)
				..controls(77.4,-24.8)and(78.5,-24.4)..(79.6,-24.1)
				..controls(80.6,-23.8)and(82,-23.2)..(82.8,-22.9)
				..controls(83.5,-22.5)and(84.8,-22)..(85.8,-21.8)
				..controls(86.7,-21.5)and(87.8,-21.1)..(88.1,-20.9)
				..controls(88.4,-20.7)and(89.3,-20.3)..(90.1,-20.1)
				..controls(90.8,-19.8)and(91.9,-19.4)..(92.4,-19.1)
				..controls(93.5,-18.5)and(94.4,-18.2)..(95.4,-18.1)
				..controls(95.7,-18)and(96.4,-17.8)..(96.9,-17.5)
				..controls(97.5,-17.3)and(98.4,-16.9)..(99,-16.6)
				..controls(99.7,-16.3)and(100.9,-15.9)..(101.9,-15.6)
				..controls(102.9,-15.3)and(104.4,-14.7)..(105.3,-14.3)
				..controls(108.1,-13.2)and(108.5,-13.1)..(110.1,-13.1)
				..controls(111.4,-13.1)and(111.5,-13.1)..(112.4,-13.5)
				..controls(113.2,-14)and(115,-14.8)..(116.6,-15.4)
				..controls(116.9,-15.5)and(117.9,-15.8)..(118.8,-16.1)
				..controls(120.9,-16.6)and(121.5,-16.9)..(122.6,-17.9)
				..controls(123,-18.3)and(123.7,-18.8)..(123.9,-18.9)
				..controls(124.6,-19.4)and(126,-19.8)..(127.9,-20.2)
				..controls(129.6,-20.5)and(129.8,-20.6)..(131.1,-21.4)
				..controls(131.5,-21.7)and(132.4,-22.1)..(133.1,-22.3)
				..controls(136.1,-23.2)and(138.3,-24.1)..(140.2,-25)
				..controls(141,-25.5)and(141.5,-25.6)..(142.5,-25.8)
				..controls(145.1,-26.3)and(146,-26.6)..(149.4,-28.1)
				..controls(151.3,-28.9)and(153.5,-29.8)..(154.3,-30.1)
				..controls(155.1,-30.4)and(156.2,-30.8)..(156.6,-31)
				..controls(157.1,-31.2)and(158.2,-31.6)..(159.1,-31.8)
				..controls(160.6,-32.2)and(160.9,-32.3)..(161.7,-32.9)
				..controls(162.4,-33.3)and(162.9,-33.6)..(163.7,-33.8)
				..controls(164.3,-34)and(164.8,-34.1)..(164.9,-34.1)
				..controls(165,-34.1)and(165.6,-34.4)..(166.4,-34.8)
				..controls(167.3,-35.2)and(168.1,-35.4)..(169,-35.6)
				..controls(170,-35.8)and(170.5,-36)..(171.7,-36.6)
				..controls(172.4,-37)and(173.3,-37.4)..(173.6,-37.4)
				..controls(175.6,-37.7)and(176.2,-38)..(176.7,-39.1)
				..controls(176.9,-39.6)and(176.9,-39.7)..(176.7,-40.4)
				..controls(176.4,-41.6)and(175.8,-42.7)..(175.2,-43.4)
				..controls(174.5,-44.1)and(172.6,-45.5)..(171.7,-45.8)
				..controls(170.6,-46.2)and(169.1,-46)..(168.2,-45.4)
				..controls(167.9,-45.2)and(167.3,-44.8)..(166.9,-44.4)
				..controls(166.4,-44.1)and(165.8,-43.6)..(165.4,-43.2)
				..controls(165.1,-42.9)and(164.4,-42.3)..(163.9,-41.9)
				..controls(163.3,-41.5)and(162.7,-40.9)..(162.4,-40.5)
				..controls(161.4,-39.2)and(159.9,-38)..(158.8,-37.8)
				..controls(158.4,-37.8)and(158.3,-37.7)..(158.2,-37.2)
				..controls(158,-36.8)and(157.4,-36.1)..(157.1,-36.1)
				..controls(157,-36.1)and(156.7,-36.1)..(156.5,-36)
				..controls(156.1,-35.9)and(155.9,-35.7)..(155.7,-35.2)
				..controls(155.3,-34.6)and(155.1,-34.4)..(154,-34.1)
				..controls(153.6,-34)and(153.2,-33.7)..(152.9,-33.5)
				..controls(152.4,-32.9)and(151.6,-32.4)..(151.3,-32.4)
				..controls(150.9,-32.4)and(150.4,-32.6)..(150.4,-32.9)
				..controls(150.4,-33.2)and(150.8,-33.4)..(151.4,-33.6)
				..controls(152.1,-33.7)and(152.8,-34.1)..(153.5,-34.9)
				..controls(154.1,-35.5)and(154.3,-35.6)..(155,-35.8)
				..controls(155.7,-36)and(155.7,-36)..(156,-36.6)
				..controls(156.3,-37.3)and(156.7,-37.7)..(157.5,-37.8)
				..controls(158.1,-37.9)and(158.2,-38)..(158.4,-38.8)
				..controls(158.5,-39.4)and(158.6,-39.6)..(159.6,-40.4)
				..controls(160.2,-41)and(160.9,-41.6)..(161.3,-41.8)
				..controls(161.7,-42)and(161.9,-42.3)..(162.2,-42.8)
				..controls(162.4,-43.2)and(162.8,-43.8)..(163.1,-44.1)
				..controls(163.9,-44.8)and(165.1,-46.7)..(165.7,-48)
				..controls(166.2,-49.3)and(166.4,-49.7)..(167.2,-50.5)
				..controls(167.5,-50.8)and(167.9,-51.4)..(168.1,-51.8)
				..controls(168.3,-52.2)and(168.8,-53)..(169.1,-53.6)
				..controls(169.6,-54.3)and(169.8,-54.9)..(169.9,-55.6)
				..controls(170.4,-57.8)and(170.5,-58.1)..(171.2,-59.5)
				..controls(171.6,-60.3)and(172,-61.5)..(172.2,-62.1)
				..controls(172.3,-62.8)and(172.7,-63.7)..(172.9,-64.1)
				..controls(173.2,-64.6)and(173.4,-65.3)..(173.4,-65.6)
				..controls(173.8,-68.8)and(174.2,-70.6)..(174.9,-72.6)
				..controls(175.5,-74.5)and(175.9,-77.1)..(176.1,-80.2)
				..controls(176.3,-83.4)and(175.5,-84.4)..(173.7,-83.4)--cycle;
				\definecolor{ezs}{RGB}{247,179,26}\fill[ezs](-5.7,-156.6)--(-3.4,-157.5)--(61.9,12.7)--(59.7,13.6)--cycle;
				\definecolor{ezs}{RGB}{247,179,22}\fill[ezs](-3.8,-157.3)--(-1.5,-158.2)--(63.9,12)--(61.6,12.9)--cycle;
				\definecolor{ezs}{RGB}{248,179,19}\fill[ezs](-1.8,-158.1)--(0.4,-158.9)--(65.8,11.2)--(63.5,12.1)--cycle;
				\definecolor{ezs}{RGB}{248,178,16}\fill[ezs](0.1,-158.8)--(2.4,-159.7)--(67.7,10.5)--(65.4,11.4)--cycle;
				\definecolor{ezs}{RGB}{249,178,15}\fill[ezs](2,-159.5)--(4.3,-160.4)--(69.6,9.8)--(67.3,10.7)--cycle;
				\definecolor{ezs}{RGB}{251,177,17}\fill[ezs](3.9,-160.3)--(6.2,-161.1)--(71.5,9)--(69.2,9.9)--cycle;
				\definecolor{ezs}{RGB}{252,179,18}\fill[ezs](5.8,-161)--(8.1,-161.9)--(73.4,8.3)--(71.1,9.2)--cycle;
				\definecolor{ezs}{RGB}{252,182,17}\fill[ezs](7.7,-161.7)--(10,-162.6)--(75.3,7.6)--(73,8.5)--cycle;
				\definecolor{ezs}{RGB}{252,185,17}\fill[ezs](9.6,-162.5)--(11.9,-163.4)--(77.2,6.8)--(75,7.7)--cycle;
				\definecolor{ezs}{RGB}{251,189,18}\fill[ezs](11.5,-163.2)--(13.8,-164.1)--(79.2,6.1)--(76.9,7)--cycle;
				\definecolor{ezs}{RGB}{251,191,22}\fill[ezs](13.4,-163.9)--(15.7,-164.8)--(81.1,5.4)--(78.8,6.2)--cycle;
				\definecolor{ezs}{RGB}{252,193,25}\fill[ezs](15.4,-164.7)--(17.6,-165.6)--(83,4.6)--(80.7,5.5)--cycle;
				\definecolor{ezs}{RGB}{253,195,27}\fill[ezs](17.3,-165.4)--(19.6,-166.3)--(84.9,3.9)--(82.6,4.8)--cycle;
				\definecolor{ezs}{RGB}{253,196,31}\fill[ezs](19.2,-166.1)--(21.5,-167)--(86.8,3.2)--(84.5,4)--cycle;
				\definecolor{ezs}{RGB}{254,199,35}\fill[ezs](21.1,-166.9)--(23.4,-167.8)--(88.7,2.4)--(86.4,3.3)--cycle;
				\definecolor{ezs}{RGB}{253,201,38}\fill[ezs](23,-167.6)--(25.3,-168.5)--(90.6,1.7)--(88.3,2.6)--cycle;
				\definecolor{ezs}{RGB}{253,203,40}\fill[ezs](24.9,-168.3)--(27.2,-169.2)--(92.5,1)--(90.3,1.8)--cycle;
				\definecolor{ezs}{RGB}{253,205,42}\fill[ezs](26.8,-169.1)--(29.1,-170)--(94.4,0.2)--(92.2,1.1)--cycle;
				\definecolor{ezs}{RGB}{253,207,45}\fill[ezs](28.7,-169.8)--(31,-170.7)--(96.4,-0.5)--(94.1,0.4)--cycle;
				\definecolor{ezs}{RGB}{253,210,49}\fill[ezs](30.7,-170.5)--(32.9,-171.4)--(98.3,-1.2)--(96,-0.4)--cycle;
				\definecolor{ezs}{RGB}{253,212,54}\fill[ezs](32.6,-171.3)--(34.9,-172.2)--(100.2,-2)--(97.9,-1.1)--cycle;
				\definecolor{ezs}{RGB}{253,214,59}\fill[ezs](34.5,-172)--(36.8,-172.9)--(102.1,-2.7)--(99.8,-1.8)--cycle;
				\definecolor{ezs}{RGB}{253,216,63}\fill[ezs](36.4,-172.7)--(38.7,-173.6)--(104,-3.4)--(101.7,-2.6)--cycle;
				\definecolor{ezs}{RGB}{253,218,65}\fill[ezs](38.3,-173.5)--(40.6,-174.4)--(105.9,-4.2)--(103.6,-3.3)--cycle;
				\definecolor{ezs}{RGB}{254,220,66}\fill[ezs](40.2,-174.2)--(42.5,-175.1)--(107.8,-4.9)--(105.5,-4)--cycle;
				\definecolor{ezs}{RGB}{254,222,68}\fill[ezs](42.1,-175)--(44.4,-175.8)--(109.7,-5.6)--(107.5,-4.8)--cycle;
				\definecolor{ezs}{RGB}{254,224,71}\fill[ezs](44,-175.7)--(46.3,-176.6)--(111.7,-6.4)--(109.4,-5.5)--cycle;
				\definecolor{ezs}{RGB}{254,226,75}\fill[ezs](46,-176.4)--(48.2,-177.3)--(113.6,-7.1)--(111.3,-6.2)--cycle;
				\definecolor{ezs}{RGB}{254,228,79}\fill[ezs](47.9,-177.2)--(50.1,-178)--(115.5,-7.8)--(113.2,-7)--cycle;
				\definecolor{ezs}{RGB}{254,230,82}\fill[ezs](49.8,-177.9)--(52.1,-178.8)--(117.4,-8.6)--(115.1,-7.7)--cycle;
				\definecolor{ezs}{RGB}{254,232,85}\fill[ezs](51.7,-178.6)--(54,-179.5)--(119.3,-9.3)--(117,-8.4)--cycle;
				\definecolor{ezs}{RGB}{254,234,88}\fill[ezs](53.6,-179.4)--(55.9,-180.2)--(121.2,-10)--(118.9,-9.2)--cycle;
				\definecolor{ezs}{RGB}{254,237,91}\fill[ezs](55.5,-180.1)--(57.8,-181)--(123.1,-10.8)--(120.8,-9.9)--cycle;
				\definecolor{ezs}{RGB}{254,239,94}\fill[ezs](57.4,-180.8)--(59.7,-181.7)--(125,-11.5)--(122.8,-10.6)--cycle;
				\definecolor{ezs}{RGB}{254,241,97}\fill[ezs](59.3,-181.6)--(61.6,-182.4)--(127,-12.2)--(124.7,-11.4)--cycle;
				\definecolor{ezs}{RGB}{254,243,101}\fill[ezs](61.2,-182.3)--(63.5,-183.2)--(128.9,-13)--(126.6,-12.1)--cycle;
				\definecolor{ezs}{RGB}{254,245,103}\fill[ezs](63.2,-183)--(65.4,-183.9)--(130.8,-13.7)--(128.5,-12.8)--cycle;
				\definecolor{ezs}{RGB}{254,245,105}\fill[ezs](65.1,-183.8)--(67.4,-184.6)--(132.7,-14.4)--(130.4,-13.6)--cycle;
				\definecolor{ezs}{RGB}{254,245,105}\fill[ezs](67,-184.5)--(69.3,-185.4)--(134.6,-15.2)--(132.3,-14.3)--cycle;
				\definecolor{ezs}{RGB}{254,244,103}\fill[ezs](68.9,-185.2)--(71.2,-186.1)--(136.5,-15.9)--(134.2,-15)--cycle;
				\definecolor{ezs}{RGB}{253,243,101}\fill[ezs](70.8,-186)--(73.1,-186.8)--(138.4,-16.6)--(136.1,-15.8)--cycle;
				\definecolor{ezs}{RGB}{253,241,98}\fill[ezs](72.7,-186.7)--(75,-187.6)--(140.3,-17.4)--(138,-16.5)--cycle;
				\definecolor{ezs}{RGB}{253,239,95}\fill[ezs](74.6,-187.4)--(76.9,-188.3)--(142.2,-18.1)--(140,-17.2)--cycle;
				\definecolor{ezs}{RGB}{253,236,92}\fill[ezs](76.5,-188.2)--(78.8,-189)--(144.2,-18.8)--(141.9,-18)--cycle;
				\definecolor{ezs}{RGB}{253,234,88}\fill[ezs](78.5,-188.9)--(80.7,-189.8)--(146.1,-19.6)--(143.8,-18.7)--cycle;
				\definecolor{ezs}{RGB}{253,232,85}\fill[ezs](80.4,-189.6)--(82.7,-190.5)--(148,-20.3)--(145.7,-19.4)--cycle;
				\definecolor{ezs}{RGB}{253,230,82}\fill[ezs](82.3,-190.4)--(84.6,-191.2)--(149.9,-21)--(147.6,-20.2)--cycle;
				\definecolor{ezs}{RGB}{253,228,78}\fill[ezs](84.2,-191.1)--(86.5,-192)--(151.8,-21.8)--(149.5,-20.9)--cycle;
				\definecolor{ezs}{RGB}{253,226,74}\fill[ezs](86.1,-191.8)--(88.4,-192.7)--(153.7,-22.5)--(151.4,-21.6)--cycle;
				\definecolor{ezs}{RGB}{253,224,70}\fill[ezs](88,-192.6)--(90.3,-193.4)--(155.6,-23.3)--(153.3,-22.4)--cycle;
				\definecolor{ezs}{RGB}{253,222,67}\fill[ezs](89.9,-193.3)--(92.2,-194.2)--(157.5,-24)--(155.3,-23.1)--cycle;
				\definecolor{ezs}{RGB}{253,220,65}\fill[ezs](91.8,-194)--(94.1,-194.9)--(159.5,-24.7)--(157.2,-23.8)--cycle;
				\definecolor{ezs}{RGB}{253,217,63}\fill[ezs](93.7,-194.8)--(96,-195.6)--(161.4,-25.5)--(159.1,-24.6)--cycle;
				\definecolor{ezs}{RGB}{253,215,60}\fill[ezs](95.7,-195.5)--(97.9,-196.4)--(163.3,-26.2)--(161,-25.3)--cycle;
				\definecolor{ezs}{RGB}{253,213,56}\fill[ezs](97.6,-196.2)--(99.9,-197.1)--(165.2,-26.9)--(162.9,-26)--cycle;
				\definecolor{ezs}{RGB}{252,211,53}\fill[ezs](99.5,-197)--(101.8,-197.8)--(167.1,-27.7)--(164.8,-26.8)--cycle;
				\definecolor{ezs}{RGB}{252,209,49}\fill[ezs](101.4,-197.7)--(103.7,-198.6)--(169,-28.4)--(166.7,-27.5)--cycle;
				\definecolor{ezs}{RGB}{252,207,45}\fill[ezs](103.3,-198.4)--(105.6,-199.3)--(170.9,-29.1)--(168.6,-28.2)--cycle;
				\definecolor{ezs}{RGB}{253,205,42}\fill[ezs](105.2,-199.2)--(107.5,-200)--(172.8,-29.9)--(170.6,-29)--cycle;
				\definecolor{ezs}{RGB}{253,202,40}\fill[ezs](107.1,-199.9)--(109.4,-200.8)--(174.7,-30.6)--(172.5,-29.7)--cycle;
				\definecolor{ezs}{RGB}{253,200,38}\fill[ezs](109,-200.6)--(111.3,-201.5)--(176.7,-31.3)--(174.4,-30.4)--cycle;
				\definecolor{ezs}{RGB}{253,198,35}\fill[ezs](111,-201.4)--(113.2,-202.2)--(178.6,-32.1)--(176.3,-31.2)--cycle;
				\definecolor{ezs}{RGB}{253,196,31}\fill[ezs](112.9,-202.1)--(115.2,-203)--(180.5,-32.8)--(178.2,-31.9)--cycle;
				\definecolor{ezs}{RGB}{253,194,28}\fill[ezs](114.8,-202.8)--(117.1,-203.7)--(182.4,-33.5)--(180.1,-32.6)--cycle;
				\definecolor{ezs}{RGB}{253,192,26}\fill[ezs](116.7,-203.6)--(119,-204.5)--(184.3,-34.3)--(182,-33.4)--cycle;
				\definecolor{ezs}{RGB}{253,189,24}\fill[ezs](118.6,-204.3)--(120.9,-205.2)--(186.2,-35)--(183.9,-34.1)--cycle;
				\definecolor{ezs}{RGB}{253,186,20}\fill[ezs](120.5,-205)--(122.8,-205.9)--(188.1,-35.7)--(185.8,-34.8)--cycle;
				\definecolor{ezs}{RGB}{253,184,17}\fill[ezs](122.4,-205.8)--(124.7,-206.7)--(190,-36.5)--(187.8,-35.6)--cycle;
				\definecolor{ezs}{RGB}{252,183,15}\fill[ezs](124.3,-206.5)--(126.6,-207.4)--(192,-37.2)--(189.7,-36.3)--cycle;
				\definecolor{ezs}{RGB}{252,181,14}\fill[ezs](126.3,-207.2)--(128.5,-208.1)--(193.9,-37.9)--(191.6,-37.1)--cycle;
				\definecolor{ezs}{RGB}{252,179,13}\fill[ezs](128.2,-208)--(130.4,-208.9)--(195.8,-38.7)--(193.5,-37.8)--cycle;
				\definecolor{ezs}{RGB}{252,178,12}\fill[ezs](130.1,-208.7)--(132.4,-209.6)--(197.7,-39.4)--(195.4,-38.5)--cycle;
				\definecolor{ezs}{RGB}{251,178,10}\fill[ezs](132,-209.4)--(134.3,-210.3)--(199.6,-40.1)--(197.3,-39.3)--cycle;
				\definecolor{ezs}{RGB}{251,178,10}\fill[ezs](133.9,-210.2)--(136.2,-211.1)--(201.5,-40.9)--(199.2,-40)--cycle;
				\definecolor{ezs}{RGB}{250,177,11}\fill[ezs](135.8,-210.9)--(138.1,-211.8)--(203.4,-41.6)--(201.1,-40.7)--cycle;
				\definecolor{ezs}{RGB}{250,178,11}\fill[ezs](137.7,-211.6)--(140,-212.5)--(205.3,-42.3)--(203.1,-41.5)--cycle;
				\definecolor{ezs}{RGB}{248,180,10}\fill[ezs](139.6,-212.4)--(141.9,-213.3)--(207.3,-43.1)--(205,-42.2)--cycle;
				\definecolor{ezs}{RGB}{247,181,9}\fill[ezs](141.5,-213.1)--(143.8,-214)--(209.2,-43.8)--(206.9,-42.9)--cycle;
				\definecolor{ezs}{RGB}{246,182,9}\fill[ezs](143.5,-213.8)--(145.7,-214.7)--(211.1,-44.5)--(208.8,-43.7)--cycle;
				\definecolor{ezs}{RGB}{246,182,10}\fill[ezs](145.4,-214.6)--(147.7,-215.5)--(213,-45.3)--(210.7,-44.4)--cycle;
			\end{scope}
			% chữ EDUZY (Times đậm, làm dày nét bằng 5 lớp lệch nhẹ; căn đúng khung chữ logo gốc)
			\foreach \dx/\dy in {-0.7/0,0.7/0,0/-0.3,0/0.3,0/0}
			\node[anchor=center,inner sep=0pt,outer sep=0pt,text=ezText,align=center] at (109.6+\dx,-193.6+\dy)
			{\resizebox{1.01cm}{0.22cm}{\fontencoding{T1}\fontfamily{ptm}\fontseries{b}\selectfont EDUZY}};
		\end{tikzpicture}%
	}%
}
%====================== THIẾT LẬP THÂN ĐỀ ======================
%--- Môi trường tự luận {bt}: ex_test 3.3.4 chưa khai báo sẵn nên tạo ở đây
%    (ex_test đã tự móc \fixshowans và \listenumerate cho bt)
\makeatletter
\@ifundefined{bt}{\newtheorem{bt}{\color{gold1}\namebt}}{}
\makeatother
\renewcommand{\sectionmark}[1]{\markright{\thesection\ #1}}
%--- Khung tiêu đề (logo + tối đa 4 dòng; dòng 3, 4 để trống {} thì bỏ qua)
%  Đề thi:    \khungtieude{THI THỬ THPT QUỐC GIA NĂM 2026}{Môn Toán -- Đề số 21}
%                         {(Đề thi gồm: ...)}{Thời gian làm bài: 90 phút (không kể thời gian phát đề).}
%  Chuyên đề: \khungtieude{CHUYÊN ĐỀ 3 -- DÃY SỐ VÀ CÔNG THỨC TRUY HỒI}
%                         {Bồi dưỡng học sinh giỏi Toán THPT}{...}{}
\newcommand{\khungtieude}[4]{%
	\begin{tcolorbox}[colframe=gold1,colback=gold5!45,boxrule=0.5mm]
		\begin{minipage}[c]{3cm}
			\centerline{\eduzylogo[2.4cm]}
		\end{minipage}\hspace{0.5cm}
		\begin{minipage}[c]{12.5cm}
			\vskip 0.15 cm
			\centerline{\ind{#1}}
			\vskip 0.15 cm
			\centerline{\textbf{#2}}
			\if\relax\detokenize{#3}\relax\else
				\vskip 0.15 cm
				\centerline{\textit{#3}}
			\fi
			\if\relax\detokenize{#4}\relax\else
				\vskip 0.15 cm
				\centerline{\textit{#4}}
			\fi
			\vskip 0.15 cm
		\end{minipage}
	\end{tcolorbox}
}
%====================== BỔ SUNG CHO BÀI HỌC, CHUYÊN ĐỀ ======================
\usepackage[shortlabels]{enumitem}
\usepackage{tabularx}
%--- Danh sách gọn: \begin{enumerate}[gon] ... \end{enumerate}
\SetEnumitemKey{gon}{leftmargin=*,itemsep=1pt,parsep=0pt,topsep=3pt}
%--- Khung lý thuyết: \begin{khung}{} ... \end{khung}; tiêu đề chỉ dùng khi thật ngắn ("Mẹo", "Lỗi thường gặp")
\newtcolorbox{khung}[1]{breakable,enhanced,colframe=gold1,colback=gold5!45,
	boxrule=0.4mm,fonttitle=\bfseries\sffamily,coltitle=white,colbacktitle=gold1,
	attach boxed title to top left={xshift=6mm,yshift=-3mm},
	boxed title style={arc=1mm,boxrule=0.3mm},title={#1}}
%--- Tiêu đề mục nhỏ trong một dạng: \muc{Dấu hiệu}, \muc{Phương pháp}, \muc{Bài tự luyện}
\makeatletter
\newcommand{\muc}[1]{\par\medskip\noindent\indam{#1}\par\nobreak\smallskip\@afterindentfalse\@afterheading}
\makeatother
%--- Ô công thức cần nhớ (căn giữa, viền vàng, vừa khít nội dung)
%     \begin{center}\congthuc{$Q(n)=(n-1)!$}\end{center}
\newtcbox{\congthuc}{on line,colframe=gold1,colback=gold5!45,boxrule=0.5pt,arc=2pt,
	left=8pt,right=8pt,top=4pt,bottom=4pt}
%--- Công tắc lời giải bài tự luyện (ví dụ mẫu luôn dùng \loigiai)
%     \loigiaiTL[{đáp số}]{lời giải đầy đủ}
%     \TLfull -> in lời giải đầy đủ;  \TLdapso -> chỉ in một dòng "Đáp số";  \TLan -> ẩn hoàn toàn
%     (\CheDo là dethi thì mọi lời giải và đáp số đều ẩn)
\makeatletter
\newif\ifTLhien   \TLhientrue
\newif\ifTLdapso  \TLdapsofalse
\newcommand{\TLfull}{\TLhientrue\TLdapsofalse}
\newcommand{\TLdapso}{\TLhientrue\TLdapsotrue}
\newcommand{\TLan}{\TLhienfalse}
\newcommand{\loigiaiTL}[2][]{%
	\ifTLhien
		\ifTLdapso
			\def\@tlarg{#1}%
			\ifx\@tlarg\@empty\else
				\ifdefstring{\kindOPTN}{dethi}{}{\par\noindent{\bfseries\color{gold1}Đáp số:} #1\par}%
			\fi
		\else
			\loigiai{#2}%
		\fi
	\fi}
\makeatother
\TLdapso% mặc định: chỉ in đáp số bài tự luyện
%====================== BẢNG ĐÚNG/SAI (\choiceTFt) ======================
%  Hàng tiêu đề nền nâu đồng chữ trắng (Phát biểu | Đ | S), các hàng tô xen kẽ kem / trắng,
%  đường kẻ vàng kim, nhãn a) b) c) d) màu nâu đồng. Bản có đáp án: cột được chọn có vòng tròn
%  nâu đồng với dấu tích trắng (câu đúng tích ở cột Đúng, câu sai tích ở cột Sai).
%  Tiêu đề ghi đủ "Đúng | Sai": \OPTN{viettat=0} sau \input (cột Đúng/Sai rộng ra, phát biểu hẹp lại).
%  Dùng bảng gốc của ex_test: \def\BangDSgoc{1} trước %============================================================
%  PREAMBLE EDUZY v2 - dùng chung cho đề thi, bài học, chuyên đề (gói ex_test)
%  Khác bản preamble-eduzy.tex ở đầu trang, chân trang: dải nền nâu đồng tràn ngang,
%  logo + Eduzy.com.vn, ô chữ vàng kem cạnh xiên, chân trang "Trang (số)";
%  logo Eduzy chìm giữa mọi trang (tắt: \logochimtat, bật lại: \logochimbat).
%  Cách dùng ở file chính (xem mau-chuan.tex, mau-chuyen-de.tex):
%      \def\CheDo{loigiai}          % loigiai | dethi | solcolor | color | book
%      \def\tenchuong{<dòng chữ trên đầu trang>}
%      \def\thuonghieu{<chữ ở chân trang lẻ>}   % không khai báo thì dùng mặc định
%      \def\tacgia{<chữ ở chân trang chẵn>}      % không khai báo thì dùng mặc định
%      \def\khauhieu{<chữ cạnh logo ở đầu trang>}  % tùy chọn, mặc định Eduzy.com.vn
%      \input{preamble-eduzy-v2}
%      \TLdapso                      % tài liệu: \TLfull | \TLdapso | \TLan
%      \begin{document} ... \end{document}
%  Logo Eduzy vẽ bằng TikZ ngay trong file (\eduzylogo), không cần file ảnh.
%  Biên dịch: pdfLaTeX (máy có gói vietnam) hoặc XeLaTeX.
%============================================================
\documentclass[12pt,a4paper,twoside]{article}%<-- twoside: bat dau trang/chan trang chan-le
\usepackage{iftex}
\usepackage{amsmath,amssymb,tikz,mathrsfs,tkz-tab}
\usepackage{graphicx}
\usepackage{pgfplots}% đồ thị hàm số
\pgfplotsset{compat=1.17}
%\usepackage{mathptmx}
\usepackage{mathpazo}
\ifXeTeX% XeLaTeX: font chữ có đủ dấu tiếng Việt, giữ nguyên font toán mathpazo
	\usepackage[no-math]{fontspec}
	\setmainfont{TeX Gyre Pagella}
	\newfontfamily\EZqag{TeX Gyre Adventor}[NFSSFamily=qag,Ligatures=TeX]
\fi
\usepackage[many]{tcolorbox}
\usetikzlibrary{math,through,calc,intersections,angles,quotes,shapes,shapes.geometric,arrows,patterns,matrix,chains,arrows.meta,decorations.shapes,decorations.fractals,decorations.markings,shadows,decorations.pathreplacing}
\usepackage[top=1.6cm, bottom=1.8cm, left=1.8cm, right=1.8cm]{geometry}% v2: top 1.6cm (bản cũ 1.5cm) để chữ không sát dải đầu trang
\providecommand{\CheDo}{loigiai}% loigiai | dethi | solcolor | color | book
\expandafter\usepackage\expandafter[\CheDo]{ex_test}
% \tenchuong: khai báo ở file chính (trước \input) hoặc \renewcommand sau \input
%============================================================
%  [GỘP TỪ eduzy-de.tex] EDUZY - MÀU NÂU ĐỒNG / VÀNG KIM
%  + ĐẦU TRANG, CHÂN TRANG KIỂU DẢI NGANG (v2), ĐỐI XỨNG TRANG CHẴN / TRANG LẺ.
%  (không cần file eduzy-de.tex riêng nữa)
%  Khối này phải đặt SAU \usepackage[...]{ex_test} và SAU \def\tenchuong{...}
%============================================================
\usepackage{lastpage}
\usepackage{needspace}
%--- DÀN TRANG ĐỀ THI: tách câu (THÊM 05/11/2026; chỉ thêm lệnh mới, KHÔNG đổi hành vi mặc định) ---
%  Hai công tắc, mặc định tắt (= preamble chạy y như cũ). Bật bằng cách đặt TRƯỚC \input{preamble-eduzy-v2}
%  trong file đề (tools/dan_trang/dan_trang.py tự ghi khi đặt \NeedspaceCau):
%    \def\TachCauDS{1}   câu Đúng/Sai được tách: câu dẫn (kèm hình) ở cuối trang, bảng Phát biểu sang trang sau
%                        (câu dẫn và hình không tách nhau, bảng không tách giữa hai trang)
%    \def\TachCauLG{1}   bản lời giải (\CheDo khác dethi): phần đề của câu cao hơn 1/3 trang được phép tách
%  Lệnh dùng trong file đề (giá trị do tools/dan_trang/dan_trang.py đo ở bản dethi):
%    \NeedspaceCau[<cao câu dẫn + hình>]{<cao cả khối câu>}   ngay trên \begin{ex} (câu đầu phần: trên tiêu đề phần)
%        công tắc tắt: đúng bằng \Needspace{<cao cả khối câu>} như cũ;
%        \TachCauDS=1 và có [..] (câu Đúng/Sai được chọn tách): chỉ giữ câu dẫn + hình;
%        \TachCauLG=1 ở bản lời giải: đòi tối đa 1/3 \textheight.
%    \NeedspaceBang{<cao bảng Phát biểu>}%   ngay trước \choiceTFt: \TachCauDS=1 thì giữ trọn bảng trên một trang;
%        công tắc tắt: không làm gì.
\providecommand{\TachCauDS}{0}
\providecommand{\TachCauLG}{0}
\newlength{\EZnsCau}
\newcommand{\NeedspaceCau}[2][]{%
	\setlength{\EZnsCau}{#2}%
	\ifnum\TachCauDS=1 \if\relax\detokenize{#1}\relax\else\setlength{\EZnsCau}{#1}\fi\fi
	\ifnum\TachCauLG=1 \ifdefstring{\kindOPTN}{dethi}{}{\ifdim\EZnsCau>0.3333\textheight\setlength{\EZnsCau}{0.3333\textheight}\fi}\fi
	\Needspace{\EZnsCau}}
\newcommand{\NeedspaceBang}[1]{\ifnum\TachCauDS=1 \Needspace{#1}\fi}

%--- Thang màu: đậm -> nhạt
\definecolor{gold1}{RGB}{126,96,0}    % nâu đồng đậm  #7E6000
\definecolor{gold2}{RGB}{192,144,3}   % vàng kim      #C09003
\definecolor{gold3}{RGB}{255,215,99}  % vàng nhạt     #FFD763
\definecolor{gold4}{RGB}{255,231,153} % vàng rất nhạt #FFE799
\definecolor{gold5}{RGB}{255,243,205} % vàng kem      #FFF3CD
%--- Ghi đè tông xanh/đỏ cũ
\definecolor{myblue}{RGB}{126,96,0}
\colorlet{blue}{gold1}
\colorlet{cyan}{gold2}
\colorlet{teal}{gold2}

%--- Tên chương mặc định (file đề nên tự khai báo lại)
\providecommand{\tenchuong}{}
%--- Chân trang
\providecommand{\thuonghieu}{TOÁN 12 -- Eduzy.com.vn}
\providecommand{\tacgia}{Tô Gia Bảo -- 0353703783}

%--- Chữ cạnh logo ở đầu trang
\providecommand{\khauhieu}{Eduzy.com.vn}

%====================== THÔNG SỐ ĐẦU TRANG, CHÂN TRANG (v2) ======================
%  Dải nền đậm tràn ngang trang; ô chữ vàng kem có cạnh xiên, kèm 2 vạch xiên chuyển màu.
%  Đầu trang lẻ : logo + \khauhieu (trái)   | ô chữ \tenchuong (phải)
%  Chân trang lẻ: ô chữ \thuonghieu (trái)  | Trang (số) (phải)
%  Trang chẵn soi gương trang lẻ; chân trang chẵn ghi \tacgia.
%--- Màu: chỉ dùng thang vàng kim Eduzy
\colorlet{ezNen}{gold1}% dải nền đậm
\colorlet{ezVach1}{gold3}% vạch xiên sát dải nền (đậm hơn)
\colorlet{ezVach2}{gold4}% vạch xiên sát ô chữ (nhạt hơn)
\colorlet{ezO}{gold5}% ô chữ vàng kem
%--- Kích thước
\newlength{\EZdai}\setlength{\EZdai}{1.1cm}% chiều cao dải nền
\newlength{\EZo}\setlength{\EZo}{0.68cm}% chiều cao ô chữ
\newlength{\EZxien}\setlength{\EZxien}{0.55cm}% độ lệch ngang của cạnh xiên
\newlength{\EZvach}\setlength{\EZvach}{0.75cm}% bề rộng mỗi vạch xiên
\newlength{\EZkhe}\setlength{\EZkhe}{0.12cm}% khe trắng giữa vạch và ô chữ
\newlength{\EZdem}\setlength{\EZdem}{0.4cm}% khoảng đệm từ chữ tới cạnh xiên
\newlength{\EZle}\setlength{\EZle}{1.8cm}% lề trái/phải của chữ (bằng lề thân bài)
\newlength{\EZrongmin}\setlength{\EZrongmin}{4.5cm}% ô chữ rộng tối thiểu
\newlength{\EZrong}% bề rộng chữ trong ô (tính tự động)
%--- Font
\newcommand{\EZfontO}{\fontsize{10.5pt}{12pt}\selectfont\bfseries}% chữ trong ô
\newcommand{\EZfontKH}{\fontfamily{qag}\fontsize{11pt}{12pt}\selectfont\bfseries}% \khauhieu cạnh logo
\newcommand{\EZfontTrang}{\fontsize{11.5pt}{12pt}\selectfont\bfseries}% chữ "Trang" và số trang
%--- Đo bề rộng chữ trong ô
\newcommand{\EZdorong}[1]{%
	\settowidth{\EZrong}{\EZfontO #1}%
	\ifdim\EZrong<\EZrongmin\setlength{\EZrong}{\EZrongmin}\fi}
%--- Logo + chữ trắng \khauhieu (dựng sẵn một lần khi vào \begin{document})
\newsavebox{\EZlogobox}
\AtBeginDocument{\sbox{\EZlogobox}{%
		\raisebox{-0.5\height}{\eduzylogo[1cm]}\hspace{0.15cm}%
		\raisebox{-0.5\height}{\color{white}\EZfontKH\khauhieu}}}
%--- Hai vạch xiên + khe trắng. #1: điểm gốc ở mép dưới, #2: hướng (1 sang phải, -1 sang trái),
%    #3: hướng nghiêng (1 kiểu "/", -1 kiểu "\")
\newcommand{\EZvachxien}[3]{%
	\fill[white] (#1) --++({(#3)*\EZxien},\EZo) --++({(#2)*\EZkhe},0) --++({-(#3)*\EZxien},-\EZo) --cycle;
	\fill[ezVach2] ([xshift=(#2)*\EZkhe]#1) --++({(#3)*\EZxien},\EZo) --++({(#2)*\EZvach},0) --++({-(#3)*\EZxien},-\EZo) --cycle;
	\fill[ezVach1] ([xshift=(#2)*(\EZkhe+\EZvach)]#1) --++({(#3)*\EZxien},\EZo) --++({(#2)*\EZvach},0) --++({-(#3)*\EZxien},-\EZo) --cycle;
	\draw[white,line width=0.6pt] ([xshift=(#2)*(\EZkhe+\EZvach)]#1) --++({(#3)*\EZxien},\EZo);
}
%--- Chữ trong ô (căn giữa theo chiều cao ô, không lệch vì dấu tiếng Việt)
\tikzset{ezchu/.style={inner sep=0pt,text=ezNen,text height=0.72em,text depth=0.22em}}

%====================== LOGO CHÌM GIỮA TRANG ======================
%  Vẽ cùng đầu trang nên nằm dưới chữ của thân bài, có trên mọi trang.
%  Tắt: \logochimtat (từ trang đó trở đi), bật lại: \logochimbat.
%  Đổi cỡ, độ đậm: khai báo trước \input, ví dụ \def\EZchimcao{9cm} \def\EZchimdo{0.1}
\providecommand{\EZchimcao}{10cm}% chiều cao logo chìm
\providecommand{\EZchimdo}{0.12}% độ đậm (0 = trong suốt, 1 = rõ hẳn)
\newif\ifEZchim \EZchimtrue
\newcommand{\logochimtat}{\EZchimfalse}
\newcommand{\logochimbat}{\EZchimtrue}
%  Dựng sẵn logo mờ một lần khi vào \begin{document}; transparency group làm mờ cả logo
%  một lượt nên không lộ vệt ở chỗ các dải màu chồng nhau.
\newsavebox{\EZchimtho}
\newsavebox{\EZchimbox}
\AtBeginDocument{%
	\sbox{\EZchimtho}{\eduzylogo[\EZchimcao]}%
	\sbox{\EZchimbox}{\begin{tikzpicture}
			\begin{scope}[transparency group,opacity=\EZchimdo]
				\node[inner sep=0pt] {\usebox{\EZchimtho}};
			\end{scope}
	\end{tikzpicture}}}
\newcommand{\EZvechim}{%
	\ifEZchim
		\node[inner sep=0pt] at (current page.center) {\usebox{\EZchimbox}};
	\fi}

%====================== ĐẦU TRANG - TRANG LẺ ======================
%  (logo bên trái, ô chữ \tenchuong bên phải)
\newcommand{\myfancyheadOdd}{%
	\EZdorong{\tenchuong}%
	\begin{tikzpicture}[remember picture,overlay]
		\coordinate (T) at (current page.north west);
		\coordinate (P) at (current page.north east);
		%--- logo chìm giữa trang
		\EZvechim
		%--- dải nền
		\fill[ezNen] (T) rectangle ([yshift=-\EZdai]P);
		%--- ô chữ: X là góc dưới bên trái, cạnh trái xiên kiểu "/"
		\coordinate (X) at ([xshift=-\EZle-\EZrong-\EZdem-\EZxien,yshift=-\EZdai]P);
		\fill[ezO] (X) --++(\EZxien,\EZo) -- ([yshift=-\EZdai+\EZo]P) -- ([yshift=-\EZdai]P) --cycle;
		\EZvachxien{X}{-1}{1}
		\node[ezchu,anchor=east] at ([xshift=-\EZle,yshift=-\EZdai+0.5*\EZo]P) {\EZfontO\tenchuong};
		%--- logo + tên
		\node[inner sep=0pt,anchor=west] at ([xshift=\EZle,yshift=-0.5*\EZdai]T) {\usebox{\EZlogobox}};
	\end{tikzpicture}%
}

%====================== ĐẦU TRANG - TRANG CHẴN =====================
%  (soi gương: ô chữ \tenchuong bên trái, logo bên phải)
\newcommand{\myfancyheadEven}{%
	\EZdorong{\tenchuong}%
	\begin{tikzpicture}[remember picture,overlay]
		\coordinate (T) at (current page.north west);
		\coordinate (P) at (current page.north east);
		%--- logo chìm giữa trang
		\EZvechim
		%--- dải nền
		\fill[ezNen] (T) rectangle ([yshift=-\EZdai]P);
		%--- ô chữ: X là góc dưới bên phải, cạnh phải xiên kiểu "\"
		\coordinate (X) at ([xshift=\EZle+\EZrong+\EZdem+\EZxien,yshift=-\EZdai]T);
		\fill[ezO] (X) --++(-\EZxien,\EZo) -- ([yshift=-\EZdai+\EZo]T) -- ([yshift=-\EZdai]T) --cycle;
		\EZvachxien{X}{1}{-1}
		\node[ezchu,anchor=west] at ([xshift=\EZle,yshift=-\EZdai+0.5*\EZo]T) {\EZfontO\tenchuong};
		%--- logo + tên
		\node[inner sep=0pt,anchor=east] at ([xshift=-\EZle,yshift=-0.5*\EZdai]P) {\usebox{\EZlogobox}};
	\end{tikzpicture}%
}

%====================== CHÂN TRANG - TRANG LẺ ======================
%  (ô chữ \thuonghieu bên trái, "Trang (số)" bên phải)
\newcommand{\myfancyfootOdd}{%
	\EZdorong{\thuonghieu}%
	\begin{tikzpicture}[remember picture,overlay]
		\coordinate (T) at (current page.south west);
		\coordinate (P) at (current page.south east);
		%--- dải nền
		\fill[ezNen] (T) rectangle ([yshift=\EZdai]P);
		%--- ô chữ sát mép trên dải: X là góc dưới bên phải, cạnh phải xiên kiểu "/"
		\coordinate (X) at ([xshift=\EZle+\EZrong+\EZdem,yshift=\EZdai-\EZo]T);
		\fill[ezO] ([yshift=\EZdai-\EZo]T) -- (X) --++(\EZxien,\EZo) -- ([yshift=\EZdai]T) --cycle;
		\EZvachxien{X}{1}{1}
		\node[ezchu,anchor=west] at ([xshift=\EZle,yshift=\EZdai-0.5*\EZo]T) {\EZfontO\thuonghieu};
		%--- số trang
		\node[circle,draw=white,line width=0.9pt,text=white,inner sep=0pt,minimum size=0.62cm,anchor=east]
		(so) at ([xshift=-\EZle,yshift=0.5*\EZdai]P) {\EZfontTrang\thepage};
		\node[inner sep=0pt,text=white,anchor=base east] (tr) at ([xshift=-0.15cm]so.base west) {\EZfontTrang Trang};
	\end{tikzpicture}%
}

%====================== CHÂN TRANG - TRANG CHẴN ====================
%  (soi gương: "Trang (số)" bên trái, ô chữ \tacgia bên phải)
\newcommand{\myfancyfootEven}{%
	\EZdorong{\tacgia}%
	\begin{tikzpicture}[remember picture,overlay]
		\coordinate (T) at (current page.south west);
		\coordinate (P) at (current page.south east);
		%--- dải nền
		\fill[ezNen] (T) rectangle ([yshift=\EZdai]P);
		%--- ô chữ sát mép trên dải: X là góc dưới bên trái, cạnh trái xiên kiểu "\"
		\coordinate (X) at ([xshift=-\EZle-\EZrong-\EZdem,yshift=\EZdai-\EZo]P);
		\fill[ezO] ([yshift=\EZdai-\EZo]P) -- (X) --++(-\EZxien,\EZo) -- ([yshift=\EZdai]P) --cycle;
		\EZvachxien{X}{-1}{-1}
		\node[ezchu,anchor=east] at ([xshift=-\EZle,yshift=\EZdai-0.5*\EZo]P) {\EZfontO\tacgia};
		%--- số trang
		%    (vòng tròn ẩn so0 nằm giữa dải để lấy đường chân chữ, giống trang lẻ)
		\node[circle,inner sep=0pt,minimum size=0.62cm] (so0) at ([yshift=0.5*\EZdai]T) {\EZfontTrang\phantom{\thepage}};
		\node[inner sep=0pt,text=white,anchor=base west] (tr) at ($(so0.base -| T)+(\EZle,0)$) {\EZfontTrang Trang};
		\node[circle,draw=white,line width=0.9pt,text=white,inner sep=0pt,minimum size=0.62cm,anchor=base west]
		(so) at ([xshift=0.15cm]tr.base east) {\EZfontTrang\thepage};
	\end{tikzpicture}%
}

%--- Giữ tương thích với các file cũ đang gọi \myfancyhead, \myfancyfoot
\newcommand{\myfancyhead}{\myfancyheadOdd}
\newcommand{\myfancyfoot}{\myfancyfootOdd}

%======================= GẮN VÀO TRANG =============================
\pagestyle{fancy}
\fancyhf{}
\fancyhead[LO]{\myfancyheadOdd}
\fancyhead[RE]{\myfancyheadEven}
\fancyfoot[LO]{\myfancyfootOdd}
\fancyfoot[RE]{\myfancyfootEven}
\renewcommand{\headrulewidth}{0pt}
\renewcommand{\footrulewidth}{0pt}
\fancypagestyle{plain}{%
	\fancyhf{}
	\fancyhead[LO]{\myfancyheadOdd}
	\fancyhead[RE]{\myfancyheadEven}
	\fancyfoot[LO]{\myfancyfootOdd}
	\fancyfoot[RE]{\myfancyfootEven}
	\renewcommand{\headrulewidth}{0pt}
	\renewcommand{\footrulewidth}{0pt}
}
%=================== HẾT PHẦN GỘP TỪ eduzy-de.tex ===================
\newcommand{\indam}[1]{
	{\fontsize{12pt}{9pt}\sffamily \selectfont\bfseries\color{gold1} #1}
}
\newcommand{\ind}[1]{
	{\fontsize{14pt}{9pt}\fontfamily{qag} \selectfont\bfseries\color{gold1} #1}
}
\renewcommand{\FalseEX}{\stepcounter{dapan}\textbf{\selectfont\color{gold1}\Alph{dapan}.}\ignorespaces}
\renewtheorem{ex}{\color{gold1}Câu}
\renewcommand{\baselinestretch}{1.1}
\newcommand{\hoac}[1]{
	\left[\begin{aligned}#1\end{aligned}\right.}
\newcommand{\heva}[1]{
	\left\{\begin{aligned}#1\end{aligned}\right.}
\usepackage{esvect}
\def\vec{\vv}
\def\vt{\vv}
\def\overrightarrow{\vv}
%====================== LOGO EDUZY VẼ BẰNG TIKZ ======================
%  Đường viền dò từ logo gốc; dải màu vàng kim đo trực tiếp trên ảnh gốc
%  (nghiêng 21 độ, tô bằng 80 dải hẹp nên hiển thị giống nhau trên mọi trình xem PDF).
%  Dùng: \eduzylogo          -> cao 2.4cm (mặc định)
%        \eduzylogo[1.5cm]   -> cao 1.5cm
\definecolor{ezText}{RGB}{245,200,46}
\newcommand{\eduzylogo}[1][2.4cm]{%
	\resizebox{!}{#1}{%
		\begin{tikzpicture}[x=0.01cm,y=0.01cm]
			\path[use as bounding box] (0,0) rectangle (224,-224);
			% mũ tốt nghiệp + vương miện
			\begin{scope}
				\clip
				(119.4,-176.1)..controls(118.9,-176.1)and(116.8,-176)..(114.6,-176)
				..controls(101.4,-176)and(85.8,-175.9)..(83.7,-175.7)
				..controls(81.8,-175.6)and(81,-175.6)..(80.2,-175.7)
				..controls(78.8,-175.9)and(77.7,-175.9)..(76.2,-175.7)
				..controls(74.9,-175.5)and(70.8,-175.5)..(64.3,-175.7)
				..controls(59.1,-175.9)and(58.3,-175.7)..(58.6,-173.7)
				..controls(58.6,-172.8)and(59.6,-169.6)..(60.3,-168)
				..controls(60.8,-166.9)and(61.4,-165.8)..(62.3,-165)
				..controls(63.1,-164.2)and(64.3,-163.4)..(65,-163.2)
				..controls(65.5,-163.1)and(72.3,-163)..(76.6,-163.1)
				..controls(77.9,-163.1)and(80.1,-163.1)..(81.4,-163)
				..controls(82.8,-162.9)and(86.3,-162.9)..(89.2,-162.9)
				..controls(93,-162.9)and(94.9,-162.8)..(96,-162.7)
				..controls(97.2,-162.5)and(101,-162.5)..(114,-162.4)
				..controls(123,-162.3)and(134.9,-162.2)..(140.2,-162.1)
				..controls(150.1,-162)and(152.2,-162.1)..(153.6,-162.4)
				..controls(154.4,-162.6)and(157.1,-163.8)..(158.1,-164.4)
				..controls(158.8,-164.8)and(159.7,-165.7)..(159.9,-166.3)
				..controls(160,-166.5)and(160.2,-167.1)..(160.2,-167.6)
				..controls(160.3,-168)and(160.6,-169.4)..(160.9,-170.6)
				..controls(161.2,-171.7)and(161.5,-173.1)..(161.6,-173.6)
				..controls(161.8,-174.4)and(161.8,-174.5)..(161.5,-174.9)
				..controls(161.2,-175.6)and(160.8,-175.9)..(159.9,-176)
				..controls(159.2,-176.1)and(154.1,-176.2)..(151.5,-176.2)
				..controls(148.1,-176.1)and(122.6,-176.1)..(121.6,-176.2)
				..controls(120.8,-176.2)and(119.8,-176.2)..(119.4,-176.1)--cycle
				(130.4,-160)..controls(129,-160)and(119.8,-160)..(110.1,-159.9)
				..controls(100.3,-159.8)and(88.5,-159.7)..(83.9,-159.6)
				..controls(77.8,-159.5)and(75.3,-159.5)..(74.5,-159.6)
				..controls(73.9,-159.7)and(72.5,-159.8)..(71.3,-159.8)
				..controls(69.6,-159.8)and(69,-159.8)..(68.5,-159.6)
				..controls(67.6,-159.3)and(66.8,-158.6)..(66.2,-157.6)
				..controls(66,-157.3)and(65.6,-156.5)..(65.3,-156.1)
				..controls(64.6,-155.1)and(64.5,-154.9)..(63.7,-152.9)
				..controls(63.2,-151.6)and(62.9,-150.9)..(61.8,-149.6)
				..controls(61.5,-149.2)and(61.1,-148.5)..(60.9,-147.9)
				..controls(60.7,-147.4)and(60.3,-146.6)..(60.1,-146.1)
				..controls(59.9,-145.5)and(59.5,-144.9)..(59,-144.3)
				..controls(58.5,-143.7)and(58.2,-143.2)..(57.9,-142.5)
				..controls(57.7,-141.9)and(57.3,-141.1)..(57,-140.7)
				..controls(56.6,-140.3)and(56.1,-139.3)..(55.8,-138.5)
				..controls(55.5,-137.8)and(55.1,-136.8)..(54.9,-136.5)
				..controls(54.7,-136.1)and(54.3,-135.2)..(53.9,-134.5)
				..controls(53.6,-133.8)and(53,-132.7)..(52.6,-132)
				..controls(51.6,-130.5)and(51.4,-130.2)..(51.1,-129)
				..controls(50.6,-127.5)and(50.6,-127.4)..(49.7,-126.3)
				..controls(49.2,-125.7)and(48.6,-124.8)..(48.3,-124.2)
				..controls(48,-123.6)and(47.5,-122.7)..(47.3,-122.2)
				..controls(47,-121.8)and(46.5,-120.7)..(46.1,-119.8)
				..controls(45.8,-118.9)and(45.2,-117.8)..(44.9,-117.2)
				..controls(44.2,-116.1)and(43.3,-114.4)..(42.3,-112.2)
				..controls(41.8,-111.3)and(41.3,-110.4)..(41.1,-110.1)
				..controls(40.5,-109.4)and(40.3,-109)..(40,-107.8)
				..controls(39.8,-107)and(39.6,-106.5)..(39.2,-106)
				..controls(38.9,-105.6)and(38.4,-104.7)..(38,-104)
				..controls(37.7,-103.4)and(37.2,-102.5)..(37,-102.1)
				..controls(36.7,-101.7)and(36.3,-100.8)..(36,-100.1)
				..controls(35.6,-98.9)and(35.3,-98.4)..(34.6,-97.6)
				..controls(34.2,-97.1)and(34.2,-96.3)..(34.6,-96)
				..controls(35,-95.7)and(35.5,-95.9)..(36.2,-96.7)
				..controls(36.5,-97)and(37.2,-97.5)..(37.6,-97.7)..controls(38,-98)and(38.6,-98.4)..(39.1,-98.9)
				..controls(39.6,-99.4)and(40.3,-100)..(40.6,-100.3)
				..controls(41.3,-100.8)and(42.4,-101.8)..(43.7,-102.9)
				..controls(44,-103.2)and(44.6,-103.7)..(45,-103.9)
				..controls(45.4,-104.2)and(45.9,-104.6)..(46.1,-104.8)
				..controls(47.1,-105.9)and(48.8,-107.1)..(50.5,-107.9)
				..controls(51.8,-108.6)and(52.3,-109)..(53.4,-110)
				..controls(54.8,-111.2)and(55.6,-111.6)..(57,-112.1)
				..controls(57.5,-112.3)and(58.5,-112.8)..(59.4,-113.3)
				..controls(60.2,-113.7)and(61.3,-114.3)..(61.9,-114.5)
				..controls(62.4,-114.7)and(63.3,-115.1)..(63.8,-115.3)
				..controls(64.4,-115.6)and(65.4,-115.9)..(66.2,-116.1)
				..controls(67,-116.3)and(68.4,-116.6)..(69.2,-116.9)
				..controls(72.3,-117.9)and(72.8,-118)..(78,-118.1)
				..controls(84,-118.3)and(86.1,-118.1)..(88.1,-117.1)
				..controls(89.4,-116.5)and(93,-115)..(93.7,-114.9)
				..controls(94.8,-114.6)and(95.4,-114.2)..(95.9,-113.4)
				..controls(96.2,-113)and(96.7,-112.5)..(97.1,-112.2)
				..controls(98,-111.6)and(101.3,-108.4)..(102.1,-107.2)
				..controls(102.4,-106.8)and(102.9,-106.2)..(103.2,-105.8)
				..controls(103.5,-105.5)and(104,-104.8)..(104.2,-104.2)
				..controls(104.5,-103.6)and(105,-102.8)..(105.3,-102.4)
				..controls(105.6,-102)and(106.1,-101.2)..(106.5,-100.6)
				..controls(107.1,-99.5)and(108.6,-97.5)..(109.5,-96.5)
				..controls(109.9,-96.1)and(110.1,-96)..(110.4,-96)
				..controls(111.1,-96)and(111.6,-96.5)..(111.8,-97.6)
				..controls(112,-98.6)and(112.3,-99.3)..(112.8,-99.9)
				..controls(113.7,-100.9)and(115.2,-103.1)..(115.7,-104.2)
				..controls(116.1,-105)and(116.5,-105.5)..(116.9,-106)
				..controls(117.3,-106.4)and(117.8,-106.9)..(118,-107.3)
				..controls(118.2,-107.6)and(118.6,-108)..(119,-108.3)
				..controls(119.6,-108.8)and(121.1,-110.3)..(121.8,-111.2)
				..controls(122.1,-111.7)and(122.5,-112)..(123.4,-112.4)
				..controls(124.8,-113.2)and(125.8,-113.8)..(127,-114.6)
				..controls(127.5,-114.9)and(128.1,-115.3)..(128.2,-115.4)
				..controls(128.4,-115.5)and(129.3,-115.8)..(130.2,-116)
				..controls(131.1,-116.2)and(132.5,-116.6)..(133.3,-116.9)
				..controls(135.7,-117.9)and(136.1,-117.9)..(142.5,-117.8)
				..controls(145.8,-117.8)and(148.4,-117.7)..(148.8,-117.6)
				..controls(149.2,-117.6)and(150.2,-117.2)..(151,-116.8)
				..controls(152,-116.4)and(152.8,-116)..(153.6,-115.9)
				..controls(154.5,-115.7)and(155.2,-115.4)..(156.3,-114.9)
				..controls(157.1,-114.5)and(158.2,-114)..(158.7,-113.8)
				..controls(159.1,-113.7)and(159.8,-113.3)..(160.2,-113)
				..controls(160.6,-112.7)and(161.3,-112.4)..(161.7,-112.2)
				..controls(162.7,-111.9)and(163.2,-111.7)..(164.1,-111.1)
				..controls(164.5,-110.8)and(165.3,-110.4)..(165.9,-110)
				..controls(166.5,-109.7)and(167.3,-109.2)..(167.6,-108.9)
				..controls(167.9,-108.6)and(168.6,-108.2)..(169,-108)
				..controls(169.4,-107.8)and(170.2,-107.3)..(170.7,-106.9)
				..controls(171.2,-106.5)and(172.1,-105.8)..(172.7,-105.4)
				..controls(173.2,-105)and(174.1,-104.4)..(174.5,-104)
				..controls(174.9,-103.7)and(175.7,-103.1)..(176.2,-102.6)
				..controls(176.7,-102.2)and(177.4,-101.5)..(177.9,-101.2)
				..controls(178.3,-100.8)and(178.8,-100.4)..(179,-100.2)
				..controls(179.2,-100)and(179.6,-99.6)..(180,-99.3)
				..controls(180.4,-99.1)and(180.9,-98.6)..(181.2,-98.4)
				..controls(181.5,-98.1)and(182.1,-97.7)..(182.5,-97.5)
				..controls(183.2,-97.2)and(183.5,-96.9)..(184.3,-95.9)
				..controls(184.6,-95.6)and(184.6,-95.6)..(184.8,-95.7)
				..controls(185,-95.8)and(185.1,-96)..(185.1,-96.2)
				..controls(185.1,-96.5)and(185,-96.7)..(184.6,-97.2)
				..controls(184.1,-97.8)and(184,-97.9)..(183.9,-98.8)
				..controls(183.7,-100.3)and(183.3,-101.3)..(182.9,-101.8)
				..controls(182.6,-102.1)and(182.3,-102.6)..(182.2,-102.9)
				..controls(182,-103.3)and(181.6,-104)..(181.2,-104.6)
				..controls(180.9,-105.2)and(180.5,-106.1)..(180.4,-106.6)
				..controls(179.9,-108)and(179.7,-108.5)..(179,-109.4)
				..controls(178.6,-109.8)and(178.2,-110.5)..(177.9,-111)
				..controls(177.7,-111.5)and(177.3,-112.3)..(177,-112.7)
				..controls(176.6,-113.2)and(176.4,-113.7)..(176.3,-114.2)
				..controls(176.2,-114.7)and(175.9,-115.2)..(175.4,-116)
				..controls(175,-116.6)and(174.4,-117.7)..(174,-118.5)
				..controls(173.7,-119.2)and(173.3,-120.1)..(173,-120.5)
				..controls(172.8,-120.9)and(172.4,-121.8)..(172.1,-122.5)
				..controls(171.8,-123.2)and(171.3,-124.1)..(171,-124.6)
				..controls(170.8,-125.1)and(170.3,-126)..(169.9,-126.6)
				..controls(169.6,-127.3)and(169.2,-128.1)..(169,-128.3)
				..controls(168.8,-128.6)and(168.5,-129.3)..(168.3,-129.8)
				..controls(168.1,-130.3)and(167.8,-131.1)..(167.5,-131.6)
				..controls(166.9,-132.5)and(164.2,-137.9)..(164.2,-138.1)
				..controls(164.2,-138.1)and(163.9,-138.8)..(163.4,-139.7)
				..controls(162.9,-140.5)and(162.4,-141.5)..(162.3,-141.9)
				..controls(161.9,-143)and(161.6,-143.4)..(160.9,-144.1)
				..controls(160.4,-144.6)and(160.3,-144.8)..(160.1,-145.4)
				..controls(160,-145.8)and(159.6,-146.8)..(159.2,-147.8)
				..controls(158.7,-148.8)and(158.1,-150)..(157.9,-150.6)
				..controls(157.6,-151.1)and(157.2,-151.9)..(156.9,-152.4)
				..controls(156.5,-152.8)and(156.1,-153.7)..(155.8,-154.4)
				..controls(155.6,-155)and(155.2,-155.9)..(154.9,-156.3)
				..controls(154.7,-156.7)and(154.4,-157.4)..(154.2,-157.9)
				..controls(154.1,-158.4)and(153.8,-159)..(153.5,-159.3)
				..controls(152.8,-160.1)and(153,-160.1)..(142.4,-160.1)
				..controls(137.2,-160.1)and(131.8,-160.1)..(130.4,-160)--cycle
				(141.5,-103.7)..controls(140.8,-103)and(139.8,-102.1)..(139.2,-101.6)
				..controls(138.5,-101.2)and(137.9,-100.6)..(137.8,-100.4)
				..controls(137,-99.5)and(134,-96.5)..(133.2,-96)
				..controls(132.7,-95.7)and(131.8,-94.9)..(131.1,-94.3)
				..controls(129.4,-92.7)and(127.9,-91.6)..(126.6,-91.1)
				..controls(125.2,-90.5)and(124.1,-90)..(122.9,-89.1)
				..controls(122.1,-88.6)and(121.3,-88.2)..(119.9,-87.7)
				..controls(118.8,-87.3)and(117.5,-86.8)..(116.9,-86.5)
				..controls(114.3,-85.5)and(113.7,-85.4)..(109.7,-85.5)
				..controls(105.3,-85.7)and(103.9,-85.9)..(101.6,-87)
				..controls(100.7,-87.4)and(99.4,-87.8)..(97.9,-88)
				..controls(97,-88.2)and(96.6,-88.3)..(95.8,-88.9)
				..controls(95.4,-89.2)and(94.7,-89.6)..(94.2,-89.8)
				..controls(92.9,-90.3)and(92.3,-90.7)..(90.9,-92.1)
				..controls(90.2,-92.7)and(89.4,-93.3)..(89,-93.5)..controls(87,-94.5)and(85.2,-96)..(83.6,-97.9)
				..controls(83.3,-98.2)and(82.5,-98.9)..(81.8,-99.4)
				..controls(80.7,-100.1)and(80.5,-100.3)..(80,-101.1)
				..controls(78.7,-103.1)and(77.2,-103.9)..(76.4,-103)
				..controls(75.7,-102.2)and(75.7,-103.2)..(75.7,-83.6)
				..controls(75.6,-73.6)and(75.6,-64.8)..(75.5,-64)
				..controls(75.5,-63.2)and(75.5,-61.4)..(75.6,-59.8)--(75.7,-57.1)--(76,-56.8)
				..controls(76.8,-56.2)and(78,-56.3)..(79.8,-57.2)
				..controls(80.4,-57.5)and(81.3,-57.8)..(81.8,-57.9)
				..controls(82.3,-58.1)and(83,-58.4)..(83.6,-58.8)..controls(84.3,-59.2)and(85,-59.5)..(86,-59.8)
				..controls(86.8,-60.1)and(88,-60.6)..(88.7,-60.9)
				..controls(89.3,-61.2)and(90.6,-61.7)..(91.4,-62)
				..controls(92.2,-62.3)and(93.7,-62.9)..(94.6,-63.3)
				..controls(95.6,-63.7)and(96.8,-64.1)..(97.4,-64.3)
				..controls(98,-64.5)and(99,-64.9)..(99.7,-65.2)
				..controls(100.4,-65.6)and(101.5,-66)..(102.3,-66.2)
				..controls(104.2,-66.8)and(105,-67.1)..(106.2,-67.8)
				..controls(107.9,-68.7)and(110.2,-69.4)..(111.3,-69.2)
				..controls(112.5,-69)and(115,-68.2)..(115.6,-67.8)
				..controls(117.3,-66.7)and(117.6,-66.6)..(119.5,-66.2)
				..controls(121.9,-65.8)and(122.1,-65.7)..(123,-65.1)
				..controls(124.1,-64.4)and(124.4,-64.3)..(126,-63.9)
				..controls(127,-63.7)and(127.7,-63.4)..(128.4,-63.1)
				..controls(129,-62.8)and(129.7,-62.5)..(130.6,-62.4)
				..controls(132.1,-62.1)and(134,-61.5)..(135.2,-61)
				..controls(135.7,-60.8)and(136.7,-60.5)..(137.3,-60.3)
				..controls(138.4,-60)and(139.3,-59.7)..(141.1,-59)
				..controls(141.6,-58.7)and(142,-58.6)..(142.5,-58.6)
				..controls(143.2,-58.6)and(143.3,-58.7)..(143.6,-59)
				..controls(144,-59.5)and(144.1,-59.9)..(144,-60.8)
				..controls(144,-61.2)and(144,-71.1)..(144,-82.8)--(143.9,-104.1)--(143.7,-104.5)
				..controls(143.5,-104.7)and(143.3,-104.8)..(143,-104.9)
				..controls(142.6,-104.9)and(142.5,-104.8)..(141.5,-103.7)--cycle
				(27.8,-91.4)..controls(27.1,-90.9)and(26.3,-90)..(26.1,-89.5)
				..controls(26,-89.2)and(25.8,-88.6)..(25.7,-87.9)
				..controls(25.4,-86.2)and(25.6,-84.8)..(26.5,-83.5)
				..controls(27.4,-82.3)and(28.3,-82)..(29.4,-82.3)
				..controls(31.1,-82.8)and(32.1,-84.9)..(32,-87.4)
				..controls(31.9,-89.3)and(31.2,-90.6)..(29.9,-91.4)
				..controls(29.1,-91.9)and(28.6,-91.9)..(27.8,-91.4)--cycle
				(190.1,-91.3)..controls(189.7,-91.2)and(189.3,-90.9)..(189.1,-90.7)
				..controls(188.3,-89.8)and(187.6,-87.7)..(187.6,-86.2)
				..controls(187.6,-85.5)and(187.7,-85.2)..(188.2,-84.4)
				..controls(189.1,-82.6)and(190.1,-81.8)..(191.3,-82.1)
				..controls(192,-82.3)and(193.5,-83.7)..(193.8,-84.4)
				..controls(194.1,-85.1)and(194.2,-87.1)..(194,-88.3)
				..controls(193.8,-89.8)and(193.3,-90.5)..(192.3,-91.3)
				..controls(191.7,-91.7)and(190.9,-91.7)..(190.1,-91.3)--cycle
				(173.7,-83.4)..controls(172.2,-82.5)and(172,-82.5)..(170.6,-82.4)
				..controls(169.7,-82.4)and(168.8,-82.4)..(168.3,-82.5)
				..controls(167.2,-82.7)and(166.9,-82.7)..(166.6,-82.2)
				..controls(166.3,-81.8)and(166.3,-81)..(166.8,-80)
				..controls(167.2,-78.9)and(167.4,-77.4)..(167.7,-74.1)
				..controls(167.8,-72.8)and(167.9,-71.8)..(167.9,-71.7)
				..controls(168,-71.7)and(168,-70.2)..(168,-68.4)..controls(168,-65.4)and(168,-64.9)..(167.7,-63)
				..controls(167.3,-60.5)and(167.1,-59.9)..(166.4,-58.4)
				..controls(166.1,-57.7)and(165.8,-56.9)..(165.7,-56.6)
				..controls(165.7,-56.2)and(165.4,-55.6)..(165.2,-55.2)
				..controls(164.6,-53.8)and(161.9,-50.7)..(161,-50.3)
				..controls(159.8,-49.7)and(157.9,-50)..(155.2,-51.2)
				..controls(154.7,-51.4)and(153.4,-51.8)..(152.2,-52.1)
				..controls(150.8,-52.4)and(149.4,-52.9)..(148.6,-53.2)
				..controls(147.8,-53.6)and(146.7,-54)..(146,-54.2)
				..controls(145.4,-54.4)and(144.5,-54.7)..(144,-54.9)
				..controls(143.6,-55.2)and(142.7,-55.5)..(141.8,-55.7)
				..controls(141,-55.9)and(139.6,-56.3)..(138.7,-56.6)
				..controls(137.8,-57)and(136.3,-57.4)..(135.3,-57.6)
				..controls(133.8,-57.9)and(133.2,-58.1)..(132.1,-58.7)
				..controls(131.1,-59.2)and(130.3,-59.4)..(129.4,-59.6)
				..controls(128.7,-59.8)and(127.8,-60)..(127.4,-60.1)
				..controls(127,-60.3)and(126,-60.5)..(125.1,-60.7)
				..controls(123.2,-61.1)and(121.7,-61.5)..(120.2,-62.1)
				..controls(119.7,-62.3)and(118.1,-62.8)..(116.7,-63.2)
				..controls(114.6,-63.9)and(114,-64.1)..(113,-64.7)
				..controls(111.9,-65.4)and(111.8,-65.4)..(111,-65.5)
				..controls(109.9,-65.5)and(109.3,-65.4)..(107.7,-64.7)
				..controls(107,-64.4)and(106,-64)..(105.3,-63.9)
				..controls(103.5,-63.5)and(102.2,-63)..(101.2,-62.5)
				..controls(100.1,-61.8)and(99.3,-61.5)..(97.8,-61.1)
				..controls(97.2,-61)and(96.2,-60.7)..(95.7,-60.4)
				..controls(95.1,-60.1)and(94.2,-59.8)..(93.5,-59.6)
				..controls(92.9,-59.4)and(91.7,-59)..(91,-58.7)..controls(90.2,-58.4)and(88.9,-58)..(88,-57.7)
				..controls(87.2,-57.5)and(86.1,-57.1)..(85.5,-56.8)
				..controls(85,-56.6)and(84.1,-56.2)..(83.5,-56)..controls(82.9,-55.7)and(81.8,-55.2)..(81,-54.9)
				..controls(80.1,-54.5)and(79.1,-54.2)..(78.7,-54.1)
				..controls(77.9,-53.8)and(77.3,-53.6)..(76.2,-52.9)
				..controls(75.7,-52.6)and(75.2,-52.4)..(74.1,-52.2)
				..controls(72.5,-51.9)and(72.4,-51.8)..(71.3,-51.1)
				..controls(69.9,-50.1)and(69.4,-49.9)..(67.9,-49.8)
				..controls(67,-49.7)and(66.6,-49.6)..(66.3,-49.4)
				..controls(66.1,-49.3)and(65.7,-49)..(65.4,-48.8)
				..controls(65.1,-48.7)and(64.2,-48.4)..(63.6,-48.2)
				..controls(62.9,-48)and(61.7,-47.6)..(60.9,-47.2)
				..controls(60.1,-46.8)and(58.9,-46.4)..(58.2,-46.2)
				..controls(57.6,-46)and(56.8,-45.7)..(56.5,-45.5)
				..controls(54.9,-44.6)and(54.6,-44.5)..(53.2,-44.2)
				..controls(51.7,-43.8)and(51,-43.6)..(49.9,-42.9)
				..controls(49.5,-42.6)and(48.9,-42.4)..(47.5,-42)
				..controls(45.7,-41.6)and(45.6,-41.5)..(44.3,-40.6)
				..controls(43.5,-40.1)and(42.8,-39.6)..(42.7,-39.4)
				..controls(42.4,-38.7)and(43,-38.1)..(44.5,-37.6)..controls(44.9,-37.5)and(46,-37)..(46.9,-36.6)
				..controls(47.9,-36.1)and(48.8,-35.8)..(49.6,-35.6)
				..controls(50.5,-35.4)and(51.3,-35.2)..(52.4,-34.6)
				..controls(53.3,-34.1)and(54.5,-33.7)..(54.9,-33.5)
				..controls(55.4,-33.4)and(56.4,-33)..(57.2,-32.6)
				..controls(58.6,-32)and(60.1,-31.5)..(61.6,-31.2)
				..controls(62.3,-31.1)and(62.7,-31)..(63.2,-30.7)
				..controls(64.1,-30)and(64.3,-29.9)..(65.3,-29.7)
				..controls(65.7,-29.6)and(66.5,-29.2)..(67.1,-28.9)
				..controls(67.6,-28.6)and(68.8,-28.2)..(69.6,-27.9)
				..controls(70.5,-27.7)and(71.6,-27.3)..(72.1,-27)
				..controls(72.6,-26.7)and(73.8,-26.2)..(74.8,-25.9)
				..controls(75.7,-25.5)and(76.7,-25.1)..(77.1,-25)
				..controls(77.4,-24.8)and(78.5,-24.4)..(79.6,-24.1)
				..controls(80.6,-23.8)and(82,-23.2)..(82.8,-22.9)
				..controls(83.5,-22.5)and(84.8,-22)..(85.8,-21.8)
				..controls(86.7,-21.5)and(87.8,-21.1)..(88.1,-20.9)
				..controls(88.4,-20.7)and(89.3,-20.3)..(90.1,-20.1)
				..controls(90.8,-19.8)and(91.9,-19.4)..(92.4,-19.1)
				..controls(93.5,-18.5)and(94.4,-18.2)..(95.4,-18.1)
				..controls(95.7,-18)and(96.4,-17.8)..(96.9,-17.5)
				..controls(97.5,-17.3)and(98.4,-16.9)..(99,-16.6)
				..controls(99.7,-16.3)and(100.9,-15.9)..(101.9,-15.6)
				..controls(102.9,-15.3)and(104.4,-14.7)..(105.3,-14.3)
				..controls(108.1,-13.2)and(108.5,-13.1)..(110.1,-13.1)
				..controls(111.4,-13.1)and(111.5,-13.1)..(112.4,-13.5)
				..controls(113.2,-14)and(115,-14.8)..(116.6,-15.4)
				..controls(116.9,-15.5)and(117.9,-15.8)..(118.8,-16.1)
				..controls(120.9,-16.6)and(121.5,-16.9)..(122.6,-17.9)
				..controls(123,-18.3)and(123.7,-18.8)..(123.9,-18.9)
				..controls(124.6,-19.4)and(126,-19.8)..(127.9,-20.2)
				..controls(129.6,-20.5)and(129.8,-20.6)..(131.1,-21.4)
				..controls(131.5,-21.7)and(132.4,-22.1)..(133.1,-22.3)
				..controls(136.1,-23.2)and(138.3,-24.1)..(140.2,-25)
				..controls(141,-25.5)and(141.5,-25.6)..(142.5,-25.8)
				..controls(145.1,-26.3)and(146,-26.6)..(149.4,-28.1)
				..controls(151.3,-28.9)and(153.5,-29.8)..(154.3,-30.1)
				..controls(155.1,-30.4)and(156.2,-30.8)..(156.6,-31)
				..controls(157.1,-31.2)and(158.2,-31.6)..(159.1,-31.8)
				..controls(160.6,-32.2)and(160.9,-32.3)..(161.7,-32.9)
				..controls(162.4,-33.3)and(162.9,-33.6)..(163.7,-33.8)
				..controls(164.3,-34)and(164.8,-34.1)..(164.9,-34.1)
				..controls(165,-34.1)and(165.6,-34.4)..(166.4,-34.8)
				..controls(167.3,-35.2)and(168.1,-35.4)..(169,-35.6)
				..controls(170,-35.8)and(170.5,-36)..(171.7,-36.6)
				..controls(172.4,-37)and(173.3,-37.4)..(173.6,-37.4)
				..controls(175.6,-37.7)and(176.2,-38)..(176.7,-39.1)
				..controls(176.9,-39.6)and(176.9,-39.7)..(176.7,-40.4)
				..controls(176.4,-41.6)and(175.8,-42.7)..(175.2,-43.4)
				..controls(174.5,-44.1)and(172.6,-45.5)..(171.7,-45.8)
				..controls(170.6,-46.2)and(169.1,-46)..(168.2,-45.4)
				..controls(167.9,-45.2)and(167.3,-44.8)..(166.9,-44.4)
				..controls(166.4,-44.1)and(165.8,-43.6)..(165.4,-43.2)
				..controls(165.1,-42.9)and(164.4,-42.3)..(163.9,-41.9)
				..controls(163.3,-41.5)and(162.7,-40.9)..(162.4,-40.5)
				..controls(161.4,-39.2)and(159.9,-38)..(158.8,-37.8)
				..controls(158.4,-37.8)and(158.3,-37.7)..(158.2,-37.2)
				..controls(158,-36.8)and(157.4,-36.1)..(157.1,-36.1)
				..controls(157,-36.1)and(156.7,-36.1)..(156.5,-36)
				..controls(156.1,-35.9)and(155.9,-35.7)..(155.7,-35.2)
				..controls(155.3,-34.6)and(155.1,-34.4)..(154,-34.1)
				..controls(153.6,-34)and(153.2,-33.7)..(152.9,-33.5)
				..controls(152.4,-32.9)and(151.6,-32.4)..(151.3,-32.4)
				..controls(150.9,-32.4)and(150.4,-32.6)..(150.4,-32.9)
				..controls(150.4,-33.2)and(150.8,-33.4)..(151.4,-33.6)
				..controls(152.1,-33.7)and(152.8,-34.1)..(153.5,-34.9)
				..controls(154.1,-35.5)and(154.3,-35.6)..(155,-35.8)
				..controls(155.7,-36)and(155.7,-36)..(156,-36.6)
				..controls(156.3,-37.3)and(156.7,-37.7)..(157.5,-37.8)
				..controls(158.1,-37.9)and(158.2,-38)..(158.4,-38.8)
				..controls(158.5,-39.4)and(158.6,-39.6)..(159.6,-40.4)
				..controls(160.2,-41)and(160.9,-41.6)..(161.3,-41.8)
				..controls(161.7,-42)and(161.9,-42.3)..(162.2,-42.8)
				..controls(162.4,-43.2)and(162.8,-43.8)..(163.1,-44.1)
				..controls(163.9,-44.8)and(165.1,-46.7)..(165.7,-48)
				..controls(166.2,-49.3)and(166.4,-49.7)..(167.2,-50.5)
				..controls(167.5,-50.8)and(167.9,-51.4)..(168.1,-51.8)
				..controls(168.3,-52.2)and(168.8,-53)..(169.1,-53.6)
				..controls(169.6,-54.3)and(169.8,-54.9)..(169.9,-55.6)
				..controls(170.4,-57.8)and(170.5,-58.1)..(171.2,-59.5)
				..controls(171.6,-60.3)and(172,-61.5)..(172.2,-62.1)
				..controls(172.3,-62.8)and(172.7,-63.7)..(172.9,-64.1)
				..controls(173.2,-64.6)and(173.4,-65.3)..(173.4,-65.6)
				..controls(173.8,-68.8)and(174.2,-70.6)..(174.9,-72.6)
				..controls(175.5,-74.5)and(175.9,-77.1)..(176.1,-80.2)
				..controls(176.3,-83.4)and(175.5,-84.4)..(173.7,-83.4)--cycle;
				\definecolor{ezs}{RGB}{247,179,26}\fill[ezs](-5.7,-156.6)--(-3.4,-157.5)--(61.9,12.7)--(59.7,13.6)--cycle;
				\definecolor{ezs}{RGB}{247,179,22}\fill[ezs](-3.8,-157.3)--(-1.5,-158.2)--(63.9,12)--(61.6,12.9)--cycle;
				\definecolor{ezs}{RGB}{248,179,19}\fill[ezs](-1.8,-158.1)--(0.4,-158.9)--(65.8,11.2)--(63.5,12.1)--cycle;
				\definecolor{ezs}{RGB}{248,178,16}\fill[ezs](0.1,-158.8)--(2.4,-159.7)--(67.7,10.5)--(65.4,11.4)--cycle;
				\definecolor{ezs}{RGB}{249,178,15}\fill[ezs](2,-159.5)--(4.3,-160.4)--(69.6,9.8)--(67.3,10.7)--cycle;
				\definecolor{ezs}{RGB}{251,177,17}\fill[ezs](3.9,-160.3)--(6.2,-161.1)--(71.5,9)--(69.2,9.9)--cycle;
				\definecolor{ezs}{RGB}{252,179,18}\fill[ezs](5.8,-161)--(8.1,-161.9)--(73.4,8.3)--(71.1,9.2)--cycle;
				\definecolor{ezs}{RGB}{252,182,17}\fill[ezs](7.7,-161.7)--(10,-162.6)--(75.3,7.6)--(73,8.5)--cycle;
				\definecolor{ezs}{RGB}{252,185,17}\fill[ezs](9.6,-162.5)--(11.9,-163.4)--(77.2,6.8)--(75,7.7)--cycle;
				\definecolor{ezs}{RGB}{251,189,18}\fill[ezs](11.5,-163.2)--(13.8,-164.1)--(79.2,6.1)--(76.9,7)--cycle;
				\definecolor{ezs}{RGB}{251,191,22}\fill[ezs](13.4,-163.9)--(15.7,-164.8)--(81.1,5.4)--(78.8,6.2)--cycle;
				\definecolor{ezs}{RGB}{252,193,25}\fill[ezs](15.4,-164.7)--(17.6,-165.6)--(83,4.6)--(80.7,5.5)--cycle;
				\definecolor{ezs}{RGB}{253,195,27}\fill[ezs](17.3,-165.4)--(19.6,-166.3)--(84.9,3.9)--(82.6,4.8)--cycle;
				\definecolor{ezs}{RGB}{253,196,31}\fill[ezs](19.2,-166.1)--(21.5,-167)--(86.8,3.2)--(84.5,4)--cycle;
				\definecolor{ezs}{RGB}{254,199,35}\fill[ezs](21.1,-166.9)--(23.4,-167.8)--(88.7,2.4)--(86.4,3.3)--cycle;
				\definecolor{ezs}{RGB}{253,201,38}\fill[ezs](23,-167.6)--(25.3,-168.5)--(90.6,1.7)--(88.3,2.6)--cycle;
				\definecolor{ezs}{RGB}{253,203,40}\fill[ezs](24.9,-168.3)--(27.2,-169.2)--(92.5,1)--(90.3,1.8)--cycle;
				\definecolor{ezs}{RGB}{253,205,42}\fill[ezs](26.8,-169.1)--(29.1,-170)--(94.4,0.2)--(92.2,1.1)--cycle;
				\definecolor{ezs}{RGB}{253,207,45}\fill[ezs](28.7,-169.8)--(31,-170.7)--(96.4,-0.5)--(94.1,0.4)--cycle;
				\definecolor{ezs}{RGB}{253,210,49}\fill[ezs](30.7,-170.5)--(32.9,-171.4)--(98.3,-1.2)--(96,-0.4)--cycle;
				\definecolor{ezs}{RGB}{253,212,54}\fill[ezs](32.6,-171.3)--(34.9,-172.2)--(100.2,-2)--(97.9,-1.1)--cycle;
				\definecolor{ezs}{RGB}{253,214,59}\fill[ezs](34.5,-172)--(36.8,-172.9)--(102.1,-2.7)--(99.8,-1.8)--cycle;
				\definecolor{ezs}{RGB}{253,216,63}\fill[ezs](36.4,-172.7)--(38.7,-173.6)--(104,-3.4)--(101.7,-2.6)--cycle;
				\definecolor{ezs}{RGB}{253,218,65}\fill[ezs](38.3,-173.5)--(40.6,-174.4)--(105.9,-4.2)--(103.6,-3.3)--cycle;
				\definecolor{ezs}{RGB}{254,220,66}\fill[ezs](40.2,-174.2)--(42.5,-175.1)--(107.8,-4.9)--(105.5,-4)--cycle;
				\definecolor{ezs}{RGB}{254,222,68}\fill[ezs](42.1,-175)--(44.4,-175.8)--(109.7,-5.6)--(107.5,-4.8)--cycle;
				\definecolor{ezs}{RGB}{254,224,71}\fill[ezs](44,-175.7)--(46.3,-176.6)--(111.7,-6.4)--(109.4,-5.5)--cycle;
				\definecolor{ezs}{RGB}{254,226,75}\fill[ezs](46,-176.4)--(48.2,-177.3)--(113.6,-7.1)--(111.3,-6.2)--cycle;
				\definecolor{ezs}{RGB}{254,228,79}\fill[ezs](47.9,-177.2)--(50.1,-178)--(115.5,-7.8)--(113.2,-7)--cycle;
				\definecolor{ezs}{RGB}{254,230,82}\fill[ezs](49.8,-177.9)--(52.1,-178.8)--(117.4,-8.6)--(115.1,-7.7)--cycle;
				\definecolor{ezs}{RGB}{254,232,85}\fill[ezs](51.7,-178.6)--(54,-179.5)--(119.3,-9.3)--(117,-8.4)--cycle;
				\definecolor{ezs}{RGB}{254,234,88}\fill[ezs](53.6,-179.4)--(55.9,-180.2)--(121.2,-10)--(118.9,-9.2)--cycle;
				\definecolor{ezs}{RGB}{254,237,91}\fill[ezs](55.5,-180.1)--(57.8,-181)--(123.1,-10.8)--(120.8,-9.9)--cycle;
				\definecolor{ezs}{RGB}{254,239,94}\fill[ezs](57.4,-180.8)--(59.7,-181.7)--(125,-11.5)--(122.8,-10.6)--cycle;
				\definecolor{ezs}{RGB}{254,241,97}\fill[ezs](59.3,-181.6)--(61.6,-182.4)--(127,-12.2)--(124.7,-11.4)--cycle;
				\definecolor{ezs}{RGB}{254,243,101}\fill[ezs](61.2,-182.3)--(63.5,-183.2)--(128.9,-13)--(126.6,-12.1)--cycle;
				\definecolor{ezs}{RGB}{254,245,103}\fill[ezs](63.2,-183)--(65.4,-183.9)--(130.8,-13.7)--(128.5,-12.8)--cycle;
				\definecolor{ezs}{RGB}{254,245,105}\fill[ezs](65.1,-183.8)--(67.4,-184.6)--(132.7,-14.4)--(130.4,-13.6)--cycle;
				\definecolor{ezs}{RGB}{254,245,105}\fill[ezs](67,-184.5)--(69.3,-185.4)--(134.6,-15.2)--(132.3,-14.3)--cycle;
				\definecolor{ezs}{RGB}{254,244,103}\fill[ezs](68.9,-185.2)--(71.2,-186.1)--(136.5,-15.9)--(134.2,-15)--cycle;
				\definecolor{ezs}{RGB}{253,243,101}\fill[ezs](70.8,-186)--(73.1,-186.8)--(138.4,-16.6)--(136.1,-15.8)--cycle;
				\definecolor{ezs}{RGB}{253,241,98}\fill[ezs](72.7,-186.7)--(75,-187.6)--(140.3,-17.4)--(138,-16.5)--cycle;
				\definecolor{ezs}{RGB}{253,239,95}\fill[ezs](74.6,-187.4)--(76.9,-188.3)--(142.2,-18.1)--(140,-17.2)--cycle;
				\definecolor{ezs}{RGB}{253,236,92}\fill[ezs](76.5,-188.2)--(78.8,-189)--(144.2,-18.8)--(141.9,-18)--cycle;
				\definecolor{ezs}{RGB}{253,234,88}\fill[ezs](78.5,-188.9)--(80.7,-189.8)--(146.1,-19.6)--(143.8,-18.7)--cycle;
				\definecolor{ezs}{RGB}{253,232,85}\fill[ezs](80.4,-189.6)--(82.7,-190.5)--(148,-20.3)--(145.7,-19.4)--cycle;
				\definecolor{ezs}{RGB}{253,230,82}\fill[ezs](82.3,-190.4)--(84.6,-191.2)--(149.9,-21)--(147.6,-20.2)--cycle;
				\definecolor{ezs}{RGB}{253,228,78}\fill[ezs](84.2,-191.1)--(86.5,-192)--(151.8,-21.8)--(149.5,-20.9)--cycle;
				\definecolor{ezs}{RGB}{253,226,74}\fill[ezs](86.1,-191.8)--(88.4,-192.7)--(153.7,-22.5)--(151.4,-21.6)--cycle;
				\definecolor{ezs}{RGB}{253,224,70}\fill[ezs](88,-192.6)--(90.3,-193.4)--(155.6,-23.3)--(153.3,-22.4)--cycle;
				\definecolor{ezs}{RGB}{253,222,67}\fill[ezs](89.9,-193.3)--(92.2,-194.2)--(157.5,-24)--(155.3,-23.1)--cycle;
				\definecolor{ezs}{RGB}{253,220,65}\fill[ezs](91.8,-194)--(94.1,-194.9)--(159.5,-24.7)--(157.2,-23.8)--cycle;
				\definecolor{ezs}{RGB}{253,217,63}\fill[ezs](93.7,-194.8)--(96,-195.6)--(161.4,-25.5)--(159.1,-24.6)--cycle;
				\definecolor{ezs}{RGB}{253,215,60}\fill[ezs](95.7,-195.5)--(97.9,-196.4)--(163.3,-26.2)--(161,-25.3)--cycle;
				\definecolor{ezs}{RGB}{253,213,56}\fill[ezs](97.6,-196.2)--(99.9,-197.1)--(165.2,-26.9)--(162.9,-26)--cycle;
				\definecolor{ezs}{RGB}{252,211,53}\fill[ezs](99.5,-197)--(101.8,-197.8)--(167.1,-27.7)--(164.8,-26.8)--cycle;
				\definecolor{ezs}{RGB}{252,209,49}\fill[ezs](101.4,-197.7)--(103.7,-198.6)--(169,-28.4)--(166.7,-27.5)--cycle;
				\definecolor{ezs}{RGB}{252,207,45}\fill[ezs](103.3,-198.4)--(105.6,-199.3)--(170.9,-29.1)--(168.6,-28.2)--cycle;
				\definecolor{ezs}{RGB}{253,205,42}\fill[ezs](105.2,-199.2)--(107.5,-200)--(172.8,-29.9)--(170.6,-29)--cycle;
				\definecolor{ezs}{RGB}{253,202,40}\fill[ezs](107.1,-199.9)--(109.4,-200.8)--(174.7,-30.6)--(172.5,-29.7)--cycle;
				\definecolor{ezs}{RGB}{253,200,38}\fill[ezs](109,-200.6)--(111.3,-201.5)--(176.7,-31.3)--(174.4,-30.4)--cycle;
				\definecolor{ezs}{RGB}{253,198,35}\fill[ezs](111,-201.4)--(113.2,-202.2)--(178.6,-32.1)--(176.3,-31.2)--cycle;
				\definecolor{ezs}{RGB}{253,196,31}\fill[ezs](112.9,-202.1)--(115.2,-203)--(180.5,-32.8)--(178.2,-31.9)--cycle;
				\definecolor{ezs}{RGB}{253,194,28}\fill[ezs](114.8,-202.8)--(117.1,-203.7)--(182.4,-33.5)--(180.1,-32.6)--cycle;
				\definecolor{ezs}{RGB}{253,192,26}\fill[ezs](116.7,-203.6)--(119,-204.5)--(184.3,-34.3)--(182,-33.4)--cycle;
				\definecolor{ezs}{RGB}{253,189,24}\fill[ezs](118.6,-204.3)--(120.9,-205.2)--(186.2,-35)--(183.9,-34.1)--cycle;
				\definecolor{ezs}{RGB}{253,186,20}\fill[ezs](120.5,-205)--(122.8,-205.9)--(188.1,-35.7)--(185.8,-34.8)--cycle;
				\definecolor{ezs}{RGB}{253,184,17}\fill[ezs](122.4,-205.8)--(124.7,-206.7)--(190,-36.5)--(187.8,-35.6)--cycle;
				\definecolor{ezs}{RGB}{252,183,15}\fill[ezs](124.3,-206.5)--(126.6,-207.4)--(192,-37.2)--(189.7,-36.3)--cycle;
				\definecolor{ezs}{RGB}{252,181,14}\fill[ezs](126.3,-207.2)--(128.5,-208.1)--(193.9,-37.9)--(191.6,-37.1)--cycle;
				\definecolor{ezs}{RGB}{252,179,13}\fill[ezs](128.2,-208)--(130.4,-208.9)--(195.8,-38.7)--(193.5,-37.8)--cycle;
				\definecolor{ezs}{RGB}{252,178,12}\fill[ezs](130.1,-208.7)--(132.4,-209.6)--(197.7,-39.4)--(195.4,-38.5)--cycle;
				\definecolor{ezs}{RGB}{251,178,10}\fill[ezs](132,-209.4)--(134.3,-210.3)--(199.6,-40.1)--(197.3,-39.3)--cycle;
				\definecolor{ezs}{RGB}{251,178,10}\fill[ezs](133.9,-210.2)--(136.2,-211.1)--(201.5,-40.9)--(199.2,-40)--cycle;
				\definecolor{ezs}{RGB}{250,177,11}\fill[ezs](135.8,-210.9)--(138.1,-211.8)--(203.4,-41.6)--(201.1,-40.7)--cycle;
				\definecolor{ezs}{RGB}{250,178,11}\fill[ezs](137.7,-211.6)--(140,-212.5)--(205.3,-42.3)--(203.1,-41.5)--cycle;
				\definecolor{ezs}{RGB}{248,180,10}\fill[ezs](139.6,-212.4)--(141.9,-213.3)--(207.3,-43.1)--(205,-42.2)--cycle;
				\definecolor{ezs}{RGB}{247,181,9}\fill[ezs](141.5,-213.1)--(143.8,-214)--(209.2,-43.8)--(206.9,-42.9)--cycle;
				\definecolor{ezs}{RGB}{246,182,9}\fill[ezs](143.5,-213.8)--(145.7,-214.7)--(211.1,-44.5)--(208.8,-43.7)--cycle;
				\definecolor{ezs}{RGB}{246,182,10}\fill[ezs](145.4,-214.6)--(147.7,-215.5)--(213,-45.3)--(210.7,-44.4)--cycle;
			\end{scope}
			% chữ EDUZY (Times đậm, làm dày nét bằng 5 lớp lệch nhẹ; căn đúng khung chữ logo gốc)
			\foreach \dx/\dy in {-0.7/0,0.7/0,0/-0.3,0/0.3,0/0}
			\node[anchor=center,inner sep=0pt,outer sep=0pt,text=ezText,align=center] at (109.6+\dx,-193.6+\dy)
			{\resizebox{1.01cm}{0.22cm}{\fontencoding{T1}\fontfamily{ptm}\fontseries{b}\selectfont EDUZY}};
		\end{tikzpicture}%
	}%
}
%====================== THIẾT LẬP THÂN ĐỀ ======================
%--- Môi trường tự luận {bt}: ex_test 3.3.4 chưa khai báo sẵn nên tạo ở đây
%    (ex_test đã tự móc \fixshowans và \listenumerate cho bt)
\makeatletter
\@ifundefined{bt}{\newtheorem{bt}{\color{gold1}\namebt}}{}
\makeatother
\renewcommand{\sectionmark}[1]{\markright{\thesection\ #1}}
%--- Khung tiêu đề (logo + tối đa 4 dòng; dòng 3, 4 để trống {} thì bỏ qua)
%  Đề thi:    \khungtieude{THI THỬ THPT QUỐC GIA NĂM 2026}{Môn Toán -- Đề số 21}
%                         {(Đề thi gồm: ...)}{Thời gian làm bài: 90 phút (không kể thời gian phát đề).}
%  Chuyên đề: \khungtieude{CHUYÊN ĐỀ 3 -- DÃY SỐ VÀ CÔNG THỨC TRUY HỒI}
%                         {Bồi dưỡng học sinh giỏi Toán THPT}{...}{}
\newcommand{\khungtieude}[4]{%
	\begin{tcolorbox}[colframe=gold1,colback=gold5!45,boxrule=0.5mm]
		\begin{minipage}[c]{3cm}
			\centerline{\eduzylogo[2.4cm]}
		\end{minipage}\hspace{0.5cm}
		\begin{minipage}[c]{12.5cm}
			\vskip 0.15 cm
			\centerline{\ind{#1}}
			\vskip 0.15 cm
			\centerline{\textbf{#2}}
			\if\relax\detokenize{#3}\relax\else
				\vskip 0.15 cm
				\centerline{\textit{#3}}
			\fi
			\if\relax\detokenize{#4}\relax\else
				\vskip 0.15 cm
				\centerline{\textit{#4}}
			\fi
			\vskip 0.15 cm
		\end{minipage}
	\end{tcolorbox}
}
%====================== BỔ SUNG CHO BÀI HỌC, CHUYÊN ĐỀ ======================
\usepackage[shortlabels]{enumitem}
\usepackage{tabularx}
%--- Danh sách gọn: \begin{enumerate}[gon] ... \end{enumerate}
\SetEnumitemKey{gon}{leftmargin=*,itemsep=1pt,parsep=0pt,topsep=3pt}
%--- Khung lý thuyết: \begin{khung}{} ... \end{khung}; tiêu đề chỉ dùng khi thật ngắn ("Mẹo", "Lỗi thường gặp")
\newtcolorbox{khung}[1]{breakable,enhanced,colframe=gold1,colback=gold5!45,
	boxrule=0.4mm,fonttitle=\bfseries\sffamily,coltitle=white,colbacktitle=gold1,
	attach boxed title to top left={xshift=6mm,yshift=-3mm},
	boxed title style={arc=1mm,boxrule=0.3mm},title={#1}}
%--- Tiêu đề mục nhỏ trong một dạng: \muc{Dấu hiệu}, \muc{Phương pháp}, \muc{Bài tự luyện}
\makeatletter
\newcommand{\muc}[1]{\par\medskip\noindent\indam{#1}\par\nobreak\smallskip\@afterindentfalse\@afterheading}
\makeatother
%--- Ô công thức cần nhớ (căn giữa, viền vàng, vừa khít nội dung)
%     \begin{center}\congthuc{$Q(n)=(n-1)!$}\end{center}
\newtcbox{\congthuc}{on line,colframe=gold1,colback=gold5!45,boxrule=0.5pt,arc=2pt,
	left=8pt,right=8pt,top=4pt,bottom=4pt}
%--- Công tắc lời giải bài tự luyện (ví dụ mẫu luôn dùng \loigiai)
%     \loigiaiTL[{đáp số}]{lời giải đầy đủ}
%     \TLfull -> in lời giải đầy đủ;  \TLdapso -> chỉ in một dòng "Đáp số";  \TLan -> ẩn hoàn toàn
%     (\CheDo là dethi thì mọi lời giải và đáp số đều ẩn)
\makeatletter
\newif\ifTLhien   \TLhientrue
\newif\ifTLdapso  \TLdapsofalse
\newcommand{\TLfull}{\TLhientrue\TLdapsofalse}
\newcommand{\TLdapso}{\TLhientrue\TLdapsotrue}
\newcommand{\TLan}{\TLhienfalse}
\newcommand{\loigiaiTL}[2][]{%
	\ifTLhien
		\ifTLdapso
			\def\@tlarg{#1}%
			\ifx\@tlarg\@empty\else
				\ifdefstring{\kindOPTN}{dethi}{}{\par\noindent{\bfseries\color{gold1}Đáp số:} #1\par}%
			\fi
		\else
			\loigiai{#2}%
		\fi
	\fi}
\makeatother
\TLdapso% mặc định: chỉ in đáp số bài tự luyện
%====================== BẢNG ĐÚNG/SAI (\choiceTFt) ======================
%  Hàng tiêu đề nền nâu đồng chữ trắng (Phát biểu | Đ | S), các hàng tô xen kẽ kem / trắng,
%  đường kẻ vàng kim, nhãn a) b) c) d) màu nâu đồng. Bản có đáp án: cột được chọn có vòng tròn
%  nâu đồng với dấu tích trắng (câu đúng tích ở cột Đúng, câu sai tích ở cột Sai).
%  Tiêu đề ghi đủ "Đúng | Sai": \OPTN{viettat=0} sau \input (cột Đúng/Sai rộng ra, phát biểu hẹp lại).
%  Dùng bảng gốc của ex_test: \def\BangDSgoc{1} trước \input{preamble-eduzy-v2}.
\providecommand{\BangDSgoc}{0}
\ifnum\BangDSgoc=0
\usepackage{colortbl,pifont}
\colorlet{ezDSke}{gold2!60!gold1}% đường kẻ bảng (nét 0,4pt như ex_test để không đổi chiều cao bảng)
\colorlet{ezDSle}{gold5!60}% hàng a), c)
\colorlet{ezDSchan}{white}% hàng b), d)
%--- dấu chọn ở cột Đúng/Sai
%    (bề rộng 0, căn giữa ô: không làm đổi bề rộng cột của ex_test, không báo tràn lề)
\newcommand{\EZdauchon}{\makebox[0pt]{\tikz[baseline=-1ex]\node[circle,fill=gold1,inner sep=0pt,
	minimum size=1.2em,text=white,font=\small]{\ding{51}};}}
\def\tickT{\EZdauchon}\def\tickF{\EZdauchon}
\AtBeginDocument{\sbox{\ccTFTab}{\EZdauchon}}
%--- nhãn a) b) c) d) màu nâu đồng (ở bản lời giải, nhãn câu đúng vẫn đóng khung như ex_test)
\renewcommand{\FalseTF}{\stepcounter{dapan}\textbf{\color{gold1}\DapAnTF\sepTF}\ignorespaces}
%--- trước khi dựng bảng: kẻ màu vàng kim (bề rộng cột giữ nguyên như ex_test); sau bảng trả lại màu đen
\newcount\EZhangTF
\newcommand{\EZbatdauTF}{%
	\global\EZhangTF=0\relax
	\arrayrulecolor{ezDSke}}
\newcommand{\EZketthucTF}{\arrayrulecolor{black}}
%--- hàng tiêu đề nền nâu đồng, chữ trắng
\makeatletter
\newcommand{\EZo@TF}[1]{\centering\bfseries\color{white}#1}
\newcommand{\EZtieudeTF}{%
	\ifcase\cotDS
		\ifnum\cotSTT=1\relax
			\def\rowHeader{\hline\rowcolor{gold1}\EZo@TF{TT}&\EZo@TF{\phatbieu}\tabularnewline}%
		\else
			\def\rowHeader{\hline\rowcolor{gold1}\EZo@TF{\phatbieu}\tabularnewline}%
		\fi
	\or
		\ifnum\cotSTT=1\relax
			\def\rowHeader{\hline\rowcolor{gold1}\EZo@TF{TT}&\EZo@TF{\phatbieu}&\EZo@TF{\Dung}\tabularnewline}%
		\else
			\def\rowHeader{\hline\rowcolor{gold1}\EZo@TF{\phatbieu}&\EZo@TF{\Dung}\tabularnewline}%
		\fi
	\or
		\ifnum\cotSTT=1\relax
			\def\rowHeader{\hline\rowcolor{gold1}\EZo@TF{TT}&\EZo@TF{\phatbieu}&\EZo@TF{\Dung}&\EZo@TF{\Sai}\tabularnewline}%
		\else
			\def\rowHeader{\hline\rowcolor{gold1}\EZo@TF{\phatbieu}&\EZo@TF{\Dung}&\EZo@TF{\Sai}\tabularnewline}%
		\fi
	\fi}
%--- tô xen kẽ từng hàng (gọi ngay sau \hline đầu mỗi hàng)
\newcommand{\EZmauhangTF}{%
	\noalign{\global\advance\EZhangTF\@ne}%
	\ifodd\EZhangTF\rowcolor{ezDSle}\else\rowcolor{ezDSchan}\fi}
%--- gắn vào \choiceTFt / \writeChoiceTFt của ex_test
\xpretocmd{\choiceTFt}{\EZbatdauTF}{}{\typeout{[preamble-eduzy-v2] Không gắn được bảng Đúng/Sai (đầu bảng)}}
\xpatchcmd{\choiceTFt}{\ifnum\dongPB=0\def\rowHeader{}\fi}{\EZtieudeTF\ifnum\dongPB=0\def\rowHeader{}\fi}
	{}{\typeout{[preamble-eduzy-v2] Không gắn được bảng Đúng/Sai (hàng tiêu đề)}}
\xpatchcmd{\writeChoiceTFt}{\hline\xdef\dotEX{\sdotEX}}{\hline\EZmauhangTF\xdef\dotEX{\sdotEX}}
	{}{\typeout{[preamble-eduzy-v2] Không gắn được bảng Đúng/Sai (tô hàng)}}
\xpatchcmd{\writeChoiceTFt}{{\end{longtable}}}{{\end{longtable}\EZketthucTF}}
	{}{\typeout{[preamble-eduzy-v2] Không gắn được bảng Đúng/Sai (cuối longtable)}}
\xpatchcmd{\writeChoiceTFt}{{\end{xtabular}\end{center}\medskip\noindent}}{{\end{xtabular}\end{center}\EZketthucTF\medskip\noindent}}
	{}{\typeout{[preamble-eduzy-v2] Không gắn được bảng Đúng/Sai (cuối xtabular)}}
\fi
%=====================================================================
\renewcommand{\parallel}{\mathrel{/\!\!/}}
.
\providecommand{\BangDSgoc}{0}
\ifnum\BangDSgoc=0
\usepackage{colortbl,pifont}
\colorlet{ezDSke}{gold2!60!gold1}% đường kẻ bảng (nét 0,4pt như ex_test để không đổi chiều cao bảng)
\colorlet{ezDSle}{gold5!60}% hàng a), c)
\colorlet{ezDSchan}{white}% hàng b), d)
%--- dấu chọn ở cột Đúng/Sai
%    (bề rộng 0, căn giữa ô: không làm đổi bề rộng cột của ex_test, không báo tràn lề)
\newcommand{\EZdauchon}{\makebox[0pt]{\tikz[baseline=-1ex]\node[circle,fill=gold1,inner sep=0pt,
	minimum size=1.2em,text=white,font=\small]{\ding{51}};}}
\def\tickT{\EZdauchon}\def\tickF{\EZdauchon}
\AtBeginDocument{\sbox{\ccTFTab}{\EZdauchon}}
%--- nhãn a) b) c) d) màu nâu đồng (ở bản lời giải, nhãn câu đúng vẫn đóng khung như ex_test)
\renewcommand{\FalseTF}{\stepcounter{dapan}\textbf{\color{gold1}\DapAnTF\sepTF}\ignorespaces}
%--- trước khi dựng bảng: kẻ màu vàng kim (bề rộng cột giữ nguyên như ex_test); sau bảng trả lại màu đen
\newcount\EZhangTF
\newcommand{\EZbatdauTF}{%
	\global\EZhangTF=0\relax
	\arrayrulecolor{ezDSke}}
\newcommand{\EZketthucTF}{\arrayrulecolor{black}}
%--- hàng tiêu đề nền nâu đồng, chữ trắng
\makeatletter
\newcommand{\EZo@TF}[1]{\centering\bfseries\color{white}#1}
\newcommand{\EZtieudeTF}{%
	\ifcase\cotDS
		\ifnum\cotSTT=1\relax
			\def\rowHeader{\hline\rowcolor{gold1}\EZo@TF{TT}&\EZo@TF{\phatbieu}\tabularnewline}%
		\else
			\def\rowHeader{\hline\rowcolor{gold1}\EZo@TF{\phatbieu}\tabularnewline}%
		\fi
	\or
		\ifnum\cotSTT=1\relax
			\def\rowHeader{\hline\rowcolor{gold1}\EZo@TF{TT}&\EZo@TF{\phatbieu}&\EZo@TF{\Dung}\tabularnewline}%
		\else
			\def\rowHeader{\hline\rowcolor{gold1}\EZo@TF{\phatbieu}&\EZo@TF{\Dung}\tabularnewline}%
		\fi
	\or
		\ifnum\cotSTT=1\relax
			\def\rowHeader{\hline\rowcolor{gold1}\EZo@TF{TT}&\EZo@TF{\phatbieu}&\EZo@TF{\Dung}&\EZo@TF{\Sai}\tabularnewline}%
		\else
			\def\rowHeader{\hline\rowcolor{gold1}\EZo@TF{\phatbieu}&\EZo@TF{\Dung}&\EZo@TF{\Sai}\tabularnewline}%
		\fi
	\fi}
%--- tô xen kẽ từng hàng (gọi ngay sau \hline đầu mỗi hàng)
\newcommand{\EZmauhangTF}{%
	\noalign{\global\advance\EZhangTF\@ne}%
	\ifodd\EZhangTF\rowcolor{ezDSle}\else\rowcolor{ezDSchan}\fi}
%--- gắn vào \choiceTFt / \writeChoiceTFt của ex_test
\xpretocmd{\choiceTFt}{\EZbatdauTF}{}{\typeout{[preamble-eduzy-v2] Không gắn được bảng Đúng/Sai (đầu bảng)}}
\xpatchcmd{\choiceTFt}{\ifnum\dongPB=0\def\rowHeader{}\fi}{\EZtieudeTF\ifnum\dongPB=0\def\rowHeader{}\fi}
	{}{\typeout{[preamble-eduzy-v2] Không gắn được bảng Đúng/Sai (hàng tiêu đề)}}
\xpatchcmd{\writeChoiceTFt}{\hline\xdef\dotEX{\sdotEX}}{\hline\EZmauhangTF\xdef\dotEX{\sdotEX}}
	{}{\typeout{[preamble-eduzy-v2] Không gắn được bảng Đúng/Sai (tô hàng)}}
\xpatchcmd{\writeChoiceTFt}{{\end{longtable}}}{{\end{longtable}\EZketthucTF}}
	{}{\typeout{[preamble-eduzy-v2] Không gắn được bảng Đúng/Sai (cuối longtable)}}
\xpatchcmd{\writeChoiceTFt}{{\end{xtabular}\end{center}\medskip\noindent}}{{\end{xtabular}\end{center}\EZketthucTF\medskip\noindent}}
	{}{\typeout{[preamble-eduzy-v2] Không gắn được bảng Đúng/Sai (cuối xtabular)}}
\fi
%=====================================================================
\renewcommand{\parallel}{\mathrel{/\!\!/}}

%  trong file đề (tools/dan_trang/dan_trang.py tự ghi khi đặt \NeedspaceCau):
%    \def\TachCauDS{1}   câu Đúng/Sai được tách: câu dẫn (kèm hình) ở cuối trang, bảng Phát biểu sang trang sau
%                        (câu dẫn và hình không tách nhau, bảng không tách giữa hai trang)
%    \def\TachCauLG{1}   bản lời giải (\CheDo khác dethi): phần đề của câu cao hơn 1/3 trang được phép tách
%  Lệnh dùng trong file đề (giá trị do tools/dan_trang/dan_trang.py đo ở bản dethi):
%    \NeedspaceCau[<cao câu dẫn + hình>]{<cao cả khối câu>}   ngay trên \begin{ex} (câu đầu phần: trên tiêu đề phần)
%        công tắc tắt: đúng bằng \Needspace{<cao cả khối câu>} như cũ;
%        \TachCauDS=1 và có [..] (câu Đúng/Sai được chọn tách): chỉ giữ câu dẫn + hình;
%        \TachCauLG=1 ở bản lời giải: đòi tối đa 1/3 \textheight.
%    \NeedspaceBang{<cao bảng Phát biểu>}%   ngay trước \choiceTFt: \TachCauDS=1 thì giữ trọn bảng trên một trang;
%        công tắc tắt: không làm gì.
\providecommand{\TachCauDS}{0}
\providecommand{\TachCauLG}{0}
\newlength{\EZnsCau}
\newcommand{\NeedspaceCau}[2][]{%
	\setlength{\EZnsCau}{#2}%
	\ifnum\TachCauDS=1 \if\relax\detokenize{#1}\relax\else\setlength{\EZnsCau}{#1}\fi\fi
	\ifnum\TachCauLG=1 \ifdefstring{\kindOPTN}{dethi}{}{\ifdim\EZnsCau>0.3333\textheight\setlength{\EZnsCau}{0.3333\textheight}\fi}\fi
	\Needspace{\EZnsCau}}
\newcommand{\NeedspaceBang}[1]{\ifnum\TachCauDS=1 \Needspace{#1}\fi}

%--- Thang màu: đậm -> nhạt
\definecolor{gold1}{RGB}{126,96,0}    % nâu đồng đậm  #7E6000
\definecolor{gold2}{RGB}{192,144,3}   % vàng kim      #C09003
\definecolor{gold3}{RGB}{255,215,99}  % vàng nhạt     #FFD763
\definecolor{gold4}{RGB}{255,231,153} % vàng rất nhạt #FFE799
\definecolor{gold5}{RGB}{255,243,205} % vàng kem      #FFF3CD
%--- Ghi đè tông xanh/đỏ cũ
\definecolor{myblue}{RGB}{126,96,0}
\colorlet{blue}{gold1}
\colorlet{cyan}{gold2}
\colorlet{teal}{gold2}

%--- Tên chương mặc định (file đề nên tự khai báo lại)
\providecommand{\tenchuong}{}
%--- Chân trang
\providecommand{\thuonghieu}{TOÁN 12 -- Eduzy.com.vn}
\providecommand{\tacgia}{Tô Gia Bảo -- 0353703783}

%--- Chữ cạnh logo ở đầu trang
\providecommand{\khauhieu}{Eduzy.com.vn}

%====================== THÔNG SỐ ĐẦU TRANG, CHÂN TRANG (v2) ======================
%  Dải nền đậm tràn ngang trang; ô chữ vàng kem có cạnh xiên, kèm 2 vạch xiên chuyển màu.
%  Đầu trang lẻ : logo + \khauhieu (trái)   | ô chữ \tenchuong (phải)
%  Chân trang lẻ: ô chữ \thuonghieu (trái)  | Trang (số) (phải)
%  Trang chẵn soi gương trang lẻ; chân trang chẵn ghi \tacgia.
%--- Màu: chỉ dùng thang vàng kim Eduzy
\colorlet{ezNen}{gold1}% dải nền đậm
\colorlet{ezVach1}{gold3}% vạch xiên sát dải nền (đậm hơn)
\colorlet{ezVach2}{gold4}% vạch xiên sát ô chữ (nhạt hơn)
\colorlet{ezO}{gold5}% ô chữ vàng kem
%--- Kích thước
\newlength{\EZdai}\setlength{\EZdai}{1.1cm}% chiều cao dải nền
\newlength{\EZo}\setlength{\EZo}{0.68cm}% chiều cao ô chữ
\newlength{\EZxien}\setlength{\EZxien}{0.55cm}% độ lệch ngang của cạnh xiên
\newlength{\EZvach}\setlength{\EZvach}{0.75cm}% bề rộng mỗi vạch xiên
\newlength{\EZkhe}\setlength{\EZkhe}{0.12cm}% khe trắng giữa vạch và ô chữ
\newlength{\EZdem}\setlength{\EZdem}{0.4cm}% khoảng đệm từ chữ tới cạnh xiên
\newlength{\EZle}\setlength{\EZle}{1.8cm}% lề trái/phải của chữ (bằng lề thân bài)
\newlength{\EZrongmin}\setlength{\EZrongmin}{4.5cm}% ô chữ rộng tối thiểu
\newlength{\EZrong}% bề rộng chữ trong ô (tính tự động)
%--- Font
\newcommand{\EZfontO}{\fontsize{10.5pt}{12pt}\selectfont\bfseries}% chữ trong ô
\newcommand{\EZfontKH}{\fontfamily{qag}\fontsize{11pt}{12pt}\selectfont\bfseries}% \khauhieu cạnh logo
\newcommand{\EZfontTrang}{\fontsize{11.5pt}{12pt}\selectfont\bfseries}% chữ "Trang" và số trang
%--- Đo bề rộng chữ trong ô
\newcommand{\EZdorong}[1]{%
	\settowidth{\EZrong}{\EZfontO #1}%
	\ifdim\EZrong<\EZrongmin\setlength{\EZrong}{\EZrongmin}\fi}
%--- Logo + chữ trắng \khauhieu (dựng sẵn một lần khi vào \begin{document})
\newsavebox{\EZlogobox}
\AtBeginDocument{\sbox{\EZlogobox}{%
		\raisebox{-0.5\height}{\eduzylogo[1cm]}\hspace{0.15cm}%
		\raisebox{-0.5\height}{\color{white}\EZfontKH\khauhieu}}}
%--- Hai vạch xiên + khe trắng. #1: điểm gốc ở mép dưới, #2: hướng (1 sang phải, -1 sang trái),
%    #3: hướng nghiêng (1 kiểu "/", -1 kiểu "\")
\newcommand{\EZvachxien}[3]{%
	\fill[white] (#1) --++({(#3)*\EZxien},\EZo) --++({(#2)*\EZkhe},0) --++({-(#3)*\EZxien},-\EZo) --cycle;
	\fill[ezVach2] ([xshift=(#2)*\EZkhe]#1) --++({(#3)*\EZxien},\EZo) --++({(#2)*\EZvach},0) --++({-(#3)*\EZxien},-\EZo) --cycle;
	\fill[ezVach1] ([xshift=(#2)*(\EZkhe+\EZvach)]#1) --++({(#3)*\EZxien},\EZo) --++({(#2)*\EZvach},0) --++({-(#3)*\EZxien},-\EZo) --cycle;
	\draw[white,line width=0.6pt] ([xshift=(#2)*(\EZkhe+\EZvach)]#1) --++({(#3)*\EZxien},\EZo);
}
%--- Chữ trong ô (căn giữa theo chiều cao ô, không lệch vì dấu tiếng Việt)
\tikzset{ezchu/.style={inner sep=0pt,text=ezNen,text height=0.72em,text depth=0.22em}}

%====================== LOGO CHÌM GIỮA TRANG ======================
%  Vẽ cùng đầu trang nên nằm dưới chữ của thân bài, có trên mọi trang.
%  Tắt: \logochimtat (từ trang đó trở đi), bật lại: \logochimbat.
%  Đổi cỡ, độ đậm: khai báo trước \input, ví dụ \def\EZchimcao{9cm} \def\EZchimdo{0.1}
\providecommand{\EZchimcao}{10cm}% chiều cao logo chìm
\providecommand{\EZchimdo}{0.12}% độ đậm (0 = trong suốt, 1 = rõ hẳn)
\newif\ifEZchim \EZchimtrue
\newcommand{\logochimtat}{\EZchimfalse}
\newcommand{\logochimbat}{\EZchimtrue}
%  Dựng sẵn logo mờ một lần khi vào \begin{document}; transparency group làm mờ cả logo
%  một lượt nên không lộ vệt ở chỗ các dải màu chồng nhau.
\newsavebox{\EZchimtho}
\newsavebox{\EZchimbox}
\AtBeginDocument{%
	\sbox{\EZchimtho}{\eduzylogo[\EZchimcao]}%
	\sbox{\EZchimbox}{\begin{tikzpicture}
			\begin{scope}[transparency group,opacity=\EZchimdo]
				\node[inner sep=0pt] {\usebox{\EZchimtho}};
			\end{scope}
	\end{tikzpicture}}}
\newcommand{\EZvechim}{%
	\ifEZchim
		\node[inner sep=0pt] at (current page.center) {\usebox{\EZchimbox}};
	\fi}

%====================== ĐẦU TRANG - TRANG LẺ ======================
%  (logo bên trái, ô chữ \tenchuong bên phải)
\newcommand{\myfancyheadOdd}{%
	\EZdorong{\tenchuong}%
	\begin{tikzpicture}[remember picture,overlay]
		\coordinate (T) at (current page.north west);
		\coordinate (P) at (current page.north east);
		%--- logo chìm giữa trang
		\EZvechim
		%--- dải nền
		\fill[ezNen] (T) rectangle ([yshift=-\EZdai]P);
		%--- ô chữ: X là góc dưới bên trái, cạnh trái xiên kiểu "/"
		\coordinate (X) at ([xshift=-\EZle-\EZrong-\EZdem-\EZxien,yshift=-\EZdai]P);
		\fill[ezO] (X) --++(\EZxien,\EZo) -- ([yshift=-\EZdai+\EZo]P) -- ([yshift=-\EZdai]P) --cycle;
		\EZvachxien{X}{-1}{1}
		\node[ezchu,anchor=east] at ([xshift=-\EZle,yshift=-\EZdai+0.5*\EZo]P) {\EZfontO\tenchuong};
		%--- logo + tên
		\node[inner sep=0pt,anchor=west] at ([xshift=\EZle,yshift=-0.5*\EZdai]T) {\usebox{\EZlogobox}};
	\end{tikzpicture}%
}

%====================== ĐẦU TRANG - TRANG CHẴN =====================
%  (soi gương: ô chữ \tenchuong bên trái, logo bên phải)
\newcommand{\myfancyheadEven}{%
	\EZdorong{\tenchuong}%
	\begin{tikzpicture}[remember picture,overlay]
		\coordinate (T) at (current page.north west);
		\coordinate (P) at (current page.north east);
		%--- logo chìm giữa trang
		\EZvechim
		%--- dải nền
		\fill[ezNen] (T) rectangle ([yshift=-\EZdai]P);
		%--- ô chữ: X là góc dưới bên phải, cạnh phải xiên kiểu "\"
		\coordinate (X) at ([xshift=\EZle+\EZrong+\EZdem+\EZxien,yshift=-\EZdai]T);
		\fill[ezO] (X) --++(-\EZxien,\EZo) -- ([yshift=-\EZdai+\EZo]T) -- ([yshift=-\EZdai]T) --cycle;
		\EZvachxien{X}{1}{-1}
		\node[ezchu,anchor=west] at ([xshift=\EZle,yshift=-\EZdai+0.5*\EZo]T) {\EZfontO\tenchuong};
		%--- logo + tên
		\node[inner sep=0pt,anchor=east] at ([xshift=-\EZle,yshift=-0.5*\EZdai]P) {\usebox{\EZlogobox}};
	\end{tikzpicture}%
}

%====================== CHÂN TRANG - TRANG LẺ ======================
%  (ô chữ \thuonghieu bên trái, "Trang (số)" bên phải)
\newcommand{\myfancyfootOdd}{%
	\EZdorong{\thuonghieu}%
	\begin{tikzpicture}[remember picture,overlay]
		\coordinate (T) at (current page.south west);
		\coordinate (P) at (current page.south east);
		%--- dải nền
		\fill[ezNen] (T) rectangle ([yshift=\EZdai]P);
		%--- ô chữ sát mép trên dải: X là góc dưới bên phải, cạnh phải xiên kiểu "/"
		\coordinate (X) at ([xshift=\EZle+\EZrong+\EZdem,yshift=\EZdai-\EZo]T);
		\fill[ezO] ([yshift=\EZdai-\EZo]T) -- (X) --++(\EZxien,\EZo) -- ([yshift=\EZdai]T) --cycle;
		\EZvachxien{X}{1}{1}
		\node[ezchu,anchor=west] at ([xshift=\EZle,yshift=\EZdai-0.5*\EZo]T) {\EZfontO\thuonghieu};
		%--- số trang
		\node[circle,draw=white,line width=0.9pt,text=white,inner sep=0pt,minimum size=0.62cm,anchor=east]
		(so) at ([xshift=-\EZle,yshift=0.5*\EZdai]P) {\EZfontTrang\thepage};
		\node[inner sep=0pt,text=white,anchor=base east] (tr) at ([xshift=-0.15cm]so.base west) {\EZfontTrang Trang};
	\end{tikzpicture}%
}

%====================== CHÂN TRANG - TRANG CHẴN ====================
%  (soi gương: "Trang (số)" bên trái, ô chữ \tacgia bên phải)
\newcommand{\myfancyfootEven}{%
	\EZdorong{\tacgia}%
	\begin{tikzpicture}[remember picture,overlay]
		\coordinate (T) at (current page.south west);
		\coordinate (P) at (current page.south east);
		%--- dải nền
		\fill[ezNen] (T) rectangle ([yshift=\EZdai]P);
		%--- ô chữ sát mép trên dải: X là góc dưới bên trái, cạnh trái xiên kiểu "\"
		\coordinate (X) at ([xshift=-\EZle-\EZrong-\EZdem,yshift=\EZdai-\EZo]P);
		\fill[ezO] ([yshift=\EZdai-\EZo]P) -- (X) --++(-\EZxien,\EZo) -- ([yshift=\EZdai]P) --cycle;
		\EZvachxien{X}{-1}{-1}
		\node[ezchu,anchor=east] at ([xshift=-\EZle,yshift=\EZdai-0.5*\EZo]P) {\EZfontO\tacgia};
		%--- số trang
		%    (vòng tròn ẩn so0 nằm giữa dải để lấy đường chân chữ, giống trang lẻ)
		\node[circle,inner sep=0pt,minimum size=0.62cm] (so0) at ([yshift=0.5*\EZdai]T) {\EZfontTrang\phantom{\thepage}};
		\node[inner sep=0pt,text=white,anchor=base west] (tr) at ($(so0.base -| T)+(\EZle,0)$) {\EZfontTrang Trang};
		\node[circle,draw=white,line width=0.9pt,text=white,inner sep=0pt,minimum size=0.62cm,anchor=base west]
		(so) at ([xshift=0.15cm]tr.base east) {\EZfontTrang\thepage};
	\end{tikzpicture}%
}

%--- Giữ tương thích với các file cũ đang gọi \myfancyhead, \myfancyfoot
\newcommand{\myfancyhead}{\myfancyheadOdd}
\newcommand{\myfancyfoot}{\myfancyfootOdd}

%======================= GẮN VÀO TRANG =============================
\pagestyle{fancy}
\fancyhf{}
\fancyhead[LO]{\myfancyheadOdd}
\fancyhead[RE]{\myfancyheadEven}
\fancyfoot[LO]{\myfancyfootOdd}
\fancyfoot[RE]{\myfancyfootEven}
\renewcommand{\headrulewidth}{0pt}
\renewcommand{\footrulewidth}{0pt}
\fancypagestyle{plain}{%
	\fancyhf{}
	\fancyhead[LO]{\myfancyheadOdd}
	\fancyhead[RE]{\myfancyheadEven}
	\fancyfoot[LO]{\myfancyfootOdd}
	\fancyfoot[RE]{\myfancyfootEven}
	\renewcommand{\headrulewidth}{0pt}
	\renewcommand{\footrulewidth}{0pt}
}
%=================== HẾT PHẦN GỘP TỪ eduzy-de.tex ===================
\newcommand{\indam}[1]{
	{\fontsize{12pt}{9pt}\sffamily \selectfont\bfseries\color{gold1} #1}
}
\newcommand{\ind}[1]{
	{\fontsize{14pt}{9pt}\fontfamily{qag} \selectfont\bfseries\color{gold1} #1}
}
\renewcommand{\FalseEX}{\stepcounter{dapan}\textbf{\selectfont\color{gold1}\Alph{dapan}.}\ignorespaces}
\renewtheorem{ex}{\color{gold1}Câu}
\renewcommand{\baselinestretch}{1.1}
\newcommand{\hoac}[1]{
	\left[\begin{aligned}#1\end{aligned}\right.}
\newcommand{\heva}[1]{
	\left\{\begin{aligned}#1\end{aligned}\right.}
\usepackage{esvect}
\def\vec{\vv}
\def\vt{\vv}
\def\overrightarrow{\vv}
%====================== LOGO EDUZY VẼ BẰNG TIKZ ======================
%  Đường viền dò từ logo gốc; dải màu vàng kim đo trực tiếp trên ảnh gốc
%  (nghiêng 21 độ, tô bằng 80 dải hẹp nên hiển thị giống nhau trên mọi trình xem PDF).
%  Dùng: \eduzylogo          -> cao 2.4cm (mặc định)
%        \eduzylogo[1.5cm]   -> cao 1.5cm
\definecolor{ezText}{RGB}{245,200,46}
\newcommand{\eduzylogo}[1][2.4cm]{%
	\resizebox{!}{#1}{%
		\begin{tikzpicture}[x=0.01cm,y=0.01cm]
			\path[use as bounding box] (0,0) rectangle (224,-224);
			% mũ tốt nghiệp + vương miện
			\begin{scope}
				\clip
				(119.4,-176.1)..controls(118.9,-176.1)and(116.8,-176)..(114.6,-176)
				..controls(101.4,-176)and(85.8,-175.9)..(83.7,-175.7)
				..controls(81.8,-175.6)and(81,-175.6)..(80.2,-175.7)
				..controls(78.8,-175.9)and(77.7,-175.9)..(76.2,-175.7)
				..controls(74.9,-175.5)and(70.8,-175.5)..(64.3,-175.7)
				..controls(59.1,-175.9)and(58.3,-175.7)..(58.6,-173.7)
				..controls(58.6,-172.8)and(59.6,-169.6)..(60.3,-168)
				..controls(60.8,-166.9)and(61.4,-165.8)..(62.3,-165)
				..controls(63.1,-164.2)and(64.3,-163.4)..(65,-163.2)
				..controls(65.5,-163.1)and(72.3,-163)..(76.6,-163.1)
				..controls(77.9,-163.1)and(80.1,-163.1)..(81.4,-163)
				..controls(82.8,-162.9)and(86.3,-162.9)..(89.2,-162.9)
				..controls(93,-162.9)and(94.9,-162.8)..(96,-162.7)
				..controls(97.2,-162.5)and(101,-162.5)..(114,-162.4)
				..controls(123,-162.3)and(134.9,-162.2)..(140.2,-162.1)
				..controls(150.1,-162)and(152.2,-162.1)..(153.6,-162.4)
				..controls(154.4,-162.6)and(157.1,-163.8)..(158.1,-164.4)
				..controls(158.8,-164.8)and(159.7,-165.7)..(159.9,-166.3)
				..controls(160,-166.5)and(160.2,-167.1)..(160.2,-167.6)
				..controls(160.3,-168)and(160.6,-169.4)..(160.9,-170.6)
				..controls(161.2,-171.7)and(161.5,-173.1)..(161.6,-173.6)
				..controls(161.8,-174.4)and(161.8,-174.5)..(161.5,-174.9)
				..controls(161.2,-175.6)and(160.8,-175.9)..(159.9,-176)
				..controls(159.2,-176.1)and(154.1,-176.2)..(151.5,-176.2)
				..controls(148.1,-176.1)and(122.6,-176.1)..(121.6,-176.2)
				..controls(120.8,-176.2)and(119.8,-176.2)..(119.4,-176.1)--cycle
				(130.4,-160)..controls(129,-160)and(119.8,-160)..(110.1,-159.9)
				..controls(100.3,-159.8)and(88.5,-159.7)..(83.9,-159.6)
				..controls(77.8,-159.5)and(75.3,-159.5)..(74.5,-159.6)
				..controls(73.9,-159.7)and(72.5,-159.8)..(71.3,-159.8)
				..controls(69.6,-159.8)and(69,-159.8)..(68.5,-159.6)
				..controls(67.6,-159.3)and(66.8,-158.6)..(66.2,-157.6)
				..controls(66,-157.3)and(65.6,-156.5)..(65.3,-156.1)
				..controls(64.6,-155.1)and(64.5,-154.9)..(63.7,-152.9)
				..controls(63.2,-151.6)and(62.9,-150.9)..(61.8,-149.6)
				..controls(61.5,-149.2)and(61.1,-148.5)..(60.9,-147.9)
				..controls(60.7,-147.4)and(60.3,-146.6)..(60.1,-146.1)
				..controls(59.9,-145.5)and(59.5,-144.9)..(59,-144.3)
				..controls(58.5,-143.7)and(58.2,-143.2)..(57.9,-142.5)
				..controls(57.7,-141.9)and(57.3,-141.1)..(57,-140.7)
				..controls(56.6,-140.3)and(56.1,-139.3)..(55.8,-138.5)
				..controls(55.5,-137.8)and(55.1,-136.8)..(54.9,-136.5)
				..controls(54.7,-136.1)and(54.3,-135.2)..(53.9,-134.5)
				..controls(53.6,-133.8)and(53,-132.7)..(52.6,-132)
				..controls(51.6,-130.5)and(51.4,-130.2)..(51.1,-129)
				..controls(50.6,-127.5)and(50.6,-127.4)..(49.7,-126.3)
				..controls(49.2,-125.7)and(48.6,-124.8)..(48.3,-124.2)
				..controls(48,-123.6)and(47.5,-122.7)..(47.3,-122.2)
				..controls(47,-121.8)and(46.5,-120.7)..(46.1,-119.8)
				..controls(45.8,-118.9)and(45.2,-117.8)..(44.9,-117.2)
				..controls(44.2,-116.1)and(43.3,-114.4)..(42.3,-112.2)
				..controls(41.8,-111.3)and(41.3,-110.4)..(41.1,-110.1)
				..controls(40.5,-109.4)and(40.3,-109)..(40,-107.8)
				..controls(39.8,-107)and(39.6,-106.5)..(39.2,-106)
				..controls(38.9,-105.6)and(38.4,-104.7)..(38,-104)
				..controls(37.7,-103.4)and(37.2,-102.5)..(37,-102.1)
				..controls(36.7,-101.7)and(36.3,-100.8)..(36,-100.1)
				..controls(35.6,-98.9)and(35.3,-98.4)..(34.6,-97.6)
				..controls(34.2,-97.1)and(34.2,-96.3)..(34.6,-96)
				..controls(35,-95.7)and(35.5,-95.9)..(36.2,-96.7)
				..controls(36.5,-97)and(37.2,-97.5)..(37.6,-97.7)..controls(38,-98)and(38.6,-98.4)..(39.1,-98.9)
				..controls(39.6,-99.4)and(40.3,-100)..(40.6,-100.3)
				..controls(41.3,-100.8)and(42.4,-101.8)..(43.7,-102.9)
				..controls(44,-103.2)and(44.6,-103.7)..(45,-103.9)
				..controls(45.4,-104.2)and(45.9,-104.6)..(46.1,-104.8)
				..controls(47.1,-105.9)and(48.8,-107.1)..(50.5,-107.9)
				..controls(51.8,-108.6)and(52.3,-109)..(53.4,-110)
				..controls(54.8,-111.2)and(55.6,-111.6)..(57,-112.1)
				..controls(57.5,-112.3)and(58.5,-112.8)..(59.4,-113.3)
				..controls(60.2,-113.7)and(61.3,-114.3)..(61.9,-114.5)
				..controls(62.4,-114.7)and(63.3,-115.1)..(63.8,-115.3)
				..controls(64.4,-115.6)and(65.4,-115.9)..(66.2,-116.1)
				..controls(67,-116.3)and(68.4,-116.6)..(69.2,-116.9)
				..controls(72.3,-117.9)and(72.8,-118)..(78,-118.1)
				..controls(84,-118.3)and(86.1,-118.1)..(88.1,-117.1)
				..controls(89.4,-116.5)and(93,-115)..(93.7,-114.9)
				..controls(94.8,-114.6)and(95.4,-114.2)..(95.9,-113.4)
				..controls(96.2,-113)and(96.7,-112.5)..(97.1,-112.2)
				..controls(98,-111.6)and(101.3,-108.4)..(102.1,-107.2)
				..controls(102.4,-106.8)and(102.9,-106.2)..(103.2,-105.8)
				..controls(103.5,-105.5)and(104,-104.8)..(104.2,-104.2)
				..controls(104.5,-103.6)and(105,-102.8)..(105.3,-102.4)
				..controls(105.6,-102)and(106.1,-101.2)..(106.5,-100.6)
				..controls(107.1,-99.5)and(108.6,-97.5)..(109.5,-96.5)
				..controls(109.9,-96.1)and(110.1,-96)..(110.4,-96)
				..controls(111.1,-96)and(111.6,-96.5)..(111.8,-97.6)
				..controls(112,-98.6)and(112.3,-99.3)..(112.8,-99.9)
				..controls(113.7,-100.9)and(115.2,-103.1)..(115.7,-104.2)
				..controls(116.1,-105)and(116.5,-105.5)..(116.9,-106)
				..controls(117.3,-106.4)and(117.8,-106.9)..(118,-107.3)
				..controls(118.2,-107.6)and(118.6,-108)..(119,-108.3)
				..controls(119.6,-108.8)and(121.1,-110.3)..(121.8,-111.2)
				..controls(122.1,-111.7)and(122.5,-112)..(123.4,-112.4)
				..controls(124.8,-113.2)and(125.8,-113.8)..(127,-114.6)
				..controls(127.5,-114.9)and(128.1,-115.3)..(128.2,-115.4)
				..controls(128.4,-115.5)and(129.3,-115.8)..(130.2,-116)
				..controls(131.1,-116.2)and(132.5,-116.6)..(133.3,-116.9)
				..controls(135.7,-117.9)and(136.1,-117.9)..(142.5,-117.8)
				..controls(145.8,-117.8)and(148.4,-117.7)..(148.8,-117.6)
				..controls(149.2,-117.6)and(150.2,-117.2)..(151,-116.8)
				..controls(152,-116.4)and(152.8,-116)..(153.6,-115.9)
				..controls(154.5,-115.7)and(155.2,-115.4)..(156.3,-114.9)
				..controls(157.1,-114.5)and(158.2,-114)..(158.7,-113.8)
				..controls(159.1,-113.7)and(159.8,-113.3)..(160.2,-113)
				..controls(160.6,-112.7)and(161.3,-112.4)..(161.7,-112.2)
				..controls(162.7,-111.9)and(163.2,-111.7)..(164.1,-111.1)
				..controls(164.5,-110.8)and(165.3,-110.4)..(165.9,-110)
				..controls(166.5,-109.7)and(167.3,-109.2)..(167.6,-108.9)
				..controls(167.9,-108.6)and(168.6,-108.2)..(169,-108)
				..controls(169.4,-107.8)and(170.2,-107.3)..(170.7,-106.9)
				..controls(171.2,-106.5)and(172.1,-105.8)..(172.7,-105.4)
				..controls(173.2,-105)and(174.1,-104.4)..(174.5,-104)
				..controls(174.9,-103.7)and(175.7,-103.1)..(176.2,-102.6)
				..controls(176.7,-102.2)and(177.4,-101.5)..(177.9,-101.2)
				..controls(178.3,-100.8)and(178.8,-100.4)..(179,-100.2)
				..controls(179.2,-100)and(179.6,-99.6)..(180,-99.3)
				..controls(180.4,-99.1)and(180.9,-98.6)..(181.2,-98.4)
				..controls(181.5,-98.1)and(182.1,-97.7)..(182.5,-97.5)
				..controls(183.2,-97.2)and(183.5,-96.9)..(184.3,-95.9)
				..controls(184.6,-95.6)and(184.6,-95.6)..(184.8,-95.7)
				..controls(185,-95.8)and(185.1,-96)..(185.1,-96.2)
				..controls(185.1,-96.5)and(185,-96.7)..(184.6,-97.2)
				..controls(184.1,-97.8)and(184,-97.9)..(183.9,-98.8)
				..controls(183.7,-100.3)and(183.3,-101.3)..(182.9,-101.8)
				..controls(182.6,-102.1)and(182.3,-102.6)..(182.2,-102.9)
				..controls(182,-103.3)and(181.6,-104)..(181.2,-104.6)
				..controls(180.9,-105.2)and(180.5,-106.1)..(180.4,-106.6)
				..controls(179.9,-108)and(179.7,-108.5)..(179,-109.4)
				..controls(178.6,-109.8)and(178.2,-110.5)..(177.9,-111)
				..controls(177.7,-111.5)and(177.3,-112.3)..(177,-112.7)
				..controls(176.6,-113.2)and(176.4,-113.7)..(176.3,-114.2)
				..controls(176.2,-114.7)and(175.9,-115.2)..(175.4,-116)
				..controls(175,-116.6)and(174.4,-117.7)..(174,-118.5)
				..controls(173.7,-119.2)and(173.3,-120.1)..(173,-120.5)
				..controls(172.8,-120.9)and(172.4,-121.8)..(172.1,-122.5)
				..controls(171.8,-123.2)and(171.3,-124.1)..(171,-124.6)
				..controls(170.8,-125.1)and(170.3,-126)..(169.9,-126.6)
				..controls(169.6,-127.3)and(169.2,-128.1)..(169,-128.3)
				..controls(168.8,-128.6)and(168.5,-129.3)..(168.3,-129.8)
				..controls(168.1,-130.3)and(167.8,-131.1)..(167.5,-131.6)
				..controls(166.9,-132.5)and(164.2,-137.9)..(164.2,-138.1)
				..controls(164.2,-138.1)and(163.9,-138.8)..(163.4,-139.7)
				..controls(162.9,-140.5)and(162.4,-141.5)..(162.3,-141.9)
				..controls(161.9,-143)and(161.6,-143.4)..(160.9,-144.1)
				..controls(160.4,-144.6)and(160.3,-144.8)..(160.1,-145.4)
				..controls(160,-145.8)and(159.6,-146.8)..(159.2,-147.8)
				..controls(158.7,-148.8)and(158.1,-150)..(157.9,-150.6)
				..controls(157.6,-151.1)and(157.2,-151.9)..(156.9,-152.4)
				..controls(156.5,-152.8)and(156.1,-153.7)..(155.8,-154.4)
				..controls(155.6,-155)and(155.2,-155.9)..(154.9,-156.3)
				..controls(154.7,-156.7)and(154.4,-157.4)..(154.2,-157.9)
				..controls(154.1,-158.4)and(153.8,-159)..(153.5,-159.3)
				..controls(152.8,-160.1)and(153,-160.1)..(142.4,-160.1)
				..controls(137.2,-160.1)and(131.8,-160.1)..(130.4,-160)--cycle
				(141.5,-103.7)..controls(140.8,-103)and(139.8,-102.1)..(139.2,-101.6)
				..controls(138.5,-101.2)and(137.9,-100.6)..(137.8,-100.4)
				..controls(137,-99.5)and(134,-96.5)..(133.2,-96)
				..controls(132.7,-95.7)and(131.8,-94.9)..(131.1,-94.3)
				..controls(129.4,-92.7)and(127.9,-91.6)..(126.6,-91.1)
				..controls(125.2,-90.5)and(124.1,-90)..(122.9,-89.1)
				..controls(122.1,-88.6)and(121.3,-88.2)..(119.9,-87.7)
				..controls(118.8,-87.3)and(117.5,-86.8)..(116.9,-86.5)
				..controls(114.3,-85.5)and(113.7,-85.4)..(109.7,-85.5)
				..controls(105.3,-85.7)and(103.9,-85.9)..(101.6,-87)
				..controls(100.7,-87.4)and(99.4,-87.8)..(97.9,-88)
				..controls(97,-88.2)and(96.6,-88.3)..(95.8,-88.9)
				..controls(95.4,-89.2)and(94.7,-89.6)..(94.2,-89.8)
				..controls(92.9,-90.3)and(92.3,-90.7)..(90.9,-92.1)
				..controls(90.2,-92.7)and(89.4,-93.3)..(89,-93.5)..controls(87,-94.5)and(85.2,-96)..(83.6,-97.9)
				..controls(83.3,-98.2)and(82.5,-98.9)..(81.8,-99.4)
				..controls(80.7,-100.1)and(80.5,-100.3)..(80,-101.1)
				..controls(78.7,-103.1)and(77.2,-103.9)..(76.4,-103)
				..controls(75.7,-102.2)and(75.7,-103.2)..(75.7,-83.6)
				..controls(75.6,-73.6)and(75.6,-64.8)..(75.5,-64)
				..controls(75.5,-63.2)and(75.5,-61.4)..(75.6,-59.8)--(75.7,-57.1)--(76,-56.8)
				..controls(76.8,-56.2)and(78,-56.3)..(79.8,-57.2)
				..controls(80.4,-57.5)and(81.3,-57.8)..(81.8,-57.9)
				..controls(82.3,-58.1)and(83,-58.4)..(83.6,-58.8)..controls(84.3,-59.2)and(85,-59.5)..(86,-59.8)
				..controls(86.8,-60.1)and(88,-60.6)..(88.7,-60.9)
				..controls(89.3,-61.2)and(90.6,-61.7)..(91.4,-62)
				..controls(92.2,-62.3)and(93.7,-62.9)..(94.6,-63.3)
				..controls(95.6,-63.7)and(96.8,-64.1)..(97.4,-64.3)
				..controls(98,-64.5)and(99,-64.9)..(99.7,-65.2)
				..controls(100.4,-65.6)and(101.5,-66)..(102.3,-66.2)
				..controls(104.2,-66.8)and(105,-67.1)..(106.2,-67.8)
				..controls(107.9,-68.7)and(110.2,-69.4)..(111.3,-69.2)
				..controls(112.5,-69)and(115,-68.2)..(115.6,-67.8)
				..controls(117.3,-66.7)and(117.6,-66.6)..(119.5,-66.2)
				..controls(121.9,-65.8)and(122.1,-65.7)..(123,-65.1)
				..controls(124.1,-64.4)and(124.4,-64.3)..(126,-63.9)
				..controls(127,-63.7)and(127.7,-63.4)..(128.4,-63.1)
				..controls(129,-62.8)and(129.7,-62.5)..(130.6,-62.4)
				..controls(132.1,-62.1)and(134,-61.5)..(135.2,-61)
				..controls(135.7,-60.8)and(136.7,-60.5)..(137.3,-60.3)
				..controls(138.4,-60)and(139.3,-59.7)..(141.1,-59)
				..controls(141.6,-58.7)and(142,-58.6)..(142.5,-58.6)
				..controls(143.2,-58.6)and(143.3,-58.7)..(143.6,-59)
				..controls(144,-59.5)and(144.1,-59.9)..(144,-60.8)
				..controls(144,-61.2)and(144,-71.1)..(144,-82.8)--(143.9,-104.1)--(143.7,-104.5)
				..controls(143.5,-104.7)and(143.3,-104.8)..(143,-104.9)
				..controls(142.6,-104.9)and(142.5,-104.8)..(141.5,-103.7)--cycle
				(27.8,-91.4)..controls(27.1,-90.9)and(26.3,-90)..(26.1,-89.5)
				..controls(26,-89.2)and(25.8,-88.6)..(25.7,-87.9)
				..controls(25.4,-86.2)and(25.6,-84.8)..(26.5,-83.5)
				..controls(27.4,-82.3)and(28.3,-82)..(29.4,-82.3)
				..controls(31.1,-82.8)and(32.1,-84.9)..(32,-87.4)
				..controls(31.9,-89.3)and(31.2,-90.6)..(29.9,-91.4)
				..controls(29.1,-91.9)and(28.6,-91.9)..(27.8,-91.4)--cycle
				(190.1,-91.3)..controls(189.7,-91.2)and(189.3,-90.9)..(189.1,-90.7)
				..controls(188.3,-89.8)and(187.6,-87.7)..(187.6,-86.2)
				..controls(187.6,-85.5)and(187.7,-85.2)..(188.2,-84.4)
				..controls(189.1,-82.6)and(190.1,-81.8)..(191.3,-82.1)
				..controls(192,-82.3)and(193.5,-83.7)..(193.8,-84.4)
				..controls(194.1,-85.1)and(194.2,-87.1)..(194,-88.3)
				..controls(193.8,-89.8)and(193.3,-90.5)..(192.3,-91.3)
				..controls(191.7,-91.7)and(190.9,-91.7)..(190.1,-91.3)--cycle
				(173.7,-83.4)..controls(172.2,-82.5)and(172,-82.5)..(170.6,-82.4)
				..controls(169.7,-82.4)and(168.8,-82.4)..(168.3,-82.5)
				..controls(167.2,-82.7)and(166.9,-82.7)..(166.6,-82.2)
				..controls(166.3,-81.8)and(166.3,-81)..(166.8,-80)
				..controls(167.2,-78.9)and(167.4,-77.4)..(167.7,-74.1)
				..controls(167.8,-72.8)and(167.9,-71.8)..(167.9,-71.7)
				..controls(168,-71.7)and(168,-70.2)..(168,-68.4)..controls(168,-65.4)and(168,-64.9)..(167.7,-63)
				..controls(167.3,-60.5)and(167.1,-59.9)..(166.4,-58.4)
				..controls(166.1,-57.7)and(165.8,-56.9)..(165.7,-56.6)
				..controls(165.7,-56.2)and(165.4,-55.6)..(165.2,-55.2)
				..controls(164.6,-53.8)and(161.9,-50.7)..(161,-50.3)
				..controls(159.8,-49.7)and(157.9,-50)..(155.2,-51.2)
				..controls(154.7,-51.4)and(153.4,-51.8)..(152.2,-52.1)
				..controls(150.8,-52.4)and(149.4,-52.9)..(148.6,-53.2)
				..controls(147.8,-53.6)and(146.7,-54)..(146,-54.2)
				..controls(145.4,-54.4)and(144.5,-54.7)..(144,-54.9)
				..controls(143.6,-55.2)and(142.7,-55.5)..(141.8,-55.7)
				..controls(141,-55.9)and(139.6,-56.3)..(138.7,-56.6)
				..controls(137.8,-57)and(136.3,-57.4)..(135.3,-57.6)
				..controls(133.8,-57.9)and(133.2,-58.1)..(132.1,-58.7)
				..controls(131.1,-59.2)and(130.3,-59.4)..(129.4,-59.6)
				..controls(128.7,-59.8)and(127.8,-60)..(127.4,-60.1)
				..controls(127,-60.3)and(126,-60.5)..(125.1,-60.7)
				..controls(123.2,-61.1)and(121.7,-61.5)..(120.2,-62.1)
				..controls(119.7,-62.3)and(118.1,-62.8)..(116.7,-63.2)
				..controls(114.6,-63.9)and(114,-64.1)..(113,-64.7)
				..controls(111.9,-65.4)and(111.8,-65.4)..(111,-65.5)
				..controls(109.9,-65.5)and(109.3,-65.4)..(107.7,-64.7)
				..controls(107,-64.4)and(106,-64)..(105.3,-63.9)
				..controls(103.5,-63.5)and(102.2,-63)..(101.2,-62.5)
				..controls(100.1,-61.8)and(99.3,-61.5)..(97.8,-61.1)
				..controls(97.2,-61)and(96.2,-60.7)..(95.7,-60.4)
				..controls(95.1,-60.1)and(94.2,-59.8)..(93.5,-59.6)
				..controls(92.9,-59.4)and(91.7,-59)..(91,-58.7)..controls(90.2,-58.4)and(88.9,-58)..(88,-57.7)
				..controls(87.2,-57.5)and(86.1,-57.1)..(85.5,-56.8)
				..controls(85,-56.6)and(84.1,-56.2)..(83.5,-56)..controls(82.9,-55.7)and(81.8,-55.2)..(81,-54.9)
				..controls(80.1,-54.5)and(79.1,-54.2)..(78.7,-54.1)
				..controls(77.9,-53.8)and(77.3,-53.6)..(76.2,-52.9)
				..controls(75.7,-52.6)and(75.2,-52.4)..(74.1,-52.2)
				..controls(72.5,-51.9)and(72.4,-51.8)..(71.3,-51.1)
				..controls(69.9,-50.1)and(69.4,-49.9)..(67.9,-49.8)
				..controls(67,-49.7)and(66.6,-49.6)..(66.3,-49.4)
				..controls(66.1,-49.3)and(65.7,-49)..(65.4,-48.8)
				..controls(65.1,-48.7)and(64.2,-48.4)..(63.6,-48.2)
				..controls(62.9,-48)and(61.7,-47.6)..(60.9,-47.2)
				..controls(60.1,-46.8)and(58.9,-46.4)..(58.2,-46.2)
				..controls(57.6,-46)and(56.8,-45.7)..(56.5,-45.5)
				..controls(54.9,-44.6)and(54.6,-44.5)..(53.2,-44.2)
				..controls(51.7,-43.8)and(51,-43.6)..(49.9,-42.9)
				..controls(49.5,-42.6)and(48.9,-42.4)..(47.5,-42)
				..controls(45.7,-41.6)and(45.6,-41.5)..(44.3,-40.6)
				..controls(43.5,-40.1)and(42.8,-39.6)..(42.7,-39.4)
				..controls(42.4,-38.7)and(43,-38.1)..(44.5,-37.6)..controls(44.9,-37.5)and(46,-37)..(46.9,-36.6)
				..controls(47.9,-36.1)and(48.8,-35.8)..(49.6,-35.6)
				..controls(50.5,-35.4)and(51.3,-35.2)..(52.4,-34.6)
				..controls(53.3,-34.1)and(54.5,-33.7)..(54.9,-33.5)
				..controls(55.4,-33.4)and(56.4,-33)..(57.2,-32.6)
				..controls(58.6,-32)and(60.1,-31.5)..(61.6,-31.2)
				..controls(62.3,-31.1)and(62.7,-31)..(63.2,-30.7)
				..controls(64.1,-30)and(64.3,-29.9)..(65.3,-29.7)
				..controls(65.7,-29.6)and(66.5,-29.2)..(67.1,-28.9)
				..controls(67.6,-28.6)and(68.8,-28.2)..(69.6,-27.9)
				..controls(70.5,-27.7)and(71.6,-27.3)..(72.1,-27)
				..controls(72.6,-26.7)and(73.8,-26.2)..(74.8,-25.9)
				..controls(75.7,-25.5)and(76.7,-25.1)..(77.1,-25)
				..controls(77.4,-24.8)and(78.5,-24.4)..(79.6,-24.1)
				..controls(80.6,-23.8)and(82,-23.2)..(82.8,-22.9)
				..controls(83.5,-22.5)and(84.8,-22)..(85.8,-21.8)
				..controls(86.7,-21.5)and(87.8,-21.1)..(88.1,-20.9)
				..controls(88.4,-20.7)and(89.3,-20.3)..(90.1,-20.1)
				..controls(90.8,-19.8)and(91.9,-19.4)..(92.4,-19.1)
				..controls(93.5,-18.5)and(94.4,-18.2)..(95.4,-18.1)
				..controls(95.7,-18)and(96.4,-17.8)..(96.9,-17.5)
				..controls(97.5,-17.3)and(98.4,-16.9)..(99,-16.6)
				..controls(99.7,-16.3)and(100.9,-15.9)..(101.9,-15.6)
				..controls(102.9,-15.3)and(104.4,-14.7)..(105.3,-14.3)
				..controls(108.1,-13.2)and(108.5,-13.1)..(110.1,-13.1)
				..controls(111.4,-13.1)and(111.5,-13.1)..(112.4,-13.5)
				..controls(113.2,-14)and(115,-14.8)..(116.6,-15.4)
				..controls(116.9,-15.5)and(117.9,-15.8)..(118.8,-16.1)
				..controls(120.9,-16.6)and(121.5,-16.9)..(122.6,-17.9)
				..controls(123,-18.3)and(123.7,-18.8)..(123.9,-18.9)
				..controls(124.6,-19.4)and(126,-19.8)..(127.9,-20.2)
				..controls(129.6,-20.5)and(129.8,-20.6)..(131.1,-21.4)
				..controls(131.5,-21.7)and(132.4,-22.1)..(133.1,-22.3)
				..controls(136.1,-23.2)and(138.3,-24.1)..(140.2,-25)
				..controls(141,-25.5)and(141.5,-25.6)..(142.5,-25.8)
				..controls(145.1,-26.3)and(146,-26.6)..(149.4,-28.1)
				..controls(151.3,-28.9)and(153.5,-29.8)..(154.3,-30.1)
				..controls(155.1,-30.4)and(156.2,-30.8)..(156.6,-31)
				..controls(157.1,-31.2)and(158.2,-31.6)..(159.1,-31.8)
				..controls(160.6,-32.2)and(160.9,-32.3)..(161.7,-32.9)
				..controls(162.4,-33.3)and(162.9,-33.6)..(163.7,-33.8)
				..controls(164.3,-34)and(164.8,-34.1)..(164.9,-34.1)
				..controls(165,-34.1)and(165.6,-34.4)..(166.4,-34.8)
				..controls(167.3,-35.2)and(168.1,-35.4)..(169,-35.6)
				..controls(170,-35.8)and(170.5,-36)..(171.7,-36.6)
				..controls(172.4,-37)and(173.3,-37.4)..(173.6,-37.4)
				..controls(175.6,-37.7)and(176.2,-38)..(176.7,-39.1)
				..controls(176.9,-39.6)and(176.9,-39.7)..(176.7,-40.4)
				..controls(176.4,-41.6)and(175.8,-42.7)..(175.2,-43.4)
				..controls(174.5,-44.1)and(172.6,-45.5)..(171.7,-45.8)
				..controls(170.6,-46.2)and(169.1,-46)..(168.2,-45.4)
				..controls(167.9,-45.2)and(167.3,-44.8)..(166.9,-44.4)
				..controls(166.4,-44.1)and(165.8,-43.6)..(165.4,-43.2)
				..controls(165.1,-42.9)and(164.4,-42.3)..(163.9,-41.9)
				..controls(163.3,-41.5)and(162.7,-40.9)..(162.4,-40.5)
				..controls(161.4,-39.2)and(159.9,-38)..(158.8,-37.8)
				..controls(158.4,-37.8)and(158.3,-37.7)..(158.2,-37.2)
				..controls(158,-36.8)and(157.4,-36.1)..(157.1,-36.1)
				..controls(157,-36.1)and(156.7,-36.1)..(156.5,-36)
				..controls(156.1,-35.9)and(155.9,-35.7)..(155.7,-35.2)
				..controls(155.3,-34.6)and(155.1,-34.4)..(154,-34.1)
				..controls(153.6,-34)and(153.2,-33.7)..(152.9,-33.5)
				..controls(152.4,-32.9)and(151.6,-32.4)..(151.3,-32.4)
				..controls(150.9,-32.4)and(150.4,-32.6)..(150.4,-32.9)
				..controls(150.4,-33.2)and(150.8,-33.4)..(151.4,-33.6)
				..controls(152.1,-33.7)and(152.8,-34.1)..(153.5,-34.9)
				..controls(154.1,-35.5)and(154.3,-35.6)..(155,-35.8)
				..controls(155.7,-36)and(155.7,-36)..(156,-36.6)
				..controls(156.3,-37.3)and(156.7,-37.7)..(157.5,-37.8)
				..controls(158.1,-37.9)and(158.2,-38)..(158.4,-38.8)
				..controls(158.5,-39.4)and(158.6,-39.6)..(159.6,-40.4)
				..controls(160.2,-41)and(160.9,-41.6)..(161.3,-41.8)
				..controls(161.7,-42)and(161.9,-42.3)..(162.2,-42.8)
				..controls(162.4,-43.2)and(162.8,-43.8)..(163.1,-44.1)
				..controls(163.9,-44.8)and(165.1,-46.7)..(165.7,-48)
				..controls(166.2,-49.3)and(166.4,-49.7)..(167.2,-50.5)
				..controls(167.5,-50.8)and(167.9,-51.4)..(168.1,-51.8)
				..controls(168.3,-52.2)and(168.8,-53)..(169.1,-53.6)
				..controls(169.6,-54.3)and(169.8,-54.9)..(169.9,-55.6)
				..controls(170.4,-57.8)and(170.5,-58.1)..(171.2,-59.5)
				..controls(171.6,-60.3)and(172,-61.5)..(172.2,-62.1)
				..controls(172.3,-62.8)and(172.7,-63.7)..(172.9,-64.1)
				..controls(173.2,-64.6)and(173.4,-65.3)..(173.4,-65.6)
				..controls(173.8,-68.8)and(174.2,-70.6)..(174.9,-72.6)
				..controls(175.5,-74.5)and(175.9,-77.1)..(176.1,-80.2)
				..controls(176.3,-83.4)and(175.5,-84.4)..(173.7,-83.4)--cycle;
				\definecolor{ezs}{RGB}{247,179,26}\fill[ezs](-5.7,-156.6)--(-3.4,-157.5)--(61.9,12.7)--(59.7,13.6)--cycle;
				\definecolor{ezs}{RGB}{247,179,22}\fill[ezs](-3.8,-157.3)--(-1.5,-158.2)--(63.9,12)--(61.6,12.9)--cycle;
				\definecolor{ezs}{RGB}{248,179,19}\fill[ezs](-1.8,-158.1)--(0.4,-158.9)--(65.8,11.2)--(63.5,12.1)--cycle;
				\definecolor{ezs}{RGB}{248,178,16}\fill[ezs](0.1,-158.8)--(2.4,-159.7)--(67.7,10.5)--(65.4,11.4)--cycle;
				\definecolor{ezs}{RGB}{249,178,15}\fill[ezs](2,-159.5)--(4.3,-160.4)--(69.6,9.8)--(67.3,10.7)--cycle;
				\definecolor{ezs}{RGB}{251,177,17}\fill[ezs](3.9,-160.3)--(6.2,-161.1)--(71.5,9)--(69.2,9.9)--cycle;
				\definecolor{ezs}{RGB}{252,179,18}\fill[ezs](5.8,-161)--(8.1,-161.9)--(73.4,8.3)--(71.1,9.2)--cycle;
				\definecolor{ezs}{RGB}{252,182,17}\fill[ezs](7.7,-161.7)--(10,-162.6)--(75.3,7.6)--(73,8.5)--cycle;
				\definecolor{ezs}{RGB}{252,185,17}\fill[ezs](9.6,-162.5)--(11.9,-163.4)--(77.2,6.8)--(75,7.7)--cycle;
				\definecolor{ezs}{RGB}{251,189,18}\fill[ezs](11.5,-163.2)--(13.8,-164.1)--(79.2,6.1)--(76.9,7)--cycle;
				\definecolor{ezs}{RGB}{251,191,22}\fill[ezs](13.4,-163.9)--(15.7,-164.8)--(81.1,5.4)--(78.8,6.2)--cycle;
				\definecolor{ezs}{RGB}{252,193,25}\fill[ezs](15.4,-164.7)--(17.6,-165.6)--(83,4.6)--(80.7,5.5)--cycle;
				\definecolor{ezs}{RGB}{253,195,27}\fill[ezs](17.3,-165.4)--(19.6,-166.3)--(84.9,3.9)--(82.6,4.8)--cycle;
				\definecolor{ezs}{RGB}{253,196,31}\fill[ezs](19.2,-166.1)--(21.5,-167)--(86.8,3.2)--(84.5,4)--cycle;
				\definecolor{ezs}{RGB}{254,199,35}\fill[ezs](21.1,-166.9)--(23.4,-167.8)--(88.7,2.4)--(86.4,3.3)--cycle;
				\definecolor{ezs}{RGB}{253,201,38}\fill[ezs](23,-167.6)--(25.3,-168.5)--(90.6,1.7)--(88.3,2.6)--cycle;
				\definecolor{ezs}{RGB}{253,203,40}\fill[ezs](24.9,-168.3)--(27.2,-169.2)--(92.5,1)--(90.3,1.8)--cycle;
				\definecolor{ezs}{RGB}{253,205,42}\fill[ezs](26.8,-169.1)--(29.1,-170)--(94.4,0.2)--(92.2,1.1)--cycle;
				\definecolor{ezs}{RGB}{253,207,45}\fill[ezs](28.7,-169.8)--(31,-170.7)--(96.4,-0.5)--(94.1,0.4)--cycle;
				\definecolor{ezs}{RGB}{253,210,49}\fill[ezs](30.7,-170.5)--(32.9,-171.4)--(98.3,-1.2)--(96,-0.4)--cycle;
				\definecolor{ezs}{RGB}{253,212,54}\fill[ezs](32.6,-171.3)--(34.9,-172.2)--(100.2,-2)--(97.9,-1.1)--cycle;
				\definecolor{ezs}{RGB}{253,214,59}\fill[ezs](34.5,-172)--(36.8,-172.9)--(102.1,-2.7)--(99.8,-1.8)--cycle;
				\definecolor{ezs}{RGB}{253,216,63}\fill[ezs](36.4,-172.7)--(38.7,-173.6)--(104,-3.4)--(101.7,-2.6)--cycle;
				\definecolor{ezs}{RGB}{253,218,65}\fill[ezs](38.3,-173.5)--(40.6,-174.4)--(105.9,-4.2)--(103.6,-3.3)--cycle;
				\definecolor{ezs}{RGB}{254,220,66}\fill[ezs](40.2,-174.2)--(42.5,-175.1)--(107.8,-4.9)--(105.5,-4)--cycle;
				\definecolor{ezs}{RGB}{254,222,68}\fill[ezs](42.1,-175)--(44.4,-175.8)--(109.7,-5.6)--(107.5,-4.8)--cycle;
				\definecolor{ezs}{RGB}{254,224,71}\fill[ezs](44,-175.7)--(46.3,-176.6)--(111.7,-6.4)--(109.4,-5.5)--cycle;
				\definecolor{ezs}{RGB}{254,226,75}\fill[ezs](46,-176.4)--(48.2,-177.3)--(113.6,-7.1)--(111.3,-6.2)--cycle;
				\definecolor{ezs}{RGB}{254,228,79}\fill[ezs](47.9,-177.2)--(50.1,-178)--(115.5,-7.8)--(113.2,-7)--cycle;
				\definecolor{ezs}{RGB}{254,230,82}\fill[ezs](49.8,-177.9)--(52.1,-178.8)--(117.4,-8.6)--(115.1,-7.7)--cycle;
				\definecolor{ezs}{RGB}{254,232,85}\fill[ezs](51.7,-178.6)--(54,-179.5)--(119.3,-9.3)--(117,-8.4)--cycle;
				\definecolor{ezs}{RGB}{254,234,88}\fill[ezs](53.6,-179.4)--(55.9,-180.2)--(121.2,-10)--(118.9,-9.2)--cycle;
				\definecolor{ezs}{RGB}{254,237,91}\fill[ezs](55.5,-180.1)--(57.8,-181)--(123.1,-10.8)--(120.8,-9.9)--cycle;
				\definecolor{ezs}{RGB}{254,239,94}\fill[ezs](57.4,-180.8)--(59.7,-181.7)--(125,-11.5)--(122.8,-10.6)--cycle;
				\definecolor{ezs}{RGB}{254,241,97}\fill[ezs](59.3,-181.6)--(61.6,-182.4)--(127,-12.2)--(124.7,-11.4)--cycle;
				\definecolor{ezs}{RGB}{254,243,101}\fill[ezs](61.2,-182.3)--(63.5,-183.2)--(128.9,-13)--(126.6,-12.1)--cycle;
				\definecolor{ezs}{RGB}{254,245,103}\fill[ezs](63.2,-183)--(65.4,-183.9)--(130.8,-13.7)--(128.5,-12.8)--cycle;
				\definecolor{ezs}{RGB}{254,245,105}\fill[ezs](65.1,-183.8)--(67.4,-184.6)--(132.7,-14.4)--(130.4,-13.6)--cycle;
				\definecolor{ezs}{RGB}{254,245,105}\fill[ezs](67,-184.5)--(69.3,-185.4)--(134.6,-15.2)--(132.3,-14.3)--cycle;
				\definecolor{ezs}{RGB}{254,244,103}\fill[ezs](68.9,-185.2)--(71.2,-186.1)--(136.5,-15.9)--(134.2,-15)--cycle;
				\definecolor{ezs}{RGB}{253,243,101}\fill[ezs](70.8,-186)--(73.1,-186.8)--(138.4,-16.6)--(136.1,-15.8)--cycle;
				\definecolor{ezs}{RGB}{253,241,98}\fill[ezs](72.7,-186.7)--(75,-187.6)--(140.3,-17.4)--(138,-16.5)--cycle;
				\definecolor{ezs}{RGB}{253,239,95}\fill[ezs](74.6,-187.4)--(76.9,-188.3)--(142.2,-18.1)--(140,-17.2)--cycle;
				\definecolor{ezs}{RGB}{253,236,92}\fill[ezs](76.5,-188.2)--(78.8,-189)--(144.2,-18.8)--(141.9,-18)--cycle;
				\definecolor{ezs}{RGB}{253,234,88}\fill[ezs](78.5,-188.9)--(80.7,-189.8)--(146.1,-19.6)--(143.8,-18.7)--cycle;
				\definecolor{ezs}{RGB}{253,232,85}\fill[ezs](80.4,-189.6)--(82.7,-190.5)--(148,-20.3)--(145.7,-19.4)--cycle;
				\definecolor{ezs}{RGB}{253,230,82}\fill[ezs](82.3,-190.4)--(84.6,-191.2)--(149.9,-21)--(147.6,-20.2)--cycle;
				\definecolor{ezs}{RGB}{253,228,78}\fill[ezs](84.2,-191.1)--(86.5,-192)--(151.8,-21.8)--(149.5,-20.9)--cycle;
				\definecolor{ezs}{RGB}{253,226,74}\fill[ezs](86.1,-191.8)--(88.4,-192.7)--(153.7,-22.5)--(151.4,-21.6)--cycle;
				\definecolor{ezs}{RGB}{253,224,70}\fill[ezs](88,-192.6)--(90.3,-193.4)--(155.6,-23.3)--(153.3,-22.4)--cycle;
				\definecolor{ezs}{RGB}{253,222,67}\fill[ezs](89.9,-193.3)--(92.2,-194.2)--(157.5,-24)--(155.3,-23.1)--cycle;
				\definecolor{ezs}{RGB}{253,220,65}\fill[ezs](91.8,-194)--(94.1,-194.9)--(159.5,-24.7)--(157.2,-23.8)--cycle;
				\definecolor{ezs}{RGB}{253,217,63}\fill[ezs](93.7,-194.8)--(96,-195.6)--(161.4,-25.5)--(159.1,-24.6)--cycle;
				\definecolor{ezs}{RGB}{253,215,60}\fill[ezs](95.7,-195.5)--(97.9,-196.4)--(163.3,-26.2)--(161,-25.3)--cycle;
				\definecolor{ezs}{RGB}{253,213,56}\fill[ezs](97.6,-196.2)--(99.9,-197.1)--(165.2,-26.9)--(162.9,-26)--cycle;
				\definecolor{ezs}{RGB}{252,211,53}\fill[ezs](99.5,-197)--(101.8,-197.8)--(167.1,-27.7)--(164.8,-26.8)--cycle;
				\definecolor{ezs}{RGB}{252,209,49}\fill[ezs](101.4,-197.7)--(103.7,-198.6)--(169,-28.4)--(166.7,-27.5)--cycle;
				\definecolor{ezs}{RGB}{252,207,45}\fill[ezs](103.3,-198.4)--(105.6,-199.3)--(170.9,-29.1)--(168.6,-28.2)--cycle;
				\definecolor{ezs}{RGB}{253,205,42}\fill[ezs](105.2,-199.2)--(107.5,-200)--(172.8,-29.9)--(170.6,-29)--cycle;
				\definecolor{ezs}{RGB}{253,202,40}\fill[ezs](107.1,-199.9)--(109.4,-200.8)--(174.7,-30.6)--(172.5,-29.7)--cycle;
				\definecolor{ezs}{RGB}{253,200,38}\fill[ezs](109,-200.6)--(111.3,-201.5)--(176.7,-31.3)--(174.4,-30.4)--cycle;
				\definecolor{ezs}{RGB}{253,198,35}\fill[ezs](111,-201.4)--(113.2,-202.2)--(178.6,-32.1)--(176.3,-31.2)--cycle;
				\definecolor{ezs}{RGB}{253,196,31}\fill[ezs](112.9,-202.1)--(115.2,-203)--(180.5,-32.8)--(178.2,-31.9)--cycle;
				\definecolor{ezs}{RGB}{253,194,28}\fill[ezs](114.8,-202.8)--(117.1,-203.7)--(182.4,-33.5)--(180.1,-32.6)--cycle;
				\definecolor{ezs}{RGB}{253,192,26}\fill[ezs](116.7,-203.6)--(119,-204.5)--(184.3,-34.3)--(182,-33.4)--cycle;
				\definecolor{ezs}{RGB}{253,189,24}\fill[ezs](118.6,-204.3)--(120.9,-205.2)--(186.2,-35)--(183.9,-34.1)--cycle;
				\definecolor{ezs}{RGB}{253,186,20}\fill[ezs](120.5,-205)--(122.8,-205.9)--(188.1,-35.7)--(185.8,-34.8)--cycle;
				\definecolor{ezs}{RGB}{253,184,17}\fill[ezs](122.4,-205.8)--(124.7,-206.7)--(190,-36.5)--(187.8,-35.6)--cycle;
				\definecolor{ezs}{RGB}{252,183,15}\fill[ezs](124.3,-206.5)--(126.6,-207.4)--(192,-37.2)--(189.7,-36.3)--cycle;
				\definecolor{ezs}{RGB}{252,181,14}\fill[ezs](126.3,-207.2)--(128.5,-208.1)--(193.9,-37.9)--(191.6,-37.1)--cycle;
				\definecolor{ezs}{RGB}{252,179,13}\fill[ezs](128.2,-208)--(130.4,-208.9)--(195.8,-38.7)--(193.5,-37.8)--cycle;
				\definecolor{ezs}{RGB}{252,178,12}\fill[ezs](130.1,-208.7)--(132.4,-209.6)--(197.7,-39.4)--(195.4,-38.5)--cycle;
				\definecolor{ezs}{RGB}{251,178,10}\fill[ezs](132,-209.4)--(134.3,-210.3)--(199.6,-40.1)--(197.3,-39.3)--cycle;
				\definecolor{ezs}{RGB}{251,178,10}\fill[ezs](133.9,-210.2)--(136.2,-211.1)--(201.5,-40.9)--(199.2,-40)--cycle;
				\definecolor{ezs}{RGB}{250,177,11}\fill[ezs](135.8,-210.9)--(138.1,-211.8)--(203.4,-41.6)--(201.1,-40.7)--cycle;
				\definecolor{ezs}{RGB}{250,178,11}\fill[ezs](137.7,-211.6)--(140,-212.5)--(205.3,-42.3)--(203.1,-41.5)--cycle;
				\definecolor{ezs}{RGB}{248,180,10}\fill[ezs](139.6,-212.4)--(141.9,-213.3)--(207.3,-43.1)--(205,-42.2)--cycle;
				\definecolor{ezs}{RGB}{247,181,9}\fill[ezs](141.5,-213.1)--(143.8,-214)--(209.2,-43.8)--(206.9,-42.9)--cycle;
				\definecolor{ezs}{RGB}{246,182,9}\fill[ezs](143.5,-213.8)--(145.7,-214.7)--(211.1,-44.5)--(208.8,-43.7)--cycle;
				\definecolor{ezs}{RGB}{246,182,10}\fill[ezs](145.4,-214.6)--(147.7,-215.5)--(213,-45.3)--(210.7,-44.4)--cycle;
			\end{scope}
			% chữ EDUZY (Times đậm, làm dày nét bằng 5 lớp lệch nhẹ; căn đúng khung chữ logo gốc)
			\foreach \dx/\dy in {-0.7/0,0.7/0,0/-0.3,0/0.3,0/0}
			\node[anchor=center,inner sep=0pt,outer sep=0pt,text=ezText,align=center] at (109.6+\dx,-193.6+\dy)
			{\resizebox{1.01cm}{0.22cm}{\fontencoding{T1}\fontfamily{ptm}\fontseries{b}\selectfont EDUZY}};
		\end{tikzpicture}%
	}%
}
%====================== THIẾT LẬP THÂN ĐỀ ======================
%--- Môi trường tự luận {bt}: ex_test 3.3.4 chưa khai báo sẵn nên tạo ở đây
%    (ex_test đã tự móc \fixshowans và \listenumerate cho bt)
\makeatletter
\@ifundefined{bt}{\newtheorem{bt}{\color{gold1}\namebt}}{}
\makeatother
\renewcommand{\sectionmark}[1]{\markright{\thesection\ #1}}
%--- Khung tiêu đề (logo + tối đa 4 dòng; dòng 3, 4 để trống {} thì bỏ qua)
%  Đề thi:    \khungtieude{THI THỬ THPT QUỐC GIA NĂM 2026}{Môn Toán -- Đề số 21}
%                         {(Đề thi gồm: ...)}{Thời gian làm bài: 90 phút (không kể thời gian phát đề).}
%  Chuyên đề: \khungtieude{CHUYÊN ĐỀ 3 -- DÃY SỐ VÀ CÔNG THỨC TRUY HỒI}
%                         {Bồi dưỡng học sinh giỏi Toán THPT}{...}{}
\newcommand{\khungtieude}[4]{%
	\begin{tcolorbox}[colframe=gold1,colback=gold5!45,boxrule=0.5mm]
		\begin{minipage}[c]{3cm}
			\centerline{\eduzylogo[2.4cm]}
		\end{minipage}\hspace{0.5cm}
		\begin{minipage}[c]{12.5cm}
			\vskip 0.15 cm
			\centerline{\ind{#1}}
			\vskip 0.15 cm
			\centerline{\textbf{#2}}
			\if\relax\detokenize{#3}\relax\else
				\vskip 0.15 cm
				\centerline{\textit{#3}}
			\fi
			\if\relax\detokenize{#4}\relax\else
				\vskip 0.15 cm
				\centerline{\textit{#4}}
			\fi
			\vskip 0.15 cm
		\end{minipage}
	\end{tcolorbox}
}
%====================== BỔ SUNG CHO BÀI HỌC, CHUYÊN ĐỀ ======================
\usepackage[shortlabels]{enumitem}
\usepackage{tabularx}
%--- Danh sách gọn: \begin{enumerate}[gon] ... \end{enumerate}
\SetEnumitemKey{gon}{leftmargin=*,itemsep=1pt,parsep=0pt,topsep=3pt}
%--- Khung lý thuyết: \begin{khung}{} ... \end{khung}; tiêu đề chỉ dùng khi thật ngắn ("Mẹo", "Lỗi thường gặp")
\newtcolorbox{khung}[1]{breakable,enhanced,colframe=gold1,colback=gold5!45,
	boxrule=0.4mm,fonttitle=\bfseries\sffamily,coltitle=white,colbacktitle=gold1,
	attach boxed title to top left={xshift=6mm,yshift=-3mm},
	boxed title style={arc=1mm,boxrule=0.3mm},title={#1}}
%--- Tiêu đề mục nhỏ trong một dạng: \muc{Dấu hiệu}, \muc{Phương pháp}, \muc{Bài tự luyện}
\makeatletter
\newcommand{\muc}[1]{\par\medskip\noindent\indam{#1}\par\nobreak\smallskip\@afterindentfalse\@afterheading}
\makeatother
%--- Ô công thức cần nhớ (căn giữa, viền vàng, vừa khít nội dung)
%     \begin{center}\congthuc{$Q(n)=(n-1)!$}\end{center}
\newtcbox{\congthuc}{on line,colframe=gold1,colback=gold5!45,boxrule=0.5pt,arc=2pt,
	left=8pt,right=8pt,top=4pt,bottom=4pt}
%--- Công tắc lời giải bài tự luyện (ví dụ mẫu luôn dùng \loigiai)
%     \loigiaiTL[{đáp số}]{lời giải đầy đủ}
%     \TLfull -> in lời giải đầy đủ;  \TLdapso -> chỉ in một dòng "Đáp số";  \TLan -> ẩn hoàn toàn
%     (\CheDo là dethi thì mọi lời giải và đáp số đều ẩn)
\makeatletter
\newif\ifTLhien   \TLhientrue
\newif\ifTLdapso  \TLdapsofalse
\newcommand{\TLfull}{\TLhientrue\TLdapsofalse}
\newcommand{\TLdapso}{\TLhientrue\TLdapsotrue}
\newcommand{\TLan}{\TLhienfalse}
\newcommand{\loigiaiTL}[2][]{%
	\ifTLhien
		\ifTLdapso
			\def\@tlarg{#1}%
			\ifx\@tlarg\@empty\else
				\ifdefstring{\kindOPTN}{dethi}{}{\par\noindent{\bfseries\color{gold1}Đáp số:} #1\par}%
			\fi
		\else
			\loigiai{#2}%
		\fi
	\fi}
\makeatother
\TLdapso% mặc định: chỉ in đáp số bài tự luyện
%====================== BẢNG ĐÚNG/SAI (\choiceTFt) ======================
%  Hàng tiêu đề nền nâu đồng chữ trắng (Phát biểu | Đ | S), các hàng tô xen kẽ kem / trắng,
%  đường kẻ vàng kim, nhãn a) b) c) d) màu nâu đồng. Bản có đáp án: cột được chọn có vòng tròn
%  nâu đồng với dấu tích trắng (câu đúng tích ở cột Đúng, câu sai tích ở cột Sai).
%  Tiêu đề ghi đủ "Đúng | Sai": \OPTN{viettat=0} sau \input (cột Đúng/Sai rộng ra, phát biểu hẹp lại).
%  Dùng bảng gốc của ex_test: \def\BangDSgoc{1} trước %============================================================
%  PREAMBLE EDUZY v2 - dùng chung cho đề thi, bài học, chuyên đề (gói ex_test)
%  Khác bản preamble-eduzy.tex ở đầu trang, chân trang: dải nền nâu đồng tràn ngang,
%  logo + Eduzy.com.vn, ô chữ vàng kem cạnh xiên, chân trang "Trang (số)";
%  logo Eduzy chìm giữa mọi trang (tắt: \logochimtat, bật lại: \logochimbat).
%  Cách dùng ở file chính (xem mau-chuan.tex, mau-chuyen-de.tex):
%      \def\CheDo{loigiai}          % loigiai | dethi | solcolor | color | book
%      \def\tenchuong{<dòng chữ trên đầu trang>}
%      \def\thuonghieu{<chữ ở chân trang lẻ>}   % không khai báo thì dùng mặc định
%      \def\tacgia{<chữ ở chân trang chẵn>}      % không khai báo thì dùng mặc định
%      \def\khauhieu{<chữ cạnh logo ở đầu trang>}  % tùy chọn, mặc định Eduzy.com.vn
%      %============================================================
%  PREAMBLE EDUZY v2 - dùng chung cho đề thi, bài học, chuyên đề (gói ex_test)
%  Khác bản preamble-eduzy.tex ở đầu trang, chân trang: dải nền nâu đồng tràn ngang,
%  logo + Eduzy.com.vn, ô chữ vàng kem cạnh xiên, chân trang "Trang (số)";
%  logo Eduzy chìm giữa mọi trang (tắt: \logochimtat, bật lại: \logochimbat).
%  Cách dùng ở file chính (xem mau-chuan.tex, mau-chuyen-de.tex):
%      \def\CheDo{loigiai}          % loigiai | dethi | solcolor | color | book
%      \def\tenchuong{<dòng chữ trên đầu trang>}
%      \def\thuonghieu{<chữ ở chân trang lẻ>}   % không khai báo thì dùng mặc định
%      \def\tacgia{<chữ ở chân trang chẵn>}      % không khai báo thì dùng mặc định
%      \def\khauhieu{<chữ cạnh logo ở đầu trang>}  % tùy chọn, mặc định Eduzy.com.vn
%      \input{preamble-eduzy-v2}
%      \TLdapso                      % tài liệu: \TLfull | \TLdapso | \TLan
%      \begin{document} ... \end{document}
%  Logo Eduzy vẽ bằng TikZ ngay trong file (\eduzylogo), không cần file ảnh.
%  Biên dịch: pdfLaTeX (máy có gói vietnam) hoặc XeLaTeX.
%============================================================
\documentclass[12pt,a4paper,twoside]{article}%<-- twoside: bat dau trang/chan trang chan-le
\usepackage{iftex}
\usepackage{amsmath,amssymb,tikz,mathrsfs,tkz-tab}
\usepackage{graphicx}
\usepackage{pgfplots}% đồ thị hàm số
\pgfplotsset{compat=1.17}
%\usepackage{mathptmx}
\usepackage{mathpazo}
\ifXeTeX% XeLaTeX: font chữ có đủ dấu tiếng Việt, giữ nguyên font toán mathpazo
	\usepackage[no-math]{fontspec}
	\setmainfont{TeX Gyre Pagella}
	\newfontfamily\EZqag{TeX Gyre Adventor}[NFSSFamily=qag,Ligatures=TeX]
\fi
\usepackage[many]{tcolorbox}
\usetikzlibrary{math,through,calc,intersections,angles,quotes,shapes,shapes.geometric,arrows,patterns,matrix,chains,arrows.meta,decorations.shapes,decorations.fractals,decorations.markings,shadows,decorations.pathreplacing}
\usepackage[top=1.6cm, bottom=1.8cm, left=1.8cm, right=1.8cm]{geometry}% v2: top 1.6cm (bản cũ 1.5cm) để chữ không sát dải đầu trang
\providecommand{\CheDo}{loigiai}% loigiai | dethi | solcolor | color | book
\expandafter\usepackage\expandafter[\CheDo]{ex_test}
% \tenchuong: khai báo ở file chính (trước \input) hoặc \renewcommand sau \input
%============================================================
%  [GỘP TỪ eduzy-de.tex] EDUZY - MÀU NÂU ĐỒNG / VÀNG KIM
%  + ĐẦU TRANG, CHÂN TRANG KIỂU DẢI NGANG (v2), ĐỐI XỨNG TRANG CHẴN / TRANG LẺ.
%  (không cần file eduzy-de.tex riêng nữa)
%  Khối này phải đặt SAU \usepackage[...]{ex_test} và SAU \def\tenchuong{...}
%============================================================
\usepackage{lastpage}
\usepackage{needspace}
%--- DÀN TRANG ĐỀ THI: tách câu (THÊM 05/11/2026; chỉ thêm lệnh mới, KHÔNG đổi hành vi mặc định) ---
%  Hai công tắc, mặc định tắt (= preamble chạy y như cũ). Bật bằng cách đặt TRƯỚC \input{preamble-eduzy-v2}
%  trong file đề (tools/dan_trang/dan_trang.py tự ghi khi đặt \NeedspaceCau):
%    \def\TachCauDS{1}   câu Đúng/Sai được tách: câu dẫn (kèm hình) ở cuối trang, bảng Phát biểu sang trang sau
%                        (câu dẫn và hình không tách nhau, bảng không tách giữa hai trang)
%    \def\TachCauLG{1}   bản lời giải (\CheDo khác dethi): phần đề của câu cao hơn 1/3 trang được phép tách
%  Lệnh dùng trong file đề (giá trị do tools/dan_trang/dan_trang.py đo ở bản dethi):
%    \NeedspaceCau[<cao câu dẫn + hình>]{<cao cả khối câu>}   ngay trên \begin{ex} (câu đầu phần: trên tiêu đề phần)
%        công tắc tắt: đúng bằng \Needspace{<cao cả khối câu>} như cũ;
%        \TachCauDS=1 và có [..] (câu Đúng/Sai được chọn tách): chỉ giữ câu dẫn + hình;
%        \TachCauLG=1 ở bản lời giải: đòi tối đa 1/3 \textheight.
%    \NeedspaceBang{<cao bảng Phát biểu>}%   ngay trước \choiceTFt: \TachCauDS=1 thì giữ trọn bảng trên một trang;
%        công tắc tắt: không làm gì.
\providecommand{\TachCauDS}{0}
\providecommand{\TachCauLG}{0}
\newlength{\EZnsCau}
\newcommand{\NeedspaceCau}[2][]{%
	\setlength{\EZnsCau}{#2}%
	\ifnum\TachCauDS=1 \if\relax\detokenize{#1}\relax\else\setlength{\EZnsCau}{#1}\fi\fi
	\ifnum\TachCauLG=1 \ifdefstring{\kindOPTN}{dethi}{}{\ifdim\EZnsCau>0.3333\textheight\setlength{\EZnsCau}{0.3333\textheight}\fi}\fi
	\Needspace{\EZnsCau}}
\newcommand{\NeedspaceBang}[1]{\ifnum\TachCauDS=1 \Needspace{#1}\fi}

%--- Thang màu: đậm -> nhạt
\definecolor{gold1}{RGB}{126,96,0}    % nâu đồng đậm  #7E6000
\definecolor{gold2}{RGB}{192,144,3}   % vàng kim      #C09003
\definecolor{gold3}{RGB}{255,215,99}  % vàng nhạt     #FFD763
\definecolor{gold4}{RGB}{255,231,153} % vàng rất nhạt #FFE799
\definecolor{gold5}{RGB}{255,243,205} % vàng kem      #FFF3CD
%--- Ghi đè tông xanh/đỏ cũ
\definecolor{myblue}{RGB}{126,96,0}
\colorlet{blue}{gold1}
\colorlet{cyan}{gold2}
\colorlet{teal}{gold2}

%--- Tên chương mặc định (file đề nên tự khai báo lại)
\providecommand{\tenchuong}{}
%--- Chân trang
\providecommand{\thuonghieu}{TOÁN 12 -- Eduzy.com.vn}
\providecommand{\tacgia}{Tô Gia Bảo -- 0353703783}

%--- Chữ cạnh logo ở đầu trang
\providecommand{\khauhieu}{Eduzy.com.vn}

%====================== THÔNG SỐ ĐẦU TRANG, CHÂN TRANG (v2) ======================
%  Dải nền đậm tràn ngang trang; ô chữ vàng kem có cạnh xiên, kèm 2 vạch xiên chuyển màu.
%  Đầu trang lẻ : logo + \khauhieu (trái)   | ô chữ \tenchuong (phải)
%  Chân trang lẻ: ô chữ \thuonghieu (trái)  | Trang (số) (phải)
%  Trang chẵn soi gương trang lẻ; chân trang chẵn ghi \tacgia.
%--- Màu: chỉ dùng thang vàng kim Eduzy
\colorlet{ezNen}{gold1}% dải nền đậm
\colorlet{ezVach1}{gold3}% vạch xiên sát dải nền (đậm hơn)
\colorlet{ezVach2}{gold4}% vạch xiên sát ô chữ (nhạt hơn)
\colorlet{ezO}{gold5}% ô chữ vàng kem
%--- Kích thước
\newlength{\EZdai}\setlength{\EZdai}{1.1cm}% chiều cao dải nền
\newlength{\EZo}\setlength{\EZo}{0.68cm}% chiều cao ô chữ
\newlength{\EZxien}\setlength{\EZxien}{0.55cm}% độ lệch ngang của cạnh xiên
\newlength{\EZvach}\setlength{\EZvach}{0.75cm}% bề rộng mỗi vạch xiên
\newlength{\EZkhe}\setlength{\EZkhe}{0.12cm}% khe trắng giữa vạch và ô chữ
\newlength{\EZdem}\setlength{\EZdem}{0.4cm}% khoảng đệm từ chữ tới cạnh xiên
\newlength{\EZle}\setlength{\EZle}{1.8cm}% lề trái/phải của chữ (bằng lề thân bài)
\newlength{\EZrongmin}\setlength{\EZrongmin}{4.5cm}% ô chữ rộng tối thiểu
\newlength{\EZrong}% bề rộng chữ trong ô (tính tự động)
%--- Font
\newcommand{\EZfontO}{\fontsize{10.5pt}{12pt}\selectfont\bfseries}% chữ trong ô
\newcommand{\EZfontKH}{\fontfamily{qag}\fontsize{11pt}{12pt}\selectfont\bfseries}% \khauhieu cạnh logo
\newcommand{\EZfontTrang}{\fontsize{11.5pt}{12pt}\selectfont\bfseries}% chữ "Trang" và số trang
%--- Đo bề rộng chữ trong ô
\newcommand{\EZdorong}[1]{%
	\settowidth{\EZrong}{\EZfontO #1}%
	\ifdim\EZrong<\EZrongmin\setlength{\EZrong}{\EZrongmin}\fi}
%--- Logo + chữ trắng \khauhieu (dựng sẵn một lần khi vào \begin{document})
\newsavebox{\EZlogobox}
\AtBeginDocument{\sbox{\EZlogobox}{%
		\raisebox{-0.5\height}{\eduzylogo[1cm]}\hspace{0.15cm}%
		\raisebox{-0.5\height}{\color{white}\EZfontKH\khauhieu}}}
%--- Hai vạch xiên + khe trắng. #1: điểm gốc ở mép dưới, #2: hướng (1 sang phải, -1 sang trái),
%    #3: hướng nghiêng (1 kiểu "/", -1 kiểu "\")
\newcommand{\EZvachxien}[3]{%
	\fill[white] (#1) --++({(#3)*\EZxien},\EZo) --++({(#2)*\EZkhe},0) --++({-(#3)*\EZxien},-\EZo) --cycle;
	\fill[ezVach2] ([xshift=(#2)*\EZkhe]#1) --++({(#3)*\EZxien},\EZo) --++({(#2)*\EZvach},0) --++({-(#3)*\EZxien},-\EZo) --cycle;
	\fill[ezVach1] ([xshift=(#2)*(\EZkhe+\EZvach)]#1) --++({(#3)*\EZxien},\EZo) --++({(#2)*\EZvach},0) --++({-(#3)*\EZxien},-\EZo) --cycle;
	\draw[white,line width=0.6pt] ([xshift=(#2)*(\EZkhe+\EZvach)]#1) --++({(#3)*\EZxien},\EZo);
}
%--- Chữ trong ô (căn giữa theo chiều cao ô, không lệch vì dấu tiếng Việt)
\tikzset{ezchu/.style={inner sep=0pt,text=ezNen,text height=0.72em,text depth=0.22em}}

%====================== LOGO CHÌM GIỮA TRANG ======================
%  Vẽ cùng đầu trang nên nằm dưới chữ của thân bài, có trên mọi trang.
%  Tắt: \logochimtat (từ trang đó trở đi), bật lại: \logochimbat.
%  Đổi cỡ, độ đậm: khai báo trước \input, ví dụ \def\EZchimcao{9cm} \def\EZchimdo{0.1}
\providecommand{\EZchimcao}{10cm}% chiều cao logo chìm
\providecommand{\EZchimdo}{0.12}% độ đậm (0 = trong suốt, 1 = rõ hẳn)
\newif\ifEZchim \EZchimtrue
\newcommand{\logochimtat}{\EZchimfalse}
\newcommand{\logochimbat}{\EZchimtrue}
%  Dựng sẵn logo mờ một lần khi vào \begin{document}; transparency group làm mờ cả logo
%  một lượt nên không lộ vệt ở chỗ các dải màu chồng nhau.
\newsavebox{\EZchimtho}
\newsavebox{\EZchimbox}
\AtBeginDocument{%
	\sbox{\EZchimtho}{\eduzylogo[\EZchimcao]}%
	\sbox{\EZchimbox}{\begin{tikzpicture}
			\begin{scope}[transparency group,opacity=\EZchimdo]
				\node[inner sep=0pt] {\usebox{\EZchimtho}};
			\end{scope}
	\end{tikzpicture}}}
\newcommand{\EZvechim}{%
	\ifEZchim
		\node[inner sep=0pt] at (current page.center) {\usebox{\EZchimbox}};
	\fi}

%====================== ĐẦU TRANG - TRANG LẺ ======================
%  (logo bên trái, ô chữ \tenchuong bên phải)
\newcommand{\myfancyheadOdd}{%
	\EZdorong{\tenchuong}%
	\begin{tikzpicture}[remember picture,overlay]
		\coordinate (T) at (current page.north west);
		\coordinate (P) at (current page.north east);
		%--- logo chìm giữa trang
		\EZvechim
		%--- dải nền
		\fill[ezNen] (T) rectangle ([yshift=-\EZdai]P);
		%--- ô chữ: X là góc dưới bên trái, cạnh trái xiên kiểu "/"
		\coordinate (X) at ([xshift=-\EZle-\EZrong-\EZdem-\EZxien,yshift=-\EZdai]P);
		\fill[ezO] (X) --++(\EZxien,\EZo) -- ([yshift=-\EZdai+\EZo]P) -- ([yshift=-\EZdai]P) --cycle;
		\EZvachxien{X}{-1}{1}
		\node[ezchu,anchor=east] at ([xshift=-\EZle,yshift=-\EZdai+0.5*\EZo]P) {\EZfontO\tenchuong};
		%--- logo + tên
		\node[inner sep=0pt,anchor=west] at ([xshift=\EZle,yshift=-0.5*\EZdai]T) {\usebox{\EZlogobox}};
	\end{tikzpicture}%
}

%====================== ĐẦU TRANG - TRANG CHẴN =====================
%  (soi gương: ô chữ \tenchuong bên trái, logo bên phải)
\newcommand{\myfancyheadEven}{%
	\EZdorong{\tenchuong}%
	\begin{tikzpicture}[remember picture,overlay]
		\coordinate (T) at (current page.north west);
		\coordinate (P) at (current page.north east);
		%--- logo chìm giữa trang
		\EZvechim
		%--- dải nền
		\fill[ezNen] (T) rectangle ([yshift=-\EZdai]P);
		%--- ô chữ: X là góc dưới bên phải, cạnh phải xiên kiểu "\"
		\coordinate (X) at ([xshift=\EZle+\EZrong+\EZdem+\EZxien,yshift=-\EZdai]T);
		\fill[ezO] (X) --++(-\EZxien,\EZo) -- ([yshift=-\EZdai+\EZo]T) -- ([yshift=-\EZdai]T) --cycle;
		\EZvachxien{X}{1}{-1}
		\node[ezchu,anchor=west] at ([xshift=\EZle,yshift=-\EZdai+0.5*\EZo]T) {\EZfontO\tenchuong};
		%--- logo + tên
		\node[inner sep=0pt,anchor=east] at ([xshift=-\EZle,yshift=-0.5*\EZdai]P) {\usebox{\EZlogobox}};
	\end{tikzpicture}%
}

%====================== CHÂN TRANG - TRANG LẺ ======================
%  (ô chữ \thuonghieu bên trái, "Trang (số)" bên phải)
\newcommand{\myfancyfootOdd}{%
	\EZdorong{\thuonghieu}%
	\begin{tikzpicture}[remember picture,overlay]
		\coordinate (T) at (current page.south west);
		\coordinate (P) at (current page.south east);
		%--- dải nền
		\fill[ezNen] (T) rectangle ([yshift=\EZdai]P);
		%--- ô chữ sát mép trên dải: X là góc dưới bên phải, cạnh phải xiên kiểu "/"
		\coordinate (X) at ([xshift=\EZle+\EZrong+\EZdem,yshift=\EZdai-\EZo]T);
		\fill[ezO] ([yshift=\EZdai-\EZo]T) -- (X) --++(\EZxien,\EZo) -- ([yshift=\EZdai]T) --cycle;
		\EZvachxien{X}{1}{1}
		\node[ezchu,anchor=west] at ([xshift=\EZle,yshift=\EZdai-0.5*\EZo]T) {\EZfontO\thuonghieu};
		%--- số trang
		\node[circle,draw=white,line width=0.9pt,text=white,inner sep=0pt,minimum size=0.62cm,anchor=east]
		(so) at ([xshift=-\EZle,yshift=0.5*\EZdai]P) {\EZfontTrang\thepage};
		\node[inner sep=0pt,text=white,anchor=base east] (tr) at ([xshift=-0.15cm]so.base west) {\EZfontTrang Trang};
	\end{tikzpicture}%
}

%====================== CHÂN TRANG - TRANG CHẴN ====================
%  (soi gương: "Trang (số)" bên trái, ô chữ \tacgia bên phải)
\newcommand{\myfancyfootEven}{%
	\EZdorong{\tacgia}%
	\begin{tikzpicture}[remember picture,overlay]
		\coordinate (T) at (current page.south west);
		\coordinate (P) at (current page.south east);
		%--- dải nền
		\fill[ezNen] (T) rectangle ([yshift=\EZdai]P);
		%--- ô chữ sát mép trên dải: X là góc dưới bên trái, cạnh trái xiên kiểu "\"
		\coordinate (X) at ([xshift=-\EZle-\EZrong-\EZdem,yshift=\EZdai-\EZo]P);
		\fill[ezO] ([yshift=\EZdai-\EZo]P) -- (X) --++(-\EZxien,\EZo) -- ([yshift=\EZdai]P) --cycle;
		\EZvachxien{X}{-1}{-1}
		\node[ezchu,anchor=east] at ([xshift=-\EZle,yshift=\EZdai-0.5*\EZo]P) {\EZfontO\tacgia};
		%--- số trang
		%    (vòng tròn ẩn so0 nằm giữa dải để lấy đường chân chữ, giống trang lẻ)
		\node[circle,inner sep=0pt,minimum size=0.62cm] (so0) at ([yshift=0.5*\EZdai]T) {\EZfontTrang\phantom{\thepage}};
		\node[inner sep=0pt,text=white,anchor=base west] (tr) at ($(so0.base -| T)+(\EZle,0)$) {\EZfontTrang Trang};
		\node[circle,draw=white,line width=0.9pt,text=white,inner sep=0pt,minimum size=0.62cm,anchor=base west]
		(so) at ([xshift=0.15cm]tr.base east) {\EZfontTrang\thepage};
	\end{tikzpicture}%
}

%--- Giữ tương thích với các file cũ đang gọi \myfancyhead, \myfancyfoot
\newcommand{\myfancyhead}{\myfancyheadOdd}
\newcommand{\myfancyfoot}{\myfancyfootOdd}

%======================= GẮN VÀO TRANG =============================
\pagestyle{fancy}
\fancyhf{}
\fancyhead[LO]{\myfancyheadOdd}
\fancyhead[RE]{\myfancyheadEven}
\fancyfoot[LO]{\myfancyfootOdd}
\fancyfoot[RE]{\myfancyfootEven}
\renewcommand{\headrulewidth}{0pt}
\renewcommand{\footrulewidth}{0pt}
\fancypagestyle{plain}{%
	\fancyhf{}
	\fancyhead[LO]{\myfancyheadOdd}
	\fancyhead[RE]{\myfancyheadEven}
	\fancyfoot[LO]{\myfancyfootOdd}
	\fancyfoot[RE]{\myfancyfootEven}
	\renewcommand{\headrulewidth}{0pt}
	\renewcommand{\footrulewidth}{0pt}
}
%=================== HẾT PHẦN GỘP TỪ eduzy-de.tex ===================
\newcommand{\indam}[1]{
	{\fontsize{12pt}{9pt}\sffamily \selectfont\bfseries\color{gold1} #1}
}
\newcommand{\ind}[1]{
	{\fontsize{14pt}{9pt}\fontfamily{qag} \selectfont\bfseries\color{gold1} #1}
}
\renewcommand{\FalseEX}{\stepcounter{dapan}\textbf{\selectfont\color{gold1}\Alph{dapan}.}\ignorespaces}
\renewtheorem{ex}{\color{gold1}Câu}
\renewcommand{\baselinestretch}{1.1}
\newcommand{\hoac}[1]{
	\left[\begin{aligned}#1\end{aligned}\right.}
\newcommand{\heva}[1]{
	\left\{\begin{aligned}#1\end{aligned}\right.}
\usepackage{esvect}
\def\vec{\vv}
\def\vt{\vv}
\def\overrightarrow{\vv}
%====================== LOGO EDUZY VẼ BẰNG TIKZ ======================
%  Đường viền dò từ logo gốc; dải màu vàng kim đo trực tiếp trên ảnh gốc
%  (nghiêng 21 độ, tô bằng 80 dải hẹp nên hiển thị giống nhau trên mọi trình xem PDF).
%  Dùng: \eduzylogo          -> cao 2.4cm (mặc định)
%        \eduzylogo[1.5cm]   -> cao 1.5cm
\definecolor{ezText}{RGB}{245,200,46}
\newcommand{\eduzylogo}[1][2.4cm]{%
	\resizebox{!}{#1}{%
		\begin{tikzpicture}[x=0.01cm,y=0.01cm]
			\path[use as bounding box] (0,0) rectangle (224,-224);
			% mũ tốt nghiệp + vương miện
			\begin{scope}
				\clip
				(119.4,-176.1)..controls(118.9,-176.1)and(116.8,-176)..(114.6,-176)
				..controls(101.4,-176)and(85.8,-175.9)..(83.7,-175.7)
				..controls(81.8,-175.6)and(81,-175.6)..(80.2,-175.7)
				..controls(78.8,-175.9)and(77.7,-175.9)..(76.2,-175.7)
				..controls(74.9,-175.5)and(70.8,-175.5)..(64.3,-175.7)
				..controls(59.1,-175.9)and(58.3,-175.7)..(58.6,-173.7)
				..controls(58.6,-172.8)and(59.6,-169.6)..(60.3,-168)
				..controls(60.8,-166.9)and(61.4,-165.8)..(62.3,-165)
				..controls(63.1,-164.2)and(64.3,-163.4)..(65,-163.2)
				..controls(65.5,-163.1)and(72.3,-163)..(76.6,-163.1)
				..controls(77.9,-163.1)and(80.1,-163.1)..(81.4,-163)
				..controls(82.8,-162.9)and(86.3,-162.9)..(89.2,-162.9)
				..controls(93,-162.9)and(94.9,-162.8)..(96,-162.7)
				..controls(97.2,-162.5)and(101,-162.5)..(114,-162.4)
				..controls(123,-162.3)and(134.9,-162.2)..(140.2,-162.1)
				..controls(150.1,-162)and(152.2,-162.1)..(153.6,-162.4)
				..controls(154.4,-162.6)and(157.1,-163.8)..(158.1,-164.4)
				..controls(158.8,-164.8)and(159.7,-165.7)..(159.9,-166.3)
				..controls(160,-166.5)and(160.2,-167.1)..(160.2,-167.6)
				..controls(160.3,-168)and(160.6,-169.4)..(160.9,-170.6)
				..controls(161.2,-171.7)and(161.5,-173.1)..(161.6,-173.6)
				..controls(161.8,-174.4)and(161.8,-174.5)..(161.5,-174.9)
				..controls(161.2,-175.6)and(160.8,-175.9)..(159.9,-176)
				..controls(159.2,-176.1)and(154.1,-176.2)..(151.5,-176.2)
				..controls(148.1,-176.1)and(122.6,-176.1)..(121.6,-176.2)
				..controls(120.8,-176.2)and(119.8,-176.2)..(119.4,-176.1)--cycle
				(130.4,-160)..controls(129,-160)and(119.8,-160)..(110.1,-159.9)
				..controls(100.3,-159.8)and(88.5,-159.7)..(83.9,-159.6)
				..controls(77.8,-159.5)and(75.3,-159.5)..(74.5,-159.6)
				..controls(73.9,-159.7)and(72.5,-159.8)..(71.3,-159.8)
				..controls(69.6,-159.8)and(69,-159.8)..(68.5,-159.6)
				..controls(67.6,-159.3)and(66.8,-158.6)..(66.2,-157.6)
				..controls(66,-157.3)and(65.6,-156.5)..(65.3,-156.1)
				..controls(64.6,-155.1)and(64.5,-154.9)..(63.7,-152.9)
				..controls(63.2,-151.6)and(62.9,-150.9)..(61.8,-149.6)
				..controls(61.5,-149.2)and(61.1,-148.5)..(60.9,-147.9)
				..controls(60.7,-147.4)and(60.3,-146.6)..(60.1,-146.1)
				..controls(59.9,-145.5)and(59.5,-144.9)..(59,-144.3)
				..controls(58.5,-143.7)and(58.2,-143.2)..(57.9,-142.5)
				..controls(57.7,-141.9)and(57.3,-141.1)..(57,-140.7)
				..controls(56.6,-140.3)and(56.1,-139.3)..(55.8,-138.5)
				..controls(55.5,-137.8)and(55.1,-136.8)..(54.9,-136.5)
				..controls(54.7,-136.1)and(54.3,-135.2)..(53.9,-134.5)
				..controls(53.6,-133.8)and(53,-132.7)..(52.6,-132)
				..controls(51.6,-130.5)and(51.4,-130.2)..(51.1,-129)
				..controls(50.6,-127.5)and(50.6,-127.4)..(49.7,-126.3)
				..controls(49.2,-125.7)and(48.6,-124.8)..(48.3,-124.2)
				..controls(48,-123.6)and(47.5,-122.7)..(47.3,-122.2)
				..controls(47,-121.8)and(46.5,-120.7)..(46.1,-119.8)
				..controls(45.8,-118.9)and(45.2,-117.8)..(44.9,-117.2)
				..controls(44.2,-116.1)and(43.3,-114.4)..(42.3,-112.2)
				..controls(41.8,-111.3)and(41.3,-110.4)..(41.1,-110.1)
				..controls(40.5,-109.4)and(40.3,-109)..(40,-107.8)
				..controls(39.8,-107)and(39.6,-106.5)..(39.2,-106)
				..controls(38.9,-105.6)and(38.4,-104.7)..(38,-104)
				..controls(37.7,-103.4)and(37.2,-102.5)..(37,-102.1)
				..controls(36.7,-101.7)and(36.3,-100.8)..(36,-100.1)
				..controls(35.6,-98.9)and(35.3,-98.4)..(34.6,-97.6)
				..controls(34.2,-97.1)and(34.2,-96.3)..(34.6,-96)
				..controls(35,-95.7)and(35.5,-95.9)..(36.2,-96.7)
				..controls(36.5,-97)and(37.2,-97.5)..(37.6,-97.7)..controls(38,-98)and(38.6,-98.4)..(39.1,-98.9)
				..controls(39.6,-99.4)and(40.3,-100)..(40.6,-100.3)
				..controls(41.3,-100.8)and(42.4,-101.8)..(43.7,-102.9)
				..controls(44,-103.2)and(44.6,-103.7)..(45,-103.9)
				..controls(45.4,-104.2)and(45.9,-104.6)..(46.1,-104.8)
				..controls(47.1,-105.9)and(48.8,-107.1)..(50.5,-107.9)
				..controls(51.8,-108.6)and(52.3,-109)..(53.4,-110)
				..controls(54.8,-111.2)and(55.6,-111.6)..(57,-112.1)
				..controls(57.5,-112.3)and(58.5,-112.8)..(59.4,-113.3)
				..controls(60.2,-113.7)and(61.3,-114.3)..(61.9,-114.5)
				..controls(62.4,-114.7)and(63.3,-115.1)..(63.8,-115.3)
				..controls(64.4,-115.6)and(65.4,-115.9)..(66.2,-116.1)
				..controls(67,-116.3)and(68.4,-116.6)..(69.2,-116.9)
				..controls(72.3,-117.9)and(72.8,-118)..(78,-118.1)
				..controls(84,-118.3)and(86.1,-118.1)..(88.1,-117.1)
				..controls(89.4,-116.5)and(93,-115)..(93.7,-114.9)
				..controls(94.8,-114.6)and(95.4,-114.2)..(95.9,-113.4)
				..controls(96.2,-113)and(96.7,-112.5)..(97.1,-112.2)
				..controls(98,-111.6)and(101.3,-108.4)..(102.1,-107.2)
				..controls(102.4,-106.8)and(102.9,-106.2)..(103.2,-105.8)
				..controls(103.5,-105.5)and(104,-104.8)..(104.2,-104.2)
				..controls(104.5,-103.6)and(105,-102.8)..(105.3,-102.4)
				..controls(105.6,-102)and(106.1,-101.2)..(106.5,-100.6)
				..controls(107.1,-99.5)and(108.6,-97.5)..(109.5,-96.5)
				..controls(109.9,-96.1)and(110.1,-96)..(110.4,-96)
				..controls(111.1,-96)and(111.6,-96.5)..(111.8,-97.6)
				..controls(112,-98.6)and(112.3,-99.3)..(112.8,-99.9)
				..controls(113.7,-100.9)and(115.2,-103.1)..(115.7,-104.2)
				..controls(116.1,-105)and(116.5,-105.5)..(116.9,-106)
				..controls(117.3,-106.4)and(117.8,-106.9)..(118,-107.3)
				..controls(118.2,-107.6)and(118.6,-108)..(119,-108.3)
				..controls(119.6,-108.8)and(121.1,-110.3)..(121.8,-111.2)
				..controls(122.1,-111.7)and(122.5,-112)..(123.4,-112.4)
				..controls(124.8,-113.2)and(125.8,-113.8)..(127,-114.6)
				..controls(127.5,-114.9)and(128.1,-115.3)..(128.2,-115.4)
				..controls(128.4,-115.5)and(129.3,-115.8)..(130.2,-116)
				..controls(131.1,-116.2)and(132.5,-116.6)..(133.3,-116.9)
				..controls(135.7,-117.9)and(136.1,-117.9)..(142.5,-117.8)
				..controls(145.8,-117.8)and(148.4,-117.7)..(148.8,-117.6)
				..controls(149.2,-117.6)and(150.2,-117.2)..(151,-116.8)
				..controls(152,-116.4)and(152.8,-116)..(153.6,-115.9)
				..controls(154.5,-115.7)and(155.2,-115.4)..(156.3,-114.9)
				..controls(157.1,-114.5)and(158.2,-114)..(158.7,-113.8)
				..controls(159.1,-113.7)and(159.8,-113.3)..(160.2,-113)
				..controls(160.6,-112.7)and(161.3,-112.4)..(161.7,-112.2)
				..controls(162.7,-111.9)and(163.2,-111.7)..(164.1,-111.1)
				..controls(164.5,-110.8)and(165.3,-110.4)..(165.9,-110)
				..controls(166.5,-109.7)and(167.3,-109.2)..(167.6,-108.9)
				..controls(167.9,-108.6)and(168.6,-108.2)..(169,-108)
				..controls(169.4,-107.8)and(170.2,-107.3)..(170.7,-106.9)
				..controls(171.2,-106.5)and(172.1,-105.8)..(172.7,-105.4)
				..controls(173.2,-105)and(174.1,-104.4)..(174.5,-104)
				..controls(174.9,-103.7)and(175.7,-103.1)..(176.2,-102.6)
				..controls(176.7,-102.2)and(177.4,-101.5)..(177.9,-101.2)
				..controls(178.3,-100.8)and(178.8,-100.4)..(179,-100.2)
				..controls(179.2,-100)and(179.6,-99.6)..(180,-99.3)
				..controls(180.4,-99.1)and(180.9,-98.6)..(181.2,-98.4)
				..controls(181.5,-98.1)and(182.1,-97.7)..(182.5,-97.5)
				..controls(183.2,-97.2)and(183.5,-96.9)..(184.3,-95.9)
				..controls(184.6,-95.6)and(184.6,-95.6)..(184.8,-95.7)
				..controls(185,-95.8)and(185.1,-96)..(185.1,-96.2)
				..controls(185.1,-96.5)and(185,-96.7)..(184.6,-97.2)
				..controls(184.1,-97.8)and(184,-97.9)..(183.9,-98.8)
				..controls(183.7,-100.3)and(183.3,-101.3)..(182.9,-101.8)
				..controls(182.6,-102.1)and(182.3,-102.6)..(182.2,-102.9)
				..controls(182,-103.3)and(181.6,-104)..(181.2,-104.6)
				..controls(180.9,-105.2)and(180.5,-106.1)..(180.4,-106.6)
				..controls(179.9,-108)and(179.7,-108.5)..(179,-109.4)
				..controls(178.6,-109.8)and(178.2,-110.5)..(177.9,-111)
				..controls(177.7,-111.5)and(177.3,-112.3)..(177,-112.7)
				..controls(176.6,-113.2)and(176.4,-113.7)..(176.3,-114.2)
				..controls(176.2,-114.7)and(175.9,-115.2)..(175.4,-116)
				..controls(175,-116.6)and(174.4,-117.7)..(174,-118.5)
				..controls(173.7,-119.2)and(173.3,-120.1)..(173,-120.5)
				..controls(172.8,-120.9)and(172.4,-121.8)..(172.1,-122.5)
				..controls(171.8,-123.2)and(171.3,-124.1)..(171,-124.6)
				..controls(170.8,-125.1)and(170.3,-126)..(169.9,-126.6)
				..controls(169.6,-127.3)and(169.2,-128.1)..(169,-128.3)
				..controls(168.8,-128.6)and(168.5,-129.3)..(168.3,-129.8)
				..controls(168.1,-130.3)and(167.8,-131.1)..(167.5,-131.6)
				..controls(166.9,-132.5)and(164.2,-137.9)..(164.2,-138.1)
				..controls(164.2,-138.1)and(163.9,-138.8)..(163.4,-139.7)
				..controls(162.9,-140.5)and(162.4,-141.5)..(162.3,-141.9)
				..controls(161.9,-143)and(161.6,-143.4)..(160.9,-144.1)
				..controls(160.4,-144.6)and(160.3,-144.8)..(160.1,-145.4)
				..controls(160,-145.8)and(159.6,-146.8)..(159.2,-147.8)
				..controls(158.7,-148.8)and(158.1,-150)..(157.9,-150.6)
				..controls(157.6,-151.1)and(157.2,-151.9)..(156.9,-152.4)
				..controls(156.5,-152.8)and(156.1,-153.7)..(155.8,-154.4)
				..controls(155.6,-155)and(155.2,-155.9)..(154.9,-156.3)
				..controls(154.7,-156.7)and(154.4,-157.4)..(154.2,-157.9)
				..controls(154.1,-158.4)and(153.8,-159)..(153.5,-159.3)
				..controls(152.8,-160.1)and(153,-160.1)..(142.4,-160.1)
				..controls(137.2,-160.1)and(131.8,-160.1)..(130.4,-160)--cycle
				(141.5,-103.7)..controls(140.8,-103)and(139.8,-102.1)..(139.2,-101.6)
				..controls(138.5,-101.2)and(137.9,-100.6)..(137.8,-100.4)
				..controls(137,-99.5)and(134,-96.5)..(133.2,-96)
				..controls(132.7,-95.7)and(131.8,-94.9)..(131.1,-94.3)
				..controls(129.4,-92.7)and(127.9,-91.6)..(126.6,-91.1)
				..controls(125.2,-90.5)and(124.1,-90)..(122.9,-89.1)
				..controls(122.1,-88.6)and(121.3,-88.2)..(119.9,-87.7)
				..controls(118.8,-87.3)and(117.5,-86.8)..(116.9,-86.5)
				..controls(114.3,-85.5)and(113.7,-85.4)..(109.7,-85.5)
				..controls(105.3,-85.7)and(103.9,-85.9)..(101.6,-87)
				..controls(100.7,-87.4)and(99.4,-87.8)..(97.9,-88)
				..controls(97,-88.2)and(96.6,-88.3)..(95.8,-88.9)
				..controls(95.4,-89.2)and(94.7,-89.6)..(94.2,-89.8)
				..controls(92.9,-90.3)and(92.3,-90.7)..(90.9,-92.1)
				..controls(90.2,-92.7)and(89.4,-93.3)..(89,-93.5)..controls(87,-94.5)and(85.2,-96)..(83.6,-97.9)
				..controls(83.3,-98.2)and(82.5,-98.9)..(81.8,-99.4)
				..controls(80.7,-100.1)and(80.5,-100.3)..(80,-101.1)
				..controls(78.7,-103.1)and(77.2,-103.9)..(76.4,-103)
				..controls(75.7,-102.2)and(75.7,-103.2)..(75.7,-83.6)
				..controls(75.6,-73.6)and(75.6,-64.8)..(75.5,-64)
				..controls(75.5,-63.2)and(75.5,-61.4)..(75.6,-59.8)--(75.7,-57.1)--(76,-56.8)
				..controls(76.8,-56.2)and(78,-56.3)..(79.8,-57.2)
				..controls(80.4,-57.5)and(81.3,-57.8)..(81.8,-57.9)
				..controls(82.3,-58.1)and(83,-58.4)..(83.6,-58.8)..controls(84.3,-59.2)and(85,-59.5)..(86,-59.8)
				..controls(86.8,-60.1)and(88,-60.6)..(88.7,-60.9)
				..controls(89.3,-61.2)and(90.6,-61.7)..(91.4,-62)
				..controls(92.2,-62.3)and(93.7,-62.9)..(94.6,-63.3)
				..controls(95.6,-63.7)and(96.8,-64.1)..(97.4,-64.3)
				..controls(98,-64.5)and(99,-64.9)..(99.7,-65.2)
				..controls(100.4,-65.6)and(101.5,-66)..(102.3,-66.2)
				..controls(104.2,-66.8)and(105,-67.1)..(106.2,-67.8)
				..controls(107.9,-68.7)and(110.2,-69.4)..(111.3,-69.2)
				..controls(112.5,-69)and(115,-68.2)..(115.6,-67.8)
				..controls(117.3,-66.7)and(117.6,-66.6)..(119.5,-66.2)
				..controls(121.9,-65.8)and(122.1,-65.7)..(123,-65.1)
				..controls(124.1,-64.4)and(124.4,-64.3)..(126,-63.9)
				..controls(127,-63.7)and(127.7,-63.4)..(128.4,-63.1)
				..controls(129,-62.8)and(129.7,-62.5)..(130.6,-62.4)
				..controls(132.1,-62.1)and(134,-61.5)..(135.2,-61)
				..controls(135.7,-60.8)and(136.7,-60.5)..(137.3,-60.3)
				..controls(138.4,-60)and(139.3,-59.7)..(141.1,-59)
				..controls(141.6,-58.7)and(142,-58.6)..(142.5,-58.6)
				..controls(143.2,-58.6)and(143.3,-58.7)..(143.6,-59)
				..controls(144,-59.5)and(144.1,-59.9)..(144,-60.8)
				..controls(144,-61.2)and(144,-71.1)..(144,-82.8)--(143.9,-104.1)--(143.7,-104.5)
				..controls(143.5,-104.7)and(143.3,-104.8)..(143,-104.9)
				..controls(142.6,-104.9)and(142.5,-104.8)..(141.5,-103.7)--cycle
				(27.8,-91.4)..controls(27.1,-90.9)and(26.3,-90)..(26.1,-89.5)
				..controls(26,-89.2)and(25.8,-88.6)..(25.7,-87.9)
				..controls(25.4,-86.2)and(25.6,-84.8)..(26.5,-83.5)
				..controls(27.4,-82.3)and(28.3,-82)..(29.4,-82.3)
				..controls(31.1,-82.8)and(32.1,-84.9)..(32,-87.4)
				..controls(31.9,-89.3)and(31.2,-90.6)..(29.9,-91.4)
				..controls(29.1,-91.9)and(28.6,-91.9)..(27.8,-91.4)--cycle
				(190.1,-91.3)..controls(189.7,-91.2)and(189.3,-90.9)..(189.1,-90.7)
				..controls(188.3,-89.8)and(187.6,-87.7)..(187.6,-86.2)
				..controls(187.6,-85.5)and(187.7,-85.2)..(188.2,-84.4)
				..controls(189.1,-82.6)and(190.1,-81.8)..(191.3,-82.1)
				..controls(192,-82.3)and(193.5,-83.7)..(193.8,-84.4)
				..controls(194.1,-85.1)and(194.2,-87.1)..(194,-88.3)
				..controls(193.8,-89.8)and(193.3,-90.5)..(192.3,-91.3)
				..controls(191.7,-91.7)and(190.9,-91.7)..(190.1,-91.3)--cycle
				(173.7,-83.4)..controls(172.2,-82.5)and(172,-82.5)..(170.6,-82.4)
				..controls(169.7,-82.4)and(168.8,-82.4)..(168.3,-82.5)
				..controls(167.2,-82.7)and(166.9,-82.7)..(166.6,-82.2)
				..controls(166.3,-81.8)and(166.3,-81)..(166.8,-80)
				..controls(167.2,-78.9)and(167.4,-77.4)..(167.7,-74.1)
				..controls(167.8,-72.8)and(167.9,-71.8)..(167.9,-71.7)
				..controls(168,-71.7)and(168,-70.2)..(168,-68.4)..controls(168,-65.4)and(168,-64.9)..(167.7,-63)
				..controls(167.3,-60.5)and(167.1,-59.9)..(166.4,-58.4)
				..controls(166.1,-57.7)and(165.8,-56.9)..(165.7,-56.6)
				..controls(165.7,-56.2)and(165.4,-55.6)..(165.2,-55.2)
				..controls(164.6,-53.8)and(161.9,-50.7)..(161,-50.3)
				..controls(159.8,-49.7)and(157.9,-50)..(155.2,-51.2)
				..controls(154.7,-51.4)and(153.4,-51.8)..(152.2,-52.1)
				..controls(150.8,-52.4)and(149.4,-52.9)..(148.6,-53.2)
				..controls(147.8,-53.6)and(146.7,-54)..(146,-54.2)
				..controls(145.4,-54.4)and(144.5,-54.7)..(144,-54.9)
				..controls(143.6,-55.2)and(142.7,-55.5)..(141.8,-55.7)
				..controls(141,-55.9)and(139.6,-56.3)..(138.7,-56.6)
				..controls(137.8,-57)and(136.3,-57.4)..(135.3,-57.6)
				..controls(133.8,-57.9)and(133.2,-58.1)..(132.1,-58.7)
				..controls(131.1,-59.2)and(130.3,-59.4)..(129.4,-59.6)
				..controls(128.7,-59.8)and(127.8,-60)..(127.4,-60.1)
				..controls(127,-60.3)and(126,-60.5)..(125.1,-60.7)
				..controls(123.2,-61.1)and(121.7,-61.5)..(120.2,-62.1)
				..controls(119.7,-62.3)and(118.1,-62.8)..(116.7,-63.2)
				..controls(114.6,-63.9)and(114,-64.1)..(113,-64.7)
				..controls(111.9,-65.4)and(111.8,-65.4)..(111,-65.5)
				..controls(109.9,-65.5)and(109.3,-65.4)..(107.7,-64.7)
				..controls(107,-64.4)and(106,-64)..(105.3,-63.9)
				..controls(103.5,-63.5)and(102.2,-63)..(101.2,-62.5)
				..controls(100.1,-61.8)and(99.3,-61.5)..(97.8,-61.1)
				..controls(97.2,-61)and(96.2,-60.7)..(95.7,-60.4)
				..controls(95.1,-60.1)and(94.2,-59.8)..(93.5,-59.6)
				..controls(92.9,-59.4)and(91.7,-59)..(91,-58.7)..controls(90.2,-58.4)and(88.9,-58)..(88,-57.7)
				..controls(87.2,-57.5)and(86.1,-57.1)..(85.5,-56.8)
				..controls(85,-56.6)and(84.1,-56.2)..(83.5,-56)..controls(82.9,-55.7)and(81.8,-55.2)..(81,-54.9)
				..controls(80.1,-54.5)and(79.1,-54.2)..(78.7,-54.1)
				..controls(77.9,-53.8)and(77.3,-53.6)..(76.2,-52.9)
				..controls(75.7,-52.6)and(75.2,-52.4)..(74.1,-52.2)
				..controls(72.5,-51.9)and(72.4,-51.8)..(71.3,-51.1)
				..controls(69.9,-50.1)and(69.4,-49.9)..(67.9,-49.8)
				..controls(67,-49.7)and(66.6,-49.6)..(66.3,-49.4)
				..controls(66.1,-49.3)and(65.7,-49)..(65.4,-48.8)
				..controls(65.1,-48.7)and(64.2,-48.4)..(63.6,-48.2)
				..controls(62.9,-48)and(61.7,-47.6)..(60.9,-47.2)
				..controls(60.1,-46.8)and(58.9,-46.4)..(58.2,-46.2)
				..controls(57.6,-46)and(56.8,-45.7)..(56.5,-45.5)
				..controls(54.9,-44.6)and(54.6,-44.5)..(53.2,-44.2)
				..controls(51.7,-43.8)and(51,-43.6)..(49.9,-42.9)
				..controls(49.5,-42.6)and(48.9,-42.4)..(47.5,-42)
				..controls(45.7,-41.6)and(45.6,-41.5)..(44.3,-40.6)
				..controls(43.5,-40.1)and(42.8,-39.6)..(42.7,-39.4)
				..controls(42.4,-38.7)and(43,-38.1)..(44.5,-37.6)..controls(44.9,-37.5)and(46,-37)..(46.9,-36.6)
				..controls(47.9,-36.1)and(48.8,-35.8)..(49.6,-35.6)
				..controls(50.5,-35.4)and(51.3,-35.2)..(52.4,-34.6)
				..controls(53.3,-34.1)and(54.5,-33.7)..(54.9,-33.5)
				..controls(55.4,-33.4)and(56.4,-33)..(57.2,-32.6)
				..controls(58.6,-32)and(60.1,-31.5)..(61.6,-31.2)
				..controls(62.3,-31.1)and(62.7,-31)..(63.2,-30.7)
				..controls(64.1,-30)and(64.3,-29.9)..(65.3,-29.7)
				..controls(65.7,-29.6)and(66.5,-29.2)..(67.1,-28.9)
				..controls(67.6,-28.6)and(68.8,-28.2)..(69.6,-27.9)
				..controls(70.5,-27.7)and(71.6,-27.3)..(72.1,-27)
				..controls(72.6,-26.7)and(73.8,-26.2)..(74.8,-25.9)
				..controls(75.7,-25.5)and(76.7,-25.1)..(77.1,-25)
				..controls(77.4,-24.8)and(78.5,-24.4)..(79.6,-24.1)
				..controls(80.6,-23.8)and(82,-23.2)..(82.8,-22.9)
				..controls(83.5,-22.5)and(84.8,-22)..(85.8,-21.8)
				..controls(86.7,-21.5)and(87.8,-21.1)..(88.1,-20.9)
				..controls(88.4,-20.7)and(89.3,-20.3)..(90.1,-20.1)
				..controls(90.8,-19.8)and(91.9,-19.4)..(92.4,-19.1)
				..controls(93.5,-18.5)and(94.4,-18.2)..(95.4,-18.1)
				..controls(95.7,-18)and(96.4,-17.8)..(96.9,-17.5)
				..controls(97.5,-17.3)and(98.4,-16.9)..(99,-16.6)
				..controls(99.7,-16.3)and(100.9,-15.9)..(101.9,-15.6)
				..controls(102.9,-15.3)and(104.4,-14.7)..(105.3,-14.3)
				..controls(108.1,-13.2)and(108.5,-13.1)..(110.1,-13.1)
				..controls(111.4,-13.1)and(111.5,-13.1)..(112.4,-13.5)
				..controls(113.2,-14)and(115,-14.8)..(116.6,-15.4)
				..controls(116.9,-15.5)and(117.9,-15.8)..(118.8,-16.1)
				..controls(120.9,-16.6)and(121.5,-16.9)..(122.6,-17.9)
				..controls(123,-18.3)and(123.7,-18.8)..(123.9,-18.9)
				..controls(124.6,-19.4)and(126,-19.8)..(127.9,-20.2)
				..controls(129.6,-20.5)and(129.8,-20.6)..(131.1,-21.4)
				..controls(131.5,-21.7)and(132.4,-22.1)..(133.1,-22.3)
				..controls(136.1,-23.2)and(138.3,-24.1)..(140.2,-25)
				..controls(141,-25.5)and(141.5,-25.6)..(142.5,-25.8)
				..controls(145.1,-26.3)and(146,-26.6)..(149.4,-28.1)
				..controls(151.3,-28.9)and(153.5,-29.8)..(154.3,-30.1)
				..controls(155.1,-30.4)and(156.2,-30.8)..(156.6,-31)
				..controls(157.1,-31.2)and(158.2,-31.6)..(159.1,-31.8)
				..controls(160.6,-32.2)and(160.9,-32.3)..(161.7,-32.9)
				..controls(162.4,-33.3)and(162.9,-33.6)..(163.7,-33.8)
				..controls(164.3,-34)and(164.8,-34.1)..(164.9,-34.1)
				..controls(165,-34.1)and(165.6,-34.4)..(166.4,-34.8)
				..controls(167.3,-35.2)and(168.1,-35.4)..(169,-35.6)
				..controls(170,-35.8)and(170.5,-36)..(171.7,-36.6)
				..controls(172.4,-37)and(173.3,-37.4)..(173.6,-37.4)
				..controls(175.6,-37.7)and(176.2,-38)..(176.7,-39.1)
				..controls(176.9,-39.6)and(176.9,-39.7)..(176.7,-40.4)
				..controls(176.4,-41.6)and(175.8,-42.7)..(175.2,-43.4)
				..controls(174.5,-44.1)and(172.6,-45.5)..(171.7,-45.8)
				..controls(170.6,-46.2)and(169.1,-46)..(168.2,-45.4)
				..controls(167.9,-45.2)and(167.3,-44.8)..(166.9,-44.4)
				..controls(166.4,-44.1)and(165.8,-43.6)..(165.4,-43.2)
				..controls(165.1,-42.9)and(164.4,-42.3)..(163.9,-41.9)
				..controls(163.3,-41.5)and(162.7,-40.9)..(162.4,-40.5)
				..controls(161.4,-39.2)and(159.9,-38)..(158.8,-37.8)
				..controls(158.4,-37.8)and(158.3,-37.7)..(158.2,-37.2)
				..controls(158,-36.8)and(157.4,-36.1)..(157.1,-36.1)
				..controls(157,-36.1)and(156.7,-36.1)..(156.5,-36)
				..controls(156.1,-35.9)and(155.9,-35.7)..(155.7,-35.2)
				..controls(155.3,-34.6)and(155.1,-34.4)..(154,-34.1)
				..controls(153.6,-34)and(153.2,-33.7)..(152.9,-33.5)
				..controls(152.4,-32.9)and(151.6,-32.4)..(151.3,-32.4)
				..controls(150.9,-32.4)and(150.4,-32.6)..(150.4,-32.9)
				..controls(150.4,-33.2)and(150.8,-33.4)..(151.4,-33.6)
				..controls(152.1,-33.7)and(152.8,-34.1)..(153.5,-34.9)
				..controls(154.1,-35.5)and(154.3,-35.6)..(155,-35.8)
				..controls(155.7,-36)and(155.7,-36)..(156,-36.6)
				..controls(156.3,-37.3)and(156.7,-37.7)..(157.5,-37.8)
				..controls(158.1,-37.9)and(158.2,-38)..(158.4,-38.8)
				..controls(158.5,-39.4)and(158.6,-39.6)..(159.6,-40.4)
				..controls(160.2,-41)and(160.9,-41.6)..(161.3,-41.8)
				..controls(161.7,-42)and(161.9,-42.3)..(162.2,-42.8)
				..controls(162.4,-43.2)and(162.8,-43.8)..(163.1,-44.1)
				..controls(163.9,-44.8)and(165.1,-46.7)..(165.7,-48)
				..controls(166.2,-49.3)and(166.4,-49.7)..(167.2,-50.5)
				..controls(167.5,-50.8)and(167.9,-51.4)..(168.1,-51.8)
				..controls(168.3,-52.2)and(168.8,-53)..(169.1,-53.6)
				..controls(169.6,-54.3)and(169.8,-54.9)..(169.9,-55.6)
				..controls(170.4,-57.8)and(170.5,-58.1)..(171.2,-59.5)
				..controls(171.6,-60.3)and(172,-61.5)..(172.2,-62.1)
				..controls(172.3,-62.8)and(172.7,-63.7)..(172.9,-64.1)
				..controls(173.2,-64.6)and(173.4,-65.3)..(173.4,-65.6)
				..controls(173.8,-68.8)and(174.2,-70.6)..(174.9,-72.6)
				..controls(175.5,-74.5)and(175.9,-77.1)..(176.1,-80.2)
				..controls(176.3,-83.4)and(175.5,-84.4)..(173.7,-83.4)--cycle;
				\definecolor{ezs}{RGB}{247,179,26}\fill[ezs](-5.7,-156.6)--(-3.4,-157.5)--(61.9,12.7)--(59.7,13.6)--cycle;
				\definecolor{ezs}{RGB}{247,179,22}\fill[ezs](-3.8,-157.3)--(-1.5,-158.2)--(63.9,12)--(61.6,12.9)--cycle;
				\definecolor{ezs}{RGB}{248,179,19}\fill[ezs](-1.8,-158.1)--(0.4,-158.9)--(65.8,11.2)--(63.5,12.1)--cycle;
				\definecolor{ezs}{RGB}{248,178,16}\fill[ezs](0.1,-158.8)--(2.4,-159.7)--(67.7,10.5)--(65.4,11.4)--cycle;
				\definecolor{ezs}{RGB}{249,178,15}\fill[ezs](2,-159.5)--(4.3,-160.4)--(69.6,9.8)--(67.3,10.7)--cycle;
				\definecolor{ezs}{RGB}{251,177,17}\fill[ezs](3.9,-160.3)--(6.2,-161.1)--(71.5,9)--(69.2,9.9)--cycle;
				\definecolor{ezs}{RGB}{252,179,18}\fill[ezs](5.8,-161)--(8.1,-161.9)--(73.4,8.3)--(71.1,9.2)--cycle;
				\definecolor{ezs}{RGB}{252,182,17}\fill[ezs](7.7,-161.7)--(10,-162.6)--(75.3,7.6)--(73,8.5)--cycle;
				\definecolor{ezs}{RGB}{252,185,17}\fill[ezs](9.6,-162.5)--(11.9,-163.4)--(77.2,6.8)--(75,7.7)--cycle;
				\definecolor{ezs}{RGB}{251,189,18}\fill[ezs](11.5,-163.2)--(13.8,-164.1)--(79.2,6.1)--(76.9,7)--cycle;
				\definecolor{ezs}{RGB}{251,191,22}\fill[ezs](13.4,-163.9)--(15.7,-164.8)--(81.1,5.4)--(78.8,6.2)--cycle;
				\definecolor{ezs}{RGB}{252,193,25}\fill[ezs](15.4,-164.7)--(17.6,-165.6)--(83,4.6)--(80.7,5.5)--cycle;
				\definecolor{ezs}{RGB}{253,195,27}\fill[ezs](17.3,-165.4)--(19.6,-166.3)--(84.9,3.9)--(82.6,4.8)--cycle;
				\definecolor{ezs}{RGB}{253,196,31}\fill[ezs](19.2,-166.1)--(21.5,-167)--(86.8,3.2)--(84.5,4)--cycle;
				\definecolor{ezs}{RGB}{254,199,35}\fill[ezs](21.1,-166.9)--(23.4,-167.8)--(88.7,2.4)--(86.4,3.3)--cycle;
				\definecolor{ezs}{RGB}{253,201,38}\fill[ezs](23,-167.6)--(25.3,-168.5)--(90.6,1.7)--(88.3,2.6)--cycle;
				\definecolor{ezs}{RGB}{253,203,40}\fill[ezs](24.9,-168.3)--(27.2,-169.2)--(92.5,1)--(90.3,1.8)--cycle;
				\definecolor{ezs}{RGB}{253,205,42}\fill[ezs](26.8,-169.1)--(29.1,-170)--(94.4,0.2)--(92.2,1.1)--cycle;
				\definecolor{ezs}{RGB}{253,207,45}\fill[ezs](28.7,-169.8)--(31,-170.7)--(96.4,-0.5)--(94.1,0.4)--cycle;
				\definecolor{ezs}{RGB}{253,210,49}\fill[ezs](30.7,-170.5)--(32.9,-171.4)--(98.3,-1.2)--(96,-0.4)--cycle;
				\definecolor{ezs}{RGB}{253,212,54}\fill[ezs](32.6,-171.3)--(34.9,-172.2)--(100.2,-2)--(97.9,-1.1)--cycle;
				\definecolor{ezs}{RGB}{253,214,59}\fill[ezs](34.5,-172)--(36.8,-172.9)--(102.1,-2.7)--(99.8,-1.8)--cycle;
				\definecolor{ezs}{RGB}{253,216,63}\fill[ezs](36.4,-172.7)--(38.7,-173.6)--(104,-3.4)--(101.7,-2.6)--cycle;
				\definecolor{ezs}{RGB}{253,218,65}\fill[ezs](38.3,-173.5)--(40.6,-174.4)--(105.9,-4.2)--(103.6,-3.3)--cycle;
				\definecolor{ezs}{RGB}{254,220,66}\fill[ezs](40.2,-174.2)--(42.5,-175.1)--(107.8,-4.9)--(105.5,-4)--cycle;
				\definecolor{ezs}{RGB}{254,222,68}\fill[ezs](42.1,-175)--(44.4,-175.8)--(109.7,-5.6)--(107.5,-4.8)--cycle;
				\definecolor{ezs}{RGB}{254,224,71}\fill[ezs](44,-175.7)--(46.3,-176.6)--(111.7,-6.4)--(109.4,-5.5)--cycle;
				\definecolor{ezs}{RGB}{254,226,75}\fill[ezs](46,-176.4)--(48.2,-177.3)--(113.6,-7.1)--(111.3,-6.2)--cycle;
				\definecolor{ezs}{RGB}{254,228,79}\fill[ezs](47.9,-177.2)--(50.1,-178)--(115.5,-7.8)--(113.2,-7)--cycle;
				\definecolor{ezs}{RGB}{254,230,82}\fill[ezs](49.8,-177.9)--(52.1,-178.8)--(117.4,-8.6)--(115.1,-7.7)--cycle;
				\definecolor{ezs}{RGB}{254,232,85}\fill[ezs](51.7,-178.6)--(54,-179.5)--(119.3,-9.3)--(117,-8.4)--cycle;
				\definecolor{ezs}{RGB}{254,234,88}\fill[ezs](53.6,-179.4)--(55.9,-180.2)--(121.2,-10)--(118.9,-9.2)--cycle;
				\definecolor{ezs}{RGB}{254,237,91}\fill[ezs](55.5,-180.1)--(57.8,-181)--(123.1,-10.8)--(120.8,-9.9)--cycle;
				\definecolor{ezs}{RGB}{254,239,94}\fill[ezs](57.4,-180.8)--(59.7,-181.7)--(125,-11.5)--(122.8,-10.6)--cycle;
				\definecolor{ezs}{RGB}{254,241,97}\fill[ezs](59.3,-181.6)--(61.6,-182.4)--(127,-12.2)--(124.7,-11.4)--cycle;
				\definecolor{ezs}{RGB}{254,243,101}\fill[ezs](61.2,-182.3)--(63.5,-183.2)--(128.9,-13)--(126.6,-12.1)--cycle;
				\definecolor{ezs}{RGB}{254,245,103}\fill[ezs](63.2,-183)--(65.4,-183.9)--(130.8,-13.7)--(128.5,-12.8)--cycle;
				\definecolor{ezs}{RGB}{254,245,105}\fill[ezs](65.1,-183.8)--(67.4,-184.6)--(132.7,-14.4)--(130.4,-13.6)--cycle;
				\definecolor{ezs}{RGB}{254,245,105}\fill[ezs](67,-184.5)--(69.3,-185.4)--(134.6,-15.2)--(132.3,-14.3)--cycle;
				\definecolor{ezs}{RGB}{254,244,103}\fill[ezs](68.9,-185.2)--(71.2,-186.1)--(136.5,-15.9)--(134.2,-15)--cycle;
				\definecolor{ezs}{RGB}{253,243,101}\fill[ezs](70.8,-186)--(73.1,-186.8)--(138.4,-16.6)--(136.1,-15.8)--cycle;
				\definecolor{ezs}{RGB}{253,241,98}\fill[ezs](72.7,-186.7)--(75,-187.6)--(140.3,-17.4)--(138,-16.5)--cycle;
				\definecolor{ezs}{RGB}{253,239,95}\fill[ezs](74.6,-187.4)--(76.9,-188.3)--(142.2,-18.1)--(140,-17.2)--cycle;
				\definecolor{ezs}{RGB}{253,236,92}\fill[ezs](76.5,-188.2)--(78.8,-189)--(144.2,-18.8)--(141.9,-18)--cycle;
				\definecolor{ezs}{RGB}{253,234,88}\fill[ezs](78.5,-188.9)--(80.7,-189.8)--(146.1,-19.6)--(143.8,-18.7)--cycle;
				\definecolor{ezs}{RGB}{253,232,85}\fill[ezs](80.4,-189.6)--(82.7,-190.5)--(148,-20.3)--(145.7,-19.4)--cycle;
				\definecolor{ezs}{RGB}{253,230,82}\fill[ezs](82.3,-190.4)--(84.6,-191.2)--(149.9,-21)--(147.6,-20.2)--cycle;
				\definecolor{ezs}{RGB}{253,228,78}\fill[ezs](84.2,-191.1)--(86.5,-192)--(151.8,-21.8)--(149.5,-20.9)--cycle;
				\definecolor{ezs}{RGB}{253,226,74}\fill[ezs](86.1,-191.8)--(88.4,-192.7)--(153.7,-22.5)--(151.4,-21.6)--cycle;
				\definecolor{ezs}{RGB}{253,224,70}\fill[ezs](88,-192.6)--(90.3,-193.4)--(155.6,-23.3)--(153.3,-22.4)--cycle;
				\definecolor{ezs}{RGB}{253,222,67}\fill[ezs](89.9,-193.3)--(92.2,-194.2)--(157.5,-24)--(155.3,-23.1)--cycle;
				\definecolor{ezs}{RGB}{253,220,65}\fill[ezs](91.8,-194)--(94.1,-194.9)--(159.5,-24.7)--(157.2,-23.8)--cycle;
				\definecolor{ezs}{RGB}{253,217,63}\fill[ezs](93.7,-194.8)--(96,-195.6)--(161.4,-25.5)--(159.1,-24.6)--cycle;
				\definecolor{ezs}{RGB}{253,215,60}\fill[ezs](95.7,-195.5)--(97.9,-196.4)--(163.3,-26.2)--(161,-25.3)--cycle;
				\definecolor{ezs}{RGB}{253,213,56}\fill[ezs](97.6,-196.2)--(99.9,-197.1)--(165.2,-26.9)--(162.9,-26)--cycle;
				\definecolor{ezs}{RGB}{252,211,53}\fill[ezs](99.5,-197)--(101.8,-197.8)--(167.1,-27.7)--(164.8,-26.8)--cycle;
				\definecolor{ezs}{RGB}{252,209,49}\fill[ezs](101.4,-197.7)--(103.7,-198.6)--(169,-28.4)--(166.7,-27.5)--cycle;
				\definecolor{ezs}{RGB}{252,207,45}\fill[ezs](103.3,-198.4)--(105.6,-199.3)--(170.9,-29.1)--(168.6,-28.2)--cycle;
				\definecolor{ezs}{RGB}{253,205,42}\fill[ezs](105.2,-199.2)--(107.5,-200)--(172.8,-29.9)--(170.6,-29)--cycle;
				\definecolor{ezs}{RGB}{253,202,40}\fill[ezs](107.1,-199.9)--(109.4,-200.8)--(174.7,-30.6)--(172.5,-29.7)--cycle;
				\definecolor{ezs}{RGB}{253,200,38}\fill[ezs](109,-200.6)--(111.3,-201.5)--(176.7,-31.3)--(174.4,-30.4)--cycle;
				\definecolor{ezs}{RGB}{253,198,35}\fill[ezs](111,-201.4)--(113.2,-202.2)--(178.6,-32.1)--(176.3,-31.2)--cycle;
				\definecolor{ezs}{RGB}{253,196,31}\fill[ezs](112.9,-202.1)--(115.2,-203)--(180.5,-32.8)--(178.2,-31.9)--cycle;
				\definecolor{ezs}{RGB}{253,194,28}\fill[ezs](114.8,-202.8)--(117.1,-203.7)--(182.4,-33.5)--(180.1,-32.6)--cycle;
				\definecolor{ezs}{RGB}{253,192,26}\fill[ezs](116.7,-203.6)--(119,-204.5)--(184.3,-34.3)--(182,-33.4)--cycle;
				\definecolor{ezs}{RGB}{253,189,24}\fill[ezs](118.6,-204.3)--(120.9,-205.2)--(186.2,-35)--(183.9,-34.1)--cycle;
				\definecolor{ezs}{RGB}{253,186,20}\fill[ezs](120.5,-205)--(122.8,-205.9)--(188.1,-35.7)--(185.8,-34.8)--cycle;
				\definecolor{ezs}{RGB}{253,184,17}\fill[ezs](122.4,-205.8)--(124.7,-206.7)--(190,-36.5)--(187.8,-35.6)--cycle;
				\definecolor{ezs}{RGB}{252,183,15}\fill[ezs](124.3,-206.5)--(126.6,-207.4)--(192,-37.2)--(189.7,-36.3)--cycle;
				\definecolor{ezs}{RGB}{252,181,14}\fill[ezs](126.3,-207.2)--(128.5,-208.1)--(193.9,-37.9)--(191.6,-37.1)--cycle;
				\definecolor{ezs}{RGB}{252,179,13}\fill[ezs](128.2,-208)--(130.4,-208.9)--(195.8,-38.7)--(193.5,-37.8)--cycle;
				\definecolor{ezs}{RGB}{252,178,12}\fill[ezs](130.1,-208.7)--(132.4,-209.6)--(197.7,-39.4)--(195.4,-38.5)--cycle;
				\definecolor{ezs}{RGB}{251,178,10}\fill[ezs](132,-209.4)--(134.3,-210.3)--(199.6,-40.1)--(197.3,-39.3)--cycle;
				\definecolor{ezs}{RGB}{251,178,10}\fill[ezs](133.9,-210.2)--(136.2,-211.1)--(201.5,-40.9)--(199.2,-40)--cycle;
				\definecolor{ezs}{RGB}{250,177,11}\fill[ezs](135.8,-210.9)--(138.1,-211.8)--(203.4,-41.6)--(201.1,-40.7)--cycle;
				\definecolor{ezs}{RGB}{250,178,11}\fill[ezs](137.7,-211.6)--(140,-212.5)--(205.3,-42.3)--(203.1,-41.5)--cycle;
				\definecolor{ezs}{RGB}{248,180,10}\fill[ezs](139.6,-212.4)--(141.9,-213.3)--(207.3,-43.1)--(205,-42.2)--cycle;
				\definecolor{ezs}{RGB}{247,181,9}\fill[ezs](141.5,-213.1)--(143.8,-214)--(209.2,-43.8)--(206.9,-42.9)--cycle;
				\definecolor{ezs}{RGB}{246,182,9}\fill[ezs](143.5,-213.8)--(145.7,-214.7)--(211.1,-44.5)--(208.8,-43.7)--cycle;
				\definecolor{ezs}{RGB}{246,182,10}\fill[ezs](145.4,-214.6)--(147.7,-215.5)--(213,-45.3)--(210.7,-44.4)--cycle;
			\end{scope}
			% chữ EDUZY (Times đậm, làm dày nét bằng 5 lớp lệch nhẹ; căn đúng khung chữ logo gốc)
			\foreach \dx/\dy in {-0.7/0,0.7/0,0/-0.3,0/0.3,0/0}
			\node[anchor=center,inner sep=0pt,outer sep=0pt,text=ezText,align=center] at (109.6+\dx,-193.6+\dy)
			{\resizebox{1.01cm}{0.22cm}{\fontencoding{T1}\fontfamily{ptm}\fontseries{b}\selectfont EDUZY}};
		\end{tikzpicture}%
	}%
}
%====================== THIẾT LẬP THÂN ĐỀ ======================
%--- Môi trường tự luận {bt}: ex_test 3.3.4 chưa khai báo sẵn nên tạo ở đây
%    (ex_test đã tự móc \fixshowans và \listenumerate cho bt)
\makeatletter
\@ifundefined{bt}{\newtheorem{bt}{\color{gold1}\namebt}}{}
\makeatother
\renewcommand{\sectionmark}[1]{\markright{\thesection\ #1}}
%--- Khung tiêu đề (logo + tối đa 4 dòng; dòng 3, 4 để trống {} thì bỏ qua)
%  Đề thi:    \khungtieude{THI THỬ THPT QUỐC GIA NĂM 2026}{Môn Toán -- Đề số 21}
%                         {(Đề thi gồm: ...)}{Thời gian làm bài: 90 phút (không kể thời gian phát đề).}
%  Chuyên đề: \khungtieude{CHUYÊN ĐỀ 3 -- DÃY SỐ VÀ CÔNG THỨC TRUY HỒI}
%                         {Bồi dưỡng học sinh giỏi Toán THPT}{...}{}
\newcommand{\khungtieude}[4]{%
	\begin{tcolorbox}[colframe=gold1,colback=gold5!45,boxrule=0.5mm]
		\begin{minipage}[c]{3cm}
			\centerline{\eduzylogo[2.4cm]}
		\end{minipage}\hspace{0.5cm}
		\begin{minipage}[c]{12.5cm}
			\vskip 0.15 cm
			\centerline{\ind{#1}}
			\vskip 0.15 cm
			\centerline{\textbf{#2}}
			\if\relax\detokenize{#3}\relax\else
				\vskip 0.15 cm
				\centerline{\textit{#3}}
			\fi
			\if\relax\detokenize{#4}\relax\else
				\vskip 0.15 cm
				\centerline{\textit{#4}}
			\fi
			\vskip 0.15 cm
		\end{minipage}
	\end{tcolorbox}
}
%====================== BỔ SUNG CHO BÀI HỌC, CHUYÊN ĐỀ ======================
\usepackage[shortlabels]{enumitem}
\usepackage{tabularx}
%--- Danh sách gọn: \begin{enumerate}[gon] ... \end{enumerate}
\SetEnumitemKey{gon}{leftmargin=*,itemsep=1pt,parsep=0pt,topsep=3pt}
%--- Khung lý thuyết: \begin{khung}{} ... \end{khung}; tiêu đề chỉ dùng khi thật ngắn ("Mẹo", "Lỗi thường gặp")
\newtcolorbox{khung}[1]{breakable,enhanced,colframe=gold1,colback=gold5!45,
	boxrule=0.4mm,fonttitle=\bfseries\sffamily,coltitle=white,colbacktitle=gold1,
	attach boxed title to top left={xshift=6mm,yshift=-3mm},
	boxed title style={arc=1mm,boxrule=0.3mm},title={#1}}
%--- Tiêu đề mục nhỏ trong một dạng: \muc{Dấu hiệu}, \muc{Phương pháp}, \muc{Bài tự luyện}
\makeatletter
\newcommand{\muc}[1]{\par\medskip\noindent\indam{#1}\par\nobreak\smallskip\@afterindentfalse\@afterheading}
\makeatother
%--- Ô công thức cần nhớ (căn giữa, viền vàng, vừa khít nội dung)
%     \begin{center}\congthuc{$Q(n)=(n-1)!$}\end{center}
\newtcbox{\congthuc}{on line,colframe=gold1,colback=gold5!45,boxrule=0.5pt,arc=2pt,
	left=8pt,right=8pt,top=4pt,bottom=4pt}
%--- Công tắc lời giải bài tự luyện (ví dụ mẫu luôn dùng \loigiai)
%     \loigiaiTL[{đáp số}]{lời giải đầy đủ}
%     \TLfull -> in lời giải đầy đủ;  \TLdapso -> chỉ in một dòng "Đáp số";  \TLan -> ẩn hoàn toàn
%     (\CheDo là dethi thì mọi lời giải và đáp số đều ẩn)
\makeatletter
\newif\ifTLhien   \TLhientrue
\newif\ifTLdapso  \TLdapsofalse
\newcommand{\TLfull}{\TLhientrue\TLdapsofalse}
\newcommand{\TLdapso}{\TLhientrue\TLdapsotrue}
\newcommand{\TLan}{\TLhienfalse}
\newcommand{\loigiaiTL}[2][]{%
	\ifTLhien
		\ifTLdapso
			\def\@tlarg{#1}%
			\ifx\@tlarg\@empty\else
				\ifdefstring{\kindOPTN}{dethi}{}{\par\noindent{\bfseries\color{gold1}Đáp số:} #1\par}%
			\fi
		\else
			\loigiai{#2}%
		\fi
	\fi}
\makeatother
\TLdapso% mặc định: chỉ in đáp số bài tự luyện
%====================== BẢNG ĐÚNG/SAI (\choiceTFt) ======================
%  Hàng tiêu đề nền nâu đồng chữ trắng (Phát biểu | Đ | S), các hàng tô xen kẽ kem / trắng,
%  đường kẻ vàng kim, nhãn a) b) c) d) màu nâu đồng. Bản có đáp án: cột được chọn có vòng tròn
%  nâu đồng với dấu tích trắng (câu đúng tích ở cột Đúng, câu sai tích ở cột Sai).
%  Tiêu đề ghi đủ "Đúng | Sai": \OPTN{viettat=0} sau \input (cột Đúng/Sai rộng ra, phát biểu hẹp lại).
%  Dùng bảng gốc của ex_test: \def\BangDSgoc{1} trước \input{preamble-eduzy-v2}.
\providecommand{\BangDSgoc}{0}
\ifnum\BangDSgoc=0
\usepackage{colortbl,pifont}
\colorlet{ezDSke}{gold2!60!gold1}% đường kẻ bảng (nét 0,4pt như ex_test để không đổi chiều cao bảng)
\colorlet{ezDSle}{gold5!60}% hàng a), c)
\colorlet{ezDSchan}{white}% hàng b), d)
%--- dấu chọn ở cột Đúng/Sai
%    (bề rộng 0, căn giữa ô: không làm đổi bề rộng cột của ex_test, không báo tràn lề)
\newcommand{\EZdauchon}{\makebox[0pt]{\tikz[baseline=-1ex]\node[circle,fill=gold1,inner sep=0pt,
	minimum size=1.2em,text=white,font=\small]{\ding{51}};}}
\def\tickT{\EZdauchon}\def\tickF{\EZdauchon}
\AtBeginDocument{\sbox{\ccTFTab}{\EZdauchon}}
%--- nhãn a) b) c) d) màu nâu đồng (ở bản lời giải, nhãn câu đúng vẫn đóng khung như ex_test)
\renewcommand{\FalseTF}{\stepcounter{dapan}\textbf{\color{gold1}\DapAnTF\sepTF}\ignorespaces}
%--- trước khi dựng bảng: kẻ màu vàng kim (bề rộng cột giữ nguyên như ex_test); sau bảng trả lại màu đen
\newcount\EZhangTF
\newcommand{\EZbatdauTF}{%
	\global\EZhangTF=0\relax
	\arrayrulecolor{ezDSke}}
\newcommand{\EZketthucTF}{\arrayrulecolor{black}}
%--- hàng tiêu đề nền nâu đồng, chữ trắng
\makeatletter
\newcommand{\EZo@TF}[1]{\centering\bfseries\color{white}#1}
\newcommand{\EZtieudeTF}{%
	\ifcase\cotDS
		\ifnum\cotSTT=1\relax
			\def\rowHeader{\hline\rowcolor{gold1}\EZo@TF{TT}&\EZo@TF{\phatbieu}\tabularnewline}%
		\else
			\def\rowHeader{\hline\rowcolor{gold1}\EZo@TF{\phatbieu}\tabularnewline}%
		\fi
	\or
		\ifnum\cotSTT=1\relax
			\def\rowHeader{\hline\rowcolor{gold1}\EZo@TF{TT}&\EZo@TF{\phatbieu}&\EZo@TF{\Dung}\tabularnewline}%
		\else
			\def\rowHeader{\hline\rowcolor{gold1}\EZo@TF{\phatbieu}&\EZo@TF{\Dung}\tabularnewline}%
		\fi
	\or
		\ifnum\cotSTT=1\relax
			\def\rowHeader{\hline\rowcolor{gold1}\EZo@TF{TT}&\EZo@TF{\phatbieu}&\EZo@TF{\Dung}&\EZo@TF{\Sai}\tabularnewline}%
		\else
			\def\rowHeader{\hline\rowcolor{gold1}\EZo@TF{\phatbieu}&\EZo@TF{\Dung}&\EZo@TF{\Sai}\tabularnewline}%
		\fi
	\fi}
%--- tô xen kẽ từng hàng (gọi ngay sau \hline đầu mỗi hàng)
\newcommand{\EZmauhangTF}{%
	\noalign{\global\advance\EZhangTF\@ne}%
	\ifodd\EZhangTF\rowcolor{ezDSle}\else\rowcolor{ezDSchan}\fi}
%--- gắn vào \choiceTFt / \writeChoiceTFt của ex_test
\xpretocmd{\choiceTFt}{\EZbatdauTF}{}{\typeout{[preamble-eduzy-v2] Không gắn được bảng Đúng/Sai (đầu bảng)}}
\xpatchcmd{\choiceTFt}{\ifnum\dongPB=0\def\rowHeader{}\fi}{\EZtieudeTF\ifnum\dongPB=0\def\rowHeader{}\fi}
	{}{\typeout{[preamble-eduzy-v2] Không gắn được bảng Đúng/Sai (hàng tiêu đề)}}
\xpatchcmd{\writeChoiceTFt}{\hline\xdef\dotEX{\sdotEX}}{\hline\EZmauhangTF\xdef\dotEX{\sdotEX}}
	{}{\typeout{[preamble-eduzy-v2] Không gắn được bảng Đúng/Sai (tô hàng)}}
\xpatchcmd{\writeChoiceTFt}{{\end{longtable}}}{{\end{longtable}\EZketthucTF}}
	{}{\typeout{[preamble-eduzy-v2] Không gắn được bảng Đúng/Sai (cuối longtable)}}
\xpatchcmd{\writeChoiceTFt}{{\end{xtabular}\end{center}\medskip\noindent}}{{\end{xtabular}\end{center}\EZketthucTF\medskip\noindent}}
	{}{\typeout{[preamble-eduzy-v2] Không gắn được bảng Đúng/Sai (cuối xtabular)}}
\fi
%=====================================================================
\renewcommand{\parallel}{\mathrel{/\!\!/}}

%      \TLdapso                      % tài liệu: \TLfull | \TLdapso | \TLan
%      \begin{document} ... \end{document}
%  Logo Eduzy vẽ bằng TikZ ngay trong file (\eduzylogo), không cần file ảnh.
%  Biên dịch: pdfLaTeX (máy có gói vietnam) hoặc XeLaTeX.
%============================================================
\documentclass[12pt,a4paper,twoside]{article}%<-- twoside: bat dau trang/chan trang chan-le
\usepackage{iftex}
\usepackage{amsmath,amssymb,tikz,mathrsfs,tkz-tab}
\usepackage{graphicx}
\usepackage{pgfplots}% đồ thị hàm số
\pgfplotsset{compat=1.17}
%\usepackage{mathptmx}
\usepackage{mathpazo}
\ifXeTeX% XeLaTeX: font chữ có đủ dấu tiếng Việt, giữ nguyên font toán mathpazo
	\usepackage[no-math]{fontspec}
	\setmainfont{TeX Gyre Pagella}
	\newfontfamily\EZqag{TeX Gyre Adventor}[NFSSFamily=qag,Ligatures=TeX]
\fi
\usepackage[many]{tcolorbox}
\usetikzlibrary{math,through,calc,intersections,angles,quotes,shapes,shapes.geometric,arrows,patterns,matrix,chains,arrows.meta,decorations.shapes,decorations.fractals,decorations.markings,shadows,decorations.pathreplacing}
\usepackage[top=1.6cm, bottom=1.8cm, left=1.8cm, right=1.8cm]{geometry}% v2: top 1.6cm (bản cũ 1.5cm) để chữ không sát dải đầu trang
\providecommand{\CheDo}{loigiai}% loigiai | dethi | solcolor | color | book
\expandafter\usepackage\expandafter[\CheDo]{ex_test}
% \tenchuong: khai báo ở file chính (trước \input) hoặc \renewcommand sau \input
%============================================================
%  [GỘP TỪ eduzy-de.tex] EDUZY - MÀU NÂU ĐỒNG / VÀNG KIM
%  + ĐẦU TRANG, CHÂN TRANG KIỂU DẢI NGANG (v2), ĐỐI XỨNG TRANG CHẴN / TRANG LẺ.
%  (không cần file eduzy-de.tex riêng nữa)
%  Khối này phải đặt SAU \usepackage[...]{ex_test} và SAU \def\tenchuong{...}
%============================================================
\usepackage{lastpage}
\usepackage{needspace}
%--- DÀN TRANG ĐỀ THI: tách câu (THÊM 05/11/2026; chỉ thêm lệnh mới, KHÔNG đổi hành vi mặc định) ---
%  Hai công tắc, mặc định tắt (= preamble chạy y như cũ). Bật bằng cách đặt TRƯỚC %============================================================
%  PREAMBLE EDUZY v2 - dùng chung cho đề thi, bài học, chuyên đề (gói ex_test)
%  Khác bản preamble-eduzy.tex ở đầu trang, chân trang: dải nền nâu đồng tràn ngang,
%  logo + Eduzy.com.vn, ô chữ vàng kem cạnh xiên, chân trang "Trang (số)";
%  logo Eduzy chìm giữa mọi trang (tắt: \logochimtat, bật lại: \logochimbat).
%  Cách dùng ở file chính (xem mau-chuan.tex, mau-chuyen-de.tex):
%      \def\CheDo{loigiai}          % loigiai | dethi | solcolor | color | book
%      \def\tenchuong{<dòng chữ trên đầu trang>}
%      \def\thuonghieu{<chữ ở chân trang lẻ>}   % không khai báo thì dùng mặc định
%      \def\tacgia{<chữ ở chân trang chẵn>}      % không khai báo thì dùng mặc định
%      \def\khauhieu{<chữ cạnh logo ở đầu trang>}  % tùy chọn, mặc định Eduzy.com.vn
%      \input{preamble-eduzy-v2}
%      \TLdapso                      % tài liệu: \TLfull | \TLdapso | \TLan
%      \begin{document} ... \end{document}
%  Logo Eduzy vẽ bằng TikZ ngay trong file (\eduzylogo), không cần file ảnh.
%  Biên dịch: pdfLaTeX (máy có gói vietnam) hoặc XeLaTeX.
%============================================================
\documentclass[12pt,a4paper,twoside]{article}%<-- twoside: bat dau trang/chan trang chan-le
\usepackage{iftex}
\usepackage{amsmath,amssymb,tikz,mathrsfs,tkz-tab}
\usepackage{graphicx}
\usepackage{pgfplots}% đồ thị hàm số
\pgfplotsset{compat=1.17}
%\usepackage{mathptmx}
\usepackage{mathpazo}
\ifXeTeX% XeLaTeX: font chữ có đủ dấu tiếng Việt, giữ nguyên font toán mathpazo
	\usepackage[no-math]{fontspec}
	\setmainfont{TeX Gyre Pagella}
	\newfontfamily\EZqag{TeX Gyre Adventor}[NFSSFamily=qag,Ligatures=TeX]
\fi
\usepackage[many]{tcolorbox}
\usetikzlibrary{math,through,calc,intersections,angles,quotes,shapes,shapes.geometric,arrows,patterns,matrix,chains,arrows.meta,decorations.shapes,decorations.fractals,decorations.markings,shadows,decorations.pathreplacing}
\usepackage[top=1.6cm, bottom=1.8cm, left=1.8cm, right=1.8cm]{geometry}% v2: top 1.6cm (bản cũ 1.5cm) để chữ không sát dải đầu trang
\providecommand{\CheDo}{loigiai}% loigiai | dethi | solcolor | color | book
\expandafter\usepackage\expandafter[\CheDo]{ex_test}
% \tenchuong: khai báo ở file chính (trước \input) hoặc \renewcommand sau \input
%============================================================
%  [GỘP TỪ eduzy-de.tex] EDUZY - MÀU NÂU ĐỒNG / VÀNG KIM
%  + ĐẦU TRANG, CHÂN TRANG KIỂU DẢI NGANG (v2), ĐỐI XỨNG TRANG CHẴN / TRANG LẺ.
%  (không cần file eduzy-de.tex riêng nữa)
%  Khối này phải đặt SAU \usepackage[...]{ex_test} và SAU \def\tenchuong{...}
%============================================================
\usepackage{lastpage}
\usepackage{needspace}
%--- DÀN TRANG ĐỀ THI: tách câu (THÊM 05/11/2026; chỉ thêm lệnh mới, KHÔNG đổi hành vi mặc định) ---
%  Hai công tắc, mặc định tắt (= preamble chạy y như cũ). Bật bằng cách đặt TRƯỚC \input{preamble-eduzy-v2}
%  trong file đề (tools/dan_trang/dan_trang.py tự ghi khi đặt \NeedspaceCau):
%    \def\TachCauDS{1}   câu Đúng/Sai được tách: câu dẫn (kèm hình) ở cuối trang, bảng Phát biểu sang trang sau
%                        (câu dẫn và hình không tách nhau, bảng không tách giữa hai trang)
%    \def\TachCauLG{1}   bản lời giải (\CheDo khác dethi): phần đề của câu cao hơn 1/3 trang được phép tách
%  Lệnh dùng trong file đề (giá trị do tools/dan_trang/dan_trang.py đo ở bản dethi):
%    \NeedspaceCau[<cao câu dẫn + hình>]{<cao cả khối câu>}   ngay trên \begin{ex} (câu đầu phần: trên tiêu đề phần)
%        công tắc tắt: đúng bằng \Needspace{<cao cả khối câu>} như cũ;
%        \TachCauDS=1 và có [..] (câu Đúng/Sai được chọn tách): chỉ giữ câu dẫn + hình;
%        \TachCauLG=1 ở bản lời giải: đòi tối đa 1/3 \textheight.
%    \NeedspaceBang{<cao bảng Phát biểu>}%   ngay trước \choiceTFt: \TachCauDS=1 thì giữ trọn bảng trên một trang;
%        công tắc tắt: không làm gì.
\providecommand{\TachCauDS}{0}
\providecommand{\TachCauLG}{0}
\newlength{\EZnsCau}
\newcommand{\NeedspaceCau}[2][]{%
	\setlength{\EZnsCau}{#2}%
	\ifnum\TachCauDS=1 \if\relax\detokenize{#1}\relax\else\setlength{\EZnsCau}{#1}\fi\fi
	\ifnum\TachCauLG=1 \ifdefstring{\kindOPTN}{dethi}{}{\ifdim\EZnsCau>0.3333\textheight\setlength{\EZnsCau}{0.3333\textheight}\fi}\fi
	\Needspace{\EZnsCau}}
\newcommand{\NeedspaceBang}[1]{\ifnum\TachCauDS=1 \Needspace{#1}\fi}

%--- Thang màu: đậm -> nhạt
\definecolor{gold1}{RGB}{126,96,0}    % nâu đồng đậm  #7E6000
\definecolor{gold2}{RGB}{192,144,3}   % vàng kim      #C09003
\definecolor{gold3}{RGB}{255,215,99}  % vàng nhạt     #FFD763
\definecolor{gold4}{RGB}{255,231,153} % vàng rất nhạt #FFE799
\definecolor{gold5}{RGB}{255,243,205} % vàng kem      #FFF3CD
%--- Ghi đè tông xanh/đỏ cũ
\definecolor{myblue}{RGB}{126,96,0}
\colorlet{blue}{gold1}
\colorlet{cyan}{gold2}
\colorlet{teal}{gold2}

%--- Tên chương mặc định (file đề nên tự khai báo lại)
\providecommand{\tenchuong}{}
%--- Chân trang
\providecommand{\thuonghieu}{TOÁN 12 -- Eduzy.com.vn}
\providecommand{\tacgia}{Tô Gia Bảo -- 0353703783}

%--- Chữ cạnh logo ở đầu trang
\providecommand{\khauhieu}{Eduzy.com.vn}

%====================== THÔNG SỐ ĐẦU TRANG, CHÂN TRANG (v2) ======================
%  Dải nền đậm tràn ngang trang; ô chữ vàng kem có cạnh xiên, kèm 2 vạch xiên chuyển màu.
%  Đầu trang lẻ : logo + \khauhieu (trái)   | ô chữ \tenchuong (phải)
%  Chân trang lẻ: ô chữ \thuonghieu (trái)  | Trang (số) (phải)
%  Trang chẵn soi gương trang lẻ; chân trang chẵn ghi \tacgia.
%--- Màu: chỉ dùng thang vàng kim Eduzy
\colorlet{ezNen}{gold1}% dải nền đậm
\colorlet{ezVach1}{gold3}% vạch xiên sát dải nền (đậm hơn)
\colorlet{ezVach2}{gold4}% vạch xiên sát ô chữ (nhạt hơn)
\colorlet{ezO}{gold5}% ô chữ vàng kem
%--- Kích thước
\newlength{\EZdai}\setlength{\EZdai}{1.1cm}% chiều cao dải nền
\newlength{\EZo}\setlength{\EZo}{0.68cm}% chiều cao ô chữ
\newlength{\EZxien}\setlength{\EZxien}{0.55cm}% độ lệch ngang của cạnh xiên
\newlength{\EZvach}\setlength{\EZvach}{0.75cm}% bề rộng mỗi vạch xiên
\newlength{\EZkhe}\setlength{\EZkhe}{0.12cm}% khe trắng giữa vạch và ô chữ
\newlength{\EZdem}\setlength{\EZdem}{0.4cm}% khoảng đệm từ chữ tới cạnh xiên
\newlength{\EZle}\setlength{\EZle}{1.8cm}% lề trái/phải của chữ (bằng lề thân bài)
\newlength{\EZrongmin}\setlength{\EZrongmin}{4.5cm}% ô chữ rộng tối thiểu
\newlength{\EZrong}% bề rộng chữ trong ô (tính tự động)
%--- Font
\newcommand{\EZfontO}{\fontsize{10.5pt}{12pt}\selectfont\bfseries}% chữ trong ô
\newcommand{\EZfontKH}{\fontfamily{qag}\fontsize{11pt}{12pt}\selectfont\bfseries}% \khauhieu cạnh logo
\newcommand{\EZfontTrang}{\fontsize{11.5pt}{12pt}\selectfont\bfseries}% chữ "Trang" và số trang
%--- Đo bề rộng chữ trong ô
\newcommand{\EZdorong}[1]{%
	\settowidth{\EZrong}{\EZfontO #1}%
	\ifdim\EZrong<\EZrongmin\setlength{\EZrong}{\EZrongmin}\fi}
%--- Logo + chữ trắng \khauhieu (dựng sẵn một lần khi vào \begin{document})
\newsavebox{\EZlogobox}
\AtBeginDocument{\sbox{\EZlogobox}{%
		\raisebox{-0.5\height}{\eduzylogo[1cm]}\hspace{0.15cm}%
		\raisebox{-0.5\height}{\color{white}\EZfontKH\khauhieu}}}
%--- Hai vạch xiên + khe trắng. #1: điểm gốc ở mép dưới, #2: hướng (1 sang phải, -1 sang trái),
%    #3: hướng nghiêng (1 kiểu "/", -1 kiểu "\")
\newcommand{\EZvachxien}[3]{%
	\fill[white] (#1) --++({(#3)*\EZxien},\EZo) --++({(#2)*\EZkhe},0) --++({-(#3)*\EZxien},-\EZo) --cycle;
	\fill[ezVach2] ([xshift=(#2)*\EZkhe]#1) --++({(#3)*\EZxien},\EZo) --++({(#2)*\EZvach},0) --++({-(#3)*\EZxien},-\EZo) --cycle;
	\fill[ezVach1] ([xshift=(#2)*(\EZkhe+\EZvach)]#1) --++({(#3)*\EZxien},\EZo) --++({(#2)*\EZvach},0) --++({-(#3)*\EZxien},-\EZo) --cycle;
	\draw[white,line width=0.6pt] ([xshift=(#2)*(\EZkhe+\EZvach)]#1) --++({(#3)*\EZxien},\EZo);
}
%--- Chữ trong ô (căn giữa theo chiều cao ô, không lệch vì dấu tiếng Việt)
\tikzset{ezchu/.style={inner sep=0pt,text=ezNen,text height=0.72em,text depth=0.22em}}

%====================== LOGO CHÌM GIỮA TRANG ======================
%  Vẽ cùng đầu trang nên nằm dưới chữ của thân bài, có trên mọi trang.
%  Tắt: \logochimtat (từ trang đó trở đi), bật lại: \logochimbat.
%  Đổi cỡ, độ đậm: khai báo trước \input, ví dụ \def\EZchimcao{9cm} \def\EZchimdo{0.1}
\providecommand{\EZchimcao}{10cm}% chiều cao logo chìm
\providecommand{\EZchimdo}{0.12}% độ đậm (0 = trong suốt, 1 = rõ hẳn)
\newif\ifEZchim \EZchimtrue
\newcommand{\logochimtat}{\EZchimfalse}
\newcommand{\logochimbat}{\EZchimtrue}
%  Dựng sẵn logo mờ một lần khi vào \begin{document}; transparency group làm mờ cả logo
%  một lượt nên không lộ vệt ở chỗ các dải màu chồng nhau.
\newsavebox{\EZchimtho}
\newsavebox{\EZchimbox}
\AtBeginDocument{%
	\sbox{\EZchimtho}{\eduzylogo[\EZchimcao]}%
	\sbox{\EZchimbox}{\begin{tikzpicture}
			\begin{scope}[transparency group,opacity=\EZchimdo]
				\node[inner sep=0pt] {\usebox{\EZchimtho}};
			\end{scope}
	\end{tikzpicture}}}
\newcommand{\EZvechim}{%
	\ifEZchim
		\node[inner sep=0pt] at (current page.center) {\usebox{\EZchimbox}};
	\fi}

%====================== ĐẦU TRANG - TRANG LẺ ======================
%  (logo bên trái, ô chữ \tenchuong bên phải)
\newcommand{\myfancyheadOdd}{%
	\EZdorong{\tenchuong}%
	\begin{tikzpicture}[remember picture,overlay]
		\coordinate (T) at (current page.north west);
		\coordinate (P) at (current page.north east);
		%--- logo chìm giữa trang
		\EZvechim
		%--- dải nền
		\fill[ezNen] (T) rectangle ([yshift=-\EZdai]P);
		%--- ô chữ: X là góc dưới bên trái, cạnh trái xiên kiểu "/"
		\coordinate (X) at ([xshift=-\EZle-\EZrong-\EZdem-\EZxien,yshift=-\EZdai]P);
		\fill[ezO] (X) --++(\EZxien,\EZo) -- ([yshift=-\EZdai+\EZo]P) -- ([yshift=-\EZdai]P) --cycle;
		\EZvachxien{X}{-1}{1}
		\node[ezchu,anchor=east] at ([xshift=-\EZle,yshift=-\EZdai+0.5*\EZo]P) {\EZfontO\tenchuong};
		%--- logo + tên
		\node[inner sep=0pt,anchor=west] at ([xshift=\EZle,yshift=-0.5*\EZdai]T) {\usebox{\EZlogobox}};
	\end{tikzpicture}%
}

%====================== ĐẦU TRANG - TRANG CHẴN =====================
%  (soi gương: ô chữ \tenchuong bên trái, logo bên phải)
\newcommand{\myfancyheadEven}{%
	\EZdorong{\tenchuong}%
	\begin{tikzpicture}[remember picture,overlay]
		\coordinate (T) at (current page.north west);
		\coordinate (P) at (current page.north east);
		%--- logo chìm giữa trang
		\EZvechim
		%--- dải nền
		\fill[ezNen] (T) rectangle ([yshift=-\EZdai]P);
		%--- ô chữ: X là góc dưới bên phải, cạnh phải xiên kiểu "\"
		\coordinate (X) at ([xshift=\EZle+\EZrong+\EZdem+\EZxien,yshift=-\EZdai]T);
		\fill[ezO] (X) --++(-\EZxien,\EZo) -- ([yshift=-\EZdai+\EZo]T) -- ([yshift=-\EZdai]T) --cycle;
		\EZvachxien{X}{1}{-1}
		\node[ezchu,anchor=west] at ([xshift=\EZle,yshift=-\EZdai+0.5*\EZo]T) {\EZfontO\tenchuong};
		%--- logo + tên
		\node[inner sep=0pt,anchor=east] at ([xshift=-\EZle,yshift=-0.5*\EZdai]P) {\usebox{\EZlogobox}};
	\end{tikzpicture}%
}

%====================== CHÂN TRANG - TRANG LẺ ======================
%  (ô chữ \thuonghieu bên trái, "Trang (số)" bên phải)
\newcommand{\myfancyfootOdd}{%
	\EZdorong{\thuonghieu}%
	\begin{tikzpicture}[remember picture,overlay]
		\coordinate (T) at (current page.south west);
		\coordinate (P) at (current page.south east);
		%--- dải nền
		\fill[ezNen] (T) rectangle ([yshift=\EZdai]P);
		%--- ô chữ sát mép trên dải: X là góc dưới bên phải, cạnh phải xiên kiểu "/"
		\coordinate (X) at ([xshift=\EZle+\EZrong+\EZdem,yshift=\EZdai-\EZo]T);
		\fill[ezO] ([yshift=\EZdai-\EZo]T) -- (X) --++(\EZxien,\EZo) -- ([yshift=\EZdai]T) --cycle;
		\EZvachxien{X}{1}{1}
		\node[ezchu,anchor=west] at ([xshift=\EZle,yshift=\EZdai-0.5*\EZo]T) {\EZfontO\thuonghieu};
		%--- số trang
		\node[circle,draw=white,line width=0.9pt,text=white,inner sep=0pt,minimum size=0.62cm,anchor=east]
		(so) at ([xshift=-\EZle,yshift=0.5*\EZdai]P) {\EZfontTrang\thepage};
		\node[inner sep=0pt,text=white,anchor=base east] (tr) at ([xshift=-0.15cm]so.base west) {\EZfontTrang Trang};
	\end{tikzpicture}%
}

%====================== CHÂN TRANG - TRANG CHẴN ====================
%  (soi gương: "Trang (số)" bên trái, ô chữ \tacgia bên phải)
\newcommand{\myfancyfootEven}{%
	\EZdorong{\tacgia}%
	\begin{tikzpicture}[remember picture,overlay]
		\coordinate (T) at (current page.south west);
		\coordinate (P) at (current page.south east);
		%--- dải nền
		\fill[ezNen] (T) rectangle ([yshift=\EZdai]P);
		%--- ô chữ sát mép trên dải: X là góc dưới bên trái, cạnh trái xiên kiểu "\"
		\coordinate (X) at ([xshift=-\EZle-\EZrong-\EZdem,yshift=\EZdai-\EZo]P);
		\fill[ezO] ([yshift=\EZdai-\EZo]P) -- (X) --++(-\EZxien,\EZo) -- ([yshift=\EZdai]P) --cycle;
		\EZvachxien{X}{-1}{-1}
		\node[ezchu,anchor=east] at ([xshift=-\EZle,yshift=\EZdai-0.5*\EZo]P) {\EZfontO\tacgia};
		%--- số trang
		%    (vòng tròn ẩn so0 nằm giữa dải để lấy đường chân chữ, giống trang lẻ)
		\node[circle,inner sep=0pt,minimum size=0.62cm] (so0) at ([yshift=0.5*\EZdai]T) {\EZfontTrang\phantom{\thepage}};
		\node[inner sep=0pt,text=white,anchor=base west] (tr) at ($(so0.base -| T)+(\EZle,0)$) {\EZfontTrang Trang};
		\node[circle,draw=white,line width=0.9pt,text=white,inner sep=0pt,minimum size=0.62cm,anchor=base west]
		(so) at ([xshift=0.15cm]tr.base east) {\EZfontTrang\thepage};
	\end{tikzpicture}%
}

%--- Giữ tương thích với các file cũ đang gọi \myfancyhead, \myfancyfoot
\newcommand{\myfancyhead}{\myfancyheadOdd}
\newcommand{\myfancyfoot}{\myfancyfootOdd}

%======================= GẮN VÀO TRANG =============================
\pagestyle{fancy}
\fancyhf{}
\fancyhead[LO]{\myfancyheadOdd}
\fancyhead[RE]{\myfancyheadEven}
\fancyfoot[LO]{\myfancyfootOdd}
\fancyfoot[RE]{\myfancyfootEven}
\renewcommand{\headrulewidth}{0pt}
\renewcommand{\footrulewidth}{0pt}
\fancypagestyle{plain}{%
	\fancyhf{}
	\fancyhead[LO]{\myfancyheadOdd}
	\fancyhead[RE]{\myfancyheadEven}
	\fancyfoot[LO]{\myfancyfootOdd}
	\fancyfoot[RE]{\myfancyfootEven}
	\renewcommand{\headrulewidth}{0pt}
	\renewcommand{\footrulewidth}{0pt}
}
%=================== HẾT PHẦN GỘP TỪ eduzy-de.tex ===================
\newcommand{\indam}[1]{
	{\fontsize{12pt}{9pt}\sffamily \selectfont\bfseries\color{gold1} #1}
}
\newcommand{\ind}[1]{
	{\fontsize{14pt}{9pt}\fontfamily{qag} \selectfont\bfseries\color{gold1} #1}
}
\renewcommand{\FalseEX}{\stepcounter{dapan}\textbf{\selectfont\color{gold1}\Alph{dapan}.}\ignorespaces}
\renewtheorem{ex}{\color{gold1}Câu}
\renewcommand{\baselinestretch}{1.1}
\newcommand{\hoac}[1]{
	\left[\begin{aligned}#1\end{aligned}\right.}
\newcommand{\heva}[1]{
	\left\{\begin{aligned}#1\end{aligned}\right.}
\usepackage{esvect}
\def\vec{\vv}
\def\vt{\vv}
\def\overrightarrow{\vv}
%====================== LOGO EDUZY VẼ BẰNG TIKZ ======================
%  Đường viền dò từ logo gốc; dải màu vàng kim đo trực tiếp trên ảnh gốc
%  (nghiêng 21 độ, tô bằng 80 dải hẹp nên hiển thị giống nhau trên mọi trình xem PDF).
%  Dùng: \eduzylogo          -> cao 2.4cm (mặc định)
%        \eduzylogo[1.5cm]   -> cao 1.5cm
\definecolor{ezText}{RGB}{245,200,46}
\newcommand{\eduzylogo}[1][2.4cm]{%
	\resizebox{!}{#1}{%
		\begin{tikzpicture}[x=0.01cm,y=0.01cm]
			\path[use as bounding box] (0,0) rectangle (224,-224);
			% mũ tốt nghiệp + vương miện
			\begin{scope}
				\clip
				(119.4,-176.1)..controls(118.9,-176.1)and(116.8,-176)..(114.6,-176)
				..controls(101.4,-176)and(85.8,-175.9)..(83.7,-175.7)
				..controls(81.8,-175.6)and(81,-175.6)..(80.2,-175.7)
				..controls(78.8,-175.9)and(77.7,-175.9)..(76.2,-175.7)
				..controls(74.9,-175.5)and(70.8,-175.5)..(64.3,-175.7)
				..controls(59.1,-175.9)and(58.3,-175.7)..(58.6,-173.7)
				..controls(58.6,-172.8)and(59.6,-169.6)..(60.3,-168)
				..controls(60.8,-166.9)and(61.4,-165.8)..(62.3,-165)
				..controls(63.1,-164.2)and(64.3,-163.4)..(65,-163.2)
				..controls(65.5,-163.1)and(72.3,-163)..(76.6,-163.1)
				..controls(77.9,-163.1)and(80.1,-163.1)..(81.4,-163)
				..controls(82.8,-162.9)and(86.3,-162.9)..(89.2,-162.9)
				..controls(93,-162.9)and(94.9,-162.8)..(96,-162.7)
				..controls(97.2,-162.5)and(101,-162.5)..(114,-162.4)
				..controls(123,-162.3)and(134.9,-162.2)..(140.2,-162.1)
				..controls(150.1,-162)and(152.2,-162.1)..(153.6,-162.4)
				..controls(154.4,-162.6)and(157.1,-163.8)..(158.1,-164.4)
				..controls(158.8,-164.8)and(159.7,-165.7)..(159.9,-166.3)
				..controls(160,-166.5)and(160.2,-167.1)..(160.2,-167.6)
				..controls(160.3,-168)and(160.6,-169.4)..(160.9,-170.6)
				..controls(161.2,-171.7)and(161.5,-173.1)..(161.6,-173.6)
				..controls(161.8,-174.4)and(161.8,-174.5)..(161.5,-174.9)
				..controls(161.2,-175.6)and(160.8,-175.9)..(159.9,-176)
				..controls(159.2,-176.1)and(154.1,-176.2)..(151.5,-176.2)
				..controls(148.1,-176.1)and(122.6,-176.1)..(121.6,-176.2)
				..controls(120.8,-176.2)and(119.8,-176.2)..(119.4,-176.1)--cycle
				(130.4,-160)..controls(129,-160)and(119.8,-160)..(110.1,-159.9)
				..controls(100.3,-159.8)and(88.5,-159.7)..(83.9,-159.6)
				..controls(77.8,-159.5)and(75.3,-159.5)..(74.5,-159.6)
				..controls(73.9,-159.7)and(72.5,-159.8)..(71.3,-159.8)
				..controls(69.6,-159.8)and(69,-159.8)..(68.5,-159.6)
				..controls(67.6,-159.3)and(66.8,-158.6)..(66.2,-157.6)
				..controls(66,-157.3)and(65.6,-156.5)..(65.3,-156.1)
				..controls(64.6,-155.1)and(64.5,-154.9)..(63.7,-152.9)
				..controls(63.2,-151.6)and(62.9,-150.9)..(61.8,-149.6)
				..controls(61.5,-149.2)and(61.1,-148.5)..(60.9,-147.9)
				..controls(60.7,-147.4)and(60.3,-146.6)..(60.1,-146.1)
				..controls(59.9,-145.5)and(59.5,-144.9)..(59,-144.3)
				..controls(58.5,-143.7)and(58.2,-143.2)..(57.9,-142.5)
				..controls(57.7,-141.9)and(57.3,-141.1)..(57,-140.7)
				..controls(56.6,-140.3)and(56.1,-139.3)..(55.8,-138.5)
				..controls(55.5,-137.8)and(55.1,-136.8)..(54.9,-136.5)
				..controls(54.7,-136.1)and(54.3,-135.2)..(53.9,-134.5)
				..controls(53.6,-133.8)and(53,-132.7)..(52.6,-132)
				..controls(51.6,-130.5)and(51.4,-130.2)..(51.1,-129)
				..controls(50.6,-127.5)and(50.6,-127.4)..(49.7,-126.3)
				..controls(49.2,-125.7)and(48.6,-124.8)..(48.3,-124.2)
				..controls(48,-123.6)and(47.5,-122.7)..(47.3,-122.2)
				..controls(47,-121.8)and(46.5,-120.7)..(46.1,-119.8)
				..controls(45.8,-118.9)and(45.2,-117.8)..(44.9,-117.2)
				..controls(44.2,-116.1)and(43.3,-114.4)..(42.3,-112.2)
				..controls(41.8,-111.3)and(41.3,-110.4)..(41.1,-110.1)
				..controls(40.5,-109.4)and(40.3,-109)..(40,-107.8)
				..controls(39.8,-107)and(39.6,-106.5)..(39.2,-106)
				..controls(38.9,-105.6)and(38.4,-104.7)..(38,-104)
				..controls(37.7,-103.4)and(37.2,-102.5)..(37,-102.1)
				..controls(36.7,-101.7)and(36.3,-100.8)..(36,-100.1)
				..controls(35.6,-98.9)and(35.3,-98.4)..(34.6,-97.6)
				..controls(34.2,-97.1)and(34.2,-96.3)..(34.6,-96)
				..controls(35,-95.7)and(35.5,-95.9)..(36.2,-96.7)
				..controls(36.5,-97)and(37.2,-97.5)..(37.6,-97.7)..controls(38,-98)and(38.6,-98.4)..(39.1,-98.9)
				..controls(39.6,-99.4)and(40.3,-100)..(40.6,-100.3)
				..controls(41.3,-100.8)and(42.4,-101.8)..(43.7,-102.9)
				..controls(44,-103.2)and(44.6,-103.7)..(45,-103.9)
				..controls(45.4,-104.2)and(45.9,-104.6)..(46.1,-104.8)
				..controls(47.1,-105.9)and(48.8,-107.1)..(50.5,-107.9)
				..controls(51.8,-108.6)and(52.3,-109)..(53.4,-110)
				..controls(54.8,-111.2)and(55.6,-111.6)..(57,-112.1)
				..controls(57.5,-112.3)and(58.5,-112.8)..(59.4,-113.3)
				..controls(60.2,-113.7)and(61.3,-114.3)..(61.9,-114.5)
				..controls(62.4,-114.7)and(63.3,-115.1)..(63.8,-115.3)
				..controls(64.4,-115.6)and(65.4,-115.9)..(66.2,-116.1)
				..controls(67,-116.3)and(68.4,-116.6)..(69.2,-116.9)
				..controls(72.3,-117.9)and(72.8,-118)..(78,-118.1)
				..controls(84,-118.3)and(86.1,-118.1)..(88.1,-117.1)
				..controls(89.4,-116.5)and(93,-115)..(93.7,-114.9)
				..controls(94.8,-114.6)and(95.4,-114.2)..(95.9,-113.4)
				..controls(96.2,-113)and(96.7,-112.5)..(97.1,-112.2)
				..controls(98,-111.6)and(101.3,-108.4)..(102.1,-107.2)
				..controls(102.4,-106.8)and(102.9,-106.2)..(103.2,-105.8)
				..controls(103.5,-105.5)and(104,-104.8)..(104.2,-104.2)
				..controls(104.5,-103.6)and(105,-102.8)..(105.3,-102.4)
				..controls(105.6,-102)and(106.1,-101.2)..(106.5,-100.6)
				..controls(107.1,-99.5)and(108.6,-97.5)..(109.5,-96.5)
				..controls(109.9,-96.1)and(110.1,-96)..(110.4,-96)
				..controls(111.1,-96)and(111.6,-96.5)..(111.8,-97.6)
				..controls(112,-98.6)and(112.3,-99.3)..(112.8,-99.9)
				..controls(113.7,-100.9)and(115.2,-103.1)..(115.7,-104.2)
				..controls(116.1,-105)and(116.5,-105.5)..(116.9,-106)
				..controls(117.3,-106.4)and(117.8,-106.9)..(118,-107.3)
				..controls(118.2,-107.6)and(118.6,-108)..(119,-108.3)
				..controls(119.6,-108.8)and(121.1,-110.3)..(121.8,-111.2)
				..controls(122.1,-111.7)and(122.5,-112)..(123.4,-112.4)
				..controls(124.8,-113.2)and(125.8,-113.8)..(127,-114.6)
				..controls(127.5,-114.9)and(128.1,-115.3)..(128.2,-115.4)
				..controls(128.4,-115.5)and(129.3,-115.8)..(130.2,-116)
				..controls(131.1,-116.2)and(132.5,-116.6)..(133.3,-116.9)
				..controls(135.7,-117.9)and(136.1,-117.9)..(142.5,-117.8)
				..controls(145.8,-117.8)and(148.4,-117.7)..(148.8,-117.6)
				..controls(149.2,-117.6)and(150.2,-117.2)..(151,-116.8)
				..controls(152,-116.4)and(152.8,-116)..(153.6,-115.9)
				..controls(154.5,-115.7)and(155.2,-115.4)..(156.3,-114.9)
				..controls(157.1,-114.5)and(158.2,-114)..(158.7,-113.8)
				..controls(159.1,-113.7)and(159.8,-113.3)..(160.2,-113)
				..controls(160.6,-112.7)and(161.3,-112.4)..(161.7,-112.2)
				..controls(162.7,-111.9)and(163.2,-111.7)..(164.1,-111.1)
				..controls(164.5,-110.8)and(165.3,-110.4)..(165.9,-110)
				..controls(166.5,-109.7)and(167.3,-109.2)..(167.6,-108.9)
				..controls(167.9,-108.6)and(168.6,-108.2)..(169,-108)
				..controls(169.4,-107.8)and(170.2,-107.3)..(170.7,-106.9)
				..controls(171.2,-106.5)and(172.1,-105.8)..(172.7,-105.4)
				..controls(173.2,-105)and(174.1,-104.4)..(174.5,-104)
				..controls(174.9,-103.7)and(175.7,-103.1)..(176.2,-102.6)
				..controls(176.7,-102.2)and(177.4,-101.5)..(177.9,-101.2)
				..controls(178.3,-100.8)and(178.8,-100.4)..(179,-100.2)
				..controls(179.2,-100)and(179.6,-99.6)..(180,-99.3)
				..controls(180.4,-99.1)and(180.9,-98.6)..(181.2,-98.4)
				..controls(181.5,-98.1)and(182.1,-97.7)..(182.5,-97.5)
				..controls(183.2,-97.2)and(183.5,-96.9)..(184.3,-95.9)
				..controls(184.6,-95.6)and(184.6,-95.6)..(184.8,-95.7)
				..controls(185,-95.8)and(185.1,-96)..(185.1,-96.2)
				..controls(185.1,-96.5)and(185,-96.7)..(184.6,-97.2)
				..controls(184.1,-97.8)and(184,-97.9)..(183.9,-98.8)
				..controls(183.7,-100.3)and(183.3,-101.3)..(182.9,-101.8)
				..controls(182.6,-102.1)and(182.3,-102.6)..(182.2,-102.9)
				..controls(182,-103.3)and(181.6,-104)..(181.2,-104.6)
				..controls(180.9,-105.2)and(180.5,-106.1)..(180.4,-106.6)
				..controls(179.9,-108)and(179.7,-108.5)..(179,-109.4)
				..controls(178.6,-109.8)and(178.2,-110.5)..(177.9,-111)
				..controls(177.7,-111.5)and(177.3,-112.3)..(177,-112.7)
				..controls(176.6,-113.2)and(176.4,-113.7)..(176.3,-114.2)
				..controls(176.2,-114.7)and(175.9,-115.2)..(175.4,-116)
				..controls(175,-116.6)and(174.4,-117.7)..(174,-118.5)
				..controls(173.7,-119.2)and(173.3,-120.1)..(173,-120.5)
				..controls(172.8,-120.9)and(172.4,-121.8)..(172.1,-122.5)
				..controls(171.8,-123.2)and(171.3,-124.1)..(171,-124.6)
				..controls(170.8,-125.1)and(170.3,-126)..(169.9,-126.6)
				..controls(169.6,-127.3)and(169.2,-128.1)..(169,-128.3)
				..controls(168.8,-128.6)and(168.5,-129.3)..(168.3,-129.8)
				..controls(168.1,-130.3)and(167.8,-131.1)..(167.5,-131.6)
				..controls(166.9,-132.5)and(164.2,-137.9)..(164.2,-138.1)
				..controls(164.2,-138.1)and(163.9,-138.8)..(163.4,-139.7)
				..controls(162.9,-140.5)and(162.4,-141.5)..(162.3,-141.9)
				..controls(161.9,-143)and(161.6,-143.4)..(160.9,-144.1)
				..controls(160.4,-144.6)and(160.3,-144.8)..(160.1,-145.4)
				..controls(160,-145.8)and(159.6,-146.8)..(159.2,-147.8)
				..controls(158.7,-148.8)and(158.1,-150)..(157.9,-150.6)
				..controls(157.6,-151.1)and(157.2,-151.9)..(156.9,-152.4)
				..controls(156.5,-152.8)and(156.1,-153.7)..(155.8,-154.4)
				..controls(155.6,-155)and(155.2,-155.9)..(154.9,-156.3)
				..controls(154.7,-156.7)and(154.4,-157.4)..(154.2,-157.9)
				..controls(154.1,-158.4)and(153.8,-159)..(153.5,-159.3)
				..controls(152.8,-160.1)and(153,-160.1)..(142.4,-160.1)
				..controls(137.2,-160.1)and(131.8,-160.1)..(130.4,-160)--cycle
				(141.5,-103.7)..controls(140.8,-103)and(139.8,-102.1)..(139.2,-101.6)
				..controls(138.5,-101.2)and(137.9,-100.6)..(137.8,-100.4)
				..controls(137,-99.5)and(134,-96.5)..(133.2,-96)
				..controls(132.7,-95.7)and(131.8,-94.9)..(131.1,-94.3)
				..controls(129.4,-92.7)and(127.9,-91.6)..(126.6,-91.1)
				..controls(125.2,-90.5)and(124.1,-90)..(122.9,-89.1)
				..controls(122.1,-88.6)and(121.3,-88.2)..(119.9,-87.7)
				..controls(118.8,-87.3)and(117.5,-86.8)..(116.9,-86.5)
				..controls(114.3,-85.5)and(113.7,-85.4)..(109.7,-85.5)
				..controls(105.3,-85.7)and(103.9,-85.9)..(101.6,-87)
				..controls(100.7,-87.4)and(99.4,-87.8)..(97.9,-88)
				..controls(97,-88.2)and(96.6,-88.3)..(95.8,-88.9)
				..controls(95.4,-89.2)and(94.7,-89.6)..(94.2,-89.8)
				..controls(92.9,-90.3)and(92.3,-90.7)..(90.9,-92.1)
				..controls(90.2,-92.7)and(89.4,-93.3)..(89,-93.5)..controls(87,-94.5)and(85.2,-96)..(83.6,-97.9)
				..controls(83.3,-98.2)and(82.5,-98.9)..(81.8,-99.4)
				..controls(80.7,-100.1)and(80.5,-100.3)..(80,-101.1)
				..controls(78.7,-103.1)and(77.2,-103.9)..(76.4,-103)
				..controls(75.7,-102.2)and(75.7,-103.2)..(75.7,-83.6)
				..controls(75.6,-73.6)and(75.6,-64.8)..(75.5,-64)
				..controls(75.5,-63.2)and(75.5,-61.4)..(75.6,-59.8)--(75.7,-57.1)--(76,-56.8)
				..controls(76.8,-56.2)and(78,-56.3)..(79.8,-57.2)
				..controls(80.4,-57.5)and(81.3,-57.8)..(81.8,-57.9)
				..controls(82.3,-58.1)and(83,-58.4)..(83.6,-58.8)..controls(84.3,-59.2)and(85,-59.5)..(86,-59.8)
				..controls(86.8,-60.1)and(88,-60.6)..(88.7,-60.9)
				..controls(89.3,-61.2)and(90.6,-61.7)..(91.4,-62)
				..controls(92.2,-62.3)and(93.7,-62.9)..(94.6,-63.3)
				..controls(95.6,-63.7)and(96.8,-64.1)..(97.4,-64.3)
				..controls(98,-64.5)and(99,-64.9)..(99.7,-65.2)
				..controls(100.4,-65.6)and(101.5,-66)..(102.3,-66.2)
				..controls(104.2,-66.8)and(105,-67.1)..(106.2,-67.8)
				..controls(107.9,-68.7)and(110.2,-69.4)..(111.3,-69.2)
				..controls(112.5,-69)and(115,-68.2)..(115.6,-67.8)
				..controls(117.3,-66.7)and(117.6,-66.6)..(119.5,-66.2)
				..controls(121.9,-65.8)and(122.1,-65.7)..(123,-65.1)
				..controls(124.1,-64.4)and(124.4,-64.3)..(126,-63.9)
				..controls(127,-63.7)and(127.7,-63.4)..(128.4,-63.1)
				..controls(129,-62.8)and(129.7,-62.5)..(130.6,-62.4)
				..controls(132.1,-62.1)and(134,-61.5)..(135.2,-61)
				..controls(135.7,-60.8)and(136.7,-60.5)..(137.3,-60.3)
				..controls(138.4,-60)and(139.3,-59.7)..(141.1,-59)
				..controls(141.6,-58.7)and(142,-58.6)..(142.5,-58.6)
				..controls(143.2,-58.6)and(143.3,-58.7)..(143.6,-59)
				..controls(144,-59.5)and(144.1,-59.9)..(144,-60.8)
				..controls(144,-61.2)and(144,-71.1)..(144,-82.8)--(143.9,-104.1)--(143.7,-104.5)
				..controls(143.5,-104.7)and(143.3,-104.8)..(143,-104.9)
				..controls(142.6,-104.9)and(142.5,-104.8)..(141.5,-103.7)--cycle
				(27.8,-91.4)..controls(27.1,-90.9)and(26.3,-90)..(26.1,-89.5)
				..controls(26,-89.2)and(25.8,-88.6)..(25.7,-87.9)
				..controls(25.4,-86.2)and(25.6,-84.8)..(26.5,-83.5)
				..controls(27.4,-82.3)and(28.3,-82)..(29.4,-82.3)
				..controls(31.1,-82.8)and(32.1,-84.9)..(32,-87.4)
				..controls(31.9,-89.3)and(31.2,-90.6)..(29.9,-91.4)
				..controls(29.1,-91.9)and(28.6,-91.9)..(27.8,-91.4)--cycle
				(190.1,-91.3)..controls(189.7,-91.2)and(189.3,-90.9)..(189.1,-90.7)
				..controls(188.3,-89.8)and(187.6,-87.7)..(187.6,-86.2)
				..controls(187.6,-85.5)and(187.7,-85.2)..(188.2,-84.4)
				..controls(189.1,-82.6)and(190.1,-81.8)..(191.3,-82.1)
				..controls(192,-82.3)and(193.5,-83.7)..(193.8,-84.4)
				..controls(194.1,-85.1)and(194.2,-87.1)..(194,-88.3)
				..controls(193.8,-89.8)and(193.3,-90.5)..(192.3,-91.3)
				..controls(191.7,-91.7)and(190.9,-91.7)..(190.1,-91.3)--cycle
				(173.7,-83.4)..controls(172.2,-82.5)and(172,-82.5)..(170.6,-82.4)
				..controls(169.7,-82.4)and(168.8,-82.4)..(168.3,-82.5)
				..controls(167.2,-82.7)and(166.9,-82.7)..(166.6,-82.2)
				..controls(166.3,-81.8)and(166.3,-81)..(166.8,-80)
				..controls(167.2,-78.9)and(167.4,-77.4)..(167.7,-74.1)
				..controls(167.8,-72.8)and(167.9,-71.8)..(167.9,-71.7)
				..controls(168,-71.7)and(168,-70.2)..(168,-68.4)..controls(168,-65.4)and(168,-64.9)..(167.7,-63)
				..controls(167.3,-60.5)and(167.1,-59.9)..(166.4,-58.4)
				..controls(166.1,-57.7)and(165.8,-56.9)..(165.7,-56.6)
				..controls(165.7,-56.2)and(165.4,-55.6)..(165.2,-55.2)
				..controls(164.6,-53.8)and(161.9,-50.7)..(161,-50.3)
				..controls(159.8,-49.7)and(157.9,-50)..(155.2,-51.2)
				..controls(154.7,-51.4)and(153.4,-51.8)..(152.2,-52.1)
				..controls(150.8,-52.4)and(149.4,-52.9)..(148.6,-53.2)
				..controls(147.8,-53.6)and(146.7,-54)..(146,-54.2)
				..controls(145.4,-54.4)and(144.5,-54.7)..(144,-54.9)
				..controls(143.6,-55.2)and(142.7,-55.5)..(141.8,-55.7)
				..controls(141,-55.9)and(139.6,-56.3)..(138.7,-56.6)
				..controls(137.8,-57)and(136.3,-57.4)..(135.3,-57.6)
				..controls(133.8,-57.9)and(133.2,-58.1)..(132.1,-58.7)
				..controls(131.1,-59.2)and(130.3,-59.4)..(129.4,-59.6)
				..controls(128.7,-59.8)and(127.8,-60)..(127.4,-60.1)
				..controls(127,-60.3)and(126,-60.5)..(125.1,-60.7)
				..controls(123.2,-61.1)and(121.7,-61.5)..(120.2,-62.1)
				..controls(119.7,-62.3)and(118.1,-62.8)..(116.7,-63.2)
				..controls(114.6,-63.9)and(114,-64.1)..(113,-64.7)
				..controls(111.9,-65.4)and(111.8,-65.4)..(111,-65.5)
				..controls(109.9,-65.5)and(109.3,-65.4)..(107.7,-64.7)
				..controls(107,-64.4)and(106,-64)..(105.3,-63.9)
				..controls(103.5,-63.5)and(102.2,-63)..(101.2,-62.5)
				..controls(100.1,-61.8)and(99.3,-61.5)..(97.8,-61.1)
				..controls(97.2,-61)and(96.2,-60.7)..(95.7,-60.4)
				..controls(95.1,-60.1)and(94.2,-59.8)..(93.5,-59.6)
				..controls(92.9,-59.4)and(91.7,-59)..(91,-58.7)..controls(90.2,-58.4)and(88.9,-58)..(88,-57.7)
				..controls(87.2,-57.5)and(86.1,-57.1)..(85.5,-56.8)
				..controls(85,-56.6)and(84.1,-56.2)..(83.5,-56)..controls(82.9,-55.7)and(81.8,-55.2)..(81,-54.9)
				..controls(80.1,-54.5)and(79.1,-54.2)..(78.7,-54.1)
				..controls(77.9,-53.8)and(77.3,-53.6)..(76.2,-52.9)
				..controls(75.7,-52.6)and(75.2,-52.4)..(74.1,-52.2)
				..controls(72.5,-51.9)and(72.4,-51.8)..(71.3,-51.1)
				..controls(69.9,-50.1)and(69.4,-49.9)..(67.9,-49.8)
				..controls(67,-49.7)and(66.6,-49.6)..(66.3,-49.4)
				..controls(66.1,-49.3)and(65.7,-49)..(65.4,-48.8)
				..controls(65.1,-48.7)and(64.2,-48.4)..(63.6,-48.2)
				..controls(62.9,-48)and(61.7,-47.6)..(60.9,-47.2)
				..controls(60.1,-46.8)and(58.9,-46.4)..(58.2,-46.2)
				..controls(57.6,-46)and(56.8,-45.7)..(56.5,-45.5)
				..controls(54.9,-44.6)and(54.6,-44.5)..(53.2,-44.2)
				..controls(51.7,-43.8)and(51,-43.6)..(49.9,-42.9)
				..controls(49.5,-42.6)and(48.9,-42.4)..(47.5,-42)
				..controls(45.7,-41.6)and(45.6,-41.5)..(44.3,-40.6)
				..controls(43.5,-40.1)and(42.8,-39.6)..(42.7,-39.4)
				..controls(42.4,-38.7)and(43,-38.1)..(44.5,-37.6)..controls(44.9,-37.5)and(46,-37)..(46.9,-36.6)
				..controls(47.9,-36.1)and(48.8,-35.8)..(49.6,-35.6)
				..controls(50.5,-35.4)and(51.3,-35.2)..(52.4,-34.6)
				..controls(53.3,-34.1)and(54.5,-33.7)..(54.9,-33.5)
				..controls(55.4,-33.4)and(56.4,-33)..(57.2,-32.6)
				..controls(58.6,-32)and(60.1,-31.5)..(61.6,-31.2)
				..controls(62.3,-31.1)and(62.7,-31)..(63.2,-30.7)
				..controls(64.1,-30)and(64.3,-29.9)..(65.3,-29.7)
				..controls(65.7,-29.6)and(66.5,-29.2)..(67.1,-28.9)
				..controls(67.6,-28.6)and(68.8,-28.2)..(69.6,-27.9)
				..controls(70.5,-27.7)and(71.6,-27.3)..(72.1,-27)
				..controls(72.6,-26.7)and(73.8,-26.2)..(74.8,-25.9)
				..controls(75.7,-25.5)and(76.7,-25.1)..(77.1,-25)
				..controls(77.4,-24.8)and(78.5,-24.4)..(79.6,-24.1)
				..controls(80.6,-23.8)and(82,-23.2)..(82.8,-22.9)
				..controls(83.5,-22.5)and(84.8,-22)..(85.8,-21.8)
				..controls(86.7,-21.5)and(87.8,-21.1)..(88.1,-20.9)
				..controls(88.4,-20.7)and(89.3,-20.3)..(90.1,-20.1)
				..controls(90.8,-19.8)and(91.9,-19.4)..(92.4,-19.1)
				..controls(93.5,-18.5)and(94.4,-18.2)..(95.4,-18.1)
				..controls(95.7,-18)and(96.4,-17.8)..(96.9,-17.5)
				..controls(97.5,-17.3)and(98.4,-16.9)..(99,-16.6)
				..controls(99.7,-16.3)and(100.9,-15.9)..(101.9,-15.6)
				..controls(102.9,-15.3)and(104.4,-14.7)..(105.3,-14.3)
				..controls(108.1,-13.2)and(108.5,-13.1)..(110.1,-13.1)
				..controls(111.4,-13.1)and(111.5,-13.1)..(112.4,-13.5)
				..controls(113.2,-14)and(115,-14.8)..(116.6,-15.4)
				..controls(116.9,-15.5)and(117.9,-15.8)..(118.8,-16.1)
				..controls(120.9,-16.6)and(121.5,-16.9)..(122.6,-17.9)
				..controls(123,-18.3)and(123.7,-18.8)..(123.9,-18.9)
				..controls(124.6,-19.4)and(126,-19.8)..(127.9,-20.2)
				..controls(129.6,-20.5)and(129.8,-20.6)..(131.1,-21.4)
				..controls(131.5,-21.7)and(132.4,-22.1)..(133.1,-22.3)
				..controls(136.1,-23.2)and(138.3,-24.1)..(140.2,-25)
				..controls(141,-25.5)and(141.5,-25.6)..(142.5,-25.8)
				..controls(145.1,-26.3)and(146,-26.6)..(149.4,-28.1)
				..controls(151.3,-28.9)and(153.5,-29.8)..(154.3,-30.1)
				..controls(155.1,-30.4)and(156.2,-30.8)..(156.6,-31)
				..controls(157.1,-31.2)and(158.2,-31.6)..(159.1,-31.8)
				..controls(160.6,-32.2)and(160.9,-32.3)..(161.7,-32.9)
				..controls(162.4,-33.3)and(162.9,-33.6)..(163.7,-33.8)
				..controls(164.3,-34)and(164.8,-34.1)..(164.9,-34.1)
				..controls(165,-34.1)and(165.6,-34.4)..(166.4,-34.8)
				..controls(167.3,-35.2)and(168.1,-35.4)..(169,-35.6)
				..controls(170,-35.8)and(170.5,-36)..(171.7,-36.6)
				..controls(172.4,-37)and(173.3,-37.4)..(173.6,-37.4)
				..controls(175.6,-37.7)and(176.2,-38)..(176.7,-39.1)
				..controls(176.9,-39.6)and(176.9,-39.7)..(176.7,-40.4)
				..controls(176.4,-41.6)and(175.8,-42.7)..(175.2,-43.4)
				..controls(174.5,-44.1)and(172.6,-45.5)..(171.7,-45.8)
				..controls(170.6,-46.2)and(169.1,-46)..(168.2,-45.4)
				..controls(167.9,-45.2)and(167.3,-44.8)..(166.9,-44.4)
				..controls(166.4,-44.1)and(165.8,-43.6)..(165.4,-43.2)
				..controls(165.1,-42.9)and(164.4,-42.3)..(163.9,-41.9)
				..controls(163.3,-41.5)and(162.7,-40.9)..(162.4,-40.5)
				..controls(161.4,-39.2)and(159.9,-38)..(158.8,-37.8)
				..controls(158.4,-37.8)and(158.3,-37.7)..(158.2,-37.2)
				..controls(158,-36.8)and(157.4,-36.1)..(157.1,-36.1)
				..controls(157,-36.1)and(156.7,-36.1)..(156.5,-36)
				..controls(156.1,-35.9)and(155.9,-35.7)..(155.7,-35.2)
				..controls(155.3,-34.6)and(155.1,-34.4)..(154,-34.1)
				..controls(153.6,-34)and(153.2,-33.7)..(152.9,-33.5)
				..controls(152.4,-32.9)and(151.6,-32.4)..(151.3,-32.4)
				..controls(150.9,-32.4)and(150.4,-32.6)..(150.4,-32.9)
				..controls(150.4,-33.2)and(150.8,-33.4)..(151.4,-33.6)
				..controls(152.1,-33.7)and(152.8,-34.1)..(153.5,-34.9)
				..controls(154.1,-35.5)and(154.3,-35.6)..(155,-35.8)
				..controls(155.7,-36)and(155.7,-36)..(156,-36.6)
				..controls(156.3,-37.3)and(156.7,-37.7)..(157.5,-37.8)
				..controls(158.1,-37.9)and(158.2,-38)..(158.4,-38.8)
				..controls(158.5,-39.4)and(158.6,-39.6)..(159.6,-40.4)
				..controls(160.2,-41)and(160.9,-41.6)..(161.3,-41.8)
				..controls(161.7,-42)and(161.9,-42.3)..(162.2,-42.8)
				..controls(162.4,-43.2)and(162.8,-43.8)..(163.1,-44.1)
				..controls(163.9,-44.8)and(165.1,-46.7)..(165.7,-48)
				..controls(166.2,-49.3)and(166.4,-49.7)..(167.2,-50.5)
				..controls(167.5,-50.8)and(167.9,-51.4)..(168.1,-51.8)
				..controls(168.3,-52.2)and(168.8,-53)..(169.1,-53.6)
				..controls(169.6,-54.3)and(169.8,-54.9)..(169.9,-55.6)
				..controls(170.4,-57.8)and(170.5,-58.1)..(171.2,-59.5)
				..controls(171.6,-60.3)and(172,-61.5)..(172.2,-62.1)
				..controls(172.3,-62.8)and(172.7,-63.7)..(172.9,-64.1)
				..controls(173.2,-64.6)and(173.4,-65.3)..(173.4,-65.6)
				..controls(173.8,-68.8)and(174.2,-70.6)..(174.9,-72.6)
				..controls(175.5,-74.5)and(175.9,-77.1)..(176.1,-80.2)
				..controls(176.3,-83.4)and(175.5,-84.4)..(173.7,-83.4)--cycle;
				\definecolor{ezs}{RGB}{247,179,26}\fill[ezs](-5.7,-156.6)--(-3.4,-157.5)--(61.9,12.7)--(59.7,13.6)--cycle;
				\definecolor{ezs}{RGB}{247,179,22}\fill[ezs](-3.8,-157.3)--(-1.5,-158.2)--(63.9,12)--(61.6,12.9)--cycle;
				\definecolor{ezs}{RGB}{248,179,19}\fill[ezs](-1.8,-158.1)--(0.4,-158.9)--(65.8,11.2)--(63.5,12.1)--cycle;
				\definecolor{ezs}{RGB}{248,178,16}\fill[ezs](0.1,-158.8)--(2.4,-159.7)--(67.7,10.5)--(65.4,11.4)--cycle;
				\definecolor{ezs}{RGB}{249,178,15}\fill[ezs](2,-159.5)--(4.3,-160.4)--(69.6,9.8)--(67.3,10.7)--cycle;
				\definecolor{ezs}{RGB}{251,177,17}\fill[ezs](3.9,-160.3)--(6.2,-161.1)--(71.5,9)--(69.2,9.9)--cycle;
				\definecolor{ezs}{RGB}{252,179,18}\fill[ezs](5.8,-161)--(8.1,-161.9)--(73.4,8.3)--(71.1,9.2)--cycle;
				\definecolor{ezs}{RGB}{252,182,17}\fill[ezs](7.7,-161.7)--(10,-162.6)--(75.3,7.6)--(73,8.5)--cycle;
				\definecolor{ezs}{RGB}{252,185,17}\fill[ezs](9.6,-162.5)--(11.9,-163.4)--(77.2,6.8)--(75,7.7)--cycle;
				\definecolor{ezs}{RGB}{251,189,18}\fill[ezs](11.5,-163.2)--(13.8,-164.1)--(79.2,6.1)--(76.9,7)--cycle;
				\definecolor{ezs}{RGB}{251,191,22}\fill[ezs](13.4,-163.9)--(15.7,-164.8)--(81.1,5.4)--(78.8,6.2)--cycle;
				\definecolor{ezs}{RGB}{252,193,25}\fill[ezs](15.4,-164.7)--(17.6,-165.6)--(83,4.6)--(80.7,5.5)--cycle;
				\definecolor{ezs}{RGB}{253,195,27}\fill[ezs](17.3,-165.4)--(19.6,-166.3)--(84.9,3.9)--(82.6,4.8)--cycle;
				\definecolor{ezs}{RGB}{253,196,31}\fill[ezs](19.2,-166.1)--(21.5,-167)--(86.8,3.2)--(84.5,4)--cycle;
				\definecolor{ezs}{RGB}{254,199,35}\fill[ezs](21.1,-166.9)--(23.4,-167.8)--(88.7,2.4)--(86.4,3.3)--cycle;
				\definecolor{ezs}{RGB}{253,201,38}\fill[ezs](23,-167.6)--(25.3,-168.5)--(90.6,1.7)--(88.3,2.6)--cycle;
				\definecolor{ezs}{RGB}{253,203,40}\fill[ezs](24.9,-168.3)--(27.2,-169.2)--(92.5,1)--(90.3,1.8)--cycle;
				\definecolor{ezs}{RGB}{253,205,42}\fill[ezs](26.8,-169.1)--(29.1,-170)--(94.4,0.2)--(92.2,1.1)--cycle;
				\definecolor{ezs}{RGB}{253,207,45}\fill[ezs](28.7,-169.8)--(31,-170.7)--(96.4,-0.5)--(94.1,0.4)--cycle;
				\definecolor{ezs}{RGB}{253,210,49}\fill[ezs](30.7,-170.5)--(32.9,-171.4)--(98.3,-1.2)--(96,-0.4)--cycle;
				\definecolor{ezs}{RGB}{253,212,54}\fill[ezs](32.6,-171.3)--(34.9,-172.2)--(100.2,-2)--(97.9,-1.1)--cycle;
				\definecolor{ezs}{RGB}{253,214,59}\fill[ezs](34.5,-172)--(36.8,-172.9)--(102.1,-2.7)--(99.8,-1.8)--cycle;
				\definecolor{ezs}{RGB}{253,216,63}\fill[ezs](36.4,-172.7)--(38.7,-173.6)--(104,-3.4)--(101.7,-2.6)--cycle;
				\definecolor{ezs}{RGB}{253,218,65}\fill[ezs](38.3,-173.5)--(40.6,-174.4)--(105.9,-4.2)--(103.6,-3.3)--cycle;
				\definecolor{ezs}{RGB}{254,220,66}\fill[ezs](40.2,-174.2)--(42.5,-175.1)--(107.8,-4.9)--(105.5,-4)--cycle;
				\definecolor{ezs}{RGB}{254,222,68}\fill[ezs](42.1,-175)--(44.4,-175.8)--(109.7,-5.6)--(107.5,-4.8)--cycle;
				\definecolor{ezs}{RGB}{254,224,71}\fill[ezs](44,-175.7)--(46.3,-176.6)--(111.7,-6.4)--(109.4,-5.5)--cycle;
				\definecolor{ezs}{RGB}{254,226,75}\fill[ezs](46,-176.4)--(48.2,-177.3)--(113.6,-7.1)--(111.3,-6.2)--cycle;
				\definecolor{ezs}{RGB}{254,228,79}\fill[ezs](47.9,-177.2)--(50.1,-178)--(115.5,-7.8)--(113.2,-7)--cycle;
				\definecolor{ezs}{RGB}{254,230,82}\fill[ezs](49.8,-177.9)--(52.1,-178.8)--(117.4,-8.6)--(115.1,-7.7)--cycle;
				\definecolor{ezs}{RGB}{254,232,85}\fill[ezs](51.7,-178.6)--(54,-179.5)--(119.3,-9.3)--(117,-8.4)--cycle;
				\definecolor{ezs}{RGB}{254,234,88}\fill[ezs](53.6,-179.4)--(55.9,-180.2)--(121.2,-10)--(118.9,-9.2)--cycle;
				\definecolor{ezs}{RGB}{254,237,91}\fill[ezs](55.5,-180.1)--(57.8,-181)--(123.1,-10.8)--(120.8,-9.9)--cycle;
				\definecolor{ezs}{RGB}{254,239,94}\fill[ezs](57.4,-180.8)--(59.7,-181.7)--(125,-11.5)--(122.8,-10.6)--cycle;
				\definecolor{ezs}{RGB}{254,241,97}\fill[ezs](59.3,-181.6)--(61.6,-182.4)--(127,-12.2)--(124.7,-11.4)--cycle;
				\definecolor{ezs}{RGB}{254,243,101}\fill[ezs](61.2,-182.3)--(63.5,-183.2)--(128.9,-13)--(126.6,-12.1)--cycle;
				\definecolor{ezs}{RGB}{254,245,103}\fill[ezs](63.2,-183)--(65.4,-183.9)--(130.8,-13.7)--(128.5,-12.8)--cycle;
				\definecolor{ezs}{RGB}{254,245,105}\fill[ezs](65.1,-183.8)--(67.4,-184.6)--(132.7,-14.4)--(130.4,-13.6)--cycle;
				\definecolor{ezs}{RGB}{254,245,105}\fill[ezs](67,-184.5)--(69.3,-185.4)--(134.6,-15.2)--(132.3,-14.3)--cycle;
				\definecolor{ezs}{RGB}{254,244,103}\fill[ezs](68.9,-185.2)--(71.2,-186.1)--(136.5,-15.9)--(134.2,-15)--cycle;
				\definecolor{ezs}{RGB}{253,243,101}\fill[ezs](70.8,-186)--(73.1,-186.8)--(138.4,-16.6)--(136.1,-15.8)--cycle;
				\definecolor{ezs}{RGB}{253,241,98}\fill[ezs](72.7,-186.7)--(75,-187.6)--(140.3,-17.4)--(138,-16.5)--cycle;
				\definecolor{ezs}{RGB}{253,239,95}\fill[ezs](74.6,-187.4)--(76.9,-188.3)--(142.2,-18.1)--(140,-17.2)--cycle;
				\definecolor{ezs}{RGB}{253,236,92}\fill[ezs](76.5,-188.2)--(78.8,-189)--(144.2,-18.8)--(141.9,-18)--cycle;
				\definecolor{ezs}{RGB}{253,234,88}\fill[ezs](78.5,-188.9)--(80.7,-189.8)--(146.1,-19.6)--(143.8,-18.7)--cycle;
				\definecolor{ezs}{RGB}{253,232,85}\fill[ezs](80.4,-189.6)--(82.7,-190.5)--(148,-20.3)--(145.7,-19.4)--cycle;
				\definecolor{ezs}{RGB}{253,230,82}\fill[ezs](82.3,-190.4)--(84.6,-191.2)--(149.9,-21)--(147.6,-20.2)--cycle;
				\definecolor{ezs}{RGB}{253,228,78}\fill[ezs](84.2,-191.1)--(86.5,-192)--(151.8,-21.8)--(149.5,-20.9)--cycle;
				\definecolor{ezs}{RGB}{253,226,74}\fill[ezs](86.1,-191.8)--(88.4,-192.7)--(153.7,-22.5)--(151.4,-21.6)--cycle;
				\definecolor{ezs}{RGB}{253,224,70}\fill[ezs](88,-192.6)--(90.3,-193.4)--(155.6,-23.3)--(153.3,-22.4)--cycle;
				\definecolor{ezs}{RGB}{253,222,67}\fill[ezs](89.9,-193.3)--(92.2,-194.2)--(157.5,-24)--(155.3,-23.1)--cycle;
				\definecolor{ezs}{RGB}{253,220,65}\fill[ezs](91.8,-194)--(94.1,-194.9)--(159.5,-24.7)--(157.2,-23.8)--cycle;
				\definecolor{ezs}{RGB}{253,217,63}\fill[ezs](93.7,-194.8)--(96,-195.6)--(161.4,-25.5)--(159.1,-24.6)--cycle;
				\definecolor{ezs}{RGB}{253,215,60}\fill[ezs](95.7,-195.5)--(97.9,-196.4)--(163.3,-26.2)--(161,-25.3)--cycle;
				\definecolor{ezs}{RGB}{253,213,56}\fill[ezs](97.6,-196.2)--(99.9,-197.1)--(165.2,-26.9)--(162.9,-26)--cycle;
				\definecolor{ezs}{RGB}{252,211,53}\fill[ezs](99.5,-197)--(101.8,-197.8)--(167.1,-27.7)--(164.8,-26.8)--cycle;
				\definecolor{ezs}{RGB}{252,209,49}\fill[ezs](101.4,-197.7)--(103.7,-198.6)--(169,-28.4)--(166.7,-27.5)--cycle;
				\definecolor{ezs}{RGB}{252,207,45}\fill[ezs](103.3,-198.4)--(105.6,-199.3)--(170.9,-29.1)--(168.6,-28.2)--cycle;
				\definecolor{ezs}{RGB}{253,205,42}\fill[ezs](105.2,-199.2)--(107.5,-200)--(172.8,-29.9)--(170.6,-29)--cycle;
				\definecolor{ezs}{RGB}{253,202,40}\fill[ezs](107.1,-199.9)--(109.4,-200.8)--(174.7,-30.6)--(172.5,-29.7)--cycle;
				\definecolor{ezs}{RGB}{253,200,38}\fill[ezs](109,-200.6)--(111.3,-201.5)--(176.7,-31.3)--(174.4,-30.4)--cycle;
				\definecolor{ezs}{RGB}{253,198,35}\fill[ezs](111,-201.4)--(113.2,-202.2)--(178.6,-32.1)--(176.3,-31.2)--cycle;
				\definecolor{ezs}{RGB}{253,196,31}\fill[ezs](112.9,-202.1)--(115.2,-203)--(180.5,-32.8)--(178.2,-31.9)--cycle;
				\definecolor{ezs}{RGB}{253,194,28}\fill[ezs](114.8,-202.8)--(117.1,-203.7)--(182.4,-33.5)--(180.1,-32.6)--cycle;
				\definecolor{ezs}{RGB}{253,192,26}\fill[ezs](116.7,-203.6)--(119,-204.5)--(184.3,-34.3)--(182,-33.4)--cycle;
				\definecolor{ezs}{RGB}{253,189,24}\fill[ezs](118.6,-204.3)--(120.9,-205.2)--(186.2,-35)--(183.9,-34.1)--cycle;
				\definecolor{ezs}{RGB}{253,186,20}\fill[ezs](120.5,-205)--(122.8,-205.9)--(188.1,-35.7)--(185.8,-34.8)--cycle;
				\definecolor{ezs}{RGB}{253,184,17}\fill[ezs](122.4,-205.8)--(124.7,-206.7)--(190,-36.5)--(187.8,-35.6)--cycle;
				\definecolor{ezs}{RGB}{252,183,15}\fill[ezs](124.3,-206.5)--(126.6,-207.4)--(192,-37.2)--(189.7,-36.3)--cycle;
				\definecolor{ezs}{RGB}{252,181,14}\fill[ezs](126.3,-207.2)--(128.5,-208.1)--(193.9,-37.9)--(191.6,-37.1)--cycle;
				\definecolor{ezs}{RGB}{252,179,13}\fill[ezs](128.2,-208)--(130.4,-208.9)--(195.8,-38.7)--(193.5,-37.8)--cycle;
				\definecolor{ezs}{RGB}{252,178,12}\fill[ezs](130.1,-208.7)--(132.4,-209.6)--(197.7,-39.4)--(195.4,-38.5)--cycle;
				\definecolor{ezs}{RGB}{251,178,10}\fill[ezs](132,-209.4)--(134.3,-210.3)--(199.6,-40.1)--(197.3,-39.3)--cycle;
				\definecolor{ezs}{RGB}{251,178,10}\fill[ezs](133.9,-210.2)--(136.2,-211.1)--(201.5,-40.9)--(199.2,-40)--cycle;
				\definecolor{ezs}{RGB}{250,177,11}\fill[ezs](135.8,-210.9)--(138.1,-211.8)--(203.4,-41.6)--(201.1,-40.7)--cycle;
				\definecolor{ezs}{RGB}{250,178,11}\fill[ezs](137.7,-211.6)--(140,-212.5)--(205.3,-42.3)--(203.1,-41.5)--cycle;
				\definecolor{ezs}{RGB}{248,180,10}\fill[ezs](139.6,-212.4)--(141.9,-213.3)--(207.3,-43.1)--(205,-42.2)--cycle;
				\definecolor{ezs}{RGB}{247,181,9}\fill[ezs](141.5,-213.1)--(143.8,-214)--(209.2,-43.8)--(206.9,-42.9)--cycle;
				\definecolor{ezs}{RGB}{246,182,9}\fill[ezs](143.5,-213.8)--(145.7,-214.7)--(211.1,-44.5)--(208.8,-43.7)--cycle;
				\definecolor{ezs}{RGB}{246,182,10}\fill[ezs](145.4,-214.6)--(147.7,-215.5)--(213,-45.3)--(210.7,-44.4)--cycle;
			\end{scope}
			% chữ EDUZY (Times đậm, làm dày nét bằng 5 lớp lệch nhẹ; căn đúng khung chữ logo gốc)
			\foreach \dx/\dy in {-0.7/0,0.7/0,0/-0.3,0/0.3,0/0}
			\node[anchor=center,inner sep=0pt,outer sep=0pt,text=ezText,align=center] at (109.6+\dx,-193.6+\dy)
			{\resizebox{1.01cm}{0.22cm}{\fontencoding{T1}\fontfamily{ptm}\fontseries{b}\selectfont EDUZY}};
		\end{tikzpicture}%
	}%
}
%====================== THIẾT LẬP THÂN ĐỀ ======================
%--- Môi trường tự luận {bt}: ex_test 3.3.4 chưa khai báo sẵn nên tạo ở đây
%    (ex_test đã tự móc \fixshowans và \listenumerate cho bt)
\makeatletter
\@ifundefined{bt}{\newtheorem{bt}{\color{gold1}\namebt}}{}
\makeatother
\renewcommand{\sectionmark}[1]{\markright{\thesection\ #1}}
%--- Khung tiêu đề (logo + tối đa 4 dòng; dòng 3, 4 để trống {} thì bỏ qua)
%  Đề thi:    \khungtieude{THI THỬ THPT QUỐC GIA NĂM 2026}{Môn Toán -- Đề số 21}
%                         {(Đề thi gồm: ...)}{Thời gian làm bài: 90 phút (không kể thời gian phát đề).}
%  Chuyên đề: \khungtieude{CHUYÊN ĐỀ 3 -- DÃY SỐ VÀ CÔNG THỨC TRUY HỒI}
%                         {Bồi dưỡng học sinh giỏi Toán THPT}{...}{}
\newcommand{\khungtieude}[4]{%
	\begin{tcolorbox}[colframe=gold1,colback=gold5!45,boxrule=0.5mm]
		\begin{minipage}[c]{3cm}
			\centerline{\eduzylogo[2.4cm]}
		\end{minipage}\hspace{0.5cm}
		\begin{minipage}[c]{12.5cm}
			\vskip 0.15 cm
			\centerline{\ind{#1}}
			\vskip 0.15 cm
			\centerline{\textbf{#2}}
			\if\relax\detokenize{#3}\relax\else
				\vskip 0.15 cm
				\centerline{\textit{#3}}
			\fi
			\if\relax\detokenize{#4}\relax\else
				\vskip 0.15 cm
				\centerline{\textit{#4}}
			\fi
			\vskip 0.15 cm
		\end{minipage}
	\end{tcolorbox}
}
%====================== BỔ SUNG CHO BÀI HỌC, CHUYÊN ĐỀ ======================
\usepackage[shortlabels]{enumitem}
\usepackage{tabularx}
%--- Danh sách gọn: \begin{enumerate}[gon] ... \end{enumerate}
\SetEnumitemKey{gon}{leftmargin=*,itemsep=1pt,parsep=0pt,topsep=3pt}
%--- Khung lý thuyết: \begin{khung}{} ... \end{khung}; tiêu đề chỉ dùng khi thật ngắn ("Mẹo", "Lỗi thường gặp")
\newtcolorbox{khung}[1]{breakable,enhanced,colframe=gold1,colback=gold5!45,
	boxrule=0.4mm,fonttitle=\bfseries\sffamily,coltitle=white,colbacktitle=gold1,
	attach boxed title to top left={xshift=6mm,yshift=-3mm},
	boxed title style={arc=1mm,boxrule=0.3mm},title={#1}}
%--- Tiêu đề mục nhỏ trong một dạng: \muc{Dấu hiệu}, \muc{Phương pháp}, \muc{Bài tự luyện}
\makeatletter
\newcommand{\muc}[1]{\par\medskip\noindent\indam{#1}\par\nobreak\smallskip\@afterindentfalse\@afterheading}
\makeatother
%--- Ô công thức cần nhớ (căn giữa, viền vàng, vừa khít nội dung)
%     \begin{center}\congthuc{$Q(n)=(n-1)!$}\end{center}
\newtcbox{\congthuc}{on line,colframe=gold1,colback=gold5!45,boxrule=0.5pt,arc=2pt,
	left=8pt,right=8pt,top=4pt,bottom=4pt}
%--- Công tắc lời giải bài tự luyện (ví dụ mẫu luôn dùng \loigiai)
%     \loigiaiTL[{đáp số}]{lời giải đầy đủ}
%     \TLfull -> in lời giải đầy đủ;  \TLdapso -> chỉ in một dòng "Đáp số";  \TLan -> ẩn hoàn toàn
%     (\CheDo là dethi thì mọi lời giải và đáp số đều ẩn)
\makeatletter
\newif\ifTLhien   \TLhientrue
\newif\ifTLdapso  \TLdapsofalse
\newcommand{\TLfull}{\TLhientrue\TLdapsofalse}
\newcommand{\TLdapso}{\TLhientrue\TLdapsotrue}
\newcommand{\TLan}{\TLhienfalse}
\newcommand{\loigiaiTL}[2][]{%
	\ifTLhien
		\ifTLdapso
			\def\@tlarg{#1}%
			\ifx\@tlarg\@empty\else
				\ifdefstring{\kindOPTN}{dethi}{}{\par\noindent{\bfseries\color{gold1}Đáp số:} #1\par}%
			\fi
		\else
			\loigiai{#2}%
		\fi
	\fi}
\makeatother
\TLdapso% mặc định: chỉ in đáp số bài tự luyện
%====================== BẢNG ĐÚNG/SAI (\choiceTFt) ======================
%  Hàng tiêu đề nền nâu đồng chữ trắng (Phát biểu | Đ | S), các hàng tô xen kẽ kem / trắng,
%  đường kẻ vàng kim, nhãn a) b) c) d) màu nâu đồng. Bản có đáp án: cột được chọn có vòng tròn
%  nâu đồng với dấu tích trắng (câu đúng tích ở cột Đúng, câu sai tích ở cột Sai).
%  Tiêu đề ghi đủ "Đúng | Sai": \OPTN{viettat=0} sau \input (cột Đúng/Sai rộng ra, phát biểu hẹp lại).
%  Dùng bảng gốc của ex_test: \def\BangDSgoc{1} trước \input{preamble-eduzy-v2}.
\providecommand{\BangDSgoc}{0}
\ifnum\BangDSgoc=0
\usepackage{colortbl,pifont}
\colorlet{ezDSke}{gold2!60!gold1}% đường kẻ bảng (nét 0,4pt như ex_test để không đổi chiều cao bảng)
\colorlet{ezDSle}{gold5!60}% hàng a), c)
\colorlet{ezDSchan}{white}% hàng b), d)
%--- dấu chọn ở cột Đúng/Sai
%    (bề rộng 0, căn giữa ô: không làm đổi bề rộng cột của ex_test, không báo tràn lề)
\newcommand{\EZdauchon}{\makebox[0pt]{\tikz[baseline=-1ex]\node[circle,fill=gold1,inner sep=0pt,
	minimum size=1.2em,text=white,font=\small]{\ding{51}};}}
\def\tickT{\EZdauchon}\def\tickF{\EZdauchon}
\AtBeginDocument{\sbox{\ccTFTab}{\EZdauchon}}
%--- nhãn a) b) c) d) màu nâu đồng (ở bản lời giải, nhãn câu đúng vẫn đóng khung như ex_test)
\renewcommand{\FalseTF}{\stepcounter{dapan}\textbf{\color{gold1}\DapAnTF\sepTF}\ignorespaces}
%--- trước khi dựng bảng: kẻ màu vàng kim (bề rộng cột giữ nguyên như ex_test); sau bảng trả lại màu đen
\newcount\EZhangTF
\newcommand{\EZbatdauTF}{%
	\global\EZhangTF=0\relax
	\arrayrulecolor{ezDSke}}
\newcommand{\EZketthucTF}{\arrayrulecolor{black}}
%--- hàng tiêu đề nền nâu đồng, chữ trắng
\makeatletter
\newcommand{\EZo@TF}[1]{\centering\bfseries\color{white}#1}
\newcommand{\EZtieudeTF}{%
	\ifcase\cotDS
		\ifnum\cotSTT=1\relax
			\def\rowHeader{\hline\rowcolor{gold1}\EZo@TF{TT}&\EZo@TF{\phatbieu}\tabularnewline}%
		\else
			\def\rowHeader{\hline\rowcolor{gold1}\EZo@TF{\phatbieu}\tabularnewline}%
		\fi
	\or
		\ifnum\cotSTT=1\relax
			\def\rowHeader{\hline\rowcolor{gold1}\EZo@TF{TT}&\EZo@TF{\phatbieu}&\EZo@TF{\Dung}\tabularnewline}%
		\else
			\def\rowHeader{\hline\rowcolor{gold1}\EZo@TF{\phatbieu}&\EZo@TF{\Dung}\tabularnewline}%
		\fi
	\or
		\ifnum\cotSTT=1\relax
			\def\rowHeader{\hline\rowcolor{gold1}\EZo@TF{TT}&\EZo@TF{\phatbieu}&\EZo@TF{\Dung}&\EZo@TF{\Sai}\tabularnewline}%
		\else
			\def\rowHeader{\hline\rowcolor{gold1}\EZo@TF{\phatbieu}&\EZo@TF{\Dung}&\EZo@TF{\Sai}\tabularnewline}%
		\fi
	\fi}
%--- tô xen kẽ từng hàng (gọi ngay sau \hline đầu mỗi hàng)
\newcommand{\EZmauhangTF}{%
	\noalign{\global\advance\EZhangTF\@ne}%
	\ifodd\EZhangTF\rowcolor{ezDSle}\else\rowcolor{ezDSchan}\fi}
%--- gắn vào \choiceTFt / \writeChoiceTFt của ex_test
\xpretocmd{\choiceTFt}{\EZbatdauTF}{}{\typeout{[preamble-eduzy-v2] Không gắn được bảng Đúng/Sai (đầu bảng)}}
\xpatchcmd{\choiceTFt}{\ifnum\dongPB=0\def\rowHeader{}\fi}{\EZtieudeTF\ifnum\dongPB=0\def\rowHeader{}\fi}
	{}{\typeout{[preamble-eduzy-v2] Không gắn được bảng Đúng/Sai (hàng tiêu đề)}}
\xpatchcmd{\writeChoiceTFt}{\hline\xdef\dotEX{\sdotEX}}{\hline\EZmauhangTF\xdef\dotEX{\sdotEX}}
	{}{\typeout{[preamble-eduzy-v2] Không gắn được bảng Đúng/Sai (tô hàng)}}
\xpatchcmd{\writeChoiceTFt}{{\end{longtable}}}{{\end{longtable}\EZketthucTF}}
	{}{\typeout{[preamble-eduzy-v2] Không gắn được bảng Đúng/Sai (cuối longtable)}}
\xpatchcmd{\writeChoiceTFt}{{\end{xtabular}\end{center}\medskip\noindent}}{{\end{xtabular}\end{center}\EZketthucTF\medskip\noindent}}
	{}{\typeout{[preamble-eduzy-v2] Không gắn được bảng Đúng/Sai (cuối xtabular)}}
\fi
%=====================================================================
\renewcommand{\parallel}{\mathrel{/\!\!/}}

%  trong file đề (tools/dan_trang/dan_trang.py tự ghi khi đặt \NeedspaceCau):
%    \def\TachCauDS{1}   câu Đúng/Sai được tách: câu dẫn (kèm hình) ở cuối trang, bảng Phát biểu sang trang sau
%                        (câu dẫn và hình không tách nhau, bảng không tách giữa hai trang)
%    \def\TachCauLG{1}   bản lời giải (\CheDo khác dethi): phần đề của câu cao hơn 1/3 trang được phép tách
%  Lệnh dùng trong file đề (giá trị do tools/dan_trang/dan_trang.py đo ở bản dethi):
%    \NeedspaceCau[<cao câu dẫn + hình>]{<cao cả khối câu>}   ngay trên \begin{ex} (câu đầu phần: trên tiêu đề phần)
%        công tắc tắt: đúng bằng \Needspace{<cao cả khối câu>} như cũ;
%        \TachCauDS=1 và có [..] (câu Đúng/Sai được chọn tách): chỉ giữ câu dẫn + hình;
%        \TachCauLG=1 ở bản lời giải: đòi tối đa 1/3 \textheight.
%    \NeedspaceBang{<cao bảng Phát biểu>}%   ngay trước \choiceTFt: \TachCauDS=1 thì giữ trọn bảng trên một trang;
%        công tắc tắt: không làm gì.
\providecommand{\TachCauDS}{0}
\providecommand{\TachCauLG}{0}
\newlength{\EZnsCau}
\newcommand{\NeedspaceCau}[2][]{%
	\setlength{\EZnsCau}{#2}%
	\ifnum\TachCauDS=1 \if\relax\detokenize{#1}\relax\else\setlength{\EZnsCau}{#1}\fi\fi
	\ifnum\TachCauLG=1 \ifdefstring{\kindOPTN}{dethi}{}{\ifdim\EZnsCau>0.3333\textheight\setlength{\EZnsCau}{0.3333\textheight}\fi}\fi
	\Needspace{\EZnsCau}}
\newcommand{\NeedspaceBang}[1]{\ifnum\TachCauDS=1 \Needspace{#1}\fi}

%--- Thang màu: đậm -> nhạt
\definecolor{gold1}{RGB}{126,96,0}    % nâu đồng đậm  #7E6000
\definecolor{gold2}{RGB}{192,144,3}   % vàng kim      #C09003
\definecolor{gold3}{RGB}{255,215,99}  % vàng nhạt     #FFD763
\definecolor{gold4}{RGB}{255,231,153} % vàng rất nhạt #FFE799
\definecolor{gold5}{RGB}{255,243,205} % vàng kem      #FFF3CD
%--- Ghi đè tông xanh/đỏ cũ
\definecolor{myblue}{RGB}{126,96,0}
\colorlet{blue}{gold1}
\colorlet{cyan}{gold2}
\colorlet{teal}{gold2}

%--- Tên chương mặc định (file đề nên tự khai báo lại)
\providecommand{\tenchuong}{}
%--- Chân trang
\providecommand{\thuonghieu}{TOÁN 12 -- Eduzy.com.vn}
\providecommand{\tacgia}{Tô Gia Bảo -- 0353703783}

%--- Chữ cạnh logo ở đầu trang
\providecommand{\khauhieu}{Eduzy.com.vn}

%====================== THÔNG SỐ ĐẦU TRANG, CHÂN TRANG (v2) ======================
%  Dải nền đậm tràn ngang trang; ô chữ vàng kem có cạnh xiên, kèm 2 vạch xiên chuyển màu.
%  Đầu trang lẻ : logo + \khauhieu (trái)   | ô chữ \tenchuong (phải)
%  Chân trang lẻ: ô chữ \thuonghieu (trái)  | Trang (số) (phải)
%  Trang chẵn soi gương trang lẻ; chân trang chẵn ghi \tacgia.
%--- Màu: chỉ dùng thang vàng kim Eduzy
\colorlet{ezNen}{gold1}% dải nền đậm
\colorlet{ezVach1}{gold3}% vạch xiên sát dải nền (đậm hơn)
\colorlet{ezVach2}{gold4}% vạch xiên sát ô chữ (nhạt hơn)
\colorlet{ezO}{gold5}% ô chữ vàng kem
%--- Kích thước
\newlength{\EZdai}\setlength{\EZdai}{1.1cm}% chiều cao dải nền
\newlength{\EZo}\setlength{\EZo}{0.68cm}% chiều cao ô chữ
\newlength{\EZxien}\setlength{\EZxien}{0.55cm}% độ lệch ngang của cạnh xiên
\newlength{\EZvach}\setlength{\EZvach}{0.75cm}% bề rộng mỗi vạch xiên
\newlength{\EZkhe}\setlength{\EZkhe}{0.12cm}% khe trắng giữa vạch và ô chữ
\newlength{\EZdem}\setlength{\EZdem}{0.4cm}% khoảng đệm từ chữ tới cạnh xiên
\newlength{\EZle}\setlength{\EZle}{1.8cm}% lề trái/phải của chữ (bằng lề thân bài)
\newlength{\EZrongmin}\setlength{\EZrongmin}{4.5cm}% ô chữ rộng tối thiểu
\newlength{\EZrong}% bề rộng chữ trong ô (tính tự động)
%--- Font
\newcommand{\EZfontO}{\fontsize{10.5pt}{12pt}\selectfont\bfseries}% chữ trong ô
\newcommand{\EZfontKH}{\fontfamily{qag}\fontsize{11pt}{12pt}\selectfont\bfseries}% \khauhieu cạnh logo
\newcommand{\EZfontTrang}{\fontsize{11.5pt}{12pt}\selectfont\bfseries}% chữ "Trang" và số trang
%--- Đo bề rộng chữ trong ô
\newcommand{\EZdorong}[1]{%
	\settowidth{\EZrong}{\EZfontO #1}%
	\ifdim\EZrong<\EZrongmin\setlength{\EZrong}{\EZrongmin}\fi}
%--- Logo + chữ trắng \khauhieu (dựng sẵn một lần khi vào \begin{document})
\newsavebox{\EZlogobox}
\AtBeginDocument{\sbox{\EZlogobox}{%
		\raisebox{-0.5\height}{\eduzylogo[1cm]}\hspace{0.15cm}%
		\raisebox{-0.5\height}{\color{white}\EZfontKH\khauhieu}}}
%--- Hai vạch xiên + khe trắng. #1: điểm gốc ở mép dưới, #2: hướng (1 sang phải, -1 sang trái),
%    #3: hướng nghiêng (1 kiểu "/", -1 kiểu "\")
\newcommand{\EZvachxien}[3]{%
	\fill[white] (#1) --++({(#3)*\EZxien},\EZo) --++({(#2)*\EZkhe},0) --++({-(#3)*\EZxien},-\EZo) --cycle;
	\fill[ezVach2] ([xshift=(#2)*\EZkhe]#1) --++({(#3)*\EZxien},\EZo) --++({(#2)*\EZvach},0) --++({-(#3)*\EZxien},-\EZo) --cycle;
	\fill[ezVach1] ([xshift=(#2)*(\EZkhe+\EZvach)]#1) --++({(#3)*\EZxien},\EZo) --++({(#2)*\EZvach},0) --++({-(#3)*\EZxien},-\EZo) --cycle;
	\draw[white,line width=0.6pt] ([xshift=(#2)*(\EZkhe+\EZvach)]#1) --++({(#3)*\EZxien},\EZo);
}
%--- Chữ trong ô (căn giữa theo chiều cao ô, không lệch vì dấu tiếng Việt)
\tikzset{ezchu/.style={inner sep=0pt,text=ezNen,text height=0.72em,text depth=0.22em}}

%====================== LOGO CHÌM GIỮA TRANG ======================
%  Vẽ cùng đầu trang nên nằm dưới chữ của thân bài, có trên mọi trang.
%  Tắt: \logochimtat (từ trang đó trở đi), bật lại: \logochimbat.
%  Đổi cỡ, độ đậm: khai báo trước \input, ví dụ \def\EZchimcao{9cm} \def\EZchimdo{0.1}
\providecommand{\EZchimcao}{10cm}% chiều cao logo chìm
\providecommand{\EZchimdo}{0.12}% độ đậm (0 = trong suốt, 1 = rõ hẳn)
\newif\ifEZchim \EZchimtrue
\newcommand{\logochimtat}{\EZchimfalse}
\newcommand{\logochimbat}{\EZchimtrue}
%  Dựng sẵn logo mờ một lần khi vào \begin{document}; transparency group làm mờ cả logo
%  một lượt nên không lộ vệt ở chỗ các dải màu chồng nhau.
\newsavebox{\EZchimtho}
\newsavebox{\EZchimbox}
\AtBeginDocument{%
	\sbox{\EZchimtho}{\eduzylogo[\EZchimcao]}%
	\sbox{\EZchimbox}{\begin{tikzpicture}
			\begin{scope}[transparency group,opacity=\EZchimdo]
				\node[inner sep=0pt] {\usebox{\EZchimtho}};
			\end{scope}
	\end{tikzpicture}}}
\newcommand{\EZvechim}{%
	\ifEZchim
		\node[inner sep=0pt] at (current page.center) {\usebox{\EZchimbox}};
	\fi}

%====================== ĐẦU TRANG - TRANG LẺ ======================
%  (logo bên trái, ô chữ \tenchuong bên phải)
\newcommand{\myfancyheadOdd}{%
	\EZdorong{\tenchuong}%
	\begin{tikzpicture}[remember picture,overlay]
		\coordinate (T) at (current page.north west);
		\coordinate (P) at (current page.north east);
		%--- logo chìm giữa trang
		\EZvechim
		%--- dải nền
		\fill[ezNen] (T) rectangle ([yshift=-\EZdai]P);
		%--- ô chữ: X là góc dưới bên trái, cạnh trái xiên kiểu "/"
		\coordinate (X) at ([xshift=-\EZle-\EZrong-\EZdem-\EZxien,yshift=-\EZdai]P);
		\fill[ezO] (X) --++(\EZxien,\EZo) -- ([yshift=-\EZdai+\EZo]P) -- ([yshift=-\EZdai]P) --cycle;
		\EZvachxien{X}{-1}{1}
		\node[ezchu,anchor=east] at ([xshift=-\EZle,yshift=-\EZdai+0.5*\EZo]P) {\EZfontO\tenchuong};
		%--- logo + tên
		\node[inner sep=0pt,anchor=west] at ([xshift=\EZle,yshift=-0.5*\EZdai]T) {\usebox{\EZlogobox}};
	\end{tikzpicture}%
}

%====================== ĐẦU TRANG - TRANG CHẴN =====================
%  (soi gương: ô chữ \tenchuong bên trái, logo bên phải)
\newcommand{\myfancyheadEven}{%
	\EZdorong{\tenchuong}%
	\begin{tikzpicture}[remember picture,overlay]
		\coordinate (T) at (current page.north west);
		\coordinate (P) at (current page.north east);
		%--- logo chìm giữa trang
		\EZvechim
		%--- dải nền
		\fill[ezNen] (T) rectangle ([yshift=-\EZdai]P);
		%--- ô chữ: X là góc dưới bên phải, cạnh phải xiên kiểu "\"
		\coordinate (X) at ([xshift=\EZle+\EZrong+\EZdem+\EZxien,yshift=-\EZdai]T);
		\fill[ezO] (X) --++(-\EZxien,\EZo) -- ([yshift=-\EZdai+\EZo]T) -- ([yshift=-\EZdai]T) --cycle;
		\EZvachxien{X}{1}{-1}
		\node[ezchu,anchor=west] at ([xshift=\EZle,yshift=-\EZdai+0.5*\EZo]T) {\EZfontO\tenchuong};
		%--- logo + tên
		\node[inner sep=0pt,anchor=east] at ([xshift=-\EZle,yshift=-0.5*\EZdai]P) {\usebox{\EZlogobox}};
	\end{tikzpicture}%
}

%====================== CHÂN TRANG - TRANG LẺ ======================
%  (ô chữ \thuonghieu bên trái, "Trang (số)" bên phải)
\newcommand{\myfancyfootOdd}{%
	\EZdorong{\thuonghieu}%
	\begin{tikzpicture}[remember picture,overlay]
		\coordinate (T) at (current page.south west);
		\coordinate (P) at (current page.south east);
		%--- dải nền
		\fill[ezNen] (T) rectangle ([yshift=\EZdai]P);
		%--- ô chữ sát mép trên dải: X là góc dưới bên phải, cạnh phải xiên kiểu "/"
		\coordinate (X) at ([xshift=\EZle+\EZrong+\EZdem,yshift=\EZdai-\EZo]T);
		\fill[ezO] ([yshift=\EZdai-\EZo]T) -- (X) --++(\EZxien,\EZo) -- ([yshift=\EZdai]T) --cycle;
		\EZvachxien{X}{1}{1}
		\node[ezchu,anchor=west] at ([xshift=\EZle,yshift=\EZdai-0.5*\EZo]T) {\EZfontO\thuonghieu};
		%--- số trang
		\node[circle,draw=white,line width=0.9pt,text=white,inner sep=0pt,minimum size=0.62cm,anchor=east]
		(so) at ([xshift=-\EZle,yshift=0.5*\EZdai]P) {\EZfontTrang\thepage};
		\node[inner sep=0pt,text=white,anchor=base east] (tr) at ([xshift=-0.15cm]so.base west) {\EZfontTrang Trang};
	\end{tikzpicture}%
}

%====================== CHÂN TRANG - TRANG CHẴN ====================
%  (soi gương: "Trang (số)" bên trái, ô chữ \tacgia bên phải)
\newcommand{\myfancyfootEven}{%
	\EZdorong{\tacgia}%
	\begin{tikzpicture}[remember picture,overlay]
		\coordinate (T) at (current page.south west);
		\coordinate (P) at (current page.south east);
		%--- dải nền
		\fill[ezNen] (T) rectangle ([yshift=\EZdai]P);
		%--- ô chữ sát mép trên dải: X là góc dưới bên trái, cạnh trái xiên kiểu "\"
		\coordinate (X) at ([xshift=-\EZle-\EZrong-\EZdem,yshift=\EZdai-\EZo]P);
		\fill[ezO] ([yshift=\EZdai-\EZo]P) -- (X) --++(-\EZxien,\EZo) -- ([yshift=\EZdai]P) --cycle;
		\EZvachxien{X}{-1}{-1}
		\node[ezchu,anchor=east] at ([xshift=-\EZle,yshift=\EZdai-0.5*\EZo]P) {\EZfontO\tacgia};
		%--- số trang
		%    (vòng tròn ẩn so0 nằm giữa dải để lấy đường chân chữ, giống trang lẻ)
		\node[circle,inner sep=0pt,minimum size=0.62cm] (so0) at ([yshift=0.5*\EZdai]T) {\EZfontTrang\phantom{\thepage}};
		\node[inner sep=0pt,text=white,anchor=base west] (tr) at ($(so0.base -| T)+(\EZle,0)$) {\EZfontTrang Trang};
		\node[circle,draw=white,line width=0.9pt,text=white,inner sep=0pt,minimum size=0.62cm,anchor=base west]
		(so) at ([xshift=0.15cm]tr.base east) {\EZfontTrang\thepage};
	\end{tikzpicture}%
}

%--- Giữ tương thích với các file cũ đang gọi \myfancyhead, \myfancyfoot
\newcommand{\myfancyhead}{\myfancyheadOdd}
\newcommand{\myfancyfoot}{\myfancyfootOdd}

%======================= GẮN VÀO TRANG =============================
\pagestyle{fancy}
\fancyhf{}
\fancyhead[LO]{\myfancyheadOdd}
\fancyhead[RE]{\myfancyheadEven}
\fancyfoot[LO]{\myfancyfootOdd}
\fancyfoot[RE]{\myfancyfootEven}
\renewcommand{\headrulewidth}{0pt}
\renewcommand{\footrulewidth}{0pt}
\fancypagestyle{plain}{%
	\fancyhf{}
	\fancyhead[LO]{\myfancyheadOdd}
	\fancyhead[RE]{\myfancyheadEven}
	\fancyfoot[LO]{\myfancyfootOdd}
	\fancyfoot[RE]{\myfancyfootEven}
	\renewcommand{\headrulewidth}{0pt}
	\renewcommand{\footrulewidth}{0pt}
}
%=================== HẾT PHẦN GỘP TỪ eduzy-de.tex ===================
\newcommand{\indam}[1]{
	{\fontsize{12pt}{9pt}\sffamily \selectfont\bfseries\color{gold1} #1}
}
\newcommand{\ind}[1]{
	{\fontsize{14pt}{9pt}\fontfamily{qag} \selectfont\bfseries\color{gold1} #1}
}
\renewcommand{\FalseEX}{\stepcounter{dapan}\textbf{\selectfont\color{gold1}\Alph{dapan}.}\ignorespaces}
\renewtheorem{ex}{\color{gold1}Câu}
\renewcommand{\baselinestretch}{1.1}
\newcommand{\hoac}[1]{
	\left[\begin{aligned}#1\end{aligned}\right.}
\newcommand{\heva}[1]{
	\left\{\begin{aligned}#1\end{aligned}\right.}
\usepackage{esvect}
\def\vec{\vv}
\def\vt{\vv}
\def\overrightarrow{\vv}
%====================== LOGO EDUZY VẼ BẰNG TIKZ ======================
%  Đường viền dò từ logo gốc; dải màu vàng kim đo trực tiếp trên ảnh gốc
%  (nghiêng 21 độ, tô bằng 80 dải hẹp nên hiển thị giống nhau trên mọi trình xem PDF).
%  Dùng: \eduzylogo          -> cao 2.4cm (mặc định)
%        \eduzylogo[1.5cm]   -> cao 1.5cm
\definecolor{ezText}{RGB}{245,200,46}
\newcommand{\eduzylogo}[1][2.4cm]{%
	\resizebox{!}{#1}{%
		\begin{tikzpicture}[x=0.01cm,y=0.01cm]
			\path[use as bounding box] (0,0) rectangle (224,-224);
			% mũ tốt nghiệp + vương miện
			\begin{scope}
				\clip
				(119.4,-176.1)..controls(118.9,-176.1)and(116.8,-176)..(114.6,-176)
				..controls(101.4,-176)and(85.8,-175.9)..(83.7,-175.7)
				..controls(81.8,-175.6)and(81,-175.6)..(80.2,-175.7)
				..controls(78.8,-175.9)and(77.7,-175.9)..(76.2,-175.7)
				..controls(74.9,-175.5)and(70.8,-175.5)..(64.3,-175.7)
				..controls(59.1,-175.9)and(58.3,-175.7)..(58.6,-173.7)
				..controls(58.6,-172.8)and(59.6,-169.6)..(60.3,-168)
				..controls(60.8,-166.9)and(61.4,-165.8)..(62.3,-165)
				..controls(63.1,-164.2)and(64.3,-163.4)..(65,-163.2)
				..controls(65.5,-163.1)and(72.3,-163)..(76.6,-163.1)
				..controls(77.9,-163.1)and(80.1,-163.1)..(81.4,-163)
				..controls(82.8,-162.9)and(86.3,-162.9)..(89.2,-162.9)
				..controls(93,-162.9)and(94.9,-162.8)..(96,-162.7)
				..controls(97.2,-162.5)and(101,-162.5)..(114,-162.4)
				..controls(123,-162.3)and(134.9,-162.2)..(140.2,-162.1)
				..controls(150.1,-162)and(152.2,-162.1)..(153.6,-162.4)
				..controls(154.4,-162.6)and(157.1,-163.8)..(158.1,-164.4)
				..controls(158.8,-164.8)and(159.7,-165.7)..(159.9,-166.3)
				..controls(160,-166.5)and(160.2,-167.1)..(160.2,-167.6)
				..controls(160.3,-168)and(160.6,-169.4)..(160.9,-170.6)
				..controls(161.2,-171.7)and(161.5,-173.1)..(161.6,-173.6)
				..controls(161.8,-174.4)and(161.8,-174.5)..(161.5,-174.9)
				..controls(161.2,-175.6)and(160.8,-175.9)..(159.9,-176)
				..controls(159.2,-176.1)and(154.1,-176.2)..(151.5,-176.2)
				..controls(148.1,-176.1)and(122.6,-176.1)..(121.6,-176.2)
				..controls(120.8,-176.2)and(119.8,-176.2)..(119.4,-176.1)--cycle
				(130.4,-160)..controls(129,-160)and(119.8,-160)..(110.1,-159.9)
				..controls(100.3,-159.8)and(88.5,-159.7)..(83.9,-159.6)
				..controls(77.8,-159.5)and(75.3,-159.5)..(74.5,-159.6)
				..controls(73.9,-159.7)and(72.5,-159.8)..(71.3,-159.8)
				..controls(69.6,-159.8)and(69,-159.8)..(68.5,-159.6)
				..controls(67.6,-159.3)and(66.8,-158.6)..(66.2,-157.6)
				..controls(66,-157.3)and(65.6,-156.5)..(65.3,-156.1)
				..controls(64.6,-155.1)and(64.5,-154.9)..(63.7,-152.9)
				..controls(63.2,-151.6)and(62.9,-150.9)..(61.8,-149.6)
				..controls(61.5,-149.2)and(61.1,-148.5)..(60.9,-147.9)
				..controls(60.7,-147.4)and(60.3,-146.6)..(60.1,-146.1)
				..controls(59.9,-145.5)and(59.5,-144.9)..(59,-144.3)
				..controls(58.5,-143.7)and(58.2,-143.2)..(57.9,-142.5)
				..controls(57.7,-141.9)and(57.3,-141.1)..(57,-140.7)
				..controls(56.6,-140.3)and(56.1,-139.3)..(55.8,-138.5)
				..controls(55.5,-137.8)and(55.1,-136.8)..(54.9,-136.5)
				..controls(54.7,-136.1)and(54.3,-135.2)..(53.9,-134.5)
				..controls(53.6,-133.8)and(53,-132.7)..(52.6,-132)
				..controls(51.6,-130.5)and(51.4,-130.2)..(51.1,-129)
				..controls(50.6,-127.5)and(50.6,-127.4)..(49.7,-126.3)
				..controls(49.2,-125.7)and(48.6,-124.8)..(48.3,-124.2)
				..controls(48,-123.6)and(47.5,-122.7)..(47.3,-122.2)
				..controls(47,-121.8)and(46.5,-120.7)..(46.1,-119.8)
				..controls(45.8,-118.9)and(45.2,-117.8)..(44.9,-117.2)
				..controls(44.2,-116.1)and(43.3,-114.4)..(42.3,-112.2)
				..controls(41.8,-111.3)and(41.3,-110.4)..(41.1,-110.1)
				..controls(40.5,-109.4)and(40.3,-109)..(40,-107.8)
				..controls(39.8,-107)and(39.6,-106.5)..(39.2,-106)
				..controls(38.9,-105.6)and(38.4,-104.7)..(38,-104)
				..controls(37.7,-103.4)and(37.2,-102.5)..(37,-102.1)
				..controls(36.7,-101.7)and(36.3,-100.8)..(36,-100.1)
				..controls(35.6,-98.9)and(35.3,-98.4)..(34.6,-97.6)
				..controls(34.2,-97.1)and(34.2,-96.3)..(34.6,-96)
				..controls(35,-95.7)and(35.5,-95.9)..(36.2,-96.7)
				..controls(36.5,-97)and(37.2,-97.5)..(37.6,-97.7)..controls(38,-98)and(38.6,-98.4)..(39.1,-98.9)
				..controls(39.6,-99.4)and(40.3,-100)..(40.6,-100.3)
				..controls(41.3,-100.8)and(42.4,-101.8)..(43.7,-102.9)
				..controls(44,-103.2)and(44.6,-103.7)..(45,-103.9)
				..controls(45.4,-104.2)and(45.9,-104.6)..(46.1,-104.8)
				..controls(47.1,-105.9)and(48.8,-107.1)..(50.5,-107.9)
				..controls(51.8,-108.6)and(52.3,-109)..(53.4,-110)
				..controls(54.8,-111.2)and(55.6,-111.6)..(57,-112.1)
				..controls(57.5,-112.3)and(58.5,-112.8)..(59.4,-113.3)
				..controls(60.2,-113.7)and(61.3,-114.3)..(61.9,-114.5)
				..controls(62.4,-114.7)and(63.3,-115.1)..(63.8,-115.3)
				..controls(64.4,-115.6)and(65.4,-115.9)..(66.2,-116.1)
				..controls(67,-116.3)and(68.4,-116.6)..(69.2,-116.9)
				..controls(72.3,-117.9)and(72.8,-118)..(78,-118.1)
				..controls(84,-118.3)and(86.1,-118.1)..(88.1,-117.1)
				..controls(89.4,-116.5)and(93,-115)..(93.7,-114.9)
				..controls(94.8,-114.6)and(95.4,-114.2)..(95.9,-113.4)
				..controls(96.2,-113)and(96.7,-112.5)..(97.1,-112.2)
				..controls(98,-111.6)and(101.3,-108.4)..(102.1,-107.2)
				..controls(102.4,-106.8)and(102.9,-106.2)..(103.2,-105.8)
				..controls(103.5,-105.5)and(104,-104.8)..(104.2,-104.2)
				..controls(104.5,-103.6)and(105,-102.8)..(105.3,-102.4)
				..controls(105.6,-102)and(106.1,-101.2)..(106.5,-100.6)
				..controls(107.1,-99.5)and(108.6,-97.5)..(109.5,-96.5)
				..controls(109.9,-96.1)and(110.1,-96)..(110.4,-96)
				..controls(111.1,-96)and(111.6,-96.5)..(111.8,-97.6)
				..controls(112,-98.6)and(112.3,-99.3)..(112.8,-99.9)
				..controls(113.7,-100.9)and(115.2,-103.1)..(115.7,-104.2)
				..controls(116.1,-105)and(116.5,-105.5)..(116.9,-106)
				..controls(117.3,-106.4)and(117.8,-106.9)..(118,-107.3)
				..controls(118.2,-107.6)and(118.6,-108)..(119,-108.3)
				..controls(119.6,-108.8)and(121.1,-110.3)..(121.8,-111.2)
				..controls(122.1,-111.7)and(122.5,-112)..(123.4,-112.4)
				..controls(124.8,-113.2)and(125.8,-113.8)..(127,-114.6)
				..controls(127.5,-114.9)and(128.1,-115.3)..(128.2,-115.4)
				..controls(128.4,-115.5)and(129.3,-115.8)..(130.2,-116)
				..controls(131.1,-116.2)and(132.5,-116.6)..(133.3,-116.9)
				..controls(135.7,-117.9)and(136.1,-117.9)..(142.5,-117.8)
				..controls(145.8,-117.8)and(148.4,-117.7)..(148.8,-117.6)
				..controls(149.2,-117.6)and(150.2,-117.2)..(151,-116.8)
				..controls(152,-116.4)and(152.8,-116)..(153.6,-115.9)
				..controls(154.5,-115.7)and(155.2,-115.4)..(156.3,-114.9)
				..controls(157.1,-114.5)and(158.2,-114)..(158.7,-113.8)
				..controls(159.1,-113.7)and(159.8,-113.3)..(160.2,-113)
				..controls(160.6,-112.7)and(161.3,-112.4)..(161.7,-112.2)
				..controls(162.7,-111.9)and(163.2,-111.7)..(164.1,-111.1)
				..controls(164.5,-110.8)and(165.3,-110.4)..(165.9,-110)
				..controls(166.5,-109.7)and(167.3,-109.2)..(167.6,-108.9)
				..controls(167.9,-108.6)and(168.6,-108.2)..(169,-108)
				..controls(169.4,-107.8)and(170.2,-107.3)..(170.7,-106.9)
				..controls(171.2,-106.5)and(172.1,-105.8)..(172.7,-105.4)
				..controls(173.2,-105)and(174.1,-104.4)..(174.5,-104)
				..controls(174.9,-103.7)and(175.7,-103.1)..(176.2,-102.6)
				..controls(176.7,-102.2)and(177.4,-101.5)..(177.9,-101.2)
				..controls(178.3,-100.8)and(178.8,-100.4)..(179,-100.2)
				..controls(179.2,-100)and(179.6,-99.6)..(180,-99.3)
				..controls(180.4,-99.1)and(180.9,-98.6)..(181.2,-98.4)
				..controls(181.5,-98.1)and(182.1,-97.7)..(182.5,-97.5)
				..controls(183.2,-97.2)and(183.5,-96.9)..(184.3,-95.9)
				..controls(184.6,-95.6)and(184.6,-95.6)..(184.8,-95.7)
				..controls(185,-95.8)and(185.1,-96)..(185.1,-96.2)
				..controls(185.1,-96.5)and(185,-96.7)..(184.6,-97.2)
				..controls(184.1,-97.8)and(184,-97.9)..(183.9,-98.8)
				..controls(183.7,-100.3)and(183.3,-101.3)..(182.9,-101.8)
				..controls(182.6,-102.1)and(182.3,-102.6)..(182.2,-102.9)
				..controls(182,-103.3)and(181.6,-104)..(181.2,-104.6)
				..controls(180.9,-105.2)and(180.5,-106.1)..(180.4,-106.6)
				..controls(179.9,-108)and(179.7,-108.5)..(179,-109.4)
				..controls(178.6,-109.8)and(178.2,-110.5)..(177.9,-111)
				..controls(177.7,-111.5)and(177.3,-112.3)..(177,-112.7)
				..controls(176.6,-113.2)and(176.4,-113.7)..(176.3,-114.2)
				..controls(176.2,-114.7)and(175.9,-115.2)..(175.4,-116)
				..controls(175,-116.6)and(174.4,-117.7)..(174,-118.5)
				..controls(173.7,-119.2)and(173.3,-120.1)..(173,-120.5)
				..controls(172.8,-120.9)and(172.4,-121.8)..(172.1,-122.5)
				..controls(171.8,-123.2)and(171.3,-124.1)..(171,-124.6)
				..controls(170.8,-125.1)and(170.3,-126)..(169.9,-126.6)
				..controls(169.6,-127.3)and(169.2,-128.1)..(169,-128.3)
				..controls(168.8,-128.6)and(168.5,-129.3)..(168.3,-129.8)
				..controls(168.1,-130.3)and(167.8,-131.1)..(167.5,-131.6)
				..controls(166.9,-132.5)and(164.2,-137.9)..(164.2,-138.1)
				..controls(164.2,-138.1)and(163.9,-138.8)..(163.4,-139.7)
				..controls(162.9,-140.5)and(162.4,-141.5)..(162.3,-141.9)
				..controls(161.9,-143)and(161.6,-143.4)..(160.9,-144.1)
				..controls(160.4,-144.6)and(160.3,-144.8)..(160.1,-145.4)
				..controls(160,-145.8)and(159.6,-146.8)..(159.2,-147.8)
				..controls(158.7,-148.8)and(158.1,-150)..(157.9,-150.6)
				..controls(157.6,-151.1)and(157.2,-151.9)..(156.9,-152.4)
				..controls(156.5,-152.8)and(156.1,-153.7)..(155.8,-154.4)
				..controls(155.6,-155)and(155.2,-155.9)..(154.9,-156.3)
				..controls(154.7,-156.7)and(154.4,-157.4)..(154.2,-157.9)
				..controls(154.1,-158.4)and(153.8,-159)..(153.5,-159.3)
				..controls(152.8,-160.1)and(153,-160.1)..(142.4,-160.1)
				..controls(137.2,-160.1)and(131.8,-160.1)..(130.4,-160)--cycle
				(141.5,-103.7)..controls(140.8,-103)and(139.8,-102.1)..(139.2,-101.6)
				..controls(138.5,-101.2)and(137.9,-100.6)..(137.8,-100.4)
				..controls(137,-99.5)and(134,-96.5)..(133.2,-96)
				..controls(132.7,-95.7)and(131.8,-94.9)..(131.1,-94.3)
				..controls(129.4,-92.7)and(127.9,-91.6)..(126.6,-91.1)
				..controls(125.2,-90.5)and(124.1,-90)..(122.9,-89.1)
				..controls(122.1,-88.6)and(121.3,-88.2)..(119.9,-87.7)
				..controls(118.8,-87.3)and(117.5,-86.8)..(116.9,-86.5)
				..controls(114.3,-85.5)and(113.7,-85.4)..(109.7,-85.5)
				..controls(105.3,-85.7)and(103.9,-85.9)..(101.6,-87)
				..controls(100.7,-87.4)and(99.4,-87.8)..(97.9,-88)
				..controls(97,-88.2)and(96.6,-88.3)..(95.8,-88.9)
				..controls(95.4,-89.2)and(94.7,-89.6)..(94.2,-89.8)
				..controls(92.9,-90.3)and(92.3,-90.7)..(90.9,-92.1)
				..controls(90.2,-92.7)and(89.4,-93.3)..(89,-93.5)..controls(87,-94.5)and(85.2,-96)..(83.6,-97.9)
				..controls(83.3,-98.2)and(82.5,-98.9)..(81.8,-99.4)
				..controls(80.7,-100.1)and(80.5,-100.3)..(80,-101.1)
				..controls(78.7,-103.1)and(77.2,-103.9)..(76.4,-103)
				..controls(75.7,-102.2)and(75.7,-103.2)..(75.7,-83.6)
				..controls(75.6,-73.6)and(75.6,-64.8)..(75.5,-64)
				..controls(75.5,-63.2)and(75.5,-61.4)..(75.6,-59.8)--(75.7,-57.1)--(76,-56.8)
				..controls(76.8,-56.2)and(78,-56.3)..(79.8,-57.2)
				..controls(80.4,-57.5)and(81.3,-57.8)..(81.8,-57.9)
				..controls(82.3,-58.1)and(83,-58.4)..(83.6,-58.8)..controls(84.3,-59.2)and(85,-59.5)..(86,-59.8)
				..controls(86.8,-60.1)and(88,-60.6)..(88.7,-60.9)
				..controls(89.3,-61.2)and(90.6,-61.7)..(91.4,-62)
				..controls(92.2,-62.3)and(93.7,-62.9)..(94.6,-63.3)
				..controls(95.6,-63.7)and(96.8,-64.1)..(97.4,-64.3)
				..controls(98,-64.5)and(99,-64.9)..(99.7,-65.2)
				..controls(100.4,-65.6)and(101.5,-66)..(102.3,-66.2)
				..controls(104.2,-66.8)and(105,-67.1)..(106.2,-67.8)
				..controls(107.9,-68.7)and(110.2,-69.4)..(111.3,-69.2)
				..controls(112.5,-69)and(115,-68.2)..(115.6,-67.8)
				..controls(117.3,-66.7)and(117.6,-66.6)..(119.5,-66.2)
				..controls(121.9,-65.8)and(122.1,-65.7)..(123,-65.1)
				..controls(124.1,-64.4)and(124.4,-64.3)..(126,-63.9)
				..controls(127,-63.7)and(127.7,-63.4)..(128.4,-63.1)
				..controls(129,-62.8)and(129.7,-62.5)..(130.6,-62.4)
				..controls(132.1,-62.1)and(134,-61.5)..(135.2,-61)
				..controls(135.7,-60.8)and(136.7,-60.5)..(137.3,-60.3)
				..controls(138.4,-60)and(139.3,-59.7)..(141.1,-59)
				..controls(141.6,-58.7)and(142,-58.6)..(142.5,-58.6)
				..controls(143.2,-58.6)and(143.3,-58.7)..(143.6,-59)
				..controls(144,-59.5)and(144.1,-59.9)..(144,-60.8)
				..controls(144,-61.2)and(144,-71.1)..(144,-82.8)--(143.9,-104.1)--(143.7,-104.5)
				..controls(143.5,-104.7)and(143.3,-104.8)..(143,-104.9)
				..controls(142.6,-104.9)and(142.5,-104.8)..(141.5,-103.7)--cycle
				(27.8,-91.4)..controls(27.1,-90.9)and(26.3,-90)..(26.1,-89.5)
				..controls(26,-89.2)and(25.8,-88.6)..(25.7,-87.9)
				..controls(25.4,-86.2)and(25.6,-84.8)..(26.5,-83.5)
				..controls(27.4,-82.3)and(28.3,-82)..(29.4,-82.3)
				..controls(31.1,-82.8)and(32.1,-84.9)..(32,-87.4)
				..controls(31.9,-89.3)and(31.2,-90.6)..(29.9,-91.4)
				..controls(29.1,-91.9)and(28.6,-91.9)..(27.8,-91.4)--cycle
				(190.1,-91.3)..controls(189.7,-91.2)and(189.3,-90.9)..(189.1,-90.7)
				..controls(188.3,-89.8)and(187.6,-87.7)..(187.6,-86.2)
				..controls(187.6,-85.5)and(187.7,-85.2)..(188.2,-84.4)
				..controls(189.1,-82.6)and(190.1,-81.8)..(191.3,-82.1)
				..controls(192,-82.3)and(193.5,-83.7)..(193.8,-84.4)
				..controls(194.1,-85.1)and(194.2,-87.1)..(194,-88.3)
				..controls(193.8,-89.8)and(193.3,-90.5)..(192.3,-91.3)
				..controls(191.7,-91.7)and(190.9,-91.7)..(190.1,-91.3)--cycle
				(173.7,-83.4)..controls(172.2,-82.5)and(172,-82.5)..(170.6,-82.4)
				..controls(169.7,-82.4)and(168.8,-82.4)..(168.3,-82.5)
				..controls(167.2,-82.7)and(166.9,-82.7)..(166.6,-82.2)
				..controls(166.3,-81.8)and(166.3,-81)..(166.8,-80)
				..controls(167.2,-78.9)and(167.4,-77.4)..(167.7,-74.1)
				..controls(167.8,-72.8)and(167.9,-71.8)..(167.9,-71.7)
				..controls(168,-71.7)and(168,-70.2)..(168,-68.4)..controls(168,-65.4)and(168,-64.9)..(167.7,-63)
				..controls(167.3,-60.5)and(167.1,-59.9)..(166.4,-58.4)
				..controls(166.1,-57.7)and(165.8,-56.9)..(165.7,-56.6)
				..controls(165.7,-56.2)and(165.4,-55.6)..(165.2,-55.2)
				..controls(164.6,-53.8)and(161.9,-50.7)..(161,-50.3)
				..controls(159.8,-49.7)and(157.9,-50)..(155.2,-51.2)
				..controls(154.7,-51.4)and(153.4,-51.8)..(152.2,-52.1)
				..controls(150.8,-52.4)and(149.4,-52.9)..(148.6,-53.2)
				..controls(147.8,-53.6)and(146.7,-54)..(146,-54.2)
				..controls(145.4,-54.4)and(144.5,-54.7)..(144,-54.9)
				..controls(143.6,-55.2)and(142.7,-55.5)..(141.8,-55.7)
				..controls(141,-55.9)and(139.6,-56.3)..(138.7,-56.6)
				..controls(137.8,-57)and(136.3,-57.4)..(135.3,-57.6)
				..controls(133.8,-57.9)and(133.2,-58.1)..(132.1,-58.7)
				..controls(131.1,-59.2)and(130.3,-59.4)..(129.4,-59.6)
				..controls(128.7,-59.8)and(127.8,-60)..(127.4,-60.1)
				..controls(127,-60.3)and(126,-60.5)..(125.1,-60.7)
				..controls(123.2,-61.1)and(121.7,-61.5)..(120.2,-62.1)
				..controls(119.7,-62.3)and(118.1,-62.8)..(116.7,-63.2)
				..controls(114.6,-63.9)and(114,-64.1)..(113,-64.7)
				..controls(111.9,-65.4)and(111.8,-65.4)..(111,-65.5)
				..controls(109.9,-65.5)and(109.3,-65.4)..(107.7,-64.7)
				..controls(107,-64.4)and(106,-64)..(105.3,-63.9)
				..controls(103.5,-63.5)and(102.2,-63)..(101.2,-62.5)
				..controls(100.1,-61.8)and(99.3,-61.5)..(97.8,-61.1)
				..controls(97.2,-61)and(96.2,-60.7)..(95.7,-60.4)
				..controls(95.1,-60.1)and(94.2,-59.8)..(93.5,-59.6)
				..controls(92.9,-59.4)and(91.7,-59)..(91,-58.7)..controls(90.2,-58.4)and(88.9,-58)..(88,-57.7)
				..controls(87.2,-57.5)and(86.1,-57.1)..(85.5,-56.8)
				..controls(85,-56.6)and(84.1,-56.2)..(83.5,-56)..controls(82.9,-55.7)and(81.8,-55.2)..(81,-54.9)
				..controls(80.1,-54.5)and(79.1,-54.2)..(78.7,-54.1)
				..controls(77.9,-53.8)and(77.3,-53.6)..(76.2,-52.9)
				..controls(75.7,-52.6)and(75.2,-52.4)..(74.1,-52.2)
				..controls(72.5,-51.9)and(72.4,-51.8)..(71.3,-51.1)
				..controls(69.9,-50.1)and(69.4,-49.9)..(67.9,-49.8)
				..controls(67,-49.7)and(66.6,-49.6)..(66.3,-49.4)
				..controls(66.1,-49.3)and(65.7,-49)..(65.4,-48.8)
				..controls(65.1,-48.7)and(64.2,-48.4)..(63.6,-48.2)
				..controls(62.9,-48)and(61.7,-47.6)..(60.9,-47.2)
				..controls(60.1,-46.8)and(58.9,-46.4)..(58.2,-46.2)
				..controls(57.6,-46)and(56.8,-45.7)..(56.5,-45.5)
				..controls(54.9,-44.6)and(54.6,-44.5)..(53.2,-44.2)
				..controls(51.7,-43.8)and(51,-43.6)..(49.9,-42.9)
				..controls(49.5,-42.6)and(48.9,-42.4)..(47.5,-42)
				..controls(45.7,-41.6)and(45.6,-41.5)..(44.3,-40.6)
				..controls(43.5,-40.1)and(42.8,-39.6)..(42.7,-39.4)
				..controls(42.4,-38.7)and(43,-38.1)..(44.5,-37.6)..controls(44.9,-37.5)and(46,-37)..(46.9,-36.6)
				..controls(47.9,-36.1)and(48.8,-35.8)..(49.6,-35.6)
				..controls(50.5,-35.4)and(51.3,-35.2)..(52.4,-34.6)
				..controls(53.3,-34.1)and(54.5,-33.7)..(54.9,-33.5)
				..controls(55.4,-33.4)and(56.4,-33)..(57.2,-32.6)
				..controls(58.6,-32)and(60.1,-31.5)..(61.6,-31.2)
				..controls(62.3,-31.1)and(62.7,-31)..(63.2,-30.7)
				..controls(64.1,-30)and(64.3,-29.9)..(65.3,-29.7)
				..controls(65.7,-29.6)and(66.5,-29.2)..(67.1,-28.9)
				..controls(67.6,-28.6)and(68.8,-28.2)..(69.6,-27.9)
				..controls(70.5,-27.7)and(71.6,-27.3)..(72.1,-27)
				..controls(72.6,-26.7)and(73.8,-26.2)..(74.8,-25.9)
				..controls(75.7,-25.5)and(76.7,-25.1)..(77.1,-25)
				..controls(77.4,-24.8)and(78.5,-24.4)..(79.6,-24.1)
				..controls(80.6,-23.8)and(82,-23.2)..(82.8,-22.9)
				..controls(83.5,-22.5)and(84.8,-22)..(85.8,-21.8)
				..controls(86.7,-21.5)and(87.8,-21.1)..(88.1,-20.9)
				..controls(88.4,-20.7)and(89.3,-20.3)..(90.1,-20.1)
				..controls(90.8,-19.8)and(91.9,-19.4)..(92.4,-19.1)
				..controls(93.5,-18.5)and(94.4,-18.2)..(95.4,-18.1)
				..controls(95.7,-18)and(96.4,-17.8)..(96.9,-17.5)
				..controls(97.5,-17.3)and(98.4,-16.9)..(99,-16.6)
				..controls(99.7,-16.3)and(100.9,-15.9)..(101.9,-15.6)
				..controls(102.9,-15.3)and(104.4,-14.7)..(105.3,-14.3)
				..controls(108.1,-13.2)and(108.5,-13.1)..(110.1,-13.1)
				..controls(111.4,-13.1)and(111.5,-13.1)..(112.4,-13.5)
				..controls(113.2,-14)and(115,-14.8)..(116.6,-15.4)
				..controls(116.9,-15.5)and(117.9,-15.8)..(118.8,-16.1)
				..controls(120.9,-16.6)and(121.5,-16.9)..(122.6,-17.9)
				..controls(123,-18.3)and(123.7,-18.8)..(123.9,-18.9)
				..controls(124.6,-19.4)and(126,-19.8)..(127.9,-20.2)
				..controls(129.6,-20.5)and(129.8,-20.6)..(131.1,-21.4)
				..controls(131.5,-21.7)and(132.4,-22.1)..(133.1,-22.3)
				..controls(136.1,-23.2)and(138.3,-24.1)..(140.2,-25)
				..controls(141,-25.5)and(141.5,-25.6)..(142.5,-25.8)
				..controls(145.1,-26.3)and(146,-26.6)..(149.4,-28.1)
				..controls(151.3,-28.9)and(153.5,-29.8)..(154.3,-30.1)
				..controls(155.1,-30.4)and(156.2,-30.8)..(156.6,-31)
				..controls(157.1,-31.2)and(158.2,-31.6)..(159.1,-31.8)
				..controls(160.6,-32.2)and(160.9,-32.3)..(161.7,-32.9)
				..controls(162.4,-33.3)and(162.9,-33.6)..(163.7,-33.8)
				..controls(164.3,-34)and(164.8,-34.1)..(164.9,-34.1)
				..controls(165,-34.1)and(165.6,-34.4)..(166.4,-34.8)
				..controls(167.3,-35.2)and(168.1,-35.4)..(169,-35.6)
				..controls(170,-35.8)and(170.5,-36)..(171.7,-36.6)
				..controls(172.4,-37)and(173.3,-37.4)..(173.6,-37.4)
				..controls(175.6,-37.7)and(176.2,-38)..(176.7,-39.1)
				..controls(176.9,-39.6)and(176.9,-39.7)..(176.7,-40.4)
				..controls(176.4,-41.6)and(175.8,-42.7)..(175.2,-43.4)
				..controls(174.5,-44.1)and(172.6,-45.5)..(171.7,-45.8)
				..controls(170.6,-46.2)and(169.1,-46)..(168.2,-45.4)
				..controls(167.9,-45.2)and(167.3,-44.8)..(166.9,-44.4)
				..controls(166.4,-44.1)and(165.8,-43.6)..(165.4,-43.2)
				..controls(165.1,-42.9)and(164.4,-42.3)..(163.9,-41.9)
				..controls(163.3,-41.5)and(162.7,-40.9)..(162.4,-40.5)
				..controls(161.4,-39.2)and(159.9,-38)..(158.8,-37.8)
				..controls(158.4,-37.8)and(158.3,-37.7)..(158.2,-37.2)
				..controls(158,-36.8)and(157.4,-36.1)..(157.1,-36.1)
				..controls(157,-36.1)and(156.7,-36.1)..(156.5,-36)
				..controls(156.1,-35.9)and(155.9,-35.7)..(155.7,-35.2)
				..controls(155.3,-34.6)and(155.1,-34.4)..(154,-34.1)
				..controls(153.6,-34)and(153.2,-33.7)..(152.9,-33.5)
				..controls(152.4,-32.9)and(151.6,-32.4)..(151.3,-32.4)
				..controls(150.9,-32.4)and(150.4,-32.6)..(150.4,-32.9)
				..controls(150.4,-33.2)and(150.8,-33.4)..(151.4,-33.6)
				..controls(152.1,-33.7)and(152.8,-34.1)..(153.5,-34.9)
				..controls(154.1,-35.5)and(154.3,-35.6)..(155,-35.8)
				..controls(155.7,-36)and(155.7,-36)..(156,-36.6)
				..controls(156.3,-37.3)and(156.7,-37.7)..(157.5,-37.8)
				..controls(158.1,-37.9)and(158.2,-38)..(158.4,-38.8)
				..controls(158.5,-39.4)and(158.6,-39.6)..(159.6,-40.4)
				..controls(160.2,-41)and(160.9,-41.6)..(161.3,-41.8)
				..controls(161.7,-42)and(161.9,-42.3)..(162.2,-42.8)
				..controls(162.4,-43.2)and(162.8,-43.8)..(163.1,-44.1)
				..controls(163.9,-44.8)and(165.1,-46.7)..(165.7,-48)
				..controls(166.2,-49.3)and(166.4,-49.7)..(167.2,-50.5)
				..controls(167.5,-50.8)and(167.9,-51.4)..(168.1,-51.8)
				..controls(168.3,-52.2)and(168.8,-53)..(169.1,-53.6)
				..controls(169.6,-54.3)and(169.8,-54.9)..(169.9,-55.6)
				..controls(170.4,-57.8)and(170.5,-58.1)..(171.2,-59.5)
				..controls(171.6,-60.3)and(172,-61.5)..(172.2,-62.1)
				..controls(172.3,-62.8)and(172.7,-63.7)..(172.9,-64.1)
				..controls(173.2,-64.6)and(173.4,-65.3)..(173.4,-65.6)
				..controls(173.8,-68.8)and(174.2,-70.6)..(174.9,-72.6)
				..controls(175.5,-74.5)and(175.9,-77.1)..(176.1,-80.2)
				..controls(176.3,-83.4)and(175.5,-84.4)..(173.7,-83.4)--cycle;
				\definecolor{ezs}{RGB}{247,179,26}\fill[ezs](-5.7,-156.6)--(-3.4,-157.5)--(61.9,12.7)--(59.7,13.6)--cycle;
				\definecolor{ezs}{RGB}{247,179,22}\fill[ezs](-3.8,-157.3)--(-1.5,-158.2)--(63.9,12)--(61.6,12.9)--cycle;
				\definecolor{ezs}{RGB}{248,179,19}\fill[ezs](-1.8,-158.1)--(0.4,-158.9)--(65.8,11.2)--(63.5,12.1)--cycle;
				\definecolor{ezs}{RGB}{248,178,16}\fill[ezs](0.1,-158.8)--(2.4,-159.7)--(67.7,10.5)--(65.4,11.4)--cycle;
				\definecolor{ezs}{RGB}{249,178,15}\fill[ezs](2,-159.5)--(4.3,-160.4)--(69.6,9.8)--(67.3,10.7)--cycle;
				\definecolor{ezs}{RGB}{251,177,17}\fill[ezs](3.9,-160.3)--(6.2,-161.1)--(71.5,9)--(69.2,9.9)--cycle;
				\definecolor{ezs}{RGB}{252,179,18}\fill[ezs](5.8,-161)--(8.1,-161.9)--(73.4,8.3)--(71.1,9.2)--cycle;
				\definecolor{ezs}{RGB}{252,182,17}\fill[ezs](7.7,-161.7)--(10,-162.6)--(75.3,7.6)--(73,8.5)--cycle;
				\definecolor{ezs}{RGB}{252,185,17}\fill[ezs](9.6,-162.5)--(11.9,-163.4)--(77.2,6.8)--(75,7.7)--cycle;
				\definecolor{ezs}{RGB}{251,189,18}\fill[ezs](11.5,-163.2)--(13.8,-164.1)--(79.2,6.1)--(76.9,7)--cycle;
				\definecolor{ezs}{RGB}{251,191,22}\fill[ezs](13.4,-163.9)--(15.7,-164.8)--(81.1,5.4)--(78.8,6.2)--cycle;
				\definecolor{ezs}{RGB}{252,193,25}\fill[ezs](15.4,-164.7)--(17.6,-165.6)--(83,4.6)--(80.7,5.5)--cycle;
				\definecolor{ezs}{RGB}{253,195,27}\fill[ezs](17.3,-165.4)--(19.6,-166.3)--(84.9,3.9)--(82.6,4.8)--cycle;
				\definecolor{ezs}{RGB}{253,196,31}\fill[ezs](19.2,-166.1)--(21.5,-167)--(86.8,3.2)--(84.5,4)--cycle;
				\definecolor{ezs}{RGB}{254,199,35}\fill[ezs](21.1,-166.9)--(23.4,-167.8)--(88.7,2.4)--(86.4,3.3)--cycle;
				\definecolor{ezs}{RGB}{253,201,38}\fill[ezs](23,-167.6)--(25.3,-168.5)--(90.6,1.7)--(88.3,2.6)--cycle;
				\definecolor{ezs}{RGB}{253,203,40}\fill[ezs](24.9,-168.3)--(27.2,-169.2)--(92.5,1)--(90.3,1.8)--cycle;
				\definecolor{ezs}{RGB}{253,205,42}\fill[ezs](26.8,-169.1)--(29.1,-170)--(94.4,0.2)--(92.2,1.1)--cycle;
				\definecolor{ezs}{RGB}{253,207,45}\fill[ezs](28.7,-169.8)--(31,-170.7)--(96.4,-0.5)--(94.1,0.4)--cycle;
				\definecolor{ezs}{RGB}{253,210,49}\fill[ezs](30.7,-170.5)--(32.9,-171.4)--(98.3,-1.2)--(96,-0.4)--cycle;
				\definecolor{ezs}{RGB}{253,212,54}\fill[ezs](32.6,-171.3)--(34.9,-172.2)--(100.2,-2)--(97.9,-1.1)--cycle;
				\definecolor{ezs}{RGB}{253,214,59}\fill[ezs](34.5,-172)--(36.8,-172.9)--(102.1,-2.7)--(99.8,-1.8)--cycle;
				\definecolor{ezs}{RGB}{253,216,63}\fill[ezs](36.4,-172.7)--(38.7,-173.6)--(104,-3.4)--(101.7,-2.6)--cycle;
				\definecolor{ezs}{RGB}{253,218,65}\fill[ezs](38.3,-173.5)--(40.6,-174.4)--(105.9,-4.2)--(103.6,-3.3)--cycle;
				\definecolor{ezs}{RGB}{254,220,66}\fill[ezs](40.2,-174.2)--(42.5,-175.1)--(107.8,-4.9)--(105.5,-4)--cycle;
				\definecolor{ezs}{RGB}{254,222,68}\fill[ezs](42.1,-175)--(44.4,-175.8)--(109.7,-5.6)--(107.5,-4.8)--cycle;
				\definecolor{ezs}{RGB}{254,224,71}\fill[ezs](44,-175.7)--(46.3,-176.6)--(111.7,-6.4)--(109.4,-5.5)--cycle;
				\definecolor{ezs}{RGB}{254,226,75}\fill[ezs](46,-176.4)--(48.2,-177.3)--(113.6,-7.1)--(111.3,-6.2)--cycle;
				\definecolor{ezs}{RGB}{254,228,79}\fill[ezs](47.9,-177.2)--(50.1,-178)--(115.5,-7.8)--(113.2,-7)--cycle;
				\definecolor{ezs}{RGB}{254,230,82}\fill[ezs](49.8,-177.9)--(52.1,-178.8)--(117.4,-8.6)--(115.1,-7.7)--cycle;
				\definecolor{ezs}{RGB}{254,232,85}\fill[ezs](51.7,-178.6)--(54,-179.5)--(119.3,-9.3)--(117,-8.4)--cycle;
				\definecolor{ezs}{RGB}{254,234,88}\fill[ezs](53.6,-179.4)--(55.9,-180.2)--(121.2,-10)--(118.9,-9.2)--cycle;
				\definecolor{ezs}{RGB}{254,237,91}\fill[ezs](55.5,-180.1)--(57.8,-181)--(123.1,-10.8)--(120.8,-9.9)--cycle;
				\definecolor{ezs}{RGB}{254,239,94}\fill[ezs](57.4,-180.8)--(59.7,-181.7)--(125,-11.5)--(122.8,-10.6)--cycle;
				\definecolor{ezs}{RGB}{254,241,97}\fill[ezs](59.3,-181.6)--(61.6,-182.4)--(127,-12.2)--(124.7,-11.4)--cycle;
				\definecolor{ezs}{RGB}{254,243,101}\fill[ezs](61.2,-182.3)--(63.5,-183.2)--(128.9,-13)--(126.6,-12.1)--cycle;
				\definecolor{ezs}{RGB}{254,245,103}\fill[ezs](63.2,-183)--(65.4,-183.9)--(130.8,-13.7)--(128.5,-12.8)--cycle;
				\definecolor{ezs}{RGB}{254,245,105}\fill[ezs](65.1,-183.8)--(67.4,-184.6)--(132.7,-14.4)--(130.4,-13.6)--cycle;
				\definecolor{ezs}{RGB}{254,245,105}\fill[ezs](67,-184.5)--(69.3,-185.4)--(134.6,-15.2)--(132.3,-14.3)--cycle;
				\definecolor{ezs}{RGB}{254,244,103}\fill[ezs](68.9,-185.2)--(71.2,-186.1)--(136.5,-15.9)--(134.2,-15)--cycle;
				\definecolor{ezs}{RGB}{253,243,101}\fill[ezs](70.8,-186)--(73.1,-186.8)--(138.4,-16.6)--(136.1,-15.8)--cycle;
				\definecolor{ezs}{RGB}{253,241,98}\fill[ezs](72.7,-186.7)--(75,-187.6)--(140.3,-17.4)--(138,-16.5)--cycle;
				\definecolor{ezs}{RGB}{253,239,95}\fill[ezs](74.6,-187.4)--(76.9,-188.3)--(142.2,-18.1)--(140,-17.2)--cycle;
				\definecolor{ezs}{RGB}{253,236,92}\fill[ezs](76.5,-188.2)--(78.8,-189)--(144.2,-18.8)--(141.9,-18)--cycle;
				\definecolor{ezs}{RGB}{253,234,88}\fill[ezs](78.5,-188.9)--(80.7,-189.8)--(146.1,-19.6)--(143.8,-18.7)--cycle;
				\definecolor{ezs}{RGB}{253,232,85}\fill[ezs](80.4,-189.6)--(82.7,-190.5)--(148,-20.3)--(145.7,-19.4)--cycle;
				\definecolor{ezs}{RGB}{253,230,82}\fill[ezs](82.3,-190.4)--(84.6,-191.2)--(149.9,-21)--(147.6,-20.2)--cycle;
				\definecolor{ezs}{RGB}{253,228,78}\fill[ezs](84.2,-191.1)--(86.5,-192)--(151.8,-21.8)--(149.5,-20.9)--cycle;
				\definecolor{ezs}{RGB}{253,226,74}\fill[ezs](86.1,-191.8)--(88.4,-192.7)--(153.7,-22.5)--(151.4,-21.6)--cycle;
				\definecolor{ezs}{RGB}{253,224,70}\fill[ezs](88,-192.6)--(90.3,-193.4)--(155.6,-23.3)--(153.3,-22.4)--cycle;
				\definecolor{ezs}{RGB}{253,222,67}\fill[ezs](89.9,-193.3)--(92.2,-194.2)--(157.5,-24)--(155.3,-23.1)--cycle;
				\definecolor{ezs}{RGB}{253,220,65}\fill[ezs](91.8,-194)--(94.1,-194.9)--(159.5,-24.7)--(157.2,-23.8)--cycle;
				\definecolor{ezs}{RGB}{253,217,63}\fill[ezs](93.7,-194.8)--(96,-195.6)--(161.4,-25.5)--(159.1,-24.6)--cycle;
				\definecolor{ezs}{RGB}{253,215,60}\fill[ezs](95.7,-195.5)--(97.9,-196.4)--(163.3,-26.2)--(161,-25.3)--cycle;
				\definecolor{ezs}{RGB}{253,213,56}\fill[ezs](97.6,-196.2)--(99.9,-197.1)--(165.2,-26.9)--(162.9,-26)--cycle;
				\definecolor{ezs}{RGB}{252,211,53}\fill[ezs](99.5,-197)--(101.8,-197.8)--(167.1,-27.7)--(164.8,-26.8)--cycle;
				\definecolor{ezs}{RGB}{252,209,49}\fill[ezs](101.4,-197.7)--(103.7,-198.6)--(169,-28.4)--(166.7,-27.5)--cycle;
				\definecolor{ezs}{RGB}{252,207,45}\fill[ezs](103.3,-198.4)--(105.6,-199.3)--(170.9,-29.1)--(168.6,-28.2)--cycle;
				\definecolor{ezs}{RGB}{253,205,42}\fill[ezs](105.2,-199.2)--(107.5,-200)--(172.8,-29.9)--(170.6,-29)--cycle;
				\definecolor{ezs}{RGB}{253,202,40}\fill[ezs](107.1,-199.9)--(109.4,-200.8)--(174.7,-30.6)--(172.5,-29.7)--cycle;
				\definecolor{ezs}{RGB}{253,200,38}\fill[ezs](109,-200.6)--(111.3,-201.5)--(176.7,-31.3)--(174.4,-30.4)--cycle;
				\definecolor{ezs}{RGB}{253,198,35}\fill[ezs](111,-201.4)--(113.2,-202.2)--(178.6,-32.1)--(176.3,-31.2)--cycle;
				\definecolor{ezs}{RGB}{253,196,31}\fill[ezs](112.9,-202.1)--(115.2,-203)--(180.5,-32.8)--(178.2,-31.9)--cycle;
				\definecolor{ezs}{RGB}{253,194,28}\fill[ezs](114.8,-202.8)--(117.1,-203.7)--(182.4,-33.5)--(180.1,-32.6)--cycle;
				\definecolor{ezs}{RGB}{253,192,26}\fill[ezs](116.7,-203.6)--(119,-204.5)--(184.3,-34.3)--(182,-33.4)--cycle;
				\definecolor{ezs}{RGB}{253,189,24}\fill[ezs](118.6,-204.3)--(120.9,-205.2)--(186.2,-35)--(183.9,-34.1)--cycle;
				\definecolor{ezs}{RGB}{253,186,20}\fill[ezs](120.5,-205)--(122.8,-205.9)--(188.1,-35.7)--(185.8,-34.8)--cycle;
				\definecolor{ezs}{RGB}{253,184,17}\fill[ezs](122.4,-205.8)--(124.7,-206.7)--(190,-36.5)--(187.8,-35.6)--cycle;
				\definecolor{ezs}{RGB}{252,183,15}\fill[ezs](124.3,-206.5)--(126.6,-207.4)--(192,-37.2)--(189.7,-36.3)--cycle;
				\definecolor{ezs}{RGB}{252,181,14}\fill[ezs](126.3,-207.2)--(128.5,-208.1)--(193.9,-37.9)--(191.6,-37.1)--cycle;
				\definecolor{ezs}{RGB}{252,179,13}\fill[ezs](128.2,-208)--(130.4,-208.9)--(195.8,-38.7)--(193.5,-37.8)--cycle;
				\definecolor{ezs}{RGB}{252,178,12}\fill[ezs](130.1,-208.7)--(132.4,-209.6)--(197.7,-39.4)--(195.4,-38.5)--cycle;
				\definecolor{ezs}{RGB}{251,178,10}\fill[ezs](132,-209.4)--(134.3,-210.3)--(199.6,-40.1)--(197.3,-39.3)--cycle;
				\definecolor{ezs}{RGB}{251,178,10}\fill[ezs](133.9,-210.2)--(136.2,-211.1)--(201.5,-40.9)--(199.2,-40)--cycle;
				\definecolor{ezs}{RGB}{250,177,11}\fill[ezs](135.8,-210.9)--(138.1,-211.8)--(203.4,-41.6)--(201.1,-40.7)--cycle;
				\definecolor{ezs}{RGB}{250,178,11}\fill[ezs](137.7,-211.6)--(140,-212.5)--(205.3,-42.3)--(203.1,-41.5)--cycle;
				\definecolor{ezs}{RGB}{248,180,10}\fill[ezs](139.6,-212.4)--(141.9,-213.3)--(207.3,-43.1)--(205,-42.2)--cycle;
				\definecolor{ezs}{RGB}{247,181,9}\fill[ezs](141.5,-213.1)--(143.8,-214)--(209.2,-43.8)--(206.9,-42.9)--cycle;
				\definecolor{ezs}{RGB}{246,182,9}\fill[ezs](143.5,-213.8)--(145.7,-214.7)--(211.1,-44.5)--(208.8,-43.7)--cycle;
				\definecolor{ezs}{RGB}{246,182,10}\fill[ezs](145.4,-214.6)--(147.7,-215.5)--(213,-45.3)--(210.7,-44.4)--cycle;
			\end{scope}
			% chữ EDUZY (Times đậm, làm dày nét bằng 5 lớp lệch nhẹ; căn đúng khung chữ logo gốc)
			\foreach \dx/\dy in {-0.7/0,0.7/0,0/-0.3,0/0.3,0/0}
			\node[anchor=center,inner sep=0pt,outer sep=0pt,text=ezText,align=center] at (109.6+\dx,-193.6+\dy)
			{\resizebox{1.01cm}{0.22cm}{\fontencoding{T1}\fontfamily{ptm}\fontseries{b}\selectfont EDUZY}};
		\end{tikzpicture}%
	}%
}
%====================== THIẾT LẬP THÂN ĐỀ ======================
%--- Môi trường tự luận {bt}: ex_test 3.3.4 chưa khai báo sẵn nên tạo ở đây
%    (ex_test đã tự móc \fixshowans và \listenumerate cho bt)
\makeatletter
\@ifundefined{bt}{\newtheorem{bt}{\color{gold1}\namebt}}{}
\makeatother
\renewcommand{\sectionmark}[1]{\markright{\thesection\ #1}}
%--- Khung tiêu đề (logo + tối đa 4 dòng; dòng 3, 4 để trống {} thì bỏ qua)
%  Đề thi:    \khungtieude{THI THỬ THPT QUỐC GIA NĂM 2026}{Môn Toán -- Đề số 21}
%                         {(Đề thi gồm: ...)}{Thời gian làm bài: 90 phút (không kể thời gian phát đề).}
%  Chuyên đề: \khungtieude{CHUYÊN ĐỀ 3 -- DÃY SỐ VÀ CÔNG THỨC TRUY HỒI}
%                         {Bồi dưỡng học sinh giỏi Toán THPT}{...}{}
\newcommand{\khungtieude}[4]{%
	\begin{tcolorbox}[colframe=gold1,colback=gold5!45,boxrule=0.5mm]
		\begin{minipage}[c]{3cm}
			\centerline{\eduzylogo[2.4cm]}
		\end{minipage}\hspace{0.5cm}
		\begin{minipage}[c]{12.5cm}
			\vskip 0.15 cm
			\centerline{\ind{#1}}
			\vskip 0.15 cm
			\centerline{\textbf{#2}}
			\if\relax\detokenize{#3}\relax\else
				\vskip 0.15 cm
				\centerline{\textit{#3}}
			\fi
			\if\relax\detokenize{#4}\relax\else
				\vskip 0.15 cm
				\centerline{\textit{#4}}
			\fi
			\vskip 0.15 cm
		\end{minipage}
	\end{tcolorbox}
}
%====================== BỔ SUNG CHO BÀI HỌC, CHUYÊN ĐỀ ======================
\usepackage[shortlabels]{enumitem}
\usepackage{tabularx}
%--- Danh sách gọn: \begin{enumerate}[gon] ... \end{enumerate}
\SetEnumitemKey{gon}{leftmargin=*,itemsep=1pt,parsep=0pt,topsep=3pt}
%--- Khung lý thuyết: \begin{khung}{} ... \end{khung}; tiêu đề chỉ dùng khi thật ngắn ("Mẹo", "Lỗi thường gặp")
\newtcolorbox{khung}[1]{breakable,enhanced,colframe=gold1,colback=gold5!45,
	boxrule=0.4mm,fonttitle=\bfseries\sffamily,coltitle=white,colbacktitle=gold1,
	attach boxed title to top left={xshift=6mm,yshift=-3mm},
	boxed title style={arc=1mm,boxrule=0.3mm},title={#1}}
%--- Tiêu đề mục nhỏ trong một dạng: \muc{Dấu hiệu}, \muc{Phương pháp}, \muc{Bài tự luyện}
\makeatletter
\newcommand{\muc}[1]{\par\medskip\noindent\indam{#1}\par\nobreak\smallskip\@afterindentfalse\@afterheading}
\makeatother
%--- Ô công thức cần nhớ (căn giữa, viền vàng, vừa khít nội dung)
%     \begin{center}\congthuc{$Q(n)=(n-1)!$}\end{center}
\newtcbox{\congthuc}{on line,colframe=gold1,colback=gold5!45,boxrule=0.5pt,arc=2pt,
	left=8pt,right=8pt,top=4pt,bottom=4pt}
%--- Công tắc lời giải bài tự luyện (ví dụ mẫu luôn dùng \loigiai)
%     \loigiaiTL[{đáp số}]{lời giải đầy đủ}
%     \TLfull -> in lời giải đầy đủ;  \TLdapso -> chỉ in một dòng "Đáp số";  \TLan -> ẩn hoàn toàn
%     (\CheDo là dethi thì mọi lời giải và đáp số đều ẩn)
\makeatletter
\newif\ifTLhien   \TLhientrue
\newif\ifTLdapso  \TLdapsofalse
\newcommand{\TLfull}{\TLhientrue\TLdapsofalse}
\newcommand{\TLdapso}{\TLhientrue\TLdapsotrue}
\newcommand{\TLan}{\TLhienfalse}
\newcommand{\loigiaiTL}[2][]{%
	\ifTLhien
		\ifTLdapso
			\def\@tlarg{#1}%
			\ifx\@tlarg\@empty\else
				\ifdefstring{\kindOPTN}{dethi}{}{\par\noindent{\bfseries\color{gold1}Đáp số:} #1\par}%
			\fi
		\else
			\loigiai{#2}%
		\fi
	\fi}
\makeatother
\TLdapso% mặc định: chỉ in đáp số bài tự luyện
%====================== BẢNG ĐÚNG/SAI (\choiceTFt) ======================
%  Hàng tiêu đề nền nâu đồng chữ trắng (Phát biểu | Đ | S), các hàng tô xen kẽ kem / trắng,
%  đường kẻ vàng kim, nhãn a) b) c) d) màu nâu đồng. Bản có đáp án: cột được chọn có vòng tròn
%  nâu đồng với dấu tích trắng (câu đúng tích ở cột Đúng, câu sai tích ở cột Sai).
%  Tiêu đề ghi đủ "Đúng | Sai": \OPTN{viettat=0} sau \input (cột Đúng/Sai rộng ra, phát biểu hẹp lại).
%  Dùng bảng gốc của ex_test: \def\BangDSgoc{1} trước %============================================================
%  PREAMBLE EDUZY v2 - dùng chung cho đề thi, bài học, chuyên đề (gói ex_test)
%  Khác bản preamble-eduzy.tex ở đầu trang, chân trang: dải nền nâu đồng tràn ngang,
%  logo + Eduzy.com.vn, ô chữ vàng kem cạnh xiên, chân trang "Trang (số)";
%  logo Eduzy chìm giữa mọi trang (tắt: \logochimtat, bật lại: \logochimbat).
%  Cách dùng ở file chính (xem mau-chuan.tex, mau-chuyen-de.tex):
%      \def\CheDo{loigiai}          % loigiai | dethi | solcolor | color | book
%      \def\tenchuong{<dòng chữ trên đầu trang>}
%      \def\thuonghieu{<chữ ở chân trang lẻ>}   % không khai báo thì dùng mặc định
%      \def\tacgia{<chữ ở chân trang chẵn>}      % không khai báo thì dùng mặc định
%      \def\khauhieu{<chữ cạnh logo ở đầu trang>}  % tùy chọn, mặc định Eduzy.com.vn
%      \input{preamble-eduzy-v2}
%      \TLdapso                      % tài liệu: \TLfull | \TLdapso | \TLan
%      \begin{document} ... \end{document}
%  Logo Eduzy vẽ bằng TikZ ngay trong file (\eduzylogo), không cần file ảnh.
%  Biên dịch: pdfLaTeX (máy có gói vietnam) hoặc XeLaTeX.
%============================================================
\documentclass[12pt,a4paper,twoside]{article}%<-- twoside: bat dau trang/chan trang chan-le
\usepackage{iftex}
\usepackage{amsmath,amssymb,tikz,mathrsfs,tkz-tab}
\usepackage{graphicx}
\usepackage{pgfplots}% đồ thị hàm số
\pgfplotsset{compat=1.17}
%\usepackage{mathptmx}
\usepackage{mathpazo}
\ifXeTeX% XeLaTeX: font chữ có đủ dấu tiếng Việt, giữ nguyên font toán mathpazo
	\usepackage[no-math]{fontspec}
	\setmainfont{TeX Gyre Pagella}
	\newfontfamily\EZqag{TeX Gyre Adventor}[NFSSFamily=qag,Ligatures=TeX]
\fi
\usepackage[many]{tcolorbox}
\usetikzlibrary{math,through,calc,intersections,angles,quotes,shapes,shapes.geometric,arrows,patterns,matrix,chains,arrows.meta,decorations.shapes,decorations.fractals,decorations.markings,shadows,decorations.pathreplacing}
\usepackage[top=1.6cm, bottom=1.8cm, left=1.8cm, right=1.8cm]{geometry}% v2: top 1.6cm (bản cũ 1.5cm) để chữ không sát dải đầu trang
\providecommand{\CheDo}{loigiai}% loigiai | dethi | solcolor | color | book
\expandafter\usepackage\expandafter[\CheDo]{ex_test}
% \tenchuong: khai báo ở file chính (trước \input) hoặc \renewcommand sau \input
%============================================================
%  [GỘP TỪ eduzy-de.tex] EDUZY - MÀU NÂU ĐỒNG / VÀNG KIM
%  + ĐẦU TRANG, CHÂN TRANG KIỂU DẢI NGANG (v2), ĐỐI XỨNG TRANG CHẴN / TRANG LẺ.
%  (không cần file eduzy-de.tex riêng nữa)
%  Khối này phải đặt SAU \usepackage[...]{ex_test} và SAU \def\tenchuong{...}
%============================================================
\usepackage{lastpage}
\usepackage{needspace}
%--- DÀN TRANG ĐỀ THI: tách câu (THÊM 05/11/2026; chỉ thêm lệnh mới, KHÔNG đổi hành vi mặc định) ---
%  Hai công tắc, mặc định tắt (= preamble chạy y như cũ). Bật bằng cách đặt TRƯỚC \input{preamble-eduzy-v2}
%  trong file đề (tools/dan_trang/dan_trang.py tự ghi khi đặt \NeedspaceCau):
%    \def\TachCauDS{1}   câu Đúng/Sai được tách: câu dẫn (kèm hình) ở cuối trang, bảng Phát biểu sang trang sau
%                        (câu dẫn và hình không tách nhau, bảng không tách giữa hai trang)
%    \def\TachCauLG{1}   bản lời giải (\CheDo khác dethi): phần đề của câu cao hơn 1/3 trang được phép tách
%  Lệnh dùng trong file đề (giá trị do tools/dan_trang/dan_trang.py đo ở bản dethi):
%    \NeedspaceCau[<cao câu dẫn + hình>]{<cao cả khối câu>}   ngay trên \begin{ex} (câu đầu phần: trên tiêu đề phần)
%        công tắc tắt: đúng bằng \Needspace{<cao cả khối câu>} như cũ;
%        \TachCauDS=1 và có [..] (câu Đúng/Sai được chọn tách): chỉ giữ câu dẫn + hình;
%        \TachCauLG=1 ở bản lời giải: đòi tối đa 1/3 \textheight.
%    \NeedspaceBang{<cao bảng Phát biểu>}%   ngay trước \choiceTFt: \TachCauDS=1 thì giữ trọn bảng trên một trang;
%        công tắc tắt: không làm gì.
\providecommand{\TachCauDS}{0}
\providecommand{\TachCauLG}{0}
\newlength{\EZnsCau}
\newcommand{\NeedspaceCau}[2][]{%
	\setlength{\EZnsCau}{#2}%
	\ifnum\TachCauDS=1 \if\relax\detokenize{#1}\relax\else\setlength{\EZnsCau}{#1}\fi\fi
	\ifnum\TachCauLG=1 \ifdefstring{\kindOPTN}{dethi}{}{\ifdim\EZnsCau>0.3333\textheight\setlength{\EZnsCau}{0.3333\textheight}\fi}\fi
	\Needspace{\EZnsCau}}
\newcommand{\NeedspaceBang}[1]{\ifnum\TachCauDS=1 \Needspace{#1}\fi}

%--- Thang màu: đậm -> nhạt
\definecolor{gold1}{RGB}{126,96,0}    % nâu đồng đậm  #7E6000
\definecolor{gold2}{RGB}{192,144,3}   % vàng kim      #C09003
\definecolor{gold3}{RGB}{255,215,99}  % vàng nhạt     #FFD763
\definecolor{gold4}{RGB}{255,231,153} % vàng rất nhạt #FFE799
\definecolor{gold5}{RGB}{255,243,205} % vàng kem      #FFF3CD
%--- Ghi đè tông xanh/đỏ cũ
\definecolor{myblue}{RGB}{126,96,0}
\colorlet{blue}{gold1}
\colorlet{cyan}{gold2}
\colorlet{teal}{gold2}

%--- Tên chương mặc định (file đề nên tự khai báo lại)
\providecommand{\tenchuong}{}
%--- Chân trang
\providecommand{\thuonghieu}{TOÁN 12 -- Eduzy.com.vn}
\providecommand{\tacgia}{Tô Gia Bảo -- 0353703783}

%--- Chữ cạnh logo ở đầu trang
\providecommand{\khauhieu}{Eduzy.com.vn}

%====================== THÔNG SỐ ĐẦU TRANG, CHÂN TRANG (v2) ======================
%  Dải nền đậm tràn ngang trang; ô chữ vàng kem có cạnh xiên, kèm 2 vạch xiên chuyển màu.
%  Đầu trang lẻ : logo + \khauhieu (trái)   | ô chữ \tenchuong (phải)
%  Chân trang lẻ: ô chữ \thuonghieu (trái)  | Trang (số) (phải)
%  Trang chẵn soi gương trang lẻ; chân trang chẵn ghi \tacgia.
%--- Màu: chỉ dùng thang vàng kim Eduzy
\colorlet{ezNen}{gold1}% dải nền đậm
\colorlet{ezVach1}{gold3}% vạch xiên sát dải nền (đậm hơn)
\colorlet{ezVach2}{gold4}% vạch xiên sát ô chữ (nhạt hơn)
\colorlet{ezO}{gold5}% ô chữ vàng kem
%--- Kích thước
\newlength{\EZdai}\setlength{\EZdai}{1.1cm}% chiều cao dải nền
\newlength{\EZo}\setlength{\EZo}{0.68cm}% chiều cao ô chữ
\newlength{\EZxien}\setlength{\EZxien}{0.55cm}% độ lệch ngang của cạnh xiên
\newlength{\EZvach}\setlength{\EZvach}{0.75cm}% bề rộng mỗi vạch xiên
\newlength{\EZkhe}\setlength{\EZkhe}{0.12cm}% khe trắng giữa vạch và ô chữ
\newlength{\EZdem}\setlength{\EZdem}{0.4cm}% khoảng đệm từ chữ tới cạnh xiên
\newlength{\EZle}\setlength{\EZle}{1.8cm}% lề trái/phải của chữ (bằng lề thân bài)
\newlength{\EZrongmin}\setlength{\EZrongmin}{4.5cm}% ô chữ rộng tối thiểu
\newlength{\EZrong}% bề rộng chữ trong ô (tính tự động)
%--- Font
\newcommand{\EZfontO}{\fontsize{10.5pt}{12pt}\selectfont\bfseries}% chữ trong ô
\newcommand{\EZfontKH}{\fontfamily{qag}\fontsize{11pt}{12pt}\selectfont\bfseries}% \khauhieu cạnh logo
\newcommand{\EZfontTrang}{\fontsize{11.5pt}{12pt}\selectfont\bfseries}% chữ "Trang" và số trang
%--- Đo bề rộng chữ trong ô
\newcommand{\EZdorong}[1]{%
	\settowidth{\EZrong}{\EZfontO #1}%
	\ifdim\EZrong<\EZrongmin\setlength{\EZrong}{\EZrongmin}\fi}
%--- Logo + chữ trắng \khauhieu (dựng sẵn một lần khi vào \begin{document})
\newsavebox{\EZlogobox}
\AtBeginDocument{\sbox{\EZlogobox}{%
		\raisebox{-0.5\height}{\eduzylogo[1cm]}\hspace{0.15cm}%
		\raisebox{-0.5\height}{\color{white}\EZfontKH\khauhieu}}}
%--- Hai vạch xiên + khe trắng. #1: điểm gốc ở mép dưới, #2: hướng (1 sang phải, -1 sang trái),
%    #3: hướng nghiêng (1 kiểu "/", -1 kiểu "\")
\newcommand{\EZvachxien}[3]{%
	\fill[white] (#1) --++({(#3)*\EZxien},\EZo) --++({(#2)*\EZkhe},0) --++({-(#3)*\EZxien},-\EZo) --cycle;
	\fill[ezVach2] ([xshift=(#2)*\EZkhe]#1) --++({(#3)*\EZxien},\EZo) --++({(#2)*\EZvach},0) --++({-(#3)*\EZxien},-\EZo) --cycle;
	\fill[ezVach1] ([xshift=(#2)*(\EZkhe+\EZvach)]#1) --++({(#3)*\EZxien},\EZo) --++({(#2)*\EZvach},0) --++({-(#3)*\EZxien},-\EZo) --cycle;
	\draw[white,line width=0.6pt] ([xshift=(#2)*(\EZkhe+\EZvach)]#1) --++({(#3)*\EZxien},\EZo);
}
%--- Chữ trong ô (căn giữa theo chiều cao ô, không lệch vì dấu tiếng Việt)
\tikzset{ezchu/.style={inner sep=0pt,text=ezNen,text height=0.72em,text depth=0.22em}}

%====================== LOGO CHÌM GIỮA TRANG ======================
%  Vẽ cùng đầu trang nên nằm dưới chữ của thân bài, có trên mọi trang.
%  Tắt: \logochimtat (từ trang đó trở đi), bật lại: \logochimbat.
%  Đổi cỡ, độ đậm: khai báo trước \input, ví dụ \def\EZchimcao{9cm} \def\EZchimdo{0.1}
\providecommand{\EZchimcao}{10cm}% chiều cao logo chìm
\providecommand{\EZchimdo}{0.12}% độ đậm (0 = trong suốt, 1 = rõ hẳn)
\newif\ifEZchim \EZchimtrue
\newcommand{\logochimtat}{\EZchimfalse}
\newcommand{\logochimbat}{\EZchimtrue}
%  Dựng sẵn logo mờ một lần khi vào \begin{document}; transparency group làm mờ cả logo
%  một lượt nên không lộ vệt ở chỗ các dải màu chồng nhau.
\newsavebox{\EZchimtho}
\newsavebox{\EZchimbox}
\AtBeginDocument{%
	\sbox{\EZchimtho}{\eduzylogo[\EZchimcao]}%
	\sbox{\EZchimbox}{\begin{tikzpicture}
			\begin{scope}[transparency group,opacity=\EZchimdo]
				\node[inner sep=0pt] {\usebox{\EZchimtho}};
			\end{scope}
	\end{tikzpicture}}}
\newcommand{\EZvechim}{%
	\ifEZchim
		\node[inner sep=0pt] at (current page.center) {\usebox{\EZchimbox}};
	\fi}

%====================== ĐẦU TRANG - TRANG LẺ ======================
%  (logo bên trái, ô chữ \tenchuong bên phải)
\newcommand{\myfancyheadOdd}{%
	\EZdorong{\tenchuong}%
	\begin{tikzpicture}[remember picture,overlay]
		\coordinate (T) at (current page.north west);
		\coordinate (P) at (current page.north east);
		%--- logo chìm giữa trang
		\EZvechim
		%--- dải nền
		\fill[ezNen] (T) rectangle ([yshift=-\EZdai]P);
		%--- ô chữ: X là góc dưới bên trái, cạnh trái xiên kiểu "/"
		\coordinate (X) at ([xshift=-\EZle-\EZrong-\EZdem-\EZxien,yshift=-\EZdai]P);
		\fill[ezO] (X) --++(\EZxien,\EZo) -- ([yshift=-\EZdai+\EZo]P) -- ([yshift=-\EZdai]P) --cycle;
		\EZvachxien{X}{-1}{1}
		\node[ezchu,anchor=east] at ([xshift=-\EZle,yshift=-\EZdai+0.5*\EZo]P) {\EZfontO\tenchuong};
		%--- logo + tên
		\node[inner sep=0pt,anchor=west] at ([xshift=\EZle,yshift=-0.5*\EZdai]T) {\usebox{\EZlogobox}};
	\end{tikzpicture}%
}

%====================== ĐẦU TRANG - TRANG CHẴN =====================
%  (soi gương: ô chữ \tenchuong bên trái, logo bên phải)
\newcommand{\myfancyheadEven}{%
	\EZdorong{\tenchuong}%
	\begin{tikzpicture}[remember picture,overlay]
		\coordinate (T) at (current page.north west);
		\coordinate (P) at (current page.north east);
		%--- logo chìm giữa trang
		\EZvechim
		%--- dải nền
		\fill[ezNen] (T) rectangle ([yshift=-\EZdai]P);
		%--- ô chữ: X là góc dưới bên phải, cạnh phải xiên kiểu "\"
		\coordinate (X) at ([xshift=\EZle+\EZrong+\EZdem+\EZxien,yshift=-\EZdai]T);
		\fill[ezO] (X) --++(-\EZxien,\EZo) -- ([yshift=-\EZdai+\EZo]T) -- ([yshift=-\EZdai]T) --cycle;
		\EZvachxien{X}{1}{-1}
		\node[ezchu,anchor=west] at ([xshift=\EZle,yshift=-\EZdai+0.5*\EZo]T) {\EZfontO\tenchuong};
		%--- logo + tên
		\node[inner sep=0pt,anchor=east] at ([xshift=-\EZle,yshift=-0.5*\EZdai]P) {\usebox{\EZlogobox}};
	\end{tikzpicture}%
}

%====================== CHÂN TRANG - TRANG LẺ ======================
%  (ô chữ \thuonghieu bên trái, "Trang (số)" bên phải)
\newcommand{\myfancyfootOdd}{%
	\EZdorong{\thuonghieu}%
	\begin{tikzpicture}[remember picture,overlay]
		\coordinate (T) at (current page.south west);
		\coordinate (P) at (current page.south east);
		%--- dải nền
		\fill[ezNen] (T) rectangle ([yshift=\EZdai]P);
		%--- ô chữ sát mép trên dải: X là góc dưới bên phải, cạnh phải xiên kiểu "/"
		\coordinate (X) at ([xshift=\EZle+\EZrong+\EZdem,yshift=\EZdai-\EZo]T);
		\fill[ezO] ([yshift=\EZdai-\EZo]T) -- (X) --++(\EZxien,\EZo) -- ([yshift=\EZdai]T) --cycle;
		\EZvachxien{X}{1}{1}
		\node[ezchu,anchor=west] at ([xshift=\EZle,yshift=\EZdai-0.5*\EZo]T) {\EZfontO\thuonghieu};
		%--- số trang
		\node[circle,draw=white,line width=0.9pt,text=white,inner sep=0pt,minimum size=0.62cm,anchor=east]
		(so) at ([xshift=-\EZle,yshift=0.5*\EZdai]P) {\EZfontTrang\thepage};
		\node[inner sep=0pt,text=white,anchor=base east] (tr) at ([xshift=-0.15cm]so.base west) {\EZfontTrang Trang};
	\end{tikzpicture}%
}

%====================== CHÂN TRANG - TRANG CHẴN ====================
%  (soi gương: "Trang (số)" bên trái, ô chữ \tacgia bên phải)
\newcommand{\myfancyfootEven}{%
	\EZdorong{\tacgia}%
	\begin{tikzpicture}[remember picture,overlay]
		\coordinate (T) at (current page.south west);
		\coordinate (P) at (current page.south east);
		%--- dải nền
		\fill[ezNen] (T) rectangle ([yshift=\EZdai]P);
		%--- ô chữ sát mép trên dải: X là góc dưới bên trái, cạnh trái xiên kiểu "\"
		\coordinate (X) at ([xshift=-\EZle-\EZrong-\EZdem,yshift=\EZdai-\EZo]P);
		\fill[ezO] ([yshift=\EZdai-\EZo]P) -- (X) --++(-\EZxien,\EZo) -- ([yshift=\EZdai]P) --cycle;
		\EZvachxien{X}{-1}{-1}
		\node[ezchu,anchor=east] at ([xshift=-\EZle,yshift=\EZdai-0.5*\EZo]P) {\EZfontO\tacgia};
		%--- số trang
		%    (vòng tròn ẩn so0 nằm giữa dải để lấy đường chân chữ, giống trang lẻ)
		\node[circle,inner sep=0pt,minimum size=0.62cm] (so0) at ([yshift=0.5*\EZdai]T) {\EZfontTrang\phantom{\thepage}};
		\node[inner sep=0pt,text=white,anchor=base west] (tr) at ($(so0.base -| T)+(\EZle,0)$) {\EZfontTrang Trang};
		\node[circle,draw=white,line width=0.9pt,text=white,inner sep=0pt,minimum size=0.62cm,anchor=base west]
		(so) at ([xshift=0.15cm]tr.base east) {\EZfontTrang\thepage};
	\end{tikzpicture}%
}

%--- Giữ tương thích với các file cũ đang gọi \myfancyhead, \myfancyfoot
\newcommand{\myfancyhead}{\myfancyheadOdd}
\newcommand{\myfancyfoot}{\myfancyfootOdd}

%======================= GẮN VÀO TRANG =============================
\pagestyle{fancy}
\fancyhf{}
\fancyhead[LO]{\myfancyheadOdd}
\fancyhead[RE]{\myfancyheadEven}
\fancyfoot[LO]{\myfancyfootOdd}
\fancyfoot[RE]{\myfancyfootEven}
\renewcommand{\headrulewidth}{0pt}
\renewcommand{\footrulewidth}{0pt}
\fancypagestyle{plain}{%
	\fancyhf{}
	\fancyhead[LO]{\myfancyheadOdd}
	\fancyhead[RE]{\myfancyheadEven}
	\fancyfoot[LO]{\myfancyfootOdd}
	\fancyfoot[RE]{\myfancyfootEven}
	\renewcommand{\headrulewidth}{0pt}
	\renewcommand{\footrulewidth}{0pt}
}
%=================== HẾT PHẦN GỘP TỪ eduzy-de.tex ===================
\newcommand{\indam}[1]{
	{\fontsize{12pt}{9pt}\sffamily \selectfont\bfseries\color{gold1} #1}
}
\newcommand{\ind}[1]{
	{\fontsize{14pt}{9pt}\fontfamily{qag} \selectfont\bfseries\color{gold1} #1}
}
\renewcommand{\FalseEX}{\stepcounter{dapan}\textbf{\selectfont\color{gold1}\Alph{dapan}.}\ignorespaces}
\renewtheorem{ex}{\color{gold1}Câu}
\renewcommand{\baselinestretch}{1.1}
\newcommand{\hoac}[1]{
	\left[\begin{aligned}#1\end{aligned}\right.}
\newcommand{\heva}[1]{
	\left\{\begin{aligned}#1\end{aligned}\right.}
\usepackage{esvect}
\def\vec{\vv}
\def\vt{\vv}
\def\overrightarrow{\vv}
%====================== LOGO EDUZY VẼ BẰNG TIKZ ======================
%  Đường viền dò từ logo gốc; dải màu vàng kim đo trực tiếp trên ảnh gốc
%  (nghiêng 21 độ, tô bằng 80 dải hẹp nên hiển thị giống nhau trên mọi trình xem PDF).
%  Dùng: \eduzylogo          -> cao 2.4cm (mặc định)
%        \eduzylogo[1.5cm]   -> cao 1.5cm
\definecolor{ezText}{RGB}{245,200,46}
\newcommand{\eduzylogo}[1][2.4cm]{%
	\resizebox{!}{#1}{%
		\begin{tikzpicture}[x=0.01cm,y=0.01cm]
			\path[use as bounding box] (0,0) rectangle (224,-224);
			% mũ tốt nghiệp + vương miện
			\begin{scope}
				\clip
				(119.4,-176.1)..controls(118.9,-176.1)and(116.8,-176)..(114.6,-176)
				..controls(101.4,-176)and(85.8,-175.9)..(83.7,-175.7)
				..controls(81.8,-175.6)and(81,-175.6)..(80.2,-175.7)
				..controls(78.8,-175.9)and(77.7,-175.9)..(76.2,-175.7)
				..controls(74.9,-175.5)and(70.8,-175.5)..(64.3,-175.7)
				..controls(59.1,-175.9)and(58.3,-175.7)..(58.6,-173.7)
				..controls(58.6,-172.8)and(59.6,-169.6)..(60.3,-168)
				..controls(60.8,-166.9)and(61.4,-165.8)..(62.3,-165)
				..controls(63.1,-164.2)and(64.3,-163.4)..(65,-163.2)
				..controls(65.5,-163.1)and(72.3,-163)..(76.6,-163.1)
				..controls(77.9,-163.1)and(80.1,-163.1)..(81.4,-163)
				..controls(82.8,-162.9)and(86.3,-162.9)..(89.2,-162.9)
				..controls(93,-162.9)and(94.9,-162.8)..(96,-162.7)
				..controls(97.2,-162.5)and(101,-162.5)..(114,-162.4)
				..controls(123,-162.3)and(134.9,-162.2)..(140.2,-162.1)
				..controls(150.1,-162)and(152.2,-162.1)..(153.6,-162.4)
				..controls(154.4,-162.6)and(157.1,-163.8)..(158.1,-164.4)
				..controls(158.8,-164.8)and(159.7,-165.7)..(159.9,-166.3)
				..controls(160,-166.5)and(160.2,-167.1)..(160.2,-167.6)
				..controls(160.3,-168)and(160.6,-169.4)..(160.9,-170.6)
				..controls(161.2,-171.7)and(161.5,-173.1)..(161.6,-173.6)
				..controls(161.8,-174.4)and(161.8,-174.5)..(161.5,-174.9)
				..controls(161.2,-175.6)and(160.8,-175.9)..(159.9,-176)
				..controls(159.2,-176.1)and(154.1,-176.2)..(151.5,-176.2)
				..controls(148.1,-176.1)and(122.6,-176.1)..(121.6,-176.2)
				..controls(120.8,-176.2)and(119.8,-176.2)..(119.4,-176.1)--cycle
				(130.4,-160)..controls(129,-160)and(119.8,-160)..(110.1,-159.9)
				..controls(100.3,-159.8)and(88.5,-159.7)..(83.9,-159.6)
				..controls(77.8,-159.5)and(75.3,-159.5)..(74.5,-159.6)
				..controls(73.9,-159.7)and(72.5,-159.8)..(71.3,-159.8)
				..controls(69.6,-159.8)and(69,-159.8)..(68.5,-159.6)
				..controls(67.6,-159.3)and(66.8,-158.6)..(66.2,-157.6)
				..controls(66,-157.3)and(65.6,-156.5)..(65.3,-156.1)
				..controls(64.6,-155.1)and(64.5,-154.9)..(63.7,-152.9)
				..controls(63.2,-151.6)and(62.9,-150.9)..(61.8,-149.6)
				..controls(61.5,-149.2)and(61.1,-148.5)..(60.9,-147.9)
				..controls(60.7,-147.4)and(60.3,-146.6)..(60.1,-146.1)
				..controls(59.9,-145.5)and(59.5,-144.9)..(59,-144.3)
				..controls(58.5,-143.7)and(58.2,-143.2)..(57.9,-142.5)
				..controls(57.7,-141.9)and(57.3,-141.1)..(57,-140.7)
				..controls(56.6,-140.3)and(56.1,-139.3)..(55.8,-138.5)
				..controls(55.5,-137.8)and(55.1,-136.8)..(54.9,-136.5)
				..controls(54.7,-136.1)and(54.3,-135.2)..(53.9,-134.5)
				..controls(53.6,-133.8)and(53,-132.7)..(52.6,-132)
				..controls(51.6,-130.5)and(51.4,-130.2)..(51.1,-129)
				..controls(50.6,-127.5)and(50.6,-127.4)..(49.7,-126.3)
				..controls(49.2,-125.7)and(48.6,-124.8)..(48.3,-124.2)
				..controls(48,-123.6)and(47.5,-122.7)..(47.3,-122.2)
				..controls(47,-121.8)and(46.5,-120.7)..(46.1,-119.8)
				..controls(45.8,-118.9)and(45.2,-117.8)..(44.9,-117.2)
				..controls(44.2,-116.1)and(43.3,-114.4)..(42.3,-112.2)
				..controls(41.8,-111.3)and(41.3,-110.4)..(41.1,-110.1)
				..controls(40.5,-109.4)and(40.3,-109)..(40,-107.8)
				..controls(39.8,-107)and(39.6,-106.5)..(39.2,-106)
				..controls(38.9,-105.6)and(38.4,-104.7)..(38,-104)
				..controls(37.7,-103.4)and(37.2,-102.5)..(37,-102.1)
				..controls(36.7,-101.7)and(36.3,-100.8)..(36,-100.1)
				..controls(35.6,-98.9)and(35.3,-98.4)..(34.6,-97.6)
				..controls(34.2,-97.1)and(34.2,-96.3)..(34.6,-96)
				..controls(35,-95.7)and(35.5,-95.9)..(36.2,-96.7)
				..controls(36.5,-97)and(37.2,-97.5)..(37.6,-97.7)..controls(38,-98)and(38.6,-98.4)..(39.1,-98.9)
				..controls(39.6,-99.4)and(40.3,-100)..(40.6,-100.3)
				..controls(41.3,-100.8)and(42.4,-101.8)..(43.7,-102.9)
				..controls(44,-103.2)and(44.6,-103.7)..(45,-103.9)
				..controls(45.4,-104.2)and(45.9,-104.6)..(46.1,-104.8)
				..controls(47.1,-105.9)and(48.8,-107.1)..(50.5,-107.9)
				..controls(51.8,-108.6)and(52.3,-109)..(53.4,-110)
				..controls(54.8,-111.2)and(55.6,-111.6)..(57,-112.1)
				..controls(57.5,-112.3)and(58.5,-112.8)..(59.4,-113.3)
				..controls(60.2,-113.7)and(61.3,-114.3)..(61.9,-114.5)
				..controls(62.4,-114.7)and(63.3,-115.1)..(63.8,-115.3)
				..controls(64.4,-115.6)and(65.4,-115.9)..(66.2,-116.1)
				..controls(67,-116.3)and(68.4,-116.6)..(69.2,-116.9)
				..controls(72.3,-117.9)and(72.8,-118)..(78,-118.1)
				..controls(84,-118.3)and(86.1,-118.1)..(88.1,-117.1)
				..controls(89.4,-116.5)and(93,-115)..(93.7,-114.9)
				..controls(94.8,-114.6)and(95.4,-114.2)..(95.9,-113.4)
				..controls(96.2,-113)and(96.7,-112.5)..(97.1,-112.2)
				..controls(98,-111.6)and(101.3,-108.4)..(102.1,-107.2)
				..controls(102.4,-106.8)and(102.9,-106.2)..(103.2,-105.8)
				..controls(103.5,-105.5)and(104,-104.8)..(104.2,-104.2)
				..controls(104.5,-103.6)and(105,-102.8)..(105.3,-102.4)
				..controls(105.6,-102)and(106.1,-101.2)..(106.5,-100.6)
				..controls(107.1,-99.5)and(108.6,-97.5)..(109.5,-96.5)
				..controls(109.9,-96.1)and(110.1,-96)..(110.4,-96)
				..controls(111.1,-96)and(111.6,-96.5)..(111.8,-97.6)
				..controls(112,-98.6)and(112.3,-99.3)..(112.8,-99.9)
				..controls(113.7,-100.9)and(115.2,-103.1)..(115.7,-104.2)
				..controls(116.1,-105)and(116.5,-105.5)..(116.9,-106)
				..controls(117.3,-106.4)and(117.8,-106.9)..(118,-107.3)
				..controls(118.2,-107.6)and(118.6,-108)..(119,-108.3)
				..controls(119.6,-108.8)and(121.1,-110.3)..(121.8,-111.2)
				..controls(122.1,-111.7)and(122.5,-112)..(123.4,-112.4)
				..controls(124.8,-113.2)and(125.8,-113.8)..(127,-114.6)
				..controls(127.5,-114.9)and(128.1,-115.3)..(128.2,-115.4)
				..controls(128.4,-115.5)and(129.3,-115.8)..(130.2,-116)
				..controls(131.1,-116.2)and(132.5,-116.6)..(133.3,-116.9)
				..controls(135.7,-117.9)and(136.1,-117.9)..(142.5,-117.8)
				..controls(145.8,-117.8)and(148.4,-117.7)..(148.8,-117.6)
				..controls(149.2,-117.6)and(150.2,-117.2)..(151,-116.8)
				..controls(152,-116.4)and(152.8,-116)..(153.6,-115.9)
				..controls(154.5,-115.7)and(155.2,-115.4)..(156.3,-114.9)
				..controls(157.1,-114.5)and(158.2,-114)..(158.7,-113.8)
				..controls(159.1,-113.7)and(159.8,-113.3)..(160.2,-113)
				..controls(160.6,-112.7)and(161.3,-112.4)..(161.7,-112.2)
				..controls(162.7,-111.9)and(163.2,-111.7)..(164.1,-111.1)
				..controls(164.5,-110.8)and(165.3,-110.4)..(165.9,-110)
				..controls(166.5,-109.7)and(167.3,-109.2)..(167.6,-108.9)
				..controls(167.9,-108.6)and(168.6,-108.2)..(169,-108)
				..controls(169.4,-107.8)and(170.2,-107.3)..(170.7,-106.9)
				..controls(171.2,-106.5)and(172.1,-105.8)..(172.7,-105.4)
				..controls(173.2,-105)and(174.1,-104.4)..(174.5,-104)
				..controls(174.9,-103.7)and(175.7,-103.1)..(176.2,-102.6)
				..controls(176.7,-102.2)and(177.4,-101.5)..(177.9,-101.2)
				..controls(178.3,-100.8)and(178.8,-100.4)..(179,-100.2)
				..controls(179.2,-100)and(179.6,-99.6)..(180,-99.3)
				..controls(180.4,-99.1)and(180.9,-98.6)..(181.2,-98.4)
				..controls(181.5,-98.1)and(182.1,-97.7)..(182.5,-97.5)
				..controls(183.2,-97.2)and(183.5,-96.9)..(184.3,-95.9)
				..controls(184.6,-95.6)and(184.6,-95.6)..(184.8,-95.7)
				..controls(185,-95.8)and(185.1,-96)..(185.1,-96.2)
				..controls(185.1,-96.5)and(185,-96.7)..(184.6,-97.2)
				..controls(184.1,-97.8)and(184,-97.9)..(183.9,-98.8)
				..controls(183.7,-100.3)and(183.3,-101.3)..(182.9,-101.8)
				..controls(182.6,-102.1)and(182.3,-102.6)..(182.2,-102.9)
				..controls(182,-103.3)and(181.6,-104)..(181.2,-104.6)
				..controls(180.9,-105.2)and(180.5,-106.1)..(180.4,-106.6)
				..controls(179.9,-108)and(179.7,-108.5)..(179,-109.4)
				..controls(178.6,-109.8)and(178.2,-110.5)..(177.9,-111)
				..controls(177.7,-111.5)and(177.3,-112.3)..(177,-112.7)
				..controls(176.6,-113.2)and(176.4,-113.7)..(176.3,-114.2)
				..controls(176.2,-114.7)and(175.9,-115.2)..(175.4,-116)
				..controls(175,-116.6)and(174.4,-117.7)..(174,-118.5)
				..controls(173.7,-119.2)and(173.3,-120.1)..(173,-120.5)
				..controls(172.8,-120.9)and(172.4,-121.8)..(172.1,-122.5)
				..controls(171.8,-123.2)and(171.3,-124.1)..(171,-124.6)
				..controls(170.8,-125.1)and(170.3,-126)..(169.9,-126.6)
				..controls(169.6,-127.3)and(169.2,-128.1)..(169,-128.3)
				..controls(168.8,-128.6)and(168.5,-129.3)..(168.3,-129.8)
				..controls(168.1,-130.3)and(167.8,-131.1)..(167.5,-131.6)
				..controls(166.9,-132.5)and(164.2,-137.9)..(164.2,-138.1)
				..controls(164.2,-138.1)and(163.9,-138.8)..(163.4,-139.7)
				..controls(162.9,-140.5)and(162.4,-141.5)..(162.3,-141.9)
				..controls(161.9,-143)and(161.6,-143.4)..(160.9,-144.1)
				..controls(160.4,-144.6)and(160.3,-144.8)..(160.1,-145.4)
				..controls(160,-145.8)and(159.6,-146.8)..(159.2,-147.8)
				..controls(158.7,-148.8)and(158.1,-150)..(157.9,-150.6)
				..controls(157.6,-151.1)and(157.2,-151.9)..(156.9,-152.4)
				..controls(156.5,-152.8)and(156.1,-153.7)..(155.8,-154.4)
				..controls(155.6,-155)and(155.2,-155.9)..(154.9,-156.3)
				..controls(154.7,-156.7)and(154.4,-157.4)..(154.2,-157.9)
				..controls(154.1,-158.4)and(153.8,-159)..(153.5,-159.3)
				..controls(152.8,-160.1)and(153,-160.1)..(142.4,-160.1)
				..controls(137.2,-160.1)and(131.8,-160.1)..(130.4,-160)--cycle
				(141.5,-103.7)..controls(140.8,-103)and(139.8,-102.1)..(139.2,-101.6)
				..controls(138.5,-101.2)and(137.9,-100.6)..(137.8,-100.4)
				..controls(137,-99.5)and(134,-96.5)..(133.2,-96)
				..controls(132.7,-95.7)and(131.8,-94.9)..(131.1,-94.3)
				..controls(129.4,-92.7)and(127.9,-91.6)..(126.6,-91.1)
				..controls(125.2,-90.5)and(124.1,-90)..(122.9,-89.1)
				..controls(122.1,-88.6)and(121.3,-88.2)..(119.9,-87.7)
				..controls(118.8,-87.3)and(117.5,-86.8)..(116.9,-86.5)
				..controls(114.3,-85.5)and(113.7,-85.4)..(109.7,-85.5)
				..controls(105.3,-85.7)and(103.9,-85.9)..(101.6,-87)
				..controls(100.7,-87.4)and(99.4,-87.8)..(97.9,-88)
				..controls(97,-88.2)and(96.6,-88.3)..(95.8,-88.9)
				..controls(95.4,-89.2)and(94.7,-89.6)..(94.2,-89.8)
				..controls(92.9,-90.3)and(92.3,-90.7)..(90.9,-92.1)
				..controls(90.2,-92.7)and(89.4,-93.3)..(89,-93.5)..controls(87,-94.5)and(85.2,-96)..(83.6,-97.9)
				..controls(83.3,-98.2)and(82.5,-98.9)..(81.8,-99.4)
				..controls(80.7,-100.1)and(80.5,-100.3)..(80,-101.1)
				..controls(78.7,-103.1)and(77.2,-103.9)..(76.4,-103)
				..controls(75.7,-102.2)and(75.7,-103.2)..(75.7,-83.6)
				..controls(75.6,-73.6)and(75.6,-64.8)..(75.5,-64)
				..controls(75.5,-63.2)and(75.5,-61.4)..(75.6,-59.8)--(75.7,-57.1)--(76,-56.8)
				..controls(76.8,-56.2)and(78,-56.3)..(79.8,-57.2)
				..controls(80.4,-57.5)and(81.3,-57.8)..(81.8,-57.9)
				..controls(82.3,-58.1)and(83,-58.4)..(83.6,-58.8)..controls(84.3,-59.2)and(85,-59.5)..(86,-59.8)
				..controls(86.8,-60.1)and(88,-60.6)..(88.7,-60.9)
				..controls(89.3,-61.2)and(90.6,-61.7)..(91.4,-62)
				..controls(92.2,-62.3)and(93.7,-62.9)..(94.6,-63.3)
				..controls(95.6,-63.7)and(96.8,-64.1)..(97.4,-64.3)
				..controls(98,-64.5)and(99,-64.9)..(99.7,-65.2)
				..controls(100.4,-65.6)and(101.5,-66)..(102.3,-66.2)
				..controls(104.2,-66.8)and(105,-67.1)..(106.2,-67.8)
				..controls(107.9,-68.7)and(110.2,-69.4)..(111.3,-69.2)
				..controls(112.5,-69)and(115,-68.2)..(115.6,-67.8)
				..controls(117.3,-66.7)and(117.6,-66.6)..(119.5,-66.2)
				..controls(121.9,-65.8)and(122.1,-65.7)..(123,-65.1)
				..controls(124.1,-64.4)and(124.4,-64.3)..(126,-63.9)
				..controls(127,-63.7)and(127.7,-63.4)..(128.4,-63.1)
				..controls(129,-62.8)and(129.7,-62.5)..(130.6,-62.4)
				..controls(132.1,-62.1)and(134,-61.5)..(135.2,-61)
				..controls(135.7,-60.8)and(136.7,-60.5)..(137.3,-60.3)
				..controls(138.4,-60)and(139.3,-59.7)..(141.1,-59)
				..controls(141.6,-58.7)and(142,-58.6)..(142.5,-58.6)
				..controls(143.2,-58.6)and(143.3,-58.7)..(143.6,-59)
				..controls(144,-59.5)and(144.1,-59.9)..(144,-60.8)
				..controls(144,-61.2)and(144,-71.1)..(144,-82.8)--(143.9,-104.1)--(143.7,-104.5)
				..controls(143.5,-104.7)and(143.3,-104.8)..(143,-104.9)
				..controls(142.6,-104.9)and(142.5,-104.8)..(141.5,-103.7)--cycle
				(27.8,-91.4)..controls(27.1,-90.9)and(26.3,-90)..(26.1,-89.5)
				..controls(26,-89.2)and(25.8,-88.6)..(25.7,-87.9)
				..controls(25.4,-86.2)and(25.6,-84.8)..(26.5,-83.5)
				..controls(27.4,-82.3)and(28.3,-82)..(29.4,-82.3)
				..controls(31.1,-82.8)and(32.1,-84.9)..(32,-87.4)
				..controls(31.9,-89.3)and(31.2,-90.6)..(29.9,-91.4)
				..controls(29.1,-91.9)and(28.6,-91.9)..(27.8,-91.4)--cycle
				(190.1,-91.3)..controls(189.7,-91.2)and(189.3,-90.9)..(189.1,-90.7)
				..controls(188.3,-89.8)and(187.6,-87.7)..(187.6,-86.2)
				..controls(187.6,-85.5)and(187.7,-85.2)..(188.2,-84.4)
				..controls(189.1,-82.6)and(190.1,-81.8)..(191.3,-82.1)
				..controls(192,-82.3)and(193.5,-83.7)..(193.8,-84.4)
				..controls(194.1,-85.1)and(194.2,-87.1)..(194,-88.3)
				..controls(193.8,-89.8)and(193.3,-90.5)..(192.3,-91.3)
				..controls(191.7,-91.7)and(190.9,-91.7)..(190.1,-91.3)--cycle
				(173.7,-83.4)..controls(172.2,-82.5)and(172,-82.5)..(170.6,-82.4)
				..controls(169.7,-82.4)and(168.8,-82.4)..(168.3,-82.5)
				..controls(167.2,-82.7)and(166.9,-82.7)..(166.6,-82.2)
				..controls(166.3,-81.8)and(166.3,-81)..(166.8,-80)
				..controls(167.2,-78.9)and(167.4,-77.4)..(167.7,-74.1)
				..controls(167.8,-72.8)and(167.9,-71.8)..(167.9,-71.7)
				..controls(168,-71.7)and(168,-70.2)..(168,-68.4)..controls(168,-65.4)and(168,-64.9)..(167.7,-63)
				..controls(167.3,-60.5)and(167.1,-59.9)..(166.4,-58.4)
				..controls(166.1,-57.7)and(165.8,-56.9)..(165.7,-56.6)
				..controls(165.7,-56.2)and(165.4,-55.6)..(165.2,-55.2)
				..controls(164.6,-53.8)and(161.9,-50.7)..(161,-50.3)
				..controls(159.8,-49.7)and(157.9,-50)..(155.2,-51.2)
				..controls(154.7,-51.4)and(153.4,-51.8)..(152.2,-52.1)
				..controls(150.8,-52.4)and(149.4,-52.9)..(148.6,-53.2)
				..controls(147.8,-53.6)and(146.7,-54)..(146,-54.2)
				..controls(145.4,-54.4)and(144.5,-54.7)..(144,-54.9)
				..controls(143.6,-55.2)and(142.7,-55.5)..(141.8,-55.7)
				..controls(141,-55.9)and(139.6,-56.3)..(138.7,-56.6)
				..controls(137.8,-57)and(136.3,-57.4)..(135.3,-57.6)
				..controls(133.8,-57.9)and(133.2,-58.1)..(132.1,-58.7)
				..controls(131.1,-59.2)and(130.3,-59.4)..(129.4,-59.6)
				..controls(128.7,-59.8)and(127.8,-60)..(127.4,-60.1)
				..controls(127,-60.3)and(126,-60.5)..(125.1,-60.7)
				..controls(123.2,-61.1)and(121.7,-61.5)..(120.2,-62.1)
				..controls(119.7,-62.3)and(118.1,-62.8)..(116.7,-63.2)
				..controls(114.6,-63.9)and(114,-64.1)..(113,-64.7)
				..controls(111.9,-65.4)and(111.8,-65.4)..(111,-65.5)
				..controls(109.9,-65.5)and(109.3,-65.4)..(107.7,-64.7)
				..controls(107,-64.4)and(106,-64)..(105.3,-63.9)
				..controls(103.5,-63.5)and(102.2,-63)..(101.2,-62.5)
				..controls(100.1,-61.8)and(99.3,-61.5)..(97.8,-61.1)
				..controls(97.2,-61)and(96.2,-60.7)..(95.7,-60.4)
				..controls(95.1,-60.1)and(94.2,-59.8)..(93.5,-59.6)
				..controls(92.9,-59.4)and(91.7,-59)..(91,-58.7)..controls(90.2,-58.4)and(88.9,-58)..(88,-57.7)
				..controls(87.2,-57.5)and(86.1,-57.1)..(85.5,-56.8)
				..controls(85,-56.6)and(84.1,-56.2)..(83.5,-56)..controls(82.9,-55.7)and(81.8,-55.2)..(81,-54.9)
				..controls(80.1,-54.5)and(79.1,-54.2)..(78.7,-54.1)
				..controls(77.9,-53.8)and(77.3,-53.6)..(76.2,-52.9)
				..controls(75.7,-52.6)and(75.2,-52.4)..(74.1,-52.2)
				..controls(72.5,-51.9)and(72.4,-51.8)..(71.3,-51.1)
				..controls(69.9,-50.1)and(69.4,-49.9)..(67.9,-49.8)
				..controls(67,-49.7)and(66.6,-49.6)..(66.3,-49.4)
				..controls(66.1,-49.3)and(65.7,-49)..(65.4,-48.8)
				..controls(65.1,-48.7)and(64.2,-48.4)..(63.6,-48.2)
				..controls(62.9,-48)and(61.7,-47.6)..(60.9,-47.2)
				..controls(60.1,-46.8)and(58.9,-46.4)..(58.2,-46.2)
				..controls(57.6,-46)and(56.8,-45.7)..(56.5,-45.5)
				..controls(54.9,-44.6)and(54.6,-44.5)..(53.2,-44.2)
				..controls(51.7,-43.8)and(51,-43.6)..(49.9,-42.9)
				..controls(49.5,-42.6)and(48.9,-42.4)..(47.5,-42)
				..controls(45.7,-41.6)and(45.6,-41.5)..(44.3,-40.6)
				..controls(43.5,-40.1)and(42.8,-39.6)..(42.7,-39.4)
				..controls(42.4,-38.7)and(43,-38.1)..(44.5,-37.6)..controls(44.9,-37.5)and(46,-37)..(46.9,-36.6)
				..controls(47.9,-36.1)and(48.8,-35.8)..(49.6,-35.6)
				..controls(50.5,-35.4)and(51.3,-35.2)..(52.4,-34.6)
				..controls(53.3,-34.1)and(54.5,-33.7)..(54.9,-33.5)
				..controls(55.4,-33.4)and(56.4,-33)..(57.2,-32.6)
				..controls(58.6,-32)and(60.1,-31.5)..(61.6,-31.2)
				..controls(62.3,-31.1)and(62.7,-31)..(63.2,-30.7)
				..controls(64.1,-30)and(64.3,-29.9)..(65.3,-29.7)
				..controls(65.7,-29.6)and(66.5,-29.2)..(67.1,-28.9)
				..controls(67.6,-28.6)and(68.8,-28.2)..(69.6,-27.9)
				..controls(70.5,-27.7)and(71.6,-27.3)..(72.1,-27)
				..controls(72.6,-26.7)and(73.8,-26.2)..(74.8,-25.9)
				..controls(75.7,-25.5)and(76.7,-25.1)..(77.1,-25)
				..controls(77.4,-24.8)and(78.5,-24.4)..(79.6,-24.1)
				..controls(80.6,-23.8)and(82,-23.2)..(82.8,-22.9)
				..controls(83.5,-22.5)and(84.8,-22)..(85.8,-21.8)
				..controls(86.7,-21.5)and(87.8,-21.1)..(88.1,-20.9)
				..controls(88.4,-20.7)and(89.3,-20.3)..(90.1,-20.1)
				..controls(90.8,-19.8)and(91.9,-19.4)..(92.4,-19.1)
				..controls(93.5,-18.5)and(94.4,-18.2)..(95.4,-18.1)
				..controls(95.7,-18)and(96.4,-17.8)..(96.9,-17.5)
				..controls(97.5,-17.3)and(98.4,-16.9)..(99,-16.6)
				..controls(99.7,-16.3)and(100.9,-15.9)..(101.9,-15.6)
				..controls(102.9,-15.3)and(104.4,-14.7)..(105.3,-14.3)
				..controls(108.1,-13.2)and(108.5,-13.1)..(110.1,-13.1)
				..controls(111.4,-13.1)and(111.5,-13.1)..(112.4,-13.5)
				..controls(113.2,-14)and(115,-14.8)..(116.6,-15.4)
				..controls(116.9,-15.5)and(117.9,-15.8)..(118.8,-16.1)
				..controls(120.9,-16.6)and(121.5,-16.9)..(122.6,-17.9)
				..controls(123,-18.3)and(123.7,-18.8)..(123.9,-18.9)
				..controls(124.6,-19.4)and(126,-19.8)..(127.9,-20.2)
				..controls(129.6,-20.5)and(129.8,-20.6)..(131.1,-21.4)
				..controls(131.5,-21.7)and(132.4,-22.1)..(133.1,-22.3)
				..controls(136.1,-23.2)and(138.3,-24.1)..(140.2,-25)
				..controls(141,-25.5)and(141.5,-25.6)..(142.5,-25.8)
				..controls(145.1,-26.3)and(146,-26.6)..(149.4,-28.1)
				..controls(151.3,-28.9)and(153.5,-29.8)..(154.3,-30.1)
				..controls(155.1,-30.4)and(156.2,-30.8)..(156.6,-31)
				..controls(157.1,-31.2)and(158.2,-31.6)..(159.1,-31.8)
				..controls(160.6,-32.2)and(160.9,-32.3)..(161.7,-32.9)
				..controls(162.4,-33.3)and(162.9,-33.6)..(163.7,-33.8)
				..controls(164.3,-34)and(164.8,-34.1)..(164.9,-34.1)
				..controls(165,-34.1)and(165.6,-34.4)..(166.4,-34.8)
				..controls(167.3,-35.2)and(168.1,-35.4)..(169,-35.6)
				..controls(170,-35.8)and(170.5,-36)..(171.7,-36.6)
				..controls(172.4,-37)and(173.3,-37.4)..(173.6,-37.4)
				..controls(175.6,-37.7)and(176.2,-38)..(176.7,-39.1)
				..controls(176.9,-39.6)and(176.9,-39.7)..(176.7,-40.4)
				..controls(176.4,-41.6)and(175.8,-42.7)..(175.2,-43.4)
				..controls(174.5,-44.1)and(172.6,-45.5)..(171.7,-45.8)
				..controls(170.6,-46.2)and(169.1,-46)..(168.2,-45.4)
				..controls(167.9,-45.2)and(167.3,-44.8)..(166.9,-44.4)
				..controls(166.4,-44.1)and(165.8,-43.6)..(165.4,-43.2)
				..controls(165.1,-42.9)and(164.4,-42.3)..(163.9,-41.9)
				..controls(163.3,-41.5)and(162.7,-40.9)..(162.4,-40.5)
				..controls(161.4,-39.2)and(159.9,-38)..(158.8,-37.8)
				..controls(158.4,-37.8)and(158.3,-37.7)..(158.2,-37.2)
				..controls(158,-36.8)and(157.4,-36.1)..(157.1,-36.1)
				..controls(157,-36.1)and(156.7,-36.1)..(156.5,-36)
				..controls(156.1,-35.9)and(155.9,-35.7)..(155.7,-35.2)
				..controls(155.3,-34.6)and(155.1,-34.4)..(154,-34.1)
				..controls(153.6,-34)and(153.2,-33.7)..(152.9,-33.5)
				..controls(152.4,-32.9)and(151.6,-32.4)..(151.3,-32.4)
				..controls(150.9,-32.4)and(150.4,-32.6)..(150.4,-32.9)
				..controls(150.4,-33.2)and(150.8,-33.4)..(151.4,-33.6)
				..controls(152.1,-33.7)and(152.8,-34.1)..(153.5,-34.9)
				..controls(154.1,-35.5)and(154.3,-35.6)..(155,-35.8)
				..controls(155.7,-36)and(155.7,-36)..(156,-36.6)
				..controls(156.3,-37.3)and(156.7,-37.7)..(157.5,-37.8)
				..controls(158.1,-37.9)and(158.2,-38)..(158.4,-38.8)
				..controls(158.5,-39.4)and(158.6,-39.6)..(159.6,-40.4)
				..controls(160.2,-41)and(160.9,-41.6)..(161.3,-41.8)
				..controls(161.7,-42)and(161.9,-42.3)..(162.2,-42.8)
				..controls(162.4,-43.2)and(162.8,-43.8)..(163.1,-44.1)
				..controls(163.9,-44.8)and(165.1,-46.7)..(165.7,-48)
				..controls(166.2,-49.3)and(166.4,-49.7)..(167.2,-50.5)
				..controls(167.5,-50.8)and(167.9,-51.4)..(168.1,-51.8)
				..controls(168.3,-52.2)and(168.8,-53)..(169.1,-53.6)
				..controls(169.6,-54.3)and(169.8,-54.9)..(169.9,-55.6)
				..controls(170.4,-57.8)and(170.5,-58.1)..(171.2,-59.5)
				..controls(171.6,-60.3)and(172,-61.5)..(172.2,-62.1)
				..controls(172.3,-62.8)and(172.7,-63.7)..(172.9,-64.1)
				..controls(173.2,-64.6)and(173.4,-65.3)..(173.4,-65.6)
				..controls(173.8,-68.8)and(174.2,-70.6)..(174.9,-72.6)
				..controls(175.5,-74.5)and(175.9,-77.1)..(176.1,-80.2)
				..controls(176.3,-83.4)and(175.5,-84.4)..(173.7,-83.4)--cycle;
				\definecolor{ezs}{RGB}{247,179,26}\fill[ezs](-5.7,-156.6)--(-3.4,-157.5)--(61.9,12.7)--(59.7,13.6)--cycle;
				\definecolor{ezs}{RGB}{247,179,22}\fill[ezs](-3.8,-157.3)--(-1.5,-158.2)--(63.9,12)--(61.6,12.9)--cycle;
				\definecolor{ezs}{RGB}{248,179,19}\fill[ezs](-1.8,-158.1)--(0.4,-158.9)--(65.8,11.2)--(63.5,12.1)--cycle;
				\definecolor{ezs}{RGB}{248,178,16}\fill[ezs](0.1,-158.8)--(2.4,-159.7)--(67.7,10.5)--(65.4,11.4)--cycle;
				\definecolor{ezs}{RGB}{249,178,15}\fill[ezs](2,-159.5)--(4.3,-160.4)--(69.6,9.8)--(67.3,10.7)--cycle;
				\definecolor{ezs}{RGB}{251,177,17}\fill[ezs](3.9,-160.3)--(6.2,-161.1)--(71.5,9)--(69.2,9.9)--cycle;
				\definecolor{ezs}{RGB}{252,179,18}\fill[ezs](5.8,-161)--(8.1,-161.9)--(73.4,8.3)--(71.1,9.2)--cycle;
				\definecolor{ezs}{RGB}{252,182,17}\fill[ezs](7.7,-161.7)--(10,-162.6)--(75.3,7.6)--(73,8.5)--cycle;
				\definecolor{ezs}{RGB}{252,185,17}\fill[ezs](9.6,-162.5)--(11.9,-163.4)--(77.2,6.8)--(75,7.7)--cycle;
				\definecolor{ezs}{RGB}{251,189,18}\fill[ezs](11.5,-163.2)--(13.8,-164.1)--(79.2,6.1)--(76.9,7)--cycle;
				\definecolor{ezs}{RGB}{251,191,22}\fill[ezs](13.4,-163.9)--(15.7,-164.8)--(81.1,5.4)--(78.8,6.2)--cycle;
				\definecolor{ezs}{RGB}{252,193,25}\fill[ezs](15.4,-164.7)--(17.6,-165.6)--(83,4.6)--(80.7,5.5)--cycle;
				\definecolor{ezs}{RGB}{253,195,27}\fill[ezs](17.3,-165.4)--(19.6,-166.3)--(84.9,3.9)--(82.6,4.8)--cycle;
				\definecolor{ezs}{RGB}{253,196,31}\fill[ezs](19.2,-166.1)--(21.5,-167)--(86.8,3.2)--(84.5,4)--cycle;
				\definecolor{ezs}{RGB}{254,199,35}\fill[ezs](21.1,-166.9)--(23.4,-167.8)--(88.7,2.4)--(86.4,3.3)--cycle;
				\definecolor{ezs}{RGB}{253,201,38}\fill[ezs](23,-167.6)--(25.3,-168.5)--(90.6,1.7)--(88.3,2.6)--cycle;
				\definecolor{ezs}{RGB}{253,203,40}\fill[ezs](24.9,-168.3)--(27.2,-169.2)--(92.5,1)--(90.3,1.8)--cycle;
				\definecolor{ezs}{RGB}{253,205,42}\fill[ezs](26.8,-169.1)--(29.1,-170)--(94.4,0.2)--(92.2,1.1)--cycle;
				\definecolor{ezs}{RGB}{253,207,45}\fill[ezs](28.7,-169.8)--(31,-170.7)--(96.4,-0.5)--(94.1,0.4)--cycle;
				\definecolor{ezs}{RGB}{253,210,49}\fill[ezs](30.7,-170.5)--(32.9,-171.4)--(98.3,-1.2)--(96,-0.4)--cycle;
				\definecolor{ezs}{RGB}{253,212,54}\fill[ezs](32.6,-171.3)--(34.9,-172.2)--(100.2,-2)--(97.9,-1.1)--cycle;
				\definecolor{ezs}{RGB}{253,214,59}\fill[ezs](34.5,-172)--(36.8,-172.9)--(102.1,-2.7)--(99.8,-1.8)--cycle;
				\definecolor{ezs}{RGB}{253,216,63}\fill[ezs](36.4,-172.7)--(38.7,-173.6)--(104,-3.4)--(101.7,-2.6)--cycle;
				\definecolor{ezs}{RGB}{253,218,65}\fill[ezs](38.3,-173.5)--(40.6,-174.4)--(105.9,-4.2)--(103.6,-3.3)--cycle;
				\definecolor{ezs}{RGB}{254,220,66}\fill[ezs](40.2,-174.2)--(42.5,-175.1)--(107.8,-4.9)--(105.5,-4)--cycle;
				\definecolor{ezs}{RGB}{254,222,68}\fill[ezs](42.1,-175)--(44.4,-175.8)--(109.7,-5.6)--(107.5,-4.8)--cycle;
				\definecolor{ezs}{RGB}{254,224,71}\fill[ezs](44,-175.7)--(46.3,-176.6)--(111.7,-6.4)--(109.4,-5.5)--cycle;
				\definecolor{ezs}{RGB}{254,226,75}\fill[ezs](46,-176.4)--(48.2,-177.3)--(113.6,-7.1)--(111.3,-6.2)--cycle;
				\definecolor{ezs}{RGB}{254,228,79}\fill[ezs](47.9,-177.2)--(50.1,-178)--(115.5,-7.8)--(113.2,-7)--cycle;
				\definecolor{ezs}{RGB}{254,230,82}\fill[ezs](49.8,-177.9)--(52.1,-178.8)--(117.4,-8.6)--(115.1,-7.7)--cycle;
				\definecolor{ezs}{RGB}{254,232,85}\fill[ezs](51.7,-178.6)--(54,-179.5)--(119.3,-9.3)--(117,-8.4)--cycle;
				\definecolor{ezs}{RGB}{254,234,88}\fill[ezs](53.6,-179.4)--(55.9,-180.2)--(121.2,-10)--(118.9,-9.2)--cycle;
				\definecolor{ezs}{RGB}{254,237,91}\fill[ezs](55.5,-180.1)--(57.8,-181)--(123.1,-10.8)--(120.8,-9.9)--cycle;
				\definecolor{ezs}{RGB}{254,239,94}\fill[ezs](57.4,-180.8)--(59.7,-181.7)--(125,-11.5)--(122.8,-10.6)--cycle;
				\definecolor{ezs}{RGB}{254,241,97}\fill[ezs](59.3,-181.6)--(61.6,-182.4)--(127,-12.2)--(124.7,-11.4)--cycle;
				\definecolor{ezs}{RGB}{254,243,101}\fill[ezs](61.2,-182.3)--(63.5,-183.2)--(128.9,-13)--(126.6,-12.1)--cycle;
				\definecolor{ezs}{RGB}{254,245,103}\fill[ezs](63.2,-183)--(65.4,-183.9)--(130.8,-13.7)--(128.5,-12.8)--cycle;
				\definecolor{ezs}{RGB}{254,245,105}\fill[ezs](65.1,-183.8)--(67.4,-184.6)--(132.7,-14.4)--(130.4,-13.6)--cycle;
				\definecolor{ezs}{RGB}{254,245,105}\fill[ezs](67,-184.5)--(69.3,-185.4)--(134.6,-15.2)--(132.3,-14.3)--cycle;
				\definecolor{ezs}{RGB}{254,244,103}\fill[ezs](68.9,-185.2)--(71.2,-186.1)--(136.5,-15.9)--(134.2,-15)--cycle;
				\definecolor{ezs}{RGB}{253,243,101}\fill[ezs](70.8,-186)--(73.1,-186.8)--(138.4,-16.6)--(136.1,-15.8)--cycle;
				\definecolor{ezs}{RGB}{253,241,98}\fill[ezs](72.7,-186.7)--(75,-187.6)--(140.3,-17.4)--(138,-16.5)--cycle;
				\definecolor{ezs}{RGB}{253,239,95}\fill[ezs](74.6,-187.4)--(76.9,-188.3)--(142.2,-18.1)--(140,-17.2)--cycle;
				\definecolor{ezs}{RGB}{253,236,92}\fill[ezs](76.5,-188.2)--(78.8,-189)--(144.2,-18.8)--(141.9,-18)--cycle;
				\definecolor{ezs}{RGB}{253,234,88}\fill[ezs](78.5,-188.9)--(80.7,-189.8)--(146.1,-19.6)--(143.8,-18.7)--cycle;
				\definecolor{ezs}{RGB}{253,232,85}\fill[ezs](80.4,-189.6)--(82.7,-190.5)--(148,-20.3)--(145.7,-19.4)--cycle;
				\definecolor{ezs}{RGB}{253,230,82}\fill[ezs](82.3,-190.4)--(84.6,-191.2)--(149.9,-21)--(147.6,-20.2)--cycle;
				\definecolor{ezs}{RGB}{253,228,78}\fill[ezs](84.2,-191.1)--(86.5,-192)--(151.8,-21.8)--(149.5,-20.9)--cycle;
				\definecolor{ezs}{RGB}{253,226,74}\fill[ezs](86.1,-191.8)--(88.4,-192.7)--(153.7,-22.5)--(151.4,-21.6)--cycle;
				\definecolor{ezs}{RGB}{253,224,70}\fill[ezs](88,-192.6)--(90.3,-193.4)--(155.6,-23.3)--(153.3,-22.4)--cycle;
				\definecolor{ezs}{RGB}{253,222,67}\fill[ezs](89.9,-193.3)--(92.2,-194.2)--(157.5,-24)--(155.3,-23.1)--cycle;
				\definecolor{ezs}{RGB}{253,220,65}\fill[ezs](91.8,-194)--(94.1,-194.9)--(159.5,-24.7)--(157.2,-23.8)--cycle;
				\definecolor{ezs}{RGB}{253,217,63}\fill[ezs](93.7,-194.8)--(96,-195.6)--(161.4,-25.5)--(159.1,-24.6)--cycle;
				\definecolor{ezs}{RGB}{253,215,60}\fill[ezs](95.7,-195.5)--(97.9,-196.4)--(163.3,-26.2)--(161,-25.3)--cycle;
				\definecolor{ezs}{RGB}{253,213,56}\fill[ezs](97.6,-196.2)--(99.9,-197.1)--(165.2,-26.9)--(162.9,-26)--cycle;
				\definecolor{ezs}{RGB}{252,211,53}\fill[ezs](99.5,-197)--(101.8,-197.8)--(167.1,-27.7)--(164.8,-26.8)--cycle;
				\definecolor{ezs}{RGB}{252,209,49}\fill[ezs](101.4,-197.7)--(103.7,-198.6)--(169,-28.4)--(166.7,-27.5)--cycle;
				\definecolor{ezs}{RGB}{252,207,45}\fill[ezs](103.3,-198.4)--(105.6,-199.3)--(170.9,-29.1)--(168.6,-28.2)--cycle;
				\definecolor{ezs}{RGB}{253,205,42}\fill[ezs](105.2,-199.2)--(107.5,-200)--(172.8,-29.9)--(170.6,-29)--cycle;
				\definecolor{ezs}{RGB}{253,202,40}\fill[ezs](107.1,-199.9)--(109.4,-200.8)--(174.7,-30.6)--(172.5,-29.7)--cycle;
				\definecolor{ezs}{RGB}{253,200,38}\fill[ezs](109,-200.6)--(111.3,-201.5)--(176.7,-31.3)--(174.4,-30.4)--cycle;
				\definecolor{ezs}{RGB}{253,198,35}\fill[ezs](111,-201.4)--(113.2,-202.2)--(178.6,-32.1)--(176.3,-31.2)--cycle;
				\definecolor{ezs}{RGB}{253,196,31}\fill[ezs](112.9,-202.1)--(115.2,-203)--(180.5,-32.8)--(178.2,-31.9)--cycle;
				\definecolor{ezs}{RGB}{253,194,28}\fill[ezs](114.8,-202.8)--(117.1,-203.7)--(182.4,-33.5)--(180.1,-32.6)--cycle;
				\definecolor{ezs}{RGB}{253,192,26}\fill[ezs](116.7,-203.6)--(119,-204.5)--(184.3,-34.3)--(182,-33.4)--cycle;
				\definecolor{ezs}{RGB}{253,189,24}\fill[ezs](118.6,-204.3)--(120.9,-205.2)--(186.2,-35)--(183.9,-34.1)--cycle;
				\definecolor{ezs}{RGB}{253,186,20}\fill[ezs](120.5,-205)--(122.8,-205.9)--(188.1,-35.7)--(185.8,-34.8)--cycle;
				\definecolor{ezs}{RGB}{253,184,17}\fill[ezs](122.4,-205.8)--(124.7,-206.7)--(190,-36.5)--(187.8,-35.6)--cycle;
				\definecolor{ezs}{RGB}{252,183,15}\fill[ezs](124.3,-206.5)--(126.6,-207.4)--(192,-37.2)--(189.7,-36.3)--cycle;
				\definecolor{ezs}{RGB}{252,181,14}\fill[ezs](126.3,-207.2)--(128.5,-208.1)--(193.9,-37.9)--(191.6,-37.1)--cycle;
				\definecolor{ezs}{RGB}{252,179,13}\fill[ezs](128.2,-208)--(130.4,-208.9)--(195.8,-38.7)--(193.5,-37.8)--cycle;
				\definecolor{ezs}{RGB}{252,178,12}\fill[ezs](130.1,-208.7)--(132.4,-209.6)--(197.7,-39.4)--(195.4,-38.5)--cycle;
				\definecolor{ezs}{RGB}{251,178,10}\fill[ezs](132,-209.4)--(134.3,-210.3)--(199.6,-40.1)--(197.3,-39.3)--cycle;
				\definecolor{ezs}{RGB}{251,178,10}\fill[ezs](133.9,-210.2)--(136.2,-211.1)--(201.5,-40.9)--(199.2,-40)--cycle;
				\definecolor{ezs}{RGB}{250,177,11}\fill[ezs](135.8,-210.9)--(138.1,-211.8)--(203.4,-41.6)--(201.1,-40.7)--cycle;
				\definecolor{ezs}{RGB}{250,178,11}\fill[ezs](137.7,-211.6)--(140,-212.5)--(205.3,-42.3)--(203.1,-41.5)--cycle;
				\definecolor{ezs}{RGB}{248,180,10}\fill[ezs](139.6,-212.4)--(141.9,-213.3)--(207.3,-43.1)--(205,-42.2)--cycle;
				\definecolor{ezs}{RGB}{247,181,9}\fill[ezs](141.5,-213.1)--(143.8,-214)--(209.2,-43.8)--(206.9,-42.9)--cycle;
				\definecolor{ezs}{RGB}{246,182,9}\fill[ezs](143.5,-213.8)--(145.7,-214.7)--(211.1,-44.5)--(208.8,-43.7)--cycle;
				\definecolor{ezs}{RGB}{246,182,10}\fill[ezs](145.4,-214.6)--(147.7,-215.5)--(213,-45.3)--(210.7,-44.4)--cycle;
			\end{scope}
			% chữ EDUZY (Times đậm, làm dày nét bằng 5 lớp lệch nhẹ; căn đúng khung chữ logo gốc)
			\foreach \dx/\dy in {-0.7/0,0.7/0,0/-0.3,0/0.3,0/0}
			\node[anchor=center,inner sep=0pt,outer sep=0pt,text=ezText,align=center] at (109.6+\dx,-193.6+\dy)
			{\resizebox{1.01cm}{0.22cm}{\fontencoding{T1}\fontfamily{ptm}\fontseries{b}\selectfont EDUZY}};
		\end{tikzpicture}%
	}%
}
%====================== THIẾT LẬP THÂN ĐỀ ======================
%--- Môi trường tự luận {bt}: ex_test 3.3.4 chưa khai báo sẵn nên tạo ở đây
%    (ex_test đã tự móc \fixshowans và \listenumerate cho bt)
\makeatletter
\@ifundefined{bt}{\newtheorem{bt}{\color{gold1}\namebt}}{}
\makeatother
\renewcommand{\sectionmark}[1]{\markright{\thesection\ #1}}
%--- Khung tiêu đề (logo + tối đa 4 dòng; dòng 3, 4 để trống {} thì bỏ qua)
%  Đề thi:    \khungtieude{THI THỬ THPT QUỐC GIA NĂM 2026}{Môn Toán -- Đề số 21}
%                         {(Đề thi gồm: ...)}{Thời gian làm bài: 90 phút (không kể thời gian phát đề).}
%  Chuyên đề: \khungtieude{CHUYÊN ĐỀ 3 -- DÃY SỐ VÀ CÔNG THỨC TRUY HỒI}
%                         {Bồi dưỡng học sinh giỏi Toán THPT}{...}{}
\newcommand{\khungtieude}[4]{%
	\begin{tcolorbox}[colframe=gold1,colback=gold5!45,boxrule=0.5mm]
		\begin{minipage}[c]{3cm}
			\centerline{\eduzylogo[2.4cm]}
		\end{minipage}\hspace{0.5cm}
		\begin{minipage}[c]{12.5cm}
			\vskip 0.15 cm
			\centerline{\ind{#1}}
			\vskip 0.15 cm
			\centerline{\textbf{#2}}
			\if\relax\detokenize{#3}\relax\else
				\vskip 0.15 cm
				\centerline{\textit{#3}}
			\fi
			\if\relax\detokenize{#4}\relax\else
				\vskip 0.15 cm
				\centerline{\textit{#4}}
			\fi
			\vskip 0.15 cm
		\end{minipage}
	\end{tcolorbox}
}
%====================== BỔ SUNG CHO BÀI HỌC, CHUYÊN ĐỀ ======================
\usepackage[shortlabels]{enumitem}
\usepackage{tabularx}
%--- Danh sách gọn: \begin{enumerate}[gon] ... \end{enumerate}
\SetEnumitemKey{gon}{leftmargin=*,itemsep=1pt,parsep=0pt,topsep=3pt}
%--- Khung lý thuyết: \begin{khung}{} ... \end{khung}; tiêu đề chỉ dùng khi thật ngắn ("Mẹo", "Lỗi thường gặp")
\newtcolorbox{khung}[1]{breakable,enhanced,colframe=gold1,colback=gold5!45,
	boxrule=0.4mm,fonttitle=\bfseries\sffamily,coltitle=white,colbacktitle=gold1,
	attach boxed title to top left={xshift=6mm,yshift=-3mm},
	boxed title style={arc=1mm,boxrule=0.3mm},title={#1}}
%--- Tiêu đề mục nhỏ trong một dạng: \muc{Dấu hiệu}, \muc{Phương pháp}, \muc{Bài tự luyện}
\makeatletter
\newcommand{\muc}[1]{\par\medskip\noindent\indam{#1}\par\nobreak\smallskip\@afterindentfalse\@afterheading}
\makeatother
%--- Ô công thức cần nhớ (căn giữa, viền vàng, vừa khít nội dung)
%     \begin{center}\congthuc{$Q(n)=(n-1)!$}\end{center}
\newtcbox{\congthuc}{on line,colframe=gold1,colback=gold5!45,boxrule=0.5pt,arc=2pt,
	left=8pt,right=8pt,top=4pt,bottom=4pt}
%--- Công tắc lời giải bài tự luyện (ví dụ mẫu luôn dùng \loigiai)
%     \loigiaiTL[{đáp số}]{lời giải đầy đủ}
%     \TLfull -> in lời giải đầy đủ;  \TLdapso -> chỉ in một dòng "Đáp số";  \TLan -> ẩn hoàn toàn
%     (\CheDo là dethi thì mọi lời giải và đáp số đều ẩn)
\makeatletter
\newif\ifTLhien   \TLhientrue
\newif\ifTLdapso  \TLdapsofalse
\newcommand{\TLfull}{\TLhientrue\TLdapsofalse}
\newcommand{\TLdapso}{\TLhientrue\TLdapsotrue}
\newcommand{\TLan}{\TLhienfalse}
\newcommand{\loigiaiTL}[2][]{%
	\ifTLhien
		\ifTLdapso
			\def\@tlarg{#1}%
			\ifx\@tlarg\@empty\else
				\ifdefstring{\kindOPTN}{dethi}{}{\par\noindent{\bfseries\color{gold1}Đáp số:} #1\par}%
			\fi
		\else
			\loigiai{#2}%
		\fi
	\fi}
\makeatother
\TLdapso% mặc định: chỉ in đáp số bài tự luyện
%====================== BẢNG ĐÚNG/SAI (\choiceTFt) ======================
%  Hàng tiêu đề nền nâu đồng chữ trắng (Phát biểu | Đ | S), các hàng tô xen kẽ kem / trắng,
%  đường kẻ vàng kim, nhãn a) b) c) d) màu nâu đồng. Bản có đáp án: cột được chọn có vòng tròn
%  nâu đồng với dấu tích trắng (câu đúng tích ở cột Đúng, câu sai tích ở cột Sai).
%  Tiêu đề ghi đủ "Đúng | Sai": \OPTN{viettat=0} sau \input (cột Đúng/Sai rộng ra, phát biểu hẹp lại).
%  Dùng bảng gốc của ex_test: \def\BangDSgoc{1} trước \input{preamble-eduzy-v2}.
\providecommand{\BangDSgoc}{0}
\ifnum\BangDSgoc=0
\usepackage{colortbl,pifont}
\colorlet{ezDSke}{gold2!60!gold1}% đường kẻ bảng (nét 0,4pt như ex_test để không đổi chiều cao bảng)
\colorlet{ezDSle}{gold5!60}% hàng a), c)
\colorlet{ezDSchan}{white}% hàng b), d)
%--- dấu chọn ở cột Đúng/Sai
%    (bề rộng 0, căn giữa ô: không làm đổi bề rộng cột của ex_test, không báo tràn lề)
\newcommand{\EZdauchon}{\makebox[0pt]{\tikz[baseline=-1ex]\node[circle,fill=gold1,inner sep=0pt,
	minimum size=1.2em,text=white,font=\small]{\ding{51}};}}
\def\tickT{\EZdauchon}\def\tickF{\EZdauchon}
\AtBeginDocument{\sbox{\ccTFTab}{\EZdauchon}}
%--- nhãn a) b) c) d) màu nâu đồng (ở bản lời giải, nhãn câu đúng vẫn đóng khung như ex_test)
\renewcommand{\FalseTF}{\stepcounter{dapan}\textbf{\color{gold1}\DapAnTF\sepTF}\ignorespaces}
%--- trước khi dựng bảng: kẻ màu vàng kim (bề rộng cột giữ nguyên như ex_test); sau bảng trả lại màu đen
\newcount\EZhangTF
\newcommand{\EZbatdauTF}{%
	\global\EZhangTF=0\relax
	\arrayrulecolor{ezDSke}}
\newcommand{\EZketthucTF}{\arrayrulecolor{black}}
%--- hàng tiêu đề nền nâu đồng, chữ trắng
\makeatletter
\newcommand{\EZo@TF}[1]{\centering\bfseries\color{white}#1}
\newcommand{\EZtieudeTF}{%
	\ifcase\cotDS
		\ifnum\cotSTT=1\relax
			\def\rowHeader{\hline\rowcolor{gold1}\EZo@TF{TT}&\EZo@TF{\phatbieu}\tabularnewline}%
		\else
			\def\rowHeader{\hline\rowcolor{gold1}\EZo@TF{\phatbieu}\tabularnewline}%
		\fi
	\or
		\ifnum\cotSTT=1\relax
			\def\rowHeader{\hline\rowcolor{gold1}\EZo@TF{TT}&\EZo@TF{\phatbieu}&\EZo@TF{\Dung}\tabularnewline}%
		\else
			\def\rowHeader{\hline\rowcolor{gold1}\EZo@TF{\phatbieu}&\EZo@TF{\Dung}\tabularnewline}%
		\fi
	\or
		\ifnum\cotSTT=1\relax
			\def\rowHeader{\hline\rowcolor{gold1}\EZo@TF{TT}&\EZo@TF{\phatbieu}&\EZo@TF{\Dung}&\EZo@TF{\Sai}\tabularnewline}%
		\else
			\def\rowHeader{\hline\rowcolor{gold1}\EZo@TF{\phatbieu}&\EZo@TF{\Dung}&\EZo@TF{\Sai}\tabularnewline}%
		\fi
	\fi}
%--- tô xen kẽ từng hàng (gọi ngay sau \hline đầu mỗi hàng)
\newcommand{\EZmauhangTF}{%
	\noalign{\global\advance\EZhangTF\@ne}%
	\ifodd\EZhangTF\rowcolor{ezDSle}\else\rowcolor{ezDSchan}\fi}
%--- gắn vào \choiceTFt / \writeChoiceTFt của ex_test
\xpretocmd{\choiceTFt}{\EZbatdauTF}{}{\typeout{[preamble-eduzy-v2] Không gắn được bảng Đúng/Sai (đầu bảng)}}
\xpatchcmd{\choiceTFt}{\ifnum\dongPB=0\def\rowHeader{}\fi}{\EZtieudeTF\ifnum\dongPB=0\def\rowHeader{}\fi}
	{}{\typeout{[preamble-eduzy-v2] Không gắn được bảng Đúng/Sai (hàng tiêu đề)}}
\xpatchcmd{\writeChoiceTFt}{\hline\xdef\dotEX{\sdotEX}}{\hline\EZmauhangTF\xdef\dotEX{\sdotEX}}
	{}{\typeout{[preamble-eduzy-v2] Không gắn được bảng Đúng/Sai (tô hàng)}}
\xpatchcmd{\writeChoiceTFt}{{\end{longtable}}}{{\end{longtable}\EZketthucTF}}
	{}{\typeout{[preamble-eduzy-v2] Không gắn được bảng Đúng/Sai (cuối longtable)}}
\xpatchcmd{\writeChoiceTFt}{{\end{xtabular}\end{center}\medskip\noindent}}{{\end{xtabular}\end{center}\EZketthucTF\medskip\noindent}}
	{}{\typeout{[preamble-eduzy-v2] Không gắn được bảng Đúng/Sai (cuối xtabular)}}
\fi
%=====================================================================
\renewcommand{\parallel}{\mathrel{/\!\!/}}
.
\providecommand{\BangDSgoc}{0}
\ifnum\BangDSgoc=0
\usepackage{colortbl,pifont}
\colorlet{ezDSke}{gold2!60!gold1}% đường kẻ bảng (nét 0,4pt như ex_test để không đổi chiều cao bảng)
\colorlet{ezDSle}{gold5!60}% hàng a), c)
\colorlet{ezDSchan}{white}% hàng b), d)
%--- dấu chọn ở cột Đúng/Sai
%    (bề rộng 0, căn giữa ô: không làm đổi bề rộng cột của ex_test, không báo tràn lề)
\newcommand{\EZdauchon}{\makebox[0pt]{\tikz[baseline=-1ex]\node[circle,fill=gold1,inner sep=0pt,
	minimum size=1.2em,text=white,font=\small]{\ding{51}};}}
\def\tickT{\EZdauchon}\def\tickF{\EZdauchon}
\AtBeginDocument{\sbox{\ccTFTab}{\EZdauchon}}
%--- nhãn a) b) c) d) màu nâu đồng (ở bản lời giải, nhãn câu đúng vẫn đóng khung như ex_test)
\renewcommand{\FalseTF}{\stepcounter{dapan}\textbf{\color{gold1}\DapAnTF\sepTF}\ignorespaces}
%--- trước khi dựng bảng: kẻ màu vàng kim (bề rộng cột giữ nguyên như ex_test); sau bảng trả lại màu đen
\newcount\EZhangTF
\newcommand{\EZbatdauTF}{%
	\global\EZhangTF=0\relax
	\arrayrulecolor{ezDSke}}
\newcommand{\EZketthucTF}{\arrayrulecolor{black}}
%--- hàng tiêu đề nền nâu đồng, chữ trắng
\makeatletter
\newcommand{\EZo@TF}[1]{\centering\bfseries\color{white}#1}
\newcommand{\EZtieudeTF}{%
	\ifcase\cotDS
		\ifnum\cotSTT=1\relax
			\def\rowHeader{\hline\rowcolor{gold1}\EZo@TF{TT}&\EZo@TF{\phatbieu}\tabularnewline}%
		\else
			\def\rowHeader{\hline\rowcolor{gold1}\EZo@TF{\phatbieu}\tabularnewline}%
		\fi
	\or
		\ifnum\cotSTT=1\relax
			\def\rowHeader{\hline\rowcolor{gold1}\EZo@TF{TT}&\EZo@TF{\phatbieu}&\EZo@TF{\Dung}\tabularnewline}%
		\else
			\def\rowHeader{\hline\rowcolor{gold1}\EZo@TF{\phatbieu}&\EZo@TF{\Dung}\tabularnewline}%
		\fi
	\or
		\ifnum\cotSTT=1\relax
			\def\rowHeader{\hline\rowcolor{gold1}\EZo@TF{TT}&\EZo@TF{\phatbieu}&\EZo@TF{\Dung}&\EZo@TF{\Sai}\tabularnewline}%
		\else
			\def\rowHeader{\hline\rowcolor{gold1}\EZo@TF{\phatbieu}&\EZo@TF{\Dung}&\EZo@TF{\Sai}\tabularnewline}%
		\fi
	\fi}
%--- tô xen kẽ từng hàng (gọi ngay sau \hline đầu mỗi hàng)
\newcommand{\EZmauhangTF}{%
	\noalign{\global\advance\EZhangTF\@ne}%
	\ifodd\EZhangTF\rowcolor{ezDSle}\else\rowcolor{ezDSchan}\fi}
%--- gắn vào \choiceTFt / \writeChoiceTFt của ex_test
\xpretocmd{\choiceTFt}{\EZbatdauTF}{}{\typeout{[preamble-eduzy-v2] Không gắn được bảng Đúng/Sai (đầu bảng)}}
\xpatchcmd{\choiceTFt}{\ifnum\dongPB=0\def\rowHeader{}\fi}{\EZtieudeTF\ifnum\dongPB=0\def\rowHeader{}\fi}
	{}{\typeout{[preamble-eduzy-v2] Không gắn được bảng Đúng/Sai (hàng tiêu đề)}}
\xpatchcmd{\writeChoiceTFt}{\hline\xdef\dotEX{\sdotEX}}{\hline\EZmauhangTF\xdef\dotEX{\sdotEX}}
	{}{\typeout{[preamble-eduzy-v2] Không gắn được bảng Đúng/Sai (tô hàng)}}
\xpatchcmd{\writeChoiceTFt}{{\end{longtable}}}{{\end{longtable}\EZketthucTF}}
	{}{\typeout{[preamble-eduzy-v2] Không gắn được bảng Đúng/Sai (cuối longtable)}}
\xpatchcmd{\writeChoiceTFt}{{\end{xtabular}\end{center}\medskip\noindent}}{{\end{xtabular}\end{center}\EZketthucTF\medskip\noindent}}
	{}{\typeout{[preamble-eduzy-v2] Không gắn được bảng Đúng/Sai (cuối xtabular)}}
\fi
%=====================================================================
\renewcommand{\parallel}{\mathrel{/\!\!/}}
.
\providecommand{\BangDSgoc}{0}
\ifnum\BangDSgoc=0
\usepackage{colortbl,pifont}
\colorlet{ezDSke}{gold2!60!gold1}% đường kẻ bảng (nét 0,4pt như ex_test để không đổi chiều cao bảng)
\colorlet{ezDSle}{gold5!60}% hàng a), c)
\colorlet{ezDSchan}{white}% hàng b), d)
%--- dấu chọn ở cột Đúng/Sai
%    (bề rộng 0, căn giữa ô: không làm đổi bề rộng cột của ex_test, không báo tràn lề)
\newcommand{\EZdauchon}{\makebox[0pt]{\tikz[baseline=-1ex]\node[circle,fill=gold1,inner sep=0pt,
	minimum size=1.2em,text=white,font=\small]{\ding{51}};}}
\def\tickT{\EZdauchon}\def\tickF{\EZdauchon}
\AtBeginDocument{\sbox{\ccTFTab}{\EZdauchon}}
%--- nhãn a) b) c) d) màu nâu đồng (ở bản lời giải, nhãn câu đúng vẫn đóng khung như ex_test)
\renewcommand{\FalseTF}{\stepcounter{dapan}\textbf{\color{gold1}\DapAnTF\sepTF}\ignorespaces}
%--- trước khi dựng bảng: kẻ màu vàng kim (bề rộng cột giữ nguyên như ex_test); sau bảng trả lại màu đen
\newcount\EZhangTF
\newcommand{\EZbatdauTF}{%
	\global\EZhangTF=0\relax
	\arrayrulecolor{ezDSke}}
\newcommand{\EZketthucTF}{\arrayrulecolor{black}}
%--- hàng tiêu đề nền nâu đồng, chữ trắng
\makeatletter
\newcommand{\EZo@TF}[1]{\centering\bfseries\color{white}#1}
\newcommand{\EZtieudeTF}{%
	\ifcase\cotDS
		\ifnum\cotSTT=1\relax
			\def\rowHeader{\hline\rowcolor{gold1}\EZo@TF{TT}&\EZo@TF{\phatbieu}\tabularnewline}%
		\else
			\def\rowHeader{\hline\rowcolor{gold1}\EZo@TF{\phatbieu}\tabularnewline}%
		\fi
	\or
		\ifnum\cotSTT=1\relax
			\def\rowHeader{\hline\rowcolor{gold1}\EZo@TF{TT}&\EZo@TF{\phatbieu}&\EZo@TF{\Dung}\tabularnewline}%
		\else
			\def\rowHeader{\hline\rowcolor{gold1}\EZo@TF{\phatbieu}&\EZo@TF{\Dung}\tabularnewline}%
		\fi
	\or
		\ifnum\cotSTT=1\relax
			\def\rowHeader{\hline\rowcolor{gold1}\EZo@TF{TT}&\EZo@TF{\phatbieu}&\EZo@TF{\Dung}&\EZo@TF{\Sai}\tabularnewline}%
		\else
			\def\rowHeader{\hline\rowcolor{gold1}\EZo@TF{\phatbieu}&\EZo@TF{\Dung}&\EZo@TF{\Sai}\tabularnewline}%
		\fi
	\fi}
%--- tô xen kẽ từng hàng (gọi ngay sau \hline đầu mỗi hàng)
\newcommand{\EZmauhangTF}{%
	\noalign{\global\advance\EZhangTF\@ne}%
	\ifodd\EZhangTF\rowcolor{ezDSle}\else\rowcolor{ezDSchan}\fi}
%--- gắn vào \choiceTFt / \writeChoiceTFt của ex_test
\xpretocmd{\choiceTFt}{\EZbatdauTF}{}{\typeout{[preamble-eduzy-v2] Không gắn được bảng Đúng/Sai (đầu bảng)}}
\xpatchcmd{\choiceTFt}{\ifnum\dongPB=0\def\rowHeader{}\fi}{\EZtieudeTF\ifnum\dongPB=0\def\rowHeader{}\fi}
	{}{\typeout{[preamble-eduzy-v2] Không gắn được bảng Đúng/Sai (hàng tiêu đề)}}
\xpatchcmd{\writeChoiceTFt}{\hline\xdef\dotEX{\sdotEX}}{\hline\EZmauhangTF\xdef\dotEX{\sdotEX}}
	{}{\typeout{[preamble-eduzy-v2] Không gắn được bảng Đúng/Sai (tô hàng)}}
\xpatchcmd{\writeChoiceTFt}{{\end{longtable}}}{{\end{longtable}\EZketthucTF}}
	{}{\typeout{[preamble-eduzy-v2] Không gắn được bảng Đúng/Sai (cuối longtable)}}
\xpatchcmd{\writeChoiceTFt}{{\end{xtabular}\end{center}\medskip\noindent}}{{\end{xtabular}\end{center}\EZketthucTF\medskip\noindent}}
	{}{\typeout{[preamble-eduzy-v2] Không gắn được bảng Đúng/Sai (cuối xtabular)}}
\fi
%=====================================================================
\renewcommand{\parallel}{\mathrel{/\!\!/}}

%=====================================================================
% preamble-vatli-hinh.tex -- BỘ MACRO HÌNH VẬT LÍ (dùng chung cho mọi đề Vật lí)
% Gọi ngay sau preamble chính:
%     %============================================================
%  PREAMBLE EDUZY v2 - dùng chung cho đề thi, bài học, chuyên đề (gói ex_test)
%  Khác bản preamble-eduzy.tex ở đầu trang, chân trang: dải nền nâu đồng tràn ngang,
%  logo + Eduzy.com.vn, ô chữ vàng kem cạnh xiên, chân trang "Trang (số)";
%  logo Eduzy chìm giữa mọi trang (tắt: \logochimtat, bật lại: \logochimbat).
%  Cách dùng ở file chính (xem mau-chuan.tex, mau-chuyen-de.tex):
%      \def\CheDo{loigiai}          % loigiai | dethi | solcolor | color | book
%      \def\tenchuong{<dòng chữ trên đầu trang>}
%      \def\thuonghieu{<chữ ở chân trang lẻ>}   % không khai báo thì dùng mặc định
%      \def\tacgia{<chữ ở chân trang chẵn>}      % không khai báo thì dùng mặc định
%      \def\khauhieu{<chữ cạnh logo ở đầu trang>}  % tùy chọn, mặc định Eduzy.com.vn
%      %============================================================
%  PREAMBLE EDUZY v2 - dùng chung cho đề thi, bài học, chuyên đề (gói ex_test)
%  Khác bản preamble-eduzy.tex ở đầu trang, chân trang: dải nền nâu đồng tràn ngang,
%  logo + Eduzy.com.vn, ô chữ vàng kem cạnh xiên, chân trang "Trang (số)";
%  logo Eduzy chìm giữa mọi trang (tắt: \logochimtat, bật lại: \logochimbat).
%  Cách dùng ở file chính (xem mau-chuan.tex, mau-chuyen-de.tex):
%      \def\CheDo{loigiai}          % loigiai | dethi | solcolor | color | book
%      \def\tenchuong{<dòng chữ trên đầu trang>}
%      \def\thuonghieu{<chữ ở chân trang lẻ>}   % không khai báo thì dùng mặc định
%      \def\tacgia{<chữ ở chân trang chẵn>}      % không khai báo thì dùng mặc định
%      \def\khauhieu{<chữ cạnh logo ở đầu trang>}  % tùy chọn, mặc định Eduzy.com.vn
%      \input{preamble-eduzy-v2}
%      \TLdapso                      % tài liệu: \TLfull | \TLdapso | \TLan
%      \begin{document} ... \end{document}
%  Logo Eduzy vẽ bằng TikZ ngay trong file (\eduzylogo), không cần file ảnh.
%  Biên dịch: pdfLaTeX (máy có gói vietnam) hoặc XeLaTeX.
%============================================================
\documentclass[12pt,a4paper,twoside]{article}%<-- twoside: bat dau trang/chan trang chan-le
\usepackage{iftex}
\usepackage{amsmath,amssymb,tikz,mathrsfs,tkz-tab}
\usepackage{graphicx}
\usepackage{pgfplots}% đồ thị hàm số
\pgfplotsset{compat=1.17}
%\usepackage{mathptmx}
\usepackage{mathpazo}
\ifXeTeX% XeLaTeX: font chữ có đủ dấu tiếng Việt, giữ nguyên font toán mathpazo
	\usepackage[no-math]{fontspec}
	\setmainfont{TeX Gyre Pagella}
	\newfontfamily\EZqag{TeX Gyre Adventor}[NFSSFamily=qag,Ligatures=TeX]
\fi
\usepackage[many]{tcolorbox}
\usetikzlibrary{math,through,calc,intersections,angles,quotes,shapes,shapes.geometric,arrows,patterns,matrix,chains,arrows.meta,decorations.shapes,decorations.fractals,decorations.markings,shadows,decorations.pathreplacing}
\usepackage[top=1.6cm, bottom=1.8cm, left=1.8cm, right=1.8cm]{geometry}% v2: top 1.6cm (bản cũ 1.5cm) để chữ không sát dải đầu trang
\providecommand{\CheDo}{loigiai}% loigiai | dethi | solcolor | color | book
\expandafter\usepackage\expandafter[\CheDo]{ex_test}
% \tenchuong: khai báo ở file chính (trước \input) hoặc \renewcommand sau \input
%============================================================
%  [GỘP TỪ eduzy-de.tex] EDUZY - MÀU NÂU ĐỒNG / VÀNG KIM
%  + ĐẦU TRANG, CHÂN TRANG KIỂU DẢI NGANG (v2), ĐỐI XỨNG TRANG CHẴN / TRANG LẺ.
%  (không cần file eduzy-de.tex riêng nữa)
%  Khối này phải đặt SAU \usepackage[...]{ex_test} và SAU \def\tenchuong{...}
%============================================================
\usepackage{lastpage}
\usepackage{needspace}
%--- DÀN TRANG ĐỀ THI: tách câu (THÊM 05/11/2026; chỉ thêm lệnh mới, KHÔNG đổi hành vi mặc định) ---
%  Hai công tắc, mặc định tắt (= preamble chạy y như cũ). Bật bằng cách đặt TRƯỚC \input{preamble-eduzy-v2}
%  trong file đề (tools/dan_trang/dan_trang.py tự ghi khi đặt \NeedspaceCau):
%    \def\TachCauDS{1}   câu Đúng/Sai được tách: câu dẫn (kèm hình) ở cuối trang, bảng Phát biểu sang trang sau
%                        (câu dẫn và hình không tách nhau, bảng không tách giữa hai trang)
%    \def\TachCauLG{1}   bản lời giải (\CheDo khác dethi): phần đề của câu cao hơn 1/3 trang được phép tách
%  Lệnh dùng trong file đề (giá trị do tools/dan_trang/dan_trang.py đo ở bản dethi):
%    \NeedspaceCau[<cao câu dẫn + hình>]{<cao cả khối câu>}   ngay trên \begin{ex} (câu đầu phần: trên tiêu đề phần)
%        công tắc tắt: đúng bằng \Needspace{<cao cả khối câu>} như cũ;
%        \TachCauDS=1 và có [..] (câu Đúng/Sai được chọn tách): chỉ giữ câu dẫn + hình;
%        \TachCauLG=1 ở bản lời giải: đòi tối đa 1/3 \textheight.
%    \NeedspaceBang{<cao bảng Phát biểu>}%   ngay trước \choiceTFt: \TachCauDS=1 thì giữ trọn bảng trên một trang;
%        công tắc tắt: không làm gì.
\providecommand{\TachCauDS}{0}
\providecommand{\TachCauLG}{0}
\newlength{\EZnsCau}
\newcommand{\NeedspaceCau}[2][]{%
	\setlength{\EZnsCau}{#2}%
	\ifnum\TachCauDS=1 \if\relax\detokenize{#1}\relax\else\setlength{\EZnsCau}{#1}\fi\fi
	\ifnum\TachCauLG=1 \ifdefstring{\kindOPTN}{dethi}{}{\ifdim\EZnsCau>0.3333\textheight\setlength{\EZnsCau}{0.3333\textheight}\fi}\fi
	\Needspace{\EZnsCau}}
\newcommand{\NeedspaceBang}[1]{\ifnum\TachCauDS=1 \Needspace{#1}\fi}

%--- Thang màu: đậm -> nhạt
\definecolor{gold1}{RGB}{126,96,0}    % nâu đồng đậm  #7E6000
\definecolor{gold2}{RGB}{192,144,3}   % vàng kim      #C09003
\definecolor{gold3}{RGB}{255,215,99}  % vàng nhạt     #FFD763
\definecolor{gold4}{RGB}{255,231,153} % vàng rất nhạt #FFE799
\definecolor{gold5}{RGB}{255,243,205} % vàng kem      #FFF3CD
%--- Ghi đè tông xanh/đỏ cũ
\definecolor{myblue}{RGB}{126,96,0}
\colorlet{blue}{gold1}
\colorlet{cyan}{gold2}
\colorlet{teal}{gold2}

%--- Tên chương mặc định (file đề nên tự khai báo lại)
\providecommand{\tenchuong}{}
%--- Chân trang
\providecommand{\thuonghieu}{TOÁN 12 -- Eduzy.com.vn}
\providecommand{\tacgia}{Tô Gia Bảo -- 0353703783}

%--- Chữ cạnh logo ở đầu trang
\providecommand{\khauhieu}{Eduzy.com.vn}

%====================== THÔNG SỐ ĐẦU TRANG, CHÂN TRANG (v2) ======================
%  Dải nền đậm tràn ngang trang; ô chữ vàng kem có cạnh xiên, kèm 2 vạch xiên chuyển màu.
%  Đầu trang lẻ : logo + \khauhieu (trái)   | ô chữ \tenchuong (phải)
%  Chân trang lẻ: ô chữ \thuonghieu (trái)  | Trang (số) (phải)
%  Trang chẵn soi gương trang lẻ; chân trang chẵn ghi \tacgia.
%--- Màu: chỉ dùng thang vàng kim Eduzy
\colorlet{ezNen}{gold1}% dải nền đậm
\colorlet{ezVach1}{gold3}% vạch xiên sát dải nền (đậm hơn)
\colorlet{ezVach2}{gold4}% vạch xiên sát ô chữ (nhạt hơn)
\colorlet{ezO}{gold5}% ô chữ vàng kem
%--- Kích thước
\newlength{\EZdai}\setlength{\EZdai}{1.1cm}% chiều cao dải nền
\newlength{\EZo}\setlength{\EZo}{0.68cm}% chiều cao ô chữ
\newlength{\EZxien}\setlength{\EZxien}{0.55cm}% độ lệch ngang của cạnh xiên
\newlength{\EZvach}\setlength{\EZvach}{0.75cm}% bề rộng mỗi vạch xiên
\newlength{\EZkhe}\setlength{\EZkhe}{0.12cm}% khe trắng giữa vạch và ô chữ
\newlength{\EZdem}\setlength{\EZdem}{0.4cm}% khoảng đệm từ chữ tới cạnh xiên
\newlength{\EZle}\setlength{\EZle}{1.8cm}% lề trái/phải của chữ (bằng lề thân bài)
\newlength{\EZrongmin}\setlength{\EZrongmin}{4.5cm}% ô chữ rộng tối thiểu
\newlength{\EZrong}% bề rộng chữ trong ô (tính tự động)
%--- Font
\newcommand{\EZfontO}{\fontsize{10.5pt}{12pt}\selectfont\bfseries}% chữ trong ô
\newcommand{\EZfontKH}{\fontfamily{qag}\fontsize{11pt}{12pt}\selectfont\bfseries}% \khauhieu cạnh logo
\newcommand{\EZfontTrang}{\fontsize{11.5pt}{12pt}\selectfont\bfseries}% chữ "Trang" và số trang
%--- Đo bề rộng chữ trong ô
\newcommand{\EZdorong}[1]{%
	\settowidth{\EZrong}{\EZfontO #1}%
	\ifdim\EZrong<\EZrongmin\setlength{\EZrong}{\EZrongmin}\fi}
%--- Logo + chữ trắng \khauhieu (dựng sẵn một lần khi vào \begin{document})
\newsavebox{\EZlogobox}
\AtBeginDocument{\sbox{\EZlogobox}{%
		\raisebox{-0.5\height}{\eduzylogo[1cm]}\hspace{0.15cm}%
		\raisebox{-0.5\height}{\color{white}\EZfontKH\khauhieu}}}
%--- Hai vạch xiên + khe trắng. #1: điểm gốc ở mép dưới, #2: hướng (1 sang phải, -1 sang trái),
%    #3: hướng nghiêng (1 kiểu "/", -1 kiểu "\")
\newcommand{\EZvachxien}[3]{%
	\fill[white] (#1) --++({(#3)*\EZxien},\EZo) --++({(#2)*\EZkhe},0) --++({-(#3)*\EZxien},-\EZo) --cycle;
	\fill[ezVach2] ([xshift=(#2)*\EZkhe]#1) --++({(#3)*\EZxien},\EZo) --++({(#2)*\EZvach},0) --++({-(#3)*\EZxien},-\EZo) --cycle;
	\fill[ezVach1] ([xshift=(#2)*(\EZkhe+\EZvach)]#1) --++({(#3)*\EZxien},\EZo) --++({(#2)*\EZvach},0) --++({-(#3)*\EZxien},-\EZo) --cycle;
	\draw[white,line width=0.6pt] ([xshift=(#2)*(\EZkhe+\EZvach)]#1) --++({(#3)*\EZxien},\EZo);
}
%--- Chữ trong ô (căn giữa theo chiều cao ô, không lệch vì dấu tiếng Việt)
\tikzset{ezchu/.style={inner sep=0pt,text=ezNen,text height=0.72em,text depth=0.22em}}

%====================== LOGO CHÌM GIỮA TRANG ======================
%  Vẽ cùng đầu trang nên nằm dưới chữ của thân bài, có trên mọi trang.
%  Tắt: \logochimtat (từ trang đó trở đi), bật lại: \logochimbat.
%  Đổi cỡ, độ đậm: khai báo trước \input, ví dụ \def\EZchimcao{9cm} \def\EZchimdo{0.1}
\providecommand{\EZchimcao}{10cm}% chiều cao logo chìm
\providecommand{\EZchimdo}{0.12}% độ đậm (0 = trong suốt, 1 = rõ hẳn)
\newif\ifEZchim \EZchimtrue
\newcommand{\logochimtat}{\EZchimfalse}
\newcommand{\logochimbat}{\EZchimtrue}
%  Dựng sẵn logo mờ một lần khi vào \begin{document}; transparency group làm mờ cả logo
%  một lượt nên không lộ vệt ở chỗ các dải màu chồng nhau.
\newsavebox{\EZchimtho}
\newsavebox{\EZchimbox}
\AtBeginDocument{%
	\sbox{\EZchimtho}{\eduzylogo[\EZchimcao]}%
	\sbox{\EZchimbox}{\begin{tikzpicture}
			\begin{scope}[transparency group,opacity=\EZchimdo]
				\node[inner sep=0pt] {\usebox{\EZchimtho}};
			\end{scope}
	\end{tikzpicture}}}
\newcommand{\EZvechim}{%
	\ifEZchim
		\node[inner sep=0pt] at (current page.center) {\usebox{\EZchimbox}};
	\fi}

%====================== ĐẦU TRANG - TRANG LẺ ======================
%  (logo bên trái, ô chữ \tenchuong bên phải)
\newcommand{\myfancyheadOdd}{%
	\EZdorong{\tenchuong}%
	\begin{tikzpicture}[remember picture,overlay]
		\coordinate (T) at (current page.north west);
		\coordinate (P) at (current page.north east);
		%--- logo chìm giữa trang
		\EZvechim
		%--- dải nền
		\fill[ezNen] (T) rectangle ([yshift=-\EZdai]P);
		%--- ô chữ: X là góc dưới bên trái, cạnh trái xiên kiểu "/"
		\coordinate (X) at ([xshift=-\EZle-\EZrong-\EZdem-\EZxien,yshift=-\EZdai]P);
		\fill[ezO] (X) --++(\EZxien,\EZo) -- ([yshift=-\EZdai+\EZo]P) -- ([yshift=-\EZdai]P) --cycle;
		\EZvachxien{X}{-1}{1}
		\node[ezchu,anchor=east] at ([xshift=-\EZle,yshift=-\EZdai+0.5*\EZo]P) {\EZfontO\tenchuong};
		%--- logo + tên
		\node[inner sep=0pt,anchor=west] at ([xshift=\EZle,yshift=-0.5*\EZdai]T) {\usebox{\EZlogobox}};
	\end{tikzpicture}%
}

%====================== ĐẦU TRANG - TRANG CHẴN =====================
%  (soi gương: ô chữ \tenchuong bên trái, logo bên phải)
\newcommand{\myfancyheadEven}{%
	\EZdorong{\tenchuong}%
	\begin{tikzpicture}[remember picture,overlay]
		\coordinate (T) at (current page.north west);
		\coordinate (P) at (current page.north east);
		%--- logo chìm giữa trang
		\EZvechim
		%--- dải nền
		\fill[ezNen] (T) rectangle ([yshift=-\EZdai]P);
		%--- ô chữ: X là góc dưới bên phải, cạnh phải xiên kiểu "\"
		\coordinate (X) at ([xshift=\EZle+\EZrong+\EZdem+\EZxien,yshift=-\EZdai]T);
		\fill[ezO] (X) --++(-\EZxien,\EZo) -- ([yshift=-\EZdai+\EZo]T) -- ([yshift=-\EZdai]T) --cycle;
		\EZvachxien{X}{1}{-1}
		\node[ezchu,anchor=west] at ([xshift=\EZle,yshift=-\EZdai+0.5*\EZo]T) {\EZfontO\tenchuong};
		%--- logo + tên
		\node[inner sep=0pt,anchor=east] at ([xshift=-\EZle,yshift=-0.5*\EZdai]P) {\usebox{\EZlogobox}};
	\end{tikzpicture}%
}

%====================== CHÂN TRANG - TRANG LẺ ======================
%  (ô chữ \thuonghieu bên trái, "Trang (số)" bên phải)
\newcommand{\myfancyfootOdd}{%
	\EZdorong{\thuonghieu}%
	\begin{tikzpicture}[remember picture,overlay]
		\coordinate (T) at (current page.south west);
		\coordinate (P) at (current page.south east);
		%--- dải nền
		\fill[ezNen] (T) rectangle ([yshift=\EZdai]P);
		%--- ô chữ sát mép trên dải: X là góc dưới bên phải, cạnh phải xiên kiểu "/"
		\coordinate (X) at ([xshift=\EZle+\EZrong+\EZdem,yshift=\EZdai-\EZo]T);
		\fill[ezO] ([yshift=\EZdai-\EZo]T) -- (X) --++(\EZxien,\EZo) -- ([yshift=\EZdai]T) --cycle;
		\EZvachxien{X}{1}{1}
		\node[ezchu,anchor=west] at ([xshift=\EZle,yshift=\EZdai-0.5*\EZo]T) {\EZfontO\thuonghieu};
		%--- số trang
		\node[circle,draw=white,line width=0.9pt,text=white,inner sep=0pt,minimum size=0.62cm,anchor=east]
		(so) at ([xshift=-\EZle,yshift=0.5*\EZdai]P) {\EZfontTrang\thepage};
		\node[inner sep=0pt,text=white,anchor=base east] (tr) at ([xshift=-0.15cm]so.base west) {\EZfontTrang Trang};
	\end{tikzpicture}%
}

%====================== CHÂN TRANG - TRANG CHẴN ====================
%  (soi gương: "Trang (số)" bên trái, ô chữ \tacgia bên phải)
\newcommand{\myfancyfootEven}{%
	\EZdorong{\tacgia}%
	\begin{tikzpicture}[remember picture,overlay]
		\coordinate (T) at (current page.south west);
		\coordinate (P) at (current page.south east);
		%--- dải nền
		\fill[ezNen] (T) rectangle ([yshift=\EZdai]P);
		%--- ô chữ sát mép trên dải: X là góc dưới bên trái, cạnh trái xiên kiểu "\"
		\coordinate (X) at ([xshift=-\EZle-\EZrong-\EZdem,yshift=\EZdai-\EZo]P);
		\fill[ezO] ([yshift=\EZdai-\EZo]P) -- (X) --++(-\EZxien,\EZo) -- ([yshift=\EZdai]P) --cycle;
		\EZvachxien{X}{-1}{-1}
		\node[ezchu,anchor=east] at ([xshift=-\EZle,yshift=\EZdai-0.5*\EZo]P) {\EZfontO\tacgia};
		%--- số trang
		%    (vòng tròn ẩn so0 nằm giữa dải để lấy đường chân chữ, giống trang lẻ)
		\node[circle,inner sep=0pt,minimum size=0.62cm] (so0) at ([yshift=0.5*\EZdai]T) {\EZfontTrang\phantom{\thepage}};
		\node[inner sep=0pt,text=white,anchor=base west] (tr) at ($(so0.base -| T)+(\EZle,0)$) {\EZfontTrang Trang};
		\node[circle,draw=white,line width=0.9pt,text=white,inner sep=0pt,minimum size=0.62cm,anchor=base west]
		(so) at ([xshift=0.15cm]tr.base east) {\EZfontTrang\thepage};
	\end{tikzpicture}%
}

%--- Giữ tương thích với các file cũ đang gọi \myfancyhead, \myfancyfoot
\newcommand{\myfancyhead}{\myfancyheadOdd}
\newcommand{\myfancyfoot}{\myfancyfootOdd}

%======================= GẮN VÀO TRANG =============================
\pagestyle{fancy}
\fancyhf{}
\fancyhead[LO]{\myfancyheadOdd}
\fancyhead[RE]{\myfancyheadEven}
\fancyfoot[LO]{\myfancyfootOdd}
\fancyfoot[RE]{\myfancyfootEven}
\renewcommand{\headrulewidth}{0pt}
\renewcommand{\footrulewidth}{0pt}
\fancypagestyle{plain}{%
	\fancyhf{}
	\fancyhead[LO]{\myfancyheadOdd}
	\fancyhead[RE]{\myfancyheadEven}
	\fancyfoot[LO]{\myfancyfootOdd}
	\fancyfoot[RE]{\myfancyfootEven}
	\renewcommand{\headrulewidth}{0pt}
	\renewcommand{\footrulewidth}{0pt}
}
%=================== HẾT PHẦN GỘP TỪ eduzy-de.tex ===================
\newcommand{\indam}[1]{
	{\fontsize{12pt}{9pt}\sffamily \selectfont\bfseries\color{gold1} #1}
}
\newcommand{\ind}[1]{
	{\fontsize{14pt}{9pt}\fontfamily{qag} \selectfont\bfseries\color{gold1} #1}
}
\renewcommand{\FalseEX}{\stepcounter{dapan}\textbf{\selectfont\color{gold1}\Alph{dapan}.}\ignorespaces}
\renewtheorem{ex}{\color{gold1}Câu}
\renewcommand{\baselinestretch}{1.1}
\newcommand{\hoac}[1]{
	\left[\begin{aligned}#1\end{aligned}\right.}
\newcommand{\heva}[1]{
	\left\{\begin{aligned}#1\end{aligned}\right.}
\usepackage{esvect}
\def\vec{\vv}
\def\vt{\vv}
\def\overrightarrow{\vv}
%====================== LOGO EDUZY VẼ BẰNG TIKZ ======================
%  Đường viền dò từ logo gốc; dải màu vàng kim đo trực tiếp trên ảnh gốc
%  (nghiêng 21 độ, tô bằng 80 dải hẹp nên hiển thị giống nhau trên mọi trình xem PDF).
%  Dùng: \eduzylogo          -> cao 2.4cm (mặc định)
%        \eduzylogo[1.5cm]   -> cao 1.5cm
\definecolor{ezText}{RGB}{245,200,46}
\newcommand{\eduzylogo}[1][2.4cm]{%
	\resizebox{!}{#1}{%
		\begin{tikzpicture}[x=0.01cm,y=0.01cm]
			\path[use as bounding box] (0,0) rectangle (224,-224);
			% mũ tốt nghiệp + vương miện
			\begin{scope}
				\clip
				(119.4,-176.1)..controls(118.9,-176.1)and(116.8,-176)..(114.6,-176)
				..controls(101.4,-176)and(85.8,-175.9)..(83.7,-175.7)
				..controls(81.8,-175.6)and(81,-175.6)..(80.2,-175.7)
				..controls(78.8,-175.9)and(77.7,-175.9)..(76.2,-175.7)
				..controls(74.9,-175.5)and(70.8,-175.5)..(64.3,-175.7)
				..controls(59.1,-175.9)and(58.3,-175.7)..(58.6,-173.7)
				..controls(58.6,-172.8)and(59.6,-169.6)..(60.3,-168)
				..controls(60.8,-166.9)and(61.4,-165.8)..(62.3,-165)
				..controls(63.1,-164.2)and(64.3,-163.4)..(65,-163.2)
				..controls(65.5,-163.1)and(72.3,-163)..(76.6,-163.1)
				..controls(77.9,-163.1)and(80.1,-163.1)..(81.4,-163)
				..controls(82.8,-162.9)and(86.3,-162.9)..(89.2,-162.9)
				..controls(93,-162.9)and(94.9,-162.8)..(96,-162.7)
				..controls(97.2,-162.5)and(101,-162.5)..(114,-162.4)
				..controls(123,-162.3)and(134.9,-162.2)..(140.2,-162.1)
				..controls(150.1,-162)and(152.2,-162.1)..(153.6,-162.4)
				..controls(154.4,-162.6)and(157.1,-163.8)..(158.1,-164.4)
				..controls(158.8,-164.8)and(159.7,-165.7)..(159.9,-166.3)
				..controls(160,-166.5)and(160.2,-167.1)..(160.2,-167.6)
				..controls(160.3,-168)and(160.6,-169.4)..(160.9,-170.6)
				..controls(161.2,-171.7)and(161.5,-173.1)..(161.6,-173.6)
				..controls(161.8,-174.4)and(161.8,-174.5)..(161.5,-174.9)
				..controls(161.2,-175.6)and(160.8,-175.9)..(159.9,-176)
				..controls(159.2,-176.1)and(154.1,-176.2)..(151.5,-176.2)
				..controls(148.1,-176.1)and(122.6,-176.1)..(121.6,-176.2)
				..controls(120.8,-176.2)and(119.8,-176.2)..(119.4,-176.1)--cycle
				(130.4,-160)..controls(129,-160)and(119.8,-160)..(110.1,-159.9)
				..controls(100.3,-159.8)and(88.5,-159.7)..(83.9,-159.6)
				..controls(77.8,-159.5)and(75.3,-159.5)..(74.5,-159.6)
				..controls(73.9,-159.7)and(72.5,-159.8)..(71.3,-159.8)
				..controls(69.6,-159.8)and(69,-159.8)..(68.5,-159.6)
				..controls(67.6,-159.3)and(66.8,-158.6)..(66.2,-157.6)
				..controls(66,-157.3)and(65.6,-156.5)..(65.3,-156.1)
				..controls(64.6,-155.1)and(64.5,-154.9)..(63.7,-152.9)
				..controls(63.2,-151.6)and(62.9,-150.9)..(61.8,-149.6)
				..controls(61.5,-149.2)and(61.1,-148.5)..(60.9,-147.9)
				..controls(60.7,-147.4)and(60.3,-146.6)..(60.1,-146.1)
				..controls(59.9,-145.5)and(59.5,-144.9)..(59,-144.3)
				..controls(58.5,-143.7)and(58.2,-143.2)..(57.9,-142.5)
				..controls(57.7,-141.9)and(57.3,-141.1)..(57,-140.7)
				..controls(56.6,-140.3)and(56.1,-139.3)..(55.8,-138.5)
				..controls(55.5,-137.8)and(55.1,-136.8)..(54.9,-136.5)
				..controls(54.7,-136.1)and(54.3,-135.2)..(53.9,-134.5)
				..controls(53.6,-133.8)and(53,-132.7)..(52.6,-132)
				..controls(51.6,-130.5)and(51.4,-130.2)..(51.1,-129)
				..controls(50.6,-127.5)and(50.6,-127.4)..(49.7,-126.3)
				..controls(49.2,-125.7)and(48.6,-124.8)..(48.3,-124.2)
				..controls(48,-123.6)and(47.5,-122.7)..(47.3,-122.2)
				..controls(47,-121.8)and(46.5,-120.7)..(46.1,-119.8)
				..controls(45.8,-118.9)and(45.2,-117.8)..(44.9,-117.2)
				..controls(44.2,-116.1)and(43.3,-114.4)..(42.3,-112.2)
				..controls(41.8,-111.3)and(41.3,-110.4)..(41.1,-110.1)
				..controls(40.5,-109.4)and(40.3,-109)..(40,-107.8)
				..controls(39.8,-107)and(39.6,-106.5)..(39.2,-106)
				..controls(38.9,-105.6)and(38.4,-104.7)..(38,-104)
				..controls(37.7,-103.4)and(37.2,-102.5)..(37,-102.1)
				..controls(36.7,-101.7)and(36.3,-100.8)..(36,-100.1)
				..controls(35.6,-98.9)and(35.3,-98.4)..(34.6,-97.6)
				..controls(34.2,-97.1)and(34.2,-96.3)..(34.6,-96)
				..controls(35,-95.7)and(35.5,-95.9)..(36.2,-96.7)
				..controls(36.5,-97)and(37.2,-97.5)..(37.6,-97.7)..controls(38,-98)and(38.6,-98.4)..(39.1,-98.9)
				..controls(39.6,-99.4)and(40.3,-100)..(40.6,-100.3)
				..controls(41.3,-100.8)and(42.4,-101.8)..(43.7,-102.9)
				..controls(44,-103.2)and(44.6,-103.7)..(45,-103.9)
				..controls(45.4,-104.2)and(45.9,-104.6)..(46.1,-104.8)
				..controls(47.1,-105.9)and(48.8,-107.1)..(50.5,-107.9)
				..controls(51.8,-108.6)and(52.3,-109)..(53.4,-110)
				..controls(54.8,-111.2)and(55.6,-111.6)..(57,-112.1)
				..controls(57.5,-112.3)and(58.5,-112.8)..(59.4,-113.3)
				..controls(60.2,-113.7)and(61.3,-114.3)..(61.9,-114.5)
				..controls(62.4,-114.7)and(63.3,-115.1)..(63.8,-115.3)
				..controls(64.4,-115.6)and(65.4,-115.9)..(66.2,-116.1)
				..controls(67,-116.3)and(68.4,-116.6)..(69.2,-116.9)
				..controls(72.3,-117.9)and(72.8,-118)..(78,-118.1)
				..controls(84,-118.3)and(86.1,-118.1)..(88.1,-117.1)
				..controls(89.4,-116.5)and(93,-115)..(93.7,-114.9)
				..controls(94.8,-114.6)and(95.4,-114.2)..(95.9,-113.4)
				..controls(96.2,-113)and(96.7,-112.5)..(97.1,-112.2)
				..controls(98,-111.6)and(101.3,-108.4)..(102.1,-107.2)
				..controls(102.4,-106.8)and(102.9,-106.2)..(103.2,-105.8)
				..controls(103.5,-105.5)and(104,-104.8)..(104.2,-104.2)
				..controls(104.5,-103.6)and(105,-102.8)..(105.3,-102.4)
				..controls(105.6,-102)and(106.1,-101.2)..(106.5,-100.6)
				..controls(107.1,-99.5)and(108.6,-97.5)..(109.5,-96.5)
				..controls(109.9,-96.1)and(110.1,-96)..(110.4,-96)
				..controls(111.1,-96)and(111.6,-96.5)..(111.8,-97.6)
				..controls(112,-98.6)and(112.3,-99.3)..(112.8,-99.9)
				..controls(113.7,-100.9)and(115.2,-103.1)..(115.7,-104.2)
				..controls(116.1,-105)and(116.5,-105.5)..(116.9,-106)
				..controls(117.3,-106.4)and(117.8,-106.9)..(118,-107.3)
				..controls(118.2,-107.6)and(118.6,-108)..(119,-108.3)
				..controls(119.6,-108.8)and(121.1,-110.3)..(121.8,-111.2)
				..controls(122.1,-111.7)and(122.5,-112)..(123.4,-112.4)
				..controls(124.8,-113.2)and(125.8,-113.8)..(127,-114.6)
				..controls(127.5,-114.9)and(128.1,-115.3)..(128.2,-115.4)
				..controls(128.4,-115.5)and(129.3,-115.8)..(130.2,-116)
				..controls(131.1,-116.2)and(132.5,-116.6)..(133.3,-116.9)
				..controls(135.7,-117.9)and(136.1,-117.9)..(142.5,-117.8)
				..controls(145.8,-117.8)and(148.4,-117.7)..(148.8,-117.6)
				..controls(149.2,-117.6)and(150.2,-117.2)..(151,-116.8)
				..controls(152,-116.4)and(152.8,-116)..(153.6,-115.9)
				..controls(154.5,-115.7)and(155.2,-115.4)..(156.3,-114.9)
				..controls(157.1,-114.5)and(158.2,-114)..(158.7,-113.8)
				..controls(159.1,-113.7)and(159.8,-113.3)..(160.2,-113)
				..controls(160.6,-112.7)and(161.3,-112.4)..(161.7,-112.2)
				..controls(162.7,-111.9)and(163.2,-111.7)..(164.1,-111.1)
				..controls(164.5,-110.8)and(165.3,-110.4)..(165.9,-110)
				..controls(166.5,-109.7)and(167.3,-109.2)..(167.6,-108.9)
				..controls(167.9,-108.6)and(168.6,-108.2)..(169,-108)
				..controls(169.4,-107.8)and(170.2,-107.3)..(170.7,-106.9)
				..controls(171.2,-106.5)and(172.1,-105.8)..(172.7,-105.4)
				..controls(173.2,-105)and(174.1,-104.4)..(174.5,-104)
				..controls(174.9,-103.7)and(175.7,-103.1)..(176.2,-102.6)
				..controls(176.7,-102.2)and(177.4,-101.5)..(177.9,-101.2)
				..controls(178.3,-100.8)and(178.8,-100.4)..(179,-100.2)
				..controls(179.2,-100)and(179.6,-99.6)..(180,-99.3)
				..controls(180.4,-99.1)and(180.9,-98.6)..(181.2,-98.4)
				..controls(181.5,-98.1)and(182.1,-97.7)..(182.5,-97.5)
				..controls(183.2,-97.2)and(183.5,-96.9)..(184.3,-95.9)
				..controls(184.6,-95.6)and(184.6,-95.6)..(184.8,-95.7)
				..controls(185,-95.8)and(185.1,-96)..(185.1,-96.2)
				..controls(185.1,-96.5)and(185,-96.7)..(184.6,-97.2)
				..controls(184.1,-97.8)and(184,-97.9)..(183.9,-98.8)
				..controls(183.7,-100.3)and(183.3,-101.3)..(182.9,-101.8)
				..controls(182.6,-102.1)and(182.3,-102.6)..(182.2,-102.9)
				..controls(182,-103.3)and(181.6,-104)..(181.2,-104.6)
				..controls(180.9,-105.2)and(180.5,-106.1)..(180.4,-106.6)
				..controls(179.9,-108)and(179.7,-108.5)..(179,-109.4)
				..controls(178.6,-109.8)and(178.2,-110.5)..(177.9,-111)
				..controls(177.7,-111.5)and(177.3,-112.3)..(177,-112.7)
				..controls(176.6,-113.2)and(176.4,-113.7)..(176.3,-114.2)
				..controls(176.2,-114.7)and(175.9,-115.2)..(175.4,-116)
				..controls(175,-116.6)and(174.4,-117.7)..(174,-118.5)
				..controls(173.7,-119.2)and(173.3,-120.1)..(173,-120.5)
				..controls(172.8,-120.9)and(172.4,-121.8)..(172.1,-122.5)
				..controls(171.8,-123.2)and(171.3,-124.1)..(171,-124.6)
				..controls(170.8,-125.1)and(170.3,-126)..(169.9,-126.6)
				..controls(169.6,-127.3)and(169.2,-128.1)..(169,-128.3)
				..controls(168.8,-128.6)and(168.5,-129.3)..(168.3,-129.8)
				..controls(168.1,-130.3)and(167.8,-131.1)..(167.5,-131.6)
				..controls(166.9,-132.5)and(164.2,-137.9)..(164.2,-138.1)
				..controls(164.2,-138.1)and(163.9,-138.8)..(163.4,-139.7)
				..controls(162.9,-140.5)and(162.4,-141.5)..(162.3,-141.9)
				..controls(161.9,-143)and(161.6,-143.4)..(160.9,-144.1)
				..controls(160.4,-144.6)and(160.3,-144.8)..(160.1,-145.4)
				..controls(160,-145.8)and(159.6,-146.8)..(159.2,-147.8)
				..controls(158.7,-148.8)and(158.1,-150)..(157.9,-150.6)
				..controls(157.6,-151.1)and(157.2,-151.9)..(156.9,-152.4)
				..controls(156.5,-152.8)and(156.1,-153.7)..(155.8,-154.4)
				..controls(155.6,-155)and(155.2,-155.9)..(154.9,-156.3)
				..controls(154.7,-156.7)and(154.4,-157.4)..(154.2,-157.9)
				..controls(154.1,-158.4)and(153.8,-159)..(153.5,-159.3)
				..controls(152.8,-160.1)and(153,-160.1)..(142.4,-160.1)
				..controls(137.2,-160.1)and(131.8,-160.1)..(130.4,-160)--cycle
				(141.5,-103.7)..controls(140.8,-103)and(139.8,-102.1)..(139.2,-101.6)
				..controls(138.5,-101.2)and(137.9,-100.6)..(137.8,-100.4)
				..controls(137,-99.5)and(134,-96.5)..(133.2,-96)
				..controls(132.7,-95.7)and(131.8,-94.9)..(131.1,-94.3)
				..controls(129.4,-92.7)and(127.9,-91.6)..(126.6,-91.1)
				..controls(125.2,-90.5)and(124.1,-90)..(122.9,-89.1)
				..controls(122.1,-88.6)and(121.3,-88.2)..(119.9,-87.7)
				..controls(118.8,-87.3)and(117.5,-86.8)..(116.9,-86.5)
				..controls(114.3,-85.5)and(113.7,-85.4)..(109.7,-85.5)
				..controls(105.3,-85.7)and(103.9,-85.9)..(101.6,-87)
				..controls(100.7,-87.4)and(99.4,-87.8)..(97.9,-88)
				..controls(97,-88.2)and(96.6,-88.3)..(95.8,-88.9)
				..controls(95.4,-89.2)and(94.7,-89.6)..(94.2,-89.8)
				..controls(92.9,-90.3)and(92.3,-90.7)..(90.9,-92.1)
				..controls(90.2,-92.7)and(89.4,-93.3)..(89,-93.5)..controls(87,-94.5)and(85.2,-96)..(83.6,-97.9)
				..controls(83.3,-98.2)and(82.5,-98.9)..(81.8,-99.4)
				..controls(80.7,-100.1)and(80.5,-100.3)..(80,-101.1)
				..controls(78.7,-103.1)and(77.2,-103.9)..(76.4,-103)
				..controls(75.7,-102.2)and(75.7,-103.2)..(75.7,-83.6)
				..controls(75.6,-73.6)and(75.6,-64.8)..(75.5,-64)
				..controls(75.5,-63.2)and(75.5,-61.4)..(75.6,-59.8)--(75.7,-57.1)--(76,-56.8)
				..controls(76.8,-56.2)and(78,-56.3)..(79.8,-57.2)
				..controls(80.4,-57.5)and(81.3,-57.8)..(81.8,-57.9)
				..controls(82.3,-58.1)and(83,-58.4)..(83.6,-58.8)..controls(84.3,-59.2)and(85,-59.5)..(86,-59.8)
				..controls(86.8,-60.1)and(88,-60.6)..(88.7,-60.9)
				..controls(89.3,-61.2)and(90.6,-61.7)..(91.4,-62)
				..controls(92.2,-62.3)and(93.7,-62.9)..(94.6,-63.3)
				..controls(95.6,-63.7)and(96.8,-64.1)..(97.4,-64.3)
				..controls(98,-64.5)and(99,-64.9)..(99.7,-65.2)
				..controls(100.4,-65.6)and(101.5,-66)..(102.3,-66.2)
				..controls(104.2,-66.8)and(105,-67.1)..(106.2,-67.8)
				..controls(107.9,-68.7)and(110.2,-69.4)..(111.3,-69.2)
				..controls(112.5,-69)and(115,-68.2)..(115.6,-67.8)
				..controls(117.3,-66.7)and(117.6,-66.6)..(119.5,-66.2)
				..controls(121.9,-65.8)and(122.1,-65.7)..(123,-65.1)
				..controls(124.1,-64.4)and(124.4,-64.3)..(126,-63.9)
				..controls(127,-63.7)and(127.7,-63.4)..(128.4,-63.1)
				..controls(129,-62.8)and(129.7,-62.5)..(130.6,-62.4)
				..controls(132.1,-62.1)and(134,-61.5)..(135.2,-61)
				..controls(135.7,-60.8)and(136.7,-60.5)..(137.3,-60.3)
				..controls(138.4,-60)and(139.3,-59.7)..(141.1,-59)
				..controls(141.6,-58.7)and(142,-58.6)..(142.5,-58.6)
				..controls(143.2,-58.6)and(143.3,-58.7)..(143.6,-59)
				..controls(144,-59.5)and(144.1,-59.9)..(144,-60.8)
				..controls(144,-61.2)and(144,-71.1)..(144,-82.8)--(143.9,-104.1)--(143.7,-104.5)
				..controls(143.5,-104.7)and(143.3,-104.8)..(143,-104.9)
				..controls(142.6,-104.9)and(142.5,-104.8)..(141.5,-103.7)--cycle
				(27.8,-91.4)..controls(27.1,-90.9)and(26.3,-90)..(26.1,-89.5)
				..controls(26,-89.2)and(25.8,-88.6)..(25.7,-87.9)
				..controls(25.4,-86.2)and(25.6,-84.8)..(26.5,-83.5)
				..controls(27.4,-82.3)and(28.3,-82)..(29.4,-82.3)
				..controls(31.1,-82.8)and(32.1,-84.9)..(32,-87.4)
				..controls(31.9,-89.3)and(31.2,-90.6)..(29.9,-91.4)
				..controls(29.1,-91.9)and(28.6,-91.9)..(27.8,-91.4)--cycle
				(190.1,-91.3)..controls(189.7,-91.2)and(189.3,-90.9)..(189.1,-90.7)
				..controls(188.3,-89.8)and(187.6,-87.7)..(187.6,-86.2)
				..controls(187.6,-85.5)and(187.7,-85.2)..(188.2,-84.4)
				..controls(189.1,-82.6)and(190.1,-81.8)..(191.3,-82.1)
				..controls(192,-82.3)and(193.5,-83.7)..(193.8,-84.4)
				..controls(194.1,-85.1)and(194.2,-87.1)..(194,-88.3)
				..controls(193.8,-89.8)and(193.3,-90.5)..(192.3,-91.3)
				..controls(191.7,-91.7)and(190.9,-91.7)..(190.1,-91.3)--cycle
				(173.7,-83.4)..controls(172.2,-82.5)and(172,-82.5)..(170.6,-82.4)
				..controls(169.7,-82.4)and(168.8,-82.4)..(168.3,-82.5)
				..controls(167.2,-82.7)and(166.9,-82.7)..(166.6,-82.2)
				..controls(166.3,-81.8)and(166.3,-81)..(166.8,-80)
				..controls(167.2,-78.9)and(167.4,-77.4)..(167.7,-74.1)
				..controls(167.8,-72.8)and(167.9,-71.8)..(167.9,-71.7)
				..controls(168,-71.7)and(168,-70.2)..(168,-68.4)..controls(168,-65.4)and(168,-64.9)..(167.7,-63)
				..controls(167.3,-60.5)and(167.1,-59.9)..(166.4,-58.4)
				..controls(166.1,-57.7)and(165.8,-56.9)..(165.7,-56.6)
				..controls(165.7,-56.2)and(165.4,-55.6)..(165.2,-55.2)
				..controls(164.6,-53.8)and(161.9,-50.7)..(161,-50.3)
				..controls(159.8,-49.7)and(157.9,-50)..(155.2,-51.2)
				..controls(154.7,-51.4)and(153.4,-51.8)..(152.2,-52.1)
				..controls(150.8,-52.4)and(149.4,-52.9)..(148.6,-53.2)
				..controls(147.8,-53.6)and(146.7,-54)..(146,-54.2)
				..controls(145.4,-54.4)and(144.5,-54.7)..(144,-54.9)
				..controls(143.6,-55.2)and(142.7,-55.5)..(141.8,-55.7)
				..controls(141,-55.9)and(139.6,-56.3)..(138.7,-56.6)
				..controls(137.8,-57)and(136.3,-57.4)..(135.3,-57.6)
				..controls(133.8,-57.9)and(133.2,-58.1)..(132.1,-58.7)
				..controls(131.1,-59.2)and(130.3,-59.4)..(129.4,-59.6)
				..controls(128.7,-59.8)and(127.8,-60)..(127.4,-60.1)
				..controls(127,-60.3)and(126,-60.5)..(125.1,-60.7)
				..controls(123.2,-61.1)and(121.7,-61.5)..(120.2,-62.1)
				..controls(119.7,-62.3)and(118.1,-62.8)..(116.7,-63.2)
				..controls(114.6,-63.9)and(114,-64.1)..(113,-64.7)
				..controls(111.9,-65.4)and(111.8,-65.4)..(111,-65.5)
				..controls(109.9,-65.5)and(109.3,-65.4)..(107.7,-64.7)
				..controls(107,-64.4)and(106,-64)..(105.3,-63.9)
				..controls(103.5,-63.5)and(102.2,-63)..(101.2,-62.5)
				..controls(100.1,-61.8)and(99.3,-61.5)..(97.8,-61.1)
				..controls(97.2,-61)and(96.2,-60.7)..(95.7,-60.4)
				..controls(95.1,-60.1)and(94.2,-59.8)..(93.5,-59.6)
				..controls(92.9,-59.4)and(91.7,-59)..(91,-58.7)..controls(90.2,-58.4)and(88.9,-58)..(88,-57.7)
				..controls(87.2,-57.5)and(86.1,-57.1)..(85.5,-56.8)
				..controls(85,-56.6)and(84.1,-56.2)..(83.5,-56)..controls(82.9,-55.7)and(81.8,-55.2)..(81,-54.9)
				..controls(80.1,-54.5)and(79.1,-54.2)..(78.7,-54.1)
				..controls(77.9,-53.8)and(77.3,-53.6)..(76.2,-52.9)
				..controls(75.7,-52.6)and(75.2,-52.4)..(74.1,-52.2)
				..controls(72.5,-51.9)and(72.4,-51.8)..(71.3,-51.1)
				..controls(69.9,-50.1)and(69.4,-49.9)..(67.9,-49.8)
				..controls(67,-49.7)and(66.6,-49.6)..(66.3,-49.4)
				..controls(66.1,-49.3)and(65.7,-49)..(65.4,-48.8)
				..controls(65.1,-48.7)and(64.2,-48.4)..(63.6,-48.2)
				..controls(62.9,-48)and(61.7,-47.6)..(60.9,-47.2)
				..controls(60.1,-46.8)and(58.9,-46.4)..(58.2,-46.2)
				..controls(57.6,-46)and(56.8,-45.7)..(56.5,-45.5)
				..controls(54.9,-44.6)and(54.6,-44.5)..(53.2,-44.2)
				..controls(51.7,-43.8)and(51,-43.6)..(49.9,-42.9)
				..controls(49.5,-42.6)and(48.9,-42.4)..(47.5,-42)
				..controls(45.7,-41.6)and(45.6,-41.5)..(44.3,-40.6)
				..controls(43.5,-40.1)and(42.8,-39.6)..(42.7,-39.4)
				..controls(42.4,-38.7)and(43,-38.1)..(44.5,-37.6)..controls(44.9,-37.5)and(46,-37)..(46.9,-36.6)
				..controls(47.9,-36.1)and(48.8,-35.8)..(49.6,-35.6)
				..controls(50.5,-35.4)and(51.3,-35.2)..(52.4,-34.6)
				..controls(53.3,-34.1)and(54.5,-33.7)..(54.9,-33.5)
				..controls(55.4,-33.4)and(56.4,-33)..(57.2,-32.6)
				..controls(58.6,-32)and(60.1,-31.5)..(61.6,-31.2)
				..controls(62.3,-31.1)and(62.7,-31)..(63.2,-30.7)
				..controls(64.1,-30)and(64.3,-29.9)..(65.3,-29.7)
				..controls(65.7,-29.6)and(66.5,-29.2)..(67.1,-28.9)
				..controls(67.6,-28.6)and(68.8,-28.2)..(69.6,-27.9)
				..controls(70.5,-27.7)and(71.6,-27.3)..(72.1,-27)
				..controls(72.6,-26.7)and(73.8,-26.2)..(74.8,-25.9)
				..controls(75.7,-25.5)and(76.7,-25.1)..(77.1,-25)
				..controls(77.4,-24.8)and(78.5,-24.4)..(79.6,-24.1)
				..controls(80.6,-23.8)and(82,-23.2)..(82.8,-22.9)
				..controls(83.5,-22.5)and(84.8,-22)..(85.8,-21.8)
				..controls(86.7,-21.5)and(87.8,-21.1)..(88.1,-20.9)
				..controls(88.4,-20.7)and(89.3,-20.3)..(90.1,-20.1)
				..controls(90.8,-19.8)and(91.9,-19.4)..(92.4,-19.1)
				..controls(93.5,-18.5)and(94.4,-18.2)..(95.4,-18.1)
				..controls(95.7,-18)and(96.4,-17.8)..(96.9,-17.5)
				..controls(97.5,-17.3)and(98.4,-16.9)..(99,-16.6)
				..controls(99.7,-16.3)and(100.9,-15.9)..(101.9,-15.6)
				..controls(102.9,-15.3)and(104.4,-14.7)..(105.3,-14.3)
				..controls(108.1,-13.2)and(108.5,-13.1)..(110.1,-13.1)
				..controls(111.4,-13.1)and(111.5,-13.1)..(112.4,-13.5)
				..controls(113.2,-14)and(115,-14.8)..(116.6,-15.4)
				..controls(116.9,-15.5)and(117.9,-15.8)..(118.8,-16.1)
				..controls(120.9,-16.6)and(121.5,-16.9)..(122.6,-17.9)
				..controls(123,-18.3)and(123.7,-18.8)..(123.9,-18.9)
				..controls(124.6,-19.4)and(126,-19.8)..(127.9,-20.2)
				..controls(129.6,-20.5)and(129.8,-20.6)..(131.1,-21.4)
				..controls(131.5,-21.7)and(132.4,-22.1)..(133.1,-22.3)
				..controls(136.1,-23.2)and(138.3,-24.1)..(140.2,-25)
				..controls(141,-25.5)and(141.5,-25.6)..(142.5,-25.8)
				..controls(145.1,-26.3)and(146,-26.6)..(149.4,-28.1)
				..controls(151.3,-28.9)and(153.5,-29.8)..(154.3,-30.1)
				..controls(155.1,-30.4)and(156.2,-30.8)..(156.6,-31)
				..controls(157.1,-31.2)and(158.2,-31.6)..(159.1,-31.8)
				..controls(160.6,-32.2)and(160.9,-32.3)..(161.7,-32.9)
				..controls(162.4,-33.3)and(162.9,-33.6)..(163.7,-33.8)
				..controls(164.3,-34)and(164.8,-34.1)..(164.9,-34.1)
				..controls(165,-34.1)and(165.6,-34.4)..(166.4,-34.8)
				..controls(167.3,-35.2)and(168.1,-35.4)..(169,-35.6)
				..controls(170,-35.8)and(170.5,-36)..(171.7,-36.6)
				..controls(172.4,-37)and(173.3,-37.4)..(173.6,-37.4)
				..controls(175.6,-37.7)and(176.2,-38)..(176.7,-39.1)
				..controls(176.9,-39.6)and(176.9,-39.7)..(176.7,-40.4)
				..controls(176.4,-41.6)and(175.8,-42.7)..(175.2,-43.4)
				..controls(174.5,-44.1)and(172.6,-45.5)..(171.7,-45.8)
				..controls(170.6,-46.2)and(169.1,-46)..(168.2,-45.4)
				..controls(167.9,-45.2)and(167.3,-44.8)..(166.9,-44.4)
				..controls(166.4,-44.1)and(165.8,-43.6)..(165.4,-43.2)
				..controls(165.1,-42.9)and(164.4,-42.3)..(163.9,-41.9)
				..controls(163.3,-41.5)and(162.7,-40.9)..(162.4,-40.5)
				..controls(161.4,-39.2)and(159.9,-38)..(158.8,-37.8)
				..controls(158.4,-37.8)and(158.3,-37.7)..(158.2,-37.2)
				..controls(158,-36.8)and(157.4,-36.1)..(157.1,-36.1)
				..controls(157,-36.1)and(156.7,-36.1)..(156.5,-36)
				..controls(156.1,-35.9)and(155.9,-35.7)..(155.7,-35.2)
				..controls(155.3,-34.6)and(155.1,-34.4)..(154,-34.1)
				..controls(153.6,-34)and(153.2,-33.7)..(152.9,-33.5)
				..controls(152.4,-32.9)and(151.6,-32.4)..(151.3,-32.4)
				..controls(150.9,-32.4)and(150.4,-32.6)..(150.4,-32.9)
				..controls(150.4,-33.2)and(150.8,-33.4)..(151.4,-33.6)
				..controls(152.1,-33.7)and(152.8,-34.1)..(153.5,-34.9)
				..controls(154.1,-35.5)and(154.3,-35.6)..(155,-35.8)
				..controls(155.7,-36)and(155.7,-36)..(156,-36.6)
				..controls(156.3,-37.3)and(156.7,-37.7)..(157.5,-37.8)
				..controls(158.1,-37.9)and(158.2,-38)..(158.4,-38.8)
				..controls(158.5,-39.4)and(158.6,-39.6)..(159.6,-40.4)
				..controls(160.2,-41)and(160.9,-41.6)..(161.3,-41.8)
				..controls(161.7,-42)and(161.9,-42.3)..(162.2,-42.8)
				..controls(162.4,-43.2)and(162.8,-43.8)..(163.1,-44.1)
				..controls(163.9,-44.8)and(165.1,-46.7)..(165.7,-48)
				..controls(166.2,-49.3)and(166.4,-49.7)..(167.2,-50.5)
				..controls(167.5,-50.8)and(167.9,-51.4)..(168.1,-51.8)
				..controls(168.3,-52.2)and(168.8,-53)..(169.1,-53.6)
				..controls(169.6,-54.3)and(169.8,-54.9)..(169.9,-55.6)
				..controls(170.4,-57.8)and(170.5,-58.1)..(171.2,-59.5)
				..controls(171.6,-60.3)and(172,-61.5)..(172.2,-62.1)
				..controls(172.3,-62.8)and(172.7,-63.7)..(172.9,-64.1)
				..controls(173.2,-64.6)and(173.4,-65.3)..(173.4,-65.6)
				..controls(173.8,-68.8)and(174.2,-70.6)..(174.9,-72.6)
				..controls(175.5,-74.5)and(175.9,-77.1)..(176.1,-80.2)
				..controls(176.3,-83.4)and(175.5,-84.4)..(173.7,-83.4)--cycle;
				\definecolor{ezs}{RGB}{247,179,26}\fill[ezs](-5.7,-156.6)--(-3.4,-157.5)--(61.9,12.7)--(59.7,13.6)--cycle;
				\definecolor{ezs}{RGB}{247,179,22}\fill[ezs](-3.8,-157.3)--(-1.5,-158.2)--(63.9,12)--(61.6,12.9)--cycle;
				\definecolor{ezs}{RGB}{248,179,19}\fill[ezs](-1.8,-158.1)--(0.4,-158.9)--(65.8,11.2)--(63.5,12.1)--cycle;
				\definecolor{ezs}{RGB}{248,178,16}\fill[ezs](0.1,-158.8)--(2.4,-159.7)--(67.7,10.5)--(65.4,11.4)--cycle;
				\definecolor{ezs}{RGB}{249,178,15}\fill[ezs](2,-159.5)--(4.3,-160.4)--(69.6,9.8)--(67.3,10.7)--cycle;
				\definecolor{ezs}{RGB}{251,177,17}\fill[ezs](3.9,-160.3)--(6.2,-161.1)--(71.5,9)--(69.2,9.9)--cycle;
				\definecolor{ezs}{RGB}{252,179,18}\fill[ezs](5.8,-161)--(8.1,-161.9)--(73.4,8.3)--(71.1,9.2)--cycle;
				\definecolor{ezs}{RGB}{252,182,17}\fill[ezs](7.7,-161.7)--(10,-162.6)--(75.3,7.6)--(73,8.5)--cycle;
				\definecolor{ezs}{RGB}{252,185,17}\fill[ezs](9.6,-162.5)--(11.9,-163.4)--(77.2,6.8)--(75,7.7)--cycle;
				\definecolor{ezs}{RGB}{251,189,18}\fill[ezs](11.5,-163.2)--(13.8,-164.1)--(79.2,6.1)--(76.9,7)--cycle;
				\definecolor{ezs}{RGB}{251,191,22}\fill[ezs](13.4,-163.9)--(15.7,-164.8)--(81.1,5.4)--(78.8,6.2)--cycle;
				\definecolor{ezs}{RGB}{252,193,25}\fill[ezs](15.4,-164.7)--(17.6,-165.6)--(83,4.6)--(80.7,5.5)--cycle;
				\definecolor{ezs}{RGB}{253,195,27}\fill[ezs](17.3,-165.4)--(19.6,-166.3)--(84.9,3.9)--(82.6,4.8)--cycle;
				\definecolor{ezs}{RGB}{253,196,31}\fill[ezs](19.2,-166.1)--(21.5,-167)--(86.8,3.2)--(84.5,4)--cycle;
				\definecolor{ezs}{RGB}{254,199,35}\fill[ezs](21.1,-166.9)--(23.4,-167.8)--(88.7,2.4)--(86.4,3.3)--cycle;
				\definecolor{ezs}{RGB}{253,201,38}\fill[ezs](23,-167.6)--(25.3,-168.5)--(90.6,1.7)--(88.3,2.6)--cycle;
				\definecolor{ezs}{RGB}{253,203,40}\fill[ezs](24.9,-168.3)--(27.2,-169.2)--(92.5,1)--(90.3,1.8)--cycle;
				\definecolor{ezs}{RGB}{253,205,42}\fill[ezs](26.8,-169.1)--(29.1,-170)--(94.4,0.2)--(92.2,1.1)--cycle;
				\definecolor{ezs}{RGB}{253,207,45}\fill[ezs](28.7,-169.8)--(31,-170.7)--(96.4,-0.5)--(94.1,0.4)--cycle;
				\definecolor{ezs}{RGB}{253,210,49}\fill[ezs](30.7,-170.5)--(32.9,-171.4)--(98.3,-1.2)--(96,-0.4)--cycle;
				\definecolor{ezs}{RGB}{253,212,54}\fill[ezs](32.6,-171.3)--(34.9,-172.2)--(100.2,-2)--(97.9,-1.1)--cycle;
				\definecolor{ezs}{RGB}{253,214,59}\fill[ezs](34.5,-172)--(36.8,-172.9)--(102.1,-2.7)--(99.8,-1.8)--cycle;
				\definecolor{ezs}{RGB}{253,216,63}\fill[ezs](36.4,-172.7)--(38.7,-173.6)--(104,-3.4)--(101.7,-2.6)--cycle;
				\definecolor{ezs}{RGB}{253,218,65}\fill[ezs](38.3,-173.5)--(40.6,-174.4)--(105.9,-4.2)--(103.6,-3.3)--cycle;
				\definecolor{ezs}{RGB}{254,220,66}\fill[ezs](40.2,-174.2)--(42.5,-175.1)--(107.8,-4.9)--(105.5,-4)--cycle;
				\definecolor{ezs}{RGB}{254,222,68}\fill[ezs](42.1,-175)--(44.4,-175.8)--(109.7,-5.6)--(107.5,-4.8)--cycle;
				\definecolor{ezs}{RGB}{254,224,71}\fill[ezs](44,-175.7)--(46.3,-176.6)--(111.7,-6.4)--(109.4,-5.5)--cycle;
				\definecolor{ezs}{RGB}{254,226,75}\fill[ezs](46,-176.4)--(48.2,-177.3)--(113.6,-7.1)--(111.3,-6.2)--cycle;
				\definecolor{ezs}{RGB}{254,228,79}\fill[ezs](47.9,-177.2)--(50.1,-178)--(115.5,-7.8)--(113.2,-7)--cycle;
				\definecolor{ezs}{RGB}{254,230,82}\fill[ezs](49.8,-177.9)--(52.1,-178.8)--(117.4,-8.6)--(115.1,-7.7)--cycle;
				\definecolor{ezs}{RGB}{254,232,85}\fill[ezs](51.7,-178.6)--(54,-179.5)--(119.3,-9.3)--(117,-8.4)--cycle;
				\definecolor{ezs}{RGB}{254,234,88}\fill[ezs](53.6,-179.4)--(55.9,-180.2)--(121.2,-10)--(118.9,-9.2)--cycle;
				\definecolor{ezs}{RGB}{254,237,91}\fill[ezs](55.5,-180.1)--(57.8,-181)--(123.1,-10.8)--(120.8,-9.9)--cycle;
				\definecolor{ezs}{RGB}{254,239,94}\fill[ezs](57.4,-180.8)--(59.7,-181.7)--(125,-11.5)--(122.8,-10.6)--cycle;
				\definecolor{ezs}{RGB}{254,241,97}\fill[ezs](59.3,-181.6)--(61.6,-182.4)--(127,-12.2)--(124.7,-11.4)--cycle;
				\definecolor{ezs}{RGB}{254,243,101}\fill[ezs](61.2,-182.3)--(63.5,-183.2)--(128.9,-13)--(126.6,-12.1)--cycle;
				\definecolor{ezs}{RGB}{254,245,103}\fill[ezs](63.2,-183)--(65.4,-183.9)--(130.8,-13.7)--(128.5,-12.8)--cycle;
				\definecolor{ezs}{RGB}{254,245,105}\fill[ezs](65.1,-183.8)--(67.4,-184.6)--(132.7,-14.4)--(130.4,-13.6)--cycle;
				\definecolor{ezs}{RGB}{254,245,105}\fill[ezs](67,-184.5)--(69.3,-185.4)--(134.6,-15.2)--(132.3,-14.3)--cycle;
				\definecolor{ezs}{RGB}{254,244,103}\fill[ezs](68.9,-185.2)--(71.2,-186.1)--(136.5,-15.9)--(134.2,-15)--cycle;
				\definecolor{ezs}{RGB}{253,243,101}\fill[ezs](70.8,-186)--(73.1,-186.8)--(138.4,-16.6)--(136.1,-15.8)--cycle;
				\definecolor{ezs}{RGB}{253,241,98}\fill[ezs](72.7,-186.7)--(75,-187.6)--(140.3,-17.4)--(138,-16.5)--cycle;
				\definecolor{ezs}{RGB}{253,239,95}\fill[ezs](74.6,-187.4)--(76.9,-188.3)--(142.2,-18.1)--(140,-17.2)--cycle;
				\definecolor{ezs}{RGB}{253,236,92}\fill[ezs](76.5,-188.2)--(78.8,-189)--(144.2,-18.8)--(141.9,-18)--cycle;
				\definecolor{ezs}{RGB}{253,234,88}\fill[ezs](78.5,-188.9)--(80.7,-189.8)--(146.1,-19.6)--(143.8,-18.7)--cycle;
				\definecolor{ezs}{RGB}{253,232,85}\fill[ezs](80.4,-189.6)--(82.7,-190.5)--(148,-20.3)--(145.7,-19.4)--cycle;
				\definecolor{ezs}{RGB}{253,230,82}\fill[ezs](82.3,-190.4)--(84.6,-191.2)--(149.9,-21)--(147.6,-20.2)--cycle;
				\definecolor{ezs}{RGB}{253,228,78}\fill[ezs](84.2,-191.1)--(86.5,-192)--(151.8,-21.8)--(149.5,-20.9)--cycle;
				\definecolor{ezs}{RGB}{253,226,74}\fill[ezs](86.1,-191.8)--(88.4,-192.7)--(153.7,-22.5)--(151.4,-21.6)--cycle;
				\definecolor{ezs}{RGB}{253,224,70}\fill[ezs](88,-192.6)--(90.3,-193.4)--(155.6,-23.3)--(153.3,-22.4)--cycle;
				\definecolor{ezs}{RGB}{253,222,67}\fill[ezs](89.9,-193.3)--(92.2,-194.2)--(157.5,-24)--(155.3,-23.1)--cycle;
				\definecolor{ezs}{RGB}{253,220,65}\fill[ezs](91.8,-194)--(94.1,-194.9)--(159.5,-24.7)--(157.2,-23.8)--cycle;
				\definecolor{ezs}{RGB}{253,217,63}\fill[ezs](93.7,-194.8)--(96,-195.6)--(161.4,-25.5)--(159.1,-24.6)--cycle;
				\definecolor{ezs}{RGB}{253,215,60}\fill[ezs](95.7,-195.5)--(97.9,-196.4)--(163.3,-26.2)--(161,-25.3)--cycle;
				\definecolor{ezs}{RGB}{253,213,56}\fill[ezs](97.6,-196.2)--(99.9,-197.1)--(165.2,-26.9)--(162.9,-26)--cycle;
				\definecolor{ezs}{RGB}{252,211,53}\fill[ezs](99.5,-197)--(101.8,-197.8)--(167.1,-27.7)--(164.8,-26.8)--cycle;
				\definecolor{ezs}{RGB}{252,209,49}\fill[ezs](101.4,-197.7)--(103.7,-198.6)--(169,-28.4)--(166.7,-27.5)--cycle;
				\definecolor{ezs}{RGB}{252,207,45}\fill[ezs](103.3,-198.4)--(105.6,-199.3)--(170.9,-29.1)--(168.6,-28.2)--cycle;
				\definecolor{ezs}{RGB}{253,205,42}\fill[ezs](105.2,-199.2)--(107.5,-200)--(172.8,-29.9)--(170.6,-29)--cycle;
				\definecolor{ezs}{RGB}{253,202,40}\fill[ezs](107.1,-199.9)--(109.4,-200.8)--(174.7,-30.6)--(172.5,-29.7)--cycle;
				\definecolor{ezs}{RGB}{253,200,38}\fill[ezs](109,-200.6)--(111.3,-201.5)--(176.7,-31.3)--(174.4,-30.4)--cycle;
				\definecolor{ezs}{RGB}{253,198,35}\fill[ezs](111,-201.4)--(113.2,-202.2)--(178.6,-32.1)--(176.3,-31.2)--cycle;
				\definecolor{ezs}{RGB}{253,196,31}\fill[ezs](112.9,-202.1)--(115.2,-203)--(180.5,-32.8)--(178.2,-31.9)--cycle;
				\definecolor{ezs}{RGB}{253,194,28}\fill[ezs](114.8,-202.8)--(117.1,-203.7)--(182.4,-33.5)--(180.1,-32.6)--cycle;
				\definecolor{ezs}{RGB}{253,192,26}\fill[ezs](116.7,-203.6)--(119,-204.5)--(184.3,-34.3)--(182,-33.4)--cycle;
				\definecolor{ezs}{RGB}{253,189,24}\fill[ezs](118.6,-204.3)--(120.9,-205.2)--(186.2,-35)--(183.9,-34.1)--cycle;
				\definecolor{ezs}{RGB}{253,186,20}\fill[ezs](120.5,-205)--(122.8,-205.9)--(188.1,-35.7)--(185.8,-34.8)--cycle;
				\definecolor{ezs}{RGB}{253,184,17}\fill[ezs](122.4,-205.8)--(124.7,-206.7)--(190,-36.5)--(187.8,-35.6)--cycle;
				\definecolor{ezs}{RGB}{252,183,15}\fill[ezs](124.3,-206.5)--(126.6,-207.4)--(192,-37.2)--(189.7,-36.3)--cycle;
				\definecolor{ezs}{RGB}{252,181,14}\fill[ezs](126.3,-207.2)--(128.5,-208.1)--(193.9,-37.9)--(191.6,-37.1)--cycle;
				\definecolor{ezs}{RGB}{252,179,13}\fill[ezs](128.2,-208)--(130.4,-208.9)--(195.8,-38.7)--(193.5,-37.8)--cycle;
				\definecolor{ezs}{RGB}{252,178,12}\fill[ezs](130.1,-208.7)--(132.4,-209.6)--(197.7,-39.4)--(195.4,-38.5)--cycle;
				\definecolor{ezs}{RGB}{251,178,10}\fill[ezs](132,-209.4)--(134.3,-210.3)--(199.6,-40.1)--(197.3,-39.3)--cycle;
				\definecolor{ezs}{RGB}{251,178,10}\fill[ezs](133.9,-210.2)--(136.2,-211.1)--(201.5,-40.9)--(199.2,-40)--cycle;
				\definecolor{ezs}{RGB}{250,177,11}\fill[ezs](135.8,-210.9)--(138.1,-211.8)--(203.4,-41.6)--(201.1,-40.7)--cycle;
				\definecolor{ezs}{RGB}{250,178,11}\fill[ezs](137.7,-211.6)--(140,-212.5)--(205.3,-42.3)--(203.1,-41.5)--cycle;
				\definecolor{ezs}{RGB}{248,180,10}\fill[ezs](139.6,-212.4)--(141.9,-213.3)--(207.3,-43.1)--(205,-42.2)--cycle;
				\definecolor{ezs}{RGB}{247,181,9}\fill[ezs](141.5,-213.1)--(143.8,-214)--(209.2,-43.8)--(206.9,-42.9)--cycle;
				\definecolor{ezs}{RGB}{246,182,9}\fill[ezs](143.5,-213.8)--(145.7,-214.7)--(211.1,-44.5)--(208.8,-43.7)--cycle;
				\definecolor{ezs}{RGB}{246,182,10}\fill[ezs](145.4,-214.6)--(147.7,-215.5)--(213,-45.3)--(210.7,-44.4)--cycle;
			\end{scope}
			% chữ EDUZY (Times đậm, làm dày nét bằng 5 lớp lệch nhẹ; căn đúng khung chữ logo gốc)
			\foreach \dx/\dy in {-0.7/0,0.7/0,0/-0.3,0/0.3,0/0}
			\node[anchor=center,inner sep=0pt,outer sep=0pt,text=ezText,align=center] at (109.6+\dx,-193.6+\dy)
			{\resizebox{1.01cm}{0.22cm}{\fontencoding{T1}\fontfamily{ptm}\fontseries{b}\selectfont EDUZY}};
		\end{tikzpicture}%
	}%
}
%====================== THIẾT LẬP THÂN ĐỀ ======================
%--- Môi trường tự luận {bt}: ex_test 3.3.4 chưa khai báo sẵn nên tạo ở đây
%    (ex_test đã tự móc \fixshowans và \listenumerate cho bt)
\makeatletter
\@ifundefined{bt}{\newtheorem{bt}{\color{gold1}\namebt}}{}
\makeatother
\renewcommand{\sectionmark}[1]{\markright{\thesection\ #1}}
%--- Khung tiêu đề (logo + tối đa 4 dòng; dòng 3, 4 để trống {} thì bỏ qua)
%  Đề thi:    \khungtieude{THI THỬ THPT QUỐC GIA NĂM 2026}{Môn Toán -- Đề số 21}
%                         {(Đề thi gồm: ...)}{Thời gian làm bài: 90 phút (không kể thời gian phát đề).}
%  Chuyên đề: \khungtieude{CHUYÊN ĐỀ 3 -- DÃY SỐ VÀ CÔNG THỨC TRUY HỒI}
%                         {Bồi dưỡng học sinh giỏi Toán THPT}{...}{}
\newcommand{\khungtieude}[4]{%
	\begin{tcolorbox}[colframe=gold1,colback=gold5!45,boxrule=0.5mm]
		\begin{minipage}[c]{3cm}
			\centerline{\eduzylogo[2.4cm]}
		\end{minipage}\hspace{0.5cm}
		\begin{minipage}[c]{12.5cm}
			\vskip 0.15 cm
			\centerline{\ind{#1}}
			\vskip 0.15 cm
			\centerline{\textbf{#2}}
			\if\relax\detokenize{#3}\relax\else
				\vskip 0.15 cm
				\centerline{\textit{#3}}
			\fi
			\if\relax\detokenize{#4}\relax\else
				\vskip 0.15 cm
				\centerline{\textit{#4}}
			\fi
			\vskip 0.15 cm
		\end{minipage}
	\end{tcolorbox}
}
%====================== BỔ SUNG CHO BÀI HỌC, CHUYÊN ĐỀ ======================
\usepackage[shortlabels]{enumitem}
\usepackage{tabularx}
%--- Danh sách gọn: \begin{enumerate}[gon] ... \end{enumerate}
\SetEnumitemKey{gon}{leftmargin=*,itemsep=1pt,parsep=0pt,topsep=3pt}
%--- Khung lý thuyết: \begin{khung}{} ... \end{khung}; tiêu đề chỉ dùng khi thật ngắn ("Mẹo", "Lỗi thường gặp")
\newtcolorbox{khung}[1]{breakable,enhanced,colframe=gold1,colback=gold5!45,
	boxrule=0.4mm,fonttitle=\bfseries\sffamily,coltitle=white,colbacktitle=gold1,
	attach boxed title to top left={xshift=6mm,yshift=-3mm},
	boxed title style={arc=1mm,boxrule=0.3mm},title={#1}}
%--- Tiêu đề mục nhỏ trong một dạng: \muc{Dấu hiệu}, \muc{Phương pháp}, \muc{Bài tự luyện}
\makeatletter
\newcommand{\muc}[1]{\par\medskip\noindent\indam{#1}\par\nobreak\smallskip\@afterindentfalse\@afterheading}
\makeatother
%--- Ô công thức cần nhớ (căn giữa, viền vàng, vừa khít nội dung)
%     \begin{center}\congthuc{$Q(n)=(n-1)!$}\end{center}
\newtcbox{\congthuc}{on line,colframe=gold1,colback=gold5!45,boxrule=0.5pt,arc=2pt,
	left=8pt,right=8pt,top=4pt,bottom=4pt}
%--- Công tắc lời giải bài tự luyện (ví dụ mẫu luôn dùng \loigiai)
%     \loigiaiTL[{đáp số}]{lời giải đầy đủ}
%     \TLfull -> in lời giải đầy đủ;  \TLdapso -> chỉ in một dòng "Đáp số";  \TLan -> ẩn hoàn toàn
%     (\CheDo là dethi thì mọi lời giải và đáp số đều ẩn)
\makeatletter
\newif\ifTLhien   \TLhientrue
\newif\ifTLdapso  \TLdapsofalse
\newcommand{\TLfull}{\TLhientrue\TLdapsofalse}
\newcommand{\TLdapso}{\TLhientrue\TLdapsotrue}
\newcommand{\TLan}{\TLhienfalse}
\newcommand{\loigiaiTL}[2][]{%
	\ifTLhien
		\ifTLdapso
			\def\@tlarg{#1}%
			\ifx\@tlarg\@empty\else
				\ifdefstring{\kindOPTN}{dethi}{}{\par\noindent{\bfseries\color{gold1}Đáp số:} #1\par}%
			\fi
		\else
			\loigiai{#2}%
		\fi
	\fi}
\makeatother
\TLdapso% mặc định: chỉ in đáp số bài tự luyện
%====================== BẢNG ĐÚNG/SAI (\choiceTFt) ======================
%  Hàng tiêu đề nền nâu đồng chữ trắng (Phát biểu | Đ | S), các hàng tô xen kẽ kem / trắng,
%  đường kẻ vàng kim, nhãn a) b) c) d) màu nâu đồng. Bản có đáp án: cột được chọn có vòng tròn
%  nâu đồng với dấu tích trắng (câu đúng tích ở cột Đúng, câu sai tích ở cột Sai).
%  Tiêu đề ghi đủ "Đúng | Sai": \OPTN{viettat=0} sau \input (cột Đúng/Sai rộng ra, phát biểu hẹp lại).
%  Dùng bảng gốc của ex_test: \def\BangDSgoc{1} trước \input{preamble-eduzy-v2}.
\providecommand{\BangDSgoc}{0}
\ifnum\BangDSgoc=0
\usepackage{colortbl,pifont}
\colorlet{ezDSke}{gold2!60!gold1}% đường kẻ bảng (nét 0,4pt như ex_test để không đổi chiều cao bảng)
\colorlet{ezDSle}{gold5!60}% hàng a), c)
\colorlet{ezDSchan}{white}% hàng b), d)
%--- dấu chọn ở cột Đúng/Sai
%    (bề rộng 0, căn giữa ô: không làm đổi bề rộng cột của ex_test, không báo tràn lề)
\newcommand{\EZdauchon}{\makebox[0pt]{\tikz[baseline=-1ex]\node[circle,fill=gold1,inner sep=0pt,
	minimum size=1.2em,text=white,font=\small]{\ding{51}};}}
\def\tickT{\EZdauchon}\def\tickF{\EZdauchon}
\AtBeginDocument{\sbox{\ccTFTab}{\EZdauchon}}
%--- nhãn a) b) c) d) màu nâu đồng (ở bản lời giải, nhãn câu đúng vẫn đóng khung như ex_test)
\renewcommand{\FalseTF}{\stepcounter{dapan}\textbf{\color{gold1}\DapAnTF\sepTF}\ignorespaces}
%--- trước khi dựng bảng: kẻ màu vàng kim (bề rộng cột giữ nguyên như ex_test); sau bảng trả lại màu đen
\newcount\EZhangTF
\newcommand{\EZbatdauTF}{%
	\global\EZhangTF=0\relax
	\arrayrulecolor{ezDSke}}
\newcommand{\EZketthucTF}{\arrayrulecolor{black}}
%--- hàng tiêu đề nền nâu đồng, chữ trắng
\makeatletter
\newcommand{\EZo@TF}[1]{\centering\bfseries\color{white}#1}
\newcommand{\EZtieudeTF}{%
	\ifcase\cotDS
		\ifnum\cotSTT=1\relax
			\def\rowHeader{\hline\rowcolor{gold1}\EZo@TF{TT}&\EZo@TF{\phatbieu}\tabularnewline}%
		\else
			\def\rowHeader{\hline\rowcolor{gold1}\EZo@TF{\phatbieu}\tabularnewline}%
		\fi
	\or
		\ifnum\cotSTT=1\relax
			\def\rowHeader{\hline\rowcolor{gold1}\EZo@TF{TT}&\EZo@TF{\phatbieu}&\EZo@TF{\Dung}\tabularnewline}%
		\else
			\def\rowHeader{\hline\rowcolor{gold1}\EZo@TF{\phatbieu}&\EZo@TF{\Dung}\tabularnewline}%
		\fi
	\or
		\ifnum\cotSTT=1\relax
			\def\rowHeader{\hline\rowcolor{gold1}\EZo@TF{TT}&\EZo@TF{\phatbieu}&\EZo@TF{\Dung}&\EZo@TF{\Sai}\tabularnewline}%
		\else
			\def\rowHeader{\hline\rowcolor{gold1}\EZo@TF{\phatbieu}&\EZo@TF{\Dung}&\EZo@TF{\Sai}\tabularnewline}%
		\fi
	\fi}
%--- tô xen kẽ từng hàng (gọi ngay sau \hline đầu mỗi hàng)
\newcommand{\EZmauhangTF}{%
	\noalign{\global\advance\EZhangTF\@ne}%
	\ifodd\EZhangTF\rowcolor{ezDSle}\else\rowcolor{ezDSchan}\fi}
%--- gắn vào \choiceTFt / \writeChoiceTFt của ex_test
\xpretocmd{\choiceTFt}{\EZbatdauTF}{}{\typeout{[preamble-eduzy-v2] Không gắn được bảng Đúng/Sai (đầu bảng)}}
\xpatchcmd{\choiceTFt}{\ifnum\dongPB=0\def\rowHeader{}\fi}{\EZtieudeTF\ifnum\dongPB=0\def\rowHeader{}\fi}
	{}{\typeout{[preamble-eduzy-v2] Không gắn được bảng Đúng/Sai (hàng tiêu đề)}}
\xpatchcmd{\writeChoiceTFt}{\hline\xdef\dotEX{\sdotEX}}{\hline\EZmauhangTF\xdef\dotEX{\sdotEX}}
	{}{\typeout{[preamble-eduzy-v2] Không gắn được bảng Đúng/Sai (tô hàng)}}
\xpatchcmd{\writeChoiceTFt}{{\end{longtable}}}{{\end{longtable}\EZketthucTF}}
	{}{\typeout{[preamble-eduzy-v2] Không gắn được bảng Đúng/Sai (cuối longtable)}}
\xpatchcmd{\writeChoiceTFt}{{\end{xtabular}\end{center}\medskip\noindent}}{{\end{xtabular}\end{center}\EZketthucTF\medskip\noindent}}
	{}{\typeout{[preamble-eduzy-v2] Không gắn được bảng Đúng/Sai (cuối xtabular)}}
\fi
%=====================================================================
\renewcommand{\parallel}{\mathrel{/\!\!/}}

%      \TLdapso                      % tài liệu: \TLfull | \TLdapso | \TLan
%      \begin{document} ... \end{document}
%  Logo Eduzy vẽ bằng TikZ ngay trong file (\eduzylogo), không cần file ảnh.
%  Biên dịch: pdfLaTeX (máy có gói vietnam) hoặc XeLaTeX.
%============================================================
\documentclass[12pt,a4paper,twoside]{article}%<-- twoside: bat dau trang/chan trang chan-le
\usepackage{iftex}
\usepackage{amsmath,amssymb,tikz,mathrsfs,tkz-tab}
\usepackage{graphicx}
\usepackage{pgfplots}% đồ thị hàm số
\pgfplotsset{compat=1.17}
%\usepackage{mathptmx}
\usepackage{mathpazo}
\ifXeTeX% XeLaTeX: font chữ có đủ dấu tiếng Việt, giữ nguyên font toán mathpazo
	\usepackage[no-math]{fontspec}
	\setmainfont{TeX Gyre Pagella}
	\newfontfamily\EZqag{TeX Gyre Adventor}[NFSSFamily=qag,Ligatures=TeX]
\fi
\usepackage[many]{tcolorbox}
\usetikzlibrary{math,through,calc,intersections,angles,quotes,shapes,shapes.geometric,arrows,patterns,matrix,chains,arrows.meta,decorations.shapes,decorations.fractals,decorations.markings,shadows,decorations.pathreplacing}
\usepackage[top=1.6cm, bottom=1.8cm, left=1.8cm, right=1.8cm]{geometry}% v2: top 1.6cm (bản cũ 1.5cm) để chữ không sát dải đầu trang
\providecommand{\CheDo}{loigiai}% loigiai | dethi | solcolor | color | book
\expandafter\usepackage\expandafter[\CheDo]{ex_test}
% \tenchuong: khai báo ở file chính (trước \input) hoặc \renewcommand sau \input
%============================================================
%  [GỘP TỪ eduzy-de.tex] EDUZY - MÀU NÂU ĐỒNG / VÀNG KIM
%  + ĐẦU TRANG, CHÂN TRANG KIỂU DẢI NGANG (v2), ĐỐI XỨNG TRANG CHẴN / TRANG LẺ.
%  (không cần file eduzy-de.tex riêng nữa)
%  Khối này phải đặt SAU \usepackage[...]{ex_test} và SAU \def\tenchuong{...}
%============================================================
\usepackage{lastpage}
\usepackage{needspace}
%--- DÀN TRANG ĐỀ THI: tách câu (THÊM 05/11/2026; chỉ thêm lệnh mới, KHÔNG đổi hành vi mặc định) ---
%  Hai công tắc, mặc định tắt (= preamble chạy y như cũ). Bật bằng cách đặt TRƯỚC %============================================================
%  PREAMBLE EDUZY v2 - dùng chung cho đề thi, bài học, chuyên đề (gói ex_test)
%  Khác bản preamble-eduzy.tex ở đầu trang, chân trang: dải nền nâu đồng tràn ngang,
%  logo + Eduzy.com.vn, ô chữ vàng kem cạnh xiên, chân trang "Trang (số)";
%  logo Eduzy chìm giữa mọi trang (tắt: \logochimtat, bật lại: \logochimbat).
%  Cách dùng ở file chính (xem mau-chuan.tex, mau-chuyen-de.tex):
%      \def\CheDo{loigiai}          % loigiai | dethi | solcolor | color | book
%      \def\tenchuong{<dòng chữ trên đầu trang>}
%      \def\thuonghieu{<chữ ở chân trang lẻ>}   % không khai báo thì dùng mặc định
%      \def\tacgia{<chữ ở chân trang chẵn>}      % không khai báo thì dùng mặc định
%      \def\khauhieu{<chữ cạnh logo ở đầu trang>}  % tùy chọn, mặc định Eduzy.com.vn
%      \input{preamble-eduzy-v2}
%      \TLdapso                      % tài liệu: \TLfull | \TLdapso | \TLan
%      \begin{document} ... \end{document}
%  Logo Eduzy vẽ bằng TikZ ngay trong file (\eduzylogo), không cần file ảnh.
%  Biên dịch: pdfLaTeX (máy có gói vietnam) hoặc XeLaTeX.
%============================================================
\documentclass[12pt,a4paper,twoside]{article}%<-- twoside: bat dau trang/chan trang chan-le
\usepackage{iftex}
\usepackage{amsmath,amssymb,tikz,mathrsfs,tkz-tab}
\usepackage{graphicx}
\usepackage{pgfplots}% đồ thị hàm số
\pgfplotsset{compat=1.17}
%\usepackage{mathptmx}
\usepackage{mathpazo}
\ifXeTeX% XeLaTeX: font chữ có đủ dấu tiếng Việt, giữ nguyên font toán mathpazo
	\usepackage[no-math]{fontspec}
	\setmainfont{TeX Gyre Pagella}
	\newfontfamily\EZqag{TeX Gyre Adventor}[NFSSFamily=qag,Ligatures=TeX]
\fi
\usepackage[many]{tcolorbox}
\usetikzlibrary{math,through,calc,intersections,angles,quotes,shapes,shapes.geometric,arrows,patterns,matrix,chains,arrows.meta,decorations.shapes,decorations.fractals,decorations.markings,shadows,decorations.pathreplacing}
\usepackage[top=1.6cm, bottom=1.8cm, left=1.8cm, right=1.8cm]{geometry}% v2: top 1.6cm (bản cũ 1.5cm) để chữ không sát dải đầu trang
\providecommand{\CheDo}{loigiai}% loigiai | dethi | solcolor | color | book
\expandafter\usepackage\expandafter[\CheDo]{ex_test}
% \tenchuong: khai báo ở file chính (trước \input) hoặc \renewcommand sau \input
%============================================================
%  [GỘP TỪ eduzy-de.tex] EDUZY - MÀU NÂU ĐỒNG / VÀNG KIM
%  + ĐẦU TRANG, CHÂN TRANG KIỂU DẢI NGANG (v2), ĐỐI XỨNG TRANG CHẴN / TRANG LẺ.
%  (không cần file eduzy-de.tex riêng nữa)
%  Khối này phải đặt SAU \usepackage[...]{ex_test} và SAU \def\tenchuong{...}
%============================================================
\usepackage{lastpage}
\usepackage{needspace}
%--- DÀN TRANG ĐỀ THI: tách câu (THÊM 05/11/2026; chỉ thêm lệnh mới, KHÔNG đổi hành vi mặc định) ---
%  Hai công tắc, mặc định tắt (= preamble chạy y như cũ). Bật bằng cách đặt TRƯỚC \input{preamble-eduzy-v2}
%  trong file đề (tools/dan_trang/dan_trang.py tự ghi khi đặt \NeedspaceCau):
%    \def\TachCauDS{1}   câu Đúng/Sai được tách: câu dẫn (kèm hình) ở cuối trang, bảng Phát biểu sang trang sau
%                        (câu dẫn và hình không tách nhau, bảng không tách giữa hai trang)
%    \def\TachCauLG{1}   bản lời giải (\CheDo khác dethi): phần đề của câu cao hơn 1/3 trang được phép tách
%  Lệnh dùng trong file đề (giá trị do tools/dan_trang/dan_trang.py đo ở bản dethi):
%    \NeedspaceCau[<cao câu dẫn + hình>]{<cao cả khối câu>}   ngay trên \begin{ex} (câu đầu phần: trên tiêu đề phần)
%        công tắc tắt: đúng bằng \Needspace{<cao cả khối câu>} như cũ;
%        \TachCauDS=1 và có [..] (câu Đúng/Sai được chọn tách): chỉ giữ câu dẫn + hình;
%        \TachCauLG=1 ở bản lời giải: đòi tối đa 1/3 \textheight.
%    \NeedspaceBang{<cao bảng Phát biểu>}%   ngay trước \choiceTFt: \TachCauDS=1 thì giữ trọn bảng trên một trang;
%        công tắc tắt: không làm gì.
\providecommand{\TachCauDS}{0}
\providecommand{\TachCauLG}{0}
\newlength{\EZnsCau}
\newcommand{\NeedspaceCau}[2][]{%
	\setlength{\EZnsCau}{#2}%
	\ifnum\TachCauDS=1 \if\relax\detokenize{#1}\relax\else\setlength{\EZnsCau}{#1}\fi\fi
	\ifnum\TachCauLG=1 \ifdefstring{\kindOPTN}{dethi}{}{\ifdim\EZnsCau>0.3333\textheight\setlength{\EZnsCau}{0.3333\textheight}\fi}\fi
	\Needspace{\EZnsCau}}
\newcommand{\NeedspaceBang}[1]{\ifnum\TachCauDS=1 \Needspace{#1}\fi}

%--- Thang màu: đậm -> nhạt
\definecolor{gold1}{RGB}{126,96,0}    % nâu đồng đậm  #7E6000
\definecolor{gold2}{RGB}{192,144,3}   % vàng kim      #C09003
\definecolor{gold3}{RGB}{255,215,99}  % vàng nhạt     #FFD763
\definecolor{gold4}{RGB}{255,231,153} % vàng rất nhạt #FFE799
\definecolor{gold5}{RGB}{255,243,205} % vàng kem      #FFF3CD
%--- Ghi đè tông xanh/đỏ cũ
\definecolor{myblue}{RGB}{126,96,0}
\colorlet{blue}{gold1}
\colorlet{cyan}{gold2}
\colorlet{teal}{gold2}

%--- Tên chương mặc định (file đề nên tự khai báo lại)
\providecommand{\tenchuong}{}
%--- Chân trang
\providecommand{\thuonghieu}{TOÁN 12 -- Eduzy.com.vn}
\providecommand{\tacgia}{Tô Gia Bảo -- 0353703783}

%--- Chữ cạnh logo ở đầu trang
\providecommand{\khauhieu}{Eduzy.com.vn}

%====================== THÔNG SỐ ĐẦU TRANG, CHÂN TRANG (v2) ======================
%  Dải nền đậm tràn ngang trang; ô chữ vàng kem có cạnh xiên, kèm 2 vạch xiên chuyển màu.
%  Đầu trang lẻ : logo + \khauhieu (trái)   | ô chữ \tenchuong (phải)
%  Chân trang lẻ: ô chữ \thuonghieu (trái)  | Trang (số) (phải)
%  Trang chẵn soi gương trang lẻ; chân trang chẵn ghi \tacgia.
%--- Màu: chỉ dùng thang vàng kim Eduzy
\colorlet{ezNen}{gold1}% dải nền đậm
\colorlet{ezVach1}{gold3}% vạch xiên sát dải nền (đậm hơn)
\colorlet{ezVach2}{gold4}% vạch xiên sát ô chữ (nhạt hơn)
\colorlet{ezO}{gold5}% ô chữ vàng kem
%--- Kích thước
\newlength{\EZdai}\setlength{\EZdai}{1.1cm}% chiều cao dải nền
\newlength{\EZo}\setlength{\EZo}{0.68cm}% chiều cao ô chữ
\newlength{\EZxien}\setlength{\EZxien}{0.55cm}% độ lệch ngang của cạnh xiên
\newlength{\EZvach}\setlength{\EZvach}{0.75cm}% bề rộng mỗi vạch xiên
\newlength{\EZkhe}\setlength{\EZkhe}{0.12cm}% khe trắng giữa vạch và ô chữ
\newlength{\EZdem}\setlength{\EZdem}{0.4cm}% khoảng đệm từ chữ tới cạnh xiên
\newlength{\EZle}\setlength{\EZle}{1.8cm}% lề trái/phải của chữ (bằng lề thân bài)
\newlength{\EZrongmin}\setlength{\EZrongmin}{4.5cm}% ô chữ rộng tối thiểu
\newlength{\EZrong}% bề rộng chữ trong ô (tính tự động)
%--- Font
\newcommand{\EZfontO}{\fontsize{10.5pt}{12pt}\selectfont\bfseries}% chữ trong ô
\newcommand{\EZfontKH}{\fontfamily{qag}\fontsize{11pt}{12pt}\selectfont\bfseries}% \khauhieu cạnh logo
\newcommand{\EZfontTrang}{\fontsize{11.5pt}{12pt}\selectfont\bfseries}% chữ "Trang" và số trang
%--- Đo bề rộng chữ trong ô
\newcommand{\EZdorong}[1]{%
	\settowidth{\EZrong}{\EZfontO #1}%
	\ifdim\EZrong<\EZrongmin\setlength{\EZrong}{\EZrongmin}\fi}
%--- Logo + chữ trắng \khauhieu (dựng sẵn một lần khi vào \begin{document})
\newsavebox{\EZlogobox}
\AtBeginDocument{\sbox{\EZlogobox}{%
		\raisebox{-0.5\height}{\eduzylogo[1cm]}\hspace{0.15cm}%
		\raisebox{-0.5\height}{\color{white}\EZfontKH\khauhieu}}}
%--- Hai vạch xiên + khe trắng. #1: điểm gốc ở mép dưới, #2: hướng (1 sang phải, -1 sang trái),
%    #3: hướng nghiêng (1 kiểu "/", -1 kiểu "\")
\newcommand{\EZvachxien}[3]{%
	\fill[white] (#1) --++({(#3)*\EZxien},\EZo) --++({(#2)*\EZkhe},0) --++({-(#3)*\EZxien},-\EZo) --cycle;
	\fill[ezVach2] ([xshift=(#2)*\EZkhe]#1) --++({(#3)*\EZxien},\EZo) --++({(#2)*\EZvach},0) --++({-(#3)*\EZxien},-\EZo) --cycle;
	\fill[ezVach1] ([xshift=(#2)*(\EZkhe+\EZvach)]#1) --++({(#3)*\EZxien},\EZo) --++({(#2)*\EZvach},0) --++({-(#3)*\EZxien},-\EZo) --cycle;
	\draw[white,line width=0.6pt] ([xshift=(#2)*(\EZkhe+\EZvach)]#1) --++({(#3)*\EZxien},\EZo);
}
%--- Chữ trong ô (căn giữa theo chiều cao ô, không lệch vì dấu tiếng Việt)
\tikzset{ezchu/.style={inner sep=0pt,text=ezNen,text height=0.72em,text depth=0.22em}}

%====================== LOGO CHÌM GIỮA TRANG ======================
%  Vẽ cùng đầu trang nên nằm dưới chữ của thân bài, có trên mọi trang.
%  Tắt: \logochimtat (từ trang đó trở đi), bật lại: \logochimbat.
%  Đổi cỡ, độ đậm: khai báo trước \input, ví dụ \def\EZchimcao{9cm} \def\EZchimdo{0.1}
\providecommand{\EZchimcao}{10cm}% chiều cao logo chìm
\providecommand{\EZchimdo}{0.12}% độ đậm (0 = trong suốt, 1 = rõ hẳn)
\newif\ifEZchim \EZchimtrue
\newcommand{\logochimtat}{\EZchimfalse}
\newcommand{\logochimbat}{\EZchimtrue}
%  Dựng sẵn logo mờ một lần khi vào \begin{document}; transparency group làm mờ cả logo
%  một lượt nên không lộ vệt ở chỗ các dải màu chồng nhau.
\newsavebox{\EZchimtho}
\newsavebox{\EZchimbox}
\AtBeginDocument{%
	\sbox{\EZchimtho}{\eduzylogo[\EZchimcao]}%
	\sbox{\EZchimbox}{\begin{tikzpicture}
			\begin{scope}[transparency group,opacity=\EZchimdo]
				\node[inner sep=0pt] {\usebox{\EZchimtho}};
			\end{scope}
	\end{tikzpicture}}}
\newcommand{\EZvechim}{%
	\ifEZchim
		\node[inner sep=0pt] at (current page.center) {\usebox{\EZchimbox}};
	\fi}

%====================== ĐẦU TRANG - TRANG LẺ ======================
%  (logo bên trái, ô chữ \tenchuong bên phải)
\newcommand{\myfancyheadOdd}{%
	\EZdorong{\tenchuong}%
	\begin{tikzpicture}[remember picture,overlay]
		\coordinate (T) at (current page.north west);
		\coordinate (P) at (current page.north east);
		%--- logo chìm giữa trang
		\EZvechim
		%--- dải nền
		\fill[ezNen] (T) rectangle ([yshift=-\EZdai]P);
		%--- ô chữ: X là góc dưới bên trái, cạnh trái xiên kiểu "/"
		\coordinate (X) at ([xshift=-\EZle-\EZrong-\EZdem-\EZxien,yshift=-\EZdai]P);
		\fill[ezO] (X) --++(\EZxien,\EZo) -- ([yshift=-\EZdai+\EZo]P) -- ([yshift=-\EZdai]P) --cycle;
		\EZvachxien{X}{-1}{1}
		\node[ezchu,anchor=east] at ([xshift=-\EZle,yshift=-\EZdai+0.5*\EZo]P) {\EZfontO\tenchuong};
		%--- logo + tên
		\node[inner sep=0pt,anchor=west] at ([xshift=\EZle,yshift=-0.5*\EZdai]T) {\usebox{\EZlogobox}};
	\end{tikzpicture}%
}

%====================== ĐẦU TRANG - TRANG CHẴN =====================
%  (soi gương: ô chữ \tenchuong bên trái, logo bên phải)
\newcommand{\myfancyheadEven}{%
	\EZdorong{\tenchuong}%
	\begin{tikzpicture}[remember picture,overlay]
		\coordinate (T) at (current page.north west);
		\coordinate (P) at (current page.north east);
		%--- logo chìm giữa trang
		\EZvechim
		%--- dải nền
		\fill[ezNen] (T) rectangle ([yshift=-\EZdai]P);
		%--- ô chữ: X là góc dưới bên phải, cạnh phải xiên kiểu "\"
		\coordinate (X) at ([xshift=\EZle+\EZrong+\EZdem+\EZxien,yshift=-\EZdai]T);
		\fill[ezO] (X) --++(-\EZxien,\EZo) -- ([yshift=-\EZdai+\EZo]T) -- ([yshift=-\EZdai]T) --cycle;
		\EZvachxien{X}{1}{-1}
		\node[ezchu,anchor=west] at ([xshift=\EZle,yshift=-\EZdai+0.5*\EZo]T) {\EZfontO\tenchuong};
		%--- logo + tên
		\node[inner sep=0pt,anchor=east] at ([xshift=-\EZle,yshift=-0.5*\EZdai]P) {\usebox{\EZlogobox}};
	\end{tikzpicture}%
}

%====================== CHÂN TRANG - TRANG LẺ ======================
%  (ô chữ \thuonghieu bên trái, "Trang (số)" bên phải)
\newcommand{\myfancyfootOdd}{%
	\EZdorong{\thuonghieu}%
	\begin{tikzpicture}[remember picture,overlay]
		\coordinate (T) at (current page.south west);
		\coordinate (P) at (current page.south east);
		%--- dải nền
		\fill[ezNen] (T) rectangle ([yshift=\EZdai]P);
		%--- ô chữ sát mép trên dải: X là góc dưới bên phải, cạnh phải xiên kiểu "/"
		\coordinate (X) at ([xshift=\EZle+\EZrong+\EZdem,yshift=\EZdai-\EZo]T);
		\fill[ezO] ([yshift=\EZdai-\EZo]T) -- (X) --++(\EZxien,\EZo) -- ([yshift=\EZdai]T) --cycle;
		\EZvachxien{X}{1}{1}
		\node[ezchu,anchor=west] at ([xshift=\EZle,yshift=\EZdai-0.5*\EZo]T) {\EZfontO\thuonghieu};
		%--- số trang
		\node[circle,draw=white,line width=0.9pt,text=white,inner sep=0pt,minimum size=0.62cm,anchor=east]
		(so) at ([xshift=-\EZle,yshift=0.5*\EZdai]P) {\EZfontTrang\thepage};
		\node[inner sep=0pt,text=white,anchor=base east] (tr) at ([xshift=-0.15cm]so.base west) {\EZfontTrang Trang};
	\end{tikzpicture}%
}

%====================== CHÂN TRANG - TRANG CHẴN ====================
%  (soi gương: "Trang (số)" bên trái, ô chữ \tacgia bên phải)
\newcommand{\myfancyfootEven}{%
	\EZdorong{\tacgia}%
	\begin{tikzpicture}[remember picture,overlay]
		\coordinate (T) at (current page.south west);
		\coordinate (P) at (current page.south east);
		%--- dải nền
		\fill[ezNen] (T) rectangle ([yshift=\EZdai]P);
		%--- ô chữ sát mép trên dải: X là góc dưới bên trái, cạnh trái xiên kiểu "\"
		\coordinate (X) at ([xshift=-\EZle-\EZrong-\EZdem,yshift=\EZdai-\EZo]P);
		\fill[ezO] ([yshift=\EZdai-\EZo]P) -- (X) --++(-\EZxien,\EZo) -- ([yshift=\EZdai]P) --cycle;
		\EZvachxien{X}{-1}{-1}
		\node[ezchu,anchor=east] at ([xshift=-\EZle,yshift=\EZdai-0.5*\EZo]P) {\EZfontO\tacgia};
		%--- số trang
		%    (vòng tròn ẩn so0 nằm giữa dải để lấy đường chân chữ, giống trang lẻ)
		\node[circle,inner sep=0pt,minimum size=0.62cm] (so0) at ([yshift=0.5*\EZdai]T) {\EZfontTrang\phantom{\thepage}};
		\node[inner sep=0pt,text=white,anchor=base west] (tr) at ($(so0.base -| T)+(\EZle,0)$) {\EZfontTrang Trang};
		\node[circle,draw=white,line width=0.9pt,text=white,inner sep=0pt,minimum size=0.62cm,anchor=base west]
		(so) at ([xshift=0.15cm]tr.base east) {\EZfontTrang\thepage};
	\end{tikzpicture}%
}

%--- Giữ tương thích với các file cũ đang gọi \myfancyhead, \myfancyfoot
\newcommand{\myfancyhead}{\myfancyheadOdd}
\newcommand{\myfancyfoot}{\myfancyfootOdd}

%======================= GẮN VÀO TRANG =============================
\pagestyle{fancy}
\fancyhf{}
\fancyhead[LO]{\myfancyheadOdd}
\fancyhead[RE]{\myfancyheadEven}
\fancyfoot[LO]{\myfancyfootOdd}
\fancyfoot[RE]{\myfancyfootEven}
\renewcommand{\headrulewidth}{0pt}
\renewcommand{\footrulewidth}{0pt}
\fancypagestyle{plain}{%
	\fancyhf{}
	\fancyhead[LO]{\myfancyheadOdd}
	\fancyhead[RE]{\myfancyheadEven}
	\fancyfoot[LO]{\myfancyfootOdd}
	\fancyfoot[RE]{\myfancyfootEven}
	\renewcommand{\headrulewidth}{0pt}
	\renewcommand{\footrulewidth}{0pt}
}
%=================== HẾT PHẦN GỘP TỪ eduzy-de.tex ===================
\newcommand{\indam}[1]{
	{\fontsize{12pt}{9pt}\sffamily \selectfont\bfseries\color{gold1} #1}
}
\newcommand{\ind}[1]{
	{\fontsize{14pt}{9pt}\fontfamily{qag} \selectfont\bfseries\color{gold1} #1}
}
\renewcommand{\FalseEX}{\stepcounter{dapan}\textbf{\selectfont\color{gold1}\Alph{dapan}.}\ignorespaces}
\renewtheorem{ex}{\color{gold1}Câu}
\renewcommand{\baselinestretch}{1.1}
\newcommand{\hoac}[1]{
	\left[\begin{aligned}#1\end{aligned}\right.}
\newcommand{\heva}[1]{
	\left\{\begin{aligned}#1\end{aligned}\right.}
\usepackage{esvect}
\def\vec{\vv}
\def\vt{\vv}
\def\overrightarrow{\vv}
%====================== LOGO EDUZY VẼ BẰNG TIKZ ======================
%  Đường viền dò từ logo gốc; dải màu vàng kim đo trực tiếp trên ảnh gốc
%  (nghiêng 21 độ, tô bằng 80 dải hẹp nên hiển thị giống nhau trên mọi trình xem PDF).
%  Dùng: \eduzylogo          -> cao 2.4cm (mặc định)
%        \eduzylogo[1.5cm]   -> cao 1.5cm
\definecolor{ezText}{RGB}{245,200,46}
\newcommand{\eduzylogo}[1][2.4cm]{%
	\resizebox{!}{#1}{%
		\begin{tikzpicture}[x=0.01cm,y=0.01cm]
			\path[use as bounding box] (0,0) rectangle (224,-224);
			% mũ tốt nghiệp + vương miện
			\begin{scope}
				\clip
				(119.4,-176.1)..controls(118.9,-176.1)and(116.8,-176)..(114.6,-176)
				..controls(101.4,-176)and(85.8,-175.9)..(83.7,-175.7)
				..controls(81.8,-175.6)and(81,-175.6)..(80.2,-175.7)
				..controls(78.8,-175.9)and(77.7,-175.9)..(76.2,-175.7)
				..controls(74.9,-175.5)and(70.8,-175.5)..(64.3,-175.7)
				..controls(59.1,-175.9)and(58.3,-175.7)..(58.6,-173.7)
				..controls(58.6,-172.8)and(59.6,-169.6)..(60.3,-168)
				..controls(60.8,-166.9)and(61.4,-165.8)..(62.3,-165)
				..controls(63.1,-164.2)and(64.3,-163.4)..(65,-163.2)
				..controls(65.5,-163.1)and(72.3,-163)..(76.6,-163.1)
				..controls(77.9,-163.1)and(80.1,-163.1)..(81.4,-163)
				..controls(82.8,-162.9)and(86.3,-162.9)..(89.2,-162.9)
				..controls(93,-162.9)and(94.9,-162.8)..(96,-162.7)
				..controls(97.2,-162.5)and(101,-162.5)..(114,-162.4)
				..controls(123,-162.3)and(134.9,-162.2)..(140.2,-162.1)
				..controls(150.1,-162)and(152.2,-162.1)..(153.6,-162.4)
				..controls(154.4,-162.6)and(157.1,-163.8)..(158.1,-164.4)
				..controls(158.8,-164.8)and(159.7,-165.7)..(159.9,-166.3)
				..controls(160,-166.5)and(160.2,-167.1)..(160.2,-167.6)
				..controls(160.3,-168)and(160.6,-169.4)..(160.9,-170.6)
				..controls(161.2,-171.7)and(161.5,-173.1)..(161.6,-173.6)
				..controls(161.8,-174.4)and(161.8,-174.5)..(161.5,-174.9)
				..controls(161.2,-175.6)and(160.8,-175.9)..(159.9,-176)
				..controls(159.2,-176.1)and(154.1,-176.2)..(151.5,-176.2)
				..controls(148.1,-176.1)and(122.6,-176.1)..(121.6,-176.2)
				..controls(120.8,-176.2)and(119.8,-176.2)..(119.4,-176.1)--cycle
				(130.4,-160)..controls(129,-160)and(119.8,-160)..(110.1,-159.9)
				..controls(100.3,-159.8)and(88.5,-159.7)..(83.9,-159.6)
				..controls(77.8,-159.5)and(75.3,-159.5)..(74.5,-159.6)
				..controls(73.9,-159.7)and(72.5,-159.8)..(71.3,-159.8)
				..controls(69.6,-159.8)and(69,-159.8)..(68.5,-159.6)
				..controls(67.6,-159.3)and(66.8,-158.6)..(66.2,-157.6)
				..controls(66,-157.3)and(65.6,-156.5)..(65.3,-156.1)
				..controls(64.6,-155.1)and(64.5,-154.9)..(63.7,-152.9)
				..controls(63.2,-151.6)and(62.9,-150.9)..(61.8,-149.6)
				..controls(61.5,-149.2)and(61.1,-148.5)..(60.9,-147.9)
				..controls(60.7,-147.4)and(60.3,-146.6)..(60.1,-146.1)
				..controls(59.9,-145.5)and(59.5,-144.9)..(59,-144.3)
				..controls(58.5,-143.7)and(58.2,-143.2)..(57.9,-142.5)
				..controls(57.7,-141.9)and(57.3,-141.1)..(57,-140.7)
				..controls(56.6,-140.3)and(56.1,-139.3)..(55.8,-138.5)
				..controls(55.5,-137.8)and(55.1,-136.8)..(54.9,-136.5)
				..controls(54.7,-136.1)and(54.3,-135.2)..(53.9,-134.5)
				..controls(53.6,-133.8)and(53,-132.7)..(52.6,-132)
				..controls(51.6,-130.5)and(51.4,-130.2)..(51.1,-129)
				..controls(50.6,-127.5)and(50.6,-127.4)..(49.7,-126.3)
				..controls(49.2,-125.7)and(48.6,-124.8)..(48.3,-124.2)
				..controls(48,-123.6)and(47.5,-122.7)..(47.3,-122.2)
				..controls(47,-121.8)and(46.5,-120.7)..(46.1,-119.8)
				..controls(45.8,-118.9)and(45.2,-117.8)..(44.9,-117.2)
				..controls(44.2,-116.1)and(43.3,-114.4)..(42.3,-112.2)
				..controls(41.8,-111.3)and(41.3,-110.4)..(41.1,-110.1)
				..controls(40.5,-109.4)and(40.3,-109)..(40,-107.8)
				..controls(39.8,-107)and(39.6,-106.5)..(39.2,-106)
				..controls(38.9,-105.6)and(38.4,-104.7)..(38,-104)
				..controls(37.7,-103.4)and(37.2,-102.5)..(37,-102.1)
				..controls(36.7,-101.7)and(36.3,-100.8)..(36,-100.1)
				..controls(35.6,-98.9)and(35.3,-98.4)..(34.6,-97.6)
				..controls(34.2,-97.1)and(34.2,-96.3)..(34.6,-96)
				..controls(35,-95.7)and(35.5,-95.9)..(36.2,-96.7)
				..controls(36.5,-97)and(37.2,-97.5)..(37.6,-97.7)..controls(38,-98)and(38.6,-98.4)..(39.1,-98.9)
				..controls(39.6,-99.4)and(40.3,-100)..(40.6,-100.3)
				..controls(41.3,-100.8)and(42.4,-101.8)..(43.7,-102.9)
				..controls(44,-103.2)and(44.6,-103.7)..(45,-103.9)
				..controls(45.4,-104.2)and(45.9,-104.6)..(46.1,-104.8)
				..controls(47.1,-105.9)and(48.8,-107.1)..(50.5,-107.9)
				..controls(51.8,-108.6)and(52.3,-109)..(53.4,-110)
				..controls(54.8,-111.2)and(55.6,-111.6)..(57,-112.1)
				..controls(57.5,-112.3)and(58.5,-112.8)..(59.4,-113.3)
				..controls(60.2,-113.7)and(61.3,-114.3)..(61.9,-114.5)
				..controls(62.4,-114.7)and(63.3,-115.1)..(63.8,-115.3)
				..controls(64.4,-115.6)and(65.4,-115.9)..(66.2,-116.1)
				..controls(67,-116.3)and(68.4,-116.6)..(69.2,-116.9)
				..controls(72.3,-117.9)and(72.8,-118)..(78,-118.1)
				..controls(84,-118.3)and(86.1,-118.1)..(88.1,-117.1)
				..controls(89.4,-116.5)and(93,-115)..(93.7,-114.9)
				..controls(94.8,-114.6)and(95.4,-114.2)..(95.9,-113.4)
				..controls(96.2,-113)and(96.7,-112.5)..(97.1,-112.2)
				..controls(98,-111.6)and(101.3,-108.4)..(102.1,-107.2)
				..controls(102.4,-106.8)and(102.9,-106.2)..(103.2,-105.8)
				..controls(103.5,-105.5)and(104,-104.8)..(104.2,-104.2)
				..controls(104.5,-103.6)and(105,-102.8)..(105.3,-102.4)
				..controls(105.6,-102)and(106.1,-101.2)..(106.5,-100.6)
				..controls(107.1,-99.5)and(108.6,-97.5)..(109.5,-96.5)
				..controls(109.9,-96.1)and(110.1,-96)..(110.4,-96)
				..controls(111.1,-96)and(111.6,-96.5)..(111.8,-97.6)
				..controls(112,-98.6)and(112.3,-99.3)..(112.8,-99.9)
				..controls(113.7,-100.9)and(115.2,-103.1)..(115.7,-104.2)
				..controls(116.1,-105)and(116.5,-105.5)..(116.9,-106)
				..controls(117.3,-106.4)and(117.8,-106.9)..(118,-107.3)
				..controls(118.2,-107.6)and(118.6,-108)..(119,-108.3)
				..controls(119.6,-108.8)and(121.1,-110.3)..(121.8,-111.2)
				..controls(122.1,-111.7)and(122.5,-112)..(123.4,-112.4)
				..controls(124.8,-113.2)and(125.8,-113.8)..(127,-114.6)
				..controls(127.5,-114.9)and(128.1,-115.3)..(128.2,-115.4)
				..controls(128.4,-115.5)and(129.3,-115.8)..(130.2,-116)
				..controls(131.1,-116.2)and(132.5,-116.6)..(133.3,-116.9)
				..controls(135.7,-117.9)and(136.1,-117.9)..(142.5,-117.8)
				..controls(145.8,-117.8)and(148.4,-117.7)..(148.8,-117.6)
				..controls(149.2,-117.6)and(150.2,-117.2)..(151,-116.8)
				..controls(152,-116.4)and(152.8,-116)..(153.6,-115.9)
				..controls(154.5,-115.7)and(155.2,-115.4)..(156.3,-114.9)
				..controls(157.1,-114.5)and(158.2,-114)..(158.7,-113.8)
				..controls(159.1,-113.7)and(159.8,-113.3)..(160.2,-113)
				..controls(160.6,-112.7)and(161.3,-112.4)..(161.7,-112.2)
				..controls(162.7,-111.9)and(163.2,-111.7)..(164.1,-111.1)
				..controls(164.5,-110.8)and(165.3,-110.4)..(165.9,-110)
				..controls(166.5,-109.7)and(167.3,-109.2)..(167.6,-108.9)
				..controls(167.9,-108.6)and(168.6,-108.2)..(169,-108)
				..controls(169.4,-107.8)and(170.2,-107.3)..(170.7,-106.9)
				..controls(171.2,-106.5)and(172.1,-105.8)..(172.7,-105.4)
				..controls(173.2,-105)and(174.1,-104.4)..(174.5,-104)
				..controls(174.9,-103.7)and(175.7,-103.1)..(176.2,-102.6)
				..controls(176.7,-102.2)and(177.4,-101.5)..(177.9,-101.2)
				..controls(178.3,-100.8)and(178.8,-100.4)..(179,-100.2)
				..controls(179.2,-100)and(179.6,-99.6)..(180,-99.3)
				..controls(180.4,-99.1)and(180.9,-98.6)..(181.2,-98.4)
				..controls(181.5,-98.1)and(182.1,-97.7)..(182.5,-97.5)
				..controls(183.2,-97.2)and(183.5,-96.9)..(184.3,-95.9)
				..controls(184.6,-95.6)and(184.6,-95.6)..(184.8,-95.7)
				..controls(185,-95.8)and(185.1,-96)..(185.1,-96.2)
				..controls(185.1,-96.5)and(185,-96.7)..(184.6,-97.2)
				..controls(184.1,-97.8)and(184,-97.9)..(183.9,-98.8)
				..controls(183.7,-100.3)and(183.3,-101.3)..(182.9,-101.8)
				..controls(182.6,-102.1)and(182.3,-102.6)..(182.2,-102.9)
				..controls(182,-103.3)and(181.6,-104)..(181.2,-104.6)
				..controls(180.9,-105.2)and(180.5,-106.1)..(180.4,-106.6)
				..controls(179.9,-108)and(179.7,-108.5)..(179,-109.4)
				..controls(178.6,-109.8)and(178.2,-110.5)..(177.9,-111)
				..controls(177.7,-111.5)and(177.3,-112.3)..(177,-112.7)
				..controls(176.6,-113.2)and(176.4,-113.7)..(176.3,-114.2)
				..controls(176.2,-114.7)and(175.9,-115.2)..(175.4,-116)
				..controls(175,-116.6)and(174.4,-117.7)..(174,-118.5)
				..controls(173.7,-119.2)and(173.3,-120.1)..(173,-120.5)
				..controls(172.8,-120.9)and(172.4,-121.8)..(172.1,-122.5)
				..controls(171.8,-123.2)and(171.3,-124.1)..(171,-124.6)
				..controls(170.8,-125.1)and(170.3,-126)..(169.9,-126.6)
				..controls(169.6,-127.3)and(169.2,-128.1)..(169,-128.3)
				..controls(168.8,-128.6)and(168.5,-129.3)..(168.3,-129.8)
				..controls(168.1,-130.3)and(167.8,-131.1)..(167.5,-131.6)
				..controls(166.9,-132.5)and(164.2,-137.9)..(164.2,-138.1)
				..controls(164.2,-138.1)and(163.9,-138.8)..(163.4,-139.7)
				..controls(162.9,-140.5)and(162.4,-141.5)..(162.3,-141.9)
				..controls(161.9,-143)and(161.6,-143.4)..(160.9,-144.1)
				..controls(160.4,-144.6)and(160.3,-144.8)..(160.1,-145.4)
				..controls(160,-145.8)and(159.6,-146.8)..(159.2,-147.8)
				..controls(158.7,-148.8)and(158.1,-150)..(157.9,-150.6)
				..controls(157.6,-151.1)and(157.2,-151.9)..(156.9,-152.4)
				..controls(156.5,-152.8)and(156.1,-153.7)..(155.8,-154.4)
				..controls(155.6,-155)and(155.2,-155.9)..(154.9,-156.3)
				..controls(154.7,-156.7)and(154.4,-157.4)..(154.2,-157.9)
				..controls(154.1,-158.4)and(153.8,-159)..(153.5,-159.3)
				..controls(152.8,-160.1)and(153,-160.1)..(142.4,-160.1)
				..controls(137.2,-160.1)and(131.8,-160.1)..(130.4,-160)--cycle
				(141.5,-103.7)..controls(140.8,-103)and(139.8,-102.1)..(139.2,-101.6)
				..controls(138.5,-101.2)and(137.9,-100.6)..(137.8,-100.4)
				..controls(137,-99.5)and(134,-96.5)..(133.2,-96)
				..controls(132.7,-95.7)and(131.8,-94.9)..(131.1,-94.3)
				..controls(129.4,-92.7)and(127.9,-91.6)..(126.6,-91.1)
				..controls(125.2,-90.5)and(124.1,-90)..(122.9,-89.1)
				..controls(122.1,-88.6)and(121.3,-88.2)..(119.9,-87.7)
				..controls(118.8,-87.3)and(117.5,-86.8)..(116.9,-86.5)
				..controls(114.3,-85.5)and(113.7,-85.4)..(109.7,-85.5)
				..controls(105.3,-85.7)and(103.9,-85.9)..(101.6,-87)
				..controls(100.7,-87.4)and(99.4,-87.8)..(97.9,-88)
				..controls(97,-88.2)and(96.6,-88.3)..(95.8,-88.9)
				..controls(95.4,-89.2)and(94.7,-89.6)..(94.2,-89.8)
				..controls(92.9,-90.3)and(92.3,-90.7)..(90.9,-92.1)
				..controls(90.2,-92.7)and(89.4,-93.3)..(89,-93.5)..controls(87,-94.5)and(85.2,-96)..(83.6,-97.9)
				..controls(83.3,-98.2)and(82.5,-98.9)..(81.8,-99.4)
				..controls(80.7,-100.1)and(80.5,-100.3)..(80,-101.1)
				..controls(78.7,-103.1)and(77.2,-103.9)..(76.4,-103)
				..controls(75.7,-102.2)and(75.7,-103.2)..(75.7,-83.6)
				..controls(75.6,-73.6)and(75.6,-64.8)..(75.5,-64)
				..controls(75.5,-63.2)and(75.5,-61.4)..(75.6,-59.8)--(75.7,-57.1)--(76,-56.8)
				..controls(76.8,-56.2)and(78,-56.3)..(79.8,-57.2)
				..controls(80.4,-57.5)and(81.3,-57.8)..(81.8,-57.9)
				..controls(82.3,-58.1)and(83,-58.4)..(83.6,-58.8)..controls(84.3,-59.2)and(85,-59.5)..(86,-59.8)
				..controls(86.8,-60.1)and(88,-60.6)..(88.7,-60.9)
				..controls(89.3,-61.2)and(90.6,-61.7)..(91.4,-62)
				..controls(92.2,-62.3)and(93.7,-62.9)..(94.6,-63.3)
				..controls(95.6,-63.7)and(96.8,-64.1)..(97.4,-64.3)
				..controls(98,-64.5)and(99,-64.9)..(99.7,-65.2)
				..controls(100.4,-65.6)and(101.5,-66)..(102.3,-66.2)
				..controls(104.2,-66.8)and(105,-67.1)..(106.2,-67.8)
				..controls(107.9,-68.7)and(110.2,-69.4)..(111.3,-69.2)
				..controls(112.5,-69)and(115,-68.2)..(115.6,-67.8)
				..controls(117.3,-66.7)and(117.6,-66.6)..(119.5,-66.2)
				..controls(121.9,-65.8)and(122.1,-65.7)..(123,-65.1)
				..controls(124.1,-64.4)and(124.4,-64.3)..(126,-63.9)
				..controls(127,-63.7)and(127.7,-63.4)..(128.4,-63.1)
				..controls(129,-62.8)and(129.7,-62.5)..(130.6,-62.4)
				..controls(132.1,-62.1)and(134,-61.5)..(135.2,-61)
				..controls(135.7,-60.8)and(136.7,-60.5)..(137.3,-60.3)
				..controls(138.4,-60)and(139.3,-59.7)..(141.1,-59)
				..controls(141.6,-58.7)and(142,-58.6)..(142.5,-58.6)
				..controls(143.2,-58.6)and(143.3,-58.7)..(143.6,-59)
				..controls(144,-59.5)and(144.1,-59.9)..(144,-60.8)
				..controls(144,-61.2)and(144,-71.1)..(144,-82.8)--(143.9,-104.1)--(143.7,-104.5)
				..controls(143.5,-104.7)and(143.3,-104.8)..(143,-104.9)
				..controls(142.6,-104.9)and(142.5,-104.8)..(141.5,-103.7)--cycle
				(27.8,-91.4)..controls(27.1,-90.9)and(26.3,-90)..(26.1,-89.5)
				..controls(26,-89.2)and(25.8,-88.6)..(25.7,-87.9)
				..controls(25.4,-86.2)and(25.6,-84.8)..(26.5,-83.5)
				..controls(27.4,-82.3)and(28.3,-82)..(29.4,-82.3)
				..controls(31.1,-82.8)and(32.1,-84.9)..(32,-87.4)
				..controls(31.9,-89.3)and(31.2,-90.6)..(29.9,-91.4)
				..controls(29.1,-91.9)and(28.6,-91.9)..(27.8,-91.4)--cycle
				(190.1,-91.3)..controls(189.7,-91.2)and(189.3,-90.9)..(189.1,-90.7)
				..controls(188.3,-89.8)and(187.6,-87.7)..(187.6,-86.2)
				..controls(187.6,-85.5)and(187.7,-85.2)..(188.2,-84.4)
				..controls(189.1,-82.6)and(190.1,-81.8)..(191.3,-82.1)
				..controls(192,-82.3)and(193.5,-83.7)..(193.8,-84.4)
				..controls(194.1,-85.1)and(194.2,-87.1)..(194,-88.3)
				..controls(193.8,-89.8)and(193.3,-90.5)..(192.3,-91.3)
				..controls(191.7,-91.7)and(190.9,-91.7)..(190.1,-91.3)--cycle
				(173.7,-83.4)..controls(172.2,-82.5)and(172,-82.5)..(170.6,-82.4)
				..controls(169.7,-82.4)and(168.8,-82.4)..(168.3,-82.5)
				..controls(167.2,-82.7)and(166.9,-82.7)..(166.6,-82.2)
				..controls(166.3,-81.8)and(166.3,-81)..(166.8,-80)
				..controls(167.2,-78.9)and(167.4,-77.4)..(167.7,-74.1)
				..controls(167.8,-72.8)and(167.9,-71.8)..(167.9,-71.7)
				..controls(168,-71.7)and(168,-70.2)..(168,-68.4)..controls(168,-65.4)and(168,-64.9)..(167.7,-63)
				..controls(167.3,-60.5)and(167.1,-59.9)..(166.4,-58.4)
				..controls(166.1,-57.7)and(165.8,-56.9)..(165.7,-56.6)
				..controls(165.7,-56.2)and(165.4,-55.6)..(165.2,-55.2)
				..controls(164.6,-53.8)and(161.9,-50.7)..(161,-50.3)
				..controls(159.8,-49.7)and(157.9,-50)..(155.2,-51.2)
				..controls(154.7,-51.4)and(153.4,-51.8)..(152.2,-52.1)
				..controls(150.8,-52.4)and(149.4,-52.9)..(148.6,-53.2)
				..controls(147.8,-53.6)and(146.7,-54)..(146,-54.2)
				..controls(145.4,-54.4)and(144.5,-54.7)..(144,-54.9)
				..controls(143.6,-55.2)and(142.7,-55.5)..(141.8,-55.7)
				..controls(141,-55.9)and(139.6,-56.3)..(138.7,-56.6)
				..controls(137.8,-57)and(136.3,-57.4)..(135.3,-57.6)
				..controls(133.8,-57.9)and(133.2,-58.1)..(132.1,-58.7)
				..controls(131.1,-59.2)and(130.3,-59.4)..(129.4,-59.6)
				..controls(128.7,-59.8)and(127.8,-60)..(127.4,-60.1)
				..controls(127,-60.3)and(126,-60.5)..(125.1,-60.7)
				..controls(123.2,-61.1)and(121.7,-61.5)..(120.2,-62.1)
				..controls(119.7,-62.3)and(118.1,-62.8)..(116.7,-63.2)
				..controls(114.6,-63.9)and(114,-64.1)..(113,-64.7)
				..controls(111.9,-65.4)and(111.8,-65.4)..(111,-65.5)
				..controls(109.9,-65.5)and(109.3,-65.4)..(107.7,-64.7)
				..controls(107,-64.4)and(106,-64)..(105.3,-63.9)
				..controls(103.5,-63.5)and(102.2,-63)..(101.2,-62.5)
				..controls(100.1,-61.8)and(99.3,-61.5)..(97.8,-61.1)
				..controls(97.2,-61)and(96.2,-60.7)..(95.7,-60.4)
				..controls(95.1,-60.1)and(94.2,-59.8)..(93.5,-59.6)
				..controls(92.9,-59.4)and(91.7,-59)..(91,-58.7)..controls(90.2,-58.4)and(88.9,-58)..(88,-57.7)
				..controls(87.2,-57.5)and(86.1,-57.1)..(85.5,-56.8)
				..controls(85,-56.6)and(84.1,-56.2)..(83.5,-56)..controls(82.9,-55.7)and(81.8,-55.2)..(81,-54.9)
				..controls(80.1,-54.5)and(79.1,-54.2)..(78.7,-54.1)
				..controls(77.9,-53.8)and(77.3,-53.6)..(76.2,-52.9)
				..controls(75.7,-52.6)and(75.2,-52.4)..(74.1,-52.2)
				..controls(72.5,-51.9)and(72.4,-51.8)..(71.3,-51.1)
				..controls(69.9,-50.1)and(69.4,-49.9)..(67.9,-49.8)
				..controls(67,-49.7)and(66.6,-49.6)..(66.3,-49.4)
				..controls(66.1,-49.3)and(65.7,-49)..(65.4,-48.8)
				..controls(65.1,-48.7)and(64.2,-48.4)..(63.6,-48.2)
				..controls(62.9,-48)and(61.7,-47.6)..(60.9,-47.2)
				..controls(60.1,-46.8)and(58.9,-46.4)..(58.2,-46.2)
				..controls(57.6,-46)and(56.8,-45.7)..(56.5,-45.5)
				..controls(54.9,-44.6)and(54.6,-44.5)..(53.2,-44.2)
				..controls(51.7,-43.8)and(51,-43.6)..(49.9,-42.9)
				..controls(49.5,-42.6)and(48.9,-42.4)..(47.5,-42)
				..controls(45.7,-41.6)and(45.6,-41.5)..(44.3,-40.6)
				..controls(43.5,-40.1)and(42.8,-39.6)..(42.7,-39.4)
				..controls(42.4,-38.7)and(43,-38.1)..(44.5,-37.6)..controls(44.9,-37.5)and(46,-37)..(46.9,-36.6)
				..controls(47.9,-36.1)and(48.8,-35.8)..(49.6,-35.6)
				..controls(50.5,-35.4)and(51.3,-35.2)..(52.4,-34.6)
				..controls(53.3,-34.1)and(54.5,-33.7)..(54.9,-33.5)
				..controls(55.4,-33.4)and(56.4,-33)..(57.2,-32.6)
				..controls(58.6,-32)and(60.1,-31.5)..(61.6,-31.2)
				..controls(62.3,-31.1)and(62.7,-31)..(63.2,-30.7)
				..controls(64.1,-30)and(64.3,-29.9)..(65.3,-29.7)
				..controls(65.7,-29.6)and(66.5,-29.2)..(67.1,-28.9)
				..controls(67.6,-28.6)and(68.8,-28.2)..(69.6,-27.9)
				..controls(70.5,-27.7)and(71.6,-27.3)..(72.1,-27)
				..controls(72.6,-26.7)and(73.8,-26.2)..(74.8,-25.9)
				..controls(75.7,-25.5)and(76.7,-25.1)..(77.1,-25)
				..controls(77.4,-24.8)and(78.5,-24.4)..(79.6,-24.1)
				..controls(80.6,-23.8)and(82,-23.2)..(82.8,-22.9)
				..controls(83.5,-22.5)and(84.8,-22)..(85.8,-21.8)
				..controls(86.7,-21.5)and(87.8,-21.1)..(88.1,-20.9)
				..controls(88.4,-20.7)and(89.3,-20.3)..(90.1,-20.1)
				..controls(90.8,-19.8)and(91.9,-19.4)..(92.4,-19.1)
				..controls(93.5,-18.5)and(94.4,-18.2)..(95.4,-18.1)
				..controls(95.7,-18)and(96.4,-17.8)..(96.9,-17.5)
				..controls(97.5,-17.3)and(98.4,-16.9)..(99,-16.6)
				..controls(99.7,-16.3)and(100.9,-15.9)..(101.9,-15.6)
				..controls(102.9,-15.3)and(104.4,-14.7)..(105.3,-14.3)
				..controls(108.1,-13.2)and(108.5,-13.1)..(110.1,-13.1)
				..controls(111.4,-13.1)and(111.5,-13.1)..(112.4,-13.5)
				..controls(113.2,-14)and(115,-14.8)..(116.6,-15.4)
				..controls(116.9,-15.5)and(117.9,-15.8)..(118.8,-16.1)
				..controls(120.9,-16.6)and(121.5,-16.9)..(122.6,-17.9)
				..controls(123,-18.3)and(123.7,-18.8)..(123.9,-18.9)
				..controls(124.6,-19.4)and(126,-19.8)..(127.9,-20.2)
				..controls(129.6,-20.5)and(129.8,-20.6)..(131.1,-21.4)
				..controls(131.5,-21.7)and(132.4,-22.1)..(133.1,-22.3)
				..controls(136.1,-23.2)and(138.3,-24.1)..(140.2,-25)
				..controls(141,-25.5)and(141.5,-25.6)..(142.5,-25.8)
				..controls(145.1,-26.3)and(146,-26.6)..(149.4,-28.1)
				..controls(151.3,-28.9)and(153.5,-29.8)..(154.3,-30.1)
				..controls(155.1,-30.4)and(156.2,-30.8)..(156.6,-31)
				..controls(157.1,-31.2)and(158.2,-31.6)..(159.1,-31.8)
				..controls(160.6,-32.2)and(160.9,-32.3)..(161.7,-32.9)
				..controls(162.4,-33.3)and(162.9,-33.6)..(163.7,-33.8)
				..controls(164.3,-34)and(164.8,-34.1)..(164.9,-34.1)
				..controls(165,-34.1)and(165.6,-34.4)..(166.4,-34.8)
				..controls(167.3,-35.2)and(168.1,-35.4)..(169,-35.6)
				..controls(170,-35.8)and(170.5,-36)..(171.7,-36.6)
				..controls(172.4,-37)and(173.3,-37.4)..(173.6,-37.4)
				..controls(175.6,-37.7)and(176.2,-38)..(176.7,-39.1)
				..controls(176.9,-39.6)and(176.9,-39.7)..(176.7,-40.4)
				..controls(176.4,-41.6)and(175.8,-42.7)..(175.2,-43.4)
				..controls(174.5,-44.1)and(172.6,-45.5)..(171.7,-45.8)
				..controls(170.6,-46.2)and(169.1,-46)..(168.2,-45.4)
				..controls(167.9,-45.2)and(167.3,-44.8)..(166.9,-44.4)
				..controls(166.4,-44.1)and(165.8,-43.6)..(165.4,-43.2)
				..controls(165.1,-42.9)and(164.4,-42.3)..(163.9,-41.9)
				..controls(163.3,-41.5)and(162.7,-40.9)..(162.4,-40.5)
				..controls(161.4,-39.2)and(159.9,-38)..(158.8,-37.8)
				..controls(158.4,-37.8)and(158.3,-37.7)..(158.2,-37.2)
				..controls(158,-36.8)and(157.4,-36.1)..(157.1,-36.1)
				..controls(157,-36.1)and(156.7,-36.1)..(156.5,-36)
				..controls(156.1,-35.9)and(155.9,-35.7)..(155.7,-35.2)
				..controls(155.3,-34.6)and(155.1,-34.4)..(154,-34.1)
				..controls(153.6,-34)and(153.2,-33.7)..(152.9,-33.5)
				..controls(152.4,-32.9)and(151.6,-32.4)..(151.3,-32.4)
				..controls(150.9,-32.4)and(150.4,-32.6)..(150.4,-32.9)
				..controls(150.4,-33.2)and(150.8,-33.4)..(151.4,-33.6)
				..controls(152.1,-33.7)and(152.8,-34.1)..(153.5,-34.9)
				..controls(154.1,-35.5)and(154.3,-35.6)..(155,-35.8)
				..controls(155.7,-36)and(155.7,-36)..(156,-36.6)
				..controls(156.3,-37.3)and(156.7,-37.7)..(157.5,-37.8)
				..controls(158.1,-37.9)and(158.2,-38)..(158.4,-38.8)
				..controls(158.5,-39.4)and(158.6,-39.6)..(159.6,-40.4)
				..controls(160.2,-41)and(160.9,-41.6)..(161.3,-41.8)
				..controls(161.7,-42)and(161.9,-42.3)..(162.2,-42.8)
				..controls(162.4,-43.2)and(162.8,-43.8)..(163.1,-44.1)
				..controls(163.9,-44.8)and(165.1,-46.7)..(165.7,-48)
				..controls(166.2,-49.3)and(166.4,-49.7)..(167.2,-50.5)
				..controls(167.5,-50.8)and(167.9,-51.4)..(168.1,-51.8)
				..controls(168.3,-52.2)and(168.8,-53)..(169.1,-53.6)
				..controls(169.6,-54.3)and(169.8,-54.9)..(169.9,-55.6)
				..controls(170.4,-57.8)and(170.5,-58.1)..(171.2,-59.5)
				..controls(171.6,-60.3)and(172,-61.5)..(172.2,-62.1)
				..controls(172.3,-62.8)and(172.7,-63.7)..(172.9,-64.1)
				..controls(173.2,-64.6)and(173.4,-65.3)..(173.4,-65.6)
				..controls(173.8,-68.8)and(174.2,-70.6)..(174.9,-72.6)
				..controls(175.5,-74.5)and(175.9,-77.1)..(176.1,-80.2)
				..controls(176.3,-83.4)and(175.5,-84.4)..(173.7,-83.4)--cycle;
				\definecolor{ezs}{RGB}{247,179,26}\fill[ezs](-5.7,-156.6)--(-3.4,-157.5)--(61.9,12.7)--(59.7,13.6)--cycle;
				\definecolor{ezs}{RGB}{247,179,22}\fill[ezs](-3.8,-157.3)--(-1.5,-158.2)--(63.9,12)--(61.6,12.9)--cycle;
				\definecolor{ezs}{RGB}{248,179,19}\fill[ezs](-1.8,-158.1)--(0.4,-158.9)--(65.8,11.2)--(63.5,12.1)--cycle;
				\definecolor{ezs}{RGB}{248,178,16}\fill[ezs](0.1,-158.8)--(2.4,-159.7)--(67.7,10.5)--(65.4,11.4)--cycle;
				\definecolor{ezs}{RGB}{249,178,15}\fill[ezs](2,-159.5)--(4.3,-160.4)--(69.6,9.8)--(67.3,10.7)--cycle;
				\definecolor{ezs}{RGB}{251,177,17}\fill[ezs](3.9,-160.3)--(6.2,-161.1)--(71.5,9)--(69.2,9.9)--cycle;
				\definecolor{ezs}{RGB}{252,179,18}\fill[ezs](5.8,-161)--(8.1,-161.9)--(73.4,8.3)--(71.1,9.2)--cycle;
				\definecolor{ezs}{RGB}{252,182,17}\fill[ezs](7.7,-161.7)--(10,-162.6)--(75.3,7.6)--(73,8.5)--cycle;
				\definecolor{ezs}{RGB}{252,185,17}\fill[ezs](9.6,-162.5)--(11.9,-163.4)--(77.2,6.8)--(75,7.7)--cycle;
				\definecolor{ezs}{RGB}{251,189,18}\fill[ezs](11.5,-163.2)--(13.8,-164.1)--(79.2,6.1)--(76.9,7)--cycle;
				\definecolor{ezs}{RGB}{251,191,22}\fill[ezs](13.4,-163.9)--(15.7,-164.8)--(81.1,5.4)--(78.8,6.2)--cycle;
				\definecolor{ezs}{RGB}{252,193,25}\fill[ezs](15.4,-164.7)--(17.6,-165.6)--(83,4.6)--(80.7,5.5)--cycle;
				\definecolor{ezs}{RGB}{253,195,27}\fill[ezs](17.3,-165.4)--(19.6,-166.3)--(84.9,3.9)--(82.6,4.8)--cycle;
				\definecolor{ezs}{RGB}{253,196,31}\fill[ezs](19.2,-166.1)--(21.5,-167)--(86.8,3.2)--(84.5,4)--cycle;
				\definecolor{ezs}{RGB}{254,199,35}\fill[ezs](21.1,-166.9)--(23.4,-167.8)--(88.7,2.4)--(86.4,3.3)--cycle;
				\definecolor{ezs}{RGB}{253,201,38}\fill[ezs](23,-167.6)--(25.3,-168.5)--(90.6,1.7)--(88.3,2.6)--cycle;
				\definecolor{ezs}{RGB}{253,203,40}\fill[ezs](24.9,-168.3)--(27.2,-169.2)--(92.5,1)--(90.3,1.8)--cycle;
				\definecolor{ezs}{RGB}{253,205,42}\fill[ezs](26.8,-169.1)--(29.1,-170)--(94.4,0.2)--(92.2,1.1)--cycle;
				\definecolor{ezs}{RGB}{253,207,45}\fill[ezs](28.7,-169.8)--(31,-170.7)--(96.4,-0.5)--(94.1,0.4)--cycle;
				\definecolor{ezs}{RGB}{253,210,49}\fill[ezs](30.7,-170.5)--(32.9,-171.4)--(98.3,-1.2)--(96,-0.4)--cycle;
				\definecolor{ezs}{RGB}{253,212,54}\fill[ezs](32.6,-171.3)--(34.9,-172.2)--(100.2,-2)--(97.9,-1.1)--cycle;
				\definecolor{ezs}{RGB}{253,214,59}\fill[ezs](34.5,-172)--(36.8,-172.9)--(102.1,-2.7)--(99.8,-1.8)--cycle;
				\definecolor{ezs}{RGB}{253,216,63}\fill[ezs](36.4,-172.7)--(38.7,-173.6)--(104,-3.4)--(101.7,-2.6)--cycle;
				\definecolor{ezs}{RGB}{253,218,65}\fill[ezs](38.3,-173.5)--(40.6,-174.4)--(105.9,-4.2)--(103.6,-3.3)--cycle;
				\definecolor{ezs}{RGB}{254,220,66}\fill[ezs](40.2,-174.2)--(42.5,-175.1)--(107.8,-4.9)--(105.5,-4)--cycle;
				\definecolor{ezs}{RGB}{254,222,68}\fill[ezs](42.1,-175)--(44.4,-175.8)--(109.7,-5.6)--(107.5,-4.8)--cycle;
				\definecolor{ezs}{RGB}{254,224,71}\fill[ezs](44,-175.7)--(46.3,-176.6)--(111.7,-6.4)--(109.4,-5.5)--cycle;
				\definecolor{ezs}{RGB}{254,226,75}\fill[ezs](46,-176.4)--(48.2,-177.3)--(113.6,-7.1)--(111.3,-6.2)--cycle;
				\definecolor{ezs}{RGB}{254,228,79}\fill[ezs](47.9,-177.2)--(50.1,-178)--(115.5,-7.8)--(113.2,-7)--cycle;
				\definecolor{ezs}{RGB}{254,230,82}\fill[ezs](49.8,-177.9)--(52.1,-178.8)--(117.4,-8.6)--(115.1,-7.7)--cycle;
				\definecolor{ezs}{RGB}{254,232,85}\fill[ezs](51.7,-178.6)--(54,-179.5)--(119.3,-9.3)--(117,-8.4)--cycle;
				\definecolor{ezs}{RGB}{254,234,88}\fill[ezs](53.6,-179.4)--(55.9,-180.2)--(121.2,-10)--(118.9,-9.2)--cycle;
				\definecolor{ezs}{RGB}{254,237,91}\fill[ezs](55.5,-180.1)--(57.8,-181)--(123.1,-10.8)--(120.8,-9.9)--cycle;
				\definecolor{ezs}{RGB}{254,239,94}\fill[ezs](57.4,-180.8)--(59.7,-181.7)--(125,-11.5)--(122.8,-10.6)--cycle;
				\definecolor{ezs}{RGB}{254,241,97}\fill[ezs](59.3,-181.6)--(61.6,-182.4)--(127,-12.2)--(124.7,-11.4)--cycle;
				\definecolor{ezs}{RGB}{254,243,101}\fill[ezs](61.2,-182.3)--(63.5,-183.2)--(128.9,-13)--(126.6,-12.1)--cycle;
				\definecolor{ezs}{RGB}{254,245,103}\fill[ezs](63.2,-183)--(65.4,-183.9)--(130.8,-13.7)--(128.5,-12.8)--cycle;
				\definecolor{ezs}{RGB}{254,245,105}\fill[ezs](65.1,-183.8)--(67.4,-184.6)--(132.7,-14.4)--(130.4,-13.6)--cycle;
				\definecolor{ezs}{RGB}{254,245,105}\fill[ezs](67,-184.5)--(69.3,-185.4)--(134.6,-15.2)--(132.3,-14.3)--cycle;
				\definecolor{ezs}{RGB}{254,244,103}\fill[ezs](68.9,-185.2)--(71.2,-186.1)--(136.5,-15.9)--(134.2,-15)--cycle;
				\definecolor{ezs}{RGB}{253,243,101}\fill[ezs](70.8,-186)--(73.1,-186.8)--(138.4,-16.6)--(136.1,-15.8)--cycle;
				\definecolor{ezs}{RGB}{253,241,98}\fill[ezs](72.7,-186.7)--(75,-187.6)--(140.3,-17.4)--(138,-16.5)--cycle;
				\definecolor{ezs}{RGB}{253,239,95}\fill[ezs](74.6,-187.4)--(76.9,-188.3)--(142.2,-18.1)--(140,-17.2)--cycle;
				\definecolor{ezs}{RGB}{253,236,92}\fill[ezs](76.5,-188.2)--(78.8,-189)--(144.2,-18.8)--(141.9,-18)--cycle;
				\definecolor{ezs}{RGB}{253,234,88}\fill[ezs](78.5,-188.9)--(80.7,-189.8)--(146.1,-19.6)--(143.8,-18.7)--cycle;
				\definecolor{ezs}{RGB}{253,232,85}\fill[ezs](80.4,-189.6)--(82.7,-190.5)--(148,-20.3)--(145.7,-19.4)--cycle;
				\definecolor{ezs}{RGB}{253,230,82}\fill[ezs](82.3,-190.4)--(84.6,-191.2)--(149.9,-21)--(147.6,-20.2)--cycle;
				\definecolor{ezs}{RGB}{253,228,78}\fill[ezs](84.2,-191.1)--(86.5,-192)--(151.8,-21.8)--(149.5,-20.9)--cycle;
				\definecolor{ezs}{RGB}{253,226,74}\fill[ezs](86.1,-191.8)--(88.4,-192.7)--(153.7,-22.5)--(151.4,-21.6)--cycle;
				\definecolor{ezs}{RGB}{253,224,70}\fill[ezs](88,-192.6)--(90.3,-193.4)--(155.6,-23.3)--(153.3,-22.4)--cycle;
				\definecolor{ezs}{RGB}{253,222,67}\fill[ezs](89.9,-193.3)--(92.2,-194.2)--(157.5,-24)--(155.3,-23.1)--cycle;
				\definecolor{ezs}{RGB}{253,220,65}\fill[ezs](91.8,-194)--(94.1,-194.9)--(159.5,-24.7)--(157.2,-23.8)--cycle;
				\definecolor{ezs}{RGB}{253,217,63}\fill[ezs](93.7,-194.8)--(96,-195.6)--(161.4,-25.5)--(159.1,-24.6)--cycle;
				\definecolor{ezs}{RGB}{253,215,60}\fill[ezs](95.7,-195.5)--(97.9,-196.4)--(163.3,-26.2)--(161,-25.3)--cycle;
				\definecolor{ezs}{RGB}{253,213,56}\fill[ezs](97.6,-196.2)--(99.9,-197.1)--(165.2,-26.9)--(162.9,-26)--cycle;
				\definecolor{ezs}{RGB}{252,211,53}\fill[ezs](99.5,-197)--(101.8,-197.8)--(167.1,-27.7)--(164.8,-26.8)--cycle;
				\definecolor{ezs}{RGB}{252,209,49}\fill[ezs](101.4,-197.7)--(103.7,-198.6)--(169,-28.4)--(166.7,-27.5)--cycle;
				\definecolor{ezs}{RGB}{252,207,45}\fill[ezs](103.3,-198.4)--(105.6,-199.3)--(170.9,-29.1)--(168.6,-28.2)--cycle;
				\definecolor{ezs}{RGB}{253,205,42}\fill[ezs](105.2,-199.2)--(107.5,-200)--(172.8,-29.9)--(170.6,-29)--cycle;
				\definecolor{ezs}{RGB}{253,202,40}\fill[ezs](107.1,-199.9)--(109.4,-200.8)--(174.7,-30.6)--(172.5,-29.7)--cycle;
				\definecolor{ezs}{RGB}{253,200,38}\fill[ezs](109,-200.6)--(111.3,-201.5)--(176.7,-31.3)--(174.4,-30.4)--cycle;
				\definecolor{ezs}{RGB}{253,198,35}\fill[ezs](111,-201.4)--(113.2,-202.2)--(178.6,-32.1)--(176.3,-31.2)--cycle;
				\definecolor{ezs}{RGB}{253,196,31}\fill[ezs](112.9,-202.1)--(115.2,-203)--(180.5,-32.8)--(178.2,-31.9)--cycle;
				\definecolor{ezs}{RGB}{253,194,28}\fill[ezs](114.8,-202.8)--(117.1,-203.7)--(182.4,-33.5)--(180.1,-32.6)--cycle;
				\definecolor{ezs}{RGB}{253,192,26}\fill[ezs](116.7,-203.6)--(119,-204.5)--(184.3,-34.3)--(182,-33.4)--cycle;
				\definecolor{ezs}{RGB}{253,189,24}\fill[ezs](118.6,-204.3)--(120.9,-205.2)--(186.2,-35)--(183.9,-34.1)--cycle;
				\definecolor{ezs}{RGB}{253,186,20}\fill[ezs](120.5,-205)--(122.8,-205.9)--(188.1,-35.7)--(185.8,-34.8)--cycle;
				\definecolor{ezs}{RGB}{253,184,17}\fill[ezs](122.4,-205.8)--(124.7,-206.7)--(190,-36.5)--(187.8,-35.6)--cycle;
				\definecolor{ezs}{RGB}{252,183,15}\fill[ezs](124.3,-206.5)--(126.6,-207.4)--(192,-37.2)--(189.7,-36.3)--cycle;
				\definecolor{ezs}{RGB}{252,181,14}\fill[ezs](126.3,-207.2)--(128.5,-208.1)--(193.9,-37.9)--(191.6,-37.1)--cycle;
				\definecolor{ezs}{RGB}{252,179,13}\fill[ezs](128.2,-208)--(130.4,-208.9)--(195.8,-38.7)--(193.5,-37.8)--cycle;
				\definecolor{ezs}{RGB}{252,178,12}\fill[ezs](130.1,-208.7)--(132.4,-209.6)--(197.7,-39.4)--(195.4,-38.5)--cycle;
				\definecolor{ezs}{RGB}{251,178,10}\fill[ezs](132,-209.4)--(134.3,-210.3)--(199.6,-40.1)--(197.3,-39.3)--cycle;
				\definecolor{ezs}{RGB}{251,178,10}\fill[ezs](133.9,-210.2)--(136.2,-211.1)--(201.5,-40.9)--(199.2,-40)--cycle;
				\definecolor{ezs}{RGB}{250,177,11}\fill[ezs](135.8,-210.9)--(138.1,-211.8)--(203.4,-41.6)--(201.1,-40.7)--cycle;
				\definecolor{ezs}{RGB}{250,178,11}\fill[ezs](137.7,-211.6)--(140,-212.5)--(205.3,-42.3)--(203.1,-41.5)--cycle;
				\definecolor{ezs}{RGB}{248,180,10}\fill[ezs](139.6,-212.4)--(141.9,-213.3)--(207.3,-43.1)--(205,-42.2)--cycle;
				\definecolor{ezs}{RGB}{247,181,9}\fill[ezs](141.5,-213.1)--(143.8,-214)--(209.2,-43.8)--(206.9,-42.9)--cycle;
				\definecolor{ezs}{RGB}{246,182,9}\fill[ezs](143.5,-213.8)--(145.7,-214.7)--(211.1,-44.5)--(208.8,-43.7)--cycle;
				\definecolor{ezs}{RGB}{246,182,10}\fill[ezs](145.4,-214.6)--(147.7,-215.5)--(213,-45.3)--(210.7,-44.4)--cycle;
			\end{scope}
			% chữ EDUZY (Times đậm, làm dày nét bằng 5 lớp lệch nhẹ; căn đúng khung chữ logo gốc)
			\foreach \dx/\dy in {-0.7/0,0.7/0,0/-0.3,0/0.3,0/0}
			\node[anchor=center,inner sep=0pt,outer sep=0pt,text=ezText,align=center] at (109.6+\dx,-193.6+\dy)
			{\resizebox{1.01cm}{0.22cm}{\fontencoding{T1}\fontfamily{ptm}\fontseries{b}\selectfont EDUZY}};
		\end{tikzpicture}%
	}%
}
%====================== THIẾT LẬP THÂN ĐỀ ======================
%--- Môi trường tự luận {bt}: ex_test 3.3.4 chưa khai báo sẵn nên tạo ở đây
%    (ex_test đã tự móc \fixshowans và \listenumerate cho bt)
\makeatletter
\@ifundefined{bt}{\newtheorem{bt}{\color{gold1}\namebt}}{}
\makeatother
\renewcommand{\sectionmark}[1]{\markright{\thesection\ #1}}
%--- Khung tiêu đề (logo + tối đa 4 dòng; dòng 3, 4 để trống {} thì bỏ qua)
%  Đề thi:    \khungtieude{THI THỬ THPT QUỐC GIA NĂM 2026}{Môn Toán -- Đề số 21}
%                         {(Đề thi gồm: ...)}{Thời gian làm bài: 90 phút (không kể thời gian phát đề).}
%  Chuyên đề: \khungtieude{CHUYÊN ĐỀ 3 -- DÃY SỐ VÀ CÔNG THỨC TRUY HỒI}
%                         {Bồi dưỡng học sinh giỏi Toán THPT}{...}{}
\newcommand{\khungtieude}[4]{%
	\begin{tcolorbox}[colframe=gold1,colback=gold5!45,boxrule=0.5mm]
		\begin{minipage}[c]{3cm}
			\centerline{\eduzylogo[2.4cm]}
		\end{minipage}\hspace{0.5cm}
		\begin{minipage}[c]{12.5cm}
			\vskip 0.15 cm
			\centerline{\ind{#1}}
			\vskip 0.15 cm
			\centerline{\textbf{#2}}
			\if\relax\detokenize{#3}\relax\else
				\vskip 0.15 cm
				\centerline{\textit{#3}}
			\fi
			\if\relax\detokenize{#4}\relax\else
				\vskip 0.15 cm
				\centerline{\textit{#4}}
			\fi
			\vskip 0.15 cm
		\end{minipage}
	\end{tcolorbox}
}
%====================== BỔ SUNG CHO BÀI HỌC, CHUYÊN ĐỀ ======================
\usepackage[shortlabels]{enumitem}
\usepackage{tabularx}
%--- Danh sách gọn: \begin{enumerate}[gon] ... \end{enumerate}
\SetEnumitemKey{gon}{leftmargin=*,itemsep=1pt,parsep=0pt,topsep=3pt}
%--- Khung lý thuyết: \begin{khung}{} ... \end{khung}; tiêu đề chỉ dùng khi thật ngắn ("Mẹo", "Lỗi thường gặp")
\newtcolorbox{khung}[1]{breakable,enhanced,colframe=gold1,colback=gold5!45,
	boxrule=0.4mm,fonttitle=\bfseries\sffamily,coltitle=white,colbacktitle=gold1,
	attach boxed title to top left={xshift=6mm,yshift=-3mm},
	boxed title style={arc=1mm,boxrule=0.3mm},title={#1}}
%--- Tiêu đề mục nhỏ trong một dạng: \muc{Dấu hiệu}, \muc{Phương pháp}, \muc{Bài tự luyện}
\makeatletter
\newcommand{\muc}[1]{\par\medskip\noindent\indam{#1}\par\nobreak\smallskip\@afterindentfalse\@afterheading}
\makeatother
%--- Ô công thức cần nhớ (căn giữa, viền vàng, vừa khít nội dung)
%     \begin{center}\congthuc{$Q(n)=(n-1)!$}\end{center}
\newtcbox{\congthuc}{on line,colframe=gold1,colback=gold5!45,boxrule=0.5pt,arc=2pt,
	left=8pt,right=8pt,top=4pt,bottom=4pt}
%--- Công tắc lời giải bài tự luyện (ví dụ mẫu luôn dùng \loigiai)
%     \loigiaiTL[{đáp số}]{lời giải đầy đủ}
%     \TLfull -> in lời giải đầy đủ;  \TLdapso -> chỉ in một dòng "Đáp số";  \TLan -> ẩn hoàn toàn
%     (\CheDo là dethi thì mọi lời giải và đáp số đều ẩn)
\makeatletter
\newif\ifTLhien   \TLhientrue
\newif\ifTLdapso  \TLdapsofalse
\newcommand{\TLfull}{\TLhientrue\TLdapsofalse}
\newcommand{\TLdapso}{\TLhientrue\TLdapsotrue}
\newcommand{\TLan}{\TLhienfalse}
\newcommand{\loigiaiTL}[2][]{%
	\ifTLhien
		\ifTLdapso
			\def\@tlarg{#1}%
			\ifx\@tlarg\@empty\else
				\ifdefstring{\kindOPTN}{dethi}{}{\par\noindent{\bfseries\color{gold1}Đáp số:} #1\par}%
			\fi
		\else
			\loigiai{#2}%
		\fi
	\fi}
\makeatother
\TLdapso% mặc định: chỉ in đáp số bài tự luyện
%====================== BẢNG ĐÚNG/SAI (\choiceTFt) ======================
%  Hàng tiêu đề nền nâu đồng chữ trắng (Phát biểu | Đ | S), các hàng tô xen kẽ kem / trắng,
%  đường kẻ vàng kim, nhãn a) b) c) d) màu nâu đồng. Bản có đáp án: cột được chọn có vòng tròn
%  nâu đồng với dấu tích trắng (câu đúng tích ở cột Đúng, câu sai tích ở cột Sai).
%  Tiêu đề ghi đủ "Đúng | Sai": \OPTN{viettat=0} sau \input (cột Đúng/Sai rộng ra, phát biểu hẹp lại).
%  Dùng bảng gốc của ex_test: \def\BangDSgoc{1} trước \input{preamble-eduzy-v2}.
\providecommand{\BangDSgoc}{0}
\ifnum\BangDSgoc=0
\usepackage{colortbl,pifont}
\colorlet{ezDSke}{gold2!60!gold1}% đường kẻ bảng (nét 0,4pt như ex_test để không đổi chiều cao bảng)
\colorlet{ezDSle}{gold5!60}% hàng a), c)
\colorlet{ezDSchan}{white}% hàng b), d)
%--- dấu chọn ở cột Đúng/Sai
%    (bề rộng 0, căn giữa ô: không làm đổi bề rộng cột của ex_test, không báo tràn lề)
\newcommand{\EZdauchon}{\makebox[0pt]{\tikz[baseline=-1ex]\node[circle,fill=gold1,inner sep=0pt,
	minimum size=1.2em,text=white,font=\small]{\ding{51}};}}
\def\tickT{\EZdauchon}\def\tickF{\EZdauchon}
\AtBeginDocument{\sbox{\ccTFTab}{\EZdauchon}}
%--- nhãn a) b) c) d) màu nâu đồng (ở bản lời giải, nhãn câu đúng vẫn đóng khung như ex_test)
\renewcommand{\FalseTF}{\stepcounter{dapan}\textbf{\color{gold1}\DapAnTF\sepTF}\ignorespaces}
%--- trước khi dựng bảng: kẻ màu vàng kim (bề rộng cột giữ nguyên như ex_test); sau bảng trả lại màu đen
\newcount\EZhangTF
\newcommand{\EZbatdauTF}{%
	\global\EZhangTF=0\relax
	\arrayrulecolor{ezDSke}}
\newcommand{\EZketthucTF}{\arrayrulecolor{black}}
%--- hàng tiêu đề nền nâu đồng, chữ trắng
\makeatletter
\newcommand{\EZo@TF}[1]{\centering\bfseries\color{white}#1}
\newcommand{\EZtieudeTF}{%
	\ifcase\cotDS
		\ifnum\cotSTT=1\relax
			\def\rowHeader{\hline\rowcolor{gold1}\EZo@TF{TT}&\EZo@TF{\phatbieu}\tabularnewline}%
		\else
			\def\rowHeader{\hline\rowcolor{gold1}\EZo@TF{\phatbieu}\tabularnewline}%
		\fi
	\or
		\ifnum\cotSTT=1\relax
			\def\rowHeader{\hline\rowcolor{gold1}\EZo@TF{TT}&\EZo@TF{\phatbieu}&\EZo@TF{\Dung}\tabularnewline}%
		\else
			\def\rowHeader{\hline\rowcolor{gold1}\EZo@TF{\phatbieu}&\EZo@TF{\Dung}\tabularnewline}%
		\fi
	\or
		\ifnum\cotSTT=1\relax
			\def\rowHeader{\hline\rowcolor{gold1}\EZo@TF{TT}&\EZo@TF{\phatbieu}&\EZo@TF{\Dung}&\EZo@TF{\Sai}\tabularnewline}%
		\else
			\def\rowHeader{\hline\rowcolor{gold1}\EZo@TF{\phatbieu}&\EZo@TF{\Dung}&\EZo@TF{\Sai}\tabularnewline}%
		\fi
	\fi}
%--- tô xen kẽ từng hàng (gọi ngay sau \hline đầu mỗi hàng)
\newcommand{\EZmauhangTF}{%
	\noalign{\global\advance\EZhangTF\@ne}%
	\ifodd\EZhangTF\rowcolor{ezDSle}\else\rowcolor{ezDSchan}\fi}
%--- gắn vào \choiceTFt / \writeChoiceTFt của ex_test
\xpretocmd{\choiceTFt}{\EZbatdauTF}{}{\typeout{[preamble-eduzy-v2] Không gắn được bảng Đúng/Sai (đầu bảng)}}
\xpatchcmd{\choiceTFt}{\ifnum\dongPB=0\def\rowHeader{}\fi}{\EZtieudeTF\ifnum\dongPB=0\def\rowHeader{}\fi}
	{}{\typeout{[preamble-eduzy-v2] Không gắn được bảng Đúng/Sai (hàng tiêu đề)}}
\xpatchcmd{\writeChoiceTFt}{\hline\xdef\dotEX{\sdotEX}}{\hline\EZmauhangTF\xdef\dotEX{\sdotEX}}
	{}{\typeout{[preamble-eduzy-v2] Không gắn được bảng Đúng/Sai (tô hàng)}}
\xpatchcmd{\writeChoiceTFt}{{\end{longtable}}}{{\end{longtable}\EZketthucTF}}
	{}{\typeout{[preamble-eduzy-v2] Không gắn được bảng Đúng/Sai (cuối longtable)}}
\xpatchcmd{\writeChoiceTFt}{{\end{xtabular}\end{center}\medskip\noindent}}{{\end{xtabular}\end{center}\EZketthucTF\medskip\noindent}}
	{}{\typeout{[preamble-eduzy-v2] Không gắn được bảng Đúng/Sai (cuối xtabular)}}
\fi
%=====================================================================
\renewcommand{\parallel}{\mathrel{/\!\!/}}

%  trong file đề (tools/dan_trang/dan_trang.py tự ghi khi đặt \NeedspaceCau):
%    \def\TachCauDS{1}   câu Đúng/Sai được tách: câu dẫn (kèm hình) ở cuối trang, bảng Phát biểu sang trang sau
%                        (câu dẫn và hình không tách nhau, bảng không tách giữa hai trang)
%    \def\TachCauLG{1}   bản lời giải (\CheDo khác dethi): phần đề của câu cao hơn 1/3 trang được phép tách
%  Lệnh dùng trong file đề (giá trị do tools/dan_trang/dan_trang.py đo ở bản dethi):
%    \NeedspaceCau[<cao câu dẫn + hình>]{<cao cả khối câu>}   ngay trên \begin{ex} (câu đầu phần: trên tiêu đề phần)
%        công tắc tắt: đúng bằng \Needspace{<cao cả khối câu>} như cũ;
%        \TachCauDS=1 và có [..] (câu Đúng/Sai được chọn tách): chỉ giữ câu dẫn + hình;
%        \TachCauLG=1 ở bản lời giải: đòi tối đa 1/3 \textheight.
%    \NeedspaceBang{<cao bảng Phát biểu>}%   ngay trước \choiceTFt: \TachCauDS=1 thì giữ trọn bảng trên một trang;
%        công tắc tắt: không làm gì.
\providecommand{\TachCauDS}{0}
\providecommand{\TachCauLG}{0}
\newlength{\EZnsCau}
\newcommand{\NeedspaceCau}[2][]{%
	\setlength{\EZnsCau}{#2}%
	\ifnum\TachCauDS=1 \if\relax\detokenize{#1}\relax\else\setlength{\EZnsCau}{#1}\fi\fi
	\ifnum\TachCauLG=1 \ifdefstring{\kindOPTN}{dethi}{}{\ifdim\EZnsCau>0.3333\textheight\setlength{\EZnsCau}{0.3333\textheight}\fi}\fi
	\Needspace{\EZnsCau}}
\newcommand{\NeedspaceBang}[1]{\ifnum\TachCauDS=1 \Needspace{#1}\fi}

%--- Thang màu: đậm -> nhạt
\definecolor{gold1}{RGB}{126,96,0}    % nâu đồng đậm  #7E6000
\definecolor{gold2}{RGB}{192,144,3}   % vàng kim      #C09003
\definecolor{gold3}{RGB}{255,215,99}  % vàng nhạt     #FFD763
\definecolor{gold4}{RGB}{255,231,153} % vàng rất nhạt #FFE799
\definecolor{gold5}{RGB}{255,243,205} % vàng kem      #FFF3CD
%--- Ghi đè tông xanh/đỏ cũ
\definecolor{myblue}{RGB}{126,96,0}
\colorlet{blue}{gold1}
\colorlet{cyan}{gold2}
\colorlet{teal}{gold2}

%--- Tên chương mặc định (file đề nên tự khai báo lại)
\providecommand{\tenchuong}{}
%--- Chân trang
\providecommand{\thuonghieu}{TOÁN 12 -- Eduzy.com.vn}
\providecommand{\tacgia}{Tô Gia Bảo -- 0353703783}

%--- Chữ cạnh logo ở đầu trang
\providecommand{\khauhieu}{Eduzy.com.vn}

%====================== THÔNG SỐ ĐẦU TRANG, CHÂN TRANG (v2) ======================
%  Dải nền đậm tràn ngang trang; ô chữ vàng kem có cạnh xiên, kèm 2 vạch xiên chuyển màu.
%  Đầu trang lẻ : logo + \khauhieu (trái)   | ô chữ \tenchuong (phải)
%  Chân trang lẻ: ô chữ \thuonghieu (trái)  | Trang (số) (phải)
%  Trang chẵn soi gương trang lẻ; chân trang chẵn ghi \tacgia.
%--- Màu: chỉ dùng thang vàng kim Eduzy
\colorlet{ezNen}{gold1}% dải nền đậm
\colorlet{ezVach1}{gold3}% vạch xiên sát dải nền (đậm hơn)
\colorlet{ezVach2}{gold4}% vạch xiên sát ô chữ (nhạt hơn)
\colorlet{ezO}{gold5}% ô chữ vàng kem
%--- Kích thước
\newlength{\EZdai}\setlength{\EZdai}{1.1cm}% chiều cao dải nền
\newlength{\EZo}\setlength{\EZo}{0.68cm}% chiều cao ô chữ
\newlength{\EZxien}\setlength{\EZxien}{0.55cm}% độ lệch ngang của cạnh xiên
\newlength{\EZvach}\setlength{\EZvach}{0.75cm}% bề rộng mỗi vạch xiên
\newlength{\EZkhe}\setlength{\EZkhe}{0.12cm}% khe trắng giữa vạch và ô chữ
\newlength{\EZdem}\setlength{\EZdem}{0.4cm}% khoảng đệm từ chữ tới cạnh xiên
\newlength{\EZle}\setlength{\EZle}{1.8cm}% lề trái/phải của chữ (bằng lề thân bài)
\newlength{\EZrongmin}\setlength{\EZrongmin}{4.5cm}% ô chữ rộng tối thiểu
\newlength{\EZrong}% bề rộng chữ trong ô (tính tự động)
%--- Font
\newcommand{\EZfontO}{\fontsize{10.5pt}{12pt}\selectfont\bfseries}% chữ trong ô
\newcommand{\EZfontKH}{\fontfamily{qag}\fontsize{11pt}{12pt}\selectfont\bfseries}% \khauhieu cạnh logo
\newcommand{\EZfontTrang}{\fontsize{11.5pt}{12pt}\selectfont\bfseries}% chữ "Trang" và số trang
%--- Đo bề rộng chữ trong ô
\newcommand{\EZdorong}[1]{%
	\settowidth{\EZrong}{\EZfontO #1}%
	\ifdim\EZrong<\EZrongmin\setlength{\EZrong}{\EZrongmin}\fi}
%--- Logo + chữ trắng \khauhieu (dựng sẵn một lần khi vào \begin{document})
\newsavebox{\EZlogobox}
\AtBeginDocument{\sbox{\EZlogobox}{%
		\raisebox{-0.5\height}{\eduzylogo[1cm]}\hspace{0.15cm}%
		\raisebox{-0.5\height}{\color{white}\EZfontKH\khauhieu}}}
%--- Hai vạch xiên + khe trắng. #1: điểm gốc ở mép dưới, #2: hướng (1 sang phải, -1 sang trái),
%    #3: hướng nghiêng (1 kiểu "/", -1 kiểu "\")
\newcommand{\EZvachxien}[3]{%
	\fill[white] (#1) --++({(#3)*\EZxien},\EZo) --++({(#2)*\EZkhe},0) --++({-(#3)*\EZxien},-\EZo) --cycle;
	\fill[ezVach2] ([xshift=(#2)*\EZkhe]#1) --++({(#3)*\EZxien},\EZo) --++({(#2)*\EZvach},0) --++({-(#3)*\EZxien},-\EZo) --cycle;
	\fill[ezVach1] ([xshift=(#2)*(\EZkhe+\EZvach)]#1) --++({(#3)*\EZxien},\EZo) --++({(#2)*\EZvach},0) --++({-(#3)*\EZxien},-\EZo) --cycle;
	\draw[white,line width=0.6pt] ([xshift=(#2)*(\EZkhe+\EZvach)]#1) --++({(#3)*\EZxien},\EZo);
}
%--- Chữ trong ô (căn giữa theo chiều cao ô, không lệch vì dấu tiếng Việt)
\tikzset{ezchu/.style={inner sep=0pt,text=ezNen,text height=0.72em,text depth=0.22em}}

%====================== LOGO CHÌM GIỮA TRANG ======================
%  Vẽ cùng đầu trang nên nằm dưới chữ của thân bài, có trên mọi trang.
%  Tắt: \logochimtat (từ trang đó trở đi), bật lại: \logochimbat.
%  Đổi cỡ, độ đậm: khai báo trước \input, ví dụ \def\EZchimcao{9cm} \def\EZchimdo{0.1}
\providecommand{\EZchimcao}{10cm}% chiều cao logo chìm
\providecommand{\EZchimdo}{0.12}% độ đậm (0 = trong suốt, 1 = rõ hẳn)
\newif\ifEZchim \EZchimtrue
\newcommand{\logochimtat}{\EZchimfalse}
\newcommand{\logochimbat}{\EZchimtrue}
%  Dựng sẵn logo mờ một lần khi vào \begin{document}; transparency group làm mờ cả logo
%  một lượt nên không lộ vệt ở chỗ các dải màu chồng nhau.
\newsavebox{\EZchimtho}
\newsavebox{\EZchimbox}
\AtBeginDocument{%
	\sbox{\EZchimtho}{\eduzylogo[\EZchimcao]}%
	\sbox{\EZchimbox}{\begin{tikzpicture}
			\begin{scope}[transparency group,opacity=\EZchimdo]
				\node[inner sep=0pt] {\usebox{\EZchimtho}};
			\end{scope}
	\end{tikzpicture}}}
\newcommand{\EZvechim}{%
	\ifEZchim
		\node[inner sep=0pt] at (current page.center) {\usebox{\EZchimbox}};
	\fi}

%====================== ĐẦU TRANG - TRANG LẺ ======================
%  (logo bên trái, ô chữ \tenchuong bên phải)
\newcommand{\myfancyheadOdd}{%
	\EZdorong{\tenchuong}%
	\begin{tikzpicture}[remember picture,overlay]
		\coordinate (T) at (current page.north west);
		\coordinate (P) at (current page.north east);
		%--- logo chìm giữa trang
		\EZvechim
		%--- dải nền
		\fill[ezNen] (T) rectangle ([yshift=-\EZdai]P);
		%--- ô chữ: X là góc dưới bên trái, cạnh trái xiên kiểu "/"
		\coordinate (X) at ([xshift=-\EZle-\EZrong-\EZdem-\EZxien,yshift=-\EZdai]P);
		\fill[ezO] (X) --++(\EZxien,\EZo) -- ([yshift=-\EZdai+\EZo]P) -- ([yshift=-\EZdai]P) --cycle;
		\EZvachxien{X}{-1}{1}
		\node[ezchu,anchor=east] at ([xshift=-\EZle,yshift=-\EZdai+0.5*\EZo]P) {\EZfontO\tenchuong};
		%--- logo + tên
		\node[inner sep=0pt,anchor=west] at ([xshift=\EZle,yshift=-0.5*\EZdai]T) {\usebox{\EZlogobox}};
	\end{tikzpicture}%
}

%====================== ĐẦU TRANG - TRANG CHẴN =====================
%  (soi gương: ô chữ \tenchuong bên trái, logo bên phải)
\newcommand{\myfancyheadEven}{%
	\EZdorong{\tenchuong}%
	\begin{tikzpicture}[remember picture,overlay]
		\coordinate (T) at (current page.north west);
		\coordinate (P) at (current page.north east);
		%--- logo chìm giữa trang
		\EZvechim
		%--- dải nền
		\fill[ezNen] (T) rectangle ([yshift=-\EZdai]P);
		%--- ô chữ: X là góc dưới bên phải, cạnh phải xiên kiểu "\"
		\coordinate (X) at ([xshift=\EZle+\EZrong+\EZdem+\EZxien,yshift=-\EZdai]T);
		\fill[ezO] (X) --++(-\EZxien,\EZo) -- ([yshift=-\EZdai+\EZo]T) -- ([yshift=-\EZdai]T) --cycle;
		\EZvachxien{X}{1}{-1}
		\node[ezchu,anchor=west] at ([xshift=\EZle,yshift=-\EZdai+0.5*\EZo]T) {\EZfontO\tenchuong};
		%--- logo + tên
		\node[inner sep=0pt,anchor=east] at ([xshift=-\EZle,yshift=-0.5*\EZdai]P) {\usebox{\EZlogobox}};
	\end{tikzpicture}%
}

%====================== CHÂN TRANG - TRANG LẺ ======================
%  (ô chữ \thuonghieu bên trái, "Trang (số)" bên phải)
\newcommand{\myfancyfootOdd}{%
	\EZdorong{\thuonghieu}%
	\begin{tikzpicture}[remember picture,overlay]
		\coordinate (T) at (current page.south west);
		\coordinate (P) at (current page.south east);
		%--- dải nền
		\fill[ezNen] (T) rectangle ([yshift=\EZdai]P);
		%--- ô chữ sát mép trên dải: X là góc dưới bên phải, cạnh phải xiên kiểu "/"
		\coordinate (X) at ([xshift=\EZle+\EZrong+\EZdem,yshift=\EZdai-\EZo]T);
		\fill[ezO] ([yshift=\EZdai-\EZo]T) -- (X) --++(\EZxien,\EZo) -- ([yshift=\EZdai]T) --cycle;
		\EZvachxien{X}{1}{1}
		\node[ezchu,anchor=west] at ([xshift=\EZle,yshift=\EZdai-0.5*\EZo]T) {\EZfontO\thuonghieu};
		%--- số trang
		\node[circle,draw=white,line width=0.9pt,text=white,inner sep=0pt,minimum size=0.62cm,anchor=east]
		(so) at ([xshift=-\EZle,yshift=0.5*\EZdai]P) {\EZfontTrang\thepage};
		\node[inner sep=0pt,text=white,anchor=base east] (tr) at ([xshift=-0.15cm]so.base west) {\EZfontTrang Trang};
	\end{tikzpicture}%
}

%====================== CHÂN TRANG - TRANG CHẴN ====================
%  (soi gương: "Trang (số)" bên trái, ô chữ \tacgia bên phải)
\newcommand{\myfancyfootEven}{%
	\EZdorong{\tacgia}%
	\begin{tikzpicture}[remember picture,overlay]
		\coordinate (T) at (current page.south west);
		\coordinate (P) at (current page.south east);
		%--- dải nền
		\fill[ezNen] (T) rectangle ([yshift=\EZdai]P);
		%--- ô chữ sát mép trên dải: X là góc dưới bên trái, cạnh trái xiên kiểu "\"
		\coordinate (X) at ([xshift=-\EZle-\EZrong-\EZdem,yshift=\EZdai-\EZo]P);
		\fill[ezO] ([yshift=\EZdai-\EZo]P) -- (X) --++(-\EZxien,\EZo) -- ([yshift=\EZdai]P) --cycle;
		\EZvachxien{X}{-1}{-1}
		\node[ezchu,anchor=east] at ([xshift=-\EZle,yshift=\EZdai-0.5*\EZo]P) {\EZfontO\tacgia};
		%--- số trang
		%    (vòng tròn ẩn so0 nằm giữa dải để lấy đường chân chữ, giống trang lẻ)
		\node[circle,inner sep=0pt,minimum size=0.62cm] (so0) at ([yshift=0.5*\EZdai]T) {\EZfontTrang\phantom{\thepage}};
		\node[inner sep=0pt,text=white,anchor=base west] (tr) at ($(so0.base -| T)+(\EZle,0)$) {\EZfontTrang Trang};
		\node[circle,draw=white,line width=0.9pt,text=white,inner sep=0pt,minimum size=0.62cm,anchor=base west]
		(so) at ([xshift=0.15cm]tr.base east) {\EZfontTrang\thepage};
	\end{tikzpicture}%
}

%--- Giữ tương thích với các file cũ đang gọi \myfancyhead, \myfancyfoot
\newcommand{\myfancyhead}{\myfancyheadOdd}
\newcommand{\myfancyfoot}{\myfancyfootOdd}

%======================= GẮN VÀO TRANG =============================
\pagestyle{fancy}
\fancyhf{}
\fancyhead[LO]{\myfancyheadOdd}
\fancyhead[RE]{\myfancyheadEven}
\fancyfoot[LO]{\myfancyfootOdd}
\fancyfoot[RE]{\myfancyfootEven}
\renewcommand{\headrulewidth}{0pt}
\renewcommand{\footrulewidth}{0pt}
\fancypagestyle{plain}{%
	\fancyhf{}
	\fancyhead[LO]{\myfancyheadOdd}
	\fancyhead[RE]{\myfancyheadEven}
	\fancyfoot[LO]{\myfancyfootOdd}
	\fancyfoot[RE]{\myfancyfootEven}
	\renewcommand{\headrulewidth}{0pt}
	\renewcommand{\footrulewidth}{0pt}
}
%=================== HẾT PHẦN GỘP TỪ eduzy-de.tex ===================
\newcommand{\indam}[1]{
	{\fontsize{12pt}{9pt}\sffamily \selectfont\bfseries\color{gold1} #1}
}
\newcommand{\ind}[1]{
	{\fontsize{14pt}{9pt}\fontfamily{qag} \selectfont\bfseries\color{gold1} #1}
}
\renewcommand{\FalseEX}{\stepcounter{dapan}\textbf{\selectfont\color{gold1}\Alph{dapan}.}\ignorespaces}
\renewtheorem{ex}{\color{gold1}Câu}
\renewcommand{\baselinestretch}{1.1}
\newcommand{\hoac}[1]{
	\left[\begin{aligned}#1\end{aligned}\right.}
\newcommand{\heva}[1]{
	\left\{\begin{aligned}#1\end{aligned}\right.}
\usepackage{esvect}
\def\vec{\vv}
\def\vt{\vv}
\def\overrightarrow{\vv}
%====================== LOGO EDUZY VẼ BẰNG TIKZ ======================
%  Đường viền dò từ logo gốc; dải màu vàng kim đo trực tiếp trên ảnh gốc
%  (nghiêng 21 độ, tô bằng 80 dải hẹp nên hiển thị giống nhau trên mọi trình xem PDF).
%  Dùng: \eduzylogo          -> cao 2.4cm (mặc định)
%        \eduzylogo[1.5cm]   -> cao 1.5cm
\definecolor{ezText}{RGB}{245,200,46}
\newcommand{\eduzylogo}[1][2.4cm]{%
	\resizebox{!}{#1}{%
		\begin{tikzpicture}[x=0.01cm,y=0.01cm]
			\path[use as bounding box] (0,0) rectangle (224,-224);
			% mũ tốt nghiệp + vương miện
			\begin{scope}
				\clip
				(119.4,-176.1)..controls(118.9,-176.1)and(116.8,-176)..(114.6,-176)
				..controls(101.4,-176)and(85.8,-175.9)..(83.7,-175.7)
				..controls(81.8,-175.6)and(81,-175.6)..(80.2,-175.7)
				..controls(78.8,-175.9)and(77.7,-175.9)..(76.2,-175.7)
				..controls(74.9,-175.5)and(70.8,-175.5)..(64.3,-175.7)
				..controls(59.1,-175.9)and(58.3,-175.7)..(58.6,-173.7)
				..controls(58.6,-172.8)and(59.6,-169.6)..(60.3,-168)
				..controls(60.8,-166.9)and(61.4,-165.8)..(62.3,-165)
				..controls(63.1,-164.2)and(64.3,-163.4)..(65,-163.2)
				..controls(65.5,-163.1)and(72.3,-163)..(76.6,-163.1)
				..controls(77.9,-163.1)and(80.1,-163.1)..(81.4,-163)
				..controls(82.8,-162.9)and(86.3,-162.9)..(89.2,-162.9)
				..controls(93,-162.9)and(94.9,-162.8)..(96,-162.7)
				..controls(97.2,-162.5)and(101,-162.5)..(114,-162.4)
				..controls(123,-162.3)and(134.9,-162.2)..(140.2,-162.1)
				..controls(150.1,-162)and(152.2,-162.1)..(153.6,-162.4)
				..controls(154.4,-162.6)and(157.1,-163.8)..(158.1,-164.4)
				..controls(158.8,-164.8)and(159.7,-165.7)..(159.9,-166.3)
				..controls(160,-166.5)and(160.2,-167.1)..(160.2,-167.6)
				..controls(160.3,-168)and(160.6,-169.4)..(160.9,-170.6)
				..controls(161.2,-171.7)and(161.5,-173.1)..(161.6,-173.6)
				..controls(161.8,-174.4)and(161.8,-174.5)..(161.5,-174.9)
				..controls(161.2,-175.6)and(160.8,-175.9)..(159.9,-176)
				..controls(159.2,-176.1)and(154.1,-176.2)..(151.5,-176.2)
				..controls(148.1,-176.1)and(122.6,-176.1)..(121.6,-176.2)
				..controls(120.8,-176.2)and(119.8,-176.2)..(119.4,-176.1)--cycle
				(130.4,-160)..controls(129,-160)and(119.8,-160)..(110.1,-159.9)
				..controls(100.3,-159.8)and(88.5,-159.7)..(83.9,-159.6)
				..controls(77.8,-159.5)and(75.3,-159.5)..(74.5,-159.6)
				..controls(73.9,-159.7)and(72.5,-159.8)..(71.3,-159.8)
				..controls(69.6,-159.8)and(69,-159.8)..(68.5,-159.6)
				..controls(67.6,-159.3)and(66.8,-158.6)..(66.2,-157.6)
				..controls(66,-157.3)and(65.6,-156.5)..(65.3,-156.1)
				..controls(64.6,-155.1)and(64.5,-154.9)..(63.7,-152.9)
				..controls(63.2,-151.6)and(62.9,-150.9)..(61.8,-149.6)
				..controls(61.5,-149.2)and(61.1,-148.5)..(60.9,-147.9)
				..controls(60.7,-147.4)and(60.3,-146.6)..(60.1,-146.1)
				..controls(59.9,-145.5)and(59.5,-144.9)..(59,-144.3)
				..controls(58.5,-143.7)and(58.2,-143.2)..(57.9,-142.5)
				..controls(57.7,-141.9)and(57.3,-141.1)..(57,-140.7)
				..controls(56.6,-140.3)and(56.1,-139.3)..(55.8,-138.5)
				..controls(55.5,-137.8)and(55.1,-136.8)..(54.9,-136.5)
				..controls(54.7,-136.1)and(54.3,-135.2)..(53.9,-134.5)
				..controls(53.6,-133.8)and(53,-132.7)..(52.6,-132)
				..controls(51.6,-130.5)and(51.4,-130.2)..(51.1,-129)
				..controls(50.6,-127.5)and(50.6,-127.4)..(49.7,-126.3)
				..controls(49.2,-125.7)and(48.6,-124.8)..(48.3,-124.2)
				..controls(48,-123.6)and(47.5,-122.7)..(47.3,-122.2)
				..controls(47,-121.8)and(46.5,-120.7)..(46.1,-119.8)
				..controls(45.8,-118.9)and(45.2,-117.8)..(44.9,-117.2)
				..controls(44.2,-116.1)and(43.3,-114.4)..(42.3,-112.2)
				..controls(41.8,-111.3)and(41.3,-110.4)..(41.1,-110.1)
				..controls(40.5,-109.4)and(40.3,-109)..(40,-107.8)
				..controls(39.8,-107)and(39.6,-106.5)..(39.2,-106)
				..controls(38.9,-105.6)and(38.4,-104.7)..(38,-104)
				..controls(37.7,-103.4)and(37.2,-102.5)..(37,-102.1)
				..controls(36.7,-101.7)and(36.3,-100.8)..(36,-100.1)
				..controls(35.6,-98.9)and(35.3,-98.4)..(34.6,-97.6)
				..controls(34.2,-97.1)and(34.2,-96.3)..(34.6,-96)
				..controls(35,-95.7)and(35.5,-95.9)..(36.2,-96.7)
				..controls(36.5,-97)and(37.2,-97.5)..(37.6,-97.7)..controls(38,-98)and(38.6,-98.4)..(39.1,-98.9)
				..controls(39.6,-99.4)and(40.3,-100)..(40.6,-100.3)
				..controls(41.3,-100.8)and(42.4,-101.8)..(43.7,-102.9)
				..controls(44,-103.2)and(44.6,-103.7)..(45,-103.9)
				..controls(45.4,-104.2)and(45.9,-104.6)..(46.1,-104.8)
				..controls(47.1,-105.9)and(48.8,-107.1)..(50.5,-107.9)
				..controls(51.8,-108.6)and(52.3,-109)..(53.4,-110)
				..controls(54.8,-111.2)and(55.6,-111.6)..(57,-112.1)
				..controls(57.5,-112.3)and(58.5,-112.8)..(59.4,-113.3)
				..controls(60.2,-113.7)and(61.3,-114.3)..(61.9,-114.5)
				..controls(62.4,-114.7)and(63.3,-115.1)..(63.8,-115.3)
				..controls(64.4,-115.6)and(65.4,-115.9)..(66.2,-116.1)
				..controls(67,-116.3)and(68.4,-116.6)..(69.2,-116.9)
				..controls(72.3,-117.9)and(72.8,-118)..(78,-118.1)
				..controls(84,-118.3)and(86.1,-118.1)..(88.1,-117.1)
				..controls(89.4,-116.5)and(93,-115)..(93.7,-114.9)
				..controls(94.8,-114.6)and(95.4,-114.2)..(95.9,-113.4)
				..controls(96.2,-113)and(96.7,-112.5)..(97.1,-112.2)
				..controls(98,-111.6)and(101.3,-108.4)..(102.1,-107.2)
				..controls(102.4,-106.8)and(102.9,-106.2)..(103.2,-105.8)
				..controls(103.5,-105.5)and(104,-104.8)..(104.2,-104.2)
				..controls(104.5,-103.6)and(105,-102.8)..(105.3,-102.4)
				..controls(105.6,-102)and(106.1,-101.2)..(106.5,-100.6)
				..controls(107.1,-99.5)and(108.6,-97.5)..(109.5,-96.5)
				..controls(109.9,-96.1)and(110.1,-96)..(110.4,-96)
				..controls(111.1,-96)and(111.6,-96.5)..(111.8,-97.6)
				..controls(112,-98.6)and(112.3,-99.3)..(112.8,-99.9)
				..controls(113.7,-100.9)and(115.2,-103.1)..(115.7,-104.2)
				..controls(116.1,-105)and(116.5,-105.5)..(116.9,-106)
				..controls(117.3,-106.4)and(117.8,-106.9)..(118,-107.3)
				..controls(118.2,-107.6)and(118.6,-108)..(119,-108.3)
				..controls(119.6,-108.8)and(121.1,-110.3)..(121.8,-111.2)
				..controls(122.1,-111.7)and(122.5,-112)..(123.4,-112.4)
				..controls(124.8,-113.2)and(125.8,-113.8)..(127,-114.6)
				..controls(127.5,-114.9)and(128.1,-115.3)..(128.2,-115.4)
				..controls(128.4,-115.5)and(129.3,-115.8)..(130.2,-116)
				..controls(131.1,-116.2)and(132.5,-116.6)..(133.3,-116.9)
				..controls(135.7,-117.9)and(136.1,-117.9)..(142.5,-117.8)
				..controls(145.8,-117.8)and(148.4,-117.7)..(148.8,-117.6)
				..controls(149.2,-117.6)and(150.2,-117.2)..(151,-116.8)
				..controls(152,-116.4)and(152.8,-116)..(153.6,-115.9)
				..controls(154.5,-115.7)and(155.2,-115.4)..(156.3,-114.9)
				..controls(157.1,-114.5)and(158.2,-114)..(158.7,-113.8)
				..controls(159.1,-113.7)and(159.8,-113.3)..(160.2,-113)
				..controls(160.6,-112.7)and(161.3,-112.4)..(161.7,-112.2)
				..controls(162.7,-111.9)and(163.2,-111.7)..(164.1,-111.1)
				..controls(164.5,-110.8)and(165.3,-110.4)..(165.9,-110)
				..controls(166.5,-109.7)and(167.3,-109.2)..(167.6,-108.9)
				..controls(167.9,-108.6)and(168.6,-108.2)..(169,-108)
				..controls(169.4,-107.8)and(170.2,-107.3)..(170.7,-106.9)
				..controls(171.2,-106.5)and(172.1,-105.8)..(172.7,-105.4)
				..controls(173.2,-105)and(174.1,-104.4)..(174.5,-104)
				..controls(174.9,-103.7)and(175.7,-103.1)..(176.2,-102.6)
				..controls(176.7,-102.2)and(177.4,-101.5)..(177.9,-101.2)
				..controls(178.3,-100.8)and(178.8,-100.4)..(179,-100.2)
				..controls(179.2,-100)and(179.6,-99.6)..(180,-99.3)
				..controls(180.4,-99.1)and(180.9,-98.6)..(181.2,-98.4)
				..controls(181.5,-98.1)and(182.1,-97.7)..(182.5,-97.5)
				..controls(183.2,-97.2)and(183.5,-96.9)..(184.3,-95.9)
				..controls(184.6,-95.6)and(184.6,-95.6)..(184.8,-95.7)
				..controls(185,-95.8)and(185.1,-96)..(185.1,-96.2)
				..controls(185.1,-96.5)and(185,-96.7)..(184.6,-97.2)
				..controls(184.1,-97.8)and(184,-97.9)..(183.9,-98.8)
				..controls(183.7,-100.3)and(183.3,-101.3)..(182.9,-101.8)
				..controls(182.6,-102.1)and(182.3,-102.6)..(182.2,-102.9)
				..controls(182,-103.3)and(181.6,-104)..(181.2,-104.6)
				..controls(180.9,-105.2)and(180.5,-106.1)..(180.4,-106.6)
				..controls(179.9,-108)and(179.7,-108.5)..(179,-109.4)
				..controls(178.6,-109.8)and(178.2,-110.5)..(177.9,-111)
				..controls(177.7,-111.5)and(177.3,-112.3)..(177,-112.7)
				..controls(176.6,-113.2)and(176.4,-113.7)..(176.3,-114.2)
				..controls(176.2,-114.7)and(175.9,-115.2)..(175.4,-116)
				..controls(175,-116.6)and(174.4,-117.7)..(174,-118.5)
				..controls(173.7,-119.2)and(173.3,-120.1)..(173,-120.5)
				..controls(172.8,-120.9)and(172.4,-121.8)..(172.1,-122.5)
				..controls(171.8,-123.2)and(171.3,-124.1)..(171,-124.6)
				..controls(170.8,-125.1)and(170.3,-126)..(169.9,-126.6)
				..controls(169.6,-127.3)and(169.2,-128.1)..(169,-128.3)
				..controls(168.8,-128.6)and(168.5,-129.3)..(168.3,-129.8)
				..controls(168.1,-130.3)and(167.8,-131.1)..(167.5,-131.6)
				..controls(166.9,-132.5)and(164.2,-137.9)..(164.2,-138.1)
				..controls(164.2,-138.1)and(163.9,-138.8)..(163.4,-139.7)
				..controls(162.9,-140.5)and(162.4,-141.5)..(162.3,-141.9)
				..controls(161.9,-143)and(161.6,-143.4)..(160.9,-144.1)
				..controls(160.4,-144.6)and(160.3,-144.8)..(160.1,-145.4)
				..controls(160,-145.8)and(159.6,-146.8)..(159.2,-147.8)
				..controls(158.7,-148.8)and(158.1,-150)..(157.9,-150.6)
				..controls(157.6,-151.1)and(157.2,-151.9)..(156.9,-152.4)
				..controls(156.5,-152.8)and(156.1,-153.7)..(155.8,-154.4)
				..controls(155.6,-155)and(155.2,-155.9)..(154.9,-156.3)
				..controls(154.7,-156.7)and(154.4,-157.4)..(154.2,-157.9)
				..controls(154.1,-158.4)and(153.8,-159)..(153.5,-159.3)
				..controls(152.8,-160.1)and(153,-160.1)..(142.4,-160.1)
				..controls(137.2,-160.1)and(131.8,-160.1)..(130.4,-160)--cycle
				(141.5,-103.7)..controls(140.8,-103)and(139.8,-102.1)..(139.2,-101.6)
				..controls(138.5,-101.2)and(137.9,-100.6)..(137.8,-100.4)
				..controls(137,-99.5)and(134,-96.5)..(133.2,-96)
				..controls(132.7,-95.7)and(131.8,-94.9)..(131.1,-94.3)
				..controls(129.4,-92.7)and(127.9,-91.6)..(126.6,-91.1)
				..controls(125.2,-90.5)and(124.1,-90)..(122.9,-89.1)
				..controls(122.1,-88.6)and(121.3,-88.2)..(119.9,-87.7)
				..controls(118.8,-87.3)and(117.5,-86.8)..(116.9,-86.5)
				..controls(114.3,-85.5)and(113.7,-85.4)..(109.7,-85.5)
				..controls(105.3,-85.7)and(103.9,-85.9)..(101.6,-87)
				..controls(100.7,-87.4)and(99.4,-87.8)..(97.9,-88)
				..controls(97,-88.2)and(96.6,-88.3)..(95.8,-88.9)
				..controls(95.4,-89.2)and(94.7,-89.6)..(94.2,-89.8)
				..controls(92.9,-90.3)and(92.3,-90.7)..(90.9,-92.1)
				..controls(90.2,-92.7)and(89.4,-93.3)..(89,-93.5)..controls(87,-94.5)and(85.2,-96)..(83.6,-97.9)
				..controls(83.3,-98.2)and(82.5,-98.9)..(81.8,-99.4)
				..controls(80.7,-100.1)and(80.5,-100.3)..(80,-101.1)
				..controls(78.7,-103.1)and(77.2,-103.9)..(76.4,-103)
				..controls(75.7,-102.2)and(75.7,-103.2)..(75.7,-83.6)
				..controls(75.6,-73.6)and(75.6,-64.8)..(75.5,-64)
				..controls(75.5,-63.2)and(75.5,-61.4)..(75.6,-59.8)--(75.7,-57.1)--(76,-56.8)
				..controls(76.8,-56.2)and(78,-56.3)..(79.8,-57.2)
				..controls(80.4,-57.5)and(81.3,-57.8)..(81.8,-57.9)
				..controls(82.3,-58.1)and(83,-58.4)..(83.6,-58.8)..controls(84.3,-59.2)and(85,-59.5)..(86,-59.8)
				..controls(86.8,-60.1)and(88,-60.6)..(88.7,-60.9)
				..controls(89.3,-61.2)and(90.6,-61.7)..(91.4,-62)
				..controls(92.2,-62.3)and(93.7,-62.9)..(94.6,-63.3)
				..controls(95.6,-63.7)and(96.8,-64.1)..(97.4,-64.3)
				..controls(98,-64.5)and(99,-64.9)..(99.7,-65.2)
				..controls(100.4,-65.6)and(101.5,-66)..(102.3,-66.2)
				..controls(104.2,-66.8)and(105,-67.1)..(106.2,-67.8)
				..controls(107.9,-68.7)and(110.2,-69.4)..(111.3,-69.2)
				..controls(112.5,-69)and(115,-68.2)..(115.6,-67.8)
				..controls(117.3,-66.7)and(117.6,-66.6)..(119.5,-66.2)
				..controls(121.9,-65.8)and(122.1,-65.7)..(123,-65.1)
				..controls(124.1,-64.4)and(124.4,-64.3)..(126,-63.9)
				..controls(127,-63.7)and(127.7,-63.4)..(128.4,-63.1)
				..controls(129,-62.8)and(129.7,-62.5)..(130.6,-62.4)
				..controls(132.1,-62.1)and(134,-61.5)..(135.2,-61)
				..controls(135.7,-60.8)and(136.7,-60.5)..(137.3,-60.3)
				..controls(138.4,-60)and(139.3,-59.7)..(141.1,-59)
				..controls(141.6,-58.7)and(142,-58.6)..(142.5,-58.6)
				..controls(143.2,-58.6)and(143.3,-58.7)..(143.6,-59)
				..controls(144,-59.5)and(144.1,-59.9)..(144,-60.8)
				..controls(144,-61.2)and(144,-71.1)..(144,-82.8)--(143.9,-104.1)--(143.7,-104.5)
				..controls(143.5,-104.7)and(143.3,-104.8)..(143,-104.9)
				..controls(142.6,-104.9)and(142.5,-104.8)..(141.5,-103.7)--cycle
				(27.8,-91.4)..controls(27.1,-90.9)and(26.3,-90)..(26.1,-89.5)
				..controls(26,-89.2)and(25.8,-88.6)..(25.7,-87.9)
				..controls(25.4,-86.2)and(25.6,-84.8)..(26.5,-83.5)
				..controls(27.4,-82.3)and(28.3,-82)..(29.4,-82.3)
				..controls(31.1,-82.8)and(32.1,-84.9)..(32,-87.4)
				..controls(31.9,-89.3)and(31.2,-90.6)..(29.9,-91.4)
				..controls(29.1,-91.9)and(28.6,-91.9)..(27.8,-91.4)--cycle
				(190.1,-91.3)..controls(189.7,-91.2)and(189.3,-90.9)..(189.1,-90.7)
				..controls(188.3,-89.8)and(187.6,-87.7)..(187.6,-86.2)
				..controls(187.6,-85.5)and(187.7,-85.2)..(188.2,-84.4)
				..controls(189.1,-82.6)and(190.1,-81.8)..(191.3,-82.1)
				..controls(192,-82.3)and(193.5,-83.7)..(193.8,-84.4)
				..controls(194.1,-85.1)and(194.2,-87.1)..(194,-88.3)
				..controls(193.8,-89.8)and(193.3,-90.5)..(192.3,-91.3)
				..controls(191.7,-91.7)and(190.9,-91.7)..(190.1,-91.3)--cycle
				(173.7,-83.4)..controls(172.2,-82.5)and(172,-82.5)..(170.6,-82.4)
				..controls(169.7,-82.4)and(168.8,-82.4)..(168.3,-82.5)
				..controls(167.2,-82.7)and(166.9,-82.7)..(166.6,-82.2)
				..controls(166.3,-81.8)and(166.3,-81)..(166.8,-80)
				..controls(167.2,-78.9)and(167.4,-77.4)..(167.7,-74.1)
				..controls(167.8,-72.8)and(167.9,-71.8)..(167.9,-71.7)
				..controls(168,-71.7)and(168,-70.2)..(168,-68.4)..controls(168,-65.4)and(168,-64.9)..(167.7,-63)
				..controls(167.3,-60.5)and(167.1,-59.9)..(166.4,-58.4)
				..controls(166.1,-57.7)and(165.8,-56.9)..(165.7,-56.6)
				..controls(165.7,-56.2)and(165.4,-55.6)..(165.2,-55.2)
				..controls(164.6,-53.8)and(161.9,-50.7)..(161,-50.3)
				..controls(159.8,-49.7)and(157.9,-50)..(155.2,-51.2)
				..controls(154.7,-51.4)and(153.4,-51.8)..(152.2,-52.1)
				..controls(150.8,-52.4)and(149.4,-52.9)..(148.6,-53.2)
				..controls(147.8,-53.6)and(146.7,-54)..(146,-54.2)
				..controls(145.4,-54.4)and(144.5,-54.7)..(144,-54.9)
				..controls(143.6,-55.2)and(142.7,-55.5)..(141.8,-55.7)
				..controls(141,-55.9)and(139.6,-56.3)..(138.7,-56.6)
				..controls(137.8,-57)and(136.3,-57.4)..(135.3,-57.6)
				..controls(133.8,-57.9)and(133.2,-58.1)..(132.1,-58.7)
				..controls(131.1,-59.2)and(130.3,-59.4)..(129.4,-59.6)
				..controls(128.7,-59.8)and(127.8,-60)..(127.4,-60.1)
				..controls(127,-60.3)and(126,-60.5)..(125.1,-60.7)
				..controls(123.2,-61.1)and(121.7,-61.5)..(120.2,-62.1)
				..controls(119.7,-62.3)and(118.1,-62.8)..(116.7,-63.2)
				..controls(114.6,-63.9)and(114,-64.1)..(113,-64.7)
				..controls(111.9,-65.4)and(111.8,-65.4)..(111,-65.5)
				..controls(109.9,-65.5)and(109.3,-65.4)..(107.7,-64.7)
				..controls(107,-64.4)and(106,-64)..(105.3,-63.9)
				..controls(103.5,-63.5)and(102.2,-63)..(101.2,-62.5)
				..controls(100.1,-61.8)and(99.3,-61.5)..(97.8,-61.1)
				..controls(97.2,-61)and(96.2,-60.7)..(95.7,-60.4)
				..controls(95.1,-60.1)and(94.2,-59.8)..(93.5,-59.6)
				..controls(92.9,-59.4)and(91.7,-59)..(91,-58.7)..controls(90.2,-58.4)and(88.9,-58)..(88,-57.7)
				..controls(87.2,-57.5)and(86.1,-57.1)..(85.5,-56.8)
				..controls(85,-56.6)and(84.1,-56.2)..(83.5,-56)..controls(82.9,-55.7)and(81.8,-55.2)..(81,-54.9)
				..controls(80.1,-54.5)and(79.1,-54.2)..(78.7,-54.1)
				..controls(77.9,-53.8)and(77.3,-53.6)..(76.2,-52.9)
				..controls(75.7,-52.6)and(75.2,-52.4)..(74.1,-52.2)
				..controls(72.5,-51.9)and(72.4,-51.8)..(71.3,-51.1)
				..controls(69.9,-50.1)and(69.4,-49.9)..(67.9,-49.8)
				..controls(67,-49.7)and(66.6,-49.6)..(66.3,-49.4)
				..controls(66.1,-49.3)and(65.7,-49)..(65.4,-48.8)
				..controls(65.1,-48.7)and(64.2,-48.4)..(63.6,-48.2)
				..controls(62.9,-48)and(61.7,-47.6)..(60.9,-47.2)
				..controls(60.1,-46.8)and(58.9,-46.4)..(58.2,-46.2)
				..controls(57.6,-46)and(56.8,-45.7)..(56.5,-45.5)
				..controls(54.9,-44.6)and(54.6,-44.5)..(53.2,-44.2)
				..controls(51.7,-43.8)and(51,-43.6)..(49.9,-42.9)
				..controls(49.5,-42.6)and(48.9,-42.4)..(47.5,-42)
				..controls(45.7,-41.6)and(45.6,-41.5)..(44.3,-40.6)
				..controls(43.5,-40.1)and(42.8,-39.6)..(42.7,-39.4)
				..controls(42.4,-38.7)and(43,-38.1)..(44.5,-37.6)..controls(44.9,-37.5)and(46,-37)..(46.9,-36.6)
				..controls(47.9,-36.1)and(48.8,-35.8)..(49.6,-35.6)
				..controls(50.5,-35.4)and(51.3,-35.2)..(52.4,-34.6)
				..controls(53.3,-34.1)and(54.5,-33.7)..(54.9,-33.5)
				..controls(55.4,-33.4)and(56.4,-33)..(57.2,-32.6)
				..controls(58.6,-32)and(60.1,-31.5)..(61.6,-31.2)
				..controls(62.3,-31.1)and(62.7,-31)..(63.2,-30.7)
				..controls(64.1,-30)and(64.3,-29.9)..(65.3,-29.7)
				..controls(65.7,-29.6)and(66.5,-29.2)..(67.1,-28.9)
				..controls(67.6,-28.6)and(68.8,-28.2)..(69.6,-27.9)
				..controls(70.5,-27.7)and(71.6,-27.3)..(72.1,-27)
				..controls(72.6,-26.7)and(73.8,-26.2)..(74.8,-25.9)
				..controls(75.7,-25.5)and(76.7,-25.1)..(77.1,-25)
				..controls(77.4,-24.8)and(78.5,-24.4)..(79.6,-24.1)
				..controls(80.6,-23.8)and(82,-23.2)..(82.8,-22.9)
				..controls(83.5,-22.5)and(84.8,-22)..(85.8,-21.8)
				..controls(86.7,-21.5)and(87.8,-21.1)..(88.1,-20.9)
				..controls(88.4,-20.7)and(89.3,-20.3)..(90.1,-20.1)
				..controls(90.8,-19.8)and(91.9,-19.4)..(92.4,-19.1)
				..controls(93.5,-18.5)and(94.4,-18.2)..(95.4,-18.1)
				..controls(95.7,-18)and(96.4,-17.8)..(96.9,-17.5)
				..controls(97.5,-17.3)and(98.4,-16.9)..(99,-16.6)
				..controls(99.7,-16.3)and(100.9,-15.9)..(101.9,-15.6)
				..controls(102.9,-15.3)and(104.4,-14.7)..(105.3,-14.3)
				..controls(108.1,-13.2)and(108.5,-13.1)..(110.1,-13.1)
				..controls(111.4,-13.1)and(111.5,-13.1)..(112.4,-13.5)
				..controls(113.2,-14)and(115,-14.8)..(116.6,-15.4)
				..controls(116.9,-15.5)and(117.9,-15.8)..(118.8,-16.1)
				..controls(120.9,-16.6)and(121.5,-16.9)..(122.6,-17.9)
				..controls(123,-18.3)and(123.7,-18.8)..(123.9,-18.9)
				..controls(124.6,-19.4)and(126,-19.8)..(127.9,-20.2)
				..controls(129.6,-20.5)and(129.8,-20.6)..(131.1,-21.4)
				..controls(131.5,-21.7)and(132.4,-22.1)..(133.1,-22.3)
				..controls(136.1,-23.2)and(138.3,-24.1)..(140.2,-25)
				..controls(141,-25.5)and(141.5,-25.6)..(142.5,-25.8)
				..controls(145.1,-26.3)and(146,-26.6)..(149.4,-28.1)
				..controls(151.3,-28.9)and(153.5,-29.8)..(154.3,-30.1)
				..controls(155.1,-30.4)and(156.2,-30.8)..(156.6,-31)
				..controls(157.1,-31.2)and(158.2,-31.6)..(159.1,-31.8)
				..controls(160.6,-32.2)and(160.9,-32.3)..(161.7,-32.9)
				..controls(162.4,-33.3)and(162.9,-33.6)..(163.7,-33.8)
				..controls(164.3,-34)and(164.8,-34.1)..(164.9,-34.1)
				..controls(165,-34.1)and(165.6,-34.4)..(166.4,-34.8)
				..controls(167.3,-35.2)and(168.1,-35.4)..(169,-35.6)
				..controls(170,-35.8)and(170.5,-36)..(171.7,-36.6)
				..controls(172.4,-37)and(173.3,-37.4)..(173.6,-37.4)
				..controls(175.6,-37.7)and(176.2,-38)..(176.7,-39.1)
				..controls(176.9,-39.6)and(176.9,-39.7)..(176.7,-40.4)
				..controls(176.4,-41.6)and(175.8,-42.7)..(175.2,-43.4)
				..controls(174.5,-44.1)and(172.6,-45.5)..(171.7,-45.8)
				..controls(170.6,-46.2)and(169.1,-46)..(168.2,-45.4)
				..controls(167.9,-45.2)and(167.3,-44.8)..(166.9,-44.4)
				..controls(166.4,-44.1)and(165.8,-43.6)..(165.4,-43.2)
				..controls(165.1,-42.9)and(164.4,-42.3)..(163.9,-41.9)
				..controls(163.3,-41.5)and(162.7,-40.9)..(162.4,-40.5)
				..controls(161.4,-39.2)and(159.9,-38)..(158.8,-37.8)
				..controls(158.4,-37.8)and(158.3,-37.7)..(158.2,-37.2)
				..controls(158,-36.8)and(157.4,-36.1)..(157.1,-36.1)
				..controls(157,-36.1)and(156.7,-36.1)..(156.5,-36)
				..controls(156.1,-35.9)and(155.9,-35.7)..(155.7,-35.2)
				..controls(155.3,-34.6)and(155.1,-34.4)..(154,-34.1)
				..controls(153.6,-34)and(153.2,-33.7)..(152.9,-33.5)
				..controls(152.4,-32.9)and(151.6,-32.4)..(151.3,-32.4)
				..controls(150.9,-32.4)and(150.4,-32.6)..(150.4,-32.9)
				..controls(150.4,-33.2)and(150.8,-33.4)..(151.4,-33.6)
				..controls(152.1,-33.7)and(152.8,-34.1)..(153.5,-34.9)
				..controls(154.1,-35.5)and(154.3,-35.6)..(155,-35.8)
				..controls(155.7,-36)and(155.7,-36)..(156,-36.6)
				..controls(156.3,-37.3)and(156.7,-37.7)..(157.5,-37.8)
				..controls(158.1,-37.9)and(158.2,-38)..(158.4,-38.8)
				..controls(158.5,-39.4)and(158.6,-39.6)..(159.6,-40.4)
				..controls(160.2,-41)and(160.9,-41.6)..(161.3,-41.8)
				..controls(161.7,-42)and(161.9,-42.3)..(162.2,-42.8)
				..controls(162.4,-43.2)and(162.8,-43.8)..(163.1,-44.1)
				..controls(163.9,-44.8)and(165.1,-46.7)..(165.7,-48)
				..controls(166.2,-49.3)and(166.4,-49.7)..(167.2,-50.5)
				..controls(167.5,-50.8)and(167.9,-51.4)..(168.1,-51.8)
				..controls(168.3,-52.2)and(168.8,-53)..(169.1,-53.6)
				..controls(169.6,-54.3)and(169.8,-54.9)..(169.9,-55.6)
				..controls(170.4,-57.8)and(170.5,-58.1)..(171.2,-59.5)
				..controls(171.6,-60.3)and(172,-61.5)..(172.2,-62.1)
				..controls(172.3,-62.8)and(172.7,-63.7)..(172.9,-64.1)
				..controls(173.2,-64.6)and(173.4,-65.3)..(173.4,-65.6)
				..controls(173.8,-68.8)and(174.2,-70.6)..(174.9,-72.6)
				..controls(175.5,-74.5)and(175.9,-77.1)..(176.1,-80.2)
				..controls(176.3,-83.4)and(175.5,-84.4)..(173.7,-83.4)--cycle;
				\definecolor{ezs}{RGB}{247,179,26}\fill[ezs](-5.7,-156.6)--(-3.4,-157.5)--(61.9,12.7)--(59.7,13.6)--cycle;
				\definecolor{ezs}{RGB}{247,179,22}\fill[ezs](-3.8,-157.3)--(-1.5,-158.2)--(63.9,12)--(61.6,12.9)--cycle;
				\definecolor{ezs}{RGB}{248,179,19}\fill[ezs](-1.8,-158.1)--(0.4,-158.9)--(65.8,11.2)--(63.5,12.1)--cycle;
				\definecolor{ezs}{RGB}{248,178,16}\fill[ezs](0.1,-158.8)--(2.4,-159.7)--(67.7,10.5)--(65.4,11.4)--cycle;
				\definecolor{ezs}{RGB}{249,178,15}\fill[ezs](2,-159.5)--(4.3,-160.4)--(69.6,9.8)--(67.3,10.7)--cycle;
				\definecolor{ezs}{RGB}{251,177,17}\fill[ezs](3.9,-160.3)--(6.2,-161.1)--(71.5,9)--(69.2,9.9)--cycle;
				\definecolor{ezs}{RGB}{252,179,18}\fill[ezs](5.8,-161)--(8.1,-161.9)--(73.4,8.3)--(71.1,9.2)--cycle;
				\definecolor{ezs}{RGB}{252,182,17}\fill[ezs](7.7,-161.7)--(10,-162.6)--(75.3,7.6)--(73,8.5)--cycle;
				\definecolor{ezs}{RGB}{252,185,17}\fill[ezs](9.6,-162.5)--(11.9,-163.4)--(77.2,6.8)--(75,7.7)--cycle;
				\definecolor{ezs}{RGB}{251,189,18}\fill[ezs](11.5,-163.2)--(13.8,-164.1)--(79.2,6.1)--(76.9,7)--cycle;
				\definecolor{ezs}{RGB}{251,191,22}\fill[ezs](13.4,-163.9)--(15.7,-164.8)--(81.1,5.4)--(78.8,6.2)--cycle;
				\definecolor{ezs}{RGB}{252,193,25}\fill[ezs](15.4,-164.7)--(17.6,-165.6)--(83,4.6)--(80.7,5.5)--cycle;
				\definecolor{ezs}{RGB}{253,195,27}\fill[ezs](17.3,-165.4)--(19.6,-166.3)--(84.9,3.9)--(82.6,4.8)--cycle;
				\definecolor{ezs}{RGB}{253,196,31}\fill[ezs](19.2,-166.1)--(21.5,-167)--(86.8,3.2)--(84.5,4)--cycle;
				\definecolor{ezs}{RGB}{254,199,35}\fill[ezs](21.1,-166.9)--(23.4,-167.8)--(88.7,2.4)--(86.4,3.3)--cycle;
				\definecolor{ezs}{RGB}{253,201,38}\fill[ezs](23,-167.6)--(25.3,-168.5)--(90.6,1.7)--(88.3,2.6)--cycle;
				\definecolor{ezs}{RGB}{253,203,40}\fill[ezs](24.9,-168.3)--(27.2,-169.2)--(92.5,1)--(90.3,1.8)--cycle;
				\definecolor{ezs}{RGB}{253,205,42}\fill[ezs](26.8,-169.1)--(29.1,-170)--(94.4,0.2)--(92.2,1.1)--cycle;
				\definecolor{ezs}{RGB}{253,207,45}\fill[ezs](28.7,-169.8)--(31,-170.7)--(96.4,-0.5)--(94.1,0.4)--cycle;
				\definecolor{ezs}{RGB}{253,210,49}\fill[ezs](30.7,-170.5)--(32.9,-171.4)--(98.3,-1.2)--(96,-0.4)--cycle;
				\definecolor{ezs}{RGB}{253,212,54}\fill[ezs](32.6,-171.3)--(34.9,-172.2)--(100.2,-2)--(97.9,-1.1)--cycle;
				\definecolor{ezs}{RGB}{253,214,59}\fill[ezs](34.5,-172)--(36.8,-172.9)--(102.1,-2.7)--(99.8,-1.8)--cycle;
				\definecolor{ezs}{RGB}{253,216,63}\fill[ezs](36.4,-172.7)--(38.7,-173.6)--(104,-3.4)--(101.7,-2.6)--cycle;
				\definecolor{ezs}{RGB}{253,218,65}\fill[ezs](38.3,-173.5)--(40.6,-174.4)--(105.9,-4.2)--(103.6,-3.3)--cycle;
				\definecolor{ezs}{RGB}{254,220,66}\fill[ezs](40.2,-174.2)--(42.5,-175.1)--(107.8,-4.9)--(105.5,-4)--cycle;
				\definecolor{ezs}{RGB}{254,222,68}\fill[ezs](42.1,-175)--(44.4,-175.8)--(109.7,-5.6)--(107.5,-4.8)--cycle;
				\definecolor{ezs}{RGB}{254,224,71}\fill[ezs](44,-175.7)--(46.3,-176.6)--(111.7,-6.4)--(109.4,-5.5)--cycle;
				\definecolor{ezs}{RGB}{254,226,75}\fill[ezs](46,-176.4)--(48.2,-177.3)--(113.6,-7.1)--(111.3,-6.2)--cycle;
				\definecolor{ezs}{RGB}{254,228,79}\fill[ezs](47.9,-177.2)--(50.1,-178)--(115.5,-7.8)--(113.2,-7)--cycle;
				\definecolor{ezs}{RGB}{254,230,82}\fill[ezs](49.8,-177.9)--(52.1,-178.8)--(117.4,-8.6)--(115.1,-7.7)--cycle;
				\definecolor{ezs}{RGB}{254,232,85}\fill[ezs](51.7,-178.6)--(54,-179.5)--(119.3,-9.3)--(117,-8.4)--cycle;
				\definecolor{ezs}{RGB}{254,234,88}\fill[ezs](53.6,-179.4)--(55.9,-180.2)--(121.2,-10)--(118.9,-9.2)--cycle;
				\definecolor{ezs}{RGB}{254,237,91}\fill[ezs](55.5,-180.1)--(57.8,-181)--(123.1,-10.8)--(120.8,-9.9)--cycle;
				\definecolor{ezs}{RGB}{254,239,94}\fill[ezs](57.4,-180.8)--(59.7,-181.7)--(125,-11.5)--(122.8,-10.6)--cycle;
				\definecolor{ezs}{RGB}{254,241,97}\fill[ezs](59.3,-181.6)--(61.6,-182.4)--(127,-12.2)--(124.7,-11.4)--cycle;
				\definecolor{ezs}{RGB}{254,243,101}\fill[ezs](61.2,-182.3)--(63.5,-183.2)--(128.9,-13)--(126.6,-12.1)--cycle;
				\definecolor{ezs}{RGB}{254,245,103}\fill[ezs](63.2,-183)--(65.4,-183.9)--(130.8,-13.7)--(128.5,-12.8)--cycle;
				\definecolor{ezs}{RGB}{254,245,105}\fill[ezs](65.1,-183.8)--(67.4,-184.6)--(132.7,-14.4)--(130.4,-13.6)--cycle;
				\definecolor{ezs}{RGB}{254,245,105}\fill[ezs](67,-184.5)--(69.3,-185.4)--(134.6,-15.2)--(132.3,-14.3)--cycle;
				\definecolor{ezs}{RGB}{254,244,103}\fill[ezs](68.9,-185.2)--(71.2,-186.1)--(136.5,-15.9)--(134.2,-15)--cycle;
				\definecolor{ezs}{RGB}{253,243,101}\fill[ezs](70.8,-186)--(73.1,-186.8)--(138.4,-16.6)--(136.1,-15.8)--cycle;
				\definecolor{ezs}{RGB}{253,241,98}\fill[ezs](72.7,-186.7)--(75,-187.6)--(140.3,-17.4)--(138,-16.5)--cycle;
				\definecolor{ezs}{RGB}{253,239,95}\fill[ezs](74.6,-187.4)--(76.9,-188.3)--(142.2,-18.1)--(140,-17.2)--cycle;
				\definecolor{ezs}{RGB}{253,236,92}\fill[ezs](76.5,-188.2)--(78.8,-189)--(144.2,-18.8)--(141.9,-18)--cycle;
				\definecolor{ezs}{RGB}{253,234,88}\fill[ezs](78.5,-188.9)--(80.7,-189.8)--(146.1,-19.6)--(143.8,-18.7)--cycle;
				\definecolor{ezs}{RGB}{253,232,85}\fill[ezs](80.4,-189.6)--(82.7,-190.5)--(148,-20.3)--(145.7,-19.4)--cycle;
				\definecolor{ezs}{RGB}{253,230,82}\fill[ezs](82.3,-190.4)--(84.6,-191.2)--(149.9,-21)--(147.6,-20.2)--cycle;
				\definecolor{ezs}{RGB}{253,228,78}\fill[ezs](84.2,-191.1)--(86.5,-192)--(151.8,-21.8)--(149.5,-20.9)--cycle;
				\definecolor{ezs}{RGB}{253,226,74}\fill[ezs](86.1,-191.8)--(88.4,-192.7)--(153.7,-22.5)--(151.4,-21.6)--cycle;
				\definecolor{ezs}{RGB}{253,224,70}\fill[ezs](88,-192.6)--(90.3,-193.4)--(155.6,-23.3)--(153.3,-22.4)--cycle;
				\definecolor{ezs}{RGB}{253,222,67}\fill[ezs](89.9,-193.3)--(92.2,-194.2)--(157.5,-24)--(155.3,-23.1)--cycle;
				\definecolor{ezs}{RGB}{253,220,65}\fill[ezs](91.8,-194)--(94.1,-194.9)--(159.5,-24.7)--(157.2,-23.8)--cycle;
				\definecolor{ezs}{RGB}{253,217,63}\fill[ezs](93.7,-194.8)--(96,-195.6)--(161.4,-25.5)--(159.1,-24.6)--cycle;
				\definecolor{ezs}{RGB}{253,215,60}\fill[ezs](95.7,-195.5)--(97.9,-196.4)--(163.3,-26.2)--(161,-25.3)--cycle;
				\definecolor{ezs}{RGB}{253,213,56}\fill[ezs](97.6,-196.2)--(99.9,-197.1)--(165.2,-26.9)--(162.9,-26)--cycle;
				\definecolor{ezs}{RGB}{252,211,53}\fill[ezs](99.5,-197)--(101.8,-197.8)--(167.1,-27.7)--(164.8,-26.8)--cycle;
				\definecolor{ezs}{RGB}{252,209,49}\fill[ezs](101.4,-197.7)--(103.7,-198.6)--(169,-28.4)--(166.7,-27.5)--cycle;
				\definecolor{ezs}{RGB}{252,207,45}\fill[ezs](103.3,-198.4)--(105.6,-199.3)--(170.9,-29.1)--(168.6,-28.2)--cycle;
				\definecolor{ezs}{RGB}{253,205,42}\fill[ezs](105.2,-199.2)--(107.5,-200)--(172.8,-29.9)--(170.6,-29)--cycle;
				\definecolor{ezs}{RGB}{253,202,40}\fill[ezs](107.1,-199.9)--(109.4,-200.8)--(174.7,-30.6)--(172.5,-29.7)--cycle;
				\definecolor{ezs}{RGB}{253,200,38}\fill[ezs](109,-200.6)--(111.3,-201.5)--(176.7,-31.3)--(174.4,-30.4)--cycle;
				\definecolor{ezs}{RGB}{253,198,35}\fill[ezs](111,-201.4)--(113.2,-202.2)--(178.6,-32.1)--(176.3,-31.2)--cycle;
				\definecolor{ezs}{RGB}{253,196,31}\fill[ezs](112.9,-202.1)--(115.2,-203)--(180.5,-32.8)--(178.2,-31.9)--cycle;
				\definecolor{ezs}{RGB}{253,194,28}\fill[ezs](114.8,-202.8)--(117.1,-203.7)--(182.4,-33.5)--(180.1,-32.6)--cycle;
				\definecolor{ezs}{RGB}{253,192,26}\fill[ezs](116.7,-203.6)--(119,-204.5)--(184.3,-34.3)--(182,-33.4)--cycle;
				\definecolor{ezs}{RGB}{253,189,24}\fill[ezs](118.6,-204.3)--(120.9,-205.2)--(186.2,-35)--(183.9,-34.1)--cycle;
				\definecolor{ezs}{RGB}{253,186,20}\fill[ezs](120.5,-205)--(122.8,-205.9)--(188.1,-35.7)--(185.8,-34.8)--cycle;
				\definecolor{ezs}{RGB}{253,184,17}\fill[ezs](122.4,-205.8)--(124.7,-206.7)--(190,-36.5)--(187.8,-35.6)--cycle;
				\definecolor{ezs}{RGB}{252,183,15}\fill[ezs](124.3,-206.5)--(126.6,-207.4)--(192,-37.2)--(189.7,-36.3)--cycle;
				\definecolor{ezs}{RGB}{252,181,14}\fill[ezs](126.3,-207.2)--(128.5,-208.1)--(193.9,-37.9)--(191.6,-37.1)--cycle;
				\definecolor{ezs}{RGB}{252,179,13}\fill[ezs](128.2,-208)--(130.4,-208.9)--(195.8,-38.7)--(193.5,-37.8)--cycle;
				\definecolor{ezs}{RGB}{252,178,12}\fill[ezs](130.1,-208.7)--(132.4,-209.6)--(197.7,-39.4)--(195.4,-38.5)--cycle;
				\definecolor{ezs}{RGB}{251,178,10}\fill[ezs](132,-209.4)--(134.3,-210.3)--(199.6,-40.1)--(197.3,-39.3)--cycle;
				\definecolor{ezs}{RGB}{251,178,10}\fill[ezs](133.9,-210.2)--(136.2,-211.1)--(201.5,-40.9)--(199.2,-40)--cycle;
				\definecolor{ezs}{RGB}{250,177,11}\fill[ezs](135.8,-210.9)--(138.1,-211.8)--(203.4,-41.6)--(201.1,-40.7)--cycle;
				\definecolor{ezs}{RGB}{250,178,11}\fill[ezs](137.7,-211.6)--(140,-212.5)--(205.3,-42.3)--(203.1,-41.5)--cycle;
				\definecolor{ezs}{RGB}{248,180,10}\fill[ezs](139.6,-212.4)--(141.9,-213.3)--(207.3,-43.1)--(205,-42.2)--cycle;
				\definecolor{ezs}{RGB}{247,181,9}\fill[ezs](141.5,-213.1)--(143.8,-214)--(209.2,-43.8)--(206.9,-42.9)--cycle;
				\definecolor{ezs}{RGB}{246,182,9}\fill[ezs](143.5,-213.8)--(145.7,-214.7)--(211.1,-44.5)--(208.8,-43.7)--cycle;
				\definecolor{ezs}{RGB}{246,182,10}\fill[ezs](145.4,-214.6)--(147.7,-215.5)--(213,-45.3)--(210.7,-44.4)--cycle;
			\end{scope}
			% chữ EDUZY (Times đậm, làm dày nét bằng 5 lớp lệch nhẹ; căn đúng khung chữ logo gốc)
			\foreach \dx/\dy in {-0.7/0,0.7/0,0/-0.3,0/0.3,0/0}
			\node[anchor=center,inner sep=0pt,outer sep=0pt,text=ezText,align=center] at (109.6+\dx,-193.6+\dy)
			{\resizebox{1.01cm}{0.22cm}{\fontencoding{T1}\fontfamily{ptm}\fontseries{b}\selectfont EDUZY}};
		\end{tikzpicture}%
	}%
}
%====================== THIẾT LẬP THÂN ĐỀ ======================
%--- Môi trường tự luận {bt}: ex_test 3.3.4 chưa khai báo sẵn nên tạo ở đây
%    (ex_test đã tự móc \fixshowans và \listenumerate cho bt)
\makeatletter
\@ifundefined{bt}{\newtheorem{bt}{\color{gold1}\namebt}}{}
\makeatother
\renewcommand{\sectionmark}[1]{\markright{\thesection\ #1}}
%--- Khung tiêu đề (logo + tối đa 4 dòng; dòng 3, 4 để trống {} thì bỏ qua)
%  Đề thi:    \khungtieude{THI THỬ THPT QUỐC GIA NĂM 2026}{Môn Toán -- Đề số 21}
%                         {(Đề thi gồm: ...)}{Thời gian làm bài: 90 phút (không kể thời gian phát đề).}
%  Chuyên đề: \khungtieude{CHUYÊN ĐỀ 3 -- DÃY SỐ VÀ CÔNG THỨC TRUY HỒI}
%                         {Bồi dưỡng học sinh giỏi Toán THPT}{...}{}
\newcommand{\khungtieude}[4]{%
	\begin{tcolorbox}[colframe=gold1,colback=gold5!45,boxrule=0.5mm]
		\begin{minipage}[c]{3cm}
			\centerline{\eduzylogo[2.4cm]}
		\end{minipage}\hspace{0.5cm}
		\begin{minipage}[c]{12.5cm}
			\vskip 0.15 cm
			\centerline{\ind{#1}}
			\vskip 0.15 cm
			\centerline{\textbf{#2}}
			\if\relax\detokenize{#3}\relax\else
				\vskip 0.15 cm
				\centerline{\textit{#3}}
			\fi
			\if\relax\detokenize{#4}\relax\else
				\vskip 0.15 cm
				\centerline{\textit{#4}}
			\fi
			\vskip 0.15 cm
		\end{minipage}
	\end{tcolorbox}
}
%====================== BỔ SUNG CHO BÀI HỌC, CHUYÊN ĐỀ ======================
\usepackage[shortlabels]{enumitem}
\usepackage{tabularx}
%--- Danh sách gọn: \begin{enumerate}[gon] ... \end{enumerate}
\SetEnumitemKey{gon}{leftmargin=*,itemsep=1pt,parsep=0pt,topsep=3pt}
%--- Khung lý thuyết: \begin{khung}{} ... \end{khung}; tiêu đề chỉ dùng khi thật ngắn ("Mẹo", "Lỗi thường gặp")
\newtcolorbox{khung}[1]{breakable,enhanced,colframe=gold1,colback=gold5!45,
	boxrule=0.4mm,fonttitle=\bfseries\sffamily,coltitle=white,colbacktitle=gold1,
	attach boxed title to top left={xshift=6mm,yshift=-3mm},
	boxed title style={arc=1mm,boxrule=0.3mm},title={#1}}
%--- Tiêu đề mục nhỏ trong một dạng: \muc{Dấu hiệu}, \muc{Phương pháp}, \muc{Bài tự luyện}
\makeatletter
\newcommand{\muc}[1]{\par\medskip\noindent\indam{#1}\par\nobreak\smallskip\@afterindentfalse\@afterheading}
\makeatother
%--- Ô công thức cần nhớ (căn giữa, viền vàng, vừa khít nội dung)
%     \begin{center}\congthuc{$Q(n)=(n-1)!$}\end{center}
\newtcbox{\congthuc}{on line,colframe=gold1,colback=gold5!45,boxrule=0.5pt,arc=2pt,
	left=8pt,right=8pt,top=4pt,bottom=4pt}
%--- Công tắc lời giải bài tự luyện (ví dụ mẫu luôn dùng \loigiai)
%     \loigiaiTL[{đáp số}]{lời giải đầy đủ}
%     \TLfull -> in lời giải đầy đủ;  \TLdapso -> chỉ in một dòng "Đáp số";  \TLan -> ẩn hoàn toàn
%     (\CheDo là dethi thì mọi lời giải và đáp số đều ẩn)
\makeatletter
\newif\ifTLhien   \TLhientrue
\newif\ifTLdapso  \TLdapsofalse
\newcommand{\TLfull}{\TLhientrue\TLdapsofalse}
\newcommand{\TLdapso}{\TLhientrue\TLdapsotrue}
\newcommand{\TLan}{\TLhienfalse}
\newcommand{\loigiaiTL}[2][]{%
	\ifTLhien
		\ifTLdapso
			\def\@tlarg{#1}%
			\ifx\@tlarg\@empty\else
				\ifdefstring{\kindOPTN}{dethi}{}{\par\noindent{\bfseries\color{gold1}Đáp số:} #1\par}%
			\fi
		\else
			\loigiai{#2}%
		\fi
	\fi}
\makeatother
\TLdapso% mặc định: chỉ in đáp số bài tự luyện
%====================== BẢNG ĐÚNG/SAI (\choiceTFt) ======================
%  Hàng tiêu đề nền nâu đồng chữ trắng (Phát biểu | Đ | S), các hàng tô xen kẽ kem / trắng,
%  đường kẻ vàng kim, nhãn a) b) c) d) màu nâu đồng. Bản có đáp án: cột được chọn có vòng tròn
%  nâu đồng với dấu tích trắng (câu đúng tích ở cột Đúng, câu sai tích ở cột Sai).
%  Tiêu đề ghi đủ "Đúng | Sai": \OPTN{viettat=0} sau \input (cột Đúng/Sai rộng ra, phát biểu hẹp lại).
%  Dùng bảng gốc của ex_test: \def\BangDSgoc{1} trước %============================================================
%  PREAMBLE EDUZY v2 - dùng chung cho đề thi, bài học, chuyên đề (gói ex_test)
%  Khác bản preamble-eduzy.tex ở đầu trang, chân trang: dải nền nâu đồng tràn ngang,
%  logo + Eduzy.com.vn, ô chữ vàng kem cạnh xiên, chân trang "Trang (số)";
%  logo Eduzy chìm giữa mọi trang (tắt: \logochimtat, bật lại: \logochimbat).
%  Cách dùng ở file chính (xem mau-chuan.tex, mau-chuyen-de.tex):
%      \def\CheDo{loigiai}          % loigiai | dethi | solcolor | color | book
%      \def\tenchuong{<dòng chữ trên đầu trang>}
%      \def\thuonghieu{<chữ ở chân trang lẻ>}   % không khai báo thì dùng mặc định
%      \def\tacgia{<chữ ở chân trang chẵn>}      % không khai báo thì dùng mặc định
%      \def\khauhieu{<chữ cạnh logo ở đầu trang>}  % tùy chọn, mặc định Eduzy.com.vn
%      \input{preamble-eduzy-v2}
%      \TLdapso                      % tài liệu: \TLfull | \TLdapso | \TLan
%      \begin{document} ... \end{document}
%  Logo Eduzy vẽ bằng TikZ ngay trong file (\eduzylogo), không cần file ảnh.
%  Biên dịch: pdfLaTeX (máy có gói vietnam) hoặc XeLaTeX.
%============================================================
\documentclass[12pt,a4paper,twoside]{article}%<-- twoside: bat dau trang/chan trang chan-le
\usepackage{iftex}
\usepackage{amsmath,amssymb,tikz,mathrsfs,tkz-tab}
\usepackage{graphicx}
\usepackage{pgfplots}% đồ thị hàm số
\pgfplotsset{compat=1.17}
%\usepackage{mathptmx}
\usepackage{mathpazo}
\ifXeTeX% XeLaTeX: font chữ có đủ dấu tiếng Việt, giữ nguyên font toán mathpazo
	\usepackage[no-math]{fontspec}
	\setmainfont{TeX Gyre Pagella}
	\newfontfamily\EZqag{TeX Gyre Adventor}[NFSSFamily=qag,Ligatures=TeX]
\fi
\usepackage[many]{tcolorbox}
\usetikzlibrary{math,through,calc,intersections,angles,quotes,shapes,shapes.geometric,arrows,patterns,matrix,chains,arrows.meta,decorations.shapes,decorations.fractals,decorations.markings,shadows,decorations.pathreplacing}
\usepackage[top=1.6cm, bottom=1.8cm, left=1.8cm, right=1.8cm]{geometry}% v2: top 1.6cm (bản cũ 1.5cm) để chữ không sát dải đầu trang
\providecommand{\CheDo}{loigiai}% loigiai | dethi | solcolor | color | book
\expandafter\usepackage\expandafter[\CheDo]{ex_test}
% \tenchuong: khai báo ở file chính (trước \input) hoặc \renewcommand sau \input
%============================================================
%  [GỘP TỪ eduzy-de.tex] EDUZY - MÀU NÂU ĐỒNG / VÀNG KIM
%  + ĐẦU TRANG, CHÂN TRANG KIỂU DẢI NGANG (v2), ĐỐI XỨNG TRANG CHẴN / TRANG LẺ.
%  (không cần file eduzy-de.tex riêng nữa)
%  Khối này phải đặt SAU \usepackage[...]{ex_test} và SAU \def\tenchuong{...}
%============================================================
\usepackage{lastpage}
\usepackage{needspace}
%--- DÀN TRANG ĐỀ THI: tách câu (THÊM 05/11/2026; chỉ thêm lệnh mới, KHÔNG đổi hành vi mặc định) ---
%  Hai công tắc, mặc định tắt (= preamble chạy y như cũ). Bật bằng cách đặt TRƯỚC \input{preamble-eduzy-v2}
%  trong file đề (tools/dan_trang/dan_trang.py tự ghi khi đặt \NeedspaceCau):
%    \def\TachCauDS{1}   câu Đúng/Sai được tách: câu dẫn (kèm hình) ở cuối trang, bảng Phát biểu sang trang sau
%                        (câu dẫn và hình không tách nhau, bảng không tách giữa hai trang)
%    \def\TachCauLG{1}   bản lời giải (\CheDo khác dethi): phần đề của câu cao hơn 1/3 trang được phép tách
%  Lệnh dùng trong file đề (giá trị do tools/dan_trang/dan_trang.py đo ở bản dethi):
%    \NeedspaceCau[<cao câu dẫn + hình>]{<cao cả khối câu>}   ngay trên \begin{ex} (câu đầu phần: trên tiêu đề phần)
%        công tắc tắt: đúng bằng \Needspace{<cao cả khối câu>} như cũ;
%        \TachCauDS=1 và có [..] (câu Đúng/Sai được chọn tách): chỉ giữ câu dẫn + hình;
%        \TachCauLG=1 ở bản lời giải: đòi tối đa 1/3 \textheight.
%    \NeedspaceBang{<cao bảng Phát biểu>}%   ngay trước \choiceTFt: \TachCauDS=1 thì giữ trọn bảng trên một trang;
%        công tắc tắt: không làm gì.
\providecommand{\TachCauDS}{0}
\providecommand{\TachCauLG}{0}
\newlength{\EZnsCau}
\newcommand{\NeedspaceCau}[2][]{%
	\setlength{\EZnsCau}{#2}%
	\ifnum\TachCauDS=1 \if\relax\detokenize{#1}\relax\else\setlength{\EZnsCau}{#1}\fi\fi
	\ifnum\TachCauLG=1 \ifdefstring{\kindOPTN}{dethi}{}{\ifdim\EZnsCau>0.3333\textheight\setlength{\EZnsCau}{0.3333\textheight}\fi}\fi
	\Needspace{\EZnsCau}}
\newcommand{\NeedspaceBang}[1]{\ifnum\TachCauDS=1 \Needspace{#1}\fi}

%--- Thang màu: đậm -> nhạt
\definecolor{gold1}{RGB}{126,96,0}    % nâu đồng đậm  #7E6000
\definecolor{gold2}{RGB}{192,144,3}   % vàng kim      #C09003
\definecolor{gold3}{RGB}{255,215,99}  % vàng nhạt     #FFD763
\definecolor{gold4}{RGB}{255,231,153} % vàng rất nhạt #FFE799
\definecolor{gold5}{RGB}{255,243,205} % vàng kem      #FFF3CD
%--- Ghi đè tông xanh/đỏ cũ
\definecolor{myblue}{RGB}{126,96,0}
\colorlet{blue}{gold1}
\colorlet{cyan}{gold2}
\colorlet{teal}{gold2}

%--- Tên chương mặc định (file đề nên tự khai báo lại)
\providecommand{\tenchuong}{}
%--- Chân trang
\providecommand{\thuonghieu}{TOÁN 12 -- Eduzy.com.vn}
\providecommand{\tacgia}{Tô Gia Bảo -- 0353703783}

%--- Chữ cạnh logo ở đầu trang
\providecommand{\khauhieu}{Eduzy.com.vn}

%====================== THÔNG SỐ ĐẦU TRANG, CHÂN TRANG (v2) ======================
%  Dải nền đậm tràn ngang trang; ô chữ vàng kem có cạnh xiên, kèm 2 vạch xiên chuyển màu.
%  Đầu trang lẻ : logo + \khauhieu (trái)   | ô chữ \tenchuong (phải)
%  Chân trang lẻ: ô chữ \thuonghieu (trái)  | Trang (số) (phải)
%  Trang chẵn soi gương trang lẻ; chân trang chẵn ghi \tacgia.
%--- Màu: chỉ dùng thang vàng kim Eduzy
\colorlet{ezNen}{gold1}% dải nền đậm
\colorlet{ezVach1}{gold3}% vạch xiên sát dải nền (đậm hơn)
\colorlet{ezVach2}{gold4}% vạch xiên sát ô chữ (nhạt hơn)
\colorlet{ezO}{gold5}% ô chữ vàng kem
%--- Kích thước
\newlength{\EZdai}\setlength{\EZdai}{1.1cm}% chiều cao dải nền
\newlength{\EZo}\setlength{\EZo}{0.68cm}% chiều cao ô chữ
\newlength{\EZxien}\setlength{\EZxien}{0.55cm}% độ lệch ngang của cạnh xiên
\newlength{\EZvach}\setlength{\EZvach}{0.75cm}% bề rộng mỗi vạch xiên
\newlength{\EZkhe}\setlength{\EZkhe}{0.12cm}% khe trắng giữa vạch và ô chữ
\newlength{\EZdem}\setlength{\EZdem}{0.4cm}% khoảng đệm từ chữ tới cạnh xiên
\newlength{\EZle}\setlength{\EZle}{1.8cm}% lề trái/phải của chữ (bằng lề thân bài)
\newlength{\EZrongmin}\setlength{\EZrongmin}{4.5cm}% ô chữ rộng tối thiểu
\newlength{\EZrong}% bề rộng chữ trong ô (tính tự động)
%--- Font
\newcommand{\EZfontO}{\fontsize{10.5pt}{12pt}\selectfont\bfseries}% chữ trong ô
\newcommand{\EZfontKH}{\fontfamily{qag}\fontsize{11pt}{12pt}\selectfont\bfseries}% \khauhieu cạnh logo
\newcommand{\EZfontTrang}{\fontsize{11.5pt}{12pt}\selectfont\bfseries}% chữ "Trang" và số trang
%--- Đo bề rộng chữ trong ô
\newcommand{\EZdorong}[1]{%
	\settowidth{\EZrong}{\EZfontO #1}%
	\ifdim\EZrong<\EZrongmin\setlength{\EZrong}{\EZrongmin}\fi}
%--- Logo + chữ trắng \khauhieu (dựng sẵn một lần khi vào \begin{document})
\newsavebox{\EZlogobox}
\AtBeginDocument{\sbox{\EZlogobox}{%
		\raisebox{-0.5\height}{\eduzylogo[1cm]}\hspace{0.15cm}%
		\raisebox{-0.5\height}{\color{white}\EZfontKH\khauhieu}}}
%--- Hai vạch xiên + khe trắng. #1: điểm gốc ở mép dưới, #2: hướng (1 sang phải, -1 sang trái),
%    #3: hướng nghiêng (1 kiểu "/", -1 kiểu "\")
\newcommand{\EZvachxien}[3]{%
	\fill[white] (#1) --++({(#3)*\EZxien},\EZo) --++({(#2)*\EZkhe},0) --++({-(#3)*\EZxien},-\EZo) --cycle;
	\fill[ezVach2] ([xshift=(#2)*\EZkhe]#1) --++({(#3)*\EZxien},\EZo) --++({(#2)*\EZvach},0) --++({-(#3)*\EZxien},-\EZo) --cycle;
	\fill[ezVach1] ([xshift=(#2)*(\EZkhe+\EZvach)]#1) --++({(#3)*\EZxien},\EZo) --++({(#2)*\EZvach},0) --++({-(#3)*\EZxien},-\EZo) --cycle;
	\draw[white,line width=0.6pt] ([xshift=(#2)*(\EZkhe+\EZvach)]#1) --++({(#3)*\EZxien},\EZo);
}
%--- Chữ trong ô (căn giữa theo chiều cao ô, không lệch vì dấu tiếng Việt)
\tikzset{ezchu/.style={inner sep=0pt,text=ezNen,text height=0.72em,text depth=0.22em}}

%====================== LOGO CHÌM GIỮA TRANG ======================
%  Vẽ cùng đầu trang nên nằm dưới chữ của thân bài, có trên mọi trang.
%  Tắt: \logochimtat (từ trang đó trở đi), bật lại: \logochimbat.
%  Đổi cỡ, độ đậm: khai báo trước \input, ví dụ \def\EZchimcao{9cm} \def\EZchimdo{0.1}
\providecommand{\EZchimcao}{10cm}% chiều cao logo chìm
\providecommand{\EZchimdo}{0.12}% độ đậm (0 = trong suốt, 1 = rõ hẳn)
\newif\ifEZchim \EZchimtrue
\newcommand{\logochimtat}{\EZchimfalse}
\newcommand{\logochimbat}{\EZchimtrue}
%  Dựng sẵn logo mờ một lần khi vào \begin{document}; transparency group làm mờ cả logo
%  một lượt nên không lộ vệt ở chỗ các dải màu chồng nhau.
\newsavebox{\EZchimtho}
\newsavebox{\EZchimbox}
\AtBeginDocument{%
	\sbox{\EZchimtho}{\eduzylogo[\EZchimcao]}%
	\sbox{\EZchimbox}{\begin{tikzpicture}
			\begin{scope}[transparency group,opacity=\EZchimdo]
				\node[inner sep=0pt] {\usebox{\EZchimtho}};
			\end{scope}
	\end{tikzpicture}}}
\newcommand{\EZvechim}{%
	\ifEZchim
		\node[inner sep=0pt] at (current page.center) {\usebox{\EZchimbox}};
	\fi}

%====================== ĐẦU TRANG - TRANG LẺ ======================
%  (logo bên trái, ô chữ \tenchuong bên phải)
\newcommand{\myfancyheadOdd}{%
	\EZdorong{\tenchuong}%
	\begin{tikzpicture}[remember picture,overlay]
		\coordinate (T) at (current page.north west);
		\coordinate (P) at (current page.north east);
		%--- logo chìm giữa trang
		\EZvechim
		%--- dải nền
		\fill[ezNen] (T) rectangle ([yshift=-\EZdai]P);
		%--- ô chữ: X là góc dưới bên trái, cạnh trái xiên kiểu "/"
		\coordinate (X) at ([xshift=-\EZle-\EZrong-\EZdem-\EZxien,yshift=-\EZdai]P);
		\fill[ezO] (X) --++(\EZxien,\EZo) -- ([yshift=-\EZdai+\EZo]P) -- ([yshift=-\EZdai]P) --cycle;
		\EZvachxien{X}{-1}{1}
		\node[ezchu,anchor=east] at ([xshift=-\EZle,yshift=-\EZdai+0.5*\EZo]P) {\EZfontO\tenchuong};
		%--- logo + tên
		\node[inner sep=0pt,anchor=west] at ([xshift=\EZle,yshift=-0.5*\EZdai]T) {\usebox{\EZlogobox}};
	\end{tikzpicture}%
}

%====================== ĐẦU TRANG - TRANG CHẴN =====================
%  (soi gương: ô chữ \tenchuong bên trái, logo bên phải)
\newcommand{\myfancyheadEven}{%
	\EZdorong{\tenchuong}%
	\begin{tikzpicture}[remember picture,overlay]
		\coordinate (T) at (current page.north west);
		\coordinate (P) at (current page.north east);
		%--- logo chìm giữa trang
		\EZvechim
		%--- dải nền
		\fill[ezNen] (T) rectangle ([yshift=-\EZdai]P);
		%--- ô chữ: X là góc dưới bên phải, cạnh phải xiên kiểu "\"
		\coordinate (X) at ([xshift=\EZle+\EZrong+\EZdem+\EZxien,yshift=-\EZdai]T);
		\fill[ezO] (X) --++(-\EZxien,\EZo) -- ([yshift=-\EZdai+\EZo]T) -- ([yshift=-\EZdai]T) --cycle;
		\EZvachxien{X}{1}{-1}
		\node[ezchu,anchor=west] at ([xshift=\EZle,yshift=-\EZdai+0.5*\EZo]T) {\EZfontO\tenchuong};
		%--- logo + tên
		\node[inner sep=0pt,anchor=east] at ([xshift=-\EZle,yshift=-0.5*\EZdai]P) {\usebox{\EZlogobox}};
	\end{tikzpicture}%
}

%====================== CHÂN TRANG - TRANG LẺ ======================
%  (ô chữ \thuonghieu bên trái, "Trang (số)" bên phải)
\newcommand{\myfancyfootOdd}{%
	\EZdorong{\thuonghieu}%
	\begin{tikzpicture}[remember picture,overlay]
		\coordinate (T) at (current page.south west);
		\coordinate (P) at (current page.south east);
		%--- dải nền
		\fill[ezNen] (T) rectangle ([yshift=\EZdai]P);
		%--- ô chữ sát mép trên dải: X là góc dưới bên phải, cạnh phải xiên kiểu "/"
		\coordinate (X) at ([xshift=\EZle+\EZrong+\EZdem,yshift=\EZdai-\EZo]T);
		\fill[ezO] ([yshift=\EZdai-\EZo]T) -- (X) --++(\EZxien,\EZo) -- ([yshift=\EZdai]T) --cycle;
		\EZvachxien{X}{1}{1}
		\node[ezchu,anchor=west] at ([xshift=\EZle,yshift=\EZdai-0.5*\EZo]T) {\EZfontO\thuonghieu};
		%--- số trang
		\node[circle,draw=white,line width=0.9pt,text=white,inner sep=0pt,minimum size=0.62cm,anchor=east]
		(so) at ([xshift=-\EZle,yshift=0.5*\EZdai]P) {\EZfontTrang\thepage};
		\node[inner sep=0pt,text=white,anchor=base east] (tr) at ([xshift=-0.15cm]so.base west) {\EZfontTrang Trang};
	\end{tikzpicture}%
}

%====================== CHÂN TRANG - TRANG CHẴN ====================
%  (soi gương: "Trang (số)" bên trái, ô chữ \tacgia bên phải)
\newcommand{\myfancyfootEven}{%
	\EZdorong{\tacgia}%
	\begin{tikzpicture}[remember picture,overlay]
		\coordinate (T) at (current page.south west);
		\coordinate (P) at (current page.south east);
		%--- dải nền
		\fill[ezNen] (T) rectangle ([yshift=\EZdai]P);
		%--- ô chữ sát mép trên dải: X là góc dưới bên trái, cạnh trái xiên kiểu "\"
		\coordinate (X) at ([xshift=-\EZle-\EZrong-\EZdem,yshift=\EZdai-\EZo]P);
		\fill[ezO] ([yshift=\EZdai-\EZo]P) -- (X) --++(-\EZxien,\EZo) -- ([yshift=\EZdai]P) --cycle;
		\EZvachxien{X}{-1}{-1}
		\node[ezchu,anchor=east] at ([xshift=-\EZle,yshift=\EZdai-0.5*\EZo]P) {\EZfontO\tacgia};
		%--- số trang
		%    (vòng tròn ẩn so0 nằm giữa dải để lấy đường chân chữ, giống trang lẻ)
		\node[circle,inner sep=0pt,minimum size=0.62cm] (so0) at ([yshift=0.5*\EZdai]T) {\EZfontTrang\phantom{\thepage}};
		\node[inner sep=0pt,text=white,anchor=base west] (tr) at ($(so0.base -| T)+(\EZle,0)$) {\EZfontTrang Trang};
		\node[circle,draw=white,line width=0.9pt,text=white,inner sep=0pt,minimum size=0.62cm,anchor=base west]
		(so) at ([xshift=0.15cm]tr.base east) {\EZfontTrang\thepage};
	\end{tikzpicture}%
}

%--- Giữ tương thích với các file cũ đang gọi \myfancyhead, \myfancyfoot
\newcommand{\myfancyhead}{\myfancyheadOdd}
\newcommand{\myfancyfoot}{\myfancyfootOdd}

%======================= GẮN VÀO TRANG =============================
\pagestyle{fancy}
\fancyhf{}
\fancyhead[LO]{\myfancyheadOdd}
\fancyhead[RE]{\myfancyheadEven}
\fancyfoot[LO]{\myfancyfootOdd}
\fancyfoot[RE]{\myfancyfootEven}
\renewcommand{\headrulewidth}{0pt}
\renewcommand{\footrulewidth}{0pt}
\fancypagestyle{plain}{%
	\fancyhf{}
	\fancyhead[LO]{\myfancyheadOdd}
	\fancyhead[RE]{\myfancyheadEven}
	\fancyfoot[LO]{\myfancyfootOdd}
	\fancyfoot[RE]{\myfancyfootEven}
	\renewcommand{\headrulewidth}{0pt}
	\renewcommand{\footrulewidth}{0pt}
}
%=================== HẾT PHẦN GỘP TỪ eduzy-de.tex ===================
\newcommand{\indam}[1]{
	{\fontsize{12pt}{9pt}\sffamily \selectfont\bfseries\color{gold1} #1}
}
\newcommand{\ind}[1]{
	{\fontsize{14pt}{9pt}\fontfamily{qag} \selectfont\bfseries\color{gold1} #1}
}
\renewcommand{\FalseEX}{\stepcounter{dapan}\textbf{\selectfont\color{gold1}\Alph{dapan}.}\ignorespaces}
\renewtheorem{ex}{\color{gold1}Câu}
\renewcommand{\baselinestretch}{1.1}
\newcommand{\hoac}[1]{
	\left[\begin{aligned}#1\end{aligned}\right.}
\newcommand{\heva}[1]{
	\left\{\begin{aligned}#1\end{aligned}\right.}
\usepackage{esvect}
\def\vec{\vv}
\def\vt{\vv}
\def\overrightarrow{\vv}
%====================== LOGO EDUZY VẼ BẰNG TIKZ ======================
%  Đường viền dò từ logo gốc; dải màu vàng kim đo trực tiếp trên ảnh gốc
%  (nghiêng 21 độ, tô bằng 80 dải hẹp nên hiển thị giống nhau trên mọi trình xem PDF).
%  Dùng: \eduzylogo          -> cao 2.4cm (mặc định)
%        \eduzylogo[1.5cm]   -> cao 1.5cm
\definecolor{ezText}{RGB}{245,200,46}
\newcommand{\eduzylogo}[1][2.4cm]{%
	\resizebox{!}{#1}{%
		\begin{tikzpicture}[x=0.01cm,y=0.01cm]
			\path[use as bounding box] (0,0) rectangle (224,-224);
			% mũ tốt nghiệp + vương miện
			\begin{scope}
				\clip
				(119.4,-176.1)..controls(118.9,-176.1)and(116.8,-176)..(114.6,-176)
				..controls(101.4,-176)and(85.8,-175.9)..(83.7,-175.7)
				..controls(81.8,-175.6)and(81,-175.6)..(80.2,-175.7)
				..controls(78.8,-175.9)and(77.7,-175.9)..(76.2,-175.7)
				..controls(74.9,-175.5)and(70.8,-175.5)..(64.3,-175.7)
				..controls(59.1,-175.9)and(58.3,-175.7)..(58.6,-173.7)
				..controls(58.6,-172.8)and(59.6,-169.6)..(60.3,-168)
				..controls(60.8,-166.9)and(61.4,-165.8)..(62.3,-165)
				..controls(63.1,-164.2)and(64.3,-163.4)..(65,-163.2)
				..controls(65.5,-163.1)and(72.3,-163)..(76.6,-163.1)
				..controls(77.9,-163.1)and(80.1,-163.1)..(81.4,-163)
				..controls(82.8,-162.9)and(86.3,-162.9)..(89.2,-162.9)
				..controls(93,-162.9)and(94.9,-162.8)..(96,-162.7)
				..controls(97.2,-162.5)and(101,-162.5)..(114,-162.4)
				..controls(123,-162.3)and(134.9,-162.2)..(140.2,-162.1)
				..controls(150.1,-162)and(152.2,-162.1)..(153.6,-162.4)
				..controls(154.4,-162.6)and(157.1,-163.8)..(158.1,-164.4)
				..controls(158.8,-164.8)and(159.7,-165.7)..(159.9,-166.3)
				..controls(160,-166.5)and(160.2,-167.1)..(160.2,-167.6)
				..controls(160.3,-168)and(160.6,-169.4)..(160.9,-170.6)
				..controls(161.2,-171.7)and(161.5,-173.1)..(161.6,-173.6)
				..controls(161.8,-174.4)and(161.8,-174.5)..(161.5,-174.9)
				..controls(161.2,-175.6)and(160.8,-175.9)..(159.9,-176)
				..controls(159.2,-176.1)and(154.1,-176.2)..(151.5,-176.2)
				..controls(148.1,-176.1)and(122.6,-176.1)..(121.6,-176.2)
				..controls(120.8,-176.2)and(119.8,-176.2)..(119.4,-176.1)--cycle
				(130.4,-160)..controls(129,-160)and(119.8,-160)..(110.1,-159.9)
				..controls(100.3,-159.8)and(88.5,-159.7)..(83.9,-159.6)
				..controls(77.8,-159.5)and(75.3,-159.5)..(74.5,-159.6)
				..controls(73.9,-159.7)and(72.5,-159.8)..(71.3,-159.8)
				..controls(69.6,-159.8)and(69,-159.8)..(68.5,-159.6)
				..controls(67.6,-159.3)and(66.8,-158.6)..(66.2,-157.6)
				..controls(66,-157.3)and(65.6,-156.5)..(65.3,-156.1)
				..controls(64.6,-155.1)and(64.5,-154.9)..(63.7,-152.9)
				..controls(63.2,-151.6)and(62.9,-150.9)..(61.8,-149.6)
				..controls(61.5,-149.2)and(61.1,-148.5)..(60.9,-147.9)
				..controls(60.7,-147.4)and(60.3,-146.6)..(60.1,-146.1)
				..controls(59.9,-145.5)and(59.5,-144.9)..(59,-144.3)
				..controls(58.5,-143.7)and(58.2,-143.2)..(57.9,-142.5)
				..controls(57.7,-141.9)and(57.3,-141.1)..(57,-140.7)
				..controls(56.6,-140.3)and(56.1,-139.3)..(55.8,-138.5)
				..controls(55.5,-137.8)and(55.1,-136.8)..(54.9,-136.5)
				..controls(54.7,-136.1)and(54.3,-135.2)..(53.9,-134.5)
				..controls(53.6,-133.8)and(53,-132.7)..(52.6,-132)
				..controls(51.6,-130.5)and(51.4,-130.2)..(51.1,-129)
				..controls(50.6,-127.5)and(50.6,-127.4)..(49.7,-126.3)
				..controls(49.2,-125.7)and(48.6,-124.8)..(48.3,-124.2)
				..controls(48,-123.6)and(47.5,-122.7)..(47.3,-122.2)
				..controls(47,-121.8)and(46.5,-120.7)..(46.1,-119.8)
				..controls(45.8,-118.9)and(45.2,-117.8)..(44.9,-117.2)
				..controls(44.2,-116.1)and(43.3,-114.4)..(42.3,-112.2)
				..controls(41.8,-111.3)and(41.3,-110.4)..(41.1,-110.1)
				..controls(40.5,-109.4)and(40.3,-109)..(40,-107.8)
				..controls(39.8,-107)and(39.6,-106.5)..(39.2,-106)
				..controls(38.9,-105.6)and(38.4,-104.7)..(38,-104)
				..controls(37.7,-103.4)and(37.2,-102.5)..(37,-102.1)
				..controls(36.7,-101.7)and(36.3,-100.8)..(36,-100.1)
				..controls(35.6,-98.9)and(35.3,-98.4)..(34.6,-97.6)
				..controls(34.2,-97.1)and(34.2,-96.3)..(34.6,-96)
				..controls(35,-95.7)and(35.5,-95.9)..(36.2,-96.7)
				..controls(36.5,-97)and(37.2,-97.5)..(37.6,-97.7)..controls(38,-98)and(38.6,-98.4)..(39.1,-98.9)
				..controls(39.6,-99.4)and(40.3,-100)..(40.6,-100.3)
				..controls(41.3,-100.8)and(42.4,-101.8)..(43.7,-102.9)
				..controls(44,-103.2)and(44.6,-103.7)..(45,-103.9)
				..controls(45.4,-104.2)and(45.9,-104.6)..(46.1,-104.8)
				..controls(47.1,-105.9)and(48.8,-107.1)..(50.5,-107.9)
				..controls(51.8,-108.6)and(52.3,-109)..(53.4,-110)
				..controls(54.8,-111.2)and(55.6,-111.6)..(57,-112.1)
				..controls(57.5,-112.3)and(58.5,-112.8)..(59.4,-113.3)
				..controls(60.2,-113.7)and(61.3,-114.3)..(61.9,-114.5)
				..controls(62.4,-114.7)and(63.3,-115.1)..(63.8,-115.3)
				..controls(64.4,-115.6)and(65.4,-115.9)..(66.2,-116.1)
				..controls(67,-116.3)and(68.4,-116.6)..(69.2,-116.9)
				..controls(72.3,-117.9)and(72.8,-118)..(78,-118.1)
				..controls(84,-118.3)and(86.1,-118.1)..(88.1,-117.1)
				..controls(89.4,-116.5)and(93,-115)..(93.7,-114.9)
				..controls(94.8,-114.6)and(95.4,-114.2)..(95.9,-113.4)
				..controls(96.2,-113)and(96.7,-112.5)..(97.1,-112.2)
				..controls(98,-111.6)and(101.3,-108.4)..(102.1,-107.2)
				..controls(102.4,-106.8)and(102.9,-106.2)..(103.2,-105.8)
				..controls(103.5,-105.5)and(104,-104.8)..(104.2,-104.2)
				..controls(104.5,-103.6)and(105,-102.8)..(105.3,-102.4)
				..controls(105.6,-102)and(106.1,-101.2)..(106.5,-100.6)
				..controls(107.1,-99.5)and(108.6,-97.5)..(109.5,-96.5)
				..controls(109.9,-96.1)and(110.1,-96)..(110.4,-96)
				..controls(111.1,-96)and(111.6,-96.5)..(111.8,-97.6)
				..controls(112,-98.6)and(112.3,-99.3)..(112.8,-99.9)
				..controls(113.7,-100.9)and(115.2,-103.1)..(115.7,-104.2)
				..controls(116.1,-105)and(116.5,-105.5)..(116.9,-106)
				..controls(117.3,-106.4)and(117.8,-106.9)..(118,-107.3)
				..controls(118.2,-107.6)and(118.6,-108)..(119,-108.3)
				..controls(119.6,-108.8)and(121.1,-110.3)..(121.8,-111.2)
				..controls(122.1,-111.7)and(122.5,-112)..(123.4,-112.4)
				..controls(124.8,-113.2)and(125.8,-113.8)..(127,-114.6)
				..controls(127.5,-114.9)and(128.1,-115.3)..(128.2,-115.4)
				..controls(128.4,-115.5)and(129.3,-115.8)..(130.2,-116)
				..controls(131.1,-116.2)and(132.5,-116.6)..(133.3,-116.9)
				..controls(135.7,-117.9)and(136.1,-117.9)..(142.5,-117.8)
				..controls(145.8,-117.8)and(148.4,-117.7)..(148.8,-117.6)
				..controls(149.2,-117.6)and(150.2,-117.2)..(151,-116.8)
				..controls(152,-116.4)and(152.8,-116)..(153.6,-115.9)
				..controls(154.5,-115.7)and(155.2,-115.4)..(156.3,-114.9)
				..controls(157.1,-114.5)and(158.2,-114)..(158.7,-113.8)
				..controls(159.1,-113.7)and(159.8,-113.3)..(160.2,-113)
				..controls(160.6,-112.7)and(161.3,-112.4)..(161.7,-112.2)
				..controls(162.7,-111.9)and(163.2,-111.7)..(164.1,-111.1)
				..controls(164.5,-110.8)and(165.3,-110.4)..(165.9,-110)
				..controls(166.5,-109.7)and(167.3,-109.2)..(167.6,-108.9)
				..controls(167.9,-108.6)and(168.6,-108.2)..(169,-108)
				..controls(169.4,-107.8)and(170.2,-107.3)..(170.7,-106.9)
				..controls(171.2,-106.5)and(172.1,-105.8)..(172.7,-105.4)
				..controls(173.2,-105)and(174.1,-104.4)..(174.5,-104)
				..controls(174.9,-103.7)and(175.7,-103.1)..(176.2,-102.6)
				..controls(176.7,-102.2)and(177.4,-101.5)..(177.9,-101.2)
				..controls(178.3,-100.8)and(178.8,-100.4)..(179,-100.2)
				..controls(179.2,-100)and(179.6,-99.6)..(180,-99.3)
				..controls(180.4,-99.1)and(180.9,-98.6)..(181.2,-98.4)
				..controls(181.5,-98.1)and(182.1,-97.7)..(182.5,-97.5)
				..controls(183.2,-97.2)and(183.5,-96.9)..(184.3,-95.9)
				..controls(184.6,-95.6)and(184.6,-95.6)..(184.8,-95.7)
				..controls(185,-95.8)and(185.1,-96)..(185.1,-96.2)
				..controls(185.1,-96.5)and(185,-96.7)..(184.6,-97.2)
				..controls(184.1,-97.8)and(184,-97.9)..(183.9,-98.8)
				..controls(183.7,-100.3)and(183.3,-101.3)..(182.9,-101.8)
				..controls(182.6,-102.1)and(182.3,-102.6)..(182.2,-102.9)
				..controls(182,-103.3)and(181.6,-104)..(181.2,-104.6)
				..controls(180.9,-105.2)and(180.5,-106.1)..(180.4,-106.6)
				..controls(179.9,-108)and(179.7,-108.5)..(179,-109.4)
				..controls(178.6,-109.8)and(178.2,-110.5)..(177.9,-111)
				..controls(177.7,-111.5)and(177.3,-112.3)..(177,-112.7)
				..controls(176.6,-113.2)and(176.4,-113.7)..(176.3,-114.2)
				..controls(176.2,-114.7)and(175.9,-115.2)..(175.4,-116)
				..controls(175,-116.6)and(174.4,-117.7)..(174,-118.5)
				..controls(173.7,-119.2)and(173.3,-120.1)..(173,-120.5)
				..controls(172.8,-120.9)and(172.4,-121.8)..(172.1,-122.5)
				..controls(171.8,-123.2)and(171.3,-124.1)..(171,-124.6)
				..controls(170.8,-125.1)and(170.3,-126)..(169.9,-126.6)
				..controls(169.6,-127.3)and(169.2,-128.1)..(169,-128.3)
				..controls(168.8,-128.6)and(168.5,-129.3)..(168.3,-129.8)
				..controls(168.1,-130.3)and(167.8,-131.1)..(167.5,-131.6)
				..controls(166.9,-132.5)and(164.2,-137.9)..(164.2,-138.1)
				..controls(164.2,-138.1)and(163.9,-138.8)..(163.4,-139.7)
				..controls(162.9,-140.5)and(162.4,-141.5)..(162.3,-141.9)
				..controls(161.9,-143)and(161.6,-143.4)..(160.9,-144.1)
				..controls(160.4,-144.6)and(160.3,-144.8)..(160.1,-145.4)
				..controls(160,-145.8)and(159.6,-146.8)..(159.2,-147.8)
				..controls(158.7,-148.8)and(158.1,-150)..(157.9,-150.6)
				..controls(157.6,-151.1)and(157.2,-151.9)..(156.9,-152.4)
				..controls(156.5,-152.8)and(156.1,-153.7)..(155.8,-154.4)
				..controls(155.6,-155)and(155.2,-155.9)..(154.9,-156.3)
				..controls(154.7,-156.7)and(154.4,-157.4)..(154.2,-157.9)
				..controls(154.1,-158.4)and(153.8,-159)..(153.5,-159.3)
				..controls(152.8,-160.1)and(153,-160.1)..(142.4,-160.1)
				..controls(137.2,-160.1)and(131.8,-160.1)..(130.4,-160)--cycle
				(141.5,-103.7)..controls(140.8,-103)and(139.8,-102.1)..(139.2,-101.6)
				..controls(138.5,-101.2)and(137.9,-100.6)..(137.8,-100.4)
				..controls(137,-99.5)and(134,-96.5)..(133.2,-96)
				..controls(132.7,-95.7)and(131.8,-94.9)..(131.1,-94.3)
				..controls(129.4,-92.7)and(127.9,-91.6)..(126.6,-91.1)
				..controls(125.2,-90.5)and(124.1,-90)..(122.9,-89.1)
				..controls(122.1,-88.6)and(121.3,-88.2)..(119.9,-87.7)
				..controls(118.8,-87.3)and(117.5,-86.8)..(116.9,-86.5)
				..controls(114.3,-85.5)and(113.7,-85.4)..(109.7,-85.5)
				..controls(105.3,-85.7)and(103.9,-85.9)..(101.6,-87)
				..controls(100.7,-87.4)and(99.4,-87.8)..(97.9,-88)
				..controls(97,-88.2)and(96.6,-88.3)..(95.8,-88.9)
				..controls(95.4,-89.2)and(94.7,-89.6)..(94.2,-89.8)
				..controls(92.9,-90.3)and(92.3,-90.7)..(90.9,-92.1)
				..controls(90.2,-92.7)and(89.4,-93.3)..(89,-93.5)..controls(87,-94.5)and(85.2,-96)..(83.6,-97.9)
				..controls(83.3,-98.2)and(82.5,-98.9)..(81.8,-99.4)
				..controls(80.7,-100.1)and(80.5,-100.3)..(80,-101.1)
				..controls(78.7,-103.1)and(77.2,-103.9)..(76.4,-103)
				..controls(75.7,-102.2)and(75.7,-103.2)..(75.7,-83.6)
				..controls(75.6,-73.6)and(75.6,-64.8)..(75.5,-64)
				..controls(75.5,-63.2)and(75.5,-61.4)..(75.6,-59.8)--(75.7,-57.1)--(76,-56.8)
				..controls(76.8,-56.2)and(78,-56.3)..(79.8,-57.2)
				..controls(80.4,-57.5)and(81.3,-57.8)..(81.8,-57.9)
				..controls(82.3,-58.1)and(83,-58.4)..(83.6,-58.8)..controls(84.3,-59.2)and(85,-59.5)..(86,-59.8)
				..controls(86.8,-60.1)and(88,-60.6)..(88.7,-60.9)
				..controls(89.3,-61.2)and(90.6,-61.7)..(91.4,-62)
				..controls(92.2,-62.3)and(93.7,-62.9)..(94.6,-63.3)
				..controls(95.6,-63.7)and(96.8,-64.1)..(97.4,-64.3)
				..controls(98,-64.5)and(99,-64.9)..(99.7,-65.2)
				..controls(100.4,-65.6)and(101.5,-66)..(102.3,-66.2)
				..controls(104.2,-66.8)and(105,-67.1)..(106.2,-67.8)
				..controls(107.9,-68.7)and(110.2,-69.4)..(111.3,-69.2)
				..controls(112.5,-69)and(115,-68.2)..(115.6,-67.8)
				..controls(117.3,-66.7)and(117.6,-66.6)..(119.5,-66.2)
				..controls(121.9,-65.8)and(122.1,-65.7)..(123,-65.1)
				..controls(124.1,-64.4)and(124.4,-64.3)..(126,-63.9)
				..controls(127,-63.7)and(127.7,-63.4)..(128.4,-63.1)
				..controls(129,-62.8)and(129.7,-62.5)..(130.6,-62.4)
				..controls(132.1,-62.1)and(134,-61.5)..(135.2,-61)
				..controls(135.7,-60.8)and(136.7,-60.5)..(137.3,-60.3)
				..controls(138.4,-60)and(139.3,-59.7)..(141.1,-59)
				..controls(141.6,-58.7)and(142,-58.6)..(142.5,-58.6)
				..controls(143.2,-58.6)and(143.3,-58.7)..(143.6,-59)
				..controls(144,-59.5)and(144.1,-59.9)..(144,-60.8)
				..controls(144,-61.2)and(144,-71.1)..(144,-82.8)--(143.9,-104.1)--(143.7,-104.5)
				..controls(143.5,-104.7)and(143.3,-104.8)..(143,-104.9)
				..controls(142.6,-104.9)and(142.5,-104.8)..(141.5,-103.7)--cycle
				(27.8,-91.4)..controls(27.1,-90.9)and(26.3,-90)..(26.1,-89.5)
				..controls(26,-89.2)and(25.8,-88.6)..(25.7,-87.9)
				..controls(25.4,-86.2)and(25.6,-84.8)..(26.5,-83.5)
				..controls(27.4,-82.3)and(28.3,-82)..(29.4,-82.3)
				..controls(31.1,-82.8)and(32.1,-84.9)..(32,-87.4)
				..controls(31.9,-89.3)and(31.2,-90.6)..(29.9,-91.4)
				..controls(29.1,-91.9)and(28.6,-91.9)..(27.8,-91.4)--cycle
				(190.1,-91.3)..controls(189.7,-91.2)and(189.3,-90.9)..(189.1,-90.7)
				..controls(188.3,-89.8)and(187.6,-87.7)..(187.6,-86.2)
				..controls(187.6,-85.5)and(187.7,-85.2)..(188.2,-84.4)
				..controls(189.1,-82.6)and(190.1,-81.8)..(191.3,-82.1)
				..controls(192,-82.3)and(193.5,-83.7)..(193.8,-84.4)
				..controls(194.1,-85.1)and(194.2,-87.1)..(194,-88.3)
				..controls(193.8,-89.8)and(193.3,-90.5)..(192.3,-91.3)
				..controls(191.7,-91.7)and(190.9,-91.7)..(190.1,-91.3)--cycle
				(173.7,-83.4)..controls(172.2,-82.5)and(172,-82.5)..(170.6,-82.4)
				..controls(169.7,-82.4)and(168.8,-82.4)..(168.3,-82.5)
				..controls(167.2,-82.7)and(166.9,-82.7)..(166.6,-82.2)
				..controls(166.3,-81.8)and(166.3,-81)..(166.8,-80)
				..controls(167.2,-78.9)and(167.4,-77.4)..(167.7,-74.1)
				..controls(167.8,-72.8)and(167.9,-71.8)..(167.9,-71.7)
				..controls(168,-71.7)and(168,-70.2)..(168,-68.4)..controls(168,-65.4)and(168,-64.9)..(167.7,-63)
				..controls(167.3,-60.5)and(167.1,-59.9)..(166.4,-58.4)
				..controls(166.1,-57.7)and(165.8,-56.9)..(165.7,-56.6)
				..controls(165.7,-56.2)and(165.4,-55.6)..(165.2,-55.2)
				..controls(164.6,-53.8)and(161.9,-50.7)..(161,-50.3)
				..controls(159.8,-49.7)and(157.9,-50)..(155.2,-51.2)
				..controls(154.7,-51.4)and(153.4,-51.8)..(152.2,-52.1)
				..controls(150.8,-52.4)and(149.4,-52.9)..(148.6,-53.2)
				..controls(147.8,-53.6)and(146.7,-54)..(146,-54.2)
				..controls(145.4,-54.4)and(144.5,-54.7)..(144,-54.9)
				..controls(143.6,-55.2)and(142.7,-55.5)..(141.8,-55.7)
				..controls(141,-55.9)and(139.6,-56.3)..(138.7,-56.6)
				..controls(137.8,-57)and(136.3,-57.4)..(135.3,-57.6)
				..controls(133.8,-57.9)and(133.2,-58.1)..(132.1,-58.7)
				..controls(131.1,-59.2)and(130.3,-59.4)..(129.4,-59.6)
				..controls(128.7,-59.8)and(127.8,-60)..(127.4,-60.1)
				..controls(127,-60.3)and(126,-60.5)..(125.1,-60.7)
				..controls(123.2,-61.1)and(121.7,-61.5)..(120.2,-62.1)
				..controls(119.7,-62.3)and(118.1,-62.8)..(116.7,-63.2)
				..controls(114.6,-63.9)and(114,-64.1)..(113,-64.7)
				..controls(111.9,-65.4)and(111.8,-65.4)..(111,-65.5)
				..controls(109.9,-65.5)and(109.3,-65.4)..(107.7,-64.7)
				..controls(107,-64.4)and(106,-64)..(105.3,-63.9)
				..controls(103.5,-63.5)and(102.2,-63)..(101.2,-62.5)
				..controls(100.1,-61.8)and(99.3,-61.5)..(97.8,-61.1)
				..controls(97.2,-61)and(96.2,-60.7)..(95.7,-60.4)
				..controls(95.1,-60.1)and(94.2,-59.8)..(93.5,-59.6)
				..controls(92.9,-59.4)and(91.7,-59)..(91,-58.7)..controls(90.2,-58.4)and(88.9,-58)..(88,-57.7)
				..controls(87.2,-57.5)and(86.1,-57.1)..(85.5,-56.8)
				..controls(85,-56.6)and(84.1,-56.2)..(83.5,-56)..controls(82.9,-55.7)and(81.8,-55.2)..(81,-54.9)
				..controls(80.1,-54.5)and(79.1,-54.2)..(78.7,-54.1)
				..controls(77.9,-53.8)and(77.3,-53.6)..(76.2,-52.9)
				..controls(75.7,-52.6)and(75.2,-52.4)..(74.1,-52.2)
				..controls(72.5,-51.9)and(72.4,-51.8)..(71.3,-51.1)
				..controls(69.9,-50.1)and(69.4,-49.9)..(67.9,-49.8)
				..controls(67,-49.7)and(66.6,-49.6)..(66.3,-49.4)
				..controls(66.1,-49.3)and(65.7,-49)..(65.4,-48.8)
				..controls(65.1,-48.7)and(64.2,-48.4)..(63.6,-48.2)
				..controls(62.9,-48)and(61.7,-47.6)..(60.9,-47.2)
				..controls(60.1,-46.8)and(58.9,-46.4)..(58.2,-46.2)
				..controls(57.6,-46)and(56.8,-45.7)..(56.5,-45.5)
				..controls(54.9,-44.6)and(54.6,-44.5)..(53.2,-44.2)
				..controls(51.7,-43.8)and(51,-43.6)..(49.9,-42.9)
				..controls(49.5,-42.6)and(48.9,-42.4)..(47.5,-42)
				..controls(45.7,-41.6)and(45.6,-41.5)..(44.3,-40.6)
				..controls(43.5,-40.1)and(42.8,-39.6)..(42.7,-39.4)
				..controls(42.4,-38.7)and(43,-38.1)..(44.5,-37.6)..controls(44.9,-37.5)and(46,-37)..(46.9,-36.6)
				..controls(47.9,-36.1)and(48.8,-35.8)..(49.6,-35.6)
				..controls(50.5,-35.4)and(51.3,-35.2)..(52.4,-34.6)
				..controls(53.3,-34.1)and(54.5,-33.7)..(54.9,-33.5)
				..controls(55.4,-33.4)and(56.4,-33)..(57.2,-32.6)
				..controls(58.6,-32)and(60.1,-31.5)..(61.6,-31.2)
				..controls(62.3,-31.1)and(62.7,-31)..(63.2,-30.7)
				..controls(64.1,-30)and(64.3,-29.9)..(65.3,-29.7)
				..controls(65.7,-29.6)and(66.5,-29.2)..(67.1,-28.9)
				..controls(67.6,-28.6)and(68.8,-28.2)..(69.6,-27.9)
				..controls(70.5,-27.7)and(71.6,-27.3)..(72.1,-27)
				..controls(72.6,-26.7)and(73.8,-26.2)..(74.8,-25.9)
				..controls(75.7,-25.5)and(76.7,-25.1)..(77.1,-25)
				..controls(77.4,-24.8)and(78.5,-24.4)..(79.6,-24.1)
				..controls(80.6,-23.8)and(82,-23.2)..(82.8,-22.9)
				..controls(83.5,-22.5)and(84.8,-22)..(85.8,-21.8)
				..controls(86.7,-21.5)and(87.8,-21.1)..(88.1,-20.9)
				..controls(88.4,-20.7)and(89.3,-20.3)..(90.1,-20.1)
				..controls(90.8,-19.8)and(91.9,-19.4)..(92.4,-19.1)
				..controls(93.5,-18.5)and(94.4,-18.2)..(95.4,-18.1)
				..controls(95.7,-18)and(96.4,-17.8)..(96.9,-17.5)
				..controls(97.5,-17.3)and(98.4,-16.9)..(99,-16.6)
				..controls(99.7,-16.3)and(100.9,-15.9)..(101.9,-15.6)
				..controls(102.9,-15.3)and(104.4,-14.7)..(105.3,-14.3)
				..controls(108.1,-13.2)and(108.5,-13.1)..(110.1,-13.1)
				..controls(111.4,-13.1)and(111.5,-13.1)..(112.4,-13.5)
				..controls(113.2,-14)and(115,-14.8)..(116.6,-15.4)
				..controls(116.9,-15.5)and(117.9,-15.8)..(118.8,-16.1)
				..controls(120.9,-16.6)and(121.5,-16.9)..(122.6,-17.9)
				..controls(123,-18.3)and(123.7,-18.8)..(123.9,-18.9)
				..controls(124.6,-19.4)and(126,-19.8)..(127.9,-20.2)
				..controls(129.6,-20.5)and(129.8,-20.6)..(131.1,-21.4)
				..controls(131.5,-21.7)and(132.4,-22.1)..(133.1,-22.3)
				..controls(136.1,-23.2)and(138.3,-24.1)..(140.2,-25)
				..controls(141,-25.5)and(141.5,-25.6)..(142.5,-25.8)
				..controls(145.1,-26.3)and(146,-26.6)..(149.4,-28.1)
				..controls(151.3,-28.9)and(153.5,-29.8)..(154.3,-30.1)
				..controls(155.1,-30.4)and(156.2,-30.8)..(156.6,-31)
				..controls(157.1,-31.2)and(158.2,-31.6)..(159.1,-31.8)
				..controls(160.6,-32.2)and(160.9,-32.3)..(161.7,-32.9)
				..controls(162.4,-33.3)and(162.9,-33.6)..(163.7,-33.8)
				..controls(164.3,-34)and(164.8,-34.1)..(164.9,-34.1)
				..controls(165,-34.1)and(165.6,-34.4)..(166.4,-34.8)
				..controls(167.3,-35.2)and(168.1,-35.4)..(169,-35.6)
				..controls(170,-35.8)and(170.5,-36)..(171.7,-36.6)
				..controls(172.4,-37)and(173.3,-37.4)..(173.6,-37.4)
				..controls(175.6,-37.7)and(176.2,-38)..(176.7,-39.1)
				..controls(176.9,-39.6)and(176.9,-39.7)..(176.7,-40.4)
				..controls(176.4,-41.6)and(175.8,-42.7)..(175.2,-43.4)
				..controls(174.5,-44.1)and(172.6,-45.5)..(171.7,-45.8)
				..controls(170.6,-46.2)and(169.1,-46)..(168.2,-45.4)
				..controls(167.9,-45.2)and(167.3,-44.8)..(166.9,-44.4)
				..controls(166.4,-44.1)and(165.8,-43.6)..(165.4,-43.2)
				..controls(165.1,-42.9)and(164.4,-42.3)..(163.9,-41.9)
				..controls(163.3,-41.5)and(162.7,-40.9)..(162.4,-40.5)
				..controls(161.4,-39.2)and(159.9,-38)..(158.8,-37.8)
				..controls(158.4,-37.8)and(158.3,-37.7)..(158.2,-37.2)
				..controls(158,-36.8)and(157.4,-36.1)..(157.1,-36.1)
				..controls(157,-36.1)and(156.7,-36.1)..(156.5,-36)
				..controls(156.1,-35.9)and(155.9,-35.7)..(155.7,-35.2)
				..controls(155.3,-34.6)and(155.1,-34.4)..(154,-34.1)
				..controls(153.6,-34)and(153.2,-33.7)..(152.9,-33.5)
				..controls(152.4,-32.9)and(151.6,-32.4)..(151.3,-32.4)
				..controls(150.9,-32.4)and(150.4,-32.6)..(150.4,-32.9)
				..controls(150.4,-33.2)and(150.8,-33.4)..(151.4,-33.6)
				..controls(152.1,-33.7)and(152.8,-34.1)..(153.5,-34.9)
				..controls(154.1,-35.5)and(154.3,-35.6)..(155,-35.8)
				..controls(155.7,-36)and(155.7,-36)..(156,-36.6)
				..controls(156.3,-37.3)and(156.7,-37.7)..(157.5,-37.8)
				..controls(158.1,-37.9)and(158.2,-38)..(158.4,-38.8)
				..controls(158.5,-39.4)and(158.6,-39.6)..(159.6,-40.4)
				..controls(160.2,-41)and(160.9,-41.6)..(161.3,-41.8)
				..controls(161.7,-42)and(161.9,-42.3)..(162.2,-42.8)
				..controls(162.4,-43.2)and(162.8,-43.8)..(163.1,-44.1)
				..controls(163.9,-44.8)and(165.1,-46.7)..(165.7,-48)
				..controls(166.2,-49.3)and(166.4,-49.7)..(167.2,-50.5)
				..controls(167.5,-50.8)and(167.9,-51.4)..(168.1,-51.8)
				..controls(168.3,-52.2)and(168.8,-53)..(169.1,-53.6)
				..controls(169.6,-54.3)and(169.8,-54.9)..(169.9,-55.6)
				..controls(170.4,-57.8)and(170.5,-58.1)..(171.2,-59.5)
				..controls(171.6,-60.3)and(172,-61.5)..(172.2,-62.1)
				..controls(172.3,-62.8)and(172.7,-63.7)..(172.9,-64.1)
				..controls(173.2,-64.6)and(173.4,-65.3)..(173.4,-65.6)
				..controls(173.8,-68.8)and(174.2,-70.6)..(174.9,-72.6)
				..controls(175.5,-74.5)and(175.9,-77.1)..(176.1,-80.2)
				..controls(176.3,-83.4)and(175.5,-84.4)..(173.7,-83.4)--cycle;
				\definecolor{ezs}{RGB}{247,179,26}\fill[ezs](-5.7,-156.6)--(-3.4,-157.5)--(61.9,12.7)--(59.7,13.6)--cycle;
				\definecolor{ezs}{RGB}{247,179,22}\fill[ezs](-3.8,-157.3)--(-1.5,-158.2)--(63.9,12)--(61.6,12.9)--cycle;
				\definecolor{ezs}{RGB}{248,179,19}\fill[ezs](-1.8,-158.1)--(0.4,-158.9)--(65.8,11.2)--(63.5,12.1)--cycle;
				\definecolor{ezs}{RGB}{248,178,16}\fill[ezs](0.1,-158.8)--(2.4,-159.7)--(67.7,10.5)--(65.4,11.4)--cycle;
				\definecolor{ezs}{RGB}{249,178,15}\fill[ezs](2,-159.5)--(4.3,-160.4)--(69.6,9.8)--(67.3,10.7)--cycle;
				\definecolor{ezs}{RGB}{251,177,17}\fill[ezs](3.9,-160.3)--(6.2,-161.1)--(71.5,9)--(69.2,9.9)--cycle;
				\definecolor{ezs}{RGB}{252,179,18}\fill[ezs](5.8,-161)--(8.1,-161.9)--(73.4,8.3)--(71.1,9.2)--cycle;
				\definecolor{ezs}{RGB}{252,182,17}\fill[ezs](7.7,-161.7)--(10,-162.6)--(75.3,7.6)--(73,8.5)--cycle;
				\definecolor{ezs}{RGB}{252,185,17}\fill[ezs](9.6,-162.5)--(11.9,-163.4)--(77.2,6.8)--(75,7.7)--cycle;
				\definecolor{ezs}{RGB}{251,189,18}\fill[ezs](11.5,-163.2)--(13.8,-164.1)--(79.2,6.1)--(76.9,7)--cycle;
				\definecolor{ezs}{RGB}{251,191,22}\fill[ezs](13.4,-163.9)--(15.7,-164.8)--(81.1,5.4)--(78.8,6.2)--cycle;
				\definecolor{ezs}{RGB}{252,193,25}\fill[ezs](15.4,-164.7)--(17.6,-165.6)--(83,4.6)--(80.7,5.5)--cycle;
				\definecolor{ezs}{RGB}{253,195,27}\fill[ezs](17.3,-165.4)--(19.6,-166.3)--(84.9,3.9)--(82.6,4.8)--cycle;
				\definecolor{ezs}{RGB}{253,196,31}\fill[ezs](19.2,-166.1)--(21.5,-167)--(86.8,3.2)--(84.5,4)--cycle;
				\definecolor{ezs}{RGB}{254,199,35}\fill[ezs](21.1,-166.9)--(23.4,-167.8)--(88.7,2.4)--(86.4,3.3)--cycle;
				\definecolor{ezs}{RGB}{253,201,38}\fill[ezs](23,-167.6)--(25.3,-168.5)--(90.6,1.7)--(88.3,2.6)--cycle;
				\definecolor{ezs}{RGB}{253,203,40}\fill[ezs](24.9,-168.3)--(27.2,-169.2)--(92.5,1)--(90.3,1.8)--cycle;
				\definecolor{ezs}{RGB}{253,205,42}\fill[ezs](26.8,-169.1)--(29.1,-170)--(94.4,0.2)--(92.2,1.1)--cycle;
				\definecolor{ezs}{RGB}{253,207,45}\fill[ezs](28.7,-169.8)--(31,-170.7)--(96.4,-0.5)--(94.1,0.4)--cycle;
				\definecolor{ezs}{RGB}{253,210,49}\fill[ezs](30.7,-170.5)--(32.9,-171.4)--(98.3,-1.2)--(96,-0.4)--cycle;
				\definecolor{ezs}{RGB}{253,212,54}\fill[ezs](32.6,-171.3)--(34.9,-172.2)--(100.2,-2)--(97.9,-1.1)--cycle;
				\definecolor{ezs}{RGB}{253,214,59}\fill[ezs](34.5,-172)--(36.8,-172.9)--(102.1,-2.7)--(99.8,-1.8)--cycle;
				\definecolor{ezs}{RGB}{253,216,63}\fill[ezs](36.4,-172.7)--(38.7,-173.6)--(104,-3.4)--(101.7,-2.6)--cycle;
				\definecolor{ezs}{RGB}{253,218,65}\fill[ezs](38.3,-173.5)--(40.6,-174.4)--(105.9,-4.2)--(103.6,-3.3)--cycle;
				\definecolor{ezs}{RGB}{254,220,66}\fill[ezs](40.2,-174.2)--(42.5,-175.1)--(107.8,-4.9)--(105.5,-4)--cycle;
				\definecolor{ezs}{RGB}{254,222,68}\fill[ezs](42.1,-175)--(44.4,-175.8)--(109.7,-5.6)--(107.5,-4.8)--cycle;
				\definecolor{ezs}{RGB}{254,224,71}\fill[ezs](44,-175.7)--(46.3,-176.6)--(111.7,-6.4)--(109.4,-5.5)--cycle;
				\definecolor{ezs}{RGB}{254,226,75}\fill[ezs](46,-176.4)--(48.2,-177.3)--(113.6,-7.1)--(111.3,-6.2)--cycle;
				\definecolor{ezs}{RGB}{254,228,79}\fill[ezs](47.9,-177.2)--(50.1,-178)--(115.5,-7.8)--(113.2,-7)--cycle;
				\definecolor{ezs}{RGB}{254,230,82}\fill[ezs](49.8,-177.9)--(52.1,-178.8)--(117.4,-8.6)--(115.1,-7.7)--cycle;
				\definecolor{ezs}{RGB}{254,232,85}\fill[ezs](51.7,-178.6)--(54,-179.5)--(119.3,-9.3)--(117,-8.4)--cycle;
				\definecolor{ezs}{RGB}{254,234,88}\fill[ezs](53.6,-179.4)--(55.9,-180.2)--(121.2,-10)--(118.9,-9.2)--cycle;
				\definecolor{ezs}{RGB}{254,237,91}\fill[ezs](55.5,-180.1)--(57.8,-181)--(123.1,-10.8)--(120.8,-9.9)--cycle;
				\definecolor{ezs}{RGB}{254,239,94}\fill[ezs](57.4,-180.8)--(59.7,-181.7)--(125,-11.5)--(122.8,-10.6)--cycle;
				\definecolor{ezs}{RGB}{254,241,97}\fill[ezs](59.3,-181.6)--(61.6,-182.4)--(127,-12.2)--(124.7,-11.4)--cycle;
				\definecolor{ezs}{RGB}{254,243,101}\fill[ezs](61.2,-182.3)--(63.5,-183.2)--(128.9,-13)--(126.6,-12.1)--cycle;
				\definecolor{ezs}{RGB}{254,245,103}\fill[ezs](63.2,-183)--(65.4,-183.9)--(130.8,-13.7)--(128.5,-12.8)--cycle;
				\definecolor{ezs}{RGB}{254,245,105}\fill[ezs](65.1,-183.8)--(67.4,-184.6)--(132.7,-14.4)--(130.4,-13.6)--cycle;
				\definecolor{ezs}{RGB}{254,245,105}\fill[ezs](67,-184.5)--(69.3,-185.4)--(134.6,-15.2)--(132.3,-14.3)--cycle;
				\definecolor{ezs}{RGB}{254,244,103}\fill[ezs](68.9,-185.2)--(71.2,-186.1)--(136.5,-15.9)--(134.2,-15)--cycle;
				\definecolor{ezs}{RGB}{253,243,101}\fill[ezs](70.8,-186)--(73.1,-186.8)--(138.4,-16.6)--(136.1,-15.8)--cycle;
				\definecolor{ezs}{RGB}{253,241,98}\fill[ezs](72.7,-186.7)--(75,-187.6)--(140.3,-17.4)--(138,-16.5)--cycle;
				\definecolor{ezs}{RGB}{253,239,95}\fill[ezs](74.6,-187.4)--(76.9,-188.3)--(142.2,-18.1)--(140,-17.2)--cycle;
				\definecolor{ezs}{RGB}{253,236,92}\fill[ezs](76.5,-188.2)--(78.8,-189)--(144.2,-18.8)--(141.9,-18)--cycle;
				\definecolor{ezs}{RGB}{253,234,88}\fill[ezs](78.5,-188.9)--(80.7,-189.8)--(146.1,-19.6)--(143.8,-18.7)--cycle;
				\definecolor{ezs}{RGB}{253,232,85}\fill[ezs](80.4,-189.6)--(82.7,-190.5)--(148,-20.3)--(145.7,-19.4)--cycle;
				\definecolor{ezs}{RGB}{253,230,82}\fill[ezs](82.3,-190.4)--(84.6,-191.2)--(149.9,-21)--(147.6,-20.2)--cycle;
				\definecolor{ezs}{RGB}{253,228,78}\fill[ezs](84.2,-191.1)--(86.5,-192)--(151.8,-21.8)--(149.5,-20.9)--cycle;
				\definecolor{ezs}{RGB}{253,226,74}\fill[ezs](86.1,-191.8)--(88.4,-192.7)--(153.7,-22.5)--(151.4,-21.6)--cycle;
				\definecolor{ezs}{RGB}{253,224,70}\fill[ezs](88,-192.6)--(90.3,-193.4)--(155.6,-23.3)--(153.3,-22.4)--cycle;
				\definecolor{ezs}{RGB}{253,222,67}\fill[ezs](89.9,-193.3)--(92.2,-194.2)--(157.5,-24)--(155.3,-23.1)--cycle;
				\definecolor{ezs}{RGB}{253,220,65}\fill[ezs](91.8,-194)--(94.1,-194.9)--(159.5,-24.7)--(157.2,-23.8)--cycle;
				\definecolor{ezs}{RGB}{253,217,63}\fill[ezs](93.7,-194.8)--(96,-195.6)--(161.4,-25.5)--(159.1,-24.6)--cycle;
				\definecolor{ezs}{RGB}{253,215,60}\fill[ezs](95.7,-195.5)--(97.9,-196.4)--(163.3,-26.2)--(161,-25.3)--cycle;
				\definecolor{ezs}{RGB}{253,213,56}\fill[ezs](97.6,-196.2)--(99.9,-197.1)--(165.2,-26.9)--(162.9,-26)--cycle;
				\definecolor{ezs}{RGB}{252,211,53}\fill[ezs](99.5,-197)--(101.8,-197.8)--(167.1,-27.7)--(164.8,-26.8)--cycle;
				\definecolor{ezs}{RGB}{252,209,49}\fill[ezs](101.4,-197.7)--(103.7,-198.6)--(169,-28.4)--(166.7,-27.5)--cycle;
				\definecolor{ezs}{RGB}{252,207,45}\fill[ezs](103.3,-198.4)--(105.6,-199.3)--(170.9,-29.1)--(168.6,-28.2)--cycle;
				\definecolor{ezs}{RGB}{253,205,42}\fill[ezs](105.2,-199.2)--(107.5,-200)--(172.8,-29.9)--(170.6,-29)--cycle;
				\definecolor{ezs}{RGB}{253,202,40}\fill[ezs](107.1,-199.9)--(109.4,-200.8)--(174.7,-30.6)--(172.5,-29.7)--cycle;
				\definecolor{ezs}{RGB}{253,200,38}\fill[ezs](109,-200.6)--(111.3,-201.5)--(176.7,-31.3)--(174.4,-30.4)--cycle;
				\definecolor{ezs}{RGB}{253,198,35}\fill[ezs](111,-201.4)--(113.2,-202.2)--(178.6,-32.1)--(176.3,-31.2)--cycle;
				\definecolor{ezs}{RGB}{253,196,31}\fill[ezs](112.9,-202.1)--(115.2,-203)--(180.5,-32.8)--(178.2,-31.9)--cycle;
				\definecolor{ezs}{RGB}{253,194,28}\fill[ezs](114.8,-202.8)--(117.1,-203.7)--(182.4,-33.5)--(180.1,-32.6)--cycle;
				\definecolor{ezs}{RGB}{253,192,26}\fill[ezs](116.7,-203.6)--(119,-204.5)--(184.3,-34.3)--(182,-33.4)--cycle;
				\definecolor{ezs}{RGB}{253,189,24}\fill[ezs](118.6,-204.3)--(120.9,-205.2)--(186.2,-35)--(183.9,-34.1)--cycle;
				\definecolor{ezs}{RGB}{253,186,20}\fill[ezs](120.5,-205)--(122.8,-205.9)--(188.1,-35.7)--(185.8,-34.8)--cycle;
				\definecolor{ezs}{RGB}{253,184,17}\fill[ezs](122.4,-205.8)--(124.7,-206.7)--(190,-36.5)--(187.8,-35.6)--cycle;
				\definecolor{ezs}{RGB}{252,183,15}\fill[ezs](124.3,-206.5)--(126.6,-207.4)--(192,-37.2)--(189.7,-36.3)--cycle;
				\definecolor{ezs}{RGB}{252,181,14}\fill[ezs](126.3,-207.2)--(128.5,-208.1)--(193.9,-37.9)--(191.6,-37.1)--cycle;
				\definecolor{ezs}{RGB}{252,179,13}\fill[ezs](128.2,-208)--(130.4,-208.9)--(195.8,-38.7)--(193.5,-37.8)--cycle;
				\definecolor{ezs}{RGB}{252,178,12}\fill[ezs](130.1,-208.7)--(132.4,-209.6)--(197.7,-39.4)--(195.4,-38.5)--cycle;
				\definecolor{ezs}{RGB}{251,178,10}\fill[ezs](132,-209.4)--(134.3,-210.3)--(199.6,-40.1)--(197.3,-39.3)--cycle;
				\definecolor{ezs}{RGB}{251,178,10}\fill[ezs](133.9,-210.2)--(136.2,-211.1)--(201.5,-40.9)--(199.2,-40)--cycle;
				\definecolor{ezs}{RGB}{250,177,11}\fill[ezs](135.8,-210.9)--(138.1,-211.8)--(203.4,-41.6)--(201.1,-40.7)--cycle;
				\definecolor{ezs}{RGB}{250,178,11}\fill[ezs](137.7,-211.6)--(140,-212.5)--(205.3,-42.3)--(203.1,-41.5)--cycle;
				\definecolor{ezs}{RGB}{248,180,10}\fill[ezs](139.6,-212.4)--(141.9,-213.3)--(207.3,-43.1)--(205,-42.2)--cycle;
				\definecolor{ezs}{RGB}{247,181,9}\fill[ezs](141.5,-213.1)--(143.8,-214)--(209.2,-43.8)--(206.9,-42.9)--cycle;
				\definecolor{ezs}{RGB}{246,182,9}\fill[ezs](143.5,-213.8)--(145.7,-214.7)--(211.1,-44.5)--(208.8,-43.7)--cycle;
				\definecolor{ezs}{RGB}{246,182,10}\fill[ezs](145.4,-214.6)--(147.7,-215.5)--(213,-45.3)--(210.7,-44.4)--cycle;
			\end{scope}
			% chữ EDUZY (Times đậm, làm dày nét bằng 5 lớp lệch nhẹ; căn đúng khung chữ logo gốc)
			\foreach \dx/\dy in {-0.7/0,0.7/0,0/-0.3,0/0.3,0/0}
			\node[anchor=center,inner sep=0pt,outer sep=0pt,text=ezText,align=center] at (109.6+\dx,-193.6+\dy)
			{\resizebox{1.01cm}{0.22cm}{\fontencoding{T1}\fontfamily{ptm}\fontseries{b}\selectfont EDUZY}};
		\end{tikzpicture}%
	}%
}
%====================== THIẾT LẬP THÂN ĐỀ ======================
%--- Môi trường tự luận {bt}: ex_test 3.3.4 chưa khai báo sẵn nên tạo ở đây
%    (ex_test đã tự móc \fixshowans và \listenumerate cho bt)
\makeatletter
\@ifundefined{bt}{\newtheorem{bt}{\color{gold1}\namebt}}{}
\makeatother
\renewcommand{\sectionmark}[1]{\markright{\thesection\ #1}}
%--- Khung tiêu đề (logo + tối đa 4 dòng; dòng 3, 4 để trống {} thì bỏ qua)
%  Đề thi:    \khungtieude{THI THỬ THPT QUỐC GIA NĂM 2026}{Môn Toán -- Đề số 21}
%                         {(Đề thi gồm: ...)}{Thời gian làm bài: 90 phút (không kể thời gian phát đề).}
%  Chuyên đề: \khungtieude{CHUYÊN ĐỀ 3 -- DÃY SỐ VÀ CÔNG THỨC TRUY HỒI}
%                         {Bồi dưỡng học sinh giỏi Toán THPT}{...}{}
\newcommand{\khungtieude}[4]{%
	\begin{tcolorbox}[colframe=gold1,colback=gold5!45,boxrule=0.5mm]
		\begin{minipage}[c]{3cm}
			\centerline{\eduzylogo[2.4cm]}
		\end{minipage}\hspace{0.5cm}
		\begin{minipage}[c]{12.5cm}
			\vskip 0.15 cm
			\centerline{\ind{#1}}
			\vskip 0.15 cm
			\centerline{\textbf{#2}}
			\if\relax\detokenize{#3}\relax\else
				\vskip 0.15 cm
				\centerline{\textit{#3}}
			\fi
			\if\relax\detokenize{#4}\relax\else
				\vskip 0.15 cm
				\centerline{\textit{#4}}
			\fi
			\vskip 0.15 cm
		\end{minipage}
	\end{tcolorbox}
}
%====================== BỔ SUNG CHO BÀI HỌC, CHUYÊN ĐỀ ======================
\usepackage[shortlabels]{enumitem}
\usepackage{tabularx}
%--- Danh sách gọn: \begin{enumerate}[gon] ... \end{enumerate}
\SetEnumitemKey{gon}{leftmargin=*,itemsep=1pt,parsep=0pt,topsep=3pt}
%--- Khung lý thuyết: \begin{khung}{} ... \end{khung}; tiêu đề chỉ dùng khi thật ngắn ("Mẹo", "Lỗi thường gặp")
\newtcolorbox{khung}[1]{breakable,enhanced,colframe=gold1,colback=gold5!45,
	boxrule=0.4mm,fonttitle=\bfseries\sffamily,coltitle=white,colbacktitle=gold1,
	attach boxed title to top left={xshift=6mm,yshift=-3mm},
	boxed title style={arc=1mm,boxrule=0.3mm},title={#1}}
%--- Tiêu đề mục nhỏ trong một dạng: \muc{Dấu hiệu}, \muc{Phương pháp}, \muc{Bài tự luyện}
\makeatletter
\newcommand{\muc}[1]{\par\medskip\noindent\indam{#1}\par\nobreak\smallskip\@afterindentfalse\@afterheading}
\makeatother
%--- Ô công thức cần nhớ (căn giữa, viền vàng, vừa khít nội dung)
%     \begin{center}\congthuc{$Q(n)=(n-1)!$}\end{center}
\newtcbox{\congthuc}{on line,colframe=gold1,colback=gold5!45,boxrule=0.5pt,arc=2pt,
	left=8pt,right=8pt,top=4pt,bottom=4pt}
%--- Công tắc lời giải bài tự luyện (ví dụ mẫu luôn dùng \loigiai)
%     \loigiaiTL[{đáp số}]{lời giải đầy đủ}
%     \TLfull -> in lời giải đầy đủ;  \TLdapso -> chỉ in một dòng "Đáp số";  \TLan -> ẩn hoàn toàn
%     (\CheDo là dethi thì mọi lời giải và đáp số đều ẩn)
\makeatletter
\newif\ifTLhien   \TLhientrue
\newif\ifTLdapso  \TLdapsofalse
\newcommand{\TLfull}{\TLhientrue\TLdapsofalse}
\newcommand{\TLdapso}{\TLhientrue\TLdapsotrue}
\newcommand{\TLan}{\TLhienfalse}
\newcommand{\loigiaiTL}[2][]{%
	\ifTLhien
		\ifTLdapso
			\def\@tlarg{#1}%
			\ifx\@tlarg\@empty\else
				\ifdefstring{\kindOPTN}{dethi}{}{\par\noindent{\bfseries\color{gold1}Đáp số:} #1\par}%
			\fi
		\else
			\loigiai{#2}%
		\fi
	\fi}
\makeatother
\TLdapso% mặc định: chỉ in đáp số bài tự luyện
%====================== BẢNG ĐÚNG/SAI (\choiceTFt) ======================
%  Hàng tiêu đề nền nâu đồng chữ trắng (Phát biểu | Đ | S), các hàng tô xen kẽ kem / trắng,
%  đường kẻ vàng kim, nhãn a) b) c) d) màu nâu đồng. Bản có đáp án: cột được chọn có vòng tròn
%  nâu đồng với dấu tích trắng (câu đúng tích ở cột Đúng, câu sai tích ở cột Sai).
%  Tiêu đề ghi đủ "Đúng | Sai": \OPTN{viettat=0} sau \input (cột Đúng/Sai rộng ra, phát biểu hẹp lại).
%  Dùng bảng gốc của ex_test: \def\BangDSgoc{1} trước \input{preamble-eduzy-v2}.
\providecommand{\BangDSgoc}{0}
\ifnum\BangDSgoc=0
\usepackage{colortbl,pifont}
\colorlet{ezDSke}{gold2!60!gold1}% đường kẻ bảng (nét 0,4pt như ex_test để không đổi chiều cao bảng)
\colorlet{ezDSle}{gold5!60}% hàng a), c)
\colorlet{ezDSchan}{white}% hàng b), d)
%--- dấu chọn ở cột Đúng/Sai
%    (bề rộng 0, căn giữa ô: không làm đổi bề rộng cột của ex_test, không báo tràn lề)
\newcommand{\EZdauchon}{\makebox[0pt]{\tikz[baseline=-1ex]\node[circle,fill=gold1,inner sep=0pt,
	minimum size=1.2em,text=white,font=\small]{\ding{51}};}}
\def\tickT{\EZdauchon}\def\tickF{\EZdauchon}
\AtBeginDocument{\sbox{\ccTFTab}{\EZdauchon}}
%--- nhãn a) b) c) d) màu nâu đồng (ở bản lời giải, nhãn câu đúng vẫn đóng khung như ex_test)
\renewcommand{\FalseTF}{\stepcounter{dapan}\textbf{\color{gold1}\DapAnTF\sepTF}\ignorespaces}
%--- trước khi dựng bảng: kẻ màu vàng kim (bề rộng cột giữ nguyên như ex_test); sau bảng trả lại màu đen
\newcount\EZhangTF
\newcommand{\EZbatdauTF}{%
	\global\EZhangTF=0\relax
	\arrayrulecolor{ezDSke}}
\newcommand{\EZketthucTF}{\arrayrulecolor{black}}
%--- hàng tiêu đề nền nâu đồng, chữ trắng
\makeatletter
\newcommand{\EZo@TF}[1]{\centering\bfseries\color{white}#1}
\newcommand{\EZtieudeTF}{%
	\ifcase\cotDS
		\ifnum\cotSTT=1\relax
			\def\rowHeader{\hline\rowcolor{gold1}\EZo@TF{TT}&\EZo@TF{\phatbieu}\tabularnewline}%
		\else
			\def\rowHeader{\hline\rowcolor{gold1}\EZo@TF{\phatbieu}\tabularnewline}%
		\fi
	\or
		\ifnum\cotSTT=1\relax
			\def\rowHeader{\hline\rowcolor{gold1}\EZo@TF{TT}&\EZo@TF{\phatbieu}&\EZo@TF{\Dung}\tabularnewline}%
		\else
			\def\rowHeader{\hline\rowcolor{gold1}\EZo@TF{\phatbieu}&\EZo@TF{\Dung}\tabularnewline}%
		\fi
	\or
		\ifnum\cotSTT=1\relax
			\def\rowHeader{\hline\rowcolor{gold1}\EZo@TF{TT}&\EZo@TF{\phatbieu}&\EZo@TF{\Dung}&\EZo@TF{\Sai}\tabularnewline}%
		\else
			\def\rowHeader{\hline\rowcolor{gold1}\EZo@TF{\phatbieu}&\EZo@TF{\Dung}&\EZo@TF{\Sai}\tabularnewline}%
		\fi
	\fi}
%--- tô xen kẽ từng hàng (gọi ngay sau \hline đầu mỗi hàng)
\newcommand{\EZmauhangTF}{%
	\noalign{\global\advance\EZhangTF\@ne}%
	\ifodd\EZhangTF\rowcolor{ezDSle}\else\rowcolor{ezDSchan}\fi}
%--- gắn vào \choiceTFt / \writeChoiceTFt của ex_test
\xpretocmd{\choiceTFt}{\EZbatdauTF}{}{\typeout{[preamble-eduzy-v2] Không gắn được bảng Đúng/Sai (đầu bảng)}}
\xpatchcmd{\choiceTFt}{\ifnum\dongPB=0\def\rowHeader{}\fi}{\EZtieudeTF\ifnum\dongPB=0\def\rowHeader{}\fi}
	{}{\typeout{[preamble-eduzy-v2] Không gắn được bảng Đúng/Sai (hàng tiêu đề)}}
\xpatchcmd{\writeChoiceTFt}{\hline\xdef\dotEX{\sdotEX}}{\hline\EZmauhangTF\xdef\dotEX{\sdotEX}}
	{}{\typeout{[preamble-eduzy-v2] Không gắn được bảng Đúng/Sai (tô hàng)}}
\xpatchcmd{\writeChoiceTFt}{{\end{longtable}}}{{\end{longtable}\EZketthucTF}}
	{}{\typeout{[preamble-eduzy-v2] Không gắn được bảng Đúng/Sai (cuối longtable)}}
\xpatchcmd{\writeChoiceTFt}{{\end{xtabular}\end{center}\medskip\noindent}}{{\end{xtabular}\end{center}\EZketthucTF\medskip\noindent}}
	{}{\typeout{[preamble-eduzy-v2] Không gắn được bảng Đúng/Sai (cuối xtabular)}}
\fi
%=====================================================================
\renewcommand{\parallel}{\mathrel{/\!\!/}}
.
\providecommand{\BangDSgoc}{0}
\ifnum\BangDSgoc=0
\usepackage{colortbl,pifont}
\colorlet{ezDSke}{gold2!60!gold1}% đường kẻ bảng (nét 0,4pt như ex_test để không đổi chiều cao bảng)
\colorlet{ezDSle}{gold5!60}% hàng a), c)
\colorlet{ezDSchan}{white}% hàng b), d)
%--- dấu chọn ở cột Đúng/Sai
%    (bề rộng 0, căn giữa ô: không làm đổi bề rộng cột của ex_test, không báo tràn lề)
\newcommand{\EZdauchon}{\makebox[0pt]{\tikz[baseline=-1ex]\node[circle,fill=gold1,inner sep=0pt,
	minimum size=1.2em,text=white,font=\small]{\ding{51}};}}
\def\tickT{\EZdauchon}\def\tickF{\EZdauchon}
\AtBeginDocument{\sbox{\ccTFTab}{\EZdauchon}}
%--- nhãn a) b) c) d) màu nâu đồng (ở bản lời giải, nhãn câu đúng vẫn đóng khung như ex_test)
\renewcommand{\FalseTF}{\stepcounter{dapan}\textbf{\color{gold1}\DapAnTF\sepTF}\ignorespaces}
%--- trước khi dựng bảng: kẻ màu vàng kim (bề rộng cột giữ nguyên như ex_test); sau bảng trả lại màu đen
\newcount\EZhangTF
\newcommand{\EZbatdauTF}{%
	\global\EZhangTF=0\relax
	\arrayrulecolor{ezDSke}}
\newcommand{\EZketthucTF}{\arrayrulecolor{black}}
%--- hàng tiêu đề nền nâu đồng, chữ trắng
\makeatletter
\newcommand{\EZo@TF}[1]{\centering\bfseries\color{white}#1}
\newcommand{\EZtieudeTF}{%
	\ifcase\cotDS
		\ifnum\cotSTT=1\relax
			\def\rowHeader{\hline\rowcolor{gold1}\EZo@TF{TT}&\EZo@TF{\phatbieu}\tabularnewline}%
		\else
			\def\rowHeader{\hline\rowcolor{gold1}\EZo@TF{\phatbieu}\tabularnewline}%
		\fi
	\or
		\ifnum\cotSTT=1\relax
			\def\rowHeader{\hline\rowcolor{gold1}\EZo@TF{TT}&\EZo@TF{\phatbieu}&\EZo@TF{\Dung}\tabularnewline}%
		\else
			\def\rowHeader{\hline\rowcolor{gold1}\EZo@TF{\phatbieu}&\EZo@TF{\Dung}\tabularnewline}%
		\fi
	\or
		\ifnum\cotSTT=1\relax
			\def\rowHeader{\hline\rowcolor{gold1}\EZo@TF{TT}&\EZo@TF{\phatbieu}&\EZo@TF{\Dung}&\EZo@TF{\Sai}\tabularnewline}%
		\else
			\def\rowHeader{\hline\rowcolor{gold1}\EZo@TF{\phatbieu}&\EZo@TF{\Dung}&\EZo@TF{\Sai}\tabularnewline}%
		\fi
	\fi}
%--- tô xen kẽ từng hàng (gọi ngay sau \hline đầu mỗi hàng)
\newcommand{\EZmauhangTF}{%
	\noalign{\global\advance\EZhangTF\@ne}%
	\ifodd\EZhangTF\rowcolor{ezDSle}\else\rowcolor{ezDSchan}\fi}
%--- gắn vào \choiceTFt / \writeChoiceTFt của ex_test
\xpretocmd{\choiceTFt}{\EZbatdauTF}{}{\typeout{[preamble-eduzy-v2] Không gắn được bảng Đúng/Sai (đầu bảng)}}
\xpatchcmd{\choiceTFt}{\ifnum\dongPB=0\def\rowHeader{}\fi}{\EZtieudeTF\ifnum\dongPB=0\def\rowHeader{}\fi}
	{}{\typeout{[preamble-eduzy-v2] Không gắn được bảng Đúng/Sai (hàng tiêu đề)}}
\xpatchcmd{\writeChoiceTFt}{\hline\xdef\dotEX{\sdotEX}}{\hline\EZmauhangTF\xdef\dotEX{\sdotEX}}
	{}{\typeout{[preamble-eduzy-v2] Không gắn được bảng Đúng/Sai (tô hàng)}}
\xpatchcmd{\writeChoiceTFt}{{\end{longtable}}}{{\end{longtable}\EZketthucTF}}
	{}{\typeout{[preamble-eduzy-v2] Không gắn được bảng Đúng/Sai (cuối longtable)}}
\xpatchcmd{\writeChoiceTFt}{{\end{xtabular}\end{center}\medskip\noindent}}{{\end{xtabular}\end{center}\EZketthucTF\medskip\noindent}}
	{}{\typeout{[preamble-eduzy-v2] Không gắn được bảng Đúng/Sai (cuối xtabular)}}
\fi
%=====================================================================
\renewcommand{\parallel}{\mathrel{/\!\!/}}

%     %=====================================================================
% preamble-vatli-hinh.tex -- BỘ MACRO HÌNH VẬT LÍ (dùng chung cho mọi đề Vật lí)
% Gọi ngay sau preamble chính:
%     %============================================================
%  PREAMBLE EDUZY v2 - dùng chung cho đề thi, bài học, chuyên đề (gói ex_test)
%  Khác bản preamble-eduzy.tex ở đầu trang, chân trang: dải nền nâu đồng tràn ngang,
%  logo + Eduzy.com.vn, ô chữ vàng kem cạnh xiên, chân trang "Trang (số)";
%  logo Eduzy chìm giữa mọi trang (tắt: \logochimtat, bật lại: \logochimbat).
%  Cách dùng ở file chính (xem mau-chuan.tex, mau-chuyen-de.tex):
%      \def\CheDo{loigiai}          % loigiai | dethi | solcolor | color | book
%      \def\tenchuong{<dòng chữ trên đầu trang>}
%      \def\thuonghieu{<chữ ở chân trang lẻ>}   % không khai báo thì dùng mặc định
%      \def\tacgia{<chữ ở chân trang chẵn>}      % không khai báo thì dùng mặc định
%      \def\khauhieu{<chữ cạnh logo ở đầu trang>}  % tùy chọn, mặc định Eduzy.com.vn
%      \input{preamble-eduzy-v2}
%      \TLdapso                      % tài liệu: \TLfull | \TLdapso | \TLan
%      \begin{document} ... \end{document}
%  Logo Eduzy vẽ bằng TikZ ngay trong file (\eduzylogo), không cần file ảnh.
%  Biên dịch: pdfLaTeX (máy có gói vietnam) hoặc XeLaTeX.
%============================================================
\documentclass[12pt,a4paper,twoside]{article}%<-- twoside: bat dau trang/chan trang chan-le
\usepackage{iftex}
\usepackage{amsmath,amssymb,tikz,mathrsfs,tkz-tab}
\usepackage{graphicx}
\usepackage{pgfplots}% đồ thị hàm số
\pgfplotsset{compat=1.17}
%\usepackage{mathptmx}
\usepackage{mathpazo}
\ifXeTeX% XeLaTeX: font chữ có đủ dấu tiếng Việt, giữ nguyên font toán mathpazo
	\usepackage[no-math]{fontspec}
	\setmainfont{TeX Gyre Pagella}
	\newfontfamily\EZqag{TeX Gyre Adventor}[NFSSFamily=qag,Ligatures=TeX]
\fi
\usepackage[many]{tcolorbox}
\usetikzlibrary{math,through,calc,intersections,angles,quotes,shapes,shapes.geometric,arrows,patterns,matrix,chains,arrows.meta,decorations.shapes,decorations.fractals,decorations.markings,shadows,decorations.pathreplacing}
\usepackage[top=1.6cm, bottom=1.8cm, left=1.8cm, right=1.8cm]{geometry}% v2: top 1.6cm (bản cũ 1.5cm) để chữ không sát dải đầu trang
\providecommand{\CheDo}{loigiai}% loigiai | dethi | solcolor | color | book
\expandafter\usepackage\expandafter[\CheDo]{ex_test}
% \tenchuong: khai báo ở file chính (trước \input) hoặc \renewcommand sau \input
%============================================================
%  [GỘP TỪ eduzy-de.tex] EDUZY - MÀU NÂU ĐỒNG / VÀNG KIM
%  + ĐẦU TRANG, CHÂN TRANG KIỂU DẢI NGANG (v2), ĐỐI XỨNG TRANG CHẴN / TRANG LẺ.
%  (không cần file eduzy-de.tex riêng nữa)
%  Khối này phải đặt SAU \usepackage[...]{ex_test} và SAU \def\tenchuong{...}
%============================================================
\usepackage{lastpage}
\usepackage{needspace}
%--- DÀN TRANG ĐỀ THI: tách câu (THÊM 05/11/2026; chỉ thêm lệnh mới, KHÔNG đổi hành vi mặc định) ---
%  Hai công tắc, mặc định tắt (= preamble chạy y như cũ). Bật bằng cách đặt TRƯỚC \input{preamble-eduzy-v2}
%  trong file đề (tools/dan_trang/dan_trang.py tự ghi khi đặt \NeedspaceCau):
%    \def\TachCauDS{1}   câu Đúng/Sai được tách: câu dẫn (kèm hình) ở cuối trang, bảng Phát biểu sang trang sau
%                        (câu dẫn và hình không tách nhau, bảng không tách giữa hai trang)
%    \def\TachCauLG{1}   bản lời giải (\CheDo khác dethi): phần đề của câu cao hơn 1/3 trang được phép tách
%  Lệnh dùng trong file đề (giá trị do tools/dan_trang/dan_trang.py đo ở bản dethi):
%    \NeedspaceCau[<cao câu dẫn + hình>]{<cao cả khối câu>}   ngay trên \begin{ex} (câu đầu phần: trên tiêu đề phần)
%        công tắc tắt: đúng bằng \Needspace{<cao cả khối câu>} như cũ;
%        \TachCauDS=1 và có [..] (câu Đúng/Sai được chọn tách): chỉ giữ câu dẫn + hình;
%        \TachCauLG=1 ở bản lời giải: đòi tối đa 1/3 \textheight.
%    \NeedspaceBang{<cao bảng Phát biểu>}%   ngay trước \choiceTFt: \TachCauDS=1 thì giữ trọn bảng trên một trang;
%        công tắc tắt: không làm gì.
\providecommand{\TachCauDS}{0}
\providecommand{\TachCauLG}{0}
\newlength{\EZnsCau}
\newcommand{\NeedspaceCau}[2][]{%
	\setlength{\EZnsCau}{#2}%
	\ifnum\TachCauDS=1 \if\relax\detokenize{#1}\relax\else\setlength{\EZnsCau}{#1}\fi\fi
	\ifnum\TachCauLG=1 \ifdefstring{\kindOPTN}{dethi}{}{\ifdim\EZnsCau>0.3333\textheight\setlength{\EZnsCau}{0.3333\textheight}\fi}\fi
	\Needspace{\EZnsCau}}
\newcommand{\NeedspaceBang}[1]{\ifnum\TachCauDS=1 \Needspace{#1}\fi}

%--- Thang màu: đậm -> nhạt
\definecolor{gold1}{RGB}{126,96,0}    % nâu đồng đậm  #7E6000
\definecolor{gold2}{RGB}{192,144,3}   % vàng kim      #C09003
\definecolor{gold3}{RGB}{255,215,99}  % vàng nhạt     #FFD763
\definecolor{gold4}{RGB}{255,231,153} % vàng rất nhạt #FFE799
\definecolor{gold5}{RGB}{255,243,205} % vàng kem      #FFF3CD
%--- Ghi đè tông xanh/đỏ cũ
\definecolor{myblue}{RGB}{126,96,0}
\colorlet{blue}{gold1}
\colorlet{cyan}{gold2}
\colorlet{teal}{gold2}

%--- Tên chương mặc định (file đề nên tự khai báo lại)
\providecommand{\tenchuong}{}
%--- Chân trang
\providecommand{\thuonghieu}{TOÁN 12 -- Eduzy.com.vn}
\providecommand{\tacgia}{Tô Gia Bảo -- 0353703783}

%--- Chữ cạnh logo ở đầu trang
\providecommand{\khauhieu}{Eduzy.com.vn}

%====================== THÔNG SỐ ĐẦU TRANG, CHÂN TRANG (v2) ======================
%  Dải nền đậm tràn ngang trang; ô chữ vàng kem có cạnh xiên, kèm 2 vạch xiên chuyển màu.
%  Đầu trang lẻ : logo + \khauhieu (trái)   | ô chữ \tenchuong (phải)
%  Chân trang lẻ: ô chữ \thuonghieu (trái)  | Trang (số) (phải)
%  Trang chẵn soi gương trang lẻ; chân trang chẵn ghi \tacgia.
%--- Màu: chỉ dùng thang vàng kim Eduzy
\colorlet{ezNen}{gold1}% dải nền đậm
\colorlet{ezVach1}{gold3}% vạch xiên sát dải nền (đậm hơn)
\colorlet{ezVach2}{gold4}% vạch xiên sát ô chữ (nhạt hơn)
\colorlet{ezO}{gold5}% ô chữ vàng kem
%--- Kích thước
\newlength{\EZdai}\setlength{\EZdai}{1.1cm}% chiều cao dải nền
\newlength{\EZo}\setlength{\EZo}{0.68cm}% chiều cao ô chữ
\newlength{\EZxien}\setlength{\EZxien}{0.55cm}% độ lệch ngang của cạnh xiên
\newlength{\EZvach}\setlength{\EZvach}{0.75cm}% bề rộng mỗi vạch xiên
\newlength{\EZkhe}\setlength{\EZkhe}{0.12cm}% khe trắng giữa vạch và ô chữ
\newlength{\EZdem}\setlength{\EZdem}{0.4cm}% khoảng đệm từ chữ tới cạnh xiên
\newlength{\EZle}\setlength{\EZle}{1.8cm}% lề trái/phải của chữ (bằng lề thân bài)
\newlength{\EZrongmin}\setlength{\EZrongmin}{4.5cm}% ô chữ rộng tối thiểu
\newlength{\EZrong}% bề rộng chữ trong ô (tính tự động)
%--- Font
\newcommand{\EZfontO}{\fontsize{10.5pt}{12pt}\selectfont\bfseries}% chữ trong ô
\newcommand{\EZfontKH}{\fontfamily{qag}\fontsize{11pt}{12pt}\selectfont\bfseries}% \khauhieu cạnh logo
\newcommand{\EZfontTrang}{\fontsize{11.5pt}{12pt}\selectfont\bfseries}% chữ "Trang" và số trang
%--- Đo bề rộng chữ trong ô
\newcommand{\EZdorong}[1]{%
	\settowidth{\EZrong}{\EZfontO #1}%
	\ifdim\EZrong<\EZrongmin\setlength{\EZrong}{\EZrongmin}\fi}
%--- Logo + chữ trắng \khauhieu (dựng sẵn một lần khi vào \begin{document})
\newsavebox{\EZlogobox}
\AtBeginDocument{\sbox{\EZlogobox}{%
		\raisebox{-0.5\height}{\eduzylogo[1cm]}\hspace{0.15cm}%
		\raisebox{-0.5\height}{\color{white}\EZfontKH\khauhieu}}}
%--- Hai vạch xiên + khe trắng. #1: điểm gốc ở mép dưới, #2: hướng (1 sang phải, -1 sang trái),
%    #3: hướng nghiêng (1 kiểu "/", -1 kiểu "\")
\newcommand{\EZvachxien}[3]{%
	\fill[white] (#1) --++({(#3)*\EZxien},\EZo) --++({(#2)*\EZkhe},0) --++({-(#3)*\EZxien},-\EZo) --cycle;
	\fill[ezVach2] ([xshift=(#2)*\EZkhe]#1) --++({(#3)*\EZxien},\EZo) --++({(#2)*\EZvach},0) --++({-(#3)*\EZxien},-\EZo) --cycle;
	\fill[ezVach1] ([xshift=(#2)*(\EZkhe+\EZvach)]#1) --++({(#3)*\EZxien},\EZo) --++({(#2)*\EZvach},0) --++({-(#3)*\EZxien},-\EZo) --cycle;
	\draw[white,line width=0.6pt] ([xshift=(#2)*(\EZkhe+\EZvach)]#1) --++({(#3)*\EZxien},\EZo);
}
%--- Chữ trong ô (căn giữa theo chiều cao ô, không lệch vì dấu tiếng Việt)
\tikzset{ezchu/.style={inner sep=0pt,text=ezNen,text height=0.72em,text depth=0.22em}}

%====================== LOGO CHÌM GIỮA TRANG ======================
%  Vẽ cùng đầu trang nên nằm dưới chữ của thân bài, có trên mọi trang.
%  Tắt: \logochimtat (từ trang đó trở đi), bật lại: \logochimbat.
%  Đổi cỡ, độ đậm: khai báo trước \input, ví dụ \def\EZchimcao{9cm} \def\EZchimdo{0.1}
\providecommand{\EZchimcao}{10cm}% chiều cao logo chìm
\providecommand{\EZchimdo}{0.12}% độ đậm (0 = trong suốt, 1 = rõ hẳn)
\newif\ifEZchim \EZchimtrue
\newcommand{\logochimtat}{\EZchimfalse}
\newcommand{\logochimbat}{\EZchimtrue}
%  Dựng sẵn logo mờ một lần khi vào \begin{document}; transparency group làm mờ cả logo
%  một lượt nên không lộ vệt ở chỗ các dải màu chồng nhau.
\newsavebox{\EZchimtho}
\newsavebox{\EZchimbox}
\AtBeginDocument{%
	\sbox{\EZchimtho}{\eduzylogo[\EZchimcao]}%
	\sbox{\EZchimbox}{\begin{tikzpicture}
			\begin{scope}[transparency group,opacity=\EZchimdo]
				\node[inner sep=0pt] {\usebox{\EZchimtho}};
			\end{scope}
	\end{tikzpicture}}}
\newcommand{\EZvechim}{%
	\ifEZchim
		\node[inner sep=0pt] at (current page.center) {\usebox{\EZchimbox}};
	\fi}

%====================== ĐẦU TRANG - TRANG LẺ ======================
%  (logo bên trái, ô chữ \tenchuong bên phải)
\newcommand{\myfancyheadOdd}{%
	\EZdorong{\tenchuong}%
	\begin{tikzpicture}[remember picture,overlay]
		\coordinate (T) at (current page.north west);
		\coordinate (P) at (current page.north east);
		%--- logo chìm giữa trang
		\EZvechim
		%--- dải nền
		\fill[ezNen] (T) rectangle ([yshift=-\EZdai]P);
		%--- ô chữ: X là góc dưới bên trái, cạnh trái xiên kiểu "/"
		\coordinate (X) at ([xshift=-\EZle-\EZrong-\EZdem-\EZxien,yshift=-\EZdai]P);
		\fill[ezO] (X) --++(\EZxien,\EZo) -- ([yshift=-\EZdai+\EZo]P) -- ([yshift=-\EZdai]P) --cycle;
		\EZvachxien{X}{-1}{1}
		\node[ezchu,anchor=east] at ([xshift=-\EZle,yshift=-\EZdai+0.5*\EZo]P) {\EZfontO\tenchuong};
		%--- logo + tên
		\node[inner sep=0pt,anchor=west] at ([xshift=\EZle,yshift=-0.5*\EZdai]T) {\usebox{\EZlogobox}};
	\end{tikzpicture}%
}

%====================== ĐẦU TRANG - TRANG CHẴN =====================
%  (soi gương: ô chữ \tenchuong bên trái, logo bên phải)
\newcommand{\myfancyheadEven}{%
	\EZdorong{\tenchuong}%
	\begin{tikzpicture}[remember picture,overlay]
		\coordinate (T) at (current page.north west);
		\coordinate (P) at (current page.north east);
		%--- logo chìm giữa trang
		\EZvechim
		%--- dải nền
		\fill[ezNen] (T) rectangle ([yshift=-\EZdai]P);
		%--- ô chữ: X là góc dưới bên phải, cạnh phải xiên kiểu "\"
		\coordinate (X) at ([xshift=\EZle+\EZrong+\EZdem+\EZxien,yshift=-\EZdai]T);
		\fill[ezO] (X) --++(-\EZxien,\EZo) -- ([yshift=-\EZdai+\EZo]T) -- ([yshift=-\EZdai]T) --cycle;
		\EZvachxien{X}{1}{-1}
		\node[ezchu,anchor=west] at ([xshift=\EZle,yshift=-\EZdai+0.5*\EZo]T) {\EZfontO\tenchuong};
		%--- logo + tên
		\node[inner sep=0pt,anchor=east] at ([xshift=-\EZle,yshift=-0.5*\EZdai]P) {\usebox{\EZlogobox}};
	\end{tikzpicture}%
}

%====================== CHÂN TRANG - TRANG LẺ ======================
%  (ô chữ \thuonghieu bên trái, "Trang (số)" bên phải)
\newcommand{\myfancyfootOdd}{%
	\EZdorong{\thuonghieu}%
	\begin{tikzpicture}[remember picture,overlay]
		\coordinate (T) at (current page.south west);
		\coordinate (P) at (current page.south east);
		%--- dải nền
		\fill[ezNen] (T) rectangle ([yshift=\EZdai]P);
		%--- ô chữ sát mép trên dải: X là góc dưới bên phải, cạnh phải xiên kiểu "/"
		\coordinate (X) at ([xshift=\EZle+\EZrong+\EZdem,yshift=\EZdai-\EZo]T);
		\fill[ezO] ([yshift=\EZdai-\EZo]T) -- (X) --++(\EZxien,\EZo) -- ([yshift=\EZdai]T) --cycle;
		\EZvachxien{X}{1}{1}
		\node[ezchu,anchor=west] at ([xshift=\EZle,yshift=\EZdai-0.5*\EZo]T) {\EZfontO\thuonghieu};
		%--- số trang
		\node[circle,draw=white,line width=0.9pt,text=white,inner sep=0pt,minimum size=0.62cm,anchor=east]
		(so) at ([xshift=-\EZle,yshift=0.5*\EZdai]P) {\EZfontTrang\thepage};
		\node[inner sep=0pt,text=white,anchor=base east] (tr) at ([xshift=-0.15cm]so.base west) {\EZfontTrang Trang};
	\end{tikzpicture}%
}

%====================== CHÂN TRANG - TRANG CHẴN ====================
%  (soi gương: "Trang (số)" bên trái, ô chữ \tacgia bên phải)
\newcommand{\myfancyfootEven}{%
	\EZdorong{\tacgia}%
	\begin{tikzpicture}[remember picture,overlay]
		\coordinate (T) at (current page.south west);
		\coordinate (P) at (current page.south east);
		%--- dải nền
		\fill[ezNen] (T) rectangle ([yshift=\EZdai]P);
		%--- ô chữ sát mép trên dải: X là góc dưới bên trái, cạnh trái xiên kiểu "\"
		\coordinate (X) at ([xshift=-\EZle-\EZrong-\EZdem,yshift=\EZdai-\EZo]P);
		\fill[ezO] ([yshift=\EZdai-\EZo]P) -- (X) --++(-\EZxien,\EZo) -- ([yshift=\EZdai]P) --cycle;
		\EZvachxien{X}{-1}{-1}
		\node[ezchu,anchor=east] at ([xshift=-\EZle,yshift=\EZdai-0.5*\EZo]P) {\EZfontO\tacgia};
		%--- số trang
		%    (vòng tròn ẩn so0 nằm giữa dải để lấy đường chân chữ, giống trang lẻ)
		\node[circle,inner sep=0pt,minimum size=0.62cm] (so0) at ([yshift=0.5*\EZdai]T) {\EZfontTrang\phantom{\thepage}};
		\node[inner sep=0pt,text=white,anchor=base west] (tr) at ($(so0.base -| T)+(\EZle,0)$) {\EZfontTrang Trang};
		\node[circle,draw=white,line width=0.9pt,text=white,inner sep=0pt,minimum size=0.62cm,anchor=base west]
		(so) at ([xshift=0.15cm]tr.base east) {\EZfontTrang\thepage};
	\end{tikzpicture}%
}

%--- Giữ tương thích với các file cũ đang gọi \myfancyhead, \myfancyfoot
\newcommand{\myfancyhead}{\myfancyheadOdd}
\newcommand{\myfancyfoot}{\myfancyfootOdd}

%======================= GẮN VÀO TRANG =============================
\pagestyle{fancy}
\fancyhf{}
\fancyhead[LO]{\myfancyheadOdd}
\fancyhead[RE]{\myfancyheadEven}
\fancyfoot[LO]{\myfancyfootOdd}
\fancyfoot[RE]{\myfancyfootEven}
\renewcommand{\headrulewidth}{0pt}
\renewcommand{\footrulewidth}{0pt}
\fancypagestyle{plain}{%
	\fancyhf{}
	\fancyhead[LO]{\myfancyheadOdd}
	\fancyhead[RE]{\myfancyheadEven}
	\fancyfoot[LO]{\myfancyfootOdd}
	\fancyfoot[RE]{\myfancyfootEven}
	\renewcommand{\headrulewidth}{0pt}
	\renewcommand{\footrulewidth}{0pt}
}
%=================== HẾT PHẦN GỘP TỪ eduzy-de.tex ===================
\newcommand{\indam}[1]{
	{\fontsize{12pt}{9pt}\sffamily \selectfont\bfseries\color{gold1} #1}
}
\newcommand{\ind}[1]{
	{\fontsize{14pt}{9pt}\fontfamily{qag} \selectfont\bfseries\color{gold1} #1}
}
\renewcommand{\FalseEX}{\stepcounter{dapan}\textbf{\selectfont\color{gold1}\Alph{dapan}.}\ignorespaces}
\renewtheorem{ex}{\color{gold1}Câu}
\renewcommand{\baselinestretch}{1.1}
\newcommand{\hoac}[1]{
	\left[\begin{aligned}#1\end{aligned}\right.}
\newcommand{\heva}[1]{
	\left\{\begin{aligned}#1\end{aligned}\right.}
\usepackage{esvect}
\def\vec{\vv}
\def\vt{\vv}
\def\overrightarrow{\vv}
%====================== LOGO EDUZY VẼ BẰNG TIKZ ======================
%  Đường viền dò từ logo gốc; dải màu vàng kim đo trực tiếp trên ảnh gốc
%  (nghiêng 21 độ, tô bằng 80 dải hẹp nên hiển thị giống nhau trên mọi trình xem PDF).
%  Dùng: \eduzylogo          -> cao 2.4cm (mặc định)
%        \eduzylogo[1.5cm]   -> cao 1.5cm
\definecolor{ezText}{RGB}{245,200,46}
\newcommand{\eduzylogo}[1][2.4cm]{%
	\resizebox{!}{#1}{%
		\begin{tikzpicture}[x=0.01cm,y=0.01cm]
			\path[use as bounding box] (0,0) rectangle (224,-224);
			% mũ tốt nghiệp + vương miện
			\begin{scope}
				\clip
				(119.4,-176.1)..controls(118.9,-176.1)and(116.8,-176)..(114.6,-176)
				..controls(101.4,-176)and(85.8,-175.9)..(83.7,-175.7)
				..controls(81.8,-175.6)and(81,-175.6)..(80.2,-175.7)
				..controls(78.8,-175.9)and(77.7,-175.9)..(76.2,-175.7)
				..controls(74.9,-175.5)and(70.8,-175.5)..(64.3,-175.7)
				..controls(59.1,-175.9)and(58.3,-175.7)..(58.6,-173.7)
				..controls(58.6,-172.8)and(59.6,-169.6)..(60.3,-168)
				..controls(60.8,-166.9)and(61.4,-165.8)..(62.3,-165)
				..controls(63.1,-164.2)and(64.3,-163.4)..(65,-163.2)
				..controls(65.5,-163.1)and(72.3,-163)..(76.6,-163.1)
				..controls(77.9,-163.1)and(80.1,-163.1)..(81.4,-163)
				..controls(82.8,-162.9)and(86.3,-162.9)..(89.2,-162.9)
				..controls(93,-162.9)and(94.9,-162.8)..(96,-162.7)
				..controls(97.2,-162.5)and(101,-162.5)..(114,-162.4)
				..controls(123,-162.3)and(134.9,-162.2)..(140.2,-162.1)
				..controls(150.1,-162)and(152.2,-162.1)..(153.6,-162.4)
				..controls(154.4,-162.6)and(157.1,-163.8)..(158.1,-164.4)
				..controls(158.8,-164.8)and(159.7,-165.7)..(159.9,-166.3)
				..controls(160,-166.5)and(160.2,-167.1)..(160.2,-167.6)
				..controls(160.3,-168)and(160.6,-169.4)..(160.9,-170.6)
				..controls(161.2,-171.7)and(161.5,-173.1)..(161.6,-173.6)
				..controls(161.8,-174.4)and(161.8,-174.5)..(161.5,-174.9)
				..controls(161.2,-175.6)and(160.8,-175.9)..(159.9,-176)
				..controls(159.2,-176.1)and(154.1,-176.2)..(151.5,-176.2)
				..controls(148.1,-176.1)and(122.6,-176.1)..(121.6,-176.2)
				..controls(120.8,-176.2)and(119.8,-176.2)..(119.4,-176.1)--cycle
				(130.4,-160)..controls(129,-160)and(119.8,-160)..(110.1,-159.9)
				..controls(100.3,-159.8)and(88.5,-159.7)..(83.9,-159.6)
				..controls(77.8,-159.5)and(75.3,-159.5)..(74.5,-159.6)
				..controls(73.9,-159.7)and(72.5,-159.8)..(71.3,-159.8)
				..controls(69.6,-159.8)and(69,-159.8)..(68.5,-159.6)
				..controls(67.6,-159.3)and(66.8,-158.6)..(66.2,-157.6)
				..controls(66,-157.3)and(65.6,-156.5)..(65.3,-156.1)
				..controls(64.6,-155.1)and(64.5,-154.9)..(63.7,-152.9)
				..controls(63.2,-151.6)and(62.9,-150.9)..(61.8,-149.6)
				..controls(61.5,-149.2)and(61.1,-148.5)..(60.9,-147.9)
				..controls(60.7,-147.4)and(60.3,-146.6)..(60.1,-146.1)
				..controls(59.9,-145.5)and(59.5,-144.9)..(59,-144.3)
				..controls(58.5,-143.7)and(58.2,-143.2)..(57.9,-142.5)
				..controls(57.7,-141.9)and(57.3,-141.1)..(57,-140.7)
				..controls(56.6,-140.3)and(56.1,-139.3)..(55.8,-138.5)
				..controls(55.5,-137.8)and(55.1,-136.8)..(54.9,-136.5)
				..controls(54.7,-136.1)and(54.3,-135.2)..(53.9,-134.5)
				..controls(53.6,-133.8)and(53,-132.7)..(52.6,-132)
				..controls(51.6,-130.5)and(51.4,-130.2)..(51.1,-129)
				..controls(50.6,-127.5)and(50.6,-127.4)..(49.7,-126.3)
				..controls(49.2,-125.7)and(48.6,-124.8)..(48.3,-124.2)
				..controls(48,-123.6)and(47.5,-122.7)..(47.3,-122.2)
				..controls(47,-121.8)and(46.5,-120.7)..(46.1,-119.8)
				..controls(45.8,-118.9)and(45.2,-117.8)..(44.9,-117.2)
				..controls(44.2,-116.1)and(43.3,-114.4)..(42.3,-112.2)
				..controls(41.8,-111.3)and(41.3,-110.4)..(41.1,-110.1)
				..controls(40.5,-109.4)and(40.3,-109)..(40,-107.8)
				..controls(39.8,-107)and(39.6,-106.5)..(39.2,-106)
				..controls(38.9,-105.6)and(38.4,-104.7)..(38,-104)
				..controls(37.7,-103.4)and(37.2,-102.5)..(37,-102.1)
				..controls(36.7,-101.7)and(36.3,-100.8)..(36,-100.1)
				..controls(35.6,-98.9)and(35.3,-98.4)..(34.6,-97.6)
				..controls(34.2,-97.1)and(34.2,-96.3)..(34.6,-96)
				..controls(35,-95.7)and(35.5,-95.9)..(36.2,-96.7)
				..controls(36.5,-97)and(37.2,-97.5)..(37.6,-97.7)..controls(38,-98)and(38.6,-98.4)..(39.1,-98.9)
				..controls(39.6,-99.4)and(40.3,-100)..(40.6,-100.3)
				..controls(41.3,-100.8)and(42.4,-101.8)..(43.7,-102.9)
				..controls(44,-103.2)and(44.6,-103.7)..(45,-103.9)
				..controls(45.4,-104.2)and(45.9,-104.6)..(46.1,-104.8)
				..controls(47.1,-105.9)and(48.8,-107.1)..(50.5,-107.9)
				..controls(51.8,-108.6)and(52.3,-109)..(53.4,-110)
				..controls(54.8,-111.2)and(55.6,-111.6)..(57,-112.1)
				..controls(57.5,-112.3)and(58.5,-112.8)..(59.4,-113.3)
				..controls(60.2,-113.7)and(61.3,-114.3)..(61.9,-114.5)
				..controls(62.4,-114.7)and(63.3,-115.1)..(63.8,-115.3)
				..controls(64.4,-115.6)and(65.4,-115.9)..(66.2,-116.1)
				..controls(67,-116.3)and(68.4,-116.6)..(69.2,-116.9)
				..controls(72.3,-117.9)and(72.8,-118)..(78,-118.1)
				..controls(84,-118.3)and(86.1,-118.1)..(88.1,-117.1)
				..controls(89.4,-116.5)and(93,-115)..(93.7,-114.9)
				..controls(94.8,-114.6)and(95.4,-114.2)..(95.9,-113.4)
				..controls(96.2,-113)and(96.7,-112.5)..(97.1,-112.2)
				..controls(98,-111.6)and(101.3,-108.4)..(102.1,-107.2)
				..controls(102.4,-106.8)and(102.9,-106.2)..(103.2,-105.8)
				..controls(103.5,-105.5)and(104,-104.8)..(104.2,-104.2)
				..controls(104.5,-103.6)and(105,-102.8)..(105.3,-102.4)
				..controls(105.6,-102)and(106.1,-101.2)..(106.5,-100.6)
				..controls(107.1,-99.5)and(108.6,-97.5)..(109.5,-96.5)
				..controls(109.9,-96.1)and(110.1,-96)..(110.4,-96)
				..controls(111.1,-96)and(111.6,-96.5)..(111.8,-97.6)
				..controls(112,-98.6)and(112.3,-99.3)..(112.8,-99.9)
				..controls(113.7,-100.9)and(115.2,-103.1)..(115.7,-104.2)
				..controls(116.1,-105)and(116.5,-105.5)..(116.9,-106)
				..controls(117.3,-106.4)and(117.8,-106.9)..(118,-107.3)
				..controls(118.2,-107.6)and(118.6,-108)..(119,-108.3)
				..controls(119.6,-108.8)and(121.1,-110.3)..(121.8,-111.2)
				..controls(122.1,-111.7)and(122.5,-112)..(123.4,-112.4)
				..controls(124.8,-113.2)and(125.8,-113.8)..(127,-114.6)
				..controls(127.5,-114.9)and(128.1,-115.3)..(128.2,-115.4)
				..controls(128.4,-115.5)and(129.3,-115.8)..(130.2,-116)
				..controls(131.1,-116.2)and(132.5,-116.6)..(133.3,-116.9)
				..controls(135.7,-117.9)and(136.1,-117.9)..(142.5,-117.8)
				..controls(145.8,-117.8)and(148.4,-117.7)..(148.8,-117.6)
				..controls(149.2,-117.6)and(150.2,-117.2)..(151,-116.8)
				..controls(152,-116.4)and(152.8,-116)..(153.6,-115.9)
				..controls(154.5,-115.7)and(155.2,-115.4)..(156.3,-114.9)
				..controls(157.1,-114.5)and(158.2,-114)..(158.7,-113.8)
				..controls(159.1,-113.7)and(159.8,-113.3)..(160.2,-113)
				..controls(160.6,-112.7)and(161.3,-112.4)..(161.7,-112.2)
				..controls(162.7,-111.9)and(163.2,-111.7)..(164.1,-111.1)
				..controls(164.5,-110.8)and(165.3,-110.4)..(165.9,-110)
				..controls(166.5,-109.7)and(167.3,-109.2)..(167.6,-108.9)
				..controls(167.9,-108.6)and(168.6,-108.2)..(169,-108)
				..controls(169.4,-107.8)and(170.2,-107.3)..(170.7,-106.9)
				..controls(171.2,-106.5)and(172.1,-105.8)..(172.7,-105.4)
				..controls(173.2,-105)and(174.1,-104.4)..(174.5,-104)
				..controls(174.9,-103.7)and(175.7,-103.1)..(176.2,-102.6)
				..controls(176.7,-102.2)and(177.4,-101.5)..(177.9,-101.2)
				..controls(178.3,-100.8)and(178.8,-100.4)..(179,-100.2)
				..controls(179.2,-100)and(179.6,-99.6)..(180,-99.3)
				..controls(180.4,-99.1)and(180.9,-98.6)..(181.2,-98.4)
				..controls(181.5,-98.1)and(182.1,-97.7)..(182.5,-97.5)
				..controls(183.2,-97.2)and(183.5,-96.9)..(184.3,-95.9)
				..controls(184.6,-95.6)and(184.6,-95.6)..(184.8,-95.7)
				..controls(185,-95.8)and(185.1,-96)..(185.1,-96.2)
				..controls(185.1,-96.5)and(185,-96.7)..(184.6,-97.2)
				..controls(184.1,-97.8)and(184,-97.9)..(183.9,-98.8)
				..controls(183.7,-100.3)and(183.3,-101.3)..(182.9,-101.8)
				..controls(182.6,-102.1)and(182.3,-102.6)..(182.2,-102.9)
				..controls(182,-103.3)and(181.6,-104)..(181.2,-104.6)
				..controls(180.9,-105.2)and(180.5,-106.1)..(180.4,-106.6)
				..controls(179.9,-108)and(179.7,-108.5)..(179,-109.4)
				..controls(178.6,-109.8)and(178.2,-110.5)..(177.9,-111)
				..controls(177.7,-111.5)and(177.3,-112.3)..(177,-112.7)
				..controls(176.6,-113.2)and(176.4,-113.7)..(176.3,-114.2)
				..controls(176.2,-114.7)and(175.9,-115.2)..(175.4,-116)
				..controls(175,-116.6)and(174.4,-117.7)..(174,-118.5)
				..controls(173.7,-119.2)and(173.3,-120.1)..(173,-120.5)
				..controls(172.8,-120.9)and(172.4,-121.8)..(172.1,-122.5)
				..controls(171.8,-123.2)and(171.3,-124.1)..(171,-124.6)
				..controls(170.8,-125.1)and(170.3,-126)..(169.9,-126.6)
				..controls(169.6,-127.3)and(169.2,-128.1)..(169,-128.3)
				..controls(168.8,-128.6)and(168.5,-129.3)..(168.3,-129.8)
				..controls(168.1,-130.3)and(167.8,-131.1)..(167.5,-131.6)
				..controls(166.9,-132.5)and(164.2,-137.9)..(164.2,-138.1)
				..controls(164.2,-138.1)and(163.9,-138.8)..(163.4,-139.7)
				..controls(162.9,-140.5)and(162.4,-141.5)..(162.3,-141.9)
				..controls(161.9,-143)and(161.6,-143.4)..(160.9,-144.1)
				..controls(160.4,-144.6)and(160.3,-144.8)..(160.1,-145.4)
				..controls(160,-145.8)and(159.6,-146.8)..(159.2,-147.8)
				..controls(158.7,-148.8)and(158.1,-150)..(157.9,-150.6)
				..controls(157.6,-151.1)and(157.2,-151.9)..(156.9,-152.4)
				..controls(156.5,-152.8)and(156.1,-153.7)..(155.8,-154.4)
				..controls(155.6,-155)and(155.2,-155.9)..(154.9,-156.3)
				..controls(154.7,-156.7)and(154.4,-157.4)..(154.2,-157.9)
				..controls(154.1,-158.4)and(153.8,-159)..(153.5,-159.3)
				..controls(152.8,-160.1)and(153,-160.1)..(142.4,-160.1)
				..controls(137.2,-160.1)and(131.8,-160.1)..(130.4,-160)--cycle
				(141.5,-103.7)..controls(140.8,-103)and(139.8,-102.1)..(139.2,-101.6)
				..controls(138.5,-101.2)and(137.9,-100.6)..(137.8,-100.4)
				..controls(137,-99.5)and(134,-96.5)..(133.2,-96)
				..controls(132.7,-95.7)and(131.8,-94.9)..(131.1,-94.3)
				..controls(129.4,-92.7)and(127.9,-91.6)..(126.6,-91.1)
				..controls(125.2,-90.5)and(124.1,-90)..(122.9,-89.1)
				..controls(122.1,-88.6)and(121.3,-88.2)..(119.9,-87.7)
				..controls(118.8,-87.3)and(117.5,-86.8)..(116.9,-86.5)
				..controls(114.3,-85.5)and(113.7,-85.4)..(109.7,-85.5)
				..controls(105.3,-85.7)and(103.9,-85.9)..(101.6,-87)
				..controls(100.7,-87.4)and(99.4,-87.8)..(97.9,-88)
				..controls(97,-88.2)and(96.6,-88.3)..(95.8,-88.9)
				..controls(95.4,-89.2)and(94.7,-89.6)..(94.2,-89.8)
				..controls(92.9,-90.3)and(92.3,-90.7)..(90.9,-92.1)
				..controls(90.2,-92.7)and(89.4,-93.3)..(89,-93.5)..controls(87,-94.5)and(85.2,-96)..(83.6,-97.9)
				..controls(83.3,-98.2)and(82.5,-98.9)..(81.8,-99.4)
				..controls(80.7,-100.1)and(80.5,-100.3)..(80,-101.1)
				..controls(78.7,-103.1)and(77.2,-103.9)..(76.4,-103)
				..controls(75.7,-102.2)and(75.7,-103.2)..(75.7,-83.6)
				..controls(75.6,-73.6)and(75.6,-64.8)..(75.5,-64)
				..controls(75.5,-63.2)and(75.5,-61.4)..(75.6,-59.8)--(75.7,-57.1)--(76,-56.8)
				..controls(76.8,-56.2)and(78,-56.3)..(79.8,-57.2)
				..controls(80.4,-57.5)and(81.3,-57.8)..(81.8,-57.9)
				..controls(82.3,-58.1)and(83,-58.4)..(83.6,-58.8)..controls(84.3,-59.2)and(85,-59.5)..(86,-59.8)
				..controls(86.8,-60.1)and(88,-60.6)..(88.7,-60.9)
				..controls(89.3,-61.2)and(90.6,-61.7)..(91.4,-62)
				..controls(92.2,-62.3)and(93.7,-62.9)..(94.6,-63.3)
				..controls(95.6,-63.7)and(96.8,-64.1)..(97.4,-64.3)
				..controls(98,-64.5)and(99,-64.9)..(99.7,-65.2)
				..controls(100.4,-65.6)and(101.5,-66)..(102.3,-66.2)
				..controls(104.2,-66.8)and(105,-67.1)..(106.2,-67.8)
				..controls(107.9,-68.7)and(110.2,-69.4)..(111.3,-69.2)
				..controls(112.5,-69)and(115,-68.2)..(115.6,-67.8)
				..controls(117.3,-66.7)and(117.6,-66.6)..(119.5,-66.2)
				..controls(121.9,-65.8)and(122.1,-65.7)..(123,-65.1)
				..controls(124.1,-64.4)and(124.4,-64.3)..(126,-63.9)
				..controls(127,-63.7)and(127.7,-63.4)..(128.4,-63.1)
				..controls(129,-62.8)and(129.7,-62.5)..(130.6,-62.4)
				..controls(132.1,-62.1)and(134,-61.5)..(135.2,-61)
				..controls(135.7,-60.8)and(136.7,-60.5)..(137.3,-60.3)
				..controls(138.4,-60)and(139.3,-59.7)..(141.1,-59)
				..controls(141.6,-58.7)and(142,-58.6)..(142.5,-58.6)
				..controls(143.2,-58.6)and(143.3,-58.7)..(143.6,-59)
				..controls(144,-59.5)and(144.1,-59.9)..(144,-60.8)
				..controls(144,-61.2)and(144,-71.1)..(144,-82.8)--(143.9,-104.1)--(143.7,-104.5)
				..controls(143.5,-104.7)and(143.3,-104.8)..(143,-104.9)
				..controls(142.6,-104.9)and(142.5,-104.8)..(141.5,-103.7)--cycle
				(27.8,-91.4)..controls(27.1,-90.9)and(26.3,-90)..(26.1,-89.5)
				..controls(26,-89.2)and(25.8,-88.6)..(25.7,-87.9)
				..controls(25.4,-86.2)and(25.6,-84.8)..(26.5,-83.5)
				..controls(27.4,-82.3)and(28.3,-82)..(29.4,-82.3)
				..controls(31.1,-82.8)and(32.1,-84.9)..(32,-87.4)
				..controls(31.9,-89.3)and(31.2,-90.6)..(29.9,-91.4)
				..controls(29.1,-91.9)and(28.6,-91.9)..(27.8,-91.4)--cycle
				(190.1,-91.3)..controls(189.7,-91.2)and(189.3,-90.9)..(189.1,-90.7)
				..controls(188.3,-89.8)and(187.6,-87.7)..(187.6,-86.2)
				..controls(187.6,-85.5)and(187.7,-85.2)..(188.2,-84.4)
				..controls(189.1,-82.6)and(190.1,-81.8)..(191.3,-82.1)
				..controls(192,-82.3)and(193.5,-83.7)..(193.8,-84.4)
				..controls(194.1,-85.1)and(194.2,-87.1)..(194,-88.3)
				..controls(193.8,-89.8)and(193.3,-90.5)..(192.3,-91.3)
				..controls(191.7,-91.7)and(190.9,-91.7)..(190.1,-91.3)--cycle
				(173.7,-83.4)..controls(172.2,-82.5)and(172,-82.5)..(170.6,-82.4)
				..controls(169.7,-82.4)and(168.8,-82.4)..(168.3,-82.5)
				..controls(167.2,-82.7)and(166.9,-82.7)..(166.6,-82.2)
				..controls(166.3,-81.8)and(166.3,-81)..(166.8,-80)
				..controls(167.2,-78.9)and(167.4,-77.4)..(167.7,-74.1)
				..controls(167.8,-72.8)and(167.9,-71.8)..(167.9,-71.7)
				..controls(168,-71.7)and(168,-70.2)..(168,-68.4)..controls(168,-65.4)and(168,-64.9)..(167.7,-63)
				..controls(167.3,-60.5)and(167.1,-59.9)..(166.4,-58.4)
				..controls(166.1,-57.7)and(165.8,-56.9)..(165.7,-56.6)
				..controls(165.7,-56.2)and(165.4,-55.6)..(165.2,-55.2)
				..controls(164.6,-53.8)and(161.9,-50.7)..(161,-50.3)
				..controls(159.8,-49.7)and(157.9,-50)..(155.2,-51.2)
				..controls(154.7,-51.4)and(153.4,-51.8)..(152.2,-52.1)
				..controls(150.8,-52.4)and(149.4,-52.9)..(148.6,-53.2)
				..controls(147.8,-53.6)and(146.7,-54)..(146,-54.2)
				..controls(145.4,-54.4)and(144.5,-54.7)..(144,-54.9)
				..controls(143.6,-55.2)and(142.7,-55.5)..(141.8,-55.7)
				..controls(141,-55.9)and(139.6,-56.3)..(138.7,-56.6)
				..controls(137.8,-57)and(136.3,-57.4)..(135.3,-57.6)
				..controls(133.8,-57.9)and(133.2,-58.1)..(132.1,-58.7)
				..controls(131.1,-59.2)and(130.3,-59.4)..(129.4,-59.6)
				..controls(128.7,-59.8)and(127.8,-60)..(127.4,-60.1)
				..controls(127,-60.3)and(126,-60.5)..(125.1,-60.7)
				..controls(123.2,-61.1)and(121.7,-61.5)..(120.2,-62.1)
				..controls(119.7,-62.3)and(118.1,-62.8)..(116.7,-63.2)
				..controls(114.6,-63.9)and(114,-64.1)..(113,-64.7)
				..controls(111.9,-65.4)and(111.8,-65.4)..(111,-65.5)
				..controls(109.9,-65.5)and(109.3,-65.4)..(107.7,-64.7)
				..controls(107,-64.4)and(106,-64)..(105.3,-63.9)
				..controls(103.5,-63.5)and(102.2,-63)..(101.2,-62.5)
				..controls(100.1,-61.8)and(99.3,-61.5)..(97.8,-61.1)
				..controls(97.2,-61)and(96.2,-60.7)..(95.7,-60.4)
				..controls(95.1,-60.1)and(94.2,-59.8)..(93.5,-59.6)
				..controls(92.9,-59.4)and(91.7,-59)..(91,-58.7)..controls(90.2,-58.4)and(88.9,-58)..(88,-57.7)
				..controls(87.2,-57.5)and(86.1,-57.1)..(85.5,-56.8)
				..controls(85,-56.6)and(84.1,-56.2)..(83.5,-56)..controls(82.9,-55.7)and(81.8,-55.2)..(81,-54.9)
				..controls(80.1,-54.5)and(79.1,-54.2)..(78.7,-54.1)
				..controls(77.9,-53.8)and(77.3,-53.6)..(76.2,-52.9)
				..controls(75.7,-52.6)and(75.2,-52.4)..(74.1,-52.2)
				..controls(72.5,-51.9)and(72.4,-51.8)..(71.3,-51.1)
				..controls(69.9,-50.1)and(69.4,-49.9)..(67.9,-49.8)
				..controls(67,-49.7)and(66.6,-49.6)..(66.3,-49.4)
				..controls(66.1,-49.3)and(65.7,-49)..(65.4,-48.8)
				..controls(65.1,-48.7)and(64.2,-48.4)..(63.6,-48.2)
				..controls(62.9,-48)and(61.7,-47.6)..(60.9,-47.2)
				..controls(60.1,-46.8)and(58.9,-46.4)..(58.2,-46.2)
				..controls(57.6,-46)and(56.8,-45.7)..(56.5,-45.5)
				..controls(54.9,-44.6)and(54.6,-44.5)..(53.2,-44.2)
				..controls(51.7,-43.8)and(51,-43.6)..(49.9,-42.9)
				..controls(49.5,-42.6)and(48.9,-42.4)..(47.5,-42)
				..controls(45.7,-41.6)and(45.6,-41.5)..(44.3,-40.6)
				..controls(43.5,-40.1)and(42.8,-39.6)..(42.7,-39.4)
				..controls(42.4,-38.7)and(43,-38.1)..(44.5,-37.6)..controls(44.9,-37.5)and(46,-37)..(46.9,-36.6)
				..controls(47.9,-36.1)and(48.8,-35.8)..(49.6,-35.6)
				..controls(50.5,-35.4)and(51.3,-35.2)..(52.4,-34.6)
				..controls(53.3,-34.1)and(54.5,-33.7)..(54.9,-33.5)
				..controls(55.4,-33.4)and(56.4,-33)..(57.2,-32.6)
				..controls(58.6,-32)and(60.1,-31.5)..(61.6,-31.2)
				..controls(62.3,-31.1)and(62.7,-31)..(63.2,-30.7)
				..controls(64.1,-30)and(64.3,-29.9)..(65.3,-29.7)
				..controls(65.7,-29.6)and(66.5,-29.2)..(67.1,-28.9)
				..controls(67.6,-28.6)and(68.8,-28.2)..(69.6,-27.9)
				..controls(70.5,-27.7)and(71.6,-27.3)..(72.1,-27)
				..controls(72.6,-26.7)and(73.8,-26.2)..(74.8,-25.9)
				..controls(75.7,-25.5)and(76.7,-25.1)..(77.1,-25)
				..controls(77.4,-24.8)and(78.5,-24.4)..(79.6,-24.1)
				..controls(80.6,-23.8)and(82,-23.2)..(82.8,-22.9)
				..controls(83.5,-22.5)and(84.8,-22)..(85.8,-21.8)
				..controls(86.7,-21.5)and(87.8,-21.1)..(88.1,-20.9)
				..controls(88.4,-20.7)and(89.3,-20.3)..(90.1,-20.1)
				..controls(90.8,-19.8)and(91.9,-19.4)..(92.4,-19.1)
				..controls(93.5,-18.5)and(94.4,-18.2)..(95.4,-18.1)
				..controls(95.7,-18)and(96.4,-17.8)..(96.9,-17.5)
				..controls(97.5,-17.3)and(98.4,-16.9)..(99,-16.6)
				..controls(99.7,-16.3)and(100.9,-15.9)..(101.9,-15.6)
				..controls(102.9,-15.3)and(104.4,-14.7)..(105.3,-14.3)
				..controls(108.1,-13.2)and(108.5,-13.1)..(110.1,-13.1)
				..controls(111.4,-13.1)and(111.5,-13.1)..(112.4,-13.5)
				..controls(113.2,-14)and(115,-14.8)..(116.6,-15.4)
				..controls(116.9,-15.5)and(117.9,-15.8)..(118.8,-16.1)
				..controls(120.9,-16.6)and(121.5,-16.9)..(122.6,-17.9)
				..controls(123,-18.3)and(123.7,-18.8)..(123.9,-18.9)
				..controls(124.6,-19.4)and(126,-19.8)..(127.9,-20.2)
				..controls(129.6,-20.5)and(129.8,-20.6)..(131.1,-21.4)
				..controls(131.5,-21.7)and(132.4,-22.1)..(133.1,-22.3)
				..controls(136.1,-23.2)and(138.3,-24.1)..(140.2,-25)
				..controls(141,-25.5)and(141.5,-25.6)..(142.5,-25.8)
				..controls(145.1,-26.3)and(146,-26.6)..(149.4,-28.1)
				..controls(151.3,-28.9)and(153.5,-29.8)..(154.3,-30.1)
				..controls(155.1,-30.4)and(156.2,-30.8)..(156.6,-31)
				..controls(157.1,-31.2)and(158.2,-31.6)..(159.1,-31.8)
				..controls(160.6,-32.2)and(160.9,-32.3)..(161.7,-32.9)
				..controls(162.4,-33.3)and(162.9,-33.6)..(163.7,-33.8)
				..controls(164.3,-34)and(164.8,-34.1)..(164.9,-34.1)
				..controls(165,-34.1)and(165.6,-34.4)..(166.4,-34.8)
				..controls(167.3,-35.2)and(168.1,-35.4)..(169,-35.6)
				..controls(170,-35.8)and(170.5,-36)..(171.7,-36.6)
				..controls(172.4,-37)and(173.3,-37.4)..(173.6,-37.4)
				..controls(175.6,-37.7)and(176.2,-38)..(176.7,-39.1)
				..controls(176.9,-39.6)and(176.9,-39.7)..(176.7,-40.4)
				..controls(176.4,-41.6)and(175.8,-42.7)..(175.2,-43.4)
				..controls(174.5,-44.1)and(172.6,-45.5)..(171.7,-45.8)
				..controls(170.6,-46.2)and(169.1,-46)..(168.2,-45.4)
				..controls(167.9,-45.2)and(167.3,-44.8)..(166.9,-44.4)
				..controls(166.4,-44.1)and(165.8,-43.6)..(165.4,-43.2)
				..controls(165.1,-42.9)and(164.4,-42.3)..(163.9,-41.9)
				..controls(163.3,-41.5)and(162.7,-40.9)..(162.4,-40.5)
				..controls(161.4,-39.2)and(159.9,-38)..(158.8,-37.8)
				..controls(158.4,-37.8)and(158.3,-37.7)..(158.2,-37.2)
				..controls(158,-36.8)and(157.4,-36.1)..(157.1,-36.1)
				..controls(157,-36.1)and(156.7,-36.1)..(156.5,-36)
				..controls(156.1,-35.9)and(155.9,-35.7)..(155.7,-35.2)
				..controls(155.3,-34.6)and(155.1,-34.4)..(154,-34.1)
				..controls(153.6,-34)and(153.2,-33.7)..(152.9,-33.5)
				..controls(152.4,-32.9)and(151.6,-32.4)..(151.3,-32.4)
				..controls(150.9,-32.4)and(150.4,-32.6)..(150.4,-32.9)
				..controls(150.4,-33.2)and(150.8,-33.4)..(151.4,-33.6)
				..controls(152.1,-33.7)and(152.8,-34.1)..(153.5,-34.9)
				..controls(154.1,-35.5)and(154.3,-35.6)..(155,-35.8)
				..controls(155.7,-36)and(155.7,-36)..(156,-36.6)
				..controls(156.3,-37.3)and(156.7,-37.7)..(157.5,-37.8)
				..controls(158.1,-37.9)and(158.2,-38)..(158.4,-38.8)
				..controls(158.5,-39.4)and(158.6,-39.6)..(159.6,-40.4)
				..controls(160.2,-41)and(160.9,-41.6)..(161.3,-41.8)
				..controls(161.7,-42)and(161.9,-42.3)..(162.2,-42.8)
				..controls(162.4,-43.2)and(162.8,-43.8)..(163.1,-44.1)
				..controls(163.9,-44.8)and(165.1,-46.7)..(165.7,-48)
				..controls(166.2,-49.3)and(166.4,-49.7)..(167.2,-50.5)
				..controls(167.5,-50.8)and(167.9,-51.4)..(168.1,-51.8)
				..controls(168.3,-52.2)and(168.8,-53)..(169.1,-53.6)
				..controls(169.6,-54.3)and(169.8,-54.9)..(169.9,-55.6)
				..controls(170.4,-57.8)and(170.5,-58.1)..(171.2,-59.5)
				..controls(171.6,-60.3)and(172,-61.5)..(172.2,-62.1)
				..controls(172.3,-62.8)and(172.7,-63.7)..(172.9,-64.1)
				..controls(173.2,-64.6)and(173.4,-65.3)..(173.4,-65.6)
				..controls(173.8,-68.8)and(174.2,-70.6)..(174.9,-72.6)
				..controls(175.5,-74.5)and(175.9,-77.1)..(176.1,-80.2)
				..controls(176.3,-83.4)and(175.5,-84.4)..(173.7,-83.4)--cycle;
				\definecolor{ezs}{RGB}{247,179,26}\fill[ezs](-5.7,-156.6)--(-3.4,-157.5)--(61.9,12.7)--(59.7,13.6)--cycle;
				\definecolor{ezs}{RGB}{247,179,22}\fill[ezs](-3.8,-157.3)--(-1.5,-158.2)--(63.9,12)--(61.6,12.9)--cycle;
				\definecolor{ezs}{RGB}{248,179,19}\fill[ezs](-1.8,-158.1)--(0.4,-158.9)--(65.8,11.2)--(63.5,12.1)--cycle;
				\definecolor{ezs}{RGB}{248,178,16}\fill[ezs](0.1,-158.8)--(2.4,-159.7)--(67.7,10.5)--(65.4,11.4)--cycle;
				\definecolor{ezs}{RGB}{249,178,15}\fill[ezs](2,-159.5)--(4.3,-160.4)--(69.6,9.8)--(67.3,10.7)--cycle;
				\definecolor{ezs}{RGB}{251,177,17}\fill[ezs](3.9,-160.3)--(6.2,-161.1)--(71.5,9)--(69.2,9.9)--cycle;
				\definecolor{ezs}{RGB}{252,179,18}\fill[ezs](5.8,-161)--(8.1,-161.9)--(73.4,8.3)--(71.1,9.2)--cycle;
				\definecolor{ezs}{RGB}{252,182,17}\fill[ezs](7.7,-161.7)--(10,-162.6)--(75.3,7.6)--(73,8.5)--cycle;
				\definecolor{ezs}{RGB}{252,185,17}\fill[ezs](9.6,-162.5)--(11.9,-163.4)--(77.2,6.8)--(75,7.7)--cycle;
				\definecolor{ezs}{RGB}{251,189,18}\fill[ezs](11.5,-163.2)--(13.8,-164.1)--(79.2,6.1)--(76.9,7)--cycle;
				\definecolor{ezs}{RGB}{251,191,22}\fill[ezs](13.4,-163.9)--(15.7,-164.8)--(81.1,5.4)--(78.8,6.2)--cycle;
				\definecolor{ezs}{RGB}{252,193,25}\fill[ezs](15.4,-164.7)--(17.6,-165.6)--(83,4.6)--(80.7,5.5)--cycle;
				\definecolor{ezs}{RGB}{253,195,27}\fill[ezs](17.3,-165.4)--(19.6,-166.3)--(84.9,3.9)--(82.6,4.8)--cycle;
				\definecolor{ezs}{RGB}{253,196,31}\fill[ezs](19.2,-166.1)--(21.5,-167)--(86.8,3.2)--(84.5,4)--cycle;
				\definecolor{ezs}{RGB}{254,199,35}\fill[ezs](21.1,-166.9)--(23.4,-167.8)--(88.7,2.4)--(86.4,3.3)--cycle;
				\definecolor{ezs}{RGB}{253,201,38}\fill[ezs](23,-167.6)--(25.3,-168.5)--(90.6,1.7)--(88.3,2.6)--cycle;
				\definecolor{ezs}{RGB}{253,203,40}\fill[ezs](24.9,-168.3)--(27.2,-169.2)--(92.5,1)--(90.3,1.8)--cycle;
				\definecolor{ezs}{RGB}{253,205,42}\fill[ezs](26.8,-169.1)--(29.1,-170)--(94.4,0.2)--(92.2,1.1)--cycle;
				\definecolor{ezs}{RGB}{253,207,45}\fill[ezs](28.7,-169.8)--(31,-170.7)--(96.4,-0.5)--(94.1,0.4)--cycle;
				\definecolor{ezs}{RGB}{253,210,49}\fill[ezs](30.7,-170.5)--(32.9,-171.4)--(98.3,-1.2)--(96,-0.4)--cycle;
				\definecolor{ezs}{RGB}{253,212,54}\fill[ezs](32.6,-171.3)--(34.9,-172.2)--(100.2,-2)--(97.9,-1.1)--cycle;
				\definecolor{ezs}{RGB}{253,214,59}\fill[ezs](34.5,-172)--(36.8,-172.9)--(102.1,-2.7)--(99.8,-1.8)--cycle;
				\definecolor{ezs}{RGB}{253,216,63}\fill[ezs](36.4,-172.7)--(38.7,-173.6)--(104,-3.4)--(101.7,-2.6)--cycle;
				\definecolor{ezs}{RGB}{253,218,65}\fill[ezs](38.3,-173.5)--(40.6,-174.4)--(105.9,-4.2)--(103.6,-3.3)--cycle;
				\definecolor{ezs}{RGB}{254,220,66}\fill[ezs](40.2,-174.2)--(42.5,-175.1)--(107.8,-4.9)--(105.5,-4)--cycle;
				\definecolor{ezs}{RGB}{254,222,68}\fill[ezs](42.1,-175)--(44.4,-175.8)--(109.7,-5.6)--(107.5,-4.8)--cycle;
				\definecolor{ezs}{RGB}{254,224,71}\fill[ezs](44,-175.7)--(46.3,-176.6)--(111.7,-6.4)--(109.4,-5.5)--cycle;
				\definecolor{ezs}{RGB}{254,226,75}\fill[ezs](46,-176.4)--(48.2,-177.3)--(113.6,-7.1)--(111.3,-6.2)--cycle;
				\definecolor{ezs}{RGB}{254,228,79}\fill[ezs](47.9,-177.2)--(50.1,-178)--(115.5,-7.8)--(113.2,-7)--cycle;
				\definecolor{ezs}{RGB}{254,230,82}\fill[ezs](49.8,-177.9)--(52.1,-178.8)--(117.4,-8.6)--(115.1,-7.7)--cycle;
				\definecolor{ezs}{RGB}{254,232,85}\fill[ezs](51.7,-178.6)--(54,-179.5)--(119.3,-9.3)--(117,-8.4)--cycle;
				\definecolor{ezs}{RGB}{254,234,88}\fill[ezs](53.6,-179.4)--(55.9,-180.2)--(121.2,-10)--(118.9,-9.2)--cycle;
				\definecolor{ezs}{RGB}{254,237,91}\fill[ezs](55.5,-180.1)--(57.8,-181)--(123.1,-10.8)--(120.8,-9.9)--cycle;
				\definecolor{ezs}{RGB}{254,239,94}\fill[ezs](57.4,-180.8)--(59.7,-181.7)--(125,-11.5)--(122.8,-10.6)--cycle;
				\definecolor{ezs}{RGB}{254,241,97}\fill[ezs](59.3,-181.6)--(61.6,-182.4)--(127,-12.2)--(124.7,-11.4)--cycle;
				\definecolor{ezs}{RGB}{254,243,101}\fill[ezs](61.2,-182.3)--(63.5,-183.2)--(128.9,-13)--(126.6,-12.1)--cycle;
				\definecolor{ezs}{RGB}{254,245,103}\fill[ezs](63.2,-183)--(65.4,-183.9)--(130.8,-13.7)--(128.5,-12.8)--cycle;
				\definecolor{ezs}{RGB}{254,245,105}\fill[ezs](65.1,-183.8)--(67.4,-184.6)--(132.7,-14.4)--(130.4,-13.6)--cycle;
				\definecolor{ezs}{RGB}{254,245,105}\fill[ezs](67,-184.5)--(69.3,-185.4)--(134.6,-15.2)--(132.3,-14.3)--cycle;
				\definecolor{ezs}{RGB}{254,244,103}\fill[ezs](68.9,-185.2)--(71.2,-186.1)--(136.5,-15.9)--(134.2,-15)--cycle;
				\definecolor{ezs}{RGB}{253,243,101}\fill[ezs](70.8,-186)--(73.1,-186.8)--(138.4,-16.6)--(136.1,-15.8)--cycle;
				\definecolor{ezs}{RGB}{253,241,98}\fill[ezs](72.7,-186.7)--(75,-187.6)--(140.3,-17.4)--(138,-16.5)--cycle;
				\definecolor{ezs}{RGB}{253,239,95}\fill[ezs](74.6,-187.4)--(76.9,-188.3)--(142.2,-18.1)--(140,-17.2)--cycle;
				\definecolor{ezs}{RGB}{253,236,92}\fill[ezs](76.5,-188.2)--(78.8,-189)--(144.2,-18.8)--(141.9,-18)--cycle;
				\definecolor{ezs}{RGB}{253,234,88}\fill[ezs](78.5,-188.9)--(80.7,-189.8)--(146.1,-19.6)--(143.8,-18.7)--cycle;
				\definecolor{ezs}{RGB}{253,232,85}\fill[ezs](80.4,-189.6)--(82.7,-190.5)--(148,-20.3)--(145.7,-19.4)--cycle;
				\definecolor{ezs}{RGB}{253,230,82}\fill[ezs](82.3,-190.4)--(84.6,-191.2)--(149.9,-21)--(147.6,-20.2)--cycle;
				\definecolor{ezs}{RGB}{253,228,78}\fill[ezs](84.2,-191.1)--(86.5,-192)--(151.8,-21.8)--(149.5,-20.9)--cycle;
				\definecolor{ezs}{RGB}{253,226,74}\fill[ezs](86.1,-191.8)--(88.4,-192.7)--(153.7,-22.5)--(151.4,-21.6)--cycle;
				\definecolor{ezs}{RGB}{253,224,70}\fill[ezs](88,-192.6)--(90.3,-193.4)--(155.6,-23.3)--(153.3,-22.4)--cycle;
				\definecolor{ezs}{RGB}{253,222,67}\fill[ezs](89.9,-193.3)--(92.2,-194.2)--(157.5,-24)--(155.3,-23.1)--cycle;
				\definecolor{ezs}{RGB}{253,220,65}\fill[ezs](91.8,-194)--(94.1,-194.9)--(159.5,-24.7)--(157.2,-23.8)--cycle;
				\definecolor{ezs}{RGB}{253,217,63}\fill[ezs](93.7,-194.8)--(96,-195.6)--(161.4,-25.5)--(159.1,-24.6)--cycle;
				\definecolor{ezs}{RGB}{253,215,60}\fill[ezs](95.7,-195.5)--(97.9,-196.4)--(163.3,-26.2)--(161,-25.3)--cycle;
				\definecolor{ezs}{RGB}{253,213,56}\fill[ezs](97.6,-196.2)--(99.9,-197.1)--(165.2,-26.9)--(162.9,-26)--cycle;
				\definecolor{ezs}{RGB}{252,211,53}\fill[ezs](99.5,-197)--(101.8,-197.8)--(167.1,-27.7)--(164.8,-26.8)--cycle;
				\definecolor{ezs}{RGB}{252,209,49}\fill[ezs](101.4,-197.7)--(103.7,-198.6)--(169,-28.4)--(166.7,-27.5)--cycle;
				\definecolor{ezs}{RGB}{252,207,45}\fill[ezs](103.3,-198.4)--(105.6,-199.3)--(170.9,-29.1)--(168.6,-28.2)--cycle;
				\definecolor{ezs}{RGB}{253,205,42}\fill[ezs](105.2,-199.2)--(107.5,-200)--(172.8,-29.9)--(170.6,-29)--cycle;
				\definecolor{ezs}{RGB}{253,202,40}\fill[ezs](107.1,-199.9)--(109.4,-200.8)--(174.7,-30.6)--(172.5,-29.7)--cycle;
				\definecolor{ezs}{RGB}{253,200,38}\fill[ezs](109,-200.6)--(111.3,-201.5)--(176.7,-31.3)--(174.4,-30.4)--cycle;
				\definecolor{ezs}{RGB}{253,198,35}\fill[ezs](111,-201.4)--(113.2,-202.2)--(178.6,-32.1)--(176.3,-31.2)--cycle;
				\definecolor{ezs}{RGB}{253,196,31}\fill[ezs](112.9,-202.1)--(115.2,-203)--(180.5,-32.8)--(178.2,-31.9)--cycle;
				\definecolor{ezs}{RGB}{253,194,28}\fill[ezs](114.8,-202.8)--(117.1,-203.7)--(182.4,-33.5)--(180.1,-32.6)--cycle;
				\definecolor{ezs}{RGB}{253,192,26}\fill[ezs](116.7,-203.6)--(119,-204.5)--(184.3,-34.3)--(182,-33.4)--cycle;
				\definecolor{ezs}{RGB}{253,189,24}\fill[ezs](118.6,-204.3)--(120.9,-205.2)--(186.2,-35)--(183.9,-34.1)--cycle;
				\definecolor{ezs}{RGB}{253,186,20}\fill[ezs](120.5,-205)--(122.8,-205.9)--(188.1,-35.7)--(185.8,-34.8)--cycle;
				\definecolor{ezs}{RGB}{253,184,17}\fill[ezs](122.4,-205.8)--(124.7,-206.7)--(190,-36.5)--(187.8,-35.6)--cycle;
				\definecolor{ezs}{RGB}{252,183,15}\fill[ezs](124.3,-206.5)--(126.6,-207.4)--(192,-37.2)--(189.7,-36.3)--cycle;
				\definecolor{ezs}{RGB}{252,181,14}\fill[ezs](126.3,-207.2)--(128.5,-208.1)--(193.9,-37.9)--(191.6,-37.1)--cycle;
				\definecolor{ezs}{RGB}{252,179,13}\fill[ezs](128.2,-208)--(130.4,-208.9)--(195.8,-38.7)--(193.5,-37.8)--cycle;
				\definecolor{ezs}{RGB}{252,178,12}\fill[ezs](130.1,-208.7)--(132.4,-209.6)--(197.7,-39.4)--(195.4,-38.5)--cycle;
				\definecolor{ezs}{RGB}{251,178,10}\fill[ezs](132,-209.4)--(134.3,-210.3)--(199.6,-40.1)--(197.3,-39.3)--cycle;
				\definecolor{ezs}{RGB}{251,178,10}\fill[ezs](133.9,-210.2)--(136.2,-211.1)--(201.5,-40.9)--(199.2,-40)--cycle;
				\definecolor{ezs}{RGB}{250,177,11}\fill[ezs](135.8,-210.9)--(138.1,-211.8)--(203.4,-41.6)--(201.1,-40.7)--cycle;
				\definecolor{ezs}{RGB}{250,178,11}\fill[ezs](137.7,-211.6)--(140,-212.5)--(205.3,-42.3)--(203.1,-41.5)--cycle;
				\definecolor{ezs}{RGB}{248,180,10}\fill[ezs](139.6,-212.4)--(141.9,-213.3)--(207.3,-43.1)--(205,-42.2)--cycle;
				\definecolor{ezs}{RGB}{247,181,9}\fill[ezs](141.5,-213.1)--(143.8,-214)--(209.2,-43.8)--(206.9,-42.9)--cycle;
				\definecolor{ezs}{RGB}{246,182,9}\fill[ezs](143.5,-213.8)--(145.7,-214.7)--(211.1,-44.5)--(208.8,-43.7)--cycle;
				\definecolor{ezs}{RGB}{246,182,10}\fill[ezs](145.4,-214.6)--(147.7,-215.5)--(213,-45.3)--(210.7,-44.4)--cycle;
			\end{scope}
			% chữ EDUZY (Times đậm, làm dày nét bằng 5 lớp lệch nhẹ; căn đúng khung chữ logo gốc)
			\foreach \dx/\dy in {-0.7/0,0.7/0,0/-0.3,0/0.3,0/0}
			\node[anchor=center,inner sep=0pt,outer sep=0pt,text=ezText,align=center] at (109.6+\dx,-193.6+\dy)
			{\resizebox{1.01cm}{0.22cm}{\fontencoding{T1}\fontfamily{ptm}\fontseries{b}\selectfont EDUZY}};
		\end{tikzpicture}%
	}%
}
%====================== THIẾT LẬP THÂN ĐỀ ======================
%--- Môi trường tự luận {bt}: ex_test 3.3.4 chưa khai báo sẵn nên tạo ở đây
%    (ex_test đã tự móc \fixshowans và \listenumerate cho bt)
\makeatletter
\@ifundefined{bt}{\newtheorem{bt}{\color{gold1}\namebt}}{}
\makeatother
\renewcommand{\sectionmark}[1]{\markright{\thesection\ #1}}
%--- Khung tiêu đề (logo + tối đa 4 dòng; dòng 3, 4 để trống {} thì bỏ qua)
%  Đề thi:    \khungtieude{THI THỬ THPT QUỐC GIA NĂM 2026}{Môn Toán -- Đề số 21}
%                         {(Đề thi gồm: ...)}{Thời gian làm bài: 90 phút (không kể thời gian phát đề).}
%  Chuyên đề: \khungtieude{CHUYÊN ĐỀ 3 -- DÃY SỐ VÀ CÔNG THỨC TRUY HỒI}
%                         {Bồi dưỡng học sinh giỏi Toán THPT}{...}{}
\newcommand{\khungtieude}[4]{%
	\begin{tcolorbox}[colframe=gold1,colback=gold5!45,boxrule=0.5mm]
		\begin{minipage}[c]{3cm}
			\centerline{\eduzylogo[2.4cm]}
		\end{minipage}\hspace{0.5cm}
		\begin{minipage}[c]{12.5cm}
			\vskip 0.15 cm
			\centerline{\ind{#1}}
			\vskip 0.15 cm
			\centerline{\textbf{#2}}
			\if\relax\detokenize{#3}\relax\else
				\vskip 0.15 cm
				\centerline{\textit{#3}}
			\fi
			\if\relax\detokenize{#4}\relax\else
				\vskip 0.15 cm
				\centerline{\textit{#4}}
			\fi
			\vskip 0.15 cm
		\end{minipage}
	\end{tcolorbox}
}
%====================== BỔ SUNG CHO BÀI HỌC, CHUYÊN ĐỀ ======================
\usepackage[shortlabels]{enumitem}
\usepackage{tabularx}
%--- Danh sách gọn: \begin{enumerate}[gon] ... \end{enumerate}
\SetEnumitemKey{gon}{leftmargin=*,itemsep=1pt,parsep=0pt,topsep=3pt}
%--- Khung lý thuyết: \begin{khung}{} ... \end{khung}; tiêu đề chỉ dùng khi thật ngắn ("Mẹo", "Lỗi thường gặp")
\newtcolorbox{khung}[1]{breakable,enhanced,colframe=gold1,colback=gold5!45,
	boxrule=0.4mm,fonttitle=\bfseries\sffamily,coltitle=white,colbacktitle=gold1,
	attach boxed title to top left={xshift=6mm,yshift=-3mm},
	boxed title style={arc=1mm,boxrule=0.3mm},title={#1}}
%--- Tiêu đề mục nhỏ trong một dạng: \muc{Dấu hiệu}, \muc{Phương pháp}, \muc{Bài tự luyện}
\makeatletter
\newcommand{\muc}[1]{\par\medskip\noindent\indam{#1}\par\nobreak\smallskip\@afterindentfalse\@afterheading}
\makeatother
%--- Ô công thức cần nhớ (căn giữa, viền vàng, vừa khít nội dung)
%     \begin{center}\congthuc{$Q(n)=(n-1)!$}\end{center}
\newtcbox{\congthuc}{on line,colframe=gold1,colback=gold5!45,boxrule=0.5pt,arc=2pt,
	left=8pt,right=8pt,top=4pt,bottom=4pt}
%--- Công tắc lời giải bài tự luyện (ví dụ mẫu luôn dùng \loigiai)
%     \loigiaiTL[{đáp số}]{lời giải đầy đủ}
%     \TLfull -> in lời giải đầy đủ;  \TLdapso -> chỉ in một dòng "Đáp số";  \TLan -> ẩn hoàn toàn
%     (\CheDo là dethi thì mọi lời giải và đáp số đều ẩn)
\makeatletter
\newif\ifTLhien   \TLhientrue
\newif\ifTLdapso  \TLdapsofalse
\newcommand{\TLfull}{\TLhientrue\TLdapsofalse}
\newcommand{\TLdapso}{\TLhientrue\TLdapsotrue}
\newcommand{\TLan}{\TLhienfalse}
\newcommand{\loigiaiTL}[2][]{%
	\ifTLhien
		\ifTLdapso
			\def\@tlarg{#1}%
			\ifx\@tlarg\@empty\else
				\ifdefstring{\kindOPTN}{dethi}{}{\par\noindent{\bfseries\color{gold1}Đáp số:} #1\par}%
			\fi
		\else
			\loigiai{#2}%
		\fi
	\fi}
\makeatother
\TLdapso% mặc định: chỉ in đáp số bài tự luyện
%====================== BẢNG ĐÚNG/SAI (\choiceTFt) ======================
%  Hàng tiêu đề nền nâu đồng chữ trắng (Phát biểu | Đ | S), các hàng tô xen kẽ kem / trắng,
%  đường kẻ vàng kim, nhãn a) b) c) d) màu nâu đồng. Bản có đáp án: cột được chọn có vòng tròn
%  nâu đồng với dấu tích trắng (câu đúng tích ở cột Đúng, câu sai tích ở cột Sai).
%  Tiêu đề ghi đủ "Đúng | Sai": \OPTN{viettat=0} sau \input (cột Đúng/Sai rộng ra, phát biểu hẹp lại).
%  Dùng bảng gốc của ex_test: \def\BangDSgoc{1} trước \input{preamble-eduzy-v2}.
\providecommand{\BangDSgoc}{0}
\ifnum\BangDSgoc=0
\usepackage{colortbl,pifont}
\colorlet{ezDSke}{gold2!60!gold1}% đường kẻ bảng (nét 0,4pt như ex_test để không đổi chiều cao bảng)
\colorlet{ezDSle}{gold5!60}% hàng a), c)
\colorlet{ezDSchan}{white}% hàng b), d)
%--- dấu chọn ở cột Đúng/Sai
%    (bề rộng 0, căn giữa ô: không làm đổi bề rộng cột của ex_test, không báo tràn lề)
\newcommand{\EZdauchon}{\makebox[0pt]{\tikz[baseline=-1ex]\node[circle,fill=gold1,inner sep=0pt,
	minimum size=1.2em,text=white,font=\small]{\ding{51}};}}
\def\tickT{\EZdauchon}\def\tickF{\EZdauchon}
\AtBeginDocument{\sbox{\ccTFTab}{\EZdauchon}}
%--- nhãn a) b) c) d) màu nâu đồng (ở bản lời giải, nhãn câu đúng vẫn đóng khung như ex_test)
\renewcommand{\FalseTF}{\stepcounter{dapan}\textbf{\color{gold1}\DapAnTF\sepTF}\ignorespaces}
%--- trước khi dựng bảng: kẻ màu vàng kim (bề rộng cột giữ nguyên như ex_test); sau bảng trả lại màu đen
\newcount\EZhangTF
\newcommand{\EZbatdauTF}{%
	\global\EZhangTF=0\relax
	\arrayrulecolor{ezDSke}}
\newcommand{\EZketthucTF}{\arrayrulecolor{black}}
%--- hàng tiêu đề nền nâu đồng, chữ trắng
\makeatletter
\newcommand{\EZo@TF}[1]{\centering\bfseries\color{white}#1}
\newcommand{\EZtieudeTF}{%
	\ifcase\cotDS
		\ifnum\cotSTT=1\relax
			\def\rowHeader{\hline\rowcolor{gold1}\EZo@TF{TT}&\EZo@TF{\phatbieu}\tabularnewline}%
		\else
			\def\rowHeader{\hline\rowcolor{gold1}\EZo@TF{\phatbieu}\tabularnewline}%
		\fi
	\or
		\ifnum\cotSTT=1\relax
			\def\rowHeader{\hline\rowcolor{gold1}\EZo@TF{TT}&\EZo@TF{\phatbieu}&\EZo@TF{\Dung}\tabularnewline}%
		\else
			\def\rowHeader{\hline\rowcolor{gold1}\EZo@TF{\phatbieu}&\EZo@TF{\Dung}\tabularnewline}%
		\fi
	\or
		\ifnum\cotSTT=1\relax
			\def\rowHeader{\hline\rowcolor{gold1}\EZo@TF{TT}&\EZo@TF{\phatbieu}&\EZo@TF{\Dung}&\EZo@TF{\Sai}\tabularnewline}%
		\else
			\def\rowHeader{\hline\rowcolor{gold1}\EZo@TF{\phatbieu}&\EZo@TF{\Dung}&\EZo@TF{\Sai}\tabularnewline}%
		\fi
	\fi}
%--- tô xen kẽ từng hàng (gọi ngay sau \hline đầu mỗi hàng)
\newcommand{\EZmauhangTF}{%
	\noalign{\global\advance\EZhangTF\@ne}%
	\ifodd\EZhangTF\rowcolor{ezDSle}\else\rowcolor{ezDSchan}\fi}
%--- gắn vào \choiceTFt / \writeChoiceTFt của ex_test
\xpretocmd{\choiceTFt}{\EZbatdauTF}{}{\typeout{[preamble-eduzy-v2] Không gắn được bảng Đúng/Sai (đầu bảng)}}
\xpatchcmd{\choiceTFt}{\ifnum\dongPB=0\def\rowHeader{}\fi}{\EZtieudeTF\ifnum\dongPB=0\def\rowHeader{}\fi}
	{}{\typeout{[preamble-eduzy-v2] Không gắn được bảng Đúng/Sai (hàng tiêu đề)}}
\xpatchcmd{\writeChoiceTFt}{\hline\xdef\dotEX{\sdotEX}}{\hline\EZmauhangTF\xdef\dotEX{\sdotEX}}
	{}{\typeout{[preamble-eduzy-v2] Không gắn được bảng Đúng/Sai (tô hàng)}}
\xpatchcmd{\writeChoiceTFt}{{\end{longtable}}}{{\end{longtable}\EZketthucTF}}
	{}{\typeout{[preamble-eduzy-v2] Không gắn được bảng Đúng/Sai (cuối longtable)}}
\xpatchcmd{\writeChoiceTFt}{{\end{xtabular}\end{center}\medskip\noindent}}{{\end{xtabular}\end{center}\EZketthucTF\medskip\noindent}}
	{}{\typeout{[preamble-eduzy-v2] Không gắn được bảng Đúng/Sai (cuối xtabular)}}
\fi
%=====================================================================
\renewcommand{\parallel}{\mathrel{/\!\!/}}

%     %=====================================================================
% preamble-vatli-hinh.tex -- BỘ MACRO HÌNH VẬT LÍ (dùng chung cho mọi đề Vật lí)
% Gọi ngay sau preamble chính:
%     \input{preamble-eduzy-v2}
%     \input{preamble-vatli-hinh}
% Tên file bắt đầu bằng "preamble" nên dong-goi.sh tự chép vào zip thành phẩm;
% tools/so_hinh.py đọc file này khi dựng ảnh so hình.
% Lưu ý màu: preamble-eduzy-v2 đổi blue, cyan, teal sang tông vàng đồng (gold1, gold2).
% Cần màu xanh thật thì dùng VLxanh (hoặc khai báo \definecolor cục bộ).
% Rút từ đề mẫu VatLi2026-CaoBang-SoGD-Lan1 (thầy duyệt 06/11/2026).
%=====================================================================

%--- Mạch điện: circuitikz kiểu châu Âu (điện trở hình chữ nhật...)
\usepackage[european]{circuitikz}

%--- Màu dùng chung
\definecolor{VLdo}{RGB}{214,72,60}
\definecolor{VLxanh}{RGB}{60,100,200}
\definecolor{VLdong}{RGB}{200,128,70}
\definecolor{VLgo}{RGB}{186,132,84}
\definecolor{VLla}{RGB}{150,200,150}

%--- Mũi tên: VLmui (mũi tên ở đầu đường), VLgiua=<vị trí> (mũi tên giữa đường, mặc định 0,55)
\tikzset{
	VLmui/.style={-{Stealth[length=4.5pt,width=3.5pt]}},
	VLgiua/.style={postaction={decorate},decoration={markings,mark=at position #1 with {\arrow{Stealth[length=6.5pt,width=5pt]}}}},
	VLgiua/.default=0.55
}
% cọc nối dây của điện kế: \VLcoc{x}{y}{màu}
\newcommand{\VLcoc}[3]{%
	\begin{scope}[shift={(#1,#2)}]
		\shade[left color=#3!45!black,right color=#3!55!black,middle color=#3!80!white] (-0.17,0) rectangle (0.17,0.28);
		\shade[left color=#3!55!black,right color=#3!65!black,middle color=#3!60!white] (-0.14,0.28) rectangle (0.14,0.5);
		\fill[#3!55!black] (0,0.5) ellipse (0.14 and 0.045);
		\fill[black!85] (0,0.5) ellipse (0.055 and 0.02);
	\end{scope}}
% phích cắm có rãnh: \VLphich{x}{y}{màu}
\newcommand{\VLphich}[3]{%
	\begin{scope}[shift={(#1,#2)}]
		\shade[left color=#3!50!black,right color=#3!60!black,middle color=#3!55!white] (-0.12,0) rectangle (0.12,0.62);
		\draw[#3!45!black,very thin] foreach \y in {0.1,0.18,...,0.5}{(-0.12,\y)--(0.12,\y)};
		\shade[left color=#3!50!black,right color=#3!60!black,middle color=#3!55!white] (-0.08,0.62) rectangle (0.08,0.85);
		\fill[#3!45!black] (0,0.85) ellipse (0.08 and 0.03);
	\end{scope}}
% điện kế nhìn chính diện, gốc ở góc trái dưới, rộng 4, cao 4: \VLdienke{1: thang số | 2: thang -, +}
\newcommand{\VLdienke}[1]{%
	\fill[black!82,rounded corners=2pt] (0,0) rectangle (4,1.35);
	\fill[black!60] (0.05,1.22) rectangle (3.95,1.35);
	\fill[VLla,draw=black!60,rounded corners=1pt] (0.28,0.25) rectangle (3.72,1.02);
	\fill[gray!8,draw=gray!65,rounded corners=1.5pt] (0.08,1.35) rectangle (3.92,4.05);
	\fill[white,draw=gray!45] (0.28,1.5) rectangle (3.72,3.85);
	\fill[gray!18] (1.25,1.52) rectangle (2.75,2.05);
	\draw[gray!40,very thin] foreach \x in {1.3,1.4,...,2.7}{(\x,1.55)--(\x,2.0)};
	\draw[black!80] (2,1.35) ++(127:1.95) arc[start angle=127,end angle=53,radius=1.95];
	\foreach \g in {53,55.4667,...,127.1}{\draw[black!80,very thin] (2,1.35)++(\g:1.95)--++(\g:0.12);}
	\foreach \g in {53,65.3333,...,127.1}{\draw[black!85] (2,1.35)++(\g:1.95)--++(\g:0.22);}
	\ifnum#1=1
		\foreach \g/\s in {127/15,114.6667/10,102.3333/5,90/0,77.6667/5,65.3333/10,53/15}{\node[font=\fontsize{5}{5}\selectfont,inner sep=0] at ($(2,1.35)+(\g:2.37)$) {\s};}
		\node[font=\scriptsize\bfseries] at (2,2.75) {G};
	\else
		\node[font=\scriptsize] at ($(2,1.35)+(127:1.7)$) {$-$};
		\node[font=\scriptsize] at ($(2,1.35)+(53:1.7)$) {$+$};
	\fi
	\draw[VLdo,semithick] (2,1.8)--(2,3.42);
	\fill[black!85] (2,1.8) circle (0.13);
	\fill[gray!55] (2,1.8) circle (0.045);
	\draw[white,line width=1.2pt,opacity=0.7] (0.2,3.95)--(0.2,1.6);
}
% nam châm thẳng 3D: \VLnamcham{dài}{cao}{màu trái}{màu phải}{chữ trái}{chữ phải}
\newcommand{\VLnamcham}[6]{%
	\fill[#3!55!white] (0,#2)--(0.5*#1,#2)--(0.5*#1+0.22,#2+0.13)--(0.22,#2+0.13)--cycle;
	\fill[#4!55!white] (0.5*#1,#2)--(#1,#2)--(#1+0.22,#2+0.13)--(0.5*#1+0.22,#2+0.13)--cycle;
	\fill[#4!70!black] (#1,0)--(#1+0.22,0.13)--(#1+0.22,#2+0.13)--(#1,#2)--cycle;
	\shade[top color=#3!75!white,bottom color=#3!95!black] (0,0) rectangle (0.5*#1,#2);
	\shade[top color=#4!75!white,bottom color=#4!95!black] (0.5*#1,0) rectangle (#1,#2);
	\node[white,font=\scriptsize\bfseries] at (0.25*#1,0.5*#2) {#5};
	\node[white,font=\scriptsize\bfseries] at (0.75*#1,0.5*#2) {#6};
}
% lưới chấm (từ trường hướng ra) / lưới dấu nhân (hướng vào)
\newcommand{\VLcham}[2]{\fill (#1,#2) circle (1.6pt);}
\newcommand{\VLnhan}[2]{\draw[semithick] (#1,#2)++(-0.06,-0.06)--++(0.12,0.12) (#1,#2)++(-0.06,0.06)--++(0.12,-0.12);}
% thanh dẫn hình trụ nằm ngang / thẳng đứng (tâm, nửa dài)
\newcommand{\VLthanhngang}[3]{%
	\shade[top color=gray!60,bottom color=gray!85,middle color=gray!35,draw=black!85] (#1-#3,#2-0.1) rectangle (#1+#3,#2+0.1);
	\filldraw[fill=gray!35,draw=black!80] (#1+#3,#2) ellipse (0.04 and 0.1);}
\newcommand{\VLthanhdung}[3]{%
	\shade[left color=gray!85,right color=gray!60,middle color=gray!35,draw=black!85] (#1-0.1,#2-#3) rectangle (#1+0.1,#2+#3);
	\filldraw[fill=gray!35,draw=black!80] (#1,#2+#3) ellipse (0.1 and 0.035);}

% Tên file bắt đầu bằng "preamble" nên dong-goi.sh tự chép vào zip thành phẩm;
% tools/so_hinh.py đọc file này khi dựng ảnh so hình.
% Lưu ý màu: preamble-eduzy-v2 đổi blue, cyan, teal sang tông vàng đồng (gold1, gold2).
% Cần màu xanh thật thì dùng VLxanh (hoặc khai báo \definecolor cục bộ).
% Rút từ đề mẫu VatLi2026-CaoBang-SoGD-Lan1 (thầy duyệt 06/11/2026).
%=====================================================================

%--- Mạch điện: circuitikz kiểu châu Âu (điện trở hình chữ nhật...)
\usepackage[european]{circuitikz}

%--- Màu dùng chung
\definecolor{VLdo}{RGB}{214,72,60}
\definecolor{VLxanh}{RGB}{60,100,200}
\definecolor{VLdong}{RGB}{200,128,70}
\definecolor{VLgo}{RGB}{186,132,84}
\definecolor{VLla}{RGB}{150,200,150}

%--- Mũi tên: VLmui (mũi tên ở đầu đường), VLgiua=<vị trí> (mũi tên giữa đường, mặc định 0,55)
\tikzset{
	VLmui/.style={-{Stealth[length=4.5pt,width=3.5pt]}},
	VLgiua/.style={postaction={decorate},decoration={markings,mark=at position #1 with {\arrow{Stealth[length=6.5pt,width=5pt]}}}},
	VLgiua/.default=0.55
}
% cọc nối dây của điện kế: \VLcoc{x}{y}{màu}
\newcommand{\VLcoc}[3]{%
	\begin{scope}[shift={(#1,#2)}]
		\shade[left color=#3!45!black,right color=#3!55!black,middle color=#3!80!white] (-0.17,0) rectangle (0.17,0.28);
		\shade[left color=#3!55!black,right color=#3!65!black,middle color=#3!60!white] (-0.14,0.28) rectangle (0.14,0.5);
		\fill[#3!55!black] (0,0.5) ellipse (0.14 and 0.045);
		\fill[black!85] (0,0.5) ellipse (0.055 and 0.02);
	\end{scope}}
% phích cắm có rãnh: \VLphich{x}{y}{màu}
\newcommand{\VLphich}[3]{%
	\begin{scope}[shift={(#1,#2)}]
		\shade[left color=#3!50!black,right color=#3!60!black,middle color=#3!55!white] (-0.12,0) rectangle (0.12,0.62);
		\draw[#3!45!black,very thin] foreach \y in {0.1,0.18,...,0.5}{(-0.12,\y)--(0.12,\y)};
		\shade[left color=#3!50!black,right color=#3!60!black,middle color=#3!55!white] (-0.08,0.62) rectangle (0.08,0.85);
		\fill[#3!45!black] (0,0.85) ellipse (0.08 and 0.03);
	\end{scope}}
% điện kế nhìn chính diện, gốc ở góc trái dưới, rộng 4, cao 4: \VLdienke{1: thang số | 2: thang -, +}
\newcommand{\VLdienke}[1]{%
	\fill[black!82,rounded corners=2pt] (0,0) rectangle (4,1.35);
	\fill[black!60] (0.05,1.22) rectangle (3.95,1.35);
	\fill[VLla,draw=black!60,rounded corners=1pt] (0.28,0.25) rectangle (3.72,1.02);
	\fill[gray!8,draw=gray!65,rounded corners=1.5pt] (0.08,1.35) rectangle (3.92,4.05);
	\fill[white,draw=gray!45] (0.28,1.5) rectangle (3.72,3.85);
	\fill[gray!18] (1.25,1.52) rectangle (2.75,2.05);
	\draw[gray!40,very thin] foreach \x in {1.3,1.4,...,2.7}{(\x,1.55)--(\x,2.0)};
	\draw[black!80] (2,1.35) ++(127:1.95) arc[start angle=127,end angle=53,radius=1.95];
	\foreach \g in {53,55.4667,...,127.1}{\draw[black!80,very thin] (2,1.35)++(\g:1.95)--++(\g:0.12);}
	\foreach \g in {53,65.3333,...,127.1}{\draw[black!85] (2,1.35)++(\g:1.95)--++(\g:0.22);}
	\ifnum#1=1
		\foreach \g/\s in {127/15,114.6667/10,102.3333/5,90/0,77.6667/5,65.3333/10,53/15}{\node[font=\fontsize{5}{5}\selectfont,inner sep=0] at ($(2,1.35)+(\g:2.37)$) {\s};}
		\node[font=\scriptsize\bfseries] at (2,2.75) {G};
	\else
		\node[font=\scriptsize] at ($(2,1.35)+(127:1.7)$) {$-$};
		\node[font=\scriptsize] at ($(2,1.35)+(53:1.7)$) {$+$};
	\fi
	\draw[VLdo,semithick] (2,1.8)--(2,3.42);
	\fill[black!85] (2,1.8) circle (0.13);
	\fill[gray!55] (2,1.8) circle (0.045);
	\draw[white,line width=1.2pt,opacity=0.7] (0.2,3.95)--(0.2,1.6);
}
% nam châm thẳng 3D: \VLnamcham{dài}{cao}{màu trái}{màu phải}{chữ trái}{chữ phải}
\newcommand{\VLnamcham}[6]{%
	\fill[#3!55!white] (0,#2)--(0.5*#1,#2)--(0.5*#1+0.22,#2+0.13)--(0.22,#2+0.13)--cycle;
	\fill[#4!55!white] (0.5*#1,#2)--(#1,#2)--(#1+0.22,#2+0.13)--(0.5*#1+0.22,#2+0.13)--cycle;
	\fill[#4!70!black] (#1,0)--(#1+0.22,0.13)--(#1+0.22,#2+0.13)--(#1,#2)--cycle;
	\shade[top color=#3!75!white,bottom color=#3!95!black] (0,0) rectangle (0.5*#1,#2);
	\shade[top color=#4!75!white,bottom color=#4!95!black] (0.5*#1,0) rectangle (#1,#2);
	\node[white,font=\scriptsize\bfseries] at (0.25*#1,0.5*#2) {#5};
	\node[white,font=\scriptsize\bfseries] at (0.75*#1,0.5*#2) {#6};
}
% lưới chấm (từ trường hướng ra) / lưới dấu nhân (hướng vào)
\newcommand{\VLcham}[2]{\fill (#1,#2) circle (1.6pt);}
\newcommand{\VLnhan}[2]{\draw[semithick] (#1,#2)++(-0.06,-0.06)--++(0.12,0.12) (#1,#2)++(-0.06,0.06)--++(0.12,-0.12);}
% thanh dẫn hình trụ nằm ngang / thẳng đứng (tâm, nửa dài)
\newcommand{\VLthanhngang}[3]{%
	\shade[top color=gray!60,bottom color=gray!85,middle color=gray!35,draw=black!85] (#1-#3,#2-0.1) rectangle (#1+#3,#2+0.1);
	\filldraw[fill=gray!35,draw=black!80] (#1+#3,#2) ellipse (0.04 and 0.1);}
\newcommand{\VLthanhdung}[3]{%
	\shade[left color=gray!85,right color=gray!60,middle color=gray!35,draw=black!85] (#1-0.1,#2-#3) rectangle (#1+0.1,#2+#3);
	\filldraw[fill=gray!35,draw=black!80] (#1,#2+#3) ellipse (0.1 and 0.035);}

% Tên file bắt đầu bằng "preamble" nên dong-goi.sh tự chép vào zip thành phẩm;
% tools/so_hinh.py đọc file này khi dựng ảnh so hình.
% Lưu ý màu: preamble-eduzy-v2 đổi blue, cyan, teal sang tông vàng đồng (gold1, gold2).
% Cần màu xanh thật thì dùng VLxanh (hoặc khai báo \definecolor cục bộ).
% Rút từ đề mẫu VatLi2026-CaoBang-SoGD-Lan1 (thầy duyệt 06/11/2026).
%=====================================================================

%--- Mạch điện: circuitikz kiểu châu Âu (điện trở hình chữ nhật...)
\usepackage[european]{circuitikz}

%--- Màu dùng chung
\definecolor{VLdo}{RGB}{214,72,60}
\definecolor{VLxanh}{RGB}{60,100,200}
\definecolor{VLdong}{RGB}{200,128,70}
\definecolor{VLgo}{RGB}{186,132,84}
\definecolor{VLla}{RGB}{150,200,150}

%--- Mũi tên: VLmui (mũi tên ở đầu đường), VLgiua=<vị trí> (mũi tên giữa đường, mặc định 0,55)
\tikzset{
	VLmui/.style={-{Stealth[length=4.5pt,width=3.5pt]}},
	VLgiua/.style={postaction={decorate},decoration={markings,mark=at position #1 with {\arrow{Stealth[length=6.5pt,width=5pt]}}}},
	VLgiua/.default=0.55
}
% cọc nối dây của điện kế: \VLcoc{x}{y}{màu}
\newcommand{\VLcoc}[3]{%
	\begin{scope}[shift={(#1,#2)}]
		\shade[left color=#3!45!black,right color=#3!55!black,middle color=#3!80!white] (-0.17,0) rectangle (0.17,0.28);
		\shade[left color=#3!55!black,right color=#3!65!black,middle color=#3!60!white] (-0.14,0.28) rectangle (0.14,0.5);
		\fill[#3!55!black] (0,0.5) ellipse (0.14 and 0.045);
		\fill[black!85] (0,0.5) ellipse (0.055 and 0.02);
	\end{scope}}
% phích cắm có rãnh: \VLphich{x}{y}{màu}
\newcommand{\VLphich}[3]{%
	\begin{scope}[shift={(#1,#2)}]
		\shade[left color=#3!50!black,right color=#3!60!black,middle color=#3!55!white] (-0.12,0) rectangle (0.12,0.62);
		\draw[#3!45!black,very thin] foreach \y in {0.1,0.18,...,0.5}{(-0.12,\y)--(0.12,\y)};
		\shade[left color=#3!50!black,right color=#3!60!black,middle color=#3!55!white] (-0.08,0.62) rectangle (0.08,0.85);
		\fill[#3!45!black] (0,0.85) ellipse (0.08 and 0.03);
	\end{scope}}
% điện kế nhìn chính diện, gốc ở góc trái dưới, rộng 4, cao 4: \VLdienke{1: thang số | 2: thang -, +}
\newcommand{\VLdienke}[1]{%
	\fill[black!82,rounded corners=2pt] (0,0) rectangle (4,1.35);
	\fill[black!60] (0.05,1.22) rectangle (3.95,1.35);
	\fill[VLla,draw=black!60,rounded corners=1pt] (0.28,0.25) rectangle (3.72,1.02);
	\fill[gray!8,draw=gray!65,rounded corners=1.5pt] (0.08,1.35) rectangle (3.92,4.05);
	\fill[white,draw=gray!45] (0.28,1.5) rectangle (3.72,3.85);
	\fill[gray!18] (1.25,1.52) rectangle (2.75,2.05);
	\draw[gray!40,very thin] foreach \x in {1.3,1.4,...,2.7}{(\x,1.55)--(\x,2.0)};
	\draw[black!80] (2,1.35) ++(127:1.95) arc[start angle=127,end angle=53,radius=1.95];
	\foreach \g in {53,55.4667,...,127.1}{\draw[black!80,very thin] (2,1.35)++(\g:1.95)--++(\g:0.12);}
	\foreach \g in {53,65.3333,...,127.1}{\draw[black!85] (2,1.35)++(\g:1.95)--++(\g:0.22);}
	\ifnum#1=1
		\foreach \g/\s in {127/15,114.6667/10,102.3333/5,90/0,77.6667/5,65.3333/10,53/15}{\node[font=\fontsize{5}{5}\selectfont,inner sep=0] at ($(2,1.35)+(\g:2.37)$) {\s};}
		\node[font=\scriptsize\bfseries] at (2,2.75) {G};
	\else
		\node[font=\scriptsize] at ($(2,1.35)+(127:1.7)$) {$-$};
		\node[font=\scriptsize] at ($(2,1.35)+(53:1.7)$) {$+$};
	\fi
	\draw[VLdo,semithick] (2,1.8)--(2,3.42);
	\fill[black!85] (2,1.8) circle (0.13);
	\fill[gray!55] (2,1.8) circle (0.045);
	\draw[white,line width=1.2pt,opacity=0.7] (0.2,3.95)--(0.2,1.6);
}
% nam châm thẳng 3D: \VLnamcham{dài}{cao}{màu trái}{màu phải}{chữ trái}{chữ phải}
\newcommand{\VLnamcham}[6]{%
	\fill[#3!55!white] (0,#2)--(0.5*#1,#2)--(0.5*#1+0.22,#2+0.13)--(0.22,#2+0.13)--cycle;
	\fill[#4!55!white] (0.5*#1,#2)--(#1,#2)--(#1+0.22,#2+0.13)--(0.5*#1+0.22,#2+0.13)--cycle;
	\fill[#4!70!black] (#1,0)--(#1+0.22,0.13)--(#1+0.22,#2+0.13)--(#1,#2)--cycle;
	\shade[top color=#3!75!white,bottom color=#3!95!black] (0,0) rectangle (0.5*#1,#2);
	\shade[top color=#4!75!white,bottom color=#4!95!black] (0.5*#1,0) rectangle (#1,#2);
	\node[white,font=\scriptsize\bfseries] at (0.25*#1,0.5*#2) {#5};
	\node[white,font=\scriptsize\bfseries] at (0.75*#1,0.5*#2) {#6};
}
% lưới chấm (từ trường hướng ra) / lưới dấu nhân (hướng vào)
\newcommand{\VLcham}[2]{\fill (#1,#2) circle (1.6pt);}
\newcommand{\VLnhan}[2]{\draw[semithick] (#1,#2)++(-0.06,-0.06)--++(0.12,0.12) (#1,#2)++(-0.06,0.06)--++(0.12,-0.12);}
% thanh dẫn hình trụ nằm ngang / thẳng đứng (tâm, nửa dài)
\newcommand{\VLthanhngang}[3]{%
	\shade[top color=gray!60,bottom color=gray!85,middle color=gray!35,draw=black!85] (#1-#3,#2-0.1) rectangle (#1+#3,#2+0.1);
	\filldraw[fill=gray!35,draw=black!80] (#1+#3,#2) ellipse (0.04 and 0.1);}
\newcommand{\VLthanhdung}[3]{%
	\shade[left color=gray!85,right color=gray!60,middle color=gray!35,draw=black!85] (#1-0.1,#2-#3) rectangle (#1+0.1,#2+#3);
	\filldraw[fill=gray!35,draw=black!80] (#1,#2+#3) ellipse (0.1 and 0.035);}

\begin{document}
	\khungtieude{ĐỀ THI THỬ THPT NGUYỄN THỊ MINH KHAI -- LẦN 1}{Môn Vật lí}{(Đề thi gồm: 18 câu trắc nghiệm, 04 câu đúng-sai, 06 câu trả lời ngắn)}{Thời gian làm bài: 50 phút (không kể thời gian phát đề).}

	\NeedspaceCau{17.7\baselineskip}
	\noindent
	\subsection*{\indam{PHẦN I. Câu trắc nghiệm nhiều phương án lựa chọn. Thí sinh trả lời từ câu 1 đến câu 18. Mỗi câu hỏi thí sinh chỉ chọn một phương án.}}
	\setcounter{ex}{0}
	\Opensolutionfile{ans}[ans/VatLi2026-HaNoi-NguyenThiMinhKhai-Lan1-1]

	\begin{ex}%[VL12-C2-TH]
		Một khối khí lí tưởng thực hiện quá trình biến đổi từ trạng thái (1) sang trạng thái (2) được biểu diễn như hình bên. Quá trình đó là
		\begin{center}
			\begin{tikzpicture}[line join=round,line cap=round,font=\footnotesize]
				\draw[VLmui,thick] (0,0)--(4,0) node[above,inner sep=3pt] {$T$ (K)};
				\draw[VLmui,thick] (0,0)--(0,3.4) node[right,inner sep=3pt] {$p$ (atm)};
				\node[below left,inner sep=2pt] at (0,0) {$O$};
				\draw[dashed,thin,gray] (0,1.12)--(1.2,1.12)--(1.2,0) (0,2.9)--(3.1,2.9)--(3.1,0);
				\draw[dashed,thin] (0,0)--(1.2,1.12);
				\draw[very thick,VLdo,VLgiua=0.55] (1.2,1.12)--(3.1,2.9);
				\foreach \p in {(1.2,1.12),(3.1,2.9)}{\filldraw[fill=white,draw=black,thick] \p circle (2pt);}
				\node[right,inner sep=3pt] at (1.2,1.05) {(1)};
				\node[above,inner sep=3pt] at (3.1,2.9) {(2)};
				\node[left,inner sep=2pt] at (0,1.12) {$p_1$};
				\node[left,inner sep=2pt] at (0,2.9) {$p_2$};
				\node[below,inner sep=2pt] at (1.2,0) {$T_1$};
				\node[below,inner sep=2pt] at (3.1,0) {$T_2$};
			\end{tikzpicture}
		\end{center}
		\choice
		{làm lạnh đẳng tích}
		{làm lạnh đẳng áp}
		{đun nóng đẳng áp}
		{\True đun nóng đẳng tích}
		\loigiai{
			Đường biểu diễn trong hệ $pOT$ có đường kéo dài đi qua gốc $O$ nên thể tích không đổi; $p$ và $T$ cùng tăng: đun nóng đẳng tích.
		}
	\end{ex}

	\NeedspaceCau{4.8\baselineskip}
	\begin{ex}%[VL12-C1-NB]
		Khi hai vật được đặt tiếp xúc với nhau, năng lượng nhiệt chỉ truyền từ vật này sang vật kia nếu chúng có
		\choice
		{khối lượng khác nhau}
		{thế năng khác nhau}
		{\True nhiệt độ khác nhau}
		{thể tích khác nhau}
		\loigiai{
			Năng lượng nhiệt tự truyền từ vật có nhiệt độ cao sang vật có nhiệt độ thấp, nên chỉ xảy ra khi hai vật có nhiệt độ khác nhau.
		}
	\end{ex}

	\NeedspaceCau{13.1\baselineskip}
	\begin{ex}%[VL12-C3-TH]
		Cho một ống dây đặt cố định và một nam châm thử được treo như hình bên. Khi đóng công tắc K, nam châm thử sẽ
		% giuhinh
		\begin{center}
			\begin{tikzpicture}[x=0.0175cm,y=-0.0175cm,line join=round,line cap=round,font=\footnotesize,
				dayvang/.style={draw=VLdong!72!black,line width=2.3pt,postaction={draw=yellow!55!VLdong,line width=1.1pt}},
				daysau/.style={draw=VLdong!45!black,line width=2.1pt,postaction={draw=VLdong!80!black,line width=0.9pt}}]
				% dây bên phải đi lên phía SAU lõi, nối vào nửa sau của vòng cuối
				\draw[dayvang,rounded corners=6pt,line cap=butt] (150.6,196)--(224,196)--(224,140);
				\draw[daysau] (224,142)--(224,118) .. controls (224,100) and (218,90) .. (209,90);
				% nửa sau các vòng dây (khuất sau lõi): từ đỉnh đi xuống, lệch sang phải
				\foreach \j in {0,...,13}{\pgfmathsetmacro{\m}{29+11*\j}
					\draw[daysau] (\m+15,90) .. controls (\m+22,90) and (\m+25,105) .. (\m+25,120) .. controls (\m+25,136) and (\m+25,149) .. (\m+21,149);}
				% lõi sắt
				\shade[top color=gray!85,bottom color=gray!90,middle color=gray!4] (20,100)--(232,100)--(232,138)--(20,138) arc[start angle=90,end angle=270,x radius=4,y radius=19];
				\shade[top color=gray!10,bottom color=gray!55,opacity=0.5] (17,124) rectangle (232,138);
				\draw[gray!45!black,line width=0.4pt] (232,100)--(20,100) arc[start angle=270,end angle=90,x radius=-4,y radius=19]--(232,138);
				\shade[left color=gray!55,right color=gray!18] (232,119) ellipse (7 and 19);
				\draw[gray!30!black,line width=0.6pt] (232,119) ellipse (7 and 19);
				\fill[white,opacity=0.6] (234,113) ellipse (2.2 and 9);
				% nửa trước các vòng dây: từ đáy đi lên, lệch sang phải
				\foreach \i in {0,...,14}{\pgfmathsetmacro{\m}{40+11*\i}
					\draw[dayvang] (\m+10,149) .. controls (\m+3,149) and (\m,135) .. (\m,120) .. controls (\m,104) and (\m+6,90) .. (\m+15,90);}
				% dây bên trái đi lên phía TRƯỚC lõi rồi vòng qua đỉnh
				\draw[dayvang,rounded corners=6pt] (62,196)--(30,196)--(30,150);
				\draw[dayvang] (30,152)--(30,108) .. controls (30,97) and (36,90) .. (44,90);
				% khóa K
				\draw[VLla!45!black,line width=1.1pt] (63,194)--(94,170);
				\draw[dayvang,line cap=butt] (100,196)--(142.6,196);
				\foreach \x in {62,100}{\shade[ball color=VLdong!75!black] (\x,196) circle (4.6);
					\draw[VLdong!40!black,line width=0.3pt] (\x,196) circle (4.6);}
				% nguồn điện: bản dài (cực dương) bên trái, bản ngắn (cực âm) bên phải
				\draw[black!80,line width=0.9pt,line cap=butt] (143,171)--(143,220);
				\draw[black!80,line width=1.8pt,line cap=butt] (150,181)--(150,211);
				\node[inner sep=0pt] at (79,222) {K};
				\node[inner sep=0pt] at (129,214) {$+$};
				\node[inner sep=0pt] at (163,214) {$-$};
				% nam châm thử treo
				\draw[line width=0.9pt] (326,25)--(376,25);
				\foreach \x in {330,340,...,370}{\draw[line width=0.8pt] (\x,25)--++(7,-9);}
				\draw[line width=0.5pt] (350,25)--(350,108);
				\fill[VLxanh!35] (301,121)--(350,107)--(350,122)--cycle;
				\fill[VLxanh!58] (301,121)--(350,122)--(350,136)--cycle;
				\fill[VLdo] (400,121)--(350,107)--(350,122)--cycle;
				\fill[VLdo!72!black] (400,121)--(350,122)--(350,136)--cycle;
				\draw[black!65,line width=0.35pt] (301,121)--(350,107)--(400,121)--(350,136)--cycle (350,107)--(350,136);
				\node[inner sep=0pt] at (312,100) {S};
				\node[inner sep=0pt] at (390,100) {N};
			\end{tikzpicture}
		\end{center}
		\choice
		{bị hút rồi bị đẩy bởi ống dây}
		{bị hút bởi ống dây}
		{vẫn đứng yên}
		{\True bị đẩy bởi ống dây}
		\loigiai{
			Dòng điện đi từ cực dương qua K vào đầu trái ống dây, ở mặt trước các vòng dòng điện đi từ dưới lên. Theo quy tắc nắm tay phải, từ trường trong lòng ống dây hướng sang trái: đầu phải ống dây là cực Nam, đối diện cực S của nam châm thử nên đẩy nam châm thử.
		}
	\end{ex}

	\NeedspaceCau{6.6\baselineskip}
	\begin{ex}%[VL12-C3-NB]
		Đặt một đoạn dây dẫn dài $L$ mang dòng điện cường độ $I$ vào trong một từ trường đều có độ lớn cảm ứng từ $B$ sao cho đoạn dây hợp với các đường sức từ một góc $\alpha$, khi đó lực từ tác dụng lên đoạn dây đo được là $F$. Độ lớn cảm ứng từ tại vị trí đặt dòng điện được tính bằng công thức nào sau đây?
		\choice
		{$B=\dfrac{F}{IL\tan\alpha}$}
		{$B=\dfrac{F}{IL\cos\alpha}$}
		{$B=FIL\sin\alpha$}
		{\True $B=\dfrac{F}{IL\sin\alpha}$}
		\loigiai{
			$F=BIL\sin\alpha\Rightarrow B=\dfrac{F}{IL\sin\alpha}$.
		}
	\end{ex}

	\NeedspaceCau{4.5\baselineskip}
	\begin{ex}%[VL12-C3-TH]
		Một đoạn dây dẫn thẳng, dài $10$~cm mang dòng điện $5$~A đặt trong từ trường đều cảm ứng từ $B=0{,}08$~T. Biết đoạn dây dẫn đặt song song với vectơ cảm ứng từ. Độ lớn lực từ tác dụng lên đoạn dây dẫn là
		\choice
		{$0{,}04$~N}
		{$0{,}06$~N}
		{\True $0$~N}
		{$0{,}08$~N}
		\loigiai{
			Dây song song với $\vv{B}$ nên $\alpha=0$: $F=BIL\sin 0^\circ=0$.
		}
	\end{ex}

	\NeedspaceCau{4.5\baselineskip}
	\begin{ex}%[VL11-C3-NB]
		Chỉ ra công thức đúng của định luật Coulomb xác định độ lớn của lực tương tác giữa hai điện tích điểm $q_1$, $q_2$ đặt cách nhau $r$ trong chân không.
		\choice
		{$F=k\dfrac{|q_1q_2|}{r}$}
		{$F=k\dfrac{q_1q_2}{r}$}
		{$F=\dfrac{q_1q_2}{kr}$}
		{\True $F=k\dfrac{|q_1q_2|}{r^2}$}
		\loigiai{
			Định luật Coulomb: $F=k\dfrac{|q_1q_2|}{r^2}$.
		}
	\end{ex}

	\NeedspaceCau{2.4\baselineskip}
	\begin{ex}%[VL11-C3-NB]
		Kí hiệu đơn vị của cường độ điện trường là
		\choice
		{\True V/m}
		{C}
		{T}
		{N$\cdot$m$^2$/C$^2$}
		\loigiai{
			Cường độ điện trường có đơn vị V/m (hay N/C).
		}
	\end{ex}

	\NeedspaceCau{13.4\baselineskip}
	\begin{ex}%[VL12-C2-NB]
		Đun nóng khối khí trong một bình kín. Các phân tử khí
		% giuhinh
		\begin{center}
			\begin{tikzpicture}[x=0.0125cm,y=-0.0125cm,line join=round,line cap=round,font=\footnotesize]
				% bóng đổ, tấm đế (nhìn hơi chếch từ trên)
				\fill[black,opacity=0.1] (120,359) ellipse (104 and 4);
				\shade[top color=gray!30,bottom color=gray!55] (22,352)--(36,341)--(218,341)--(212,352)--cycle;
				\shade[top color=gray!65,bottom color=gray!85] (22,352)--(212,352)--(212,357)--(22,357)--cycle;
				\fill[gray!75] (212,352)--(218,341)--(218,346)--(212,357)--cycle;
				\draw[gray!35!black,line width=0.35pt] (22,352)--(36,341)--(218,341)--(218,346)--(212,357)--(22,357)--cycle (22,352)--(212,352)--(218,341) (212,352)--(212,357);
				% thanh trụ
				\fill[gray!70!black] (41,346.5) ellipse (4.5 and 1.6);
				\shade[left color=gray!80,right color=gray!60,middle color=gray!8] (38.5,12) rectangle (43.5,346);
				\draw[gray!40!black,line width=0.3pt] (38.5,346)--(38.5,12) arc[start angle=180,end angle=0,x radius=2.5,y radius=-2]--(43.5,346);
				% thanh ngang, khớp nối
				\shade[top color=gray!60,bottom color=gray!85,middle color=gray!10] (12,88.5) rectangle (112,93.5);
				\draw[gray!40!black,line width=0.3pt] (12,88.5) rectangle (112,93.5);
				\shade[top color=black!55,bottom color=black!90,rounded corners=0.8pt] (34,84) rectangle (48,98);
				\fill[gray!25] (41,91) circle (2.4);
				\draw[black!80,line width=0.5pt] (39.5,89.5)--(42.5,92.5);
				\shade[top color=black!60,bottom color=black!88] (35,82)--(47,82)--(46,84)--(36,84)--cycle;
				% ống nghiệm (bình kín)
				\shade[left color=black!88,right color=black!80,middle color=gray!22] (134,76)--(134,201) .. controls (134,211) and (141,216) .. (149,216) .. controls (157,216) and (164,211) .. (164,201)--(164,76)--cycle;
				\fill[white,opacity=0.55,rounded corners=0.6pt] (145,100) rectangle (148.5,200);
				\fill[white,opacity=0.18,rounded corners=0.6pt] (157,100) rectangle (159,198);
				\draw[black!85,line width=0.4pt] (134,76)--(134,201) .. controls (134,211) and (141,216) .. (149,216) .. controls (157,216) and (164,211) .. (164,201)--(164,76);
				\fill[gray!55,opacity=0.6] (149,212) ellipse (9 and 2.2);
				% nút đậy
				\shade[left color=black!85,right color=black!75,middle color=black!50] (131,48)--(167,48)--(164,77)--(134,77)--cycle;
				\shade[top color=black!45,bottom color=black!62] (149,48) ellipse (18 and 3.2);
				\draw[black!90,line width=0.3pt] (131,48)--(134,77)--(164,77)--(167,48);
				% vòng kẹp ống nghiệm, đầu kẹp có vít
				\shade[left color=black!90,right color=black!85,middle color=black!55] (132.5,87) rectangle (165.5,96);
				\shade[top color=black!60,bottom color=black!92,rounded corners=1pt] (107,84)--(133,86)--(133,97)--(107,99)--cycle;
				\draw[gray!15,line width=0.9pt] (113,88)--(118,96);
				\fill[gray!20] (121.5,91.5) circle (1.6);
				% đèn cồn: thân dưới có gân, vai đèn, cổ đèn, bấc, quai
				\draw[black!80,line width=1pt] (198,318) .. controls (214,308) and (217,336) .. (199,338);
				\draw[gray!40,line width=0.3pt] (199,316.5) .. controls (212,309) and (214,334) .. (200,336.5);
				\shade[left color=gray!75,right color=gray!70,middle color=gray!5] (107,322)--(107,343) .. controls (120,349) and (185,349) .. (198,343)--(198,322)--cycle;
				\foreach \x in {111,115.5,...,196}{\draw[gray!50!black,line width=0.5pt] (\x,327.5)--(\x,{347.3-0.0019*(\x-152.5)*(\x-152.5)});}
				\draw[gray!35!black,line width=0.35pt] (107,322)--(107,343) .. controls (120,349) and (185,349) .. (198,343)--(198,322);
				\shade[top color=gray!45,bottom color=black!88] (104,323) .. controls (106,311) and (128,303) .. (152.5,303) .. controls (177,303) and (199,311) .. (201,323) .. controls (185,329) and (120,329) .. (104,323)--cycle;
				\fill[white,opacity=0.35] (135,309) ellipse (17 and 3.5);
				\draw[black!80,line width=0.35pt] (104,323) .. controls (106,311) and (128,303) .. (152.5,303) .. controls (177,303) and (199,311) .. (201,323) .. controls (185,329) and (120,329) .. (104,323)--cycle;
				\shade[left color=gray!55,right color=gray!45,middle color=white] (134,280)--(134,301) .. controls (142,305) and (168,305) .. (176,301)--(176,280)--cycle;
				\foreach \x in {141,148,162,169}{\draw[gray!55,line width=0.3pt] (\x,283)--(\x,302.5);}
				\shade[top color=gray!15,bottom color=gray!40] (155,280) ellipse (21 and 3.5);
				\draw[gray!45!black,line width=0.3pt] (134,280)--(134,301) .. controls (142,305) and (168,305) .. (176,301)--(176,280) (155,280) ellipse (21 and 3.5);
				\draw[black!88,line width=1.5pt] (155,281) .. controls (154,276) and (150,275) .. (148,271);
				% ngọn lửa
				\begin{scope}[shift={(148,273)},x=1cm,y=1cm,rotate=4]\VLngonlua{0}{0}{0.95}\end{scope}
			\end{tikzpicture}
		\end{center}
		\choice
		{xích lại gần nhau hơn}
		{nở ra lớn hơn}
		{\True có tốc độ trung bình lớn hơn}
		{liên kết lại với nhau}
		\loigiai{
			Nhiệt độ tăng thì các phân tử khí chuyển động nhiệt nhanh hơn; kích thước phân tử không đổi, bình kín nên khoảng cách trung bình giữa các phân tử không đổi.
		}
	\end{ex}

	\NeedspaceCau{2.7\baselineskip}
	\begin{ex}%[VL12-C1-NB]
		Trong hệ đơn vị SI, đơn vị của nhiệt dung riêng là
		\choice
		{\True J/(kg.K)}
		{J/kg}
		{J/$^\circ$C}
		{J/K}
		\loigiai{
			$c=\dfrac{Q}{m\Delta T}$ nên đơn vị của nhiệt dung riêng là J/(kg.K).
		}
	\end{ex}

	\NeedspaceCau{5.3\baselineskip}
	\begin{ex}%[VL12-C2-NB]
		Một khối khí lí tưởng có $n$ mol, áp suất $p$, thể tích $V$, nhiệt độ tuyệt đối $T$ và hằng số khí lí tưởng $R$. Phương trình Clapeyron của khối khí lí tưởng trên là phương trình nào sau đây?
		\choice
		{\True $\dfrac{pV}{T}=nR$}
		{$\dfrac{pV}{T}=\dfrac{R}{n}$}
		{$pV=\dfrac{nR}{T}$}
		{$\dfrac{pT}{V}=nR$}
		\loigiai{
			$pV=nRT\Rightarrow\dfrac{pV}{T}=nR$.
		}
	\end{ex}

	\NeedspaceCau{5.8\baselineskip}
	\begin{ex}%[VL12-C2-TH]
		Một khối khí lí tưởng xác định có áp suất bằng áp suất khí quyển. Nếu giữ nhiệt độ của khối khí đó không đổi và làm cho áp suất của nó bằng một nửa áp suất khí quyển thì thể tích của khối khí
		\choice
		{bằng một nửa giá trị ban đầu}
		{\True bằng hai lần giá trị ban đầu}
		{bằng giá trị ban đầu}
		{bằng bốn lần giá trị ban đầu}
		\loigiai{
			Đẳng nhiệt: $p_1V_1=p_2V_2\Rightarrow V_2=\dfrac{p_1}{p_2}V_1=2V_1$.
		}
	\end{ex}

	\NeedspaceCau{5.9\baselineskip}
	\begin{ex}%[VL12-C3-NB]
		Một vòng dây dẫn kín được đặt trong từ trường có vectơ cảm ứng từ $\vv{B}$. Trong khoảng thời gian $\Delta t$, từ thông qua vòng dây biến thiên một lượng $\Delta\Phi$. Độ lớn của suất điện động cảm ứng xuất hiện trong vòng dây được xác định bằng biểu thức nào dưới đây?
		\choice
		{$|e_\text{c}|=-\dfrac{\Delta\Phi}{\Delta t}$}
		{$|e_\text{c}|=-\Delta\Phi\cdot\Delta t$}
		{$|e_\text{c}|=|\Delta\Phi|\cdot\Delta t$}
		{\True $|e_\text{c}|=\left|\dfrac{\Delta\Phi}{\Delta t}\right|$}
		\loigiai{
			Định luật Faraday: $|e_\text{c}|=\left|\dfrac{\Delta\Phi}{\Delta t}\right|$.
		}
	\end{ex}

	\NeedspaceCau{13\baselineskip}
	\begin{ex}%[VL12-C3-TH]
		Hai dòng điện $I_1$ và $I_2$ chạy trong hai dây dẫn thẳng, nằm trong mặt phẳng hình vẽ và trực giao nhau. Lực từ do dòng điện $I_1$ tác dụng lên đoạn dây dẫn AB của dòng điện $I_2$
		\begin{center}
			\begin{tikzpicture}[line join=round,line cap=round,font=\footnotesize]
				\draw[VLdo,very thick,VLgiua=0.62] (0,-1.2)--(0,1.2);
				\draw[VLdo,very thick,VLgiua=0.62] (0.6,0)--(2.1,0);
				\fill (0.6,0) circle (1.8pt) (2.1,0) circle (1.8pt);
				\node[VLdo,left,inner sep=3pt] at (0,0.35) {$I_1$};
				\node[VLdo,above,inner sep=4pt] at (1.4,0) {$I_2$};
				\node[below,inner sep=4pt] at (0.6,0) {A};
				\node[below,inner sep=4pt] at (2.1,0) {B};
			\end{tikzpicture}
		\end{center}
		\choice
		{có phương vuông góc với mặt phẳng hình vẽ, chiều từ ngoài vào trong}
		{ngược hướng với dòng điện $I_2$}
		{có phương vuông góc với mặt phẳng hình vẽ, chiều từ trong ra ngoài}
		{\True cùng hướng với dòng điện $I_1$}
		\loigiai{
			Tại AB, từ trường của $I_1$ vuông góc với mặt phẳng hình vẽ, hướng vào trong (quy tắc nắm tay phải). Theo quy tắc bàn tay trái, lực từ lên AB ($I_2$ hướng sang phải) hướng lên, cùng hướng $I_1$.
		}
	\end{ex}

	\NeedspaceCau{7.6\baselineskip}
	\begin{ex}%[VL12-C2-NB]
		Xét khối khí chứa trong một bình kín, gồm $N$ phân tử, chuyển động hỗn loạn với các tốc độ là $v_1$, $v_2$, \ldots, $v_N$. Trung bình của các bình phương tốc độ phân tử được xác định theo công thức
		\choice
		{\True $\overline{v^2}=\dfrac{v_1^2+v_2^2+\ldots+v_N^2}{N}$}
		{$\overline{v^2}=\dfrac{\left(v_1+v_2+\ldots+v_N\right)^2}{N}$}
		{$\overline{v^2}=v_1^2+v_2^2+\ldots+v_N^2$}
		{$\overline{v}=\sqrt{\dfrac{v_1^2+v_2^2+\ldots+v_N^2}{N}}$}
		\loigiai{
			Trung bình của các bình phương tốc độ: $\overline{v^2}=\dfrac{v_1^2+v_2^2+\ldots+v_N^2}{N}$.
		}
	\end{ex}

	\NeedspaceCau{15.1\baselineskip}
	\begin{ex}%[VL12-C3-TH]
		Một vòng dây dẫn hình tròn được đặt trong từ trường đều có đường sức từ vuông góc với mặt phẳng khung dây (mặt phẳng hình vẽ) hướng từ ngoài vào trong, có độ lớn cảm ứng từ $B$ phụ thuộc thời gian. Trong khoảng thời gian từ $0$~s đến $T$~s, dòng điện cảm ứng $I$ có chiều không đổi theo thời gian và đã được chỉ ra như trên hình vẽ. Đồ thị diễn tả sự biến đổi của cảm ứng từ $B$ theo thời gian có thể là hình
		\begin{center}
			\begin{tikzpicture}[line join=round,line cap=round,font=\footnotesize]
				% vòng dây
				\begin{scope}[shift={(-0.4,0.1)}]
					\fill[VLxanh!8] (-1.05,-1.05) rectangle (1.05,1.05);
					\draw[dashed,thin] (-1.05,-1.05) rectangle (1.05,1.05);
					\draw[thick,VLdong!80!black] (0,0) circle (0.8);
					\draw[VLdo,very thick,VLmui] (8:0.8) arc[start angle=8,end angle=-12,radius=0.8];
					\node[VLdo,right,inner sep=2pt] at (5:0.8) {$I$};
					\VLnhantron{-0.42}{0}
					\node[VLxanh!85!black,right,inner sep=1pt] at (-0.3,0) {$\vv{B}$};
				\end{scope}
				\foreach \sx/\sy/\k in {1.25/0.75/1,4.25/0.75/2,1.25/-1.25/3,4.25/-1.25/4}{
					\begin{scope}[shift={(\sx,\sy)}]
						\draw[VLmui,thick] (0,0)--(2.3,0) node[above,inner sep=2pt] {$t$};
						\draw[VLmui,thick] (0,0)--(0,1.05) node[right,inner sep=1pt] {$B$};
						\node[left,inner sep=1.5pt] at (0,-0.1) {$O$};
						\node[draw,thin,inner sep=1.5pt] at (0.95,-0.62) {Hình \k};
						\node[below,inner sep=2pt] at (1.9,0) {$T$};
					\end{scope}}
				\begin{scope}[shift={(1.25,0.75)}]
					\draw[very thick,VLxanh] (0,0)--(0.95,0.8)--(1.9,0);
					\draw[densely dotted,gray] (0.95,0)--(0.95,0.8);
					\node[below,inner sep=2pt] at (0.95,0) {$T/2$};
				\end{scope}
				\begin{scope}[shift={(4.25,0.75)}]
					\draw[very thick,VLxanh] (0,0.85)--(1.9,0.4);
					\draw[densely dotted,gray] (1.9,0)--(1.9,0.4);
				\end{scope}
				\begin{scope}[shift={(1.25,-1.25)}]
					\draw[very thick,VLxanh] (0,0.8)--(0.95,0)--(1.9,0.75);
					\draw[densely dotted,gray] (1.9,0)--(1.9,0.75);
					\node[below,inner sep=2pt] at (0.95,0) {$T/2$};
				\end{scope}
				\begin{scope}[shift={(4.25,-1.25)}]
					\draw[very thick,VLxanh] (0,0)--(1.9,0.75);
					\draw[densely dotted,gray] (1.9,0)--(1.9,0.75);
					\node[below,inner sep=2pt] at (0.95,0) {$T/2$};
				\end{scope}
			\end{tikzpicture}
		\end{center}
		\choice
		{\True Hình 2}
		{Hình 1}
		{Hình 3}
		{Hình 4}
		\loigiai{
			Dòng điện cảm ứng chạy cùng chiều kim đồng hồ sinh ra từ trường trong lòng vòng dây hướng vào trong, cùng chiều $\vv{B}$, nên $B$ đang giảm (định luật Lenz). Chiều dòng điện không đổi trong suốt thời gian từ $0$ đến $T$ nên $B$ giảm liên tục: Hình 2.
		}
	\end{ex}

	\NeedspaceCau{20.3\baselineskip}
	\begin{ex}%[VL12-C1-VD]
		Một nhóm học sinh sử dụng bộ thí nghiệm có sơ đồ như hình dưới để đo nhiệt nóng chảy riêng $\lambda$ của nước đá. Cho nước đá có khối lượng $m=28$~g và một ít nước lạnh vào bình nhiệt lượng kế (4), sao cho toàn bộ điện trở chìm hoàn toàn trong hỗn hợp nước và nước đá. Dùng nhiệt kế (3) có thang đo từ $-10\,^\circ\text{C}$ đến $110\,^\circ\text{C}$ để đo nhiệt độ của hỗn hợp nước và nước đá trong bình. Công suất trung bình tính được từ các số chỉ trên oát kế (2) là $14{,}24$~W. Bật nguồn điện (1). Khuấy liên tục nước đá, cứ sau mỗi $2$ phút, đọc nhiệt độ trên nhiệt kế, họ vẽ được đồ thị sự phụ thuộc của nhiệt độ $t$ ($^\circ$C) vào thời gian đun $\tau$ (s) như sau:
		% giuhinh
		\begin{center}
			\begin{tikzpicture}[x=0.0066cm,y=2.45cm,line join=round,line cap=round,font=\footnotesize]
				\draw[gray!55,thin] (0,0) grid[xstep=120,ystep=0.1] (1080,1.6);
				\draw[VLmui,thick] (0,0)--(1170,0) node[right,inner sep=2pt] {$\tau$ (s)};
				\draw[VLmui,thick] (0,0)--(0,1.75) node[left,inner sep=2pt] {$t$ ($^\circ$C)};
				\foreach \x in {0,120,...,1080}{\node[below,inner sep=2pt,font=\fontsize{6.5}{6.5}\selectfont] at (\x,0) {$\x$};}
				\foreach \y/\t in {0/0,0.1/0{,}1,0.2/0{,}2,0.3/0{,}3,0.4/0{,}4,0.5/0{,}5,0.6/0{,}6,0.7/0{,}7,0.8/0{,}8,0.9/0{,}9,1.0/1{,}0,1.1/1{,}1,1.2/1{,}2,1.3/1{,}3,1.4/1{,}4,1.5/1{,}5,1.6/1{,}6}{\node[left,inner sep=2pt,font=\fontsize{6.5}{6.5}\selectfont] at (0,\y) {$\t$};}
				\draw[very thick,VLdo] (0,0)--(666.7,0);
				\draw[very thick,VLdo] (666.7,0)--(970,1.517);
				\foreach \x in {0,120,...,600}{\filldraw[fill=white,draw=black,thick] (\x,0) circle (1.6pt);}
				\foreach \x/\y in {720/0.3,840/0.82,960/1.5}{\filldraw[fill=white,draw=black,thick] (\x,\y) circle (1.6pt);}
				\node[fill=white,inner sep=0.5pt] at (270,0.13) {$t_1=0$};
				\node[fill=white,inner sep=1pt] at (600,1.15) {$t_2=\frac{1}{200}\tau-\frac{10}{3}$};
				\node[fill=white,inner sep=1pt] at (520,0.38) {$\tau_M$};
				\draw[VLmui,thin] (540,0.3)--(655,0.03);
				% bộ thí nghiệm (dùng macro chung \VLbienthenguon, \VLnhietluongke)
				\begin{scope}[shift={(1243,0.51)},x=1cm,y=1cm]
					% đầu dây nguồn (ống cao su có gân) ở cạnh phải biến thế nguồn
					\shade[top color=black!55,bottom color=black!92] (3.44,0.4) rectangle (3.84,0.98);
					\foreach \x in {3.54,3.62,3.7,3.78}{\draw[black!45,line width=0.3pt] (\x,0.41)--(\x,0.97);}
					% biến thế nguồn (1)
					\VLbienthenguon{0}{0}{1}
					% dây nối từ hai phích cắm, luồn dưới biến thế nguồn sang oát kế
					\draw[black!85,line width=1pt] (VLnguon-do) .. controls +(0.02,-0.5) and +(-0.7,-0.02) .. (1.55,-0.3) .. controls +(0.5,0.02) and +(-0.3,-0.12) .. (2.42,-0.05);
					\draw[black!85,line width=1pt] (VLnguon-xanh) .. controls +(0.05,-0.42) and +(-0.5,-0.02) .. (1.7,-0.2) .. controls +(0.35,0.02) and +(-0.2,-0.08) .. (2.3,-0.04);
					\draw[black!85,line width=1pt] (3.46,0.22) .. controls +(0.22,-0.12) and +(-0.02,0.2) .. (3.78,-0.22);
					% nhiệt lượng kế (4) có nhiệt kế (3)
					\VLnhietluongke{4.57}{-0.05}{0.9}
					% dây nối oát kế -- nhiệt lượng kế
					\draw[black!85,line width=1pt] (4.52,-0.62) .. controls +(0.5,-0.25) and +(0.25,-0.45) .. (5.18,-0.02);
					\draw[black!85,line width=1pt] (4.5,-0.72) .. controls +(0.75,-0.3) and +(0.3,-0.35) .. (5.32,0.1);
					% oát kế (2): hộp nhỏ nhìn chếch, màn hình để trống, bốn ổ cắm
					\fill[black,opacity=0.1] (4.05,-0.9) ellipse (0.62 and 0.05);
					\fill[black!85] (3.56,-0.92) rectangle (3.7,-0.88) (4.14,-0.92) rectangle (4.28,-0.88) (4.5,-0.74)--(4.6,-0.67)--(4.6,-0.63)--(4.5,-0.7)--cycle;
					\shade[top color=gray!45,bottom color=gray!20] (3.51,-0.41)--(4.33,-0.41)--(4.66,-0.19)--(3.84,-0.19)--cycle;
					\shade[top color=gray!45,bottom color=gray!70] (4.33,-0.88)--(4.66,-0.66)--(4.66,-0.19)--(4.33,-0.41)--cycle;
					\shade[top color=gray!6,bottom color=gray!32] (3.51,-0.88) rectangle (4.33,-0.41);
					\draw[gray!40!black,line width=0.4pt] (3.51,-0.88) rectangle (4.33,-0.41) (3.51,-0.41)--(3.84,-0.19)--(4.66,-0.19)--(4.33,-0.41) (4.66,-0.19)--(4.66,-0.66)--(4.33,-0.88);
					\draw[white,line width=0.4pt] (3.53,-0.43)--(4.31,-0.43);
					\fill[black!80,rounded corners=0.4pt] (3.58,-0.63) rectangle (4.26,-0.47);
					\shade[top color=VLla!40!black,bottom color=VLla!62!black] (3.6,-0.61) rectangle (4.24,-0.49);
					\fill[white,opacity=0.16] (3.74,-0.49)--(3.88,-0.49)--(3.8,-0.61)--(3.66,-0.61)--cycle;
					\foreach \x/\m in {3.68/VLdo!80!black,3.84/VLxanh!65!black,4.0/VLdo!80!black,4.16/black!85}{
						\fill[gray!60] (\x,-0.76) circle (0.062);
						\shade[ball color=\m] (\x,-0.76) circle (0.05);
						\fill[black!85] (\x,-0.765) circle (0.018);}
					% nhãn
					\draw[thin] (VLnguon-quai)--(1.72,2.28) node[above,inner sep=1.5pt] {(1)};
					\draw[thin] (3.62,-0.93)--(3.42,-1.2) node[below,inner sep=1.5pt] {(2)};
					\draw[thin] (VLnlk-nhietke)--++(-0.28,0.22) node[above left,inner sep=1pt] {(3)};
					\draw[thin] (VLnlk-than)--(5.3,-1.15) node[below,inner sep=1.5pt] {(4)};
				\end{scope}
			\end{tikzpicture}
		\end{center}
		Nhiệt nóng chảy riêng của nước đá từ thí nghiệm này bằng
		\choice
		{$3{,}33\cdot 10^5$~J/kg}
		{$3{,}43\cdot 10^5$~J/kg}
		{$3{,}46\cdot 10^5$~J/kg}
		{\True $3{,}39\cdot 10^5$~J/kg}
		\loigiai{
			Nước đá tan hết khi nhiệt độ bắt đầu tăng: $\dfrac{\tau_M}{200}-\dfrac{10}{3}=0\Rightarrow\tau_M=\dfrac{2000}{3}$~s.\\
			$\lambda=\dfrac{P\tau_M}{m}=\dfrac{14{,}24\cdot\frac{2000}{3}}{0{,}028}\approx 3{,}39\cdot 10^5$~J/kg.
		}
	\end{ex}

	\NeedspaceCau{7.2\baselineskip}
	\begin{ex}%[VL12-C3-NB]
		Các đường sức từ xung quanh dây dẫn thẳng có dòng điện không đổi chạy qua có dạng là
		\choice
		{những đường thẳng song song với dòng điện}
		{\True những đường tròn đồng tâm với tâm nằm tại vị trí nơi dòng điện chạy qua}
		{những đường thẳng vuông góc với dòng điện}
		{những đường xoắn ốc đồng trục với trục là dòng điện}
		\loigiai{
			Đường sức từ của dòng điện thẳng là những đường tròn nằm trong mặt phẳng vuông góc với dây, có tâm nằm trên dây dẫn.
		}
	\end{ex}

	\NeedspaceCau{12.7\baselineskip}
	\begin{ex}%[VL12-C3-VD]
		Khung dây hình vuông ABCD có diện tích $S=1~\text{m}^2$ được làm từ một sợi dây dẫn đồng chất, tiết diện đều như hình vẽ. Khi khung dây chuyển động thẳng đều sang phải với tốc độ $v=10$~m/s và đang đi vào từ trường đều có độ lớn cảm ứng từ là $B=0{,}1$~T thì hiệu điện thế giữa hai đầu đoạn dây CB có giá trị là bao nhiêu vôn (V)? Biết khung dây có điện trở $R=2~\Omega$.
		\begin{center}
			\begin{tikzpicture}[line join=round,line cap=round,font=\footnotesize]
				\fill[VLxanh!8] (2.5,-0.45) rectangle (4.9,1.85);
				\draw[dashed,thin] (2.5,-0.45) rectangle (4.9,1.85);
				\foreach \x in {2.9,3.7,4.5}{\foreach \y in {0.05,0.7,1.35}{\VLchamtron{\x}{\y}}}
				\node[VLxanh!85!black,right,inner sep=1pt] at (3.8,0.85) {$\vv{B}$};
				\draw[VLdong!85!black,very thick] (0,0) rectangle (1.4,1.4);
				\fill (1.4,0.7) circle (1.6pt);
				\draw[VLluc,thick,VLmui] (1.4,0.7)--(2.15,0.7) node[above,inner sep=3pt,pos=0.8] {$\vv{v}$};
				\node[above left,inner sep=1pt] at (0,1.4) {A};
				\node[above right,inner sep=1pt] at (1.4,1.4) {B};
				\node[below right,inner sep=1pt] at (1.4,0) {C};
				\node[below left,inner sep=1pt] at (0,0) {D};
				\node[above,inner sep=2pt] at (2.5,1.85) {E};
				\node[above,inner sep=2pt] at (4.9,1.85) {F};
				\node[below,inner sep=2pt] at (4.9,-0.45) {G};
				\node[below,inner sep=2pt] at (2.5,-0.45) {H};
			\end{tikzpicture}
		\end{center}
		\choice
		{$0{,}5$~V}
		{$0{,}25$~V}
		{$1{,}0$~V}
		{\True $0{,}75$~V}
		\loigiai{
			Cạnh khung $a=1$~m; chỉ cạnh BC chuyển động trong từ trường: $e=Bav=0{,}1\cdot 1\cdot 10=1$~V, $I=\dfrac{e}{R}=0{,}5$~A.\\
			Cạnh BC đóng vai trò nguồn có điện trở $\dfrac{R}{4}=0{,}5~\Omega$: $U_\text{CB}=e-I\dfrac{R}{4}=1-0{,}5\cdot 0{,}5=0{,}75$~V.
		}
	\end{ex}
	\Closesolutionfile{ans}

	\NeedspaceCau{22\baselineskip}
	\subsection*{\indam{PHẦN II. Câu trắc nghiệm đúng sai. Thí sinh trả lời từ câu 1 đến câu 4. Trong mỗi ý a), b), c), d) ở mỗi câu, thí sinh chọn đúng hoặc sai.}}
	\setcounter{ex}{0}
	\Opensolutionfile{ans}[ans/VatLi2026-HaNoi-NguyenThiMinhKhai-Lan1-2]

	\begin{ex}%[VL12-C1-VD]
		Vào những ngày lạnh, cơ thể con người cần tiêu hao năng lượng qua đường hô hấp nhiều hơn những ngày ấm do cơ thể phải làm ấm không khí trước khi đưa vào phổi. Xét vào một ngày lạnh, không khí bên ngoài có nhiệt độ $7{,}0\,^\circ\text{C}$. Một người có nhiệt độ cơ thể $37{,}0\,^\circ\text{C}$, hít vào $20$ lần trong một phút, mỗi lần hít vào $0{,}50$ lít không khí vào phổi (không khí khi vào phổi được làm ấm đến nhiệt độ cơ thể). Năng lượng cơ thể người đó tiêu hao để làm ấm không khí hít vào trong ngày đó ($24$ giờ) chiếm $7{,}23\%$ tổng năng lượng hoạt động trong $24$ giờ. Biết không khí có khối lượng riêng $1{,}30~\text{kg/m}^3$ và có nhiệt dung riêng $1{,}02\cdot 10^3$~J/(kg.K).
		\NeedspaceBang{10.9\baselineskip}%
		\choiceTFt
		{Nhiệt độ cơ thể người đó theo thang Kelvin bằng $300$~K}
		{Nhiệt lượng cơ thể tiêu hao do làm ấm không khí sau mỗi lần hít vào bằng $20{,}5$~J}
		{\True Nhiệt lượng cần cung cấp để làm nóng $1$~kg không khí tăng thêm $1\,^\circ\text{C}$ bằng $1{,}02\cdot 10^3$~J}
		{\True Tổng năng lượng hoạt động trong $24$ giờ là $7{,}92$~MJ (kết quả được làm tròn và viết đến hai chữ số thập phân)}
		\loigiai{
			\begin{itemchoice}
				\itemch \textbf{Sai.} $T=37+273=310$~K.
				\itemch \textbf{Sai.} $Q_1=DVc\Delta t=1{,}30\cdot 0{,}5\cdot 10^{-3}\cdot 1{,}02\cdot 10^3\cdot(37-7)\approx 19{,}9$~J.
				\itemch \textbf{Đúng.} $Q=mc\Delta t=1\cdot 1{,}02\cdot 10^3\cdot 1=1{,}02\cdot 10^3$~J.
				\itemch \textbf{Đúng.} Trong $24$ giờ: $Q=19{,}89\cdot 20\cdot 60\cdot 24\approx 5{,}728\cdot 10^5$~J; $W=\dfrac{Q}{0{,}0723}\approx 7{,}92\cdot 10^6~\text{J}=7{,}92$~MJ.
			\end{itemchoice}
		}
	\end{ex}

	\NeedspaceCau{24.2\baselineskip}
	\begin{ex}%[VL12-C3-VD]
		Để đo cảm ứng từ trong lòng một nam châm, một bạn học sinh bố trí thí nghiệm như Hình a. Ban đầu, khung dây được giữ cân bằng trên nêm kim loại bằng thanh cách điện.
		% giuhinh
		\begin{center}
			\begin{tikzpicture}[line join=round,line cap=round,font=\footnotesize]
				% ===== Hình a (toạ độ theo điểm ảnh của hình gốc) =====
				\begin{scope}[x=0.0195cm,y=-0.0195cm,
					khung/.style={draw=VLdong!70!black,line width=1.2pt},
					mach/.style={draw=black,line width=0.55pt}]
					% tấm đế
					\fill[black,opacity=0.07] (48,131)--(257,136)--(301,82)--(297,78)--(254,129)--cycle;
					\shade[top color=VLxanh!10!gray!12,bottom color=VLxanh!4!gray!5] (45,119.5)--(88,69)--(297,74)--(254,124.5)--cycle;
					\shade[top color=gray!30,bottom color=gray!50] (45,119.5)--(254,124.5)--(254,129)--(45,124)--cycle;
					\fill[gray!58] (254,124.5)--(297,74)--(297,78.5)--(254,129)--cycle;
					\draw[gray!55!black,line width=0.3pt] (45,119.5)--(88,69)--(297,74)--(254,124.5)--cycle (45,119.5)--(45,124)--(254,129)--(254,124.5) (254,129)--(297,78.5)--(297,74);
					\draw[white,line width=0.4pt] (47,120.6)--(253,125.5);
					% nam châm chữ U nằm ngửa: quét mặt cắt chữ U theo chiều sâu (từ sau ra trước)
					\def\nuaN{(59,80) .. controls (58,104) and (70,119) .. (87,119)--(86,105) .. controls (78,105) and (72,97) .. (72,82)--cycle}
					\def\nuaS{(87,119) .. controls (104,119) and (113,104) .. (112,82)--(100,82) .. controls (100,97) and (94,105) .. (86,105)--cycle}
					\fill[black,opacity=0.12] (66,118) .. controls (80,124) and (100,123) .. (108,116)--(146,85)--(140,80)--cycle;
					\foreach \t in {1,0.97,...,0.02}{
						\pgfmathsetmacro{\k}{62-22*\t}
						\begin{scope}[shift={(37*\t,-31*\t)}]
							\shade[top color=VLdo!\k!white,bottom color=VLdo!78!black] \nuaN;
							\shade[top color=VLxanh!55!white,bottom color=VLxanh!85!black] \nuaS;
						\end{scope}}
					% đường bao phía sau
					\draw[black!65,line width=0.3pt] (101.7,114)--(138.7,83) .. controls (146,76) and (149.6,65) .. (149,51);
					% mặt trên hai nhánh
					\fill[VLdo!42!white] (59,80)--(72,82)--(109,51)--(96,49)--cycle;
					\fill[VLxanh!28!white] (100,82)--(112,82)--(149,51)--(137,51)--cycle;
					\draw[black!65,line width=0.3pt] (59,80)--(96,49)--(109,51)--(72,82) (100,82)--(137,51)--(149,51)--(112,82);
					% mặt trước
					\shade[left color=VLdo!78!black,right color=VLdo!90!white] \nuaN;
					\shade[left color=VLxanh!50!white,right color=VLxanh!85] \nuaS;
					\draw[black!65,line width=0.35pt] (59,80) .. controls (58,104) and (70,119) .. (87,119) .. controls (104,119) and (113,104) .. (112,82)--(100,82) .. controls (100,97) and (94,105) .. (86,105) .. controls (78,105) and (72,97) .. (72,82)--cycle;
					\draw[white,opacity=0.45,line width=0.5pt] (61.5,84) .. controls (61.5,103) and (71,115) .. (85,116.5);
					\node[rotate=40,inner sep=0pt,font=\fontsize{6}{6}\selectfont\bfseries] at (74.3,72.6) {N};
					\node[rotate=40,inner sep=0pt,font=\fontsize{6}{6}\selectfont\bfseries] at (115.3,74) {S};
					% nêm kim loại phía sau, dây nối ra nguồn
					\shade[left color=gray!70,right color=gray!55,middle color=white] (232.7,86)--(232.7,40)--(236,34)--(239.3,40)--(239.3,86) arc[start angle=0,end angle=180,x radius=3.3,y radius=2.2];
					\draw[gray!35!black,line width=0.3pt] (232.7,86)--(232.7,40)--(236,34)--(239.3,40)--(239.3,86) arc[start angle=0,end angle=180,x radius=3.3,y radius=2.2];
					\shade[top color=gray!15,bottom color=gray!65,rounded corners=0.5pt] (230.6,75.5) rectangle (241.6,80.5);
					\draw[gray!35!black,line width=0.3pt,rounded corners=0.5pt] (230.6,75.5) rectangle (241.6,80.5);
					\draw[mach] (241.6,78)--(315.2,78) (319.7,78)--(323,78) (328.4,78)--(331.2,78) (335.7,78)--(348,78)--(348,103) (348,130)--(348,156)--(338,156) (322,156)--(298.5,156) (281.5,156)--(229,156)--(229,105.5)--(213.5,105.5);
					% nguồn hai pin: bản ngắn (cực âm) bên trái, bản dài (cực dương) bên phải
					\draw[black,line width=1.3pt,line cap=butt] (315.8,73)--(315.8,83.5) (331.8,73)--(331.8,83.5);
					\draw[black,line width=0.6pt,line cap=butt] (319.7,66)--(319.7,90) (335.7,66)--(335.7,90);
					\fill (325.7,78) circle (0.5pt);
					% biến trở, khóa K, ampe kế
					\filldraw[fill=white,draw=black,line width=0.55pt] (344.3,103) rectangle (351.7,130);
					\draw[line width=0.55pt,VLmui] (340.5,128.5)--(357,108.5);
					\draw[mach] (322,156)--(336.5,146.5);
					\filldraw[fill=white,draw=black,line width=0.4pt] (322,156) circle (1.1) (338,156) circle (1.1);
					\filldraw[fill=white,draw=black,line width=0.6pt] (290,156) circle (8.5);
					\node[inner sep=0pt,font=\scriptsize\sffamily\bfseries] at (290,156) {A};
					\node[inner sep=0pt] at (331,167.5) {K};
					% khung dây: cạnh X nằm trong lòng nam châm
					\draw[khung] (86,85.5)--(135,47.5)--(135,30)--(266,32);
					\draw[khung] (86,85.5)--(86,66)--(222,67);
					\draw[VLdong!70!black,line width=1.7pt] (86,85.5)--(135,47.5);
					\node[inner sep=0pt,font=\scriptsize\sffamily\bfseries] at (85.5,92.5) {X};
					% thanh cách điện
					\fill[gray!50] (230.5,69.5)--(275,34)--(275,37.5)--(230.5,73)--cycle;
					\fill[gray!28] (221.5,65.5)--(230.5,69.5)--(230.5,73)--(221.5,69)--cycle;
					\shade[left color=white,right color=gray!18] (221.5,65.5)--(266,31)--(275,34)--(230.5,69.5)--cycle;
					\draw[gray!30!black,line width=0.3pt] (221.5,65.5)--(266,31)--(275,34)--(275,37.5)--(230.5,73)--(221.5,69)--cycle (221.5,65.5)--(230.5,69.5)--(275,34) (230.5,69.5)--(230.5,73);
					% nêm kim loại phía trước
					\shade[left color=gray!70,right color=gray!55,middle color=white] (204.7,117)--(204.7,77)--(208,71)--(211.3,77)--(211.3,117) arc[start angle=0,end angle=180,x radius=3.3,y radius=2.2];
					\draw[gray!35!black,line width=0.3pt] (204.7,117)--(204.7,77)--(208,71)--(211.3,77)--(211.3,117) arc[start angle=0,end angle=180,x radius=3.3,y radius=2.2];
					\shade[top color=gray!15,bottom color=gray!65,rounded corners=0.5pt] (202.5,103) rectangle (213.5,108);
					\draw[gray!35!black,line width=0.3pt,rounded corners=0.5pt] (202.5,103) rectangle (213.5,108);
					\draw[khung] (200,66.8)--(216,67);
					% nhãn
					\draw[line width=0.4pt] (160,31.5)--(167,21.5);
					\node[inner sep=0pt,anchor=south] at (171,19.5) {Dây dẫn};
					\draw[line width=0.4pt] (274,31.5)--(283,25);
					\node[inner sep=0pt,anchor=west] at (285,23) {Thanh cách điện};
					\draw[line width=0.4pt] (206.5,113)--(186,136);
					\node[inner sep=0pt,anchor=north] at (160,138.5) {Nêm kim loại};
					\draw[line width=0.4pt] (62.5,88)--(49,95);
					\node[inner sep=0pt,anchor=east,align=center] at (47,95) {Nam\\châm};
					\node[font=\footnotesize\itshape] at (190,187) {Hình a};
				\end{scope}
				% ===== Hình b =====
				\begin{scope}[shift={(1.35cm,-0.63cm)},x=0.015cm,y=-0.015cm]
					% nam châm chữ U nhìn chính diện
					\shade[left color=VLdo!78!black,right color=VLdo!90!white] (392.7,82)--(392.7,122.4) arc[start angle=180,end angle=90,radius=43.3]--(436,148.1) arc[start angle=90,end angle=180,radius=25.7]--(410.3,82)--cycle;
					\shade[left color=VLxanh!50!white,right color=VLxanh!85] (479.3,82)--(479.3,122.4) arc[start angle=0,end angle=90,radius=43.3]--(436,148.1) arc[start angle=90,end angle=0,radius=25.7]--(461.7,82)--cycle;
					\draw[black!65,line width=0.35pt] (392.7,82)--(392.7,122.4) arc[start angle=180,end angle=0,radius=43.3]--(479.3,82)--(461.7,82)--(461.7,122.4) arc[start angle=0,end angle=180,radius=25.7]--(410.3,82)--cycle;
					\draw[white,opacity=0.45,line width=0.5pt] (396,86)--(396,122.4) arc[start angle=180,end angle=100,radius=40];
					\node[white,inner sep=0pt,font=\scriptsize\bfseries] at (401.5,96) {N};
					\node[inner sep=0pt,font=\scriptsize\bfseries] at (470.5,96) {S};
					% nêm kim loại có cọc nối dây
					\draw[line width=0.55pt] (632,137.3)--(695,137.3);
					\shade[top color=gray!15,bottom color=gray!65] (609.5,133) rectangle (633,141.6);
					\draw[gray!35!black,line width=0.3pt] (609.5,133) rectangle (633,141.6);
					\shade[left color=gray!70,right color=gray!55,middle color=white] (614.8,158.3)--(614.8,72)--(621,63.5)--(627.2,72)--(627.2,158.3)--cycle;
					\draw[gray!35!black,line width=0.35pt] (614.8,158.3)--(614.8,72)--(621,63.5)--(627.2,72)--(627.2,158.3)--cycle;
					% kích thước 8 cm
					\draw[gray!60!black,line width=0.3pt] (436,17)--(436,55) (621,17)--(621,59);
					\draw[line width=0.4pt,{Stealth[length=4pt,width=3pt]}-{Stealth[length=4pt,width=3pt]}] (437,34)--(620,34);
					\node[inner sep=0pt,anchor=south] at (528,29) {$8$~cm};
					% khung dây, thanh cách điện, vật nhỏ
					\draw[VLdong!70!black,line width=1.7pt,line cap=butt,line join=miter] (436,96)--(436,60)--(672,60);
					\node[inner sep=0pt,font=\footnotesize\sffamily\bfseries] at (436,107.5) {X};
					\shade[top color=white,bottom color=gray!22] (671,52.3) rectangle (693.3,72.7);
					\draw[gray!30!black,line width=0.35pt] (671,52.3) rectangle (693.3,72.7);
					\shade[top color=gray!45,bottom color=gray!80,rounded corners=0.6pt] (678.3,45) rectangle (685.3,52.3);
					\draw[gray!25!black,line width=0.3pt,rounded corners=0.6pt] (678.3,45) rectangle (685.3,52.3);
					% nhãn
					\draw[line width=0.4pt] (686,47)--(706,39);
					\node[inner sep=0pt,anchor=west] at (708,37.5) {Vật nhỏ};
					\draw[line width=0.4pt] (693.5,65)--(710,73.5);
					\node[inner sep=0pt,anchor=west] at (712,77) {Thanh cách điện};
					\draw[line width=0.4pt] (624,116)--(647,108);
					\node[inner sep=0pt,anchor=west] at (649,106) {Nêm kim loại};
					\node[font=\footnotesize\itshape] at (560,201.5) {Hình b};
				\end{scope}
			\end{tikzpicture}
		\end{center}
		\NeedspaceBang{13.9\baselineskip}%
		\choiceTFt
		{Khi đóng khóa $K$, phần khung dây (X) trong lòng nam châm chịu tác dụng của lực từ. Đoạn dây X di chuyển đi lên}
		{Với cách bố trí thí nghiệm như Hình a, để tăng độ lớn của lực từ tác dụng vào đoạn dây X, bạn học sinh đó tăng điện trở của biến trở}
		{Để lấy lại cân bằng cho khung dây, bạn học sinh đó đặt lên trên thanh cách điện một vật nhỏ (Hình b). Biết đoạn dây X dài $5$~cm và ampe kế chỉ $1{,}5$~A. Vật nhỏ gây ra một moment có độ lớn $4\cdot 10^{-4}$~N$\cdot$m. Độ lớn của cảm ứng từ trong lòng nam châm bằng $0{,}08$~T}
		{\True Từ trường trong lòng nam châm ở thí nghiệm nêu trên có các đường sức từ song song và cách đều nhau (từ trường đều)}
		\loigiai{
			\begin{itemchoice}
				\itemch \textbf{Sai.} Vật nhỏ đặt trên thanh cách điện (phía bên kia nêm so với X) làm khung cân bằng trở lại, nên lực từ kéo đoạn dây X đi xuống.
				\itemch \textbf{Sai.} Tăng điện trở của biến trở thì $I$ giảm nên $F=BIl$ giảm.
				\itemch \textbf{Sai.} Cân bằng moment đối với nêm: $BIl\cdot d=M\Rightarrow B=\dfrac{M}{Ild}=\dfrac{4\cdot 10^{-4}}{1{,}5\cdot 0{,}05\cdot 0{,}08}\approx 0{,}067$~T.
				\itemch \textbf{Đúng.} Trong khe giữa hai cực của nam châm chữ U, các đường sức từ song song, cách đều: từ trường đều.
			\end{itemchoice}
		}
	\end{ex}

	\NeedspaceCau{25.7\baselineskip}
	\begin{ex}%[VL12-C3-VD]
		Khung dây gồm $200$ vòng, có điện trở $r=1~\Omega$, phần diện tích khung dây được đặt trong từ trường vuông góc với mặt phẳng trang giấy và hướng vào trong (như hình vẽ dưới) là $0{,}3~\text{m}^2$. Độ lớn cảm ứng từ thay đổi theo thời gian theo biểu thức $B=0{,}05t$ ($B$ tính theo T; $t$ tính theo s). Khung dây mắc với điện trở $R=5~\Omega$ nằm trong bình cách nhiệt nằm ngang có thể tích bình là $1{,}0$ lít (ở miệng bình có vòng chặn). Pit-tông nhẹ bịt kín một khối khí lí tưởng trong bình và có thể chuyển động không ma sát với thành bình. Lúc đầu $t=0$ pit-tông đứng yên ở vị trí chính giữa của bình, áp suất khí quyển là $10^5$~Pa và nhiệt độ khí là $250$~K. Kể từ $t=0$, khí được làm nóng từ từ và pit-tông chuyển động chậm. Tại $t=\tau$, pit-tông dừng lại khi vừa chạm vào vòng chặn, nội năng của khí trong bình tăng thêm $100$~J và nhiệt độ khí trong bình là $T$. Coi toàn bộ nhiệt lượng do $R$ tỏa ra được khí hấp thụ.
		\begin{center}
			\begin{tikzpicture}[line join=round,line cap=round,font=\footnotesize]
				\ctikzset{bipoles/length=0.55cm}
				% từ trường
				\foreach \x in {-0.25,0.5,1.25,2.0}{\foreach \y in {0.0,0.8,1.6,2.4}{\VLnhan{\x}{\y}}}
				\node[VLxanh!85!black] at (0.13,2.02) {$\vv{B}$};
				% khung dây
				\draw[VLdong!85!black,thick] ([shift={(10:0.85)}]1.05,1.2) arc[start angle=10,end angle=350,radius=0.85];
				\draw[VLdong!85!black,thick] ([shift={(10:0.85)}]1.05,1.2)--(2.65,1.35)--(2.65,1.9)--(3.6,1.9);
				\draw[VLdong!85!black,thick] ([shift={(-10:0.85)}]1.05,1.2)--(2.65,1.05)--(2.65,0.5)--(3.6,0.5);
				% bình
				\fill[VLxanh!10] (3.35,0.15) rectangle (4.6,2.1);
				\foreach \x/\y in {3.45/0.3,3.85/0.25,4.2/0.4,4.45/0.22,3.95/0.55,4.4/0.7,3.9/0.85,4.25/1.0,3.8/1.2,4.45/1.25,4.05/1.4,3.9/1.65,4.3/1.6,4.5/1.9,4.1/1.95,3.45/1.95,3.85/1.8}{\fill[black!60] (\x,\y) circle (0.6pt);}
				\fill[gray!55] (3.2,0)--(6.1,0)--(6.1,0.6)--(5.95,0.6)--(5.95,0.15)--(3.35,0.15)--(3.35,2.1)--(5.95,2.1)--(5.95,1.65)--(6.1,1.65)--(6.1,2.25)--(3.2,2.25)--cycle;
				\draw[black!75,thin] (3.2,0)--(6.1,0)--(6.1,0.6)--(5.95,0.6)--(5.95,0.15)--(3.35,0.15)--(3.35,2.1)--(5.95,2.1)--(5.95,1.65)--(6.1,1.65)--(6.1,2.25)--(3.2,2.25)--cycle;
				% pit-tông
				\shade[left color=gray!75,right color=gray!45,middle color=gray!15] (4.6,0.15) rectangle (4.75,2.1);
				\draw[black!75,thin] (4.6,0.15) rectangle (4.75,2.1);
				% điện trở
				\draw (3.6,1.9) to[R,l=$R$] (3.6,0.5);
				% nhãn
				\node[draw,thin,fill=white,inner sep=3pt] (kd) at (1.4,2.95) {Khung dây};
				\draw[VLmui,thin] (kd.south)--(1.15,2.08);
				\node[draw,thin,fill=white,inner sep=3pt] (kh) at (4.0,2.95) {Khí lí tưởng};
				\draw[VLmui,thin] (kh.south)--(4.15,1.45);
				\node[draw,thin,fill=white,inner sep=3pt] (vc) at (6.45,2.95) {Vòng chặn};
				\draw[VLmui,thin] (vc.south)--(6.05,1.9);
				\node[draw,thin,fill=white,inner sep=3pt] (pt) at (7.0,1.1) {Pit-tông};
				\draw[VLmui,thin] (pt.west)--(4.75,1.1);
			\end{tikzpicture}
		\end{center}
		\NeedspaceBang{8.1\baselineskip}%
		\choiceTFt
		{Dòng điện cảm ứng chạy trong phần vòng dây tròn cùng chiều kim đồng hồ}
		{\True Độ lớn cường độ dòng điện cảm ứng chạy qua mạch là $0{,}5$~A}
		{\True Giá trị $T=500$~K}
		{Tại $t=160$~s, nhiệt độ của khí trong bình là $T'=600$~K}
		\loigiai{
			Suất điện động cảm ứng: $e=NS\left|\dfrac{\Delta B}{\Delta t}\right|=200\cdot 0{,}3\cdot 0{,}05=3$~V.
			\begin{itemchoice}
				\itemch \textbf{Sai.} $\vv{B}$ hướng vào trong và tăng, dòng điện cảm ứng sinh ra từ trường trong lòng vòng dây hướng ra ngoài: dòng điện chạy ngược chiều kim đồng hồ.
				\itemch \textbf{Đúng.} $I=\dfrac{e}{R+r}=\dfrac{3}{5+1}=0{,}5$~A.
				\itemch \textbf{Đúng.} Đẳng áp ($p=10^5$~Pa) từ $V_1=0{,}5$~lít đến $V_2=1{,}0$~lít: $T=T_1\dfrac{V_2}{V_1}=250\cdot 2=500$~K.
				\itemch \textbf{Sai.} $P=I^2R=1{,}25$~W. Giai đoạn đẳng áp: $Q=\Delta U+p\Delta V=100+10^5\cdot 0{,}5\cdot 10^{-3}=150$~J, $\tau=\dfrac{150}{1{,}25}=120$~s.\\
				Từ $120$~s đến $160$~s khí nhận thêm $1{,}25\cdot 40=50$~J, đẳng tích nên $\Delta U'=50$~J. Nội năng tỉ lệ với $T$: $T'=500+250\cdot\dfrac{50}{100}=625$~K.
			\end{itemchoice}
		}
	\end{ex}

	\NeedspaceCau[15.1\baselineskip]{23.8\baselineskip}
	\begin{ex}%[VL12-C1-VD]
		Hình ảnh dưới đây là ống phát tia X, khi sợi nung FF$'$ được đốt nóng, electron phát xạ nhiệt ra khỏi cathode (K), khi ra khỏi cathode electron được tăng tốc bởi điện trường mạnh giữa cathode và anode (A), electron thu được động năng lớn và đập vào anode tạo ra tia X. Do hiệu suất phát tia X của ống rất thấp chỉ là $1{,}0\%$, phần năng lượng còn lại làm nóng anode vì vậy luôn có dòng nước chảy vào và chảy ra để làm nguội anode. Biết hiệu điện thế $U_\text{AK}=1{,}0\cdot 10^4$~V, cường độ dòng điện qua ống $I=10$~A, nhiệt độ nước chảy vào $25\,^\circ\text{C}$, nhiệt độ nước chảy ra $55\,^\circ\text{C}$, nhiệt dung riêng của nước là $4200$~J/(kg.K), khối lượng riêng của nước $1000~\text{kg/m}^3$.
		% giuhinh
		\begin{center}
			\begin{tikzpicture}[line join=round,line cap=round,font=\footnotesize]
				% ===== sơ đồ ống phát tia X (toạ độ theo điểm ảnh của hình gốc) =====
				\begin{scope}[x=0.0128cm,y=-0.0128cm]
					% vỏ thủy tinh: cổ cathode -- bầu -- cổ anode
					\filldraw[fill=VLxanh!6,draw=black!85,line width=0.7pt] (60,113)--(156,113) arc[start angle=194.6,end angle=342.2,radius=95]--(431,108) arc[start angle=-90,end angle=90,radius=24]--(340.5,156) arc[start angle=11.5,end angle=166,radius=95]--(60,160) arc[start angle=90,end angle=270,radius=23.5]--cycle;
					% cực âm, dây nối cathode K
					\draw[line width=0.5pt] (20,25) circle (6.5) (14,34)--(26,16);
					\node[inner sep=0pt] at (44,26) {$-$};
					\draw[line width=0.5pt,rounded corners=3pt] (20,31.5)--(20,137)--(102,137);
					\node[inner sep=0pt] at (117,137) {$K$};
					\draw[line width=0.5pt] (131,137)--(154.5,137);
					\draw[line width=0.8pt] (159,119.5) .. controls (153,130) and (153,144) .. (159,154.5);
					% sợi nung FF'
					\foreach \y in {129.3,137.5,145.7}{\draw[VLdo,line width=0.6pt] (175,\y) ellipse (6.5 and 4.6);}
					\draw[VLdo,line width=0.6pt] (176,125)--(176,96)--(160,78) (176,150)--(176,180)--(155,197);
					\draw[line width=0.5pt,VLmui] (160,78)--(128,78);
					\draw[line width=0.5pt,VLmui] (155,197)--(128,197);
					\node[inner sep=0pt] at (133,57) {$F$};
					\node[inner sep=0pt] at (119,177) {$F'$};
					% chùm electron
					\draw[densely dashed,gray!70!black,line width=0.45pt] (185,124)--(236,134) (185,151)--(240,141);
					% ống dẫn nước làm nguội trong thân anode
					\filldraw[fill=white,draw=black!85,line width=0.5pt] (277,124)--(467,124)--(467,138.5)--(277,138.5);
					\draw[black!85,line width=0.5pt,rounded corners=2.5pt] (285,129.3)--(481,129.3)--(488,121)--(488,100);
					\draw[line width=0.5pt,VLmui] (488,84)--(488,105);
					\draw[line width=0.5pt,VLmui,rounded corners=3pt] (467,134.2)--(497,134.2)--(497,166);
					% anode A
					\filldraw[pattern=north east lines,draw=black,line width=0.7pt] (230,117)--(277,117)--(277,152)--(256,152)--cycle;
					\node[inner sep=0pt] at (258,101) {$A$};
					% cực dương
					\draw[line width=0.5pt] (462,22) circle (6.5) (456,31)--(468,13);
					\node[inner sep=0pt] at (434,22) {$+$};
					\draw[line width=0.5pt] (460.5,28.3)--(460.5,124);
					% tia X
					\foreach \p in {(163,222),(200,247),(245,255),(296,237)}{\draw[line width=0.55pt,VLmui] (246.5,140.5)--\p;}
					\node[inner sep=0pt] at (239,284) {Tia X};
					\node[inner sep=0pt,align=right,anchor=north east] at (488,176) {Nước làm\\nguội};
				\end{scope}
				% ===== mô hình ống phát tia X (toạ độ theo điểm ảnh của ảnh gốc) =====
				\begin{scope}[shift={(6.57cm,-0.04cm)},x=0.0155cm,y=-0.0155cm]
					\definecolor{TXtuong}{RGB}{192,214,58}\definecolor{TXtuongtoi}{RGB}{128,148,22}
					\definecolor{TXsan}{RGB}{226,166,88}\definecolor{TXsansang}{RGB}{243,208,136}
					\definecolor{TXkinh}{RGB}{206,226,150}\definecolor{TXkinhdam}{RGB}{96,120,38}
					\definecolor{TXnau}{RGB}{150,72,24}\definecolor{TXgo}{RGB}{104,58,26}
					\begin{scope}
						\clip[rounded corners=1.5pt] (5,5.5) rectangle (425,241);
						% tường, sàn gỗ
						\shade[inner color=TXtuong,outer color=TXtuongtoi] (-110,-150) rectangle (520,260);
						\shade[top color=TXsan,bottom color=TXsansang] (5,160) rectangle (425,241);
						\shade[top color=black,bottom color=TXsan,opacity=0.22] (5,160) rectangle (425,167);
						\foreach \y in {181,201,222}{\draw[TXgo,opacity=0.22,line width=0.3pt] (5,\y)--(425,\y);}
						% bóng đổ trên sàn
						\fill[black,opacity=0.22] (203,188) ellipse (64 and 5.5);
						\fill[black,opacity=0.15] (60,204) ellipse (34 and 3) (120,192) ellipse (20 and 3);
						% dây cáp trên sàn (luồn dưới bóng thủy tinh)
						\draw[TXgo!60!black,line width=0.55pt] (124,153) .. controls (134,180) and (165,189) .. (205,190) .. controls (245,191) and (285,188) .. (302,181);
						% bàn nhỏ phía dưới + hộp đen
						\fill[TXgo!75!black] (41,184) rectangle (44,198) (71,184) rectangle (74,198);
						\shade[top color=black!60,bottom color=black!95] (41,167) rectangle (73,180);
						\fill[gray!55] (50,172.5) rectangle (64,175.5);
						\shade[top color=TXgo!80!white,bottom color=TXgo] (29,180) rectangle (84,184);
						\shade[left color=TXgo,right color=TXgo!70!black] (30,184) rectangle (34,203) (79,184) rectangle (83,203);
						% bàn phía trên
						\fill[TXgo!75!black] (53,153.5) rectangle (56,168) (113,153.5) rectangle (116,186);
						\shade[left color=TXgo,right color=TXgo!70!black] (35,153.5) rectangle (39.5,181) (120,153.5) rectangle (124.5,190);
						% ống thủy tinh phía cathode (sau hộp nguồn)
						\shade[top color=TXkinh!70!white,bottom color=TXkinh!70!TXkinhdam,opacity=0.75,rounded corners=3pt] (84,112.5) rectangle (150,136);
						\draw[TXkinhdam,line width=0.35pt,rounded corners=3pt] (150,112.5)--(84,112.5)--(84,136)--(150,136);
						\fill[white,opacity=0.5,rounded corners=0.8pt] (89,115) rectangle (128,117.5);
						% hộp nguồn trên bàn, cọc nối và dây
						\draw[black!80,line width=0.5pt] (92,119) .. controls (100,110) and (118,108) .. (131,116) (70,119) .. controls (66,112) and (72,110) .. (76,114);
						\fill[black!85] (68.5,117) rectangle (71.5,125) (90.5,117) rectangle (93.5,125);
						\shade[top color=black!58,bottom color=black!96] (55,124.5) rectangle (109,149.5);
						\draw[gray!55,line width=0.3pt] (55.5,125)--(108.5,125);
						\fill[gray!60] (71,135.5) rectangle (88,138.5);
						\shade[top color=TXgo!80!white,bottom color=TXgo] (33,149.5) rectangle (128,153.5);
						% giá đỡ ống cathode
						\shade[left color=black!70,right color=black!45] (120.6,140) rectangle (123,194.5);
						\shade[top color=black!45,bottom color=black!85] (122,195.5) ellipse (12 and 3.4);
						\fill[TXkinhdam!70!black] (115.5,136.5)--(128,136.5)--(125,140.5)--(118.5,140.5)--cycle;
						% bóng thủy tinh: nền
						\begin{scope}
							\clip (202,124.5) circle (59);
							\shade[top color=TXkinh!92!white,bottom color=TXtuong!55!TXkinhdam] (140,64) rectangle (264,143);
							\shade[top color=TXtuong!55!TXkinhdam,bottom color=black!90!TXkinhdam] (140,142.8) rectangle (264,180);
							\fill[black!90!TXkinhdam] (140,179.8) rectangle (264,185);
						\end{scope}
						% cathode: trụ kim loại bóng (nửa trên sáng, nửa dưới tối)
						\shade[top color=gray!22,bottom color=white] (131,112.3) rectangle (172,123.5);
						\shade[top color=gray!70,bottom color=black] (131,123.5) rectangle (172,136.3);
						\draw[black!55,line width=0.3pt] (148,112.3)--(148,136.3);
						\shade[left color=gray!55,right color=gray!15] (172,124.3) ellipse (2.6 and 12);
						\draw[black!60,line width=0.25pt] (172,124.3) ellipse (2.6 and 12);
						% chùm electron
						\shade[left color=green!85!yellow,right color=green!50!TXkinh,opacity=0.85,rounded corners=1.2pt] (174,120.5) rectangle (213,127.5);
						\draw[green!30!white,line width=0.5pt,opacity=0.9] (175,124)--(210,124);
						% anode (bia) và thanh anode trong bóng
						\shade[top color=TXnau!85!white,bottom color=TXnau!80!black] (209,109)--(236,109)--(236,136)--(222,136)--cycle;
						\draw[TXnau!45!black,line width=0.3pt] (209,109)--(236,109)--(236,136)--(222,136)--cycle;
						\shade[top color=TXnau!90!white,bottom color=TXnau!65!black] (236,119.5) rectangle (264,130);
						% bóng thủy tinh: lớp kính phủ, viền, vệt sáng
						\shade[top color=white,bottom color=TXkinhdam,opacity=0.14] (202,124.5) circle (59);
						\draw[TXkinhdam,line width=0.7pt] (202,124.5) circle (59);
						\draw[white,opacity=0.55,line width=0.6pt] ($(202,124.5)+(215:55)$) arc[start angle=215,end angle=258,radius=55];
						\fill[white,opacity=0.3] (178,84) ellipse (14 and 6);
						% cột thủy tinh đứng, ống nghiêng (bình nước làm nguội)
						\shade[left color=black!60,right color=gray!45,middle color=gray!12] (272.5,158.5)--(279,148.5)--(354,180)--(356,188) .. controls (352,195) and (344,195) .. (339,191)--cycle;
						\draw[black!55,line width=0.3pt] (272.5,158.5)--(279,148.5)--(354,180) (272.5,158.5)--(339,191);
						\draw[white,opacity=0.7,line width=0.5pt] (281,154)--(338,179.5);
						\shade[left color=TXkinh!60!TXkinhdam,right color=TXkinh!75!TXkinhdam,middle color=white,opacity=0.85] (344,47)--(344,181) .. controls (344,190) and (357.5,190) .. (357.5,181)--(357.5,47)--cycle;
						\draw[TXkinhdam!80!black,line width=0.35pt] (344,47)--(344,181) .. controls (344,190) and (357.5,190) .. (357.5,181)--(357.5,47);
						\shade[top color=black!35,bottom color=black!75,opacity=0.75] (344,160)--(344,181) .. controls (344,190) and (357.5,190) .. (357.5,181)--(357.5,160)--cycle;
						\fill[TXkinhdam!85!black] (350.7,46.5) ellipse (8.3 and 2.2);
						% giá treo có cần cong
						\draw[TXkinhdam!85!black,line width=0.9pt] (385,193)--(385,37) .. controls (385,12) and (353,13) .. (351.5,44.5);
						\draw[TXkinh,opacity=0.6,line width=0.25pt] (384.4,190)--(384.4,38);
						\shade[top color=TXkinhdam!60!TXkinh,bottom color=TXkinhdam!80!black] (385,194.5) ellipse (16.5 and 4);
						\fill[black,opacity=0.15] (387,198) ellipse (18 and 2.5);
						% giá đỡ ống anode
						\shade[left color=black!70,right color=black!45] (285,139) rectangle (287.2,195);
						\shade[top color=black!45,bottom color=black!85] (286,196.2) ellipse (12.5 and 3.2);
						% thanh anode ngoài bóng, ống thủy tinh cổ anode, khối nối đen
						\shade[top color=TXnau!90!white,bottom color=TXnau!65!black] (258,119.5) rectangle (368,130);
						\shade[top color=white,middle color=TXkinh!80!white,bottom color=TXkinh!60!TXkinhdam,opacity=0.55,rounded corners=3.5pt] (255,112)--(322,112)--(322,137.5)--(255,137.5);
						\fill[yellow!25!white,opacity=0.75,rounded corners=1.2pt] (264,113.2) rectangle (318,118.2);
						\draw[TXkinhdam,line width=0.35pt,rounded corners=3.5pt] (259.5,112)--(322,112)--(322,137.5)--(259.5,137.5);
						\fill[TXkinhdam!70!black] (279.5,137.5)--(292.5,137.5)--(289.5,141)--(282.5,141)--cycle;
						\shade[top color=black!60,bottom color=black!96,rounded corners=0.6pt] (322,115) rectangle (340,131.5);
						\fill[black!88] (328.5,110.5) rectangle (333.5,115.2);
						\draw[TXkinhdam!50!black,line width=0.6pt] (331,131.5) .. controls (333,155) and (322,172) .. (305,179);
					\end{scope}
					\draw[gray!60,line width=0.3pt,rounded corners=1.5pt] (5,5.5) rectangle (425,241);
				\end{scope}
			\end{tikzpicture}
		\end{center}
		\NeedspaceBang{9\baselineskip}%
		\choiceTFt
		{Độ tăng nhiệt độ của nước khi ra khỏi anode là $303$~K}
		{\True Công suất tiêu thụ điện năng của ống là $10^5$~W}
		{Nhiệt lượng mà anode truyền cho dòng nước trong mỗi giây là $100$~kJ}
		{\True Lưu lượng nước chảy vào làm mát anode là $0{,}79$ lít/giây (\textit{làm tròn và lấy hai chữ số có nghĩa})}
		\loigiai{
			\begin{itemchoice}
				\itemch \textbf{Sai.} $\Delta t=55-25=30\,^\circ\text{C}$, tức độ tăng nhiệt độ là $30$~K.
				\itemch \textbf{Đúng.} $P=U_\text{AK}I=10^4\cdot 10=10^5$~W.
				\itemch \textbf{Sai.} Mỗi giây nước nhận $Q=99\%\cdot P\cdot 1~\text{s}=0{,}99\cdot 10^5~\text{J}=99$~kJ.
				\itemch \textbf{Đúng.} Mỗi giây: $m=\dfrac{Q}{c\Delta t}=\dfrac{99\,000}{4200\cdot 30}\approx 0{,}786$~kg, ứng với $0{,}79$~lít.
			\end{itemchoice}
		}
	\end{ex}
	\Closesolutionfile{ans}

	\NeedspaceCau{7.4\baselineskip}
	\subsection*{\indam{PHẦN III. Câu trắc nghiệm trả lời ngắn. Thí sinh trả lời từ câu 1 đến câu 6.}}
	\setcounter{ex}{0}
	\Opensolutionfile{ans}[ans/VatLi2026-HaNoi-NguyenThiMinhKhai-Lan1-3]

	\begin{ex}%[VL12-C1-TH]
		Người ta cung cấp một nhiệt lượng $1{,}5$~J cho chất khí đựng trong một xilanh đặt nằm ngang. Khí nở ra đẩy pit-tông di chuyển đều một đoạn $5$~cm. Biết lực ma sát giữa pit-tông và xilanh có độ lớn $20$~N. Độ biến thiên nội năng của khí trong xilanh là bao nhiêu J?
		\par
		\shortans[oly]{0{,}5}
		\loigiai{
			Pit-tông chuyển động đều, khí thực hiện công $A=F_\text{ms}s=20\cdot 0{,}05=1$~J.\\
			$\Delta U=Q-A=1{,}5-1=0{,}5$~J.
		}
	\end{ex}

	\NeedspaceCau{14.9\baselineskip}
	\begin{ex}%[VL12-C3-VD]
		Một khung dây dẫn kín, phẳng có diện tích $20~\text{cm}^2$, được đặt trong từ trường có vectơ cảm ứng từ vuông góc với mặt phẳng khung dây. Biết độ lớn cảm ứng từ biến đổi theo thời gian như đồ thị hình bên. Biết cường độ dòng điện cảm ứng trong khung dây trong khoảng thời gian từ $0{,}4$~s đến $1$~s đo được là $0{,}2$~mA. Điện trở của khung dây bằng bao nhiêu $\Omega$?
		\begin{center}
			\begin{tikzpicture}[x=4.6cm,y=0.33cm,line join=round,line cap=round,font=\footnotesize]
				\draw[VLmui,thick] (0,0)--(1.1,0) node[above,inner sep=3pt] {$t$ (s)};
				\draw[VLmui,thick] (0,0)--(0,11) node[right,inner sep=2pt] {$B$ (mT)};
				\draw[dashed,thin,gray] (0,9)--(0.2,9) (0.2,0)--(0.2,9) (0.4,0)--(0.4,9);
				\draw[very thick,VLxanh] (0,3)--(0.2,9)--(0.4,9)--(1,0);
				\foreach \p in {(0,3),(0,9),(0.2,9),(0.4,9),(0.2,0),(0.4,0),(1,0),(0,0)}{\fill \p circle (1.5pt);}
				\node[left,inner sep=3pt] at (0,9) {$9$};
				\node[left,inner sep=3pt] at (0,3) {$3$};
				\node[below left,inner sep=2pt] at (0,0) {$0$};
				\node[below,inner sep=3pt] at (0.2,0) {$0{,}2$};
				\node[below,inner sep=3pt] at (0.4,0) {$0{,}4$};
				\node[below,inner sep=3pt] at (1,0) {$1$};
			\end{tikzpicture}
		\end{center}
		\shortans[oly]{0{,}15}
		\loigiai{
			Từ $0{,}4$~s đến $1$~s: $e=S\left|\dfrac{\Delta B}{\Delta t}\right|=20\cdot 10^{-4}\cdot\dfrac{9\cdot 10^{-3}}{0{,}6}=3\cdot 10^{-5}$~V.\\
			$R=\dfrac{e}{I}=\dfrac{3\cdot 10^{-5}}{0{,}2\cdot 10^{-3}}=0{,}15~\Omega$.
		}
	\end{ex}

	\NeedspaceCau{4.6\baselineskip}
	\begin{ex}%[VL12-C3-TH]
		Một đoạn dây dẫn thẳng dài $1$~m mang dòng điện $10$~A, đặt trong một từ trường đều có cảm ứng từ $0{,}1$~T thì chịu một lực $0{,}5\sqrt{2}$~N. Góc lệch giữa cảm ứng từ và chiều dòng điện trong dây dẫn là bao nhiêu độ?
		\par
		\shortans[oly]{45}
		\loigiai{
			$\sin\alpha=\dfrac{F}{BIl}=\dfrac{0{,}5\sqrt{2}}{0{,}1\cdot 10\cdot 1}=\dfrac{\sqrt{2}}{2}\Rightarrow\alpha=45^\circ$.
		}
	\end{ex}

	\NeedspaceCau{4.6\baselineskip}
	\begin{ex}%[VL11-C3-TH]
		Một điện tích điểm $Q=-2\cdot 10^{-7}$~C, đặt tại điểm A trong môi trường có hằng số điện môi $\varepsilon=2$, với $k=9\cdot 10^9$~N$\cdot$m$^2$/C$^2$. Cường độ điện trường $E$ do điện tích $Q$ gây ra tại điểm B với $\text{AB}=7{,}5$~cm có độ lớn bằng bao nhiêu (theo đơn vị là $10^5$~V/m)?
		\par
		\shortans[oly]{1{,}6}
		\loigiai{
			$E=\dfrac{k|Q|}{\varepsilon r^2}=\dfrac{9\cdot 10^9\cdot 2\cdot 10^{-7}}{2\cdot 0{,}075^2}=1{,}6\cdot 10^5$~V/m.
		}
	\end{ex}

	\NeedspaceCau{8.9\baselineskip}
	\begin{ex}%[VL12-C1-VD]
		Vào mùa hè, người Hà Nội thường có thói quen thưởng thức trà đá trong các quán vỉa hè. Để có một cốc trà đá chất lượng, người chủ quán rót khoảng $0{,}25$~kg trà nóng $80\,^\circ\text{C}$ vào cốc, sau đó cho tiếp $m$~kg nước đá $0\,^\circ\text{C}$. Cuối cùng được cốc trà đá ở nhiệt độ phù hợp nhất là $10\,^\circ\text{C}$ (hệ vừa đạt đến trạng thái cân bằng nhiệt). Biết phần nhiệt lượng mà hệ (nước và nước đá) nhận thêm của môi trường xung quanh bằng $10\%$ nhiệt lượng mà các cục nước đá nhận để làm tăng nội năng của chúng. Nhiệt dung riêng của nước là $4{,}20$~kJ/(kg.K); nhiệt nóng chảy riêng của nước đá là $3{,}33\cdot 10^5$~J/kg. Tính $m$ (\textit{làm tròn và lấy hai chữ số ở phần thập phân}).
		\par
		\shortans[oly]{0{,}22}
		\loigiai{
			Trà tỏa nhiệt: $Q_1=0{,}25\cdot 4200\cdot(80-10)=73\,500$~J.\\
			Nước đá thu nhiệt: $Q_2=m(\lambda+c\cdot 10)=m(3{,}33\cdot 10^5+42\,000)=375\,000m$.\\
			$Q_2=Q_1+0{,}1Q_2\Rightarrow 0{,}9\cdot 375\,000m=73\,500\Rightarrow m\approx 0{,}218$~kg.
		}
	\end{ex}

	\NeedspaceCau{13.4\baselineskip}
	\begin{ex}%[VL12-C2-VD]
		Hai bình có thể tích bằng nhau, chứa cùng một loại khí (coi là khí lí tưởng) có khối lượng $m_1$ và $m_2$. Đồ thị biến đổi áp suất theo nhiệt độ tuyệt đối của khí trong hai bình biểu diễn như hình bên. Tỉ số $\dfrac{m_1}{m_2}$ bằng bao nhiêu (\textit{kết quả làm tròn và viết đến hai chữ số thập phân})?
		\begin{center}
			\begin{tikzpicture}[line join=round,line cap=round,font=\footnotesize]
				\draw[VLmui,thick] (0,0)--(4.2,0) node[right,inner sep=2pt] {$T$};
				\draw[VLmui,thick] (0,0)--(0,3.2) node[left,inner sep=2pt] {$p$};
				\node[below left,inner sep=2pt] at (0,0) {$O$};
				\draw[very thick,VLxanh,VLmui] (0,0)--(45:3.6) node[right,inner sep=3pt,black] {$m_2$};
				\draw[very thick,VLdo,VLmui] (0,0)--(30:3.9) node[right,inner sep=3pt,black] {$m_1$};
				\draw[dash dot,thin] (0.95,0) arc[start angle=0,end angle=30,radius=0.95];
				\node at (13:1.32) {$30^\circ$};
				\draw[dashed,thin,<->,>=stealth] (1.9,0) arc[start angle=0,end angle=45,radius=1.9];
				\node at (22:2.3) {$45^\circ$};
			\end{tikzpicture}
		\end{center}
		\shortans[oly]{0{,}58}
		\loigiai{
			$pV=\dfrac{m}{M}RT\Rightarrow p=\dfrac{mR}{MV}T$: hệ số góc của đường biểu diễn tỉ lệ với $m$.\\
			$\dfrac{m_1}{m_2}=\dfrac{\tan 30^\circ}{\tan 45^\circ}\approx 0{,}58$.
		}
	\end{ex}

	\centerline{--- HẾT ---}
	\Closesolutionfile{ans}

	%===================== BẢNG ĐÁP ÁN =====================
	%\newpage
	%\begin{center}\indam{BẢNG ĐÁP ÁN}\end{center}
	%\noindent\indam{PHẦN I.}\\
	%\inputansbox{6}{ans/VatLi2026-HaNoi-NguyenThiMinhKhai-Lan1-1}
	%\noindent\indam{PHẦN II.}\\
	%\inputansbox{2}{ans/VatLi2026-HaNoi-NguyenThiMinhKhai-Lan1-2}
	%\noindent\indam{PHẦN III.}\\
	%\inputansbox{6}{ans/VatLi2026-HaNoi-NguyenThiMinhKhai-Lan1-3}
\end{document}
